% 本文件由 LaTeX 排版 Agent 根据有效 translation.json 自动生成。
% author=John Chrysostom; work_id=2401; title=若望福音讲道集

\chapter{\BibleBook{若望}{福音}讲道集}

% structure: p0001 source=h1 target=omitted reason=page-title-already-rendered-by-parent

\Emph{讲道1} 前言\newline{}\Emph{讲道2} \BibleRef{若 1:1}\newline{}\Emph{讲道3} \BibleRef{若 1:1}\newline{}\Emph{讲道4} \BibleRef{若 1:1}\newline{}\Emph{讲道5} \BibleRef{若 1:3}\newline{}\Emph{讲道6} \BibleRef{若 1:6-8}\newline{}\Emph{讲道7} \BibleRef{若 1:9}\newline{}\Emph{讲道8} \BibleRef{若 1:9-10}\newline{}\Emph{讲道9} \BibleRef{若 1:11}\newline{}\Emph{讲道10} \BibleRef{若 1:11-13}\newline{}\Emph{讲道11} \BibleRef{若 1:14}\newline{}\Emph{讲道12} \BibleRef{若 1:14}\newline{}\Emph{讲道13} \BibleRef{若 1:15}\newline{}\Emph{讲道14} \BibleRef{若 1:16-17}\newline{}\Emph{讲道15} \BibleRef{若 1:18}\newline{}\Emph{讲道16} \BibleRef{若 1:21-27}\newline{}\Emph{讲道17} \BibleRef{若 1:28-34}\newline{}\Emph{讲道18} \BibleRef{若 1:35-40}\newline{}\Emph{讲道19} \BibleRef{若 1:41-42}\newline{}\Emph{讲道20} \BibleRef{若 1:43-49}\newline{}\Emph{讲道21} \BibleRef{若 1:49-2:4}\newline{}\Emph{讲道22} \BibleRef{若 2:4-10}\newline{}\Emph{讲道23} \BibleRef{若 2:11-22}\newline{}\Emph{讲道24} \BibleRef{若 2:23-3:4}\newline{}\Emph{讲道25} \BibleRef{若 3:5}\newline{}\Emph{讲道26} \BibleRef{若 3:6-11}\newline{}\Emph{讲道27} \BibleRef{若 3:12-16}\newline{}\Emph{讲道28} \BibleRef{若 3:17-20}\newline{}\Emph{讲道29} \BibleRef{若 3:22-30}\newline{}\Emph{讲道30} \BibleRef{若 3:31-34}\newline{}\Emph{讲道31} \BibleRef{若 3:35-4:12}\newline{}\Emph{讲道32} \BibleRef{若 4:13-22}\newline{}\Emph{讲道33} \BibleRef{若 4:23-27}\newline{}\Emph{讲道34} \BibleRef{若 4:28-39}\newline{}\Emph{讲道35} \BibleRef{若 4:40-53}\newline{}\Emph{讲道36} \BibleRef{若 4:54-5:3}\newline{}\Emph{讲道37} \BibleRef{若 5:6-13}\newline{}\Emph{讲道38} \BibleRef{若 5:14-21}\newline{}\Emph{讲道39} \BibleRef{若 5:23-30}\newline{}\Emph{讲道40} \BibleRef{若 5:31-39}\newline{}\Emph{讲道41} \BibleRef{若 5:39-47}\newline{}\Emph{讲道42} \BibleRef{若 6:1-7}\newline{}\Emph{讲道43} \BibleRef{若 6:16-25}\newline{}\Emph{讲道44} \BibleRef{若 6:26-27}\newline{}\Emph{讲道45} \BibleRef{若 6:28-40}\newline{}\Emph{讲道46} \BibleRef{若 6:41-49}\newline{}\Emph{讲道47} \BibleRef{若 6:53-70}\newline{}\Emph{讲道48} \BibleRef{若 7:1-8}\newline{}\Emph{讲道49} \BibleRef{若 7:9-24}\newline{}\Emph{讲道50} \BibleRef{若 7:25-35}\newline{}\Emph{讲道51} \BibleRef{若 7:37-44}\newline{}\Emph{讲道52} \BibleRef{若 7:45-8:19}\newline{}\Emph{讲道53} \BibleRef{若 8:20-31}\newline{}\Emph{讲道54} \BibleRef{若 8:31-46}\newline{}\Emph{讲道55} \BibleRef{若 8:48-59}\newline{}\Emph{讲道56} \BibleRef{若 9:1-5}\newline{}\Emph{讲道57} \BibleRef{若 9:6-16}\newline{}\Emph{讲道58} \BibleRef{若 9:17-34}\newline{}\Emph{讲道59} \BibleRef{若 9:34-10:12}\newline{}\Emph{讲道60} \BibleRef{若 10:14-19}\newline{}\Emph{讲道61} \BibleRef{若 10:22-42}\newline{}\Emph{讲道62} \BibleRef{若 11:1-27}\newline{}\Emph{讲道63} \BibleRef{若 11:30-40}\newline{}\Emph{讲道64} \BibleRef{若 11:41-48}\newline{}\Emph{讲道65} \BibleRef{若 11:49-12:8}\newline{}\Emph{讲道66} \BibleRef{若 12:8-24}\newline{}\Emph{讲道67} \BibleRef{若 12:25-32}\newline{}\Emph{讲道68} \BibleRef{若 12:34-41}\newline{}\Emph{讲道69} \BibleRef{若 12:42-50}\newline{}\Emph{讲道70} \BibleRef{若 13:1-2}\newline{}\Emph{讲道71} \BibleRef{若 13:1-18}\newline{}\Emph{讲道72} \BibleRef{若 13:20-31}\newline{}\Emph{讲道73} \BibleRef{若 13:36-14:7}\newline{}\Emph{讲道74} \BibleRef{若 14:8-14}\newline{}\Emph{讲道75} \BibleRef{若 14:15-31}\newline{}\Emph{讲道76} \BibleRef{若 14:31-15:10}\newline{}\Emph{讲道77} \BibleRef{若 15:11-16:4}\newline{}\Emph{讲道78} \BibleRef{若 16:4-15}\newline{}\Emph{讲道79} \BibleRef{若 16:16-33}\newline{}\Emph{讲道80} \BibleRef{若 17:1-5}\newline{}\Emph{讲道81} \BibleRef{若 17:6-13}\newline{}\Emph{讲道82} \BibleRef{若 17:14-26}\newline{}\Emph{讲道83} \BibleRef{若 18:1-37}\newline{}\Emph{讲道84} \BibleRef{若 18:37-19:15}\newline{}\Emph{讲道85} \BibleRef{若 19:16-20:2}\newline{}\Emph{讲道86} \BibleRef{若 20:10-23}\newline{}\Emph{讲道87} \BibleRef{若 20:24-21:14}\newline{}\Emph{讲道88} \BibleRef{若 21:15-25}\par

\SourceRecord{Translated by Charles Marriott.；Nicene and Post-Nicene Fathers, First Series；Vol. 14.；Edited by Philip Schaff.；Buffalo, NY: Christian Literature Publishing Co.,；1889.；<http://www.newadvent.org/fathers/2401.htm>.}

% structure: p0001 source=h1 target=section reason=page-title

\section{讲道集第一 论若望福音（序言）}

1. 那些观看外教竞技的人，当他们听说有一位卓越的运动员和胜利者从某处来到时，便都跑去观看他的角力、技巧和力量；你可以在那成千上万的剧场里，看到他们全都聚精会神，目不转睛，唯恐错过任何细节。同样，这些人若听说有某位杰出的音乐家来到他们中间，便放下手中一切事务——往往是必需且紧迫的事务——登上阶梯，坐下来专注地聆听歌词与伴奏，并评论两者的和谐。这就是多数人的所作所为。\par

再者，那些精于修辞学的人对诡辩家也是如此；因为他们也有自己的剧场、听众、掌声、喧哗，以及对所说的话进行最严密的批评。\par

若在修辞学家、音乐家和运动员的情况下，人们坐着观看或聆听都如此热切，那么当现在出现的不是音乐家或辩论家来比试技艺，而是一个从天上发言、发出比雷霆更清晰声音的人时，你们按理应当表现出何等热忱与认真呢？因为他的声音响彻了大地，又充满和占据了它——不是靠呼喊的响亮，而是靠天主的恩宠推动他的舌头。\par

令人惊奇的是，这声音虽然宏大，却既不刺耳也不令人不快，反而比一切音乐的和谐更甜美、更悦耳，更善于抚慰人心；此外，它至圣至威，充满了如此伟大的奥迹，并带来了如此巨大的福祉，以至于如果人们准确而心甘情愿地接受并持守它们，他们便不再是凡人，也不停留在地上，而是超脱此生的一切，使自己适应天使的境况，在地上生活就如在天上一样。\par

2. 因为这位雷霆之子，基督的爱徒，普世万教会的柱石，掌握着天国的钥匙，饮了基督的杯，并以基督的圣洗受了洗，曾满怀信心地倚靠在主的胸膛上，现在来到我们面前；他不是上演一出戏剧，不戴面具藏头（因为他要说另一种言语），也不登上讲坛，不用脚踏响舞台，不穿金饰的衣服，而是穿着一种难以形容的美丽外袍。因为他将向我们显现，已\BibleQuote{穿上基督}\BibleRef{罗 13:14}，\BibleQuote{用和平的福音当作预备穿在脚上}\BibleRef{弗 6:15}；腰间束带，不是束在腰上，而是束在腰胯间，不是用染红的皮革制成，也不是外面涂金，而是用真理本身编织和构成的。现在他将出现在我们面前，不是演戏（因为他里面没有虚伪、虚构或神话），而是以坦然的头向我們宣告坦然的真理；他的举止、神情、声音都不使听众误以为他是另一个人，他传递信息不需要乐器，如竖琴、里拉琴或其他类似之物，因为他只用舌头成就一切，发出的声音比任何琴师或音乐都更甜美、更有益。整个天堂是他的舞台；他的剧场是可居住的世界；他的听众是所有天使，以及那些已经是天使或渴望成为天使的人，因为只有这些人才能正确聆听那和谐之音，并以行为将其展现出来；其余众人，如同幼童，虽也听到，却因沉迷于糖果和孩童的玩物而不知所闻；同样，他们在喜乐和奢华之中，只为财富、权力和感官享受而活，时而也听一些言论，但行为上却显不出任何伟大或高尚之处，因为他们永远把自己束缚在砖窑的泥土上。这位宗徒身旁，站着天上的掌权者，惊叹于他灵魂的美丽、他的理解力，以及那吸引基督亲自到他身边、使他获得圣神恩宠的德行之花。因为他已使自己的灵魂如同一把精心制作、镶嵌宝石、金弦的竖琴，并把它交托给圣神，让他藉以发出伟大而崇高的声音。\par

3. 既然现在弹奏这竖琴的不再是载伯德的儿子——那渔夫，而是那知晓\BibleQuote{天主深奥的事}\BibleRef{格前 2:10}的圣神，让我们如此聆听。因为他将不对我们讲凡俗之事，而是从圣神的深处，从那些在发生之前连天使也不知道的隐密事中发言；因为天使也藉着若望的声音与我们一同学习了我们所知道的事。另一位宗徒也曾宣告说：\BibleQuote{为要将天主各样的智慧，藉着教会，使天上的率领者和掌权者得知}\BibleRef{弗 3:10}。如果率领者、掌权者、革鲁宾和色辣芬藉着教会学习了这些事，那么显然他们极其热切地聆听这一教训；而我们也因此蒙受了不小的尊荣，因为天使与我们一同学习了他们先前不知道的事——我不说他们也藉着我们学习。所以，我们要表现出极大的静默和有序的行为；不仅今天如此，不仅在我们听讲的这一天，而是在我们整个一生，因为任何时候聆听他都是好的。如果我们渴望知道王宫里正在发生什么，比如国王说了什么、做了什么、正在对他的臣民筹划什么——虽然这些事大多与我们无关——那么，更应该渴望听到天主说了什么，尤其是当一切都关涉我们的时候。而所有这一切，这人将准确地告诉我们，因为他就是王的朋友，更确切地说，是他自己内在有王在说话，并从他那里听到他从父那里听到的一切。\BibleQuote{我称你们为朋友，因为我从我父听来的一切，都显示给你们了}\BibleRef{若 15:15}。\par

4. 因此，如同我们若看见有人从天上高处 \Quote{忽然} 俯身，应许准确描述那里的一切，就应一同奔跑；现在我们也要如此预备。因为这位从那里对我们说话；他不属于这世界，正如基督亲自宣称的：\BibleQuote{你们不属于世界} \BibleRef{若 15:19}，并且他内里有说话的护慰者、无所不在者、祂认识天主的事，如同人的灵魂认识自己的一切一样，就是圣洁之神、义德之神、引导之神，祂亲手引领人升天，赐给他们另一双眼睛，使他们能看见未来的事如同现在，并让他们在肉身中得见天上的事。因此，让我们一生在极大的宁静中顺服祂。不要让任何迟钝、昏睡、污秽的人进入并停留在这里；而要让我们自己升到天上，因为他在那里向那里的居民说这些话。如果我们停留在地上，我们将从那里得不到什么大的益处。因为若望的话对于那些不愿脱离这猪一般生活的人毫无意义，正如这世界的事对他也毫无意义。雷声震惊我们的灵魂，只有声音而无意义；但这人的声音并不困扰任何信徒，反而使他们摆脱困扰和混乱；它只使魔鬼和它们的奴仆感到惊骇。因此，为了知道如何使它们惊骇，让我们保持深沉的静默，无论是外在的还是内心的，尤其是后者；因为如果灵魂被搅动、充满不安，口舌沉默又有何益？我所寻求的是心灵的平静、灵魂的平静，因为我需要的是灵魂的听觉。因此，不要有财富的欲望搅扰我们，不要有虚荣的贪恋，不要有愤怒的暴虐，也不要还有其他情欲的群队；因为耳朵若不洁净，就不能恰当地领会所言之事的崇高；也不能正确地理解这些奥秘的可畏与不可言喻的本性，以及这些神圣神谕中的一切其他德性。如果一个人不能很好地学会笛子或竖琴的旋律，除非他竭尽全力集中注意力；那么，一个坐着聆听神秘声音的人，怎能以漫不经心的灵魂听见呢？\par

5. 因此基督亲自劝诫说：\BibleQuote{不要把圣物给狗，也不要把你们的珍珠投在猪前} \BibleRef{玛 7:6}。祂称这些话为\Quote{珍珠}，虽然实际上它们比珍珠更宝贵，因为我们没有比珍珠更宝贵的物质。为此，祂也常常将它们的甘甜比作蜂蜜，并非因为它们的甘甜仅此而已，而是因为在我们中间没有比蜂蜜更甜的东西。现在，为表明它们远远超乎宝石的本性和任何蜂蜜的甘甜，请听先知论及它们时宣告的这种超越：\BibleQuote{你的话比黄金和许多宝石更可羡慕；比蜂蜜和蜂巢更甘甜。} \BibleRef{咏 18:10} 但（只）对健康的人而言；因此他加上了：\BibleQuote{因为你的仆人遵守它们。} 并且在另一处称它们为甘甜时，他加上了\BibleQuote{于我的喉咙。} 因为他写道：\BibleQuote{你的言语对我何等甘甜，于我的喉咙。} \BibleRef{咏 118:103} 他又坚持这种超越，说：\BibleQuote{比蜂蜜和蜂巢更甘甜于我的口。} 因为他非常健康。我们也不应在我们患病时接近这些，而应在治愈我们的灵魂之后，才接受所供应的食物。\par

正因为如此，在如此长的开场白之后，我还没有尝试探究这些（圣若望的）言辞，为的是使每个人除去各种软弱，如同进入天堂本身一样，洁净地进入这里，摆脱愤怒、忧虑和今生的焦虑，以及所有其他情欲。因为除非一个人首先如此重新洁净他的灵魂，否则他不可能从这些话语中获得任何大的益处。但愿没有人说距离即将来临的共融时间很短，因为不仅可以在五天内，而且可以在瞬间改变整个生活方向。请告诉我，有什么比强盗和杀人犯更坏？这不是极端的邪恶吗？然而，这样的人一下子达到了卓越的顶峰，进入了乐园本身，不需要几天，也不需要半天，只需短短的一刻。因此，人可以突然改变，从泥土变成黄金。因为属于德行和属于恶习并非出于本性，改变是容易的，因为它不受任何必然性的束缚。\BibleQuote{如果你们愿意，并听从，你们将享用大地的美物。} \BibleRef{依 1:19} 你看见只需要意志吗？意志——不是众人普通的愿望——而是认真的意志。因为我知道现在所有人都希望飞上天堂；但必须用行为表明愿望。商人也希望致富；但他不让他的愿望停留在念头上；不，他装备一艘船，召集水手，聘请舵手，为船只装备所有其他补给，借钱，渡海，前往异国，忍受许多危险，以及所有航海者知道的其他事情。同样，我们也必须表明我们的意志；因为我们也航行一次旅程，不是从陆地到陆地，而是从地上到天上。因此，让我们如此安排我们的理性，使它有助于引导我们向上，也让我们的水手顺服它，并使我们的船只坚固，以免它被今生的逆境和沮丧淹没，或被虚荣的狂风抬起，而成为一艘快速平稳的船。如果我们这样安排我们的船，这样安排我们的舵手和船员，我们将顺风航行，并将吸引天主子——真正的舵手——降临我们，祂不会让我们的船被吞没，尽管万风可能吹起，祂会斥责风和海，代替汹涌的波浪，带来极大的宁静。\par

6. 故此，你们当先整顿自己，然后来参加我们下一次的聚会，至少如果你们还渴望听到有益的事，并将所言存于你们的灵魂中。但你们中不要有人成为\Quote{路旁}的，不要有人成为\Quote{石地}，也不要有人成为\Quote{荆棘丛生}的。\BibleRef{玛 13:4} 让我们自己成为沃土吧。因为当我们看见土地洁净时，我们（宣讲者）就会乐意播下种子；但如果土地是石头或粗糙的，请原谅我们不愿徒劳。因为如果我们停止播种而开始砍除荆棘，那么将种子撒在未耕作的土地上，无疑是极度的愚蠢。\par

享受这种听道好处的人，不该参与魔鬼的宴席。\Quote{因为义与不义有什么相通？}\BibleRef{格后 6:14} 你站着听若望，并藉着他学习圣神的事；而之后你却离开去听妓女说污秽的话，做更污秽的事，以及听那些互相打斗的柔媚之伶。如果你在如此泥沼中打滚，又如何能洗净自身？我何必详细列举那里所有的无耻行为？那里全是笑声，全是羞耻，全是耻辱，辱骂与嘲笑，全是纵欲，全是毁灭。看，我预先警告你们大家。凡享受此宴席福分的人，不要让那些有害的表演毁灭自己的灵魂。那里所有言与行都是撒殚的虚华。但你们这些已领受入门圣事的人，知道你们与我们立了什么盟约，或者更确切地说，当基督引导你们进入祂的奥迹时，你们向祂立了什么约，你们对祂说了什么话，关于撒殚的虚华，你们与祂谈了些什么；你们如何也为此与撒殚及其天使绝交，并承诺甚至不会朝那个方向看一眼。因此，并非轻微的理由害怕，一旦有人对这些承诺变得粗心，就会使自己不配领受这些奥迹。\par

7. 你们难道看不见，在国王的宫殿里，被召分享特殊恩宠并被列为国王朋友的不是那些得罪国王的人，而是那些备受尊敬、有特殊功勋的人吗？一位使者从天堂来到我们这里，由天主亲自派遣，与我们谈论某些必要的事情，而你们却离开去听祂的旨意和祂带给你们的信息，却坐着听戏子。这行径该受多大的雷霆、多大的天火啊！因为正如不宜参与魔鬼的宴席，同样也不宜听魔鬼的言论；也不宜以污秽的衣袍出现在那荣美的筵席前，那席上摆满了天主亲自预备的如此众多的美物。它的力量如此之大，只要我们以清醒的心神接近它，就能立刻提升我们到天堂。因为那不断受天主话语影响的人，不可能仍停留在此世卑微的状态，反而必须立刻展翅，飞向那上界之境，降落在无限福乐的宝藏上；愿我们众人全都能达到那境界，借着我们的主耶稣基督的恩宠和慈爱，愿光荣藉着祂并同着祂归于圣父和至圣圣神，从今世直到永远。阿们。\par

\SourceRecord{Translated by Charles Marriott.；Nicene and Post-Nicene Fathers, First Series；Vol. 14.；Edited by Philip Schaff.；Buffalo, NY: Christian Literature Publishing Co.,；1889.；<http://www.newadvent.org/fathers/240101.htm>.}

% structure: p0001 source=h1 target=section reason=page-title

\section{\WorkTitle{若望福音讲道第二篇}}

% structure: p0002 source=h2 target=subsection reason=relative-heading-rank; base=h2

\subsection{\BibleRef{若 1:1}}

假如若望要与我们交谈，并对我们说他自己之言，我们便必须描述他的家族、他的乡土和他的教育。然而，既然并不是他，而是天主藉着他向人类说话，我认为探询这些问题实属多余且令人分心。但即便如此，这并非多余，反而十分必要。因为当你明白他是谁，来自何处，他的父母是谁，他的品性如何，然后听到他的声音和他所有的天上智慧，你便会确知这些（道理）并不属于他，而是属于那感动他灵魂的天主之能。\par

那么，他来自哪个地方呢？不是来自任何地方，而是来自一个贫穷的村庄，来自一块不为人所重视、不出产任何好东西的土地。因为经师们诽谤加里肋亚说：\Quote{你们查考一下，就会知道：从加里肋亚不会出先知的。}\BibleRef{若 7:52} 而\Quote{真正的以色列人}也诽谤它，说：\Quote{从纳匝肋能出什么好事吗？} 他属于这个地方，甚至不属于其中任何著名之处，而是一个连名字都不为人知的村庄。他来自那里，他的父亲是个穷渔夫，穷得让他儿子们也从事同样职业。现在你们都知道，没有工匠会选择让儿子继承他的手艺，除非贫困所迫，尤其当这门手艺是低贱的时候。但没有什么比渔夫更贫穷、更低贱、更无知的了。然而，即使在渔夫中，也有大小之别；在那里，我们的宗徒也居于较低的地位——因为他不是从海中捕鱼，而是在一个小湖上度日。当他与父亲和兄弟雅各伯在湖边，修补他们破损的渔网——这本身就标志着极端的贫穷——基督就这样召叫了他。\par

至于世俗的学识，我们从这些事实可以得知他一无所知。此外，路加也见证这一点，他不仅记载他是无知的，而且完全是\Quote{没有学识的}\BibleRef{宗 4:13}。这是很可能的。因为一个如此贫穷的人，从不参加公众集会，也不与有声望的人来往，而是像钉在渔船上一样，即使偶尔遇见什么人，也只是与鱼贩和厨子交谈，我请问，他怎么可能比无理性的动物好到哪里去呢？他岂能不效法他所捕的鱼的哑默呢？\par

2. 那么，这位渔夫——他终日与湖泊、渔网和鱼打交道；他是加里肋亚贝特赛达人；是穷渔夫的儿子，而且穷到了极点；他无知，而且无知到了极点，在跟随基督之前和之后都从未学过文字——让我们看看他发表了什么言论，他跟我们谈论什么话题。是关于田里的事吗？是关于河里的事吗？是关于卖鱼的生意吗？因为人们或许期望从渔夫口中听到这些。但不要害怕；我们将听不到这些；我们会听到天上的事，是以前没有人学过的。因为，正如一个从圣神的宝库中发言的人所料想的那样，他给我们带来了崇高的教义、最佳的生活方式与智慧，仿佛刚从天上降临；甚至天上的许多事物也未必是人所知的，如前所述。这些事属于渔夫吗？告诉我吧。它们属于任何一个修辞学家吗？属于诡辩家或哲学家吗？属于每个受过外邦智慧训练的人吗？绝对不是。人的灵魂根本无法如此思维关于那纯粹而蒙福的本体、关于仅次于它的权能、关于不朽和永生、关于将来会不朽的必死身体的本性、关于惩罚和将来的审判、关于对行为言语及思念意念的追究。它无法说明什么是人，什么是世界；什么是真正的人，什么是貌似人却非人者；什么是德性的本质，什么是邪恶的本质。\par

3. 的确，柏拉图与毕达哥拉斯的门徒们探讨了其中一些事。至于其他哲学家的名字，我们无需提及；他们在这一点上全都极其可笑；而其中比其余人更受敬重、被视为这一学科带头人的那些人，更是如此；他们撰写了一些关于政体与学说的著作，但在所有方面都比孩童更为可耻地荒谬。因为他们在整个一生中都在倡导女子共有、颠覆生活秩序、废除婚姻的尊荣，并制定其他类似的可笑律法。至于灵魂的教义，他们没有留下任何可耻的话；他们断言人的灵魂变成了苍蝇、蚊虫和灌木，甚至天主自己就是一个灵魂，以及其他类似的猥亵之词。\par

他们不仅在这方面应受指责，而且他们那永远变换不定的言辞也应受指责；因为他们根据不确定和谬误的论据主张一切，就像在\OriginalTerm{尤里普斯海峡}{Euripus}中载浮载沉的人，永不能停留在同一位置。\par

这位渔夫却不然：他所言一切皆无误；他如同站在磐石上，绝不改变立场。因为他已被认为配得进入至圣之处，且万有的主在他内发言，故他不受任何人事的影响。至于那些人，则像连做梦都不配进入王宫、却与其他人一起在广场上度日的人，凭自己的想象猜测他们看不见的事物，便犯了大错，如同瞎眼或醉酒的人，彼此冲突；不仅彼此冲突，也自相矛盾，不断改变观点，且总是针对同样的事。\par

4. 但这未受过教育的、无知的、\Place{贝特赛达}人、\Person{载伯德}的儿子，（虽然希腊人千百次嘲笑这些名字的粗鄙，我仍将更大胆地述说它们。）因为他的民族在他们看来越是蛮族，他越显得远离希腊的训导，我们所拥有的也就显得越加光明。当一个蛮族和未受过教育的人说出地上无人知晓的事物，并且不只是说出，（如果仅此而已，这已是伟大奇迹，）而且还提供另一个更强的证据，证明他所说的是受天主启示的，即说服所有时代的一切听众；谁不惊叹那住在他内的德能呢？因为正如我所说，这是最强有力的证据，证明他并未制定自己的法律。那么，这位蛮族，凭借他的福音著作，占据了整个可居住的世界。他的身体占据了亚细亚的中心，那里是往昔所有希腊派别哲学化的地方，在敌人中间照耀，驱散他们的黑暗，摧毁鬼魔的堡垒；但他的灵魂已退隐到适合如此行事之人的地方。\par

5. 至于希腊人的著作，它们都已消失沉寂，但此人的著作却日益闪耀。因为自从他（及其余渔夫）以来，\Person{毕达哥拉斯}和\Person{柏拉图}的学说，它们先前似乎盛行，如今已不再被提及，大多数人甚至不知道它们的名字。然而，\Person{柏拉图}，他们说是国王的受邀同伴，有许多朋友，曾航行到\Place{西西里}。而\Person{毕达哥拉斯}占据了\Place{大希腊}，并在那里施行了成千上万种巫术。因为与牛交谈，（他们说他曾如此行，）不过是巫术的一种。这一点最为清楚：那与野兽交谈的人并未在任何人身上造福人类，反而对他们造成最大的伤害。然而，人的本性确实更适合哲学的推理；但他仍如他们所说，与鹰和牛交谈，使用巫术。因为他并未使它们的非理性本性变成理性，（这对人是不可能的，）而是通过他的魔术伎俩欺骗愚人。他忽略教导人类任何有益的事，却教导他们可以吃生养他们之人的头颅，如同吃豆子一样。他还说服那些与他交往的人，他们老师的灵魂实际上曾是一丛灌木，又曾是一个少女，又曾是一条鱼。\par

这些事岂非理所当然地消灭了，彻底消失了？理所当然，也合情合理。但这位无知和未受教育的人的话却不然；因为叙利亚人、埃及人、印度人、波斯人、埃塞俄比亚人，以及其他成千上万的民族，将他们自己的语言翻译成他所引入的教义，虽然是蛮族，却学会了哲学。我并非无端地说整个世界成了他的剧场。因为他没有离开自己的同类，将劳力浪费在无理性的受造物上（这是极度虚荣和极端愚蠢的行为），而是同样免于这种情欲和其他情欲，只专注于一点：使整个世界都能学到一些有益的事物，并能将它从地上转移到天上。\par

为此缘故，他也未将自己的教导隐藏在迷雾和黑暗中，像那些人一样，用言语的晦涩（如同面纱）遮掩其内藏的祸害。但此人的教义比阳光更清晰，因此它们已向全世界的所有人展开。因为他并未像\Person{毕达哥拉斯}那样教导，命令前来的人沉默五年，或像无感觉的石头一样坐着；他也未编造神话，将宇宙定义为由数字构成；而是摒弃所有这些魔鬼般的垃圾和祸害，使他的话语如此简朴，以致他所说的一切不仅对智者，而且对妇女和少年都是明白的。因为他确信这些话是真实的，对所有聆听它们的人有益。他之后的所有时代都是他的见证；因为他吸引了整个世界，当我们聆听这些话时，他使我们的生命脱离了所有怪异的智慧炫耀；因此，我们这些聆听的人宁愿舍弃生命，也不愿舍弃他传给我们的教义。\par

6. 由此，并从其他一切情况来看，显而易见，此人的一切都不是人的，而是神圣和天上的，教训通过这个神圣的灵魂临到我们。因为我们不会注意到动听的句子，也不会有华丽的措辞，也不会过分和无用的词句顺序和安排，（这些远离一切真智慧，）而是无敌的神圣力量，以及正确教义不可抗拒的力量，以及无数好事丰富的供应。因为他们对表达方式的过度关心是如此过分，如此只配得上诡辩家，甚至不是诡辩家，而是愚蠢的少年，以致他们自己的首席哲学家介绍他的导师为这门艺术深感羞愧，并对审判官说，他们从他那里听到的将是直白的、未经事先准备的，没有修辞的修饰，也没有（华丽的）句子和词语的装饰；因为，他说，诸位，像我这般年纪的人，像少年一样来到你们面前编造演说，这肯定是不合适的。而观察这件事的极端荒谬性：他所描述的他的导师避免作为可耻、不配哲学和少年工作的事，正是他本人最着力培育的。他们完全沉溺于对名望的单纯爱好。\par

就如你揭开那些外面粉刷的坟墓，会发现里面充满腐败、恶臭和朽骨；同样，哲学家的学说，你若剥去其华丽辞藻，便会看到充满许多可憎之事，尤其是当他论及灵魂时，既无度地尊崇又无度地诋毁。这正是魔鬼的陷阱：从不保持适当比例，而是通过两端的过度，将陷入其中的人引入恶言。他一时说灵魂是天主的本质，一时又在如此无度而不敬地高举灵魂之后，又以另一种方式过分，侮辱它，使它变为猪、驴，以及比这些更卑微的动物。\par

但就此打住；甚至这也已过分。因为若可能从这些事学到任何有益之物，我们本应多费些时间；但若只是观察它们的不雅和荒谬，我们已说了多余的话。因此我们将放下他们的寓言，专注于我们自己的教义，这些教义是从高处由这位渔夫的唇舌传授给我们的，毫无属人的成分。\par

7. 那么让我们呈现这些话，现在我已提醒你们，正如我一开始就劝勉你们，要热切注意所说的。那么这位福音作者在他开始时立刻说了什么呢？\par

\BibleQuote{在起初已有圣言，圣言与天主同在。} \BibleRef{若 1:1} 你们看到这话语是何等的大胆和权能吗？他说话毫不犹疑或猜测，而是清楚地宣告一切？因为这是教师的本分，在他所说的一切中毫不摇摆，因为若一位作他人向导的人还需要另一个人能坚定地建立他，那他就应该被列在门徒之中，而非教师。\par

但若有人说：\Quote{他忽略了第一因，而立刻向我们谈及第二因，是什么原因？}我们将避免谈论\Quote{第一}和\Quote{第二}，因为天主性超越数目和时间的先后。因此我们避免这些说法，但我们宣认圣父非由谁而生，圣子是由圣父所生。是的，也许有人说，但为什么他撇下圣父而谈论圣子呢？为什么？因为前者对所有人是显明的，即使不以为父，至少以为天主；但\OriginalTerm{独生子}{Only-Begotten}却不为人知；因此合理的是，他从一开始就急忙将认识祂的知识植入那些不认识祂的人心中。\par

此外，他在这些点上关于圣父的论述并非沉默。我恳请你们留意他的属灵智慧。他知道人们最尊崇那在万有之先的至古老者，并以此为天主。因此他首先从这一点开始，随着进展，宣告天主是存在的，而不像\Person{柏拉图}那样，有时断言祂是理智，有时是灵魂；因为这些事物远离那神圣而纯粹的本性，这本性与我们毫无共同之处，就实体而言，与受造物毫无关联，尽管就关系而言并非如此。\par

因此他称祂为\Quote{圣言}。因为他即将教导这\Quote{圣言}是天主的独生子，为了使无人以为祂的诞生是易受苦难的，通过给予祂\Quote{圣言}的称号，他预先防范并消除了邪恶的猜疑，表明圣子来自圣父，且祂没有受苦（改变）。\par

8. 那么你看到正如我说的，他在论及圣子时并非对圣父沉默吗？如果这些实例不足以完全解释整件事，不要惊奇，因为我们的论题是天主，祂是不可描述的，也无法相称地想象；因此这人从未指定祂本质的名字（因为不可能说出天主是什么，就本质而言），而是处处通过祂的作为向我们宣告祂。因为这\Quote{圣言}很快就被称为\Quote{光}，而\Quote{光}又被称为\Quote{生命}。\par

虽然并非仅为此缘故他如此称呼祂；这是第一个原因，第二个是因为祂即将向我们宣告父的事。因为\Quote{一切}，祂说，\BibleQuote{凡我从父所听见的，我都已告诉你们了。} \BibleRef{若 15:15} 祂称祂为\Quote{光}和\Quote{生命}，因为祂将来自知识的光和随之而来的生命白白赐给了我们。总之，一个名字不够，两个、三个乃至更多也不足以教导我们属于天主的事。但我们必须满足于能够通过许多事物来领会，尽管模糊不清，祂的属性。\par

他并非简单地称祂为\Quote{圣言}，而是加了冠词，以这种方式也将祂与其余者区分开来。那么你看到我说这位福音作者从天向我们说话并非无因吗？只看从一开始，他将灵魂提升到哪里，给了它翅膀，并带走了听者的心智。因为将它置于高于一切感官之物、高于地、高于海、高于天之后，他牵着它越过天使、革鲁宾和色辣芬、上座者、宰制者和掌权者；总之，说服它超越一切受造之物旅行。\par

9. 那么怎样？当他把我们带到如此高度时，他真的能让我们停在那里吗？绝不可能；正如人把站在海滩上、眺望城市、海滩和海港的人带到海中央，确实使他离开了先前的对象，却不使他的视线停留在任何地方，而是带他到无边的视野；同样，这位福音作者，将我们带到所有受造物之上，并护送我们朝向超越这些的永恒世代，使视线悬置，不允许它达到任何向上的界限，因为确实没有界限。\par

因为理智上升到\Quote{元始}，追问什么是\Quote{元始}；然后发现那\Quote{存在}总在超越其想象，没有可以停止思想的点；但专注地向前看，不能在任何点停下，它就疲倦了，转回下面的事物。因为这\Quote{在元始已有}无非是表达永远存在和无限存在。\par

你看到真正的哲学和神圣的道理了吗？不像那些希腊人，他们指定时间，说有些神更年轻，有些更年长。在我们这里没有这样的事。因为如果天主是存在的，正如祂确实存在，那么在祂之前没有东西。如果祂是万物的创造者，祂必是首先的；如果祂是万物的主宰和主，那么万物，包括受造物和世代，都在祂之后。\par

10. 我本想再论及其他疑难，但也许我们的心智已经疲倦了；因此，当我已就这些对聆听有益的事，无论对已说的还是将要说的，向你们提出劝告后，我将再次静默。那么，这些点是什么呢？我知道许多人因所讲的话太长而感到困惑。当灵魂被今生的许多重担压得沉重时，便会发生这种情况。因为眼睛清澈透明时，视力也敏锐，甚至辨别极微小的物体也不易疲倦；但是，当某种不良的体液从头部注入眼中，或者像烟雾一样的气体从下面升到眼中，在眼球前形成一种浓密的云时，它甚至连较大的物体也无法清晰感知；灵魂的情况也是如此。因为当灵魂被净化，没有扰乱它的情欲时，它就坚定地注视其应关注的对象；但是，当灵魂被许多情欲所蒙蔽，失去其本有的优越时，它就难以胜任任何崇高之事，很快就疲倦，退却，转向睡眠和怠惰，放过了那些关乎美德和由此而来的生命之事，而不是以极大的热忱去接受。\par

为使你们不受此苦（我将不住地这样提醒你们），要坚固你们的心智，免得你们听到希伯来人中的信者从保禄那里听到的话。因为他曾对他们说，他有\Quote{许多话要说，且难以陈说} \BibleRef{希 5:11}；并非因为这些话本身如此，而是因为，他说，\Quote{你们听觉迟钝。}因为软弱有病之人的本性就是被少数话甚至许多话弄糊涂，把清楚容易的事看作难以理解。但愿这里没有人是这种人，而是赶走一切世俗挂虑，如此来听这些道理。\par

因为当贪财的欲望占据听者时，听道的欲望就不能同样占据他；因为灵魂是单一的，不能同时满足多种欲望；但两者中的一个被另一个所伤，并且由于分割，随着其竞争者得胜而变得更弱，将全部精力消耗于自身。\par

这通常发生在子女身上。当一个人只有一个孩子时，他极其爱那一个。但当他成为许多孩子的父亲时，他的感情分配变得分散而减弱。\par

如果这在有自然绝对规则和力量、且所爱之物彼此相通的情况下发生，那么对于按照自由选择而来的欲望和意向，我们能说什么呢？特别是当这些欲望彼此直接对立时；因为爱财富与爱这种听道是彼此对立的。当我们进入这里时，我们进入天堂——我指的是，不是地方，而是意向；因为一个在地上的人有可能站到天堂里，看见那里的事物，并听见从那里来的话。\par

11. 因此，谁也不要把地上的事带进天堂；这里站着的人不要挂虑自己家里的事。因为他应当随身携带，并在家中和事务中保存他从这里所获的益处，不让它被家务和市场的重担压垮。我们进入教诲席位的理由，是为从那里洗净外界的污秽；但如果我们在这一点点时间里还可能因外界的言行而受伤，那我们根本不要进入更好。因此，在集会中，谁也不要想家务事，而要在家里活跃于在集会中所听的事。让这些事比任何事都更宝贵。这些事涉及灵魂，那些事涉及身体；或者更确切地说，这里所说的事涉及灵魂和身体二者。因此，让这些事成为我们的主要事务，其他一切只是偶尔的操劳；因为这些事既属于将来的生命，也属于现在的生命，但其余的事则不属于二者，除非它们按这些事的规范管理。因为从这些事不可能只学到我们将来的状况和如何生存，也能学到如何正确引导现世的生活。\par

因为这屋宇是一间属灵的诊所，使我们在外所受的任何创伤，都能在这里得到医治，而不是让我们收集新的创伤带出去。然而，如果我们不留意对我们发言的圣神，我们不仅不能清除从前的创伤，反而会增添新的。\par

因此，让我们以极大的热忱专注于那正为我们展开的书卷；因为如果我们准确学会其中的首要原理和基础道理，此后就不需要太多辛苦研究，而是在开头稍加努力后，便如保禄所说，也能教导别人。\BibleRef{罗 15:14}因为这位宗徒非常崇高，充满许多道理，他在这些事上比其他事更着力。\par

那么，我们不可做粗心的听众。这正是我们循序渐进地讲解给你们的原因，好使一切都容易理解，又不致从你们的记忆中溜走。因此，我们当害怕，免得我们落在那句话的谴责之下，那句话是说：\BibleQuote{若我没有来向他们说话，他们就没有罪。}\BibleRef{若 15:22} 因为如果我们听了之后，回家时什么也不带，只是对所讲的话感到惊奇，那么比起那些没有听过的人，我们又能多得什么益处呢？\par

那就让我们撒在好地里吧；让我们这样做，好使你们更加吸引我们。若有人有荆棘，让他用圣神的火焚烧它们。若有人有坚硬顽固的心，让他用同一把火使它柔软而顺从。若有人像路旁一样被各种思想践踏，让他进入更隐蔽的地方，不要暴露在那些愿意入侵掠夺的人面前：这样我们就能看见你们的庄稼波浪起伏。此外，如果我们对自己如此用心，勤奋地从事这属灵的听道，即使不能立刻，也会逐渐地，我们必定会从一切生活的挂虑中获得释放。\par

所以我们要小心，免得有人论我们说，我们的耳朵如同聋蝮。\BibleRef{咏 57:4} 请告诉我，这样的听众与野兽有什么不同？当天主说话时，那人不留心听，岂不是比任何无理性的动物还要无理吗？如果取悦天主真是做人的本分，那么连如何在这事上成功都不愿意听的人，除了野兽还能是什么？那么想一想，当基督愿意将人类提升到与天使同等时，我们却自愿从人的地位堕落到野兽的地位，这是何等的不幸！因为服侍肚腹、被财富的欲望所占有、动怒、咬人、踢人，这些都不是人的行为，而是野兽的行径。不，甚至野兽，我们可以说，每一种只有一个单一的欲望，而且是出于自然的。但人，当他抛弃了理性的统治，从天主所建立的社会中撕裂自己，便委身于所有的欲望，不再仅仅是野兽，而是一种多形杂乱的怪物；他甚至连自然的借口都没有，因为他的一切邪恶都源于蓄意的选择和决心。\par

愿我们永远不必对基督的教会有所怀疑。事实上，我们对于你们所确信的是更美好的事，是属于救恩的事；但是我们越这样确信，就越要谨慎，不可停止告诫的话。好使我们登上德行的顶峰，获得所应许的福分。愿我们众人都能达到这一境界，借着我们的主耶稣基督的恩宠和慈爱，愿光荣借着祂并同着祂归于圣父及圣神，永世无尽。阿们。\par

\SourceRecord{Translated by Charles Marriott.；Nicene and Post-Nicene Fathers, First Series；Vol. 14.；Edited by Philip Schaff.；Buffalo, NY: Christian Literature Publishing Co.,；1889.；<http://www.newadvent.org/fathers/240102.htm>.}

% structure: p0001 source=h1 target=section reason=page-title

\section{\WorkTitle{\BibleBook{若望}{福音}第三篇讲道}}

% structure: p0002 source=h2 target=subsection reason=relative-heading-rank; base=h2

\subsection{\BibleRef{若 1:1}}

关于聆听时注意的事，再劝勉你们已是多余，因你们已如此迅速地以行动表明了我劝告的成果。你们争先恐后地聚集、专注的姿态、急于占据内圈位置以便更清楚听见我的声音、直到这属神聚会结束才肯离开人群、鼓掌与喝彩——总之，这一切都可视为你们灵魂热忱和你们渴望聆听的明证。因此，在这点上再劝勉你们是多余的。然而，有一件事我们必须命令和恳求：你们要继续保持同样的热心，不仅在这里表现出来，而且回到家中，丈夫与妻子、父亲与儿子，也要谈论这些事。自己也要讲说一些，并要他们回报；这样，众人一起为这美好的盛筵作出贡献。\par

谁也不应对我说，\Quote{我们的孩子不应从事这些事}；他们不仅应当从事，而且应当唯独对这些事热心。虽然因你们的软弱，我不坚持这一点，也不将孩子们从世俗学问中抽离，正如我也不将你们从世俗事务中拉出来一样；然而，在这七日中，我要求你们奉献一日给我们众人共同的主。因为，我们吩咐仆役终身为我们劳役，而自己却不愿划出一点闲暇给天主；这岂不奇怪？况且，我们的一切事奉并不增添天主什么（因为天主性一无所缺），反而为我们自己带来益处。然而，当你们带孩子们去剧院时，既不提他们的数学课，也不提任何类似的事；但如果要求获得或收集任何属神的事物，你们就称之为\Quote{浪费时间}。如果你们为其他一切事都找得出时间和时节，却认为让孩子们从事与天主相关的事是麻烦和不合时宜，你们怎能不惹天主发怒呢？\par

不要这样，弟兄们，不要这样。正是这个年龄最需要聆听这些事；因为因其柔嫩，容易记住所说的话；孩子们所听到的，如同印记印在蜡块上一样印在心灵上。此外，正是在那时，他们的生命开始倾向邪恶或德行；如果从一开始就引领他们远离不义，并引导他们走上最好的道路，他就会使他们将来养成一种习惯和本性，即使他们愿意，也不容易变坏，因为习惯的力量吸引他们行善。这样，我们将看到他们比老年人更值得尊敬，在世俗事务中也更有用，在年轻时便显出老年人的品质。\par

因为，如我先前所说，享受聆听这些事、并陪伴如此一位宗徒的人，无论是男、是女、是青年，凡分享此筵席的，不可能不获得某种伟大而显著的益处。如果我们用言语训导我们所拥有的动物，从而驯服它们，那么借着这属神的教导，我们岂不更能对人产生这样的效果？何况在这两种情况下，治疗方法和被治疗的对象都有很大差异。因为我们的凶猛不如野兽，野兽的凶猛出于本性，我们的出于选择；而且言语的能力也不同，前者的能力是人类理智的能力，后者的能力是圣神的能力与恩宠。那么，那对自己绝望的人，让他看看驯服的动物，他就不再如此受影响；让他不断来到这治疗之所，让他时刻聆听圣神的律法，回家后将所听见的写在心中；这样，当他从经验中感受到进步时，他的希望将是美好的，信心也是极大的。因为当魔鬼看见天主的法律写在灵魂里，心变成了书写它的版，他就不再靠近。因为凡有君王文书之处——不是刻在铜柱上，而是由圣神印在爱主的心灵上，并因丰盛的恩宠而闪耀——那恶者甚至不能直视它，而要远远地转身离弃我们。因为没有什么比一颗关心天主之事的心灵、一个永远沉浸在这泉水中的灵魂，更令他（及其所暗示的思想）感到恐惧。这样的人，即使不顺心的事也不能搅扰他，即使顺心的事也不能使他自满或骄傲；相反，在这一切风波和骇浪中，他反会享有极大的平静。\par

因为混乱出自我们内在，并非源于境遇的本质，而是出于我们心灵的软弱；因为如果我们是因所遭遇之事而如此受影响，那么（既然我们都航行在同一片海上，不可能避免波浪和水花）所有人都必然受苦；但如果有人能置身于风暴和怒涛的影响之外，那么显然，制造风暴的不是环境，而是我们自己心灵的状态。因此，如果我们调整心灵，使之能甘愿承受万事，我们就不会有风暴，甚至连涟漪也没有，而是始终一片平静。\par

在我声称不再谈这些事之后，不知何故，我竟陷入了如此冗长的劝勉。请原谅我的冗长；因为我害怕，是的，我非常害怕，唯恐我们的这份热心会变得冷淡。如果我对它充满信心，我现在就不会对你们说这些事了，因为那已足以使你们万事轻松。但现在是时候在接下来的部分转向今天准备思考的议题了；免得你们带着疲惫进入争战。因为我们有争战，对抗真理的仇敌，对抗那些用尽一切伎俩摧毁天主圣子（或不如说他们自己的荣耀）的人。天主圣子的荣耀永远不变，毫不会因亵渎的舌头而减少；但那些人，由于急切地想要贬低他们所声称敬拜的那一位，反使自己满面羞惭，灵魂充满惩罚。\par

那么，当我们坚持我们已坚持的主张时，他们又说什么呢？\Quote{\InnerQuote{在起初已有圣言}这句话并不绝对表示永恒，因为同样的表述也用于天地。}这是何等的厚颜无耻和不敬！我在对你谈论天主，而你却将地和属于地的人扯进来？照此推论，既然基督被称为天主子，而被称为天主子的世人也是天主了？因为，\BibleQuote{我曾说：你们是神，你们都是至高者的子民。}\BibleRef{咏 82:6}你要与独生子争辩子嗣身份，并断言祂在这方面并不比你享有更多吗？回答是：\Quote{绝无可能。}但你却这样做了，即使你口头上没这么说。\Quote{如何？}因为你声称你藉恩宠分沾了义子名分，而祂也同样。因为你说祂不是本性的儿子，你只是使祂成为恩宠的儿子。\par

然而，让我们看看他们向我们提出的证据。据说，\Quote{在起初，} \BibleQuote{在起初天主创造了天地，大地还是混沌空虚}\BibleRef{创 1:2}并且，\BibleQuote{有一位是在辣玛的一个祖斐人}\BibleRef{撒上 1:1}这就是他们认为有力的论证，而且它们的确有力；但这是为了证明我们所主张的道理的正确性，而它们完全无力确立他们的亵渎。因为告诉我，\Quote{是}一词与\Quote{创造}一词有什么共同之处？天主与人有什么共同之处？你为什么把不能混合的东西混合？为什么混淆不同的事物？为什么把上面的贬低？在那里，表示永恒的不仅是\Quote{是}这个表达，而是那一位\Quote{在起初是}。而另一个，\Quote{圣言是}；因为正如\Quote{存在}一词，用于人时只区分现在时，但用于天主时则指永恒；同样，\Quote{是}用于我们的本性时，向我们表示过去的时间，而且是有限的，但用于天主时则宣告永恒。那么，当人听到\Quote{地}和\Quote{人}这些词时，本应足够想象关于它们的任何更多东西，只要恰当地认为它们是被造的本性；因为那成为存在的，无论是什么，要么是在时间中成为的，要么是在时间之前的世代中，但天主子不仅超越了时间，而且超越了所有先前的世代，因为祂是它们的创造者和制造者，如宗徒所说：\BibleQuote{藉着祂也创造了世代。}现在，制造者必然先于被造之物。然而，既然有些人如此愚昧，甚至在这之后还对受造物抱有超出其本分的更高观念，于是藉着\Quote{祂创造了}这个表达，以及另一个\Quote{有一个人}，他预先抓住了听众的思想，并从根部切除了所有无耻。因为所有被造的，无论是天还是地，都是在时间中被造的，并且有时间的开始，没有一个是无始的，因为是被造的；所以当你听到\Quote{祂创造了地}和\Quote{有一个人}时，你是在无谓地浪费时间，编织一堆无用的愚蠢。\par

因为我甚至可以再提一件事来更进一步。是什么呢？那就是，如果曾对地说过\Quote{在起初是地}，对人也说过\Quote{在起初是人}，那么即使那样，我们也绝不可能想象关于它们比我们现在所确定的更伟大的事物。因为\Quote{地}和\Quote{人}这些术语，由于它们被预先假定，无论关于它们说了什么，都不容许心灵想象任何比我们现在所知道的更伟大的事物。正如\Quote{圣言}，即使关于它说得很少，也不容许我们（关于它）思考任何低下或贫乏的东西。因为在进展中，他论到地说：\Quote{地是看不见且不成形的。}因为已经说了\Quote{祂造了}它，并确定了它的适当界限，他之后便毫无畏惧地宣告后续内容，因为他知道没有一个人会愚蠢到认为它是无始的、非受造的，因为\Quote{地}这个词，以及另一个\Quote{被造}，足以使即使非常单纯的人也相信它不是永恒的、不是非受造的，而是在时间中被造的受造物之一。\par

3. 此外，\Quote{是}这个表达，应用于地和人的时候，并不表示绝对的存在。但在人的情况下，（它表示）他处于某地的状态；在地的情况下，它表示其处于某种状态。因为他没有绝对地说\Quote{地是}然后就沉默，而是教导了地受造之后的状态，即它是\Quote{看不见且不成形的}，仍被水覆盖而混沌。同样，关于厄耳卡纳，他不仅说\Quote{有一个人}，还补充了他来自何处，\Quote{在辣玛的一个祖斐人。}但在\Quote{圣言}的情况下，则不然。我羞于将这些情况互相比较，因为如果我们责备那些在人的情况下这样做的人，而实际上那些被这样比较的人在美德上差异巨大，尽管他们的实体的确相同；那么在本性和其他一切事物的差异都如此无限的情况下，提出这样的问题岂不是极端的疯狂？但愿那被他们亵渎的主怜悯我们。因为不是我们发明了这些讨论的必要性，而是那些与自己的救恩作战的人把它强加给我们。\par

那么我说什么呢？这第一个\Quote{是}，应用于\Quote{圣言}，只表示祂永恒的存在（因为他说：\Quote{在起初，圣言就\OriginalTerm{是}{was}}），而第二个\Quote{是}（\Quote{圣言\OriginalTerm{是与}{was with}天主同在}）则表示祂的相对存在。既然永恒且无始是天主最特有的属性，祂首先提出这一点；然后，以免有人听见祂\Quote{在起初}，就断言祂也是\OriginalTerm{无生}{unbegotten}，他立即在说明祂是什么之前补充说，祂\Quote{\OriginalTerm{与天主同在}{with God}}。他藉着加定冠词（正如我前面所说）以及这第二个表达，防止任何人以为这位\Quote{圣言}仅仅是口里说出或心里想出的那样的话。因为他不是说\Quote{在天主内}，而是说\Quote{与天主同在}，向我们宣告祂位格方面的永恒。然后，随着论述的推进，他更清楚地揭示了这一点，又加上说，这位\Quote{圣言}也\Quote{是天主}。\par

\Quote{可是，是受造的，}有人会说。那么，有什么阻止他说：\Quote{在起初，天主创造了圣言}呢？至少梅瑟谈及大地时，不是说\Quote{在起初，大地是，}而是说\Quote{祂创造了它，}然后它就有了。那么，有什么阻止若望同样说：\Quote{在起初，天主创造了圣言}呢？因为如果梅瑟担心有人会断言大地是无生的，那么若望对于圣子，如果祂真是受造的，就更应该担心了。世界是可见的，正因如此，它就宣告了它的创造者（圣咏作者说：\Quote{诸天\BibleQuote{述说天主的荣耀}}\BibleRef{咏 18:1}），但圣子是看不见的，且远远超越所有受造之物。如果在一方面，我们无需论证或教导就知道世界是受造的，但先知却清晰且优先地陈述了这一事实；那么对于圣子，如果祂真是受造的，若望就更应该宣布同样的事了。\par

\Quote{是的，}有人会说，\Quote{但伯多禄已清楚明白地肯定了这一点。}在哪里，什么时候？\Quote{当他对犹太人说，\InnerQuote{天主已立祂为主、为基督了。}}\BibleRef{宗 2:36}你为什么不加上后面的话：\Quote{这位耶稣，就是你们钉在十字架上的}？难道你不知道，这些话一部分涉及祂无混合的本性，一部分涉及祂的降生成人吗？但如果不是这样，而你执意要将一切都理解为指天主性，那么你就使天主性成了可受苦的；但若不可受苦，则不是受造的。因为如果血是从那神圣而无可言喻的本性流出的，并且那本性（而非肉体）在十字架上被钉子撕裂割伤，那么在这种假设下你的诡辩才有道理；但如果连魔鬼自己也不能说出这样的亵渎话，你为什么假装无知，无知得如此不可原谅，甚至恶神们也无法装作这样呢？此外，\Quote{主}和\Quote{基督}这些称谓不属于祂的本质，而是属于祂的尊位；因为一个指祂的权能，另一个指祂受傅油。那么，关于天主子你要说什么呢？因为即使祂像你断言的那样是受造的，这个论证也不能成立。因为祂并非先被创造，然后天主拣选了祂；祂所拥有的也不是一个可被丢弃的王国，而是本质上属于祂本性的王国；因为当被问祂是否是君王时，祂回答说：\Quote{我为此而生。}\BibleRef{若 18:37}但伯多禄是这样说的，如同论及一位被拣选者，因为他的论述全然涉及那\OriginalTerm{降生奥迹}{Dispensation}。\par

4. 你若因伯多禄如此说而感惊奇，有何不可？因为保禄与雅典人辩论时，只称祂为\Quote{人}，说：\BibleQuote{藉着祂所指定的那个人，祂已给众人一个保证，因为祂使祂从死人中复活了。}\BibleRef{宗 17:31} 他丝毫未提及\BibleQuote{天主的形体}\BibleRef{斐 2:6}，也未说祂\BibleQuote{与祂同等}，或祂是\BibleQuote{祂光荣的反映}\BibleRef{希 1:3}。这是有原因的。说这样的话的时候尚未来到；他所满足的，只是他们当时能接受祂是人，且从死人中复活了。基督自己也同样行事，保禄从祂学得，便如此谨慎。因为祂并未立刻向我们启示祂的天主性，而起初被视为先知和善人；但后来藉着祂的作为和言语，祂的真实本性才显明出来。为此，伯多禄起初也采用这方法（因为这是他对犹太人所作的首次讲道）；由于他们还未能清楚理解任何关于祂天主性的事，他便专注于与祂降生成人相关的论证，以致他们藉着这些操练了耳朵，为他其余的教导打开通路。若有人从头细查整篇讲道，会发现我所说的十分明显，因为他（伯多禄）称祂为\Quote{人}，并详述祂的受难、复活以及按血肉的诞生。保禄说：\BibleQuote{按血肉是达味的后裔}\BibleRef{罗 1:3}，也只教导我们\Quote{受造}一词是就祂降生成人而言，正如我们所承认的。可是，雷霆之子如今向我们讲述祂的不可言说与永恒存在，因此他撇下\Quote{受造}一词，而采用\Quote{有}（存在）；然而，若祂是受造的，这一点本是他必须最确切论定的。因为若保禄害怕有些愚昧的人会以为祂比圣父更大，并使生祂的那一位服从于祂（这就是为什么宗徒在致格林多人书信中写道：\BibleQuote{但当祂说万物都服在祂以下时，显然，那使万物服于祂的，是在例外。}\BibleRef{格前 15:27} 然而谁能想像圣父，甚至与万物一同，将服从于圣子？），若我说，他仍害怕这些愚昧的想象，并说：\BibleQuote{那使万物服于祂的，是在例外}；那么，若天主子真是受造的，若望岂不更应害怕有人会假定祂是非受造的，并应首先教训这一点吗？\par

但如今，由于祂是受生的，若望或任何其他宗徒或先知，有充分理由从未断言祂是受造的。倘若祂是受造的，独生子自己也不会对此缄默不言。因为那出于迁就而如此谦卑地谈论自己的那一位，必然不会在这件事上沉默。我认为这样设想并非不合理：祂更可能拥有更高的本性，却对此缄默不语，而不是没有这本性，却略过不提，而不使人知道祂没有它。因为在第一种情形中，有一个沉默的好借口，即祂希望藉着沉默于祂属性的伟大来教导人谦卑；但在第二种情形中，你找不到任何正义的借口为沉默辩解。因为那拒绝许多自己真实属性的那一位，若祂是受造的，为何要对祂曾受造保持沉默？祂为了教导谦卑，常常说出低微的表达，那些本不属于祂，那么若祂真是受造的，岂会不谈论这一点呢？你岂不见祂，为使没有人以为祂不是受生的，而做和说一切事来显示祂是受生的，说出与祂的尊威和本质不相称的话，并降至先知的卑微角色？因为\BibleQuote{我凭我自己不能作什么，我怎样听见，就怎样审判}\BibleRef{若 5:30}；以及\BibleQuote{祂告诉我该说什么，该讲什么}\BibleRef{若 12:49}，诸如此类，仅属于一位先知。若祂出于消除这猜疑的愿望，而不耻于说出这样低微的话，那么若祂是受造的，岂不更会说许多类似的话，使没有人假定祂是非受造的，例如：\Quote{不要以为我是由父所生；我是受造的，而非受生的，也不分享祂的本质。}但事实却相反，祂所做的恰恰相反，说出的话甚至迫使人违背自己的意愿和愿望，而接受相反的意见。例如：\BibleQuote{我在父内，父在我内}\BibleRef{若 14:11}；以及\BibleQuote{斐理伯！这么长久的时间，我和你们在一起，而你还不认识我吗？谁看见了我，就是看见了父}\BibleRef{若 14:9}。以及\BibleQuote{为使众人尊敬子，如同尊敬父一样}\BibleRef{若 5:23}。\BibleQuote{父怎样叫死人复活，使他们生活，子也照样随自己的意愿使人生活}\BibleRef{若 5:21}。\BibleQuote{我父作工到现在，我也作工}\BibleRef{若 5:17}。\BibleQuote{就如父认识我，我也认识父一样}\BibleRef{若 10:15}。\BibleQuote{我与父原是一体}\BibleRef{若 10:30}。祂处处藉着\Quote{如同}、\Quote{照样}以及与父同在等措辞，宣告祂与父不偏不倚的相似。祂在自己内的权能，藉着这些以及许多其他话语彰显出来；例如祂说：\BibleQuote{安静！平静吧！}\BibleRef{谷 4:39}，\BibleQuote{我愿意，你洁净了吧！}\BibleRef{玛 8:3}，\BibleQuote{你又聋又哑的邪魔，我命你从他身上出去！}\BibleRef{谷 9:25}。又说：\BibleQuote{你们一向听过对古人说：不可杀人；但我告诉你们：凡向弟兄发怒的，就要受裁判}\BibleRef{玛 5:21}。祂所颁布的一切其他法律，以及祂所行的奇事，足以显示祂的权能；或者更准确地说，其中极小的一部分，也足以说服任何人，除非是那些完全愚钝的人。\par

5. 但虚荣是一种强大的东西，甚至能使被它迷惑的人对显而易见的真理视而不见，并说服他们去反对别人所认可的事；不仅如此，它还怂恿一些明知并确信真理的人假装无知并进行反对。正如在犹太人身上发生的那样，他们并非由于无知而否认天主之子，而是为了从群众中获得尊荣；福音作者说，\Quote{他们相信了}，却害怕\Quote{被逐出会堂}。\BibleRef{若 12:40} 于是他们将救恩交给了别人。因为，一个人若是如此热切地作今世荣耀的奴隶，就不可能获得来自天主的荣耀。因此，主责备他们说：\Quote{你们既然彼此领受光荣，而不寻求出自信仰的唯一者——天主的光荣，怎么能信呢？} \BibleRef{若 5:44} 这种情欲是一种深深的沉醉，使被其制服的人难以恢复。它将其俘虏的灵魂从天上之事分离，把它们钉在地上，不让它们仰望真光，反而说服它们永远在泥沼中打滚，给予它们如此强大的主人，以致它们无须命令就能统治它们。因为患此病的人自动地、不待吩咐就做他认为会取悦主人一切事。为了他们，他穿着华服，美化面容，这些辛苦不是为自己，而是为别人；他带领一群随从穿过市场，好让别人羡慕他，他所做的一切仅仅是为了取悦世人。还有什么心态比这更可悲呢？为了让别人羡慕他，他不断地奔向毁灭。\par

你愿意了解它行使怎样的专制统治吗？当然，基督的话足以显示这一切。但请再听听以下补充。如果你问任何一个参与国事并花费巨资的人，为什么他们挥霍这么多金钱，如此巨大的开支有何意义；你会从他们那里听到，除了取悦民众之外没有别的。如果你再问民众是什么，他们会说，那是一群混乱、动荡的东西，大部分由愚昧组成，像海浪一样盲目地翻腾，常常由各种相反的意见构成。有这样主人的人岂不是比任何人都可怜？然而，尽管奇怪，世俗之人热衷于这些事并不那么奇怪；但那些声称已从世界抽身的人却患上同样的病，或者更严重的一种，这才是最奇怪的。因为对前者，损失只限于金钱；但对后者，危险涉及灵魂。当人们为了声誉而改变正确的信仰，并侮辱天主以使自己享有声望时，告诉我，他们的行为中难道没有极端的愚昧和疯狂吗？其他的情欲，即使非常有害，至少也带来一些快乐，虽然是短暂即逝的；那些爱金钱、爱酒、爱女人的人，在伤害中也有快乐，虽然是短暂的。但被这种情欲俘虏的人过的是一种不断痛苦、被剥夺享受的生活，因为他们没有得到他们渴望的东西，即来自众人的荣耀。他们以为自己享受它，但实际上并没有，因为他们追求的目标根本不是荣耀。因此，他们的心态不被称为荣耀，而是某种没有荣耀的东西——虚荣，古人这样称呼它，并且很有道理；因为它完全是空的，里面没有任何明亮或荣耀的东西，正像演员的面具表面明亮可爱，但里面是空的（因此，尽管它们比真脸更美，却从未吸引任何人去爱它们），同样，或者更可怜的是，众人的喝彩为我们装饰了这种情欲，它是个危险的对手，残忍的主人。它的外表只是光鲜，但里面不仅是面具的空虚，还充满羞辱，充满野蛮的暴政。那么，有人会问，这种如此不合理、毫无快乐的情欲是从哪里产生的呢？除了来自一个卑劣、渺小的灵魂，还能从哪里呢？被爱名声所俘的人不可能轻易地想象任何伟大或崇高的事；他必定是卑鄙、庸俗、可耻、渺小的。他不是为了德性而做任何事，而是为了取悦那些不值得考虑的人，并且总是考虑他们错误且谬误的意见，这样的人还能有什么价值呢？想想看：如果有人问他，你怎么看待众人？他显然会说，\Quote{他们是无思想的，不值得重视。} 如果有人再问他，\Quote{你愿意像他们一样吗？} 我想他不可能希望像他们一样。那么，寻求那些你绝不愿意与之相似的人的好感，岂不是非常可笑吗？\par

6. 你说他们人数众多，是一个集体？这正是你最应该鄙视他们的原因。如果单独看他们每人都是可鄙的，那么当他们在一起时更是如此；因为当他们聚在一起时，个人的愚昧因数量而增加，变得更大。所以，一个人或许可以单独纠正他们中的一个，但不能在他们聚集时这样做，因为那时他们的愚昧变得强烈，他们像羊一样被带领，彼此跟随各种意见。告诉我，你还要寻求这种世俗的荣耀吗？不要，我恳求你。它使一切颠倒；它是贪婪、诽谤、作假见证、背叛之母；它煽动并激怒那些没有受伤害的人去攻击那些没有伤害他们的人。陷入这种病的人既不知道友谊，也不记得旧日同伴，也不懂得尊重任何人；他把一切良善从灵魂中抛弃，与每个人为敌，不稳定，没有亲情。\par

再者，愤怒的激情虽然专横且难以承受，但并非总是搅扰人，只有当遇到激起它的人时才如此；但\OriginalTerm{虚荣}{vainglory}的激情却始终活跃，可以说，没有一时片刻能止息，因为理智既不阻止也不克制它，它常与我们同在，不仅劝我们犯罪，还从我们手中夺去任何我们可能偶然做对的事，甚至有时根本不让我们做对。如果保禄称\OriginalTerm{贪婪}{covetousness}为偶像崇拜，那么我们该如何称呼它的母亲、根源和源头，即虚荣呢？我们不可能找到任何与它的邪恶相称的术语。亲爱的，如今让我们恢复理智；让我们脱去这污秽的衣裳，将其撕碎并割断，让我们在某个时刻真正自由，并认识到天主赐予我们的高贵；让我们藐视粗俗的喝彩。因为没有什么比这激情更可笑和可耻，没有什么更充满羞愧和侮辱。人们可以从许多方面看到，爱慕荣誉就是耻辱；真正的荣誉在于漠视荣誉，不将其放在心上，而是一切言语行为都按照天主所喜悦的去行。这样，我们就能从那位洞察我们一切行为的天主那里获得赏报，只要我们能满足于唯独祂作观众。当那将要赐予奖赏的天主时刻注视我们的行为时，我们何需其他眼睛呢？这岂不是很奇怪：一个仆人无论做什么，都应为讨主人喜欢而做，只求主人看见，不想引起他人（即使是显贵旁观者）对他行为的注意，只求主人能注意到他；而我们有如此伟大的主，却寻求其他观众，他们观察我们不但不能有益，反而伤害我们，使我们的一切辛劳化为乌有？不可如此，我恳求你们。让我们呼求那位我们将从祂那里领受赏报者来赞许和观看我们的行为。让我们与人的眼睛毫无关系。因为即使我们渴望得到这荣誉，也唯有当我们寻求那唯独来自天主的荣誉时，才会获得它。因为，祂说：\BibleQuote{尊重我的，我必尊重他。} \BibleRef{撒上 2:30} 同样，当我们藐视财富，只寻求那来自天主的财富时，我们便最富有（\BibleQuote{你们先该寻求}，他说，\BibleQuote{天主的国，这一切自会加给你们。} \BibleRef{玛 6:33}）；荣誉也是如此。当日财富或荣誉的赐予不再对我们有危险时，天主便自由地赐予我们；而它们没有危险的时候，就是它们不统治我们、没有权柄支配我们、不像奴隶一样命令我们，而是属于我们如同主人和自由人。因为祂不希望我们贪爱它们的原因，是使我们不被它们统治；如果我们在这件事上成功，祂会极其慷慨地赐予我们。请告诉我，当保禄说：\BibleQuote{我们不寻求人的荣誉，无论是从你们，还是从别人。} \BibleRef{得前 2:6} 时，有什么比他更光荣呢？那么，有谁比那一无所有却拥有一切的人更富有呢？因为如我所说，当我们不被它们控制时，我们就将控制它们，我们就将得到它们。如果我们渴望获得荣誉，就让我们逃避荣誉，这样我们才能在完成天主的法律之后，因基督的恩宠，既获得今世的美物，也获得那所应许的。愿光荣归于基督，及圣父、圣神，直到永远。阿们。\par

\SourceRecord{Translated by Charles Marriott.；Nicene and Post-Nicene Fathers, First Series；Vol. 14.；Edited by Philip Schaff.；Buffalo, NY: Christian Literature Publishing Co.,；1889.；<http://www.newadvent.org/fathers/240103.htm>.}

% structure: p0001 source=h1 target=section reason=page-title

\section{\WorkTitle{若望福音第四篇讲道}}

% structure: p0002 source=h2 target=subsection reason=relative-heading-rank; base=h2

\subsection{\BibleBook{若望}{福音}1:1}

1. 孩童初入学时，教师不给他们连续许多功课，也不一次全布置，而是常常向他们重复同样简短的内容，以便所说的话容易印在他们心里，免得他们一开始就被数量繁多、记忆困难所恼，因困难而生懈怠，变得不那么勤于接收所交托的功课。我愿在你们中间做成同样的事，愿减轻你们的辛劳，便从这神圣的筵席上取食物，一点一点地灌注到你们的灵魂中。因此，我要再次讲解同样的言词，并非重述同样的事，而是只把尚未说出的摆在你们面前。来，让我们再次将我们的讲论应用于开端。\par

\BibleQuote{在起初已有圣言，圣言与天主同在。} 为何其他所有福音作者都以降生奥迹开始？因为玛窦说：\BibleQuote{耶稣基督的族谱，达味之子}；路加也在他的福音开头向我们叙述了有关玛利亚的事；同样，马尔谷也专注于相同的叙述，从那里向我们详细讲述了若翰的史实——当他们以这些事开始时，若望却简短地在稍后的地方暗示它们，说：\BibleQuote{圣言成了血肉} \BibleRef{若 1:14} ，并略过其他一切——祂的受孕、祂的诞生、祂的抚育、祂的成长——立刻向我们谈论祂的永恒受生吗？\par

我现在要告诉你们原因何在。因为其他福音作者最着力于叙述祂在血肉中的来临，恐怕有些人，因心性卑下，为此原因只停留于这些道理，正如撒摩撒特的保禄那样。因此，为了使那些容易陷入这种贪爱尘世之情的人脱离它，并将他们引向天上，他有充分的理由从他的叙述从高处开始，从永恒的存在开始。因为玛窦从黑落德王开始记述，路加从提庇留凯撒开始，马尔谷从若翰的洗礼开始，而这位宗徒却撇下这一切，超越所有时间或年代。他将听者的想象射向\Quote{在起初}，不让它停留在任何一点，也不设立任何界限，像他们在黑落德、提庇留和若翰那里所做的那样。\par

此外，特别值得我们赞叹的是：若望虽然专心于更高的道理，却没有忽略降生奥迹；其他作者虽然专注于叙述降生，也没有缄默不语关于万世以前的存在。有充分的理由：因为推动所有灵魂的是同一圣神；因此他们在叙述中表现出极大的同心。但是，亲爱的，当你听到\Quote{圣言}时，不要容忍那些说祂是受造物的人，也不要容忍那些认为祂仅仅是言语的人。因为天主的言语有很多，是天使执行的，但那些言语中没有一个是神；它们全都是预言或命令（因为在圣经中，通常称天主的法律为祂的命令，预言为言语；所以在论及天使时，祂说：\BibleQuote{大有能力，成就祂的言语} \BibleRef{咏 103:20} ），但这圣言是一个自立体，是从圣父本身无情感地发出的。因为正如我先前所说，祂用\Quote{圣言}这个词表明了这一点。因此，正如\Quote{在起初已有圣言}这句话显明了祂的永恒，同样，\Quote{圣言在起初就与天主同在}这句话向我们宣告了祂的同永恒性。因为当你听到\Quote{在起初已有圣言}时，恐怕你以为祂是永恒的，却又认为圣父的生命与祂有某种间隔和更长的延续，从而将起始归于独生子，于是祂加上\Quote{圣言在起初就与天主同在}；这样，圣父本身是永恒的，因为圣父从未没有圣言，而是天主常与天主同在，然而每位在自己的位格中。\par

那么，有人说，若望怎么会断言祂在世界内，既然祂与天主同在？因为祂既与天主同在，也在世界内。因为圣父和圣子都不受任何限制。既然\BibleQuote{祂的伟大不可测量} \BibleRef{咏 144:3} ，并且\BibleQuote{祂的智慧无数} \BibleRef{咏 146:5} ，显然祂的本质不可能有时序上的开始。你已听到：\BibleQuote{在起初天主创造了天地} \BibleRef{创 1:1} ；你从这\Quote{起初}中理解到什么？显然，它们是在一切可见之物受造之前被创造的。同样，关于独生子，当你听到祂\Quote{在起初}时，要认为祂在一切理智之物和万世之前。\par

但如果有人说：\Quote{祂怎么是子，却不比圣父年轻？因为从某物出来的，必然晚于它所出的那物。} 我们要说，严格讲来，这只是人的推理；在这事上提问的人，也会在其他更不恰当的事上提问；这样的人我们甚至不应该倾听。因为我们现在谈论的是天主，不是谈论人的本性，人的本性是受这些推理的顺序和必然结论支配的。然而，为了软弱之人的确信，我们也要对这些点作些回答。\par

2. 那么，请告诉我，太阳的光辉是从太阳本身的本体发出，还是从其他源头？任何没有丧失理智的人都必须承认，它是从本体本身发出的。然而，虽然光辉是从太阳本身发出的，但我们不能说它在时间上晚于那个物体的本体，因为太阳从未不带光芒出现。现在，如果对于这些可见可感之物，已经表明有某种东西从另一种东西中发出，却并不晚于它所从出的东西；那么，对于那不可见、不可言说的本性，你们为何不信呢？同样的事在那里发生，但以适合那本体的方式。正是为此，保禄也称祂为\Quote{光辉}\BibleRef{希 1:3}，以此表明祂出自祂且与祂\OriginalTerm{同永恒}{Co-eternity}。再者，请告诉我，难道不是所有的世代和一切间隔都是由祂创造的？任何不丧失理智的人必然承认这一点。因此，圣子和圣父之间没有间隔；如果没有间隔，那么祂就不在祂之后，而是与祂同永恒。因为\Quote{先}和\Quote{后}是暗示时间的概念，因为没有时代或时间，没有人能想象这些词语；但天主超越时间和时代。\par

但是，如果在任何情况下你说你找到了圣子的一个开始，那么看看你是否因同样的理由和论据，不被迫也将圣父归为一个开始，虽较早，却仍是一个开始。因为当你为圣子指定了存在的界限和开始时，难道你不是从那一点向上追溯，说圣父在它之前吗？显然你是这样做的。那么，请告诉我，圣父的\OriginalTerm{先行存在}{prior subsistence}有多大？因为无论你说间隔小，还是说间隔大，你同样将圣父带到了一个开始。因为显然，你是通过丈量空间才说它小或大；然而，除非两端都有一个开始，否则无法丈量；因此，就你而言，你已给了圣父一个开始，而且从此以后，根据你的论证，连圣父也不是无始的了。你们看，救主所说的话是真实的，这句话处处显出其力量。那句话是什么？就是\Quote{不尊敬圣子的，就是不尊敬圣父。}\BibleRef{若 5:23}\par

我确实知道刚才所说的话很多人不能领会，因此我们在很多地方避免搅动人类理性的问题，因为其余的人不能跟随这类论证，即使能，它们也没有任何稳固或确定之处。\Quote{因为凡人的思念都是可怜的，我们的计谋都是不确定的。}\BibleRef{智 9:14} 我仍然想问我们的反对者：先知所说的\Quote{在我以前没有别的神被形成，在我以后也没有}\BibleRef{依 43:10}是什么意思？因为如果圣子比圣父年轻，那么祂说\Quote{在我以后也没有}，又是什么意思？你们要取消那\OriginalTerm{独生子}{Only-Begotten}本身的存在吗？你们要么敢于这样做，要么承认圣父和圣子具有不同\OriginalTerm{位格}{Persons}的一个天主性。\par

3. 最后，\Quote{万物是藉着祂造成的}怎可能是真的呢？因为如果有一个时代比祂更古，那么在祂之前的东西怎能由祂造成呢？你们看到，一旦真理被动摇，论证将他们带到了何等大胆的地步？为什么福音作者不说祂是由不存在之物造成的，如同保禄论及万物时所说的\Quote{叫那不存在的成为存在的}\BibleRef{罗 4:17}；却说了\Quote{在起初就有}？这与那说法相反，而且是合理的。因为天主既不是受造的，也没有任何更古老的东西；这些都是外邦人的话。也请告诉我：难道你们不说，那创造者远超祂的工程，无可比拟吗？但由于那从无中生出的与它们相似，那无可比拟的优越性在哪里？\Quote{我是元始，我是终末}\BibleRef{依 44:6}和\Quote{在我以前没有别的神被形成}\BibleRef{依 43:10}这些说法又是什么意思？因为如果圣子不是同一\OriginalTerm{本质}{Essence}，那么就有另一位神；如果祂不是\OriginalTerm{同永恒}{Co-eternal}，那么祂就在祂之后；如果祂不是从祂的本质发出，显然祂是受造的。但如果他们断言，说这些话是为了将祂与偶像区分，那么为什么他们不承认，说\Quote{唯一真神}\BibleRef{若 17:3}也是为了将祂与偶像区分？此外，如果这是为了将祂与偶像区分，你们如何解释整句话？祂说\Quote{在我以后}，\Quote{没有别的神}。这样说，祂并未排除圣子，而是说\Quote{在我以后没有偶像之神}，不是说\Quote{没有圣子}。反对者说：允许这样。那么又如何？\Quote{在我以前没有别的神被形成}这句话，你们会理解为：确实没有偶像之神在祂以前被形成，但一个儿子在祂以前被形成了？有什么恶神会断言这事？我想连撒殚自己也不会这样做。\par

此外，如果祂不与圣父\OriginalTerm{同永恒}{Co-eternal}，你们怎能说祂的生命是无限的？因为如果它有一个先前的开始，尽管它无尽，却不是无限的；因为无限必须是双向无限的。正如保禄也曾宣告，当他说\Quote{祂无始无终}\BibleRef{希 7:3}，用这话表明祂既无开始也无终结。因为正如一端没有限制，另一端也没有。一个方向没有终结，另一个方向没有开始。\par

3. 再者，既然祂是\Quote{\OriginalTerm{生命}{Life}}，怎会有祂不存在的时候呢？因为所有人都必须承认，生命总是存在的，既无开始也无终结，如果它确实是生命，正如它确实是的。因为如果有它不存在的时候，它怎能成为他人的生命，当它自己都不存在时？\par

\Quote{那么，}有人说，\Quote{若望是如何以说\BibleQuote{在起初已有}来定出一个开始呢？}请告诉我，你有没有留意到\Quote{在起初}和\Quote{已有}，难道你不明白\BibleQuote{圣言是}这个说法吗？什么！当先知说：\BibleQuote{从永远到永远，你是天主}\BibleRef{咏 89:2}，他这样说是在给他设定界限吗？不，而是为了宣告祂的永恒。如今请想，此处的情况相同。他并非用此说法来设定界限，因为他没有说\Quote{有一个开始}，而是说\Quote{在起初已存在}；借着\Quote{已有}这个词，他引导你想到圣子是无始的。\Quote{然而请注意，}他说，\Quote{圣父在提到时带有冠词，而圣子却没有。}那么，当宗徒说：\BibleQuote{伟大的天主和我们的救主耶稣基督}\BibleRef{铎 2:13}，以及又说：\BibleQuote{祂是在万有之上，永远受赞美的天主}\BibleRef{罗 9:5}，这又如何呢？的确，此处他提到圣子时没有冠词，但在致斐理伯人书中，他对圣父也是如此，至少他说：\BibleQuote{祂虽具有天主的形象，却没有以自己与天主同等为强夺的}\BibleRef{斐 2:6}；再者，致罗马人书中说：\BibleQuote{愿恩宠与平安由天主我们的父和主耶稣基督赐与你们}\BibleRef{罗 1:7}。此外，在那个地方附上冠词是多余的，因为前面紧接之处已经连续附于\Quote{圣言}。因为正如在论及圣父时，他说：\BibleQuote{天主是神}\BibleRef{若 4:24}，我们并不因为冠词未附于\Quote{神}就否认天主的神性；同样，此处虽然冠词未加于圣子，圣子并不因此就成为较低的天主。何以如此？因为他说\Quote{天主}，又说\Quote{天主}，并没有向我们揭示这神性中的任何差别，反而相反；因为他事先已说过\BibleQuote{圣言就是天主}，以免有人以为圣子的神性较低，他立即加上真正神性的特征，包括永恒（因为他说\BibleQuote{祂在起初就与天主同在}），并将造物主的职责归于祂。因为\BibleQuote{万物是借着祂造的，凡受造的，没有一样不是借着祂造的}；圣父也借着先知在各地宣告这正是祂本体的特别特征。先知们不断地从事这种论证，不仅凭其本身，而且当他们反对向偶像表示尊崇时，有人说：\BibleQuote{愿那没有创造天地的神祇从地上消灭}\BibleRef{耶 10:11}；又说：\BibleQuote{我亲手铺张了诸天}\BibleRef{依 44:24}；他随处以此作为神性的标记来宣布。圣史本人并不满足于这些话，又称祂为\Quote{生命}和\Quote{光}。如果祂永远与圣父同在，如果祂亲自创造了万物，如果祂使一切存在，并维系一切（因为他用\Quote{生命}来意指这个），如果祂光照一切，谁还能愚昧到说圣史欲借着那些最足以表达其平等与无区别的表达方式，反而教导神性较低呢？让我们不要混淆受造物与造物主，免得我们也听到有人对我们说：\BibleQuote{他们敬奉受造物而不敬奉造物主}\BibleRef{罗 1:25}；因为虽然有人声称这是指着诸天说的，但在谈论诸天时，他明确地说，我们不可敬奉受造物，因为这是外教人的行为。\par

4. 因此，让我们不要使自己处于这诅咒之下。因为天主圣子为此而来，为要把我们从这种敬奉中解救出来；为此祂取了奴仆的形象，为要释放我们脱离这种奴役；为此祂受人唾污，为此祂受人掌掴，为此祂忍受了羞耻的死亡。我恳求你们，不要使这一切归于徒然，不要回到我们先前的罪恶，或者更确切地说，回到更严重的罪恶；因为敬奉受造物并不等同于以我们所能的方式将造物主贬低到受造物的卑贱。因为祂继续是祂所是的样子；正如圣咏作者所说：\BibleQuote{你仍是一样，你的岁月没有穷尽}\BibleRef{咏 102:27}。让我们照着从我们祖先所领受的来光荣祂，让我们以信德和作为来光荣祂；因为如果我们生活腐化，纯正的教义对我们得救毫无益处。让我们按照中悦天主的方式来安排生活，远离一切污秽、不义和贪婪，如同在世上的客旅、外方人和寄居者。如果有人拥有许多财富和财产，让他像一个旅居者那样使用它们，他无论愿意与否，不久都要离开它们。如果有人受了他人的伤害，让他不要永远发怒，不，甚至不要一时发怒。因为宗徒只允许我们发泄怒气不超过一天。\par

他说：\Quote{不要让太阳落在你们的愤怒上}\BibleRef{弗 4:26}，这是有理由的；因为即使在这么短的时间内没有不愉快的事发生，也足可欣慰；但如果黑夜也赶上我们，所发生的事就变得更加痛苦，因为记忆将我们的愤怒之火放大了一万倍，我们在闲暇中更加痛苦地追究。因此，在我们获得这种有害的闲暇并燃起更旺的火之前，他嘱咐我们事先制止并扑灭祸害。因为愤怒的情绪是猛烈的，比任何火焰更猛烈；所以我们需要赶紧防止火焰，不让它燃烧高涨，因为这种弊病会变成许多邪恶的原因。它推倒了整座房屋，瓦解了旧日的友谊，并在极短的时间内造成了无法补救的悲剧。经上说：\Quote{愤怒的支配，是他的毁灭}\BibleRef{德 1:22}。那么，我们不要留下这样一头野兽没有缰绳，而要给它套上一个各方面都坚固的嚼子，就是对将来审判的畏惧。每当有朋友使你伤心，或家人激怒你，就该想想你得罪天主所犯的罪，并想到你对他的仁慈会使那审判对你更为宽大（他说：\Quote{宽恕，你们就会受宽恕}\BibleRef{路 6:37}），你的愤怒就会迅速逃遁。\par

此外，还要考虑这一点：是否曾有一次你被带入残暴之中，但控制了自己；另一次你被情绪拖着走？比较这两种境遇，你必从中获得大益处。因为请告诉我，你何时称赞自己？是当你被击败时，还是当你有克制力时？我们在第一种情形下不是痛责自己，即使无人责备也感到羞愧，并且为所说所行生出许多悔恨吗？但当我们获得克制时，我们不是骄傲自大，像征服者一样欢欣吗？因为在愤怒方面的胜利，不是以牙还牙（那是彻底失败），而是温顺地忍受虐待和辱骂。获胜不是施加恶事，而是忍受恶事。因此当发怒时，不要说：\Quote{我一定要报复，} \Quote{我一定要复仇}；不要坚持对那些劝你获胜的人说：\Quote{我不能容忍那人嘲笑我，然后逍遥法外。}他绝不会嘲笑你，除非你报复；或者即使他嘲笑你，他也会像愚人一样。你在得胜时不要向愚人寻求荣誉，而要认为从有见识的人那里得来的荣誉就够了。不仅如此，我为何将一小撮低贱的观众摆在你面前，而把观众局限于人呢？直接仰望天主：祂必赞美你，而蒙祂认可的人不应寻求来自凡人的荣誉。凡人的荣誉常来自他人的奉承或仇恨，毫无益处；但天主的判断没有这种不公，并给祂所认可的人带来极大益处。那么，让我们追求这种赞美。\par

你想知道愤怒是何等邪恶吗？请站在一旁观看别人在市场中争吵。你自己不容易看到事情的丢脸之处，因为你的理性被蒙蔽、醉倒；但当你摆脱了情绪，判断力健全时，在别人身上查看自己的情形。我请你观察，人群聚拢，愤怒的人像疯子一样在当中做出可耻的事。因为当情绪在胸中沸腾，变得激动和野蛮时，口吐火焰，眼冒火光，整个脸肿胀，双手胡乱伸展，双脚可笑地乱跳，并扑向那些阻拦他们的人，与疯子毫无区别，对这一切毫无知觉；不，与野驴无异，又踢又咬。的确，爱发怒的人并不优雅。\par

然后，当他们做出如此极其可笑的行为后，回到家中醒过神来，便感到更大的痛苦和极度的恐惧，思量着他们发怒时有哪些人在场。因为他们如同狂乱之人，当时并不认识旁观者，但等到恢复理智后，才考虑：那些人是朋友吗？是仇敌和敌人看着吗？他们对两者都同样害怕：对前者，因为他们会谴责他们，使他们更加羞愧；对后者，因为他们会因此幸灾乐祸。如果他们甚至动了手，那么恐惧就更加紧迫；例如，生怕受害者发生任何非常严重的事情，如发烧导致死亡，或危险的肿痛使他有生命之虞。并且他们说：\Quote{有什么必要}（他们说）\Quote{我去打架、施暴、争吵？愿这类事消失。}然后他们诅咒那不幸的事务——致使他们开端的缘由，而更愚蠢的人则将所做之事的罪责推给\Quote{邪灵，}和\Quote{凶时，}；但这些并非来自凶时（因为没有凶时这样的事），也不是来自邪灵，而是来自被情欲控制之人的邪恶；他们吸引邪灵到自己身上，并给自己招来一切可怕之事。\Quote{但人心膨胀，}有人说，\Quote{被侮辱刺痛。}我知道，这便是我钦佩那些制伏这可怕野兽之人的原因；然而，如果我们愿意，就有可能击退这情欲。因为当统治者侮辱我们时，我们为何不感到它？因为恐惧平衡了情欲，使我们畏惧它，完全不允许它萌发。同样，我们的仆人，虽然被我们以千万种方式侮辱，为何能默默忍受一切？因为他们也有同样的约束加于他们。你们不仅要想到对天主的畏惧，更要想到是天主亲自侮辱你们，他命令你们缄默，这样你们就会温顺地承受一切，并对侵犯者说：我怎能向你发怒？另有别者约束了我的手和舌头；这句话对你们自己和对方都是一种明智的劝告。现今，我们为了人而忍受不可忍之事，常对侮辱我们的人说：\Quote{某人侮辱了我，不是你。}难道我们对天主就不采用同样的谨慎吗？我们如何还能期望得到宽恕？我们应向自己的灵魂说：\Quote{是天主握着我们的手，他现在侮辱我们；我们不要反抗，不要让天主在我们眼中比人更不受尊敬。}你们听了这话战栗吗？我希望你们不仅在言语上战栗，更要在行为上战栗。因为天主命令我们，当被掌掴时，不仅要忍受，甚至要献上自己去承受更坏的事；而我们在反抗他时是如此激烈，以至于不仅拒绝献上自己去受苦，反而复仇，甚至常常首先发起攻击，并认为若不以其人之道还治其人之身便是耻辱。是的，祸患在于：当我们完全被击败时，却自以为得胜；当我们被压在下面，受了魔鬼千万次打击时，却想象着我们制伏了他。因此，我劝你们，要明白这胜利的本质，并追求这种本质。受苦就是获得冠冕。如果我们希望被天主宣布为胜利者，就不要在这些竞赛中遵守异教赛会的规则，而要遵守天主的规则，并学会以忍耐忍受一切；这样我们就能胜过对手，并获得当下及承诺的美物，借着我们主耶稣基督的恩宠与慈爱，他同圣父及圣神，永受光荣、权能和尊威，自今而始，以至世世，无穷世。阿们。\par

\SourceRecord{Translated by Charles Marriott.；Nicene and Post-Nicene Fathers, First Series；Vol. 14.；Edited by Philip Schaff.；Buffalo, NY: Christian Literature Publishing Co.,；1889.；<http://www.newadvent.org/fathers/240104.htm>.}

% structure: p0001 source=h1 target=section reason=page-title

\section{若望福音讲道 第五篇}

% structure: p0002 source=h2 target=subsection reason=relative-heading-rank; base=h2

\subsection{\BibleRef{若 1:3}}

1. \Person{梅瑟}在\WorkTitle{旧约}的历史和著述的开端，向我们讲述感官所及的事物，并详细列举它们。因为他说：\BibleQuote{在起初，天主创造了天地}，接着他又加上，光被创造，第二层天和星辰，各类生物，以及——我们无需细数——其他一切。但这福音作者，一笔带过，用一句话包括这些事和超越这些的事；有理由，因为它们为听众所熟知，并且他正赶往一个更伟大的主题，把他的全部论述建立起来，为要谈论的不是工程，而是创造者，即产生万物的那一位。因此，\Person{梅瑟}虽然选择了受造物中较小的部分（因为他没有对我们提及不可见的能力），却详述这些事；而\Person{若望}，因急于上升到创造者本身，越过这些事以及\Person{梅瑟}保持沉默的那些事，将它们概括在一句简短的话里：\BibleQuote{万物是藉着祂造成的}。并且为使你不以为他只是谈论\Person{梅瑟}所提及的一切，他补充说：\BibleQuote{凡受造的，没有一样不是藉着祂造成的}。这就是说，在受造物中，没有一样，无论是可见的还是可理解的，不是藉着圣子的能力而存在的。\par

因为我们不会像异端者那样在\Quote{没有一样}之后点句号。他们因为要使圣神成为受造的，便说：\Quote{凡受造的，在祂内就是\OriginalTerm{生命}{Life}}；但这样所说的话就变得不可理解了。首先，此处并不是提及圣神的时候，如果他有意如此，为何说得这样含糊？因为这句话与圣神有何关系，哪里清楚？此外，我们将会发现，按照这种论点，受造的并非圣神，而是圣子自己藉着自己受造的。但请振奋，以免所说的话从你们面前溜走；来，让我们暂时照他们的方式读，这样其荒谬之处对我们会更清楚。\Quote{凡受造的，在祂内就是\OriginalTerm{生命}{Life}}。他们说圣神被称为\Quote{\OriginalTerm{生命}{Life}}。但这\Quote{\OriginalTerm{生命}{Life}}也被发现是\Quote{\OriginalTerm{光}{Light}}，因为他又说：\BibleQuote{那生命就是人的光}\BibleRef{若 1:4}。因此，按照他们，这里的\Quote{人的光}是指圣神。然而，当他接着说：\BibleQuote{有一个人，由天主派遣，来为此光作证}\EditorialNote{6–7节}，他们必须断言，这也是论圣神说的；因为他前面称为\Quote{\OriginalTerm{圣言}{Word}}的那位，在后面他称为\Quote{天主}和\Quote{\OriginalTerm{生命}{Life}}和\Quote{\OriginalTerm{光}{Light}}。他说这\Quote{\OriginalTerm{圣言}{Word}}就是\Quote{\OriginalTerm{生命}{Life}}，这\Quote{\OriginalTerm{生命}{Life}}就是\Quote{\OriginalTerm{光}{Light}}。如果现在这圣言是生命，并且这圣言和这生命成了血肉，那么生命——即圣言——\BibleQuote{成了血肉，我们见了祂的荣耀，正如父独生者的荣耀}。那么，如果他们说此处圣神被称为\Quote{\OriginalTerm{生命}{Life}}，请考虑会有什么奇怪的后果：成了血肉的将是圣神，而非圣子；圣神将成为独生子。\par

那些如此读这段经文的人，若不坠入此，也会在躲避此中坠入另一最奇怪的结论。如果他们承认这些话是对圣子说的，却不停顿或不照我们的读法去读，那么他们就会断言圣子是藉着自己被造的。因为，如果\Quote{\OriginalTerm{圣言}{Word}就是\OriginalTerm{生命}{Life}}，并且\Quote{凡在祂内受造的，就是\OriginalTerm{生命}{Life}}，按照这种读法，祂就是在自己内、藉着自己受造的。然后，在他插入一些话之后，他又加上：\BibleQuote{我们见了祂的荣耀，正如父独生者的荣耀}\BibleRef{若 1:14}。看，按照那些主张这些事的人的读法，圣神也被发现是个独生子，因为关于祂的一切宣告都是由他发出的。看，当话语偏离真理时，它被歪曲到何处，产生何等奇怪的后果！\par

有人会说：难道圣神不是\Quote{\OriginalTerm{光}{Light}}吗？祂是光：但此处并未提及圣神。因为就连天主（圣父）也被称为\Quote{\OriginalTerm{神}{Spirit}}——即无形体——但并非一提到\Quote{\OriginalTerm{神}{Spirit}}就一定是指天主（圣父）。我们这样论圣父，有什么好奇怪的呢？我们甚至不能论护慰者说，凡提到\Quote{\OriginalTerm{神}{Spirit}}就一定是指护慰者，尽管这是祂最特有的名字；但并非总是提到\Quote{\OriginalTerm{神}{Spirit}}就是指护慰者。所以，\Person{基督}被称为\BibleQuote{天主的德能}\BibleRef{格前 1:24} 和\BibleQuote{天主的智慧}；但并非一提到\Quote{天主的德能}和\Quote{天主的智慧}就是指基督；同样，在此段中，尽管圣神确实赐予\Quote{\OriginalTerm{光}{Light}}，但福音作者现在并非在论圣神。\par

当我们把他们从这些奇怪的意见中排除出去时，那些竭尽全力反对真理的人说（仍然固守同一种读法）：\Quote{凡藉着祂存在的，就是\OriginalTerm{生命}{life}，因为}有人说：\Quote{凡存在的，就是\OriginalTerm{生命}{life}}。那么，你如何解释\Place{索多玛}人的惩罚、洪水、地狱之火以及无数类似的事？但有人说：\Quote{我们是在谈论物质受造物}。那么，这些也完全属于物质受造物。但为了我们多多驳斥他们的论据，我们要问他们：\Quote{木头是生命吗}，请告诉我？\Quote{石头是生命吗}？这些没有生命、不能动弹的东西？不，人是绝对的生命吗？谁会这样说？他不是纯粹的生命，而是能够接受生命的。\par

2. 请看，这里又有一个荒谬之处；通过同样的推论链条，我们将把论证引向这样的地步，以致连你们也能从中看出他们的愚蠢。他们以此方式断言一些绝不适宜于圣神的事。当他们从另一个立足点被驱逐后，便将那些他们先前认为适宜于圣神的话应用到人身上。然而，我们也这样来审视这段经文本身。受造物现在被称为\Quote{生命}，因此，它也是\Quote{光}，而若翰来为它作证。为什么他也不被称为\Quote{光}呢？他说他\BibleQuote{不是那光}\BibleRef{若 1:8}，而他却属于受造之物？那么他为什么不是\Quote{光}呢？他\BibleQuote{在世界，世界也藉着他造成的}\BibleRef{若 1:10}，这又如何理解？受造物在受造物中，受造物由受造物造成？但\BibleQuote{世界却不认识他}\BibleRef{若 1:10}又是什么意思？受造物如何不认识受造物？\BibleQuote{凡接待他的，他就赐他们权能，作天主的儿女}\BibleRef{若 1:12}。但笑已足够。其余的事，我留给你们去攻击这些怪诞的推理，免得我们似乎是故意为笑而笑，无端浪费时间。因为如果这些话既不是指着圣神说的（已经证明不是），也不是指着任何受造之物说的，而他们仍然坚持同样的读法，那么就会得出一个比我们先前提到的任何结论都更荒谬的结论：即圣子是由他自己造成的。因为如果圣子是真光，而这光就是生命，这生命是在他内造成的，那么按照他们自己的读法，必然得出这个结果。让我们放弃这种读法，回到公认的读法和解释上来。\par

那是什么呢？就是把句子断在\Quote{受造}，下一句以\BibleQuote{在他内有生命}开始。福音所说是这样的：\BibleQuote{凡受造的，没有一样不是藉着他受造的}；他说，凡受造的受造之物，没有一样不是藉着他受造的。你们看，通过这一简短的补充，他如何纠正了所有令人困惑的难题；因为说了\BibleQuote{没有一样不是藉着他受造的}，然后又加上\BibleQuote{凡受造的}，这就包括了理智所能认知的事物，却排除了圣神。因为在他说了\BibleQuote{万物是藉着他造的}和\BibleQuote{没有一样不是藉着他受造的}之后，他需要这个补充，免得有人说：\Quote{如果万物都是藉着他造的，那么圣神也是受造的。}他说：\Quote{我，}他回答说，\Quote{主张凡受造的，无论是看不见的，无形的，还是在天上的，都是藉着他造的。为此，我没有绝对地说\InnerQuote{万物}，而是\InnerQuote{凡受造的}，即\InnerQuote{受造之物}，但圣神是非受造的。}\par

你们看见他教导的精确性吗？他先暗示了物质世界的受造（因为梅瑟在他之前已论及这些），然后引导我们从那里上升到更高的事物，即无形无象者，把圣神从一切受造中分别出来。同样，保禄受同一恩宠的感动，说道：\BibleQuote{因为万有都是藉着他造的}\BibleRef{哥 1:16}。你们也再次看到同样的精确性。因为同一圣神也感动了这位灵魂。为了不让任何人因受造物不可见而将其从天主的工作中排除，也不致将护慰者与它们混淆，在列举了众人所知的感官事物之后，他也列举了天上的事物，说道：\BibleQuote{无论是坐位的，主治的，执政的，掌权的}\BibleRef{哥 1:16}；因为每一个前面加的\Quote{无论是}，向我们表明的无非就是：\BibleQuote{万有都是藉着他造的，凡受造的，没有一样不是藉着他造的。}\par

但如果你们认为\Quote{藉着}一词是卑下的标记（即使基督成为工具），那么请听他说道：\BibleQuote{主啊，你在起初奠定了大地，诸天是你手的化工}\BibleRef{咏 102:25}。他以创造者的身份把论及圣父的话用在圣子身上；除非他认为圣子是创造者且不隶属于任何人，否则他不会这样说。而如果这里用了\Quote{藉着他}，那唯一的理由就是防止任何人以为圣子是无始者。因为就创造者的称号而言，他丝毫不逊于圣父；请听他亲自说：\BibleQuote{父怎样叫死人起来，使他们活着，子也照样随自己的意愿使人活着}\BibleRef{若 5:21}。如果在旧约中论及圣子的话是\BibleQuote{主啊，你在起初奠定了大地}，那么他创造者的称号就很明显。但如果你们说先知这话是指圣父说的，而保禄把论及圣父的话用在了圣子身上，结论仍是一样的。因为保禄不会擅自认为同一表达适合于圣子，除非他确信圣父与圣子之间有同等的荣耀；因为把适合于无可比拟的本性的话引用到低于它、不及它的本性上，那是极端鲁莽的行为。但圣子并不低于或不及圣父的本体；因此保禄不仅敢于用这些词语论及他，而且还用了其他类似的词语。至于\Quote{从谁而来}这一说法，你们认为它只适合圣父，但保禄也用之于圣子，他说：\BibleQuote{全身都靠着关节和筋络得到滋养和联络，因着天主的增长而增长}\BibleRef{哥 2:19}。\par

3. 他并不满足于此，还以另一种方式堵住你们的口，将你们所说作为低等标记的\Quote{藉着谁}这一表达应用于圣父。因为他说：\BibleQuote{天主是信实的，你们原是被祂所召，进入祂圣子的共融中。}\BibleRef{格前 1:9}；又说：\BibleQuote{因祂的旨意}\BibleRef{格前 1:1}；在另一处：\BibleQuote{因为万有都出于祂，藉着祂，归于祂。}\BibleRef{罗 11:26}。而且，\Quote{从谁}这一表达不仅归于圣子，也归于圣神；因为天使对若瑟说：\BibleQuote{不要害怕娶你的妻子玛利亚，因为那在她内受孕的是出于圣神。}\BibleRef{玛 1:20}。同样，先知也不认为将属于圣神的\Quote{在谁内}这一表达应用于圣父是不合适的，当他（先知）说：\BibleQuote{在天主内，我们将奋勇行事。}\BibleRef{咏 59:12}。保禄也说：\BibleQuote{祈望或许藉着天主的旨意，终能顺利来到你们那里。}\BibleRef{罗 1:10}。他又将同一表达用于基督，说：\BibleQuote{在基督耶稣内。}\BibleRef{罗 6:11}。总之，我们常常并持续地发现这些表达彼此互换；这不可能发生，除非在每一处它们所指的是同一\OriginalTerm{本性}{Essence}。并且，为免你们以为\Quote{万有是藉着祂造的}这话在此处是指祂的奇迹（因为其他福音作者已论述了这些），他继续往下说：\Quote{祂在世界，世界也是藉着祂造的}；（但并非圣神，因为圣神不属于受造物之列，而是超越一切受造物的。）\par

现在让我们注意下文。若望已经论及创造之工：\Quote{万有是藉着祂造的，凡受造的，没有一样不是藉着祂造的}，接着论及祂的\OriginalTerm{上智照顾}{Providence}，他说：\Quote{在祂内有生命}。为使无人疑惑如此众多伟大之物是\Quote{藉着祂造的}，他补充说：\Quote{在祂内有生命}。因为正如那作为深渊之母的泉源，无论你汲取多少，泉源并不减少；同样，\OriginalTerm{独生子}{Only-Begotten}的能力，无论你相信已由它产生和造出了多少，它也丝毫未减。或者，用一个更熟悉的例子，我将举光的例子，宗徒本人立即添加了这句话，说：\Quote{这生命就是光}。正如光，无论它照亮多少千万人，自身的光明并不减损；同样，天主在开始祂的工作之前和完成它之后，始终不变，毫无亏损，也不因受造物的伟大而疲惫。不仅如此，如果需要造一万个，甚至无限多个这样的世界，祂仍然是一样的，足以不仅创造它们，而且在它们被造之后管理它们。因为这里的\Quote{生命}一词不仅指创造的行动，也指关于受造物长存的上智照顾；它还预先奠定了复活的理论，是这些奇妙喜讯的开始。既然\Quote{生命}已来到我们中间，死亡的权势就被瓦解；既然\Quote{光}已照耀我们，就不再黑暗，而是生命永远存留于我们内，死亡不能战胜它。因此，凡归于圣父的，也绝对归于祂（基督），即\BibleQuote{在祂内我们生活、行动、存在。}\BibleRef{宗 17:28}正如保禄所表明的，他说：\BibleQuote{万有都是藉着祂造的}\BibleRef{哥 1:16}，以及\BibleQuote{万有也藉着祂而维系}\BibleRef{哥 1:17}；为此祂也被称为\Quote{根}和\Quote{基础}。\par

但当你们听到\Quote{在祂内有生命}时，不要以为祂是一个复合的存在，因为在后面祂也论及圣父说：\BibleQuote{就如父在自己内有生命，祂也赐给子在自己内有生命}\BibleRef{若 5:26}；正如你们不会因这一表达而说圣父是复合的，同样你们也不能这样说圣子。因此，在另一处祂说：\BibleQuote{天主就是光}\BibleRef{若一 1:5}，并且在别处说祂\BibleQuote{住在不可接近的光中}\BibleRef{弟前 6:16}；然而这些表达并非要我们假定一个复合的本性，而是要一步一步地将我们引向最高的教义。因为既然群众中的人不容易理解祂的生命如何是\OriginalTerm{位格生命}{Life Impersonate}，祂首先用了那个较为卑下的表达，然后引导（如此）受训的人走向更高的教义。因为那曾说过\BibleQuote{祂赐给祂（圣子）在自己内有生命}\BibleRef{若 5:26}的同一者，在另一处说：\BibleQuote{我是生命}\BibleRef{若 14:6}；又在另一处说：\BibleQuote{我是光}\BibleRef{若 8:12}。请问，这\Quote{光}的本性是什么？这种（光）不是感官的对象，而是理智的对象，光照灵魂本身。既然基督以后要说：\BibleQuote{除非父吸引他，没有人能到我这里来}\BibleRef{若 6:44}；宗徒在此处预料到一个反驳，并宣告是祂（圣子）\BibleQuote{赐光}\BibleRef{若 1:9}；这样，即使你们听到关于圣父的类似说法，也不会说它只属于圣父，也属于圣子。因为，祂说：\BibleQuote{凡父所有的，都是我的}\BibleRef{若 16:15}。\par

首先，圣史已经教导我们关于创造的事，之后他告诉我们祂藉着来临所供给我们的灵魂的福祉；这些他在一句话中暗昧地描述了，当他说：\BibleQuote{生命是人的光。}\BibleRef{若 1:4}他没有说：\Quote{是犹太人的光}，而是普遍地\Quote{人的光}：不独犹太人，连希腊人也得来到这认识，而这光是一个向所有人提供的共同恩赐。\Quote{他为何不加\InnerQuote{天使}，却说\InnerQuote{人？}？}因为当下他的言论是关于人的本性，并且祂来到他们中间，传报喜讯。\par

\BibleQuote{光在黑暗中照耀。}\BibleRef{若 1:5}祂称死亡和错误为\Quote{黑暗}。因为感官所感知的光并不在黑暗中照耀，而是远离它；但基督的宣讲却在盛行的错误中照耀出来，并使之消失。而祂借着忍受死亡，如此地战胜了死亡，以致祂拯救了那些已被死亡所掌控的人。既然死亡没有战胜它，错误也没有，因为它处处光明，并以其本身的力量照耀，因此他说：\par

\Quote{黑暗没有把握住它。}因为它不能被战胜，也不会住在那些不愿被光照的灵魂中。\par

4. 但你不要因它没有摄取所有而烦恼，因为天主不是借着必然和强制，而是借着意愿和同意，将我们带到祂自己面前。因此你不要向这光关上你的门，你将会享受极大的幸福。但这光是由信德而来，当它来到时，它大大光照那接受它的人；并且你若显示出纯全的生活（适合它），它就会持续内住不离。\Quote{因为，}祂说，\BibleQuote{那爱我的，必遵守我的命令；我和我父要到祂那里去，并要在祂那里作我们的住所。}\BibleRef{若 14:23}就如一个人若不开眼，就不能正确地享受阳光；同样，一个人若没有扩展他灵魂的眼睛，并使之在各方面锐利，也不能大量分享这光辉。\par

但这如何成就呢？当我们从所有情欲中洁净了灵魂之时。因为罪是黑暗，且是深沉的黑暗；正如清楚的，因为人无意识地秘密地行暗昧之事。因为\BibleQuote{凡作恶的，恨光，也不来就光。}\BibleRef{若 3:20}并且\BibleQuote{他们暗中所行的事，连提起来也是可耻的。}\BibleRef{弗 5:12}因为就如一个人在黑暗中既不知友也不知敌，而且不能察觉任何物体的特性；在罪中也是如此。因为那贪求更多利益的人，在友敌之间不作区别；嫉妒的人以敌意的眼光看待与他非常亲密的人；阴谋者与所有人一同陷入必死的争吵。简言之，就识别物体的性质而言，犯罪的人并不比醉酒或疯狂的人好多少。又如夜间，木头、铅、铁、银、金、宝石，对我们来说都显得一样，因为缺少显示它们区别的光；同样，过不洁生活的人，既不知道节制的卓越，也不知道哲学的美。因为黑暗中，正如我刚才所说，纵使宝石陈列出来，也不显示它们的光泽，不是由于它们自己的本性，而是由于观看者缺乏辨别力。而这对我们这些在罪中的人，并非唯一的祸患，还有这个：我们生活在不断的恐惧中：正如人在无月之夜行走而颤抖，尽管无人惊吓他们；同样，行不义的人也不能有信心，尽管无人控告他们；但他们害怕一切，猜疑一切，被自己的良心刺痛：对他们来说，一切都充满了恐惧和痛苦，他们观察周围的一切，被一切事物惊吓。那么，让我们逃离如此痛苦的生活，尤其因为在这痛苦之后，死亡将接踵而至；一种不死的死亡，因为那地方的惩罚将永无止境；而在今世中，他们（犯罪的人）不比疯子好多少，因为他们梦见不存在的事物。他们认为自己富有而实际不富有，认为自己享乐而实际不享乐，直到他们从疯狂中解脱并摆脱睡眠，才正确察觉这种欺骗。因此，\Person{保禄}劝勉众人要清醒和警醒；基督也命令同样的事。因为清醒和警醒的人，纵使被罪俘获，也能迅速击退它；而睡着并失去理智的人，却察觉不到自己是如何被罪囚禁的。\par

那么，我们不要睡觉。现在不是黑夜，而是白昼。因此让我们\BibleQuote{行事为人要端正，如同在白昼}\BibleRef{罗 13:13}；没有什么比罪更不当的了。就失当而言，裸体行走并不像犯罪和行恶那样严重。裸体可能因贫穷而引起，算不上多大的责备；但没有什么比罪人更可耻、更不受尊敬的。让我们想一想那些因敲诈或欺骗而来到法庭的人：他们的无耻、谎言和狂妄显得多么卑鄙可笑。但我们是如此可怜可悲，连自己穿衣服歪斜都忍受不了；甚至看见别人这样，我们还会帮他纠正；然而，尽管我们和邻人都头朝下走路，我们却毫不知觉。请问，还有什么比一个去找妓女的人更可耻？还有什么比一个傲慢、污言秽语或嫉妒的人更可鄙？那么，为什么这些事不像裸体行走那样被视为丢脸呢？仅仅是出于习惯。没有人愿意忍受裸体；但人们却不断毫无畏惧地犯这些事。如果一个人进入天使的集会，那里从未发生这类事，他就会清楚地看到这种行为的巨大可笑。我为什么说天使的集会呢？即使在我们中间的宫殿里，如果有人带进一个妓女并与之行乐，或饮酒过量，或做出其他任何类似的无耻之事，他都会受到严厉的惩罚。如果人们在宫殿中做这些事是不可容忍的，那么当国王无处不在、观察所行的一切时，我们若敢做这些事，就要受最严厉的惩罚。因此，我劝你们，要在生活中彰显极大的温良和纯洁，因为我们有一位不断注视我们一切行为的国王。为了这光永远丰盛地照亮我们，让我们欣然接受这些明亮的射线，如此，我们将享受今世和来世的美好事物，借着我们的主耶稣基督的恩宠和慈爱，祂和圣父、圣神同享荣耀，永远永远。阿们。\par

\SourceRecord{Translated by Charles Marriott.；Nicene and Post-Nicene Fathers, First Series；Vol. 14.；Edited by Philip Schaff.；Buffalo, NY: Christian Literature Publishing Co.,；1889.；<http://www.newadvent.org/fathers/240105.htm>.}

% structure: p0001 source=h1 target=section reason=page-title

\section{若望福音 第六篇讲道}

% structure: p0002 source=h2 target=subsection reason=relative-heading-rank; base=h2

\subsection{若望福音 1:6}

1. 在序言中向我们谈论了关于天主圣言至关重要的事情之后，（福音作者）继续他的进程，接着按次序来到圣言的前驱，即与他同名的若翰。现在你听到他\BibleQuote{由天主派遣}，今后不要以为他所说的话只是人的话；因为他所说的一切不是出于他自己，而是出于那派遣他的那位。因此他被称为\BibleQuote{使者}\BibleRef{拉 3:1}，因为使者的优秀之处在于他不说自己的话。但此处的\OriginalTerm{是}{was}一词并非指他存在，而是指他作为使者的职务；因为\BibleQuote{\InnerQuote{有}一个人由天主派遣}被用来代替\Quote{一个人\InnerQuote{被派}由天主}。\par

那么，有些人怎么说，\BibleQuote{具有天主的形体}\BibleRef{斐 2:6}这句话不是指祂与圣父不变相似，因为未加冠词？请注意，这里任何地方都没有加冠词。那么这些话难道不是关于圣父的吗？那么我们对那先知说什么呢？他说：\BibleQuote{看，我派遣我的使者在你面前，他要预备你的道路}\BibleRef{拉 3:1}？因为\Quote{我的}和\Quote{你的}这两个词表达了两个位格。\par

% structure: p0005 source=h2 target=subsection reason=relative-heading-rank; base=h2

\subsection{若望福音 1:7}

这是什么？或许有人会说，仆人为主作证？当你看到祂不仅被祂的仆人作证，甚至来到他那里，与犹太人一起被他施洗，你难道不更加惊讶和困惑吗？但你应当不烦恼、不混乱，反而为这种不可言说的美善感到惊叹。如果还有人继续困惑混乱，祂会对这样的人说祂对若翰所说的：\BibleQuote{你暂且容许吧，因为这样应当，以满全所有正义}\BibleRef{玛 3:15}；如果一个人更进一步困扰，祂又会对他说祂对犹太人所说的：\BibleQuote{但我并不接受来自人的见证}\BibleRef{若 5:34}。如果现在祂不需要这个见证，为什么若翰由天主派遣？并不是因为祂需要他的见证——这是极端的亵渎。那么为什么？若翰自己告诉我们，当他说：\par

\BibleQuote{好使众人藉着他而相信。}\par

基督也已经说了\BibleQuote{我不接受来自人的见证}\BibleRef{若 5:34}，为了不让愚昧的人以为祂自相矛盾——祂曾在别时说\BibleQuote{另有别位为我作证，我也知道他的见证是真实的}\BibleRef{若 5:32}（祂是指若翰）；在另一处又说\BibleQuote{我不接受来自人的见证}\BibleRef{若 5:34}；祂立刻加上疑团的解答：\BibleQuote{但我这样说是为了你们，好使你们得救}\BibleRef{若 5:34}。仿佛祂曾说：\Quote{我是天主，是真正受生的圣子，是那单纯、有福的本质，我不需要任何人作证；即使没有人作证，我的本质也不因此有丝毫减损；但因为我对许多人的得救关心，我屈尊俯就到如此程度，以致将为我作证的职责交给了一个人。}因为犹太人卑劣和软弱的性情，对祂的信仰以此方式更容易被接受和更合口味。所以正如祂穿上血肉之躯，免得祂以未遮掩的天主性与人相遇而毁灭他们；同样，祂派遣了一个人作祂的使者，使听众听到同类的声音时更易接近。因为（证明）祂不需要那（使者的）见证，只要祂在未遮掩的本质中显示祂是谁，就足以迷惑他们所有。但祂没有这样做，原因我已经提到。祂会消灭所有人，因为无人能承受那不可接近之光的相遇。因此，如同我所说，祂穿上血肉，将（对自己的）见证交托给我们的一位同仆，因为祂为人的得救安排一切，不仅顾及祂自己的光荣，也顾及听众易于接受和有益的事。祂在说\BibleQuote{我说这些话是为了你们，好使你们得救}\BibleRef{若 5:34}时暗示了这一点。福音作者使用同样的言语，在说了\Quote{为那光作证}之后，又说：\par

\BibleQuote{好使众人藉着祂而相信。}\par

% structure: p0010 source=h2 target=subsection reason=relative-heading-rank; base=h2

\subsection{若望福音 1:8}

如果他没有这样说是为了驳斥这种猜疑，那么这句话就完全是多余的，是赘言而非对其教导的阐明。因为，在说了他\Quote{被派来为那光作证}之后，为什么他又说\BibleQuote{他不是那光}？他这样说并非随意或无故；而是因为，在我们当中，通常作证的人被认为比被见证者更伟大、更可信；因此，为了不让任何人在这方面猜疑若翰，他从一开始就立即消除了这种恶意的猜疑，并将其连根拔起，表明这位作证者是谁，以及那位被见证者是谁，还有被见证者与作证者之间的差距有多大。做完这些并显明祂无与伦比的卓越之后，他随后无畏地继续叙述其余部分；并且仔细地清除了那些可能在较单纯者心中秘密滋生的奇怪想法，从而轻松无碍地将教义之言的恰当顺序灌输给所有人。\par

那么，让我们祈祷，从今以后，借着这些思想的启示和教义的正直，我们也能拥有纯洁的生活和光明的品行，因为如果没有善工伴随着我们，这些事便毫无益处。因为我们纵然有全部的信德和全部的圣经知识，但如果我们赤身露体，缺乏由（圣洁的）生活而来的保护，便没有什么能阻止我们被投入地狱之火，永远在无法熄灭的火焰中焚烧。因为正如行善的人复活进入永生，同样，那些胆敢行恶的人也要复活进入永罚，这永罚永无穷尽。因此，让我们表现出全部的热忱，不要因我们行为的卑劣而破坏我们从正信中所获得的益处，而是也要通过这些行为使祂喜悦，从而有胆量仰望基督。没有比这更大的幸福了。但愿我们所有人都能获得上述恩赐，并为了天主的荣耀而行一切；愿光荣归于祂，及独生子与圣神，从今世直到永远。阿们。\par

\SourceRecord{Translated by Charles Marriott.；Nicene and Post-Nicene Fathers, First Series；Vol. 14.；Edited by Philip Schaff.；Buffalo, NY: Christian Literature Publishing Co.,；1889.；<http://www.newadvent.org/fathers/240106.htm>.}

% structure: p0001 source=h1 target=section reason=page-title

\section{\WorkTitle{若望福音讲道集 第七篇}}

% structure: p0002 source=h2 target=subsection reason=relative-heading-rank; base=h2

\subsection{若 1:9}

1. 诸位极蒙爱戴的子女，我们之所以将取自圣经的思想逐段供你们享用，而不一下子全盘倾泻给你们，是为了使你们易于记住先后呈现的内容。因为即使在建筑中，若有人在第一块石头尚未稳固之前就堆砌其他石头，他所建造的墙壁便会全然朽坏，且容易倒塌；而若有人等待灰泥先变硬，再一点一点地添加其余部分，他便能坚固地完成整座房屋，使其坚固，不致短期便损坏。我们效法这些建造者，用同样的方式建造你们的灵魂。因为我们担心，若初奠基尚未稳固，便加上后续的思考，可能因头脑无法同时容纳全部而损害前者。\par

那么，今日读给我们的是什么？\par

\BibleQuote{那才是真光，照亮那来到世上的每一个人。}因为上文提到若翰时说，他来\BibleQuote{为真光作证}；且他在我们这些日子被派遣；以免有人听到此事后，因见证者近期的来临而对所见证者产生类似的怀疑，他便将这想象提升，转移到那在一切起始之前、无始无终的存在。\par

\Quote{怎么可能？}有人说，\Quote{祂既然是圣子，怎能拥有此性体？}我们在谈论天主，你却在问怎么可能？难道你不畏惧、不战栗吗？如果有人问你：\Quote{我们的灵魂和身体如何在来世拥有永生？}你会嘲笑这问题，因为这并非人的理智所能探究，他只应相信，而不应过度好奇于所提及的话题，因为说话者的权能已提供了充分的证据。若我们说，那创造我们灵魂与身体、无比超越一切受造物的祂，是无始的，你是否会要求我们说\Quote{如何？}谁能断言这是有序灵魂或健全理智的行为？你已听到\BibleQuote{那才是真光}，为何你徒劳而鲁莽地试图以推理之力超越这无限的生命？你做不到。为何寻求不可寻求的？为何好奇于不可理解的？为何探查不可探查的？注视那阳光的源头吧。你做不到；却也不为自己的软弱而苦恼或不耐烦；在更大的事上，你为何变得如此大胆和鲁莽？雷霆之子若望，那吹响属灵号角的，当他听到圣神所说的\Quote{那}时，便不再追问。而你这不分享他的恩宠、只凭自己可怜理性说话的人，竟妄图超越他的知识尺度？正因如此，你永远无法达到他的知识尺度。因为魔鬼的诡计在于：它把那些服从它的人从天主设定的界限中引开，仿佛要带他们去更伟大的事；但当我们被这些希望引诱，被它逐出天主的恩宠后，它不仅不赐予更多（作为魔鬼，它如何能？），甚至不允许我们返回那曾安全稳固居住的先前处所，反领我们四处游荡，毫无立足之地。它就这样使第一个受造的人被驱逐出乐园的居所。因使人膨胀，妄想获得更大的知识和尊荣，它便把他从他已安全拥有的事物中逐出。因为他不仅没有像魔鬼所许诺的那样成为神，反而落入了死亡的权下；他因盼望更大的知识，不仅没有因吃树上的果子获得更多，反而丢失了他所拥有的不少知识。因为羞耻感和因赤身而想隐藏的欲望临到了他，他在受骗前是超越这一切羞耻的；而正是这看到自己赤身、将来需要衣服遮盖以及许多其他软弱，从此成了他的本性。为使我们不致如此，让我们服从天主，遵守祂的诫命，不忙碌于任何超出这些的事，以免我们从已得到的恩赐中被逐出。那些我们所谈论的人便是如此。因为他们寻求寻找那无始生命的开端，便失去了他们本可保留的。他们没有找到他们所寻求的（这是不可能的），反而从关于独生子的真信仰中堕落了。\par

我们不可挪移祖先所立的永久界碑，而要永远顺从圣神的法则；当我们听到\BibleQuote{那才是真光}时，不要再寻求什么更多。因为超越这句话是不可能的。如果祂的生成像人那样，在生者和受生者之间必然有时间间隔；但既然这生成是以一种不可言说且合乎天主的方式，你就要放弃\Quote{前}和\Quote{后}，因为这些都是时间的称谓，而圣子甚至是一切时间的创造者。\par

2. \Quote{那么}，有人说，\Quote{祂不是圣父，而是兄弟。}请问为何需要这样说？假如我们主张圣父与圣子来自不同的根源，你这样说或许有道理。但若我们避开这种不敬，并说圣父除了无始之外，也是非受生的，而圣子虽然无始，却由圣父受生，那么由此观念引出那种不敬的断言，有何必要？绝无必要。因为祂是光辉；但光辉包含在其本性的概念中，这光辉是那本性的光辉。为此，保禄称祂为光辉，为叫你想象圣父与圣子之间没有间隔。\BibleRef{希 1:3}因此，这一措辞是在宣告这点；但所引用证据的后续部分，纠正了可能困扰小人的错误见解。因为宗徒说：你不要因为听说祂是光辉，就以为祂缺少自有位格；这是不敬的，属于撒贝略主义者和马塞鲁的追随者的疯狂。我们不这样说，而是说祂也在自有位格中。为此，保禄在称祂为\BibleQuote{光辉}之后，又加上祂是\BibleQuote{祂位格的确切形象}\BibleRef{希 1:3}，以显明祂的固有位格，并显明祂属于同一本体，那本体的确切形象也就是祂。因为正如我以前所说，仅用一个词句不足以向人陈述关于天主的道理，而需要汇集许多词句，从每句中选取合适的。这样，我们才能以与我们能力相称的方式，对祂的光荣作相称的讲述——我是说，就我们的能力而言；因为若有人以为自己能说出与祂本质相称的话，并以此自许，说他知道天主如同天主知道自己一样，那正是最不认识天主的人。\par

因此，知道了这事，让我们持守\BibleQuote{他们传授给我们的，那些从起初亲眼见过并成为圣言服务者的人}\BibleRef{路 1:2}，且不要好奇过头；因为染上这种病（好奇心）的人将遭受两种恶果：一是徒劳地寻找不可能找到的东西而自寻烦恼，二是试图推翻天主所立的界限而激怒天主。这激起了怎样的愤怒，你们知道的人无须再听我们讲。因此，让我们避开他们的疯狂，战栗地聆听祂的话语，好使祂不断建筑我们。因为\BibleQuote{我看谁呢？}\BibleRef{依 66:2}，祂说，\BibleQuote{只看那谦卑、安静、敬畏我话语的人。}让我们抛弃这有害的好奇心，让我们痛彻心扉，让我们为我们的罪哀哭，正如基督所命令的，让我们为我们的过犯而心碎，让我们精确地数算过去我们曾胆敢犯下的一切恶行，并努力以各种方式将它们洗刷干净。\par

为此，天主为我们开启了许多途径。因为，\BibleQuote{你先说明你的罪，好使你成义}\BibleRef{依 43:26}；又说，\BibleQuote{我说：我要向祢承认我的过犯，祢就赦免了我心中的不义}\BibleRef{咏 31:5}；因为不断地认罪和追忆罪过，对减轻其程度不无小补。但还有比这更有效的方法：不对任何得罪我们的人怀恨，宽恕所有冒犯我们的人的过犯。你想知道第三种吗？倾听达尼尔：\BibleQuote{藉着施舍赎你的罪，藉着怜悯穷人赎你的不义。}\BibleRef{达 4:27}除此之外，还有另一种：恒心祈祷，坚持不懈地向天主转求。同样，禁食也为我们带来一些，而且不小的安慰和从所犯罪过中的解脱，只要它与善待他人相伴，并熄灭天主怒火的猛烈。\BibleRef{弟前 2:1}因为\BibleQuote{水能熄灭烈火，施舍能赎罪过}\BibleRef{德 3:30}。\par

因此，让我们沿着所有这些途径行进；因为如果我们全心投入这些工作，把时间用在这些事上，我们不仅会洗净过去的过犯，而且会为未来获得极大的益处。因为我们不会容许魔鬼有闲暇以懒惰的生活或有害的好奇心来攻击我们，因为魔鬼正是通过这些和出于这些，见我们悠闲、空闲、不为生活的卓越作任何预备时，将我们引向愚昧的问题和有害的争论。让我们堵住这攻入之路，让我们警醒，让我们清醒，好使我们在短暂的时间内稍作劳苦，在无尽的岁月中获得永恒的美善，藉着我们的主耶稣基督的恩宠和慈爱；愿光荣归于祂及圣父和圣神，从今世直到永远。阿们。\par

\SourceRecord{Translated by Charles Marriott.；Nicene and Post-Nicene Fathers, First Series；Vol. 14.；Edited by Philip Schaff.；Buffalo, NY: Christian Literature Publishing Co.,；1889.；<http://www.newadvent.org/fathers/240107.htm>.}

% structure: p0001 source=h1 target=section reason=page-title

\section{若望福音第八篇讲道}

% structure: p0002 source=h2 target=subsection reason=relative-heading-rank; base=h2

\subsection{若望福音 1:9}

1. 今天再次处理同一段话，并无妨碍，因为之前我们因阐述教义而未能考虑所读的全部内容。如今那些否认祂是真天主的人在哪里？因为在这里祂被称为\Quote{真光}\BibleRef{若 14:6}，在别处又被称作\Quote{真理}和\Quote{生命}。这句话我们将在适当的地方更清楚地讨论；但目前，我们必须在向你们的爱德讲话时，暂且搁下那件事。\par

如果祂\Quote{光照来到世界上的每一个人}，为什么仍有许多人未受光照？因为并非所有人都认识基督的威严。那么祂如何\Quote{光照每一个人}呢？祂光照每个人，尽自己所能。但如果有些人故意闭上心灵的眼睛，不接受那光的光线，他们的黑暗并非来自光的本质，而是来自他们自身的邪恶，他们故意剥夺了自己的恩赐。因为恩宠倾注于所有人，既不避开犹太人，也不避开希腊人，不避开蛮族、西徐亚人、自由人、奴隶、男人、女人、老人、青年，而是平等地接纳所有人，以同等的关怀邀请。那些不愿享受这恩赐的人，应公正地将自己的盲目归咎于自己；因为如果大门向所有人敞开，没有人阻止，而有人因邪恶故意留在外面，他们的毁灭并非由于别人，而只是由于他们自己的邪恶。\par

% structure: p0005 source=h2 target=subsection reason=relative-heading-rank; base=h2

\subsection{若望福音 1:10}

但并不与世界同久长。决无此理。因此他接着说：\Quote{世界是由祂造成的}，从而将你们再次引向对独生者永恒存在的思考。因为听到这宇宙是祂的工程的人，即使非常迟钝，即使是个仇敌，即使是天主荣耀的敌对者，也必会，不论愿意与否，被迫承认造物主先于祂的工程。因此，我总是对撒摩撒塔的保禄的疯狂感到惊奇，他竟敢正视如此明显的真理，并自愿坠入悬崖。因为他的错误并非出于无知，而是出于明知故犯，与犹太人处于同样的境地。因为犹太人看着人，放弃了纯正的信仰，明知祂是天主的独生子，却不承认祂，因为怕被逐出会堂；据说他也是为了取悦某个女人，出卖了自己的救恩。这是一个强大的力量，确实强大，虚浮光荣的暴政；它甚至能使聪明人的眼睛变暗，除非他们保持警醒；因为如果接受礼物能做到这一点，那么这种更强烈的情绪就更足以做到了。因此耶稣对犹太人说：\Quote{你们既然彼此接受光荣，而不寻求唯独来自天主的荣耀，怎能相信呢？}\BibleRef{若 5:44}\par

\Quote{世界却不认识祂。}此处\Quote{世界}指败坏的、紧贴世俗事物的群众，即普通浮躁愚昧之人。因为天主的友人和宠儿都认识祂，甚至早在祂降生肉身之前。关于圣祖，基督亲自提名说：\Quote{你们的父亲亚巴郎曾欢欣喜乐地期望看见我的日子，他看见了，且喜乐了。}\BibleRef{若 8:56} 关于达味，耶稣驳斥犹太人说：\Quote{那么，达味在圣神内怎么称祂为主，说：上主对我主说：坐在我的右边。}\BibleRef{玛 22:43} 在许多地方，祂与他们辩论时，提及梅瑟；宗徒也提及其余的先知。因为伯多禄说：\Quote{从撒慕尔起，及以后所有发言的先知，也都预告了这些日子。}\BibleRef{宗 3:24} 但是祂曾显现给雅各伯和他的父亲及祖父，与他们谈话，并应许赐给他们许多大恩惠，这些祂也成就了。\par

有人说：\Quote{那么，祂自己怎么说：\InnerQuote{许多先知曾愿意看到你们所看见的，却没有看到；愿意听到你们所听到的，却没有听到。}\BibleRef{路 10:24} 难道他们不曾分享对祂的认识吗？}当然他们认识；我将努力从这句圣言本身说明这一点，因为有些人以为这句话剥夺了他们的认识。因为祂说：\Quote{许多曾愿意看见你们所看见的。}所以他们知道祂将从天而来，生活并教导如同祂生活教导的那样；因为如果他们不知道，他们就不可能有渴望，因为没有人能渴望他完全不知道的事物。因此，他们认识人子，并且知道祂将来到人间。那么，他们没听到的是什么？他们不知道的是什么？就是你们现在所看见和所听到的。因为如果他们在过去也听到祂的声音、看到祂，那不是在肉体中，也不是在人间，也不是祂如此亲密地生活、坦率地与他们交谈的时候。事实上，为了表明这一点，祂并不是简单地说\Quote{看见我}，而是\Quote{看见你们所看见的}；也不是\Quote{听见我}，而是\Quote{听见你们所听到的}。所以，即使他们没有目睹祂在肉体中的来临，他们仍然知道它会发生，并渴望它，并且在没有见过祂在肉体的情况下相信了祂。\par

当希腊人向我们提出这样的指控，说：\Quote{那么基督先前做了什么，以至于祂没有眷顾人类？又因什么可能的理由，祂在长期忽视我们之后，终于来协助我们的救恩？}我们将回答说，在此之前祂就在世界上，眷顾祂的工程，并为所有配得的人所认识。但如果你们说，因为那时并非所有人都认识祂，祂只为那些高尚杰出的人所认识，因此祂未被承认；照此说法，你们现在也不会承认祂被敬拜，因为即使现在也不是所有人都认识祂。但正如现在，没有人会因那些不认识祂的人而拒绝相信那些认识祂的人，同样，关于先前时代，我们不可怀疑祂为许多人，甚至为所有那些高尚可敬的人所认识。\par

2. 若有人说：\Quote{为何不是所有人都听从祂？也不是所有人都敬拜祂，而只有义人才如此？}我也要问：为何如今也不是所有人都认识祂？但我为何提到基督呢？因为那时或是现在，并非所有人都认识祂的圣父。因为有些人说，万物皆是偶然运行，而另一些人则将宇宙的眷顾归于邪魔。还有人除祂之外捏造另一个神，有些亵渎地断言，祂是一种敌对的力量，并认为祂的律法是邪恶鬼魔的律法。那么，我们是否因有人如此说而就说祂不是天主呢？我们是否要承认祂是邪恶的呢？因为甚至有人如此亵渎祂。摒弃如此的心神涣散、如此的完全疯狂吧。若我们按疯子的判断来描绘教义，就无法阻止我们自己也陷入最严重的疯狂。没有人会因视力弱的人而断言太阳对眼睛有害，反而会说太阳适合赐予光明，从健康之人作出判断。没有人会称蜂蜜苦，只因它对于病人的味觉如此。那么，有人会因（心智）患病之人的想象而断定天主不存在，或是邪恶的，或是祂有时确实施展眷顾，有时却全然没有吗？谁能说这样的人心神健全，而不是否认他们神志不清、谵妄、完全疯狂呢？\par

\Quote{世界}，他说，\Quote{不认识祂}；但那些世界不配有的人认识祂。在说了那些不认识祂的人之后，他随即指出了他们无知的原因；因为他并没有绝对地说没有人认识祂，而是说\Quote{世界不认识祂}；也就是说，那些似乎被钉在世界上、只思念世俗之事的人。因为基督惯常这样称呼他们；如祂所说：\BibleQuote{圣父啊，世界没有认识祢。}\BibleRef{若 17:25} 那么，世界不仅不认识祂，也不认识祂的圣父，正如我们所说；因为没有什么比紧紧依附于现世之事更能使心智昏暗的了。\par

因此，你们既知道这事，就当远离世界，尽量摆脱世俗之事，因为由此而来的损失不是寻常的事，而是关乎至善的损失。因为紧抓现世生活的人，不可能真正把握天上的事；但热衷于其中之一的人，必然失去另一个。祂说：\BibleQuote{你们不能事奉天主而又事奉玛门}（\BibleRef{玛 6:24}），因为你必须依附一个而憎恨另一个。这一点就连事物的经验本身也大声宣告。例如，那些嘲笑贪财的人，正是按应该的方式爱天主的人；而那些尊崇（玛门的）主权的人，正是对祂爱得最淡薄的人。因为灵魂一旦被贪欲俘虏，就很难轻易拒绝做出或说出任何令天主愤怒的事，因为它成了另一位主人的奴隶，而那位主人的命令恰恰与天主相反。所以，你们终于要清醒过来，振奋精神，想起我们是谁的仆人，让我们只爱祂的国；让我们为过去侍奉玛门的时光哭泣哀号；让我们永远摆脱他那不堪忍受、沉重无比的轭，继续背负基督轻省而容易的轭。因为祂不像玛门那样向我们发号施令。玛门命我们与所有人为敌，而基督却相反，命我们拥抱并爱所有人。前者将我们钉在泥土和砖瓦上（因为黄金就是泥土），甚至不让我们在夜间稍作喘息；后者将我们从这过分且无谓的忧虑中释放出来，命我们积攒财宝在天上，不是靠对别人的不义，而是靠我们自己的义德。前者在我们因他的律法而受罚、受苦时，不能在我们多番劳苦之后帮助我们，反而增加火焰；后者即使只命令我们给一杯凉水，也绝不让我们失去为此的赏报和酬劳，反而以极大的丰盛回报我们。那么，轻视如此温和、充满一切美善的统治，去侍奉一个忘恩负义、既不能在今生也不能在来世帮助那些顺从和听从他的暴君，岂不是极端的愚昧吗？这不只是可怕的事，也不只是惩罚，即当他受罚时不能保护他们；而且，除此之外，正如我先前所说，他还使顺从他的人陷入千万种灾祸。因为人们可以看到，在那地方受罚的人，多半是因贪财、爱黄金、不肯帮助有需要的人而受罚。为了避免这种情况，让我们施舍，周济穷人，救我们的灵魂脱离今生有害的忧虑，以及因此为我们在那地方预备的惩罚。让我们在天上储存义德。不要积攒地上的财富，而要收集不可攻破的宝藏，那些可以伴随我们上天堂、帮助我们度过危险、在那时刻使审判者仁慈的宝藏。愿我们所有人都能从祂那里获得恩宠，无论是现在还是在那一天，并以极大的信心享受为那些按应该的方式爱祂的人在天上预备的美善，藉着我们的主耶稣基督的恩宠和慈爱，祂与圣父及圣神，是荣耀，自今至永远，及世之世。阿们。\par

\SourceRecord{Translated by Charles Marriott.；Nicene and Post-Nicene Fathers, First Series；Vol. 14.；Edited by Philip Schaff.；Buffalo, NY: Christian Literature Publishing Co.,；1889.；<http://www.newadvent.org/fathers/240108.htm>.}

% structure: p0001 source=h1 target=section reason=page-title

\section{若望福音讲道第九篇}

% structure: p0002 source=h2 target=subsection reason=relative-heading-rank; base=h2

\subsection{若 1:11}

如果你们还记得我们从前所讲的，我们就会更加热心地继续建造余下的部分，当作是极大的收获。因为这样，我们的话语对于你们这些记得已讲内容的人会更加明白，我们也无需太多劳力，因为你们因着热爱学习，能更清楚地看清余下的部分。一个总是丢失所给予之物的人，将永远需要一位老师，永远什么也不会知道；但一个持守他所领受的，并因此额外接受余下部分的人，很快就会从学习者变成老师，不仅对自己有益，对所有人也有益；正如我从你们极其渴望听讲所猜测的，这个集会将会特别如此。那么，来吧，让我们把主的钱财存放在你们的灵魂中，如同在安全的宝库里，并按照圣神的恩宠所能赐予我们的力量，展开今日摆在面前的话语。\par

他（圣若望）在谈到古代时曾说，\BibleQuote{世界不认识祂}\BibleRef{若 1:10}；随后他在叙述中来到宣告（福音）的时代，并说：\BibleQuote{祂来到自己的地方，自己的人却没有接纳祂}\BibleRef{若 1:11}，此处称犹太人为\Quote{自己的人}，作为祂特选的子民，或者甚至全人类，因为都是祂所造的。正如上文，当对众人的愚昧感到困惑，并为我们的共同本性感到羞愧时，他说了\Quote{世界是藉着祂造的}，并且被造后却不认识自己的造物主；同样，这里他因犹太人和众人的愚蠢而无法忍受，便以更显著的方式陈述了指控，说\Quote{自己的人却没有接纳祂}，而且\Quote{祂来到他们那里}时也是如此。不仅是他，先知们也同样惊异，说了完全相同的话，后来保禄也为同样的事感到惊奇。先知们曾以基督的身份高声呼喊说：\BibleQuote{不认识我的民族，竟事奉了我；他们一听见我，就服从了我；外方的子民对我撒谎。外方的子民衰老了，在他们的道路上跛行了。}\BibleRef{咏 17:43} 又说：\BibleQuote{那些未曾被告知有关祂的事的人，将要看见；那些未曾听见的，将要明白。}\BibleRef{依 52:15} 和：\BibleQuote{我让那些不寻找我的人找到我}\BibleRef{依 52:15}；\BibleQuote{我向那些不求问我的人显现}\BibleRef{依 45:1}。保禄在致罗马人的书信中曾说：\BibleQuote{那么，以色列人没有获得他们所追求的，但选民却获得了。}\BibleRef{罗 11:7} 又说：\BibleQuote{我们可说什么呢？那未追求义的外邦人，反倒得了义；但以色列人追求义的法律，却没有达到义的法律。}\BibleRef{罗 9:30}\par

这实在是值得我们惊奇的事：他们是在先知书的知识中长大的，每天听见梅瑟告诉他们无数关于基督来临的事，以及后来的其他先知，而且他们亲眼看见基督每日在他们中间行奇迹，专为他们花时间，还不允许祂的门徒们进入外邦人的道路，或进入撒玛黎雅人的城，祂自己也不这样做，反而处处声明祂被派遣到以色列家迷失的羊那里去\BibleRef{玛 10:5}；（我说）他们既看见了征兆，又听见了先知，还有基督亲自不断提醒他们，他们却使自己变成如此盲目和迟钝，以至于没有一件这样的事能使他们相信基督\BibleRef{玛 15:24}。相反，外邦人从未享受过这些，从未听过天主的圣言，可说连梦中都未曾有过，而是永远徘徊在疯人的寓言中（因为外教哲学即是如此），手中常持着他们的诗人的愚蠢，被拴在木石上，无论在教义还是行为上都没有任何良善或健全之处。（因为他们的生活方式比他们的教义更不洁、更可咒。这是很自然的；当他们看见自己的神祇以一切邪恶为乐，以可耻的话语和更可耻的行为受敬拜，且以此为庆典和赞美，还以污秽的谋杀和杀婴为尊荣，他们又怎能不效法这些呢？）尽管如此，他们虽然堕落到邪恶的极深处，却突然如同藉着某种机械一般，从高天之上照耀我们，从天上最高之处显现出来。\par

那么，这事如何并因何发生？请听\Person{保禄}告诉你们。因为那位有福之人精确地探究了这些事，直到寻得原因并向众人宣告为止。那么原因是什么？这种盲目如何临到犹太人身上？请听那位受托管理这职责的人如何宣告。那么，他在解除许多人的疑惑时说什么呢？\BibleRef{格前 9:17} 他说：\BibleQuote{因为他们不认识天主的正义，想确立自己的正义，就不服从天主的正义。}\BibleRef{罗 10:3} 因此他们遭受了这事。再者，他用别的说法解释同一件事时说：\BibleQuote{那么，我们说什么呢？那未追求正义的外邦人却得到了正义，即由于信德而来的正义；但以色列人追求正义的法律，却未达到这正义的法律。为什么？因为他们不是出于信德去寻求，反而因那绊脚石而绊倒了。}\BibleRef{罗 9:30} 他的意思是：\Quote{这些人的不信是他们不幸的原因，而他们的傲慢是不信的根源。}因为当他们在\Person{基督}来临之前，因领受法律、认识天主以及\Person{保禄}枚举的其他特权，比外邦人享有更大的恩惠时，在\Person{基督}来临之后，他们看见外邦人和自己借着信德被同样召唤，并且信德之后，受割礼者毫不比外邦人更受优待，于是他们心生嫉妒，被傲慢刺痛，无法承受主那不可言喻且超乎寻常的慈爱。所以，这事临到他们，无非是出于骄傲、邪恶和无情。\par

2. 因为，最愚蠢的人啊，关心别人怎样伤害了你？因他人共享同样的事物，你的福乐就减少了吗？但实在，邪恶是盲目的，不能轻易看出任何该看的事。因此，因预见他人共享同样的信心而刺痛，他们便用剑刺向自己，将自己逐出天主的慈爱。这合情合理。因为他说：\BibleQuote{朋友，我没有亏待你，我要给\InnerQuote{这些人}也如同给你一样。}\BibleRef{玛 20:14} 或者说，这些犹太人甚至不配听到这些话。因为比喻中的人，即使不满意，尚可提到一整天劳苦、炎热和汗流浃背。但这些人能说什么呢？毫无此类，而是懒惰、放荡和成千上万恶事，所有先知不断以此指责他们，且他们如外邦人一样因这些事得罪了天主。\Person{保禄}宣告这一点时说：\BibleQuote{因为犹太人与希腊人之间没有分别：所有人都犯了罪，都缺少天主的荣耀，却因他的恩宠白白地成了义。}\BibleRef{罗 10:12} 他在那封书信中很中肯且非常智慧地论述了这一点。但在书信的前面部分，他证明他们应受更大的惩罚。\BibleQuote{因为凡在律法之下犯了罪的，将按律法受审判} \BibleRef{罗 2:12}；也就是说，更严厉，因为除了本性之外，还有律法作为控告者。不仅如此，他们还使天主在外邦人中被亵渎：他说：\BibleQuote{我的名因你们而在外邦人中受亵渎。}\BibleRef{罗 2:24}\par

由于这是最刺痛他们的事（因为即使受割礼而信主的人也认为这事难以置信，因此他们责问\Person{伯多禄}，当他从\Place{凯撒勒雅}来到他们那里时，说他\BibleQuote{进了未受割礼的人的家，且与他们一同吃饭}\BibleRef{宗 11:3}；此后他们虽然明白了天主的安排，仍然惊奇\BibleQuote{圣神的恩赐也倾注在外邦人身上}\BibleRef{宗 10:45}：以他们的惊异表明他们从未料到如此难以置信的事），既然他知道这一点最触动他们，请看他是如何挫败他们的骄傲，放松他们高度膨胀的傲慢。因为他在论述了外邦人的情况，表明他们没有任何借口或得救的希望，并明确指控他们教义的扭曲和生活的污秽之后，便将论证转向犹太人；在列举了先知所有的话，其中说到他们是污秽、奸诈、虚伪的人，\BibleQuote{全都成了无益的}，他们中\BibleQuote{没有寻求天主的}，他们\BibleQuote{全都偏离了正道}\BibleRef{罗 3:12} 等等之后，他补充说：\BibleQuote{我们知道，凡律法所说的，都是对在律法之下的人说的：好叫每一个口都被堵住，全世界都成为在天主面前有罪的。}\BibleRef{罗 3:19} \BibleQuote{因为所有人都犯了罪，都缺少天主的荣耀。}\BibleRef{罗 3:23}\par

那么，犹太人啊，你为何自高自大？为何心高气傲？因为你的口也被堵住了，你的胆量也被夺去了，你和全世界一同成为有罪，像其他人一样，需要白白地成义。即使你曾站得正直，在天主面前大有胆量，你也不应嫉妒那些应受怜悯并通过祂的慈爱得救的人。这是极度的邪恶：为他人的得益而烦恼。尤其当这不会给你带来任何损失时。如果他人的救恩确实损害了你的益处，你的悲伤或许情有可原；尽管即使如此，对于一个学会真正智慧的人来说，也不应如此。但你的赏报并不因他人的惩罚而增加，也不因他人的福祉而减少，你为何因他人白白得救而自叹自怜呢？如我所说，即使你是被认可的人，你也不应痛苦于恩宠临到外邦人。然而，当你自己在主面前与他们犯同样的罪，并且自己也冒犯了主，却对他人的好处不悦，并且自高自大，仿佛只有你才该分享恩宠时，你不仅犯了嫉妒和傲慢的罪，而且极其愚蠢，并可能遭受一切最严厉的惩罚；因为你已在内心种下万恶之根——骄傲。\par

因此，智者说：\BibleQuote{骄傲是罪恶的开端}\BibleRef{德 10:13}；也就是说，它是根、源、母。第一个受造者因此被逐出那幸福的居所；欺骗他的魔鬼也因此从尊严的高位坠落；那个受诅咒者知道这种罪的性质足以使人从天庭坠落，所以当他试图使\Person{亚当}从如此高的尊荣中堕落时，便由此下手。因为他用承诺使\Person{亚当}自高自大，说他将如同神一样，就这样他击垮了他，并将他投入地狱的最深处。因为没有什么比骄傲的暴政更能使人远离天主的慈爱，并将他们交付于深渊之火。因为当骄傲在我们身上时，即使我们实践节制、贞洁、斋戒、祈祷、施舍等一切善行，我们的整个生活也会变得不洁。正如智者说：\BibleQuote{凡心里骄傲的，为天主所憎恶。}\BibleRef{箴 16:5}让我们抑制灵魂的膨胀，铲除这骄傲的肿块，如果我们希望洁净自己，并逃避为魔鬼预备的惩罚。因为骄傲者必与那邪恶者同受惩罚，请听\Person{保禄}宣告：\Quote{不可选初信的人，恐怕他自高自大，就落在审判和魔鬼的网罗里。}什么是\Quote{审判}？他指的是同一\Quote{判决}，同一惩罚。那么，他说人如何能避免这可怕之事？藉着反省自己的本性、自己罪恶的数目、那地方刑罚的严重性、以及今世看似光明的事物之短暂性——它们与草无异，比春天的花朵更易凋谢。如果我们不断在心中激起这些思虑，并记念那些行走最正直的人，那么魔鬼即使千方百计，也不能使我们自高，甚至不能绊倒我们。愿那位谦卑人的天主，良善仁慈的天主，赐给你和我一颗破碎而谦卑的心，使我们能容易地妥善安排其余之事，以光荣我们的主耶稣基督，借着祂并同着祂，偕同圣父及圣神，获享光荣于无穷世。阿们。\par

\SourceRecord{Translated by Charles Marriott.；Nicene and Post-Nicene Fathers, First Series；Vol. 14.；Edited by Philip Schaff.；Buffalo, NY: Christian Literature Publishing Co.,；1889.；<http://www.newadvent.org/fathers/240109.htm>.}

% structure: p0001 source=h1 target=section reason=page-title

\section{若望福音讲道集第十篇}

% structure: p0002 source=h2 target=subsection reason=relative-heading-rank; base=h2

\subsection{\BibleRef{若 1:11}}

1. 亲爱的各位，天主爱人且施恩惠，祂做一切并计划一切，为使我们在德行上发光，并渴望我们得祂嘉许。为此，祂不强迫或勉强任何人，而是以劝诱和恩惠吸引所有愿意的人，将他们赢归自己。因此，当祂来时，有些人接待了祂，有些人不接待。因为祂不愿意有不愿、勉强的仆人，而是要所有人出于自己的意愿和选择，并因服事而感恩于祂。人因需要仆人的服务，依所有权法，甚至违背他们的意愿将许多人保持在那种状态；但天主却一无所缺，不须我们任何事物，只做一切为我们的救恩，使我们在这一事上完全自主，因此祂不强迫或强迫任何不愿的人。因为祂只注视我们的益处；被迫服事这样的主，与不服事无异。\par

有人说：\Quote{那么，为什么祂惩罚那些不愿听从祂的人，又为什么祂威胁那些不服从祂命令的人要下地狱？}因为，极为良善的祂甚至关心那些不听从祂的人，也不离开那些退缩逃避祂的人。但是当我们拒绝了祂恩惠的第一种方式，拒绝走上劝诱和善待的道路时，祂就将另一种方式，即改正和惩罚的方式施行于我们；这方式确实极其痛苦，但在前者被忽视时仍然必要。现在的立法者也为违法者设立了许多严厉的刑罚，然而我们并不因此厌恶他们；我们甚至因他们所制定的刑罚而更尊敬他们，因为他们虽然不需要我们任何东西，又常常不知道谁应享受他们书面法律所提供的帮助，却仍关心我们生活的良好秩序，奖赏那些生活有德的人，用刑罚制止那些放纵者和破坏他人安宁的人。如果我们钦佩和爱戴这些人，难道不更应为天主如此伟大的关怀而惊奇和爱慕祂吗？因为他们的眷顾与祂对我们的眷顾之间的差距是无限的。实在，天主良善的富足是说不尽的，超越一切限度。请想：\BibleQuote{祂来到自己的地方}，不是出于祂个人的需要（因为，如我所说，天主性一无所缺），而是为了向祂的子民行善。但即使如此，当祂来到自己人中为他们谋利时，祂自己的人却不接待祂，反而排斥祂，不仅如此，他们甚至把祂赶出葡萄园，杀害了祂。但即使这样，祂也没有将他们摒于悔改之外，反而允许他们，若愿意，在如此邪恶之后，借信德洗去一切过犯，并与那些从未做过此事、而是祂的特别朋友的人同等。我说这番话并非随意或为劝服，真福保禄的全部历史都大声宣告。因为当他，那位在十字架后迫害基督、用那许多只手用石头砸死祂的殉道者斯德望的人，悔改了，谴责了自己以往的罪，奔向他曾迫害的那位时，祂立刻将他列入祂的朋友中，而且是朋友中的首领，任命他为全世界的传令官和导师，而他曾是\BibleQuote{亵渎者、迫害者和施暴者。}\BibleRef{弟前 1:13}正如他因天主的慈爱而欢喜，大声宣扬，并不羞愧，反而将从前他所妄行的事记录在他的著作中，如同刻在石柱上，展示给众人；他认为与其因犹豫展示自己的错误而掩盖天主说不尽的难以形容的慈爱，不如将自己的过去生活公之于众，以便天赐的宏大彰显。因此，他不断谈论他的迫害、他的阴谋、他对教会的攻击，时而说：\BibleQuote{我原不配称为宗徒，因为我迫害过天主的教会。}\BibleRef{格前 15:9}时而说：\BibleQuote{耶稣来到世上拯救罪人，在罪人中我是最厉害的。}\BibleRef{弟前 1:15}又说：\BibleQuote{你们曾听说过我从前在犹太教中的生活，怎样过度地迫害天主的教会，并加以摧毁。}\BibleRef{迦 1:13}\par

2. 因为他好似对基督向他所显示的容忍作出某种回报，通过显明他所拯救的是何等憎恨者和仇敌，他以极大的坦率宣布了他最初以全部热诚对基督发动的战争；并且借此向那些对自己的境况绝望的人提出美好的希望。因为他声称，基督接纳了他，为要在他身上首先\BibleQuote{显示祂全部的容忍}\BibleRef{弟前 1:16}，以及祂良善的丰富富足，作为\BibleQuote{后来要相信祂而得永生之人的榜样}。因为那些他们所敢做的事，对于任何宽恕都太过重大，正如福音作者所宣告的，他说：\par

\BibleQuote{祂来到自己的人，自己的人却没有接受祂。}祂充满万有、处处临在，又从何处而来？祂手中掌管并托住万有，又空出何处离开祂的临在？祂并非从一个地方转移到另一个地方；祂怎能如此？但祂借着下降至我们中间成就了此事。因为祂虽然已在世界，却似乎不在那里，因为尚未被人认识，但后来祂屈尊取了我们的肉躯而显现自己，他（\Person{圣若望}）称这显现与下降为\Quote{来到}。有人或许会惊奇这位门徒不以师傅的羞辱为耻，反而记下了人们对师傅所行的侮慢；然而这恰恰是他热爱真理品格的明证。此外，感到羞愧的人应该为那些施加侮辱的人羞愧，而不是为受侮辱者。实际上，正因如此，祂更显光辉，甚至在受辱之后仍然如此关怀那些施加侮辱的人；而那些人在众人眼中却显得忘恩负义、该受诅咒，因为他们拒绝了那前来带给他们如此巨大福益者，仿佛视祂为仇敌。他们不仅因此受害，而且也未得到那些接受祂的人所得到的。那些人得到了什么呢？\par

% structure: p0007 source=h2 target=subsection reason=relative-heading-rank; base=h2

\subsection{\BibleRef{若 1:12}}

\Quote{为何，蒙福者，你也不告诉我们那些不接受祂之人的惩罚？你已说他们\InnerQuote{是祂自己的人}，并且当\InnerQuote{祂来到自己的人时，他们却没有接受祂}；但他们会为此遭受什么、受什么惩罚，你却没有继续补充。然而，你若如此威胁，岂不更能使他们惧怕，并缓和他们疯狂的刚硬？为何你沉默不语？} \Quote{还有什么惩罚，}他会说，\Quote{比这更大呢：当他们有机会成为天主的子女时，他们却不这样做，反而自愿剥夺自己如此高贵和尊荣？}虽然他们的惩罚不仅止于此，即他们得不到福益，此外还有不灭的火等待他们，正如他随后更清楚地揭示的。但目前他讲述那些接受祂者不可言喻的福益，并向我们简要陈述这些话，说：\BibleQuote{凡接受祂的，祂赐给他们权能成为天主的子女。}无论是为奴的或是自由的，无论是希腊人、蛮族、西徐亚人，无论无知或有学识，无论男女，儿童或老人，无论受尊敬或卑贱，富足或贫穷，为首领或平民，祂说，所有人都被视作配得同一特权；因为信德与圣神的恩宠，消除了世俗事物造成的不平等，将所有人塑造为同一模样，并印上同一印记，即君王的印记。有什么能与这慈爱相比？一位与我们同出于泥土的君王，并不屈尊将同样具有相同本性的同仆（且无论在品性上常常比他自己更好）登记在王军中，如果他们碰巧是奴隶；但天主的独生子却不屑于把税吏、巫术师和奴隶，甚至比他们更卑微、更贫穷的人、身体残缺、遭受万般痛苦的人，都算作祂子女的行列。如此就是信德于祂的能力，如此就是祂恩宠的丰盈。正如火元素遇到矿石，立刻将泥土变为金子，同样，而且更胜一筹，圣洗圣事使受洗者变为金子而非泥土；圣神那时像火一样落入我们的灵魂，烧毁\BibleQuote{属土的肖像}\BibleRef{格前 15:49}，并产生\BibleQuote{属天的肖像}，崭新、明亮、熠熠生辉，如同出自熔炉之模。\par

那么，他为何不说\Quote{祂使他们成为天主的子女}，而说\Quote{祂赐给他们权能成为天主的子女}？为要表明我们需要极大的热忱，才能在圣洗中印在我们身上的儿子名分肖像始终保持毫无玷污或污秽；同时也表明，除非我们首先自行剥夺，否则无人能从这个权能中夺走我们。因为在人类中，那些接受某件事物绝对控制权的人，几乎与委托他们的人拥有同等的权能；更何况我们，从天主获得了如此荣耀，若我们不做任何与此权能不相称的事，我们将比所有人都坚强；因为将那尊荣交在我们手中的那一位，比万有更伟大、更美好。同时，他也愿意显示，恩宠并非无区别地降临在人身上，而是降临在那些渴望并为之努力的人身上。因为成为（祂的）子女在于这些人的权能；如果他们自己首先不作出选择，恩赐就不会降临到他们身上，也不会产生任何效果。\par

3. 因此，既然处处排除了强制，并指出了人的自愿选择和自由权能，他现在也说了同样的话。因为即使在神妙的祝福中，一方面是天主赐予恩宠，另一方面是人提供信德；在后来，为了保持纯洁，我们需要极大的热忱。仅仅受洗并相信，并不足以使我们保持纯洁，我们还必须，若我们愿意持续享有这光明，活出与此相称的生活。这便是天主在我们内的工作。以奥秘的诞生出生，并从所有以前的罪恶中洁净，来自圣洗圣事；但将来保持纯洁，此后不再沾染任何污秽，则属于我们自己的权能和勤奋。正是因此，他提醒我们这诞生的方式，并通过与肉身的痛苦相比较，显示其卓越，当他这样说时：\par

% structure: p0011 source=h2 target=subsection reason=relative-heading-rank; base=h2

\subsection{\BibleRef{若 1:13}}

他这样做，是为了让我们思量第一次诞生的卑贱低下，即\Quote{由血}和\Quote{由肉欲的意志}，同时察觉第二次诞生（由恩宠而来）的崇高尊贵，从而对其形成某种伟大的看法，一种与那生了我们者的恩赐相称的看法，并在将来表现出极大的热忱。\par

因为并非没有重大的恐惧，只怕我们日后因懈怠和过犯玷污了那美丽的礼服，就会被逐出内室和洞房，如同那五个愚拙的童女，或那没有穿婚宴礼服的人一样。\newline{}\EditorialNote{参见：\BibleRef{玛 25}；\BibleRef{玛 22}}。他也是宾客之一，因为他曾被邀请；但因为在邀请和如此大的尊荣之后，他对那邀请他的人傲慢无礼，请听他所受的惩罚，何等可怜，堪当许多眼泪。因为当他来分享那丰盛的筵席时，不仅被禁止享用最小的，反而手脚都被捆绑，被带到外面的黑暗里，去承受永恒无尽哀号和切齿的痛苦。因此，亲爱的诸位，我们也不要以为信德就足以使我们得救；因为如果我们不显示出纯洁的生活，反而穿着与这有福的召叫不相称的衣服而来，那么没有什么能阻止我们遭受那可怜之人的同样命运。令人惊奇的是，祂，既是天主又是君王，不以卑贱、乞讨、无名之辈为耻，甚至从十字路口将他们领到那筵席；而我们却如此麻木，连这么大的尊荣也不能使我们变好，甚至在蒙召之后仍停留在昔日的邪恶中，傲慢地滥用那召叫我们者的不可言喻的慈爱。因为祂召叫我们参加祂奥迹的灵性与可畏的共融，并非为了让我们带着昔日的邪恶进入；而是为了让我们脱去污秽，换上适合那些在宫殿中受款待者的衣袍。但如果我们行事与那召叫不相称，这不再在于那尊荣我们者，而在于我们自己；不是祂将我们从那卓越的宾客团体中逐出，而是我们逐出自己。\par

祂已完成了祂的一切部分。祂预备了婚宴，祂摆设了筵席，祂派人来召叫我们，我们来到时祂接纳了我们，并以一切其他尊荣待我们；但我们，当我们以污秽的衣袍——即我们不洁的行为——冒犯了祂、宾客和婚宴时，就理当被逐出。这是为了尊荣婚宴和宾客，祂才驱赶那些鲁莽无耻的人；因为若祂容忍那穿着如此衣袍的人，就显得是在侮辱其余的人。但愿我们中，或其他人中，绝不要从召叫我们者那里得到这种待遇！因为这一切事在发生之前就被记录下来，好让我们因圣经中的威胁而醒悟，不使这种羞辱和惩罚化为实际行动，而仅止于言语，并各人带着光明的衣袍前来赴那召叫；但愿我们都能享受，借着我们的主耶稣基督的恩宠和慈爱，藉着祂并同祂，与圣父和圣神，永受光荣于无穷世。阿们。\par

\SourceRecord{Translated by Charles Marriott.；Nicene and Post-Nicene Fathers, First Series；Vol. 14.；Edited by Philip Schaff.；Buffalo, NY: Christian Literature Publishing Co.,；1889.；<http://www.newadvent.org/fathers/240110.htm>.}

% structure: p0001 source=h1 target=section reason=page-title

\section{若望福音讲道集第十一篇}

% structure: p0002 source=h2 target=subsection reason=relative-heading-rank; base=h2

\subsection{\BibleRef{若 1:14}}

我愿在开始讲解福音经文之前，向诸位求一个恩惠，你们不要拒绝我的请求，因为我所求的并不沉重或烦难，而且倘若应允，它不仅有益于我这个接受者，也更有益于你们这些施予者，甚至对你们更有益处。我所要求的是什么呢？就是你们每人拿起那一部分福音，即要在主日或安息日在你们中间宣读的那一段，在日子来临之前，先在家中坐下通读一遍，并常常仔细思考其内容，好好审察每一部分：哪些是清楚的，哪些是晦涩的，哪些似乎有利于反对者而实际并非如此。总之，当你们尝试了每一点之后，再去听人宣读。因为这样的热忱对你们和我们都大有裨益。我们无需费多大力气就能使所讲的意义明朗，因为你们的心思已熟悉经文的含义，你们不仅会变得更敏锐、更有洞察力去听和学，而且还能教导他人。而现今大多数来这里听道的人，他们一下子要同时领会所有意思，包括经文和我们所作的评论，即使我们这样做一整年，他们也不会有多大收获。怎能呢？他们把听道当作副业，只在此时此地，且只有这短短的时间。若有人将责任推给事务、挂虑和公私事务的不断缠身，首先，这本身就是一个不小的指责：他们为如此多的事务所包围，如此不断地钉在今生之事上，以致不能为那比一切更需要的事情腾出一丁点闲暇。此外，这只是一个借口和托词，他们的朋友聚会、在剧场的闲逛以及看赛马的聚会都证明他们并非真忙，他们往往在那里度过整天，却从未有人抱怨事务繁忙。因此，对于琐事，你们总能毫不推诿地找到充分的闲暇；但到了该关心天主之事时，这些事对你们来说竟如此多余和卑下，以致你们认为连一点闲暇都无需分给它们吗？有这样性情的人，怎配呼吸或仰望这太阳呢？\par

这些懒人还有一个最愚蠢的借口，就是他们手头没有书。对于富人来说，我们若针对这个借口来指责他们，是荒唐的；但我猜想许多穷人也经常用这个借口，我很乐意问一问：难道他们每一个人不都拥有其职业所需的一切工具，虽然极其贫困挡在面前，却仍备齐了吗？那么，在那种情况下，他们不责怪贫困，而想尽办法不使任何方面成为障碍；但在我们可能获得如此巨大益处之时，却哀叹没有闲暇和贫困，岂非怪事？\par

此外，即使有人如此贫穷，他们仍有可能借着此地不断诵读的圣经，不忽略其中任何内容。或者，如果这对你们来说似乎不可能，那也确有理由，因为许多人并未怀着炽热的热情来聆听所讲的话，而只是为我们做了这一件事，随即回家。即使有人留下，他们的心态也不比那些离开的人好多少，因为他们只在身体上与我们同在此处。但为免我们以过多指责来累着你们，并花所有时间来挑剔，让我们继续讲福音的话，因为现在是时候把余下的话导向摆在面前的主题了。所以，请振奋精神，以免所讲的话有任何遗漏。\par

他说：\BibleQuote{圣言成了血肉}，\BibleQuote{住在我们中间}。\par

在宣告了接待祂的人是\BibleQuote{由天主而生}，并成为\BibleQuote{天主的子女}之后，他又补充了这不可言喻的尊荣的原因和理由：那就是\BibleQuote{圣言成了血肉}，主人取了奴仆的形状。因为祂原是天主的独生子，却成了人子，为使众人子成为天主的子女。因为崇高者与卑微者交往，丝毫不损及自身的尊荣，反而将另一个从其极度的卑微中提升起来；主正是如此。祂以此屈尊，丝毫没有减损自己的本性，却将我们这些常处于耻辱和黑暗中的人，提升到不可言喻的光荣。就像一个国王，以关爱和仁慈与卑微的穷人交谈，并不使自己蒙羞，反而使对方受到众人注目和显赫。如果对于人的附加尊严而言，与卑微者交往并不损害更尊贵者，那么对于那单纯而美善的本体——祂毫无附加或生长衰败，一切美善都稳固而永恒——就更不会如此了。因此，当你听见\BibleQuote{圣言成了血肉}时，不要惊惶失措，也不要沮丧。因为那本体并未变成血肉（想象这一点是亵渎），而是继续祂之所是，就这样取了奴仆的形状。\par

2. 那么，他为何使用\Quote{成了}这个说法？为堵住异端者的口。因为有些人说降生奥迹的一切情形都只是幻象、演戏、比喻，他便预先除去他们的亵渎，用了\Quote{成了}一词，以此表明不是本质的改变（切勿有此想法），而是对真正肉身的取用。因为就如（圣保禄）说：\Quote{基督救赎我们脱离了法律的诅咒，为我们成了诅咒}，他并不是说祂的本质离开其固有的荣耀而取了一个被诅咒之物的存在（这事连魔鬼也不能想象，连最愚昧的人或理智丧失者也不能，其中含有如此的不敬与疯狂），正如（圣保禄）不是说这个，而是说祂将加于我们的诅咒加于自身，使我们不再在诅咒之下；同样，这里（圣若望）也说祂\Quote{成了肉身}，不是将祂的本质变为肉身，而是将肉身娶为己有，祂的本质仍保持不变。\par

如果他们说：既然祂是天主，祂是全能者，所以祂能使自己降低到肉身的本体，那么我们要答复他们：祂是全能者，但条件是祂继续保持为天主。但若祂容许改变，即变得更坏，祂怎能是天主呢？因为改变远离那单纯的本性。因此先知说：\Quote{他们都像衣服一样渐渐陈旧，祢要卷起他们如同外套，他们就被更换；但祢永远相同，祢的岁月没有穷尽。}\BibleRef{咏 102:27}因为那本质超乎一切改变。没有比祂更好的东西，祂可以前进和达到。我说更好吗？不，也没有同等，也没有丝毫接近祂的。因此，如果祂改变，祂必须接受更坏的改变；而这就不成其为天主。但让亵渎回归到那些说这话的人头上。不仅如此，为表明他使用\Quote{成了}这一说法，只是要你们不要以为是单纯的幻象，请听他如何从后面的话来澄清论证，并推翻那恶意的猜度。因为他添加了什么？\Quote{并居住在我们中间。}他几乎是在说：\Quote{不要从\InnerQuote{成了}一词中想象什么不当的事；我所说的并非那不变本性的任何改变，而是它的居住和驻留。但居住者不可能与它所居住的相同，而是不同的；一个东西住在不同的东西中，否则就不是居住；因为没有任何东西能居住在自己内。我指的是本质上的不同；因为通过一种联合与结合，天主圣言与肉身成为一体，不是由于本体的混乱或消失，而是藉着一种不可言说、超乎理解的结合。不要问如何；因为它是如此的，正如祂所知道的一样。}\par

那么，祂所居住的帐幕是什么呢？请听先知说：\Quote{我要重建达味已倒塌的帐幕。}\BibleRef{亚 9:11}它确已倒塌，我们的本性遭受了不可治愈的堕落，只缺那只大能的手。若非起初造它的那一位向它伸出祂的手，藉着水与圣神的重生重新印上祂的肖像，它绝无可能再被竖立。并且，我请你们注意，这奥迹可敬而不可言说的性质。祂永远居住在这帐幕内，因为祂披上我们的肉身，并非为了再次脱去，而是为永远拥有它。若非如此，祂就不会认为它配得君王的宝座，也不会在穿戴它时受到天上万军的朝拜——天使、总领天使、上座者、宰制者、掌权者、异能者。什么言语，什么思想，能够表达给予我们人类如此伟大、真正可敬可畏的尊荣？哪一位天使？哪一位总领天使？没有，无论在天上或地上，没有一位。因为天主的作为如此伟大，祂的恩惠如此浩大，以致正确的描述不仅超越人的口舌，甚至超越天使的能力。\par

因此，我们要暂时结束我们的讲论，静默不语；只向你交付这项嘱咐：你要回报我们这位如此伟大的施恩者，以一种再次为我们带来全部益处的回报。这回报就是：我们以全部小心关注我们灵魂的状况。因为这也是祂慈爱的工程：祂无须我们任何东西，却说当我们照顾自己的灵魂时，祂就得了回报。因此，如果我们，当如此尊荣倾注于我们时，竟连我们所能的也不愿贡献，而况这贡献又能使益处回到我们身上，并且在此条件下为我们摆上一万种福分，这乃是极端的愚妄，应受一万次的惩罚。为这一切，让我们将光荣归于我们仁慈的天主，不单用言语，更要用行动，好使我们获得将来的福乐。愿我们众人都能得此福乐，藉着我们的主耶稣基督的恩宠和慈爱。藉着祂，并与祂一同，同圣父及圣神，永受光荣，世世无尽。阿们。\par

\SourceRecord{Translated by Charles Marriott.；Nicene and Post-Nicene Fathers, First Series；Vol. 14.；Edited by Philip Schaff.；Buffalo, NY: Christian Literature Publishing Co.,；1889.；<http://www.newadvent.org/fathers/240111.htm>.}

% structure: p0001 source=h1 target=section reason=page-title

\section{若望福音讲道第十二篇}

% structure: p0002 source=h2 target=subsection reason=relative-heading-rank; base=h2

\subsection{若望福音 1:14}

1. 也许日前我们对你们显得过分严厉和沉重，用语过于尖锐，对我们的谴责过于冗长，针对许多人的懈怠。如果我们这样做只是为了激怒你们，那么你们每个人都有理由生气；但如果是为了你们的益处，我们在言语中忽略了可能取悦你们的东西，你们若不愿赞许我们的远见，至少也应为如此温柔的慈爱而宽恕我们。因为我们实在非常害怕，如果我们努力，而你们不愿在聆听中表现出同样的勤奋，那么你们未来的交账可能更加严厉。因此我们不得不不断唤醒和激励你们，以免所说的任何话从你们眼前溜走。因为这样你们现在就能怀着极大的信心生活，并在那日于基督的审判台前展现出来。既然我们最近已经充分触动了你们，今天就让我们从这些语句本身开始吧。\par

\BibleQuote{我们看见了，}他说，\BibleQuote{祂的光荣，正如父独生子的光荣。}\par

他宣布了我们成为了\BibleQuote{天主的子女}，并表明了方式，即藉着\BibleQuote{圣言成了血肉}，他再次提到我们从同一情况中获得的另一个好处。是什么？\BibleQuote{我们看见了祂的光荣，正如父独生子的光荣}；若没有藉着与我们相似的身体向我们显示，我们就无法看见。因为如果古时的人甚至无法承受观看梅瑟那荣光焕发的面容，他与我们分享同一本性，如果那位义人需要一块面纱来遮蔽他光荣的纯净，向他们显示他们的先知面容温和；那么，我们这些泥土和尘土所造之物，怎能承受那甚至高天之上的能力也无法接近的、未遮盖的天主性呢？所以祂在我们中间支搭帐幕，使我们能够无所畏惧地接近祂，与祂说话，并与祂交谈。\par

但是，\BibleQuote{正如父独生子的光荣}是什么意思？既然许多先知也得了光荣，如这位梅瑟、厄里亚、厄里叟，前者被火战车环绕（见\BibleRef{列下 6:17}），后者被它提升；在他们之后，有达尼尔和三青年，以及许多显示奇迹的其他人；又有天使在人前显现，并部分向观者展示他们本性的闪光；而且，不只是天使，甚至革鲁宾也被先知看见在极大的光荣中，还有色辣芬也是如此：福音作者引导我们离开这一切，将我们的思想从受造物和同仆的光辉中挪开，把我们置于至善的顶峰。因为，他说：\Quote{不是先知，}\Quote{也不是天使，也不是总领天使，也不是更高的权能，也不是任何其他受造的本性}，如果还有别的，而是主本身、君王本身、真正的独生子本身、万有的主本身，我们\BibleQuote{看见了光荣。}\par

因为\Quote{正如}一词在这里不属于相似或比较，而是确认和毋庸置疑的定义；就好像他说：\Quote{我们看见了光荣，祂是独生子和真天主子，万有的君王，祂理应拥有的那种光荣。}这种说话习惯是普遍的，因为我不会拒绝甚至从普通习俗中加强我的论证，因为现在我的目的不是根据言辞的优美或文笔的优雅来说话，而只是为了你们的益处；因此，以普通习俗为例来建立我的论证，毫无阻碍。那么，大多数人的习惯是什么？常常当有人看见一个国王衣着华丽，周身闪耀着宝石，后来向别人描述那美丽、装饰和辉煌时，他们尽量列举：紫袍的鲜艳色彩、宝石的大小、白骡的洁白、轭上的黄金、柔软闪亮的床榻。但是，在列举了这些以及其他事物之后，他们仍然无法用言语充分表达那辉煌，于是立即说：\Quote{但何必多说呢；总而言之，他就像一个国王；}并不是希望用\Quote{像}这个说法来表明他们所说的人像一个国王，而是说他是一个真正的国王。同样，现在福音作者使用了\Quote{正如}一词，意欲表达祂光荣的超然本性和无可比拟的卓越。\par

实在说来，所有其他受造物，无论是天使还是总领天使，以及先知，都如同奉命行事；但祂却以君王和主人相称的权威行事，连群众也感到惊奇，因为祂教导他们\Quote{好像有权威的人}\BibleRef{玛 7:29}。如我所说，天使也曾带着极大的荣耀显现于地上，例如在\Person{达尼尔}、\Person{达味}、\Person{梅瑟}的事迹中，但他们行事都像有主人的仆人。但祂却像万有的主和宰制者，即便在祂显现于贫穷卑微的形像时也是如此，但受造物仍认出了自己的主。如今，从天上召唤贤士来朝拜祂的那颗星，那遍布各处、侍立于主前并颂扬祂的庞大天使队伍，此外还有许多其他的传报者突然出现，所有人相遇时，都互相传报这不可言喻的奥秘的喜讯：天使向牧羊人传报，牧羊人向城里的人传报，\Person{加俾额尔}向\Person{玛利亚}和\Person{依撒伯尔}传报，\Person{亚纳}和\Person{西默盎}向来到圣殿的人传报。不但男人和女人满心欢喜，就连那尚未出世的婴儿——我说的是旷野中的居民、这位圣史的同名者——也在母腹中欢喜跳跃，所有人都对未来充满希望。这一切发生在祂诞生之后不久。但当祂更进一步彰显自己时，又出现了比先前更伟大的奇事。因为不再是星辰或天空，不再是天使或总领天使，不再是\Person{加俾额尔}或\Person{弥额尔}，而是圣父自高天亲自宣告祂，并与圣父同在的护慰者，在发出声音时飞降而下，停留在祂身上。因此祂确实说：\Quote{我们见了祂的光荣，正如父独生者的光荣。}\par

2. 然而，祂这样说不仅因为上述之事，也因随之而来的事；因为如今不再只有牧羊人、寡妇和老人向我们传报喜讯，而是事物本身的声音，比任何号角更响亮，如此洪亮，以致那声音甚至立刻传到这片土地上。因为经上说：\Quote{祂的名声传遍了整个\Place{叙利亚}}\BibleRef{玛 4:24}；祂向所有的人启示了自己，各处的一切都在高呼：天国的君王已经来临。恶神处处逃避祂、远离祂，\Person{撒殚}掩面退却，死亡在那时在祂面前退缩，后来完全消失；各种疾病都被解除，坟墓释放了死者，魔鬼释放了被它们迷惑的人，疾病释放了病人。人们可以看见奇异而奇妙的事，正如先知们理当渴望看见、却没有看见的。人们可以看见眼睛被塑造（\BibleRef{若 9:6}），看见祂在短时间内，在身体较高贵的部分向众人展示了那令人赞叹的事，即所有人都渴望看见的：天主如何用泥土塑造了\Person{亚当}；瘫痪扭曲的肢体被接合、调整，死者的手在动，瘫痪的脚在跳跃，被堵塞的耳朵重新张开，那先前因不能说话而缄默的舌头大声作响。因为祂拿起了人类的共同本性，如同一位优秀的工匠拿起一座被时间毁坏的房屋，祂补上了断裂之处，连接了裂缝和摇动之处，并重新建起了完全倒塌的房屋。\par

至于灵魂的塑造，比身体的塑造更令人赞叹，又该说什么呢？我们身体的健康是伟大的，但灵魂的健康更为伟大，正如灵魂优于身体一样。不仅如此，还因为我们的身体本性随造物主的引导而往，毫无抵抗；但灵魂是自己的主人，拥有对自己行动的能力，除非它愿意，否则不会在一切事上服从天主。因为如果灵魂不情愿、如同被强迫一般，天主不会使它美丽卓越，因为这不是德行；祂必须说服灵魂心甘情愿地如此行。因此，这种医治比前者更加困难；然而连这也成功了，各种邪恶都被驱逐。正如祂所治愈的身体，不仅恢复健康，而且达到最强健的状态，同样，祂不仅把灵魂从极端的邪恶中解救出来，而且使它们达到德行的顶峰。一个税吏成了宗徒，一个迫害者、亵渎者和凌辱者成了向世界的传报者，贤士成了犹太人的导师，一个盗贼被宣布为乐园的公民，一个娼妓因信德伟大而闪耀，还有客纳罕和撒玛黎雅的两个妇人，后者原也是娼妓，竟向同乡宣讲福音，用她的网网住整座城，把他们带到基督面前；前者则因信德和恒心，为她女儿的灵魂驱逐了恶神。还有许多比这些更坏的人，立刻被列入门徒的行列，他们身体的各种软弱和灵魂的疾病都立即被转变，重新塑造为健康和极致的德行。这些事不是两三个人，不是五个人，不是十个人，不是二十人，不是一百人，而是整座城和整个民族都很容易地被重塑。何必再提诫命的智慧、天上法律的卓越、天使般的生活方式的美好秩序呢？因为祂为我们提出了这样的生活，为我们制定了这样的法律，建立了这样的生活方式，使那些付诸实践的人，即使他们比所有人都更坏，也立刻成为天使，并尽可能相似于天主。\par

3. 圣史既汇集了这一切——我们身体和灵魂中的奇事、信仰的基本要素、诫命、那些不可言说且高于诸天的恩赐、法律、政体、劝诫、未来的应许、祂的苦难——便发出了那如此奇妙、充满崇高教义的声音，说：\BibleQuote{我们见了祂的光荣，正如圣父独生者的光荣，满溢恩宠和真理。} 因为我们不仅因奇迹而钦佩祂，也因苦难而钦佩祂；就如祂被钉在十字架上，祂被鞭打，祂被掌掴，祂被吐唾沫，祂被那些祂曾施恩的人打脸。因为就连那些看似可耻的事，也有理由重复同样的说法，因为祂自己把那行动称为\Quote{光荣}。因为当时所发生的，不仅是良善和爱德的明证，也是不可言喻的权能的明证。那时，死亡被废除，诅咒被解除，魔鬼受羞辱、被俘虏、被示众，我们罪债的字据被钉在十字架上。而且，当这些奇迹无形中进行时，其他的奇迹有形地发生了，显示祂确实是天主的独生子、一切受造的主宰。因为当那蒙福的身体还挂在树上时，太阳遮蔽了它的光芒，整个大地震动变暗，坟墓裂开，地面摇动，无数的死人涌出，进入城中。当祂坟墓的石头被固定在穹顶上，封条还在上面时，死者复活了，那被钉的、被刺穿的复活了，祂以祂的大能充满了祂的十一位门徒，派遣他们走遍世界，成为全人类共同的医治者，修正他们的生活方式，将天上教义的知识传播到地极，摧毁魔鬼的暴政，教导那些伟大而不可言说的福乐，向我们传报灵魂不死、身体永生、超越想象且永无穷尽的报酬的喜讯。蒙福的圣史心怀这些事，甚至更多的事，虽然他知晓，却无法写下，因为世界容不下（因为若一一写下，\BibleQuote{我想，所写的书，世界也容不下}——\BibleRef{若 21:25})。因此，他反思这一切，高呼：\BibleQuote{我们见了祂的光荣，正如圣父独生者的光荣，满溢恩宠和真理。}\par

因此，那些被堪当看见并听见如此之事，并享受如此伟大恩赐的人，也应该活出与教义相称的生活，好能享受那里所预备的美善。因为我们的主耶稣基督来临，不仅是为了让我们在此看见祂的光荣，也是为了那将来的光荣。因此祂说：\BibleQuote{我愿我在那里，这些人也同我在那里，为叫他们看见我的光荣。}\BibleRef{若 17:24}。如果这里的光荣如此明亮辉煌，那么将来的光荣又该如何言说呢？因为它将不在这可朽坏的大地上显现，也不在我们这必朽的身体内，而是在不朽、永不陈旧、其光辉无法言喻的受造界中。有福的，三倍有福的，甚至多次有福的，是那些被堪当瞻仰那光荣的人！关于此，先知说：\BibleQuote{不义的人要被除去，免得他看见主的光荣。}\BibleRef{依 26:10}。但愿我们中无人被除去，也永不被排除于瞻仰那光荣之外。因为若我们将来不能享受它，那么说我们\Quote{还不如未曾出生}就为时已晚。因为我们活着呼吸是为了什么？若我们错失那景象，若无人允许我们那时瞻仰我们的主，我们算什么？若那些看不见太阳光的人忍受着比任何死亡更痛苦的生活，那些被剥夺了那光明的人必会遭受什么？因为在前者，损失仅限于这一剥夺；但在后者，却不限于此（尽管如果这是唯一可怕的事，惩罚的程度也不会相等，而是一个比另一个严厉得多，正如那太阳与此太阳无可比拟一样），但现在我们还必须期待其他的报应；因为看不见那光明的人不仅被引入黑暗，还要不断被焚烧，消瘦，咬牙切齿，遭受一万种其他可怕的事。所以，我们不要因在这短暂的时间内疏忽懈怠而落入永罚，而要警醒谨慎，竭尽全力，以一切作为去追求那福乐，远离那在可怕审判宝座前咆哮奔涌的火河。因为一旦被投入其中，就必须永留；无人能将他从惩罚中救出，没有父亲，没有母亲，没有兄弟。先知们自己也大声宣告了这一点：一个说：\BibleQuote{兄弟不能赎回兄弟。人能赎回吗？}\BibleRef{咏 48:7}。而厄则克耳更进一步说：\BibleQuote{纵然诺厄、达尼尔和约伯在其中，也不能救他们的子女。}\BibleRef{则 14:16}。因为只有一种辩护，即通过行为；若被剥夺了它，就无任何其他方法得救。因此，反复思想这些事，并不断反思，让我们洁净自己的生命，使它发光，好使我们能坦然无惧地看见主，并获得所应许的美善；藉着我们的主耶稣基督的恩宠和慈爱，愿光荣归于圣父和圣神，至于无穷世之世。阿们。\par

\SourceRecord{Translated by Charles Marriott.；Nicene and Post-Nicene Fathers, First Series；Vol. 14.；Edited by Philip Schaff.；Buffalo, NY: Christian Literature Publishing Co.,；1889.；<http://www.newadvent.org/fathers/240112.htm>.}

% structure: p0001 source=h1 target=section reason=page-title

\section{\WorkTitle{若望福音讲道集 第十三篇}}

% structure: p0002 source=h2 target=subsection reason=relative-heading-rank; base=h2

\subsection{\BibleRef{若 1:15}}

我们岂是徒然奔跑劳苦？我们在磐石上撒种吗？种子落在磐石上吗？种子落在我们不知的路旁和荆棘中吗？我十分忧虑和害怕，唯恐我们的耕耘没有益处；这并不是说我会和你们一样受到亏损，关乎这劳苦的赏报。因为教导人的和农夫并不相同。农夫常在一年的辛劳、苦工和汗水之后，如果土地没有为他的劳苦产生相应的回报，他就无法从别处得到安慰，只能羞愧沮丧地从谷仓回到家中，回到妻子儿女身边，不能要求任何人为他长时间的辛苦付出报酬。但在我们的情况中，并非如此。即使我们所耕种的土壤不结果实，只要我们显出了一切的勤勉，它和我们的主人就不会让我们失望地离去，反而会赐给我们赏报；因为圣保禄说：\BibleQuote{每人要按自己的劳苦领取自己的赏报}\BibleRef{格前 3:8}，不是按事情的结果。事实确是如此，请听经上说：\BibleQuote{你，他说，人子，向这百姓作证，听不听由他们，明白不明白也由他们。}\BibleRef{则 2:5} 厄则克耳又说：\BibleQuote{如果守望的人警告当逃避什么、当选择什么，他就救了自己的灵魂，即使没有人留意。}\BibleRef{则 3:18} 虽然我们有这坚强的安慰，并且确信为我们预备的赏报，但当我们看到你们身上的工作没有进展时，我们的光景并不比那些哀叹悲恸、掩面羞愧的农夫好多少。这是教师的同情，这是父亲天然的关怀。因为梅瑟也曾如此，那时他有能力脱离犹太人的忘恩负义，并为另一个更伟大的民族奠定更光荣的基础——（天主说：\BibleQuote{你放开我，让我消灭他们，并兴起一个比你更强大的民族。}\BibleRef{出 32:10}）——但因为他是一个圣人，天主的仆人，真诚慷慨的朋友，他连听也不愿听这话，反而宁愿与曾经分配给他人一同灭亡，也不愿没有他们而独自得救并享受更大的荣耀。凡负有灵魂责任的人都应当如此。因为奇怪的是，一个人若拥有软弱的孩子，就不会被称为除自己亲生的孩子之外的任何人的父亲；而手里有门徒的人，却要不断地从一个羊群换到另一个羊群，我们一会儿抓住这些人的管理权，一会儿又抓住那些人的，再一会儿又是其他人的，对任何人都没有真正的爱心。愿我们永远不必怀疑你们中有这等事。我们相信你们在我们的主耶稣基督内，在彼此相爱和向众人施爱上，越发丰盛。我们说这话，是希望你们的热情更加增长，你们品行的卓越更加增进。因为只有这样，你们才能深入到摆在我们面前的圣言的深处，如果没有邪恶的薄膜遮蔽你们理智的眼睛，并搅乱其清晰和敏锐。\par

那么，今天摆在我们面前的是什么？\BibleQuote{若翰为他作证，呼喊说：\InnerQuote{这就是我曾说的：那在我以后来的，成了在我以前的，因他原先在我以前。}} 福音作者频繁提到若翰，并屡次引用他的见证。这不是没有理由的，而是非常明智的；因为所有的犹太人都十分钦佩这个人（甚至约瑟夫斯也将战争归咎于他的死，并指出，因他的缘故，昔日母城如今已不复存在，并且继续长篇大论地赞扬他）。因此，作者想借着他使犹太人羞愧，就不断提醒他们前驱的见证。其他福音作者提到更古老的先知，每当有关基督的事发生时，他们就引导听众去参考他们。例如，当婴孩诞生时，他们说：\BibleQuote{这一切事的发生，是为应验上主借依撒意亚先知所说的话：\InnerQuote{看，一位贞女将怀孕生子。}}\BibleRef{玛 1:22}；当他被谋害并到处被追捕，以至黑落德屠杀了幼小的婴儿时，他们引用了耶肋米亚的话：\BibleQuote{在辣玛听到了声音，痛哭哀号，极大的悲泣，辣黑耳为她的儿女哭泣。}\BibleRef{玛 2:18}；而当他从埃及上来时，他们提到欧瑟亚的话：\BibleQuote{我从埃及召回了我的儿子。}\BibleRef{玛 2:15}；他们处处如此行。但是若翰提供了更清晰、更新鲜的见证，发出了比其他人更荣耀的声音，他不仅不断引出那些已经去世的人，也引出一个活着在场的人，这个人指认并给他付洗，他不断介绍他，不是想借仆人为主人增光，而是迁就听众的软弱。因为正如除非他取了仆人的形体，否则不会轻易被接受；同样，除非他借着仆人的声音预备了众同仆的耳朵，否则犹太人中的大多数人不会接纳圣言。\par

2. 除此之外，还有另一个伟大而奇妙的安排。因为一个人说关于自己的大话会使他的见证受到怀疑，并且常常成为许多听者的障碍，所以另一个人来为他作证。除此之外，众人在某种程度上更倾向于接受更熟悉和自然的声音，因为他们比其他人更容易认出它；因此天上的声音只发出了一两次，而\Person{若翰}的声音则多次不断地发出。对于那些超越本性软弱、从感官事物中解脱出来的人，他们能够听到天上的声音，不需要太多人的声音，但在一切事情上服从那声音，并被它引导；但那些仍然在低处移动、被许多帷幕包裹的人，需要那较低的声音。同样，\Person{若翰}因为他已从感官事物中完全剪除自己，不需要其他导师，而是从天上受教。他说：\BibleQuote{那派遣我以水施洗的，曾对我说：\InnerQuote{你看见圣神降下，停在谁身上，谁就是那一位。}} \BibleRef{若 1:33} 但犹太人仍然像孩子一样，还不能达到那样的高度，他们有一个作为老师的人，这个人没有对他们说自己的话，而是从上面带来信息。\par

他怎么说呢？他\BibleQuote{为他作证，并呼喊说}。\Quote{呼喊}这个词是什么意思？意思是勇敢地、自由地、毫无保留地宣告。他宣告什么？他\Quote{作证}并\Quote{呼喊}什么？\BibleQuote{这就是我曾说过的那一位：在我以后来的，成了在我以前的，因他原先在我以前。}这个证言是隐晦的，并且还包含许多卑微之处。因为他并没有说：\Quote{这是天主子，独生子，真子}；而是说：\BibleQuote{在我以后来的，成了在我以前的，因他原先在我以前。}正如母鸟不会一下子教会幼鸟飞翔，也不会在一天内完成教导，而是先把它们带到巢外，让它们稍作休息，再让它们飞翔，第二天继续更远的飞行，就这样温柔地、一点一点地把它们带到应有的高度；同样，有福的\Person{若翰}没有立即把犹太人引向崇高的事物，而是暂时教导他们在稍微高于地面的地方飞翔，说基督比他更大。然而，这在当时也不是小事：能够说服听众，一个尚未出现、未行任何奇事的人，竟大于一个人（我是指\Person{若翰}本人），他是如此奇妙、如此著名，众人都跑向他，并认为他是一位天使。因此，他暂时努力在听众心中确立这一点：被见证的那一位大于见证者；后来的大于先前的；尚未显现的大于显赫著名的。注意他是如何谨慎地引入他的见证：他不仅在他显现时指出他，甚至在他显现之前就宣告他。因为\BibleQuote{这就是我曾说过的那一位}这句话，是一个宣告此事的人所说的。同样，\Person{玛窦}也说，当众人都来到他那里时，他说：\BibleQuote{我固然以水洗你们，但在我以后来的那一位，比我更强，我连提他的鞋也不配。}那么，为什么他在基督显现之前就这样做呢？为了当基督显现时，这见证能容易被接受，听众的心已被关于他的话预先占据，而卑微的外表不会败坏它。因为如果他们在完全没听说过关于他的任何事之前就看见了主，并且在看见他的同时就听到了\Person{若翰}的话如此奇妙伟大的见证，那么他外表的卑微就会立即成为那些伟大话语的障碍。因为基督取了如此卑微平凡的外表，甚至\Place{撒玛黎雅}妇人、娼妓和税吏都能大胆地接近并与他交谈。因此，正如我所说，如果他们立刻听到这些话并看见他本人，他们也许会嘲笑\Person{若翰}的见证；但现在，因为即使在基督显现之前，他们就已经多次听到并习惯了关于他的话，他们反而受到了相反的影响，不是因被见证者的外貌而拒绝话语的教导，而是因他们已经相信了先前被告知的话，甚至更加尊崇他。\par

\Quote{那以後來的}一語，是指\Quote{那}在\Quote{我以後}講道的，並不是\Quote{那}在\Quote{我以後}出生的。瑪竇也暗示了這一點，他說：\BibleQuote{在我以後要來一個人}，並不是指祂由瑪利亞所生，而是指祂前來講道（福音）；因為若他是指誕生，他就不會說\Quote{要來}，而會說\Quote{已經來了}；因為當若翰說這話時，祂已經誕生了。那麼，\Quote{在我以前}是什麼意思呢？是更光榮、更尊貴的意思。他說：\Quote{不要因為我是先來講道，就以為我比祂大；我遠遠不如祂，甚至不配被算作僕人的行列。}這就是\Quote{在我以前}的意思，瑪竇以另一種方式表達，說：\BibleQuote{我連替祂解鞋帶也不配。}\BibleRef{路 3:16} 再者，\Quote{在我以前}這說法並不是指祂的存有，從下文可以清楚看出；因為若他是指這一點，那麼接下來的\BibleQuote{因為祂在我以前}就是多餘的。因為誰會如此遲鈍和愚蠢，以至不知道那\Quote{生在他以前}的\Quote{在他以前}呢？或者，若這些話是指祂在萬世之前就已存在，那麼所說的無非就是\BibleQuote{那在我以後來的，在我以前就已存在}。此外，這樣的事是無法理解的，而原因也無必要地插入；因為若他想宣告這一點，他本該說相反的話：\BibleQuote{那在我以後來的，在我以前，因為祂也在我以前出生。}因為人有理由將\Quote{生在我以前}作為\Quote{在我以前}的原因，但並非將\Quote{在我以前}作為\Quote{生在我以前}的原因。而我們所主張的則非常合理。因為你們至少都知道，總是那些不確定而非明顯的事，才需要為它們指出原因。現在，如果論證涉及實體的產生，那麼先出生的就必然先存在這點不可能是未知的；但因為他是在論及尊榮，所以他合理地解釋了一個看似困難的問題。因為許多人可能會追問，為什麼那位在後的，卻成了在先的，即顯現出極大的尊榮？為回答這個問題，他立即給出了原因；而這原因就是祂的首先存在。他並沒有說\Quote{祂以某種進步將我這個最先的拋在後面，如此成了在我以前}，而是說\BibleQuote{祂在我以前}，即使祂是在我以後才到來的。\par

但有人會說：如果福音作者指的是祂在人間的顯現，以及隨之而來的榮耀，他為何講論那尚未完成的事，如同已經發生一樣呢？因為他不是說\Quote{將要}，而是說\Quote{已經}。這是古代先知們的習慣：把未來的事當作過去的事來講。因此，依撒意亞論到祂被宰殺時，並沒有說\Quote{祂將被牽去宰殺}（這會表示將來），而是說\BibleQuote{祂被牽去宰殺，如同羊羔}\BibleRef{依 53:7}；然而祂尚未降生成人，先知卻說那將要發生的事如同已經發生。同樣，達味指出十字架時，不是說\Quote{他們將要刺透我的手和我的腳}，而是說\BibleQuote{他們刺透了我的手和我的腳，分了我的衣裳，為我的外衣拈鬮}\BibleRef{詠 22:16}；論到尚未出生的叛徒，他說\BibleQuote{吃我飯的人，舉腳踢我}\BibleRef{詠 40:9}；以及關於十字架的情景：\BibleQuote{他們拿苦膽給我當食物，我渴了，他們拿醋給我喝。}\BibleRef{詠 68:21}\par

4. 你願意我們再舉出更多的例子，還是這些已經足夠？就我而言，我認為已經足夠；因為即使我們沒有在整個地面上翻耕，至少也深挖到了根底；而後一種工作並不比前者少費力；我們擔心過度緊張你的注意力，會使你們後退。\par

那么，让我们给我们的讲论一个合宜的结论。什么是合宜的结论呢？就是适当地将光荣归于天主；那才是适当的，不仅是言语，更是行动。因为祂说：\Quote{你们的光也当这样在人前照耀，好使他们看见你们的善行，光荣你们在天之父。}\BibleRef{玛 5:16} 没有比最卓越的品行更充满光了。正如一位智者所说：\Quote{义人的途径像光一样明亮}\BibleRef{箴 4:18}；它们不仅为那些以工作点燃火焰、作正义之路向导的人发光，也为他们的邻人发光。那么，让我们将这些灯添上油，使火焰更高，发出丰富的光。因为这油如今不仅有强大的力量，甚至在祭献盛行之时，它也比祭献更蒙悦纳。祂说：\Quote{我要的是怜悯，而不是祭献。}\BibleRef{玛 12:7} 这很有理由；因为那是一个无生命的祭台，这是一个有生命的祭台；凡放在那祭台上的，都成为火的燃料，化为尘土，成为灰烬，其烟雾消散于空气中；但这里没有这样的事，它所结的果实是不同的。正如保禄的话所表明的；他在描述格林多人所积蓄的对穷人的慈爱宝藏时写道：\Quote{因为这一供应的事，不但补足了圣徒的缺乏，而且还因许多人感谢天主而丰裕。}\BibleRef{格后 9:12} 又写道：\Quote{他们因你们的服从福音的宣认，以及你们慷慨地救济他们和众人，而赞美天主；他们也为你们祈祷，因对你们的思念。}\BibleRef{格后 9:13} 你们看见这爱德行为化为对天主的感恩和赞美，以及受惠者的不断祈祷和更炽热的爱德吗？那么，亲爱的，让我们在这些祭台上每天祭献。因为这祭献比祈祷、禁食和其他许多事更大，只要求自诚实的收入、诚实的劳动，并远离一切贪婪、掠夺和暴力。因为天主接纳这样的奉献，其他的祂却拒绝和憎恨；祂不愿从别人的灾祸中得光荣，这样的祭献是不洁和亵渎的，反而会惹动天主的愤怒，而不是平息祂。因此我们必须十分谨慎，免得在服侍时反而侮辱了我们所要尊敬的那位。因为如果加音只因献了次等的祭品，没有犯别的错，就受了严厉的惩罚，何况我们若献上掠夺和贪婪得来的东西，岂不更要受更重的惩罚？正是为此，天主给我们显示了这条诫命的模式，好使我们怜悯人，不要严苛对待同仆；但若有人夺取属于一人的东西而给另一人，他并没有显示怜悯，而是造成了伤害，行了极不义的事。正如石头不能出油，残忍也不能生出仁爱；因为施舍若出自这样的根源，就不再是施舍了。因此我劝诫我们不要只看这一点：我们给需要的人施舍，也要注意不要从别人的掠夺中施舍。\Quote{一人祈祷，另一人诅咒，主将俯听哪一个的声音？}\BibleRef{德 34:24} 如果我们这样严格地引导自己，我们就能藉天主的恩宠获得极大的慈爱、怜悯和宽恕，宽恕我们在这漫长岁月中所犯的过错，并逃脱火河；愿我们都能从那里被救出，上升天国，藉我们的主耶稣基督的恩宠和慈爱，愿光荣归于祂，及圣父及圣神，直到永远。阿们。\par

\SourceRecord{Translated by Charles Marriott.；Nicene and Post-Nicene Fathers, First Series；Vol. 14.；Edited by Philip Schaff.；Buffalo, NY: Christian Literature Publishing Co.,；1889.；<http://www.newadvent.org/fathers/240113.htm>.}

% structure: p0001 source=h1 target=section reason=page-title

\section{若望福音讲道集第十四篇}

% structure: p0002 source=h2 target=subsection reason=relative-heading-rank; base=h2

\subsection{若望福音 1:16}

1. 前几日我曾说过，若翰为了解决那些自问主的讲道虽然在他之后，却比他更先、更光荣之人的疑惑，又补充说：\Quote{因为祂在我以前就存在了。}这诚然是一个原因。但他不满足于此，又补充了第二个原因，如今他宣告出来。是什么呢？他说：\Quote{从祂的丰满中，我们都领受了，恩宠上加恩宠。}在此他又提及另一原因。这是什么？就是\par

% structure: p0004 source=h2 target=subsection reason=relative-heading-rank; base=h2

\subsection{若望福音 1:17}

他说：\Quote{从祂的丰满中，我们都领受了}，这是什么意思？因为我们必须暂时将我们的讨论转向这一点。他说，祂不是凭参与而拥有恩赐，祂自身就是一切美善的源泉和根源，就是生命、就是光、就是真理，并非将祂美善的财富保留在自身内，而是将其满溢出来倾注于一切其他受造物，并且在满溢之后仍然保持丰满，因供给他人而毫无减损，反而不断涌流，将一份这些恩惠分施给他人，祂仍保持完美的同一性。我所拥有的，是凭参与（因为我是从他人领受的），并且只是整体中的一小部分，仿佛一滴可怜的水滴相对于无底的深渊或无边的大海；或者更确切地说，就连这个例子也不能完全表达我们试图要说的话，因为如果你从海中取出一滴水，你就减少了海本身，尽管这种减少是察觉不到的。但关于那源泉，我们不能这样说：无论一个人汲取多少，它仍然不见减损。因此，我们有必要转向另一个例子，虽然它也很微弱，不能确定我们所寻求的，但它比前一个例子更能引导我们接近现在向我们提出的思想。让我们假设有一座火的源泉；从那个源泉点燃了一万盏灯，两万盏、三万盏、无数次；难道在把它能力分施给那么多成员之后，火仍然保持在同样的丰满程度吗？人人都清楚它是的。如果在那由部分构成、因抽象而减少的物体中，我们发现有一种具有这样的本性：把某些东西供应给其他物体后自身毫无损失，那么对于那无形体的、非复合的权能，这种情况就更会发生了。如果在所给的例子中，被传递的是实体和物体，它被分开却不受分割，当我们的讨论涉及一种能力，而且是一种无形体实体的能力时，更有可能它不会遭受任何类似的情况。因此若望说：\Quote{从祂的丰满中，我们都领受了，}并将他自己的见证与洗者的见证结合起来；因为\Quote{从祂的丰满中我们都领受了}这个说法不属于前驱，而属于门徒；其意思大致是这样的：\Quote{不要以为，}他说，\Quote{我们这些长期与祂为伴、分享祂食物和餐桌的人，是因偏私而作证，}因为甚至连若翰，他先前根本不认识祂，从未与祂相处，只不过在施洗时在人群中看到祂，就呼喊说：\Quote{祂在我以前就存在了，}也是从那个源头领受了一切；并且我们十二人、三百人、三千人、五千人、无数犹太人，所有那时、现在和将来信者的圆满，都\Quote{从祂的丰满中领受了。}我们领受了什么？他说：\Quote{恩宠上加恩宠。}什么恩宠，相对于什么？相对于旧的，新的。\par

因为既有正义，又有正义，（保禄说：\Quote{按照法律所规定的正义，}\Quote{我是无可指摘的。}）\BibleRef{斐 3:6} 既有信德，又有信德。（\Quote{由信德到信德。}）\BibleRef{罗 1:17} 既有义子名分，又有义子名分。（\Quote{义子名分属于他们。}）\BibleRef{罗 9:4} 既有光荣，又有光荣。（\Quote{因为那先前有光荣的，现在因这更超越的光荣，已算不得光荣了。}）\BibleRef{格后 3:11} 既有法律，又有法律。（\Quote{因为生命之圣神的法律在基督耶稣内释放了我。}）\BibleRef{罗 8:2} 既有服务，又有服务。（\Quote{服务属于他们，}）\BibleRef{罗 9:4} 并且又说：\Quote{以圣神事奉天主。}\BibleRef{斐 3:3} 既有盟约，又有盟约。（\Quote{我要与你们订立一个新盟约，不像我从前与你们祖先所立的盟约。}）\BibleRef{耶 31:31} 既有圣化，又有圣化；既有圣洗，又有圣洗；既有祭献，又有祭献；既有圣殿，又有圣殿；既有割损，又有割损；同样，既有\Quote{恩宠}，又有\Quote{恩宠}。但在前一种情况下，这些词用作预象，在后一种情况下，用作实体，保持相同的发音，尽管意义不同。同样，在模式与形象中，用白线在黑底上划出的人形同被称为人，与那有了正确色彩的人形一样；在雕像的情况下，无论金子或石膏的造型，都同样被称为雕像，尽管在前者是模型，在后者是实体。\par

2. 因此，不要因为用词相同，就以为事物本身相同，也不要以为它们彼此相异；因为它们作为模型，并不与真理相异；但它们仅仅保留了轮廓，故逊于真理。在所有这些例子中，有什么区别呢？你是否愿意我们着手审查其中一两个例子？这样其余的例子对你就会一目了然；并且我们将看到，前者是为孩童的教训，后者是为高尚的成年人的教训；前者是如同为有死之人所定的法律，后者是为天使所定的法律。\par

那么我们从哪里开始呢？是从义子身份本身吗？那么第一个和第二个之间的区别是什么？第一个是名字的荣耀，第二个是实际伴随的事物。关于前者，先知说：\BibleQuote{我曾说：你们是神，你们都是至高者的儿子}\BibleRef{咏 81:6}；关于后者，说他们\Quote{由天主而生}。怎样和以何种方式？藉着重生的洗和圣神的更新。因为他们，即使获得了儿子的称号，仍保有奴仆之灵（因为当他们仍是奴隶时，他们被赐予了这个称谓）；但我们获得释放，得到了这荣耀，不是在名字上，而是在实际上。保禄也宣告说：\BibleQuote{因为你们没有领受奴仆的灵，以致再恐惧；相反，你们领受了义子的圣神，藉着祂我们呼叫：阿爸，父啊！}\BibleRef{罗 8:15}因为重生，且可以说，彻底再造，我们才被称为\Quote{儿子}。若有人考虑圣德的特征，第一个是什么，第二个是什么，他也会发现极大的区别。他们当不崇拜偶像、不犯奸淫或通奸时，就被称为这个名；但我们成圣，不仅因戒除这些恶习，更因获得更大的事物。这恩赐我们首先藉着圣神降临在我们身上获得；其次，藉着比犹太人远为全面的生活准则。为了证明这些话不是空谈，请听祂对他们说的话：\Quote{你们不可占卜，也不可洁净你们的儿女，因为你们是神圣的人民。}因此，他们的圣德在于远离偶像崇拜的习俗；但我们却不是这样。保禄说：\BibleQuote{使她身体和灵魂都圣洁}\BibleRef{格前 7:34}；\BibleQuote{追求和平与圣德，没有圣德，没有人能见到主}\BibleRef{希 12:14}；并说：\BibleQuote{在敬畏天主中成全圣德}\BibleRef{格后 7:1}因为\Quote{圣}这个词并非在所有场合都具有相同的力量；因为天主被称为\Quote{圣}，但不像我们那样。例如，当先知听到那由飞翔的色辣芬所发出的呼声时，他说：\BibleQuote{我完了！因为我是个口唇不洁的人，且住在口唇不洁的人民中间}\BibleRef{依 6:5}；虽然他是圣洁和纯洁的；但若我们与上天的圣德相比，我们是不洁的。天使是圣的，总领天使是圣的，革鲁宾和色辣芬本身是圣的，但这圣德又有双重区别：即，相对于我们，以及相对于更高的权能。我们可以继续所有其他点，但讨论会变得太长，范围会太大。因此我们不再继续深入，留待你们亲自处理其余部分，因为在家中你们有能力将这些事放在一起，考察它们的区别，并以同样的方式查考剩下的。\Quote{给，}有人说，\Quote{给智者一个起点，他就会变得更智慧。}\BibleRef{箴 9:9}起点在于我们，但终点在于你们。我们现在必须接续上下文。\par

在说了\BibleQuote{我们都从祂的丰满中领受了}之后，他又加上了\BibleQuote{恩宠上加恩宠}。因为犹太人是因恩宠得救的：天主说：\BibleQuote{我拣选了你们，不是因为你们人数众多，而是因为你们的祖先。}\BibleRef{申 7:7}如果他们被天主拣选不是由于他们自己的善行，显然他们是因恩宠获得了这荣耀。我们众人也都是因恩宠得救的，但方式不同；不是为了相同的事物，而是为了更大更高的事物。因此，我们所有的恩宠与他们的不同。因为不仅赦罪赐给了我们（这一点我们与他们相同，因为众人都犯了罪），而且还有义德、圣化、义子身份，以及更光荣更丰盛的圣神的恩赐。藉着这恩宠，我们成为天主所爱的，不再为仆人，而是为子女和朋友。因此他说：\BibleQuote{恩宠上加恩宠}。因为连法律的事也是出于恩宠，而且人从无中被创造这一事实本身（因为我们并不是因过去的善行而获得此恩，我们如何能呢？当我们尚不存在时？而是出于那位永远先施恩惠的天主；）不仅我们是从无中被创造，而且当被造时，我们立刻知道什么是该做的，什么是不该做的，并且我们在本性中领受了这法律，我们的造物主将良心的公正准则托付给了我们，这些，我说，都是最伟大的恩宠和不可言喻的仁慈的证明。而这法律在败落后又借着成文的法律得以恢复，这也是恩宠的工作。因为原本预期的是，那曾歪曲过已赐法律的人应受纠正和惩罚；但实际发生的并非如此，相反，是我们本性的修正，以及不是出于亏欠，而是通过仁慈和恩宠的赦免。为了表明这是恩宠和仁慈，请听达味所说：\BibleQuote{主为一切受压迫的人施行正义和审判；祂向梅瑟启示了祂的道路，向以色列子民显示祂的作为}\BibleRef{咏 103:6}；又说：\BibleQuote{主是良善正直的，因此祂必教训罪人走正路。}\BibleRef{咏 24:8}\par

3. 因此，人领受法律是由于怜悯、仁慈和恩宠；为此他说：\BibleQuote{恩宠上加恩宠}。但为了更热切地表达恩赐的伟大，他继续说：\par

第 17 节：\BibleQuote{法律是由梅瑟所传授，恩宠和真理是由耶稣基督而来的。}\BibleRef{若 1:17}\par

你岂不见，洗者若翰与门徒若翰如何温和地、用一句话并逐步地引领听众达至最高知识，先在卑微之事上操练他们？前者将那位远超众人的、无可比拟者与自己相比后，随后以\Quote{在我以前就已存在}之言显示祂的超越，又加上\Quote{祂在我以前就已存在}之语；而后者所做的远超过前者，虽仍不足以匹配独生子的尊荣，因为他并非以若翰，而是以犹太人比若翰更为尊敬的一位——梅瑟——来作比较。他说：\Quote{因为法律是藉梅瑟传授的，恩宠与真理却由耶稣基督而来。}\par

请注意他的智慧。他并非探究身份，而是探究事物；因为一旦这些被证实，即使愚昧之人也必然会由此对基督产生更高更正确的认识。当事实——既不会因偏袒某人，也不会因恶意而受怀疑——作证时，它们便提供了一种连愚人也不可质疑的判断方式；因为事实按行为者的安排留存于眼前，因此其证据最无可疑。请看，他如何使比较对较弱者也变得容易：他不藉推理证明优越性，而是藉单纯的言辞指出差异，将\Quote{恩宠与真理}对立于\Quote{法律}，将\Quote{来到}对立于\Quote{传授}。每对之间都有巨大差别：因为\Quote{传授}属于受命服务之事——一人从他人领受，再给与受命的对象；而\Quote{恩宠与真理却来到}则适宜于一位君王，有权柄地赦免一切过犯，并亲自提供恩赐。因此祂说：\BibleQuote{你的罪赦了}\BibleRef{玛 9:2}；又说：\BibleQuote{为叫你们知道，人子在世上有赦罪的权柄（祂对瘫子说），起来，拿起你的床，回家去吧。}\BibleRef{玛 5:6}\par

你可见\Quote{恩宠}如何藉祂而来？也看\Quote{真理}。祂的\Quote{恩宠}有上述事例，以及发生在右盗身上的事、圣洗圣事的恩赐、由祂所赐的圣神之恩，以及许多其他事都显明了。而祂的\Quote{真理}，我们若明白预像，便更清楚：因为预像如同模型，预先描摹并预演了在新盟约下应完成的安排，而基督来了并使之应验。现在让我们略述预像，因我们此刻无法详谈与之相关的一切；但当你从我即将提出的几个例子中学到一些要点，便会知道其余的了。\par

你是否愿意我们以苦难本身开始？那么预像说什么？\BibleQuote{一家取一只羔羊，宰杀它，并按祂所命所定的去做。}\BibleRef{出 12:3}但基督并非如此：祂不命令这事，而是自己成为羔羊，将自己作为祭献与奉献献给圣父。\par

4. 请看预像如何\Quote{由梅瑟传授}，而\Quote{真理却由耶稣基督而来}。\BibleRef{出 17:12}\par

再者，当阿玛肋克人在西乃山作战时，梅瑟的双手被扶住，由亚郎与胡尔站在他两旁托住\BibleRef{出 17:12}；但当基督来时，祂自己伸开双手在十字架上。你是否已注意到预像如何\Quote{传授}，而\Quote{真理却来了}？\par

再者，法律说：\BibleQuote{凡不持守这律法书所记载的一切事，是应受诅咒的。}\BibleRef{申 27:26}但恩宠怎么说？\BibleQuote{凡劳苦和负重担的，你们都到我面前来，我要使你们安息。}\BibleRef{玛 11:28}；而保禄说：\BibleQuote{基督赎回了我们脱离法律的诅咒，自己成了为我们受诅咒的。}\BibleRef{迦 3:13}\par

既然我们享有了这样的\Quote{恩宠}和\Quote{真理}，我劝告你们，不要因为恩赐的伟大而更加懈怠；因为我们所受的荣耀越大，所欠的德行债也越大。因为那领受微小恩惠的人，即使回报微小，也不该受同样的谴责；但那被提升至荣耀最高峰的人，却表现出卑劣和庸俗的性情，将配受更大的惩罚。但愿我永远不必对你们存有这种疑虑。因为我们依靠主，深信你们已使自己的灵魂振翅高飞向天堂，已远离尘世，虽身在世间却不处理世俗之事；虽然我们如此确信，我们仍不停止这样不断地劝告你们。在外教人的竞技中，观众们鼓励的不是那些跌倒躺卧的人，而是那些正在努力奔跑的人；对于其他人（因为这样做毫无用处，而且不能凭借鼓励使那些已经与胜利无缘的人重新站起），他们便不再理会。但在这种情况下，不仅对你们这些清醒的人，甚至对那些跌倒的人，只要他们愿意悔改，我们仍可期待一些善果。因此我们使用一切方法：劝勉、谴责、鼓励、赞美，为要促成你们的救恩。那么，不要因为我们不断劝告关于基督徒言行的道理而反感，因为这些话不是指控你们懈怠，而是对你们怀有极大希望。不仅对你们，对我们这些讲话的人，这些话也是说的，是的，将来还要说，因为我们同样需要这教导；所以，虽然由我们说出，却不妨碍它们也是对我们说的（因为圣言，当它发现人有过错时，就纠正他；当人清白自由时，就尽可能将他与过错隔离），而我们自己也不是毫无过犯的。治疗的历程对众人相同，药物为众人备好，只有敷用的方法不一样，而是按照用药者的选择：那正确使用药物的人，从敷用中获得益处；而那没有将药物放在伤口上的人，却使病情加重，并导致最痛苦的结局。那么，当我们被治疗时，不要烦恼，反而要大大喜乐，即使这管教之道带来苦痛，因为日后它会为我们结出比任何事物都更甜美的果实。让我们为此全力以赴，好使我们能离世时，洗净灵魂上罪的牙齿所留下的伤口和伤痕，从而配得瞻仰\Person{基督}的圣容，在那一天，我们不会被交付给报复性和残忍的权势，而是交付给那些能带我们进入那为爱祂的人所预备的天上产业的人；愿我们众人，因\Person{主耶稣基督}的恩宠和慈爱，都能达到那产业；愿光荣与权能归于祂，至于无穷之世。阿们。\par

\SourceRecord{Translated by Charles Marriott.；Nicene and Post-Nicene Fathers, First Series；Vol. 14.；Edited by Philip Schaff.；Buffalo, NY: Christian Literature Publishing Co.,；1889.；<http://www.newadvent.org/fathers/240114.htm>.}

% structure: p0001 source=h1 target=section reason=page-title

\section{若望福音讲道词第十五篇}

% structure: p0002 source=h2 target=subsection reason=relative-heading-rank; base=h2

\subsection{若望福音 1:18}

1. 天主不愿我们漫不经心地听圣经中的言辞和句子，而要以极大的注意。因此，真福达味多次在其圣咏前面加上了\Quote{为明瞭}的标题，并说：\BibleQuote{求你开启我的眼目，好使我看出你法律中的奇异之事。}见：\BibleRef{咏 32:42} 此后，他的儿子又显示我们应当\BibleQuote{寻求智慧如寻求银子，搜求她胜于搜求黄金。}见：\BibleRef{箴 2:4} 主在劝勉犹太人\Quote{查考圣经}时，更催促我们探求，因为如果可能初读便立即领悟，祂就不会这样说。没有人会搜寻明显而随手可得之物，而是搜寻笼罩在阴影中、须经大量探求才能寻获之物；为了激励我们搜寻，祂称之为\BibleQuote{隐藏的宝藏}见：\BibleRef{箴 2:4}。这些话对我们说，是要我们不致漫不经心或随意地对待圣经的话，而是极其精确。因为若有人听其中所说而不探求其意义，并按\OriginalTerm{字面}{the letter}接受一切，就会对天主想入非非，认为祂是人，是铜造的，易怒，暴怒，以及其他更糟的看法。但如果他充分学习底下\OriginalTerm{的意义}{the sense}，就会摆脱所有这些不当之见。见：\BibleRef{默 1:15} 我们面前的这段经文说，天主有一个怀抱，这是物质实体所特有的，但没有人会疯狂到认为那无形体者是一个形体。那么，为了按属灵意义正确解释整段，让我们从其开头搜索。\par

\BibleQuote{从来没有人见过天主。}见：\BibleRef{若 1:18} 宗徒是在什么思想脉络下说这句话的呢？在表明基督恩赐的无比伟大以及与梅瑟所施予的之间的无限差异之后，他加上了差异的合理原因。梅瑟作为仆人，是低级事物的仆役，但基督作为主、君王、君王的儿子，为我们带来了远为伟大之物，因为祂永远与父同在，并不断瞻仰祂；因此祂说：\Quote{从来没有人见过天主。} 那么，我们该如何回答那声音宏亮的依撒意亚，他说：\BibleQuote{我看见主坐在高高的宝座上。}见：\BibleRef{依 6:1} 以及若望自己为祂作证说：\Quote{他说这些话，是因为他看见了祂的光荣。}见：\BibleRef{若 12:41} 还有厄则克耳？因为他也在革鲁宾之上看见祂坐着。\EditorialNote{厄则克耳先知书第一和第十章} 还有达尼尔？因为他也说：\BibleQuote{万古者坐下了。}见：\BibleRef{达 7:9} 还有梅瑟本人，他说：\BibleQuote{求你将你的光荣显示给我，好使我能认识你。}见：\BibleRef{出 33:13} 雅各伯也由此得名，被称为\Person{以色列}；因为以色列是\Quote{看见天主的人}。还有其他人也看见过祂。那么若望何以说\Quote{从来没有人见过天主}？这是要表明，所有这些都只是（祂）\OriginalTerm{屈尊俯就}{condescension}的实例，而非\OriginalTerm{本质}{Essence}本身赤裸的显现。因为如果他们真的看见了那\OriginalTerm{本性}{Nature}，他们就不会在不同形式下看见它，因为它是单纯的，无形、无部分、无界限。它不坐，不站，不走；这些都是身体的事。但祂是怎样的，只有祂知道。祂曾通过一位先知这样宣布：\BibleQuote{我用了许多异象，并通过先知之手使用了比喻。}见：\BibleRef{欧 12:10} 也就是说，\Quote{我屈尊俯就了，我并没有显现我的本来面目。} 因为祂的圣子将要在真肉身中显现，祂从古时就预备他们，使他们尽可能看见天主的实体；但天主到底是什么，不仅先知没有看见，而且连天使、总领天使也没有看见。如果你问他们，你不会听到他们关于祂的本质回答什么，只会听到：\BibleQuote{至高天主受光荣，地上平安，善人受恩宠。}见：\BibleRef{路 2:14} 如果你想从\OriginalTerm{革鲁宾}{Cherubim}或\OriginalTerm{色辣芬}{Seraphim}那里学习什么，你会听到那奥秘的圣善之歌，以及\BibleQuote{天地充满祂的光荣。}见：\BibleRef{依 6:3} 如果你询问更高的权能者，你只会发现他们的工作是赞美天主。\Quote{你们要赞美祂，}达味说，\Quote{所有祂的军旅。}见：\BibleRef{咏 147:2} 然而，只有圣子和圣神瞻仰祂。一个受造的本性怎能看见那非受造者呢？如果我们完全无法清楚地辨别任何无形体能力（即使是受造的），正如经常在天使的例子中所证明的，那么我们更不能辨别那无形体、非受造的\OriginalTerm{本质}{Essence}。因此保禄说：\BibleQuote{独者不死，住于不可接近的光中，无人见过，也不能见。}见：\BibleRef{弟前 6:16} 那么，这种特殊属性只属于圣父，而不属于圣子吗？想都别想。它也属于圣子；为证明这一点，请听保禄声明这一点，他说：祂是\Quote{不可见的天主的\OriginalTerm{肖像}{Image}}。见：\BibleRef{哥 1:15} 如果祂是不可见者的肖像，祂自己就必须是不可见的，否则祂就不是一个\Quote{肖像}。并且不要惊奇，保禄在另一处说\BibleQuote{天主在\OriginalTerm{肉体}{Flesh}内显现}见：\BibleRef{弟前 3:16}；因为显现是通过肉体进行的，不是按（祂的）本质。此外，保禄表明祂不仅是人不能见，也是天上权能者不能见，因为说了\Quote{在肉体内显现}之后，他又加上\Quote{被天使看见}。\par

2. 因此，即使对天使，祂也在穿上肉身时才成为可见的；但在那之前，他们并未这样看见祂，因为即使对他们，祂的本质也是不可见的。\par

有人问说：\Quote{那么基督怎么说：\InnerQuote{不要轻看这些小子中的一个，因为我告诉你们：他们的天使在天上，常见我在天之父的面}？\BibleRef{玛 18:10} 难道天主有面孔，并受诸天限制吗？} 谁会如此疯狂而断言这事呢？那么这些话是什么意思？正如祂说：\Quote{心里洁净的人是有福的，因为他们要看见天主} \BibleRef{玛 5:8}，祂是指我们所能有的理智上的看见，以及心中想念天主；同样，关于天使，我们必须明白，因着他们纯洁和不眠的本性，他们不做别的，只常把天主放在自己想象中。因此基督说：\Quote{除了圣子，没有人认识父。} \BibleRef{玛 10:27} 那么，难道我们都在无知中吗？决不是；但没有人像圣子认识父那样认识祂。正如许多人在容许的看见方式中看见了祂，却没有人看见祂的本质；同样，我们许多人认识天主，但祂的实体是什么，除了那由祂所生的，没有人知道。因为祂在这里所说的\Quote{认识}是指精确的观念和理解，如同父对子的认识。\Quote{正如父认识我，我也认识父。} \BibleRef{若 10:15}\par

请注意，因此，福音作者说话时是多么圆满；因为在说了\Quote{从来没有人看见过天主}之后，他并没有接着说\Quote{那看见过的圣子，已将他显示出来}，而是在\Quote{看见}之外又加了\Quote{祂在父怀里}；因为\Quote{住在怀里}远超过\Quote{看见}。因为那仅仅\Quote{看见}的，对于对象并没有完全准确的认识；但那\Quote{住在怀里}的，便一无所不知。现在，免得你听到\Quote{除了圣子，没有人认识父}时，就断言说祂虽然比众人更认识父，却不知道父有多么伟大，福音作者就说祂住在父怀里；而基督亲自宣告，祂认识父如同父认识子一样。因此问那反驳者说：\Quote{请告诉我，父认识子吗？}如果他不疯狂，他必定回答\Quote{是的。}然后再问：\Quote{祂是以精确的看见和知识来看祂认识祂吗？祂清楚地知道祂是什么吗？}他也必定承认这点。由此你便可收集到子对父的精确认识。因为祂说：\Quote{正如父认识我，我也认识父} \BibleRef{若 10:15}；又在另一处说：\Quote{这不是说有人看见过父，惟独从天主来的，祂看见过父。} \BibleRef{若 6:46} 因此，正如我所说的，福音作者提起\Quote{怀里}，是用那一个词向我们显示这一切：那就是本质的密切和亲近，认识毫无差别，能力相等。因为父不会把另一个本质的放在自己的怀里；如果祂是众仆之一，也不敢住在主的怀里，因为这是真儿子才有的，即对祂的父怀有极大信心、并在一切方面不次于父的人所有的。\par

你也要学习祂的永恒吗？请听梅瑟如何论及父。当他问当他被犹太人问是谁派遣他时，他应该如何回答，他听到了这些话：\Quote{你要说：\InnerQuote{我是自有者}派遣了我。} \BibleRef{出 3:14} 而\Quote{自有者}这一表达，意指永远存在，无始存在，真实而绝对的存在。并且这也是\Quote{在起初就有的}这一表达所宣告的，它指永远存在；所以若望用这个词来显示圣子从永远到永远都在父的怀里。为使你不因名称相同，就以为祂是那些因恩宠而成为儿子中的一个，首先，加上了冠词，将祂与那些因恩宠而成为儿子的区分开来。但如果这还不能使你满足，如果你仍然向地上看，请听比这更绝对的名字：\Quote{\OriginalTerm{独生子}{Only-Begotten}。}如果在这之后你仍往地下看，\Quote{我并不拒绝，}（圣若望）说，\Quote{将属于人的词用于天主，我指的是\InnerQuote{怀里}一词，只是不要怀疑有任何卑下。}你看到主的慈爱和关心吗？天主将不配的表达用于自己，使你可以通过这些词语看见祂，并对祂有伟大而崇高的思想；而你仍停留在地上吗？请告诉我，为什么在此处使用了这个粗俗的、物质的词\Quote{怀里}？是为了让我们以为天主是形体吗？绝非如此，祂绝不这样说。那么为什么说它呢？如果藉它既不能确立儿子的真实性，也不能说天主不是形体，那么这个词因为没有用处，就是多余加进来的。那么为什么说它呢？因为我不会停止问你这个问题。除了藉此我们可理解独生子的真实性以及祂与父的\OriginalTerm{同永恒}{Co-eternity}之外，难道还有别的理由吗？\par

3. \Quote{祂已宣示了祂}，若望说。祂宣示了什么？难道是说\Quote{从来没有人见过天主}？或是\Quote{天主是独一的}？但所有其他先知都已为此作证，梅瑟不断呼喊：\Quote{上主你的天主是唯一的上主}\BibleRef{申 6:4}；依撒意亚也说：\Quote{在我以前没有神被塑造，在我以后也没有。}\BibleRef{依 43:10}那么，我们从\Quote{在父怀里的子}，从\Quote{独生子}那里学到了什么更多的东西呢？首先，这些话本身就是由祂的工作所发出的；其次，我们接受了更为清晰的教导，得知\Quote{天主是神，朝拜祂的人应当以心神和真理朝拜祂}\BibleRef{若 4:24}；并且，人不可能看见天主；\Quote{除了子，没有人认识祂}\BibleRef{玛 11:27}；祂是真独生子的父，以及关于祂的一切其他教导。但\Quote{宣示}一词表明祂所给予的更为明白清晰的教导，不单是对犹太人，而是对全世界，并确立了它。先知们甚至未能使所有犹太人听从，但天主的独生子却使全世界都顺服听从。因此，这里的\Quote{宣示}显示祂教导的更大清晰性，故此祂也被称为\OriginalTerm{圣言}{Word}和\OriginalTerm{伟大谋士天使}{Angel of great Counsel}。\par

既然我们蒙受了更宏大、更完美的教导，天主不再藉先知发言，而是\Quote{在这末期内藉着祂的子对我们说了话}\BibleRef{希 1:1}，就让我们展示远比他们更高超的生活，与赐予我们的尊荣相称。如果祂如此屈尊俯就，不再藉仆人，而是亲口对我们说话，而我们却没有比古人表现出更多，那将是多么奇怪的事。他们以梅瑟为导师，我们则以梅瑟的主为导师。那么，让我们展示一种与这尊荣相称的天上智慧，不要与地上的事有任何牵连。为此，祂从天上带来了祂的教导，为了将我们的思想提升到那里，使我们能按自己的能力效法我们的导师。但我们如何能成为基督的效法者呢？通过凡事为公共利益而行，而不只寻求自己的益处。\Quote{因为就连基督}，保禄说，\Quote{没有寻求自己的喜悦，反而正如经上所载：\InnerQuote{辱骂祢者的辱骂都落在我身上。}}\BibleRef{罗 15:3}因此，谁也不要寻求自己的益处。事实上，一个人真正寻求自己的利益时，正是当他关注邻人的利益时。他们的好处就是我们的好处；我们是一个身体，彼此互为肢体。那么，我们不要如同被撕裂一般。谁也不要这样说：\Quote{那个人不是我的朋友，也不是亲戚或邻居，我与他没有关系，我怎能接近他，怎样对他说话呢？}即使他不是亲戚也不是朋友，他却是一个人，与你有同样的本性，属于同一主，是你的同仆，是同路人，因为他生于同一世界。而且，如果他也分享同一信仰，看，他也成了你的肢体：因为什么友谊能造成像信仰关系这样的合一呢？我们彼此间的亲密不仅应当是朋友对朋友那样的亲近，而应当是肢体与肢体之间的亲近，因为没有人能发现比这种友谊和共融更亲密的了。正如你不能说：\Quote{我与这个肢体的亲密和联系从何而来？}（那是可笑的）；同样，你也不能对你的弟兄这样说。\Quote{我们众人都在同一个身体内受了洗}\BibleRef{格前 12:13}，保禄说。\Quote{为何进入一个身体？}为了我们不被撕裂，而是藉着彼此的交往和友谊，保持那一个身体的比例。\par

那么，我们不要彼此轻视，免得我们忽略了自己。\Quote{从来没有人恨过自己的肉身，反而养育和照顾它。}\BibleRef{弗 5:29}因此，天主赐给我们同一个住所——大地，公平地分配一切，为我们点燃同一个太阳，在我们头上伸展同一个屋顶——穹苍，使大地成为同一张桌子，为我们提供食物。祂还赐给我们比这张桌子更美好的另一张桌子，但那也是一张（明白我们奥秘的人懂得我的话），祂赐给所有人同一种出生方式——属神的，我们都有同一个家乡——天上的，我们都饮自同一杯。祂没有赐给富人更丰富、更尊贵的礼物，而赐给穷人更卑微、更微小的，而是同样召叫了所有人。祂给予属世之物时平等地对待所有人，属神之物也同样。那么，生活中为什么会有巨大的不平等呢？来自富人的贪婪和骄傲。但是，弟兄们，不要让这种事再发生了；当具有普遍意义和更大迫切性的事情把我们聚集在一起时，我们不要被世俗和微不足道的事物分裂：我指的是财富和贫穷、肉体的关系、敌意和友谊；因为所有这一切，对于拥有来自高处的爱德联系的人来说，不过是影子，甚至比影子更不真实。那么，让我们保持这联系不被破坏，那些导致如此完美合一分裂的邪灵就没有一个能够侵入；愿我们都能藉着我们的主耶稣基督的恩宠和慈爱达到这一点，愿光荣归于祂，及父和圣神，从今世直到永远。阿们。\par

\SourceRecord{Translated by Charles Marriott.；Nicene and Post-Nicene Fathers, First Series；Vol. 14.；Edited by Philip Schaff.；Buffalo, NY: Christian Literature Publishing Co.,；1889.；<http://www.newadvent.org/fathers/240115.htm>.}

% structure: p0001 source=h1 target=section reason=page-title

\section{若望福音讲道集第十六篇}

% structure: p0002 source=h2 target=subsection reason=relative-heading-rank; base=h2

\subsection{若望福音 1:19}

1. 亲爱的，嫉妒是一件可怕的事，可怕而且有害，对嫉妒者而言，并非对被嫉妒者。因为它首先伤害并消耗他们自己，如同某种致命的毒素深植于他们的灵魂中；即使它偶然伤害了它的对象，所造成的损害也是微小而不足道的，且带来的收益大于损失。的确，不仅是在嫉妒的事上，在所有其他事上，受到伤害的不是那受痛苦的人，而是那做错事的人。若非如此，保禄就不会命令门徒宁可忍受冤屈也不可加害于人，他说：\BibleQuote{为什么不宁愿受冤枉？为什么不宁愿被欺骗？} \BibleRef{格前 6:7}他深知，毁灭常常跟随的不是受害的一方，而是加害的一方。我所说这一切，是因为犹太人的嫉妒。因为那些从城市涌向若翰，并谴责自己的罪，使自己受洗的人，似乎在圣洗之后又后悔了，派人去问他：\Quote{你是谁？}他们实在是毒蛇的种类，蛇，甚至可能比这更坏。邪恶、淫乱、乖谬的世代啊，受洗之后，你们竟变得好奇，质问洗者若翰？还有什么比这更大的愚妄？你们是怎样出来的？你们承认自己的罪，跑到洗者若翰那里去？你们问他该做什么？当你们在做这一切时，你们的行为是毫无理智的，因为你们不知道他来的原则和目的。然而，有福的若翰对此什么也没说，也没有指责或责备他们，而是极其温和地回答了他们。\par

值得学习他为什么这样做。这是为了使他们的邪恶向所有人显明。若翰曾多次向犹太人见证基督，当他为他们施洗时，他不断向众人提起基督，说：\BibleQuote{我固然以水洗你们，但那在我以后来的，比我更强，他要以圣神和火洗你们。} \BibleRef{玛 3:11}他们对他怀有属人的情感；因为战战兢兢地在意世人的意见，并看\BibleQuote{外表的事} \BibleRef{格后 10:7}，他们认为若翰服从基督是不值得的事。因为有许多事指出若翰是一位杰出的人物。首先，他卓越而高贵的出身；因为他是大司祭的儿子。其次，他的生活方式，他的苦修，他对一切属人之事的轻视；因为轻视衣着、餐桌、房屋甚至食物本身，他早先在旷野中度过。在基督的例子中，一切正好相反。他的出身卑微（他们常以此反对他，说：\BibleQuote{这不是木匠的儿子吗？他的母亲不是叫玛利亚吗？他的兄弟不是雅各伯、若瑟……吗？} \BibleRef{玛 13:55}）；而且那被认为是他的家乡的地方名声如此之坏，以至于纳塔乃耳也说：\BibleQuote{纳匝肋能出什么好事吗？} \BibleRef{若 1:46}他的生活方式很普通，他的衣着也不比众人好。因为他没有腰束皮带，也没有穿骆驼毛的衣服，也没有吃蝗虫和野蜜。他像所有人一样生活，也出现在恶人和税吏的筵席上，为吸引他们归向他。犹太人因不明白而以此责备他，正如他自己所说：\BibleQuote{人子来了，也吃也喝，他们便说：看，一个贪吃好酒的人，一个税吏和罪人的朋友。} \BibleRef{玛 11:19}于是当若翰不断将他们从自己那里遣送到耶稣那里时（耶稣在他们看来是更卑微的人），他们为此感到羞愧和烦恼，更愿以若翰为师傅，他们不敢直接这样说，却派人到他那里去，以为能通过谄媚诱使他承认自己是基督。因此，他们不像在基督的事上那样，派卑微的人去抓他时，曾派仆人、黑落德党人等；但这次，他们派了\BibleQuote{司祭和肋未人}，而且不只是\BibleQuote{司祭}，而是\BibleQuote{从耶路撒冷来的}，即更尊贵的人。因为圣史不是无缘无故提及这点。他们派人去问：\Quote{你是谁？}然而，他的出生方式众人皆知，以至于所有人都说：\BibleQuote{这孩子将成为何等人物呢？} \BibleRef{路 1:66}；这消息传遍了整个山地。后来当他来到约旦河时，所有城市都轰动，从耶路撒冷和全犹太地都来到他那里受洗。那么他们现在为什么问呢？不是因为不认识他（既然他已经以多种方式显明，怎么会呢？），而是因为他们想诱使他做我上面提到的事。\par

2. 那么请听，这位有福的人如何回答他们提问的意图，而非问题本身。当他们说：\Quote{你是谁？}他并没有立即给他们直接的答案：\BibleQuote{我是旷野里呼喊者的声音。}而是做了什么？他消除了他们所形成的疑虑；因为，圣史说，当被问及\Quote{你是谁？}时，\par

% structure: p0006 source=h2 target=subsection reason=relative-heading-rank; base=h2

\subsection{若望福音 1:20}

请观察福音作者的智慧。他第三次提到这事，是为彰显洗者若翰的卓越，以及他们的邪恶与愚昧。路加也说，当群众以为他是基督时，他再次消除他们的怀疑。这是忠实仆人的本分：不仅不窃取主人的荣耀，而且当众人将这荣耀给他时，他也拒绝。但群众出于纯朴与无知而产生这种猜想；这些人是出于我所提到的恶意来问他，正如我所说，他们期望用谄媚来拉拢他。倘若他们没有这种期望，就不会立刻转到另一个问题，反而会因他给了与问题无关的回答而恼怒，并会说：\Quote{怎么，我们这样猜想吗？我们来找你问这个吗？}但现在，他们被当场捉住，便转到另一个问题，说：\par

% structure: p0008 source=h2 target=subsection reason=relative-heading-rank; base=h2

\subsection{\BibleBook{若望}{福音} 1:21}

因为他们也期待厄里亚要来，正如基督所宣告的；因为当祂的门徒问：\BibleQuote{那么，经师们怎么说厄里亚必须先来呢？}\BibleRef{玛 17:10} 祂回答说：\BibleQuote{厄里亚的确要先来，并恢复一切。}\BibleRef{玛 17:10} 然后他们问：\BibleQuote{你是那位先知吗？他回答说：不是。}\BibleRef{玛 17:10} 然而他确实是一位先知。那么他为什么否认呢？因为他再次看到问者的意图。他们期待某位特殊的先知到来，因为梅瑟说过：\BibleQuote{上主你的天主要从你们弟兄中给你兴起一位像我的先知，你们要听从祂。}\BibleRef{申 18:15} 这位就是基督。因此他们不说：\Quote{你是先知吗？}（意指普通的先知）；而是用带定冠词的表达：\Quote{你是那先知吗？}意思是\Quote{你是梅瑟所预言的那位先知吗？}因此他不是否认自己是一位先知，而是否认自己是\Quote{那位先知。}\par

% structure: p0010 source=h2 target=subsection reason=relative-heading-rank; base=h2

\subsection{\BibleBook{若望}{福音} 1:22}

请观察他们怎样更激烈地逼问他，催促他，反复提问，不肯罢休；而他却先仁慈地除去关于自己的错误看法，然后向他们提出真实的看法。因为他说：\par

% structure: p0012 source=h2 target=subsection reason=relative-heading-rank; base=h2

\subsection{\BibleBook{若望}{福音} 1:23}

他说了一些关于基督的高超而崇高的话后，好像（回应）他们的想法，就立刻转向那位先知，从那里引出对他主张的证实。\par

% structure: p0014 source=h2 target=subsection reason=relative-heading-rank; base=h2

\subsection{\BibleBook{若望}{福音} 1:24-25}

你岂不见我并非毫无理由地说他们想把他引向这一点吗？他们起初不那样说的原因，是怕被众人识破。然后当他说：\Quote{我不是基督，}他们因想掩饰内心的阴谋，就转到\Quote{厄里亚}，和\Quote{那位先知}。但当他说自己也不是这两位之一时，他们困惑之后，便抛下面具，毫不掩饰地显露出他们的恶意，说：\Quote{你若不是那基督，为什么施洗呢？}然后，又想给这事蒙上一点模糊，他们又加上其他那些：\Quote{厄里亚}和\Quote{那位先知}。因为当他们无法用谄媚绊倒他时，就以为能通过指控迫使他言不由衷。\par

何等愚昧，何等傲慢，何等不合时宜的多管闲事！你们被派来是要了解他是谁、从哪里来，而不是要为他制定法律。这也是那些想迫使他承认自己是基督的人的作为。尽管如此，他仍不生气，也没有像人们可能期望的那样对他们说这类话：\Quote{你们向我发号施令、制定法律吗？}反而再次向他们显出极大的温和。\par

% structure: p0017 source=h2 target=subsection reason=relative-heading-rank; base=h2

\subsection{\BibleBook{若望}{福音} 1:26-27}

3. 犹太人对此还能说什么？即使从这一点来看，对他们的指控也无法回避，针对他们的判决不容宽恕，他们已定了自己的罪。怎么定的？以何种方式？他们视\Person{若翰}为可信之人，如此诚实，以至于不仅在他为他人作证时相信他，甚至在他为自己作证时也相信。因为他们若不是这样，就不会派人去向他学习与他自身相关的事。因为你知道，我们只相信那些我们认为比任何人都更诚实的人，特别是当他们谈论自己的时候。不仅如此，使他们哑口无言的还有他们接近他时的态度；因为他们起初怀着极大的热忱到他那里去，尽管后来改变了。基督在说\Quote{他是点着（且发光）的明灯，你们情愿暂时喜欢他的光}时，同时指出了这两点。此外，他的回答使他更值得信赖。因为（基督）说\Quote{那不寻求自己荣耀的，他是真实的，没有不义在他里面}。而这个人并不寻求自己的荣耀，而是把犹太人指向另一位。而且派去的人是他们中最值得信赖、地位最高的人，因此他们对自己向基督表现出的不信没有任何托辞或借口。为什么你们不接受若翰关于基督所说的话？你们派出了你们中居首位的人，通过他们询问，听到了洗者若翰的回答，他们表现出了一切可能的殷勤，探究了每一个细节，列出了你们怀疑他的一切人：但他最公开、最清楚地承认自己既不是\Quote{基督}，也不是\Quote{厄里亚}，也不是\Quote{那位先知}。他甚至没有停在那里，还告诉他们他是谁，并谈到了他自己圣洗的性质，说那是一件轻微而卑微的事，不过是些水，并讲述了基督所施圣洗的优越性；他还引用了依撒意亚先知，很早以前就作证，称基督为\Quote{主}\BibleRef{依 40:3}，但也给了他\Quote{仆人}和\Quote{奴仆}的名字。此后他们该做什么？难道不应该相信被见证的那一位，敬拜他，承认他是天主吗？因为见证人的品格和属天智慧表明，他的见证不是出于谄媚，而是出于真理；这一点也可以从他不将自己的邻舍置于自己之上看出，当他可以合法地将荣耀归于自己时，他宁愿将它让给另一个人，特别是当这个荣耀如此之大时，就像我们所说的那样。所以，若翰不会放弃这个（属于基督的）见证，如果基督不是天主的话。因为虽然他可能因为对自己的本性太大而拒绝这个见证，但他不会将它归于一个次于它的本性。\par

\Quote{但有一位站在你们中间，你们却不认识他。} 基督混在人群中，像一个普通人，这是合理的，因为他处处教导人不可骄傲自大。而此处，洗者\Person{若翰}用\Quote{认识}这个词，是指对他（基督）的完全了解，即他是谁，从何而来。紧接着，他说出\Quote{在我以后来的}，言下之意是\Quote{不要以为我圣洗包含了一切，因为如果那是完美的，就不会有另一位在我以后兴起，向你们提供另一位；但我的圣洗是一种准备，为那另一位开路。我的不过是影子和预象，但必须有一位来赋予这个实体。所以，他的非常来临\InnerQuote{在我以后}特别彰显他的尊威：因为如果第一个是完美的，就不需要第二个了。} \Quote{在我以前}，更高贵、更光明。然后，免得他们以为他的优越是通过比较发现的，为了确立他的无可比拟，他说：\Quote{我连给他解鞋带也不配}；也就是说，他不仅是\Quote{在我以前}，而且是在这样的意义上在我以前，以至于我不配列在他最卑微的仆人之中。因为解鞋带是最低贱的服务。\par

现在，如果\Person{若翰}不配\Quote{解鞋带}，若翰，\BibleQuote{妇女所生的，没有一位比若翰更大的}\BibleRef{玛 11:11}，我们将把自己置于何处？如果他与全世界相若，甚至比全世界更大（因为\Person{保禄}说：\BibleQuote{世界不配有他们}\BibleRef{希 11:38}），却宣称自己不配被视为基督最卑微的仆人，我们这些满身千万罪恶、离若翰的美德如同天地相隔的人，又能说什么呢？\par

4. 他随后说他甚至连\Quote{解他的鞋带也不配}；然而真理的仇敌却疯狂到如此地步，竟然宣称他们配知道他，如同他知道自己一样。有什么比这更疯狂，比这更傲慢的吗？智者也说得好：\Quote{骄傲的开端是不认识主。}\par

\OriginalTerm{魔鬼}{devil}就不会被贬抑而成为魔鬼（它原来并不是魔鬼），倘若它没有患上这疾病。正是这疾病将它从那自信中驱逐出去，将它投入火坑，成为它一切灾祸的缘由。因为这疾病本身足以摧毁灵魂的一切卓越，无论是施舍、祈祷、守斋或任何事。因为福音作者说：\BibleQuote{那在人们眼中崇高的，在主前却是可憎的。}\BibleRef{路 16:15}因此，不仅是奸淫或淫乱常玷污行这些事的人，\OriginalTerm{骄傲}{pride}也一样，而且远超过那些恶习。为什么呢？因为奸淫虽是不可赦之罪，人尚可推脱于情欲；而骄傲却绝不可能找到任何理由或托词，甚至一丝一毫的借口也没有；它纯粹是灵魂的扭曲和最严重的疾病，其根源无非是愚妄。因为再也没有比骄傲的人更愚妄的了，即使他拥有财富，拥有许多世俗的智慧，身居王位，怀有一切在人们眼中可羡慕的东西。\par

因为若那以真正的好事而\OriginalTerm{骄傲}{pride}的人是可怜和不幸的，且失去这一切事物的赏报，那么那因虚无之物而自高自大，因阴影或花草（因为现世的荣耀就是如此）而自夸的人，岂不是比任何人都更可笑吗？他正如同一个贫穷困乏的人，终生因饥饿而憔悴，若偶然一夜梦见好运，便因此而充满傲慢。\par

可怜而可悲啊！当你的灵魂因最严重的疾病而丧亡，当你贫穷到一无所有时，你竟因拥有若干塔冷通的金子而心高气傲吗？因拥有众多奴隶和牲畜吗？然而这些并不属于你；如果你不相信我的话，就从先你而去的人的经历学习吧。如果你如此沉醉，甚至无法从别人的遭遇中得到教诲，那么稍等片刻，你将从自己的遭遇知道，这些东西对你毫无益处：当你奄奄一息，连一个时辰，甚至一刻钟都不能掌握时，你将不情愿地把它留给身边的人，而这些或许是你所不愿给的人。因为许多人甚至未被允许处理它们，就突然离世，渴望享受它们，却不被允许，被拖离它们，被迫将之让给他人，强迫让给那些他们不愿给的人。为避免我们遭遇同样的事，让我们趁还有力量和健康时，将我们的财富从此生送往我们本城；因为只有如此，别无他法，我们才能享受它们。这样我们将把它们存入一个不可侵犯且安全的地方。因为那里没有任何东西，没有任何东西能从我们手中夺走它们；没有死亡，没有经过公证的遗嘱，没有继承遗产的后继者，没有虚假告发，没有针对我们的阴谋；而那个带着巨大财富从此世离去的人可以在那里永远享受它。那么，还有谁如此可悲，竟不愿意永久享受真正属于自己的财富呢？让我们转移我们的财富，并迁往那里。我们不需要驴、骆驼、车或船来进行这样的迁移（天主甚至为我们免除了这一困难），我们只需要穷人、跛子、残疾人、病弱者。他们被托付进行此转移，他们将我们的财富运送到天堂，他们将这些财富的主人引入永恒产业的继承。愿我们全体都能通过我们的主耶稣基督的恩宠与慈爱获得此产业，愿光荣藉着祂、偕同祂，归于圣父及圣神，从现在直到永远，永世无尽。阿们。\par

\SourceRecord{Translated by Charles Marriott.；Nicene and Post-Nicene Fathers, First Series；Vol. 14.；Edited by Philip Schaff.；Buffalo, NY: Christian Literature Publishing Co.,；1889.；<http://www.newadvent.org/fathers/240116.htm>.}

% structure: p0001 source=h1 target=section reason=page-title

\section{论若望福音 第十七篇讲道}

% structure: p0002 source=h2 target=subsection reason=relative-heading-rank; base=h2

\subsection{若望福音 1:28-29}

勇敢和言论的自由，以及把一切次要于承认基督，这是一项伟大的美德；如此伟大而可佩服，以至天主的独生子在圣父面前宣布这样的人。\BibleRef{路 12:8} 但赏报更是公义的，因为你在世上承认，祂在天上承认；你在人前承认，祂在圣父及众天使前承认。\par

这样的人就是若翰，他不顾及群众、意见，也不顾及任何属于人的事物，而是将这一切践踏在脚下，以合宜的自由向众人宣讲关于基督的事。因此，福音作者特别标明地点，以显示这位高呼先驱的勇敢。因为不是在屋里，不是在角落里，也不是在旷野，而是在群众中间，在他占据了约但河之后，当所有受他洗礼的人都到场时（因为犹太人趁他施洗时来到他那里），就在那里，他大声作出了那奇妙的关于基督的宣认，充满了那些崇高、伟大、奥秘的教义，并说他不配解祂的鞋带。因此他说：\Quote{这些事发生在伯达尼，}或者，正如所有更准确的抄本所载，\Quote{在伯大巴喇。}因为伯达尼并不在\Quote{约但河外}，也不靠近旷野，而是在耶路撒冷附近。\par

他标明地点也出于另一个原因。因为他不是要叙述古老的事情，而是那些不久前刚刚发生的事，所以他使那些在场并目睹的人成为他话语的证人，并从地点本身提供证据。因为他确信自己所言毫无添加，而是诚实地按实际情况描述事物，所以他从地点获取见证，正如我所说，这将成为他真实性的一个不寻常的证明。\par

\Quote{次日，他看见耶稣向他走来，便说：\InnerQuote{看，天主的羔羊，除免世罪的。}}\par

福音作者们在他们之间分配了时期；玛窦缩短了在洗者若翰被拘禁之前的时间记载，急忙转向后来的事，而福音作者若望不仅没有缩短这个时期，反而详细叙述。玛窦在耶稣从旷野回来之后，对中间的情况只字未提，如若翰说了什么，以及犹太人派遣和说的话，并缩短了其余所有，直接转到监牢。\Quote{因为，}他说，\Quote{耶稣听见}若翰被出卖，\Quote{就离开那里。} \BibleRef{玛 14:13} 但若望并非如此。他对于旷野之行保持沉默，因为玛窦已经描述过；他叙述了从山上下来之后的事，并在经历了许多之后，补充说：\Quote{因为若翰还没有被投进监里。} \BibleRef{若 3:24}\par

有人问，耶稣现在为什么到他这里来？为什么他不只是来一次，而是第二次又来？因为玛窦说，祂来是由于洗礼的必要：因为耶稣补充说，\Quote{这样，我们应当这样履行全义。} \BibleRef{玛 3:15} 但若望说，祂在受洗后又来了，并在此时宣告，因为他说：\Quote{我看见}，他说，\Quote{圣神像鸽子从天上降下，停在祂身上。} 那么，祂为什么来到若翰那里？因为祂不是偶然来的，而是特意到他那里。\Quote{若翰}，福音作者说，\Quote{看见耶稣向他走来。} 那么祂为什么来？为的是，既然若翰与其他许多人一同给祂施洗，免得有人以为祂到若翰那里去，是像其他人一样为了认罪，并在河中洗忏悔的澡。祂来正是为此，给若翰一个机会来纠正这个看法，因为若翰说：\Quote{看，天主的羔羊，除免世罪的。} 祂就消除了所有的嫌疑。因为很明显，这样一位纯洁到足以洗涤他人罪恶的，并非来认罪，而是为了使那位奇妙的先驱有机会将此前说过的话更确定地印在那些听过他先前言语的人心上，并加以补充。使用\Quote{看}这个词，是因为许多人因之前所说的话而长期寻找祂。为此，当祂出现时，他指着祂说：\Quote{看，}这就是长期寻找的这位，这就是\Quote{羔羊。} 祂称祂为\Quote{羔羊}，为提醒犹太人依撒意亚的预言，以及在梅瑟法律下的预象，以便更好地引导他们从预象走向现实。梅瑟的那只羔羊并没有立即除去任何人的罪，但这只却除去了全世界的罪；因为当世界濒临毁灭时，祂迅速将其从天主的义怒中解救出来。\par

% structure: p0009 source=h2 target=subsection reason=relative-heading-rank; base=h2

\subsection{若望福音 1:30}

2. 你在这里也看到他如何解释\Quote{在我以前}这个词？因为称祂为\Quote{羔羊}，又说祂\Quote{除免世罪}，然后他说\Quote{祂在我以前就已存在，因为祂在我之前}；声明这就是\Quote{在我以前}的意思，即承担世罪，\Quote{并以圣神施洗。} \Quote{因为我的来临，除了宣告世界的共同恩主，并提供水的洗礼之外，别无目的；但祂却是为了洁净众人，并赐给他们护慰者的能力。} \Quote{祂在我以前就已存在}，也就是说，祂已显得比我更光荣，因为\Quote{祂在我之前。} 愿那些接受了撒摩沙塔的保禄的疯狂的人，在反对如此明显真理时感到羞愧。\par

% structure: p0011 source=h2 target=subsection reason=relative-heading-rank; base=h2

\subsection{若望福音 1:31}

在此，他通过表明自己的见证并非出于人情，而是出于天主的启示，使他的见证免于嫌疑。他说：\Quote{我先前不认识祂。} 那么，你怎么能成为可信的证人？你自己尚且无知，又怎能教导别人？他没有说\Quote{我不认识祂}，而是说\Quote{我先前不认识祂}；这样，他就显得最可信；因为何以要对一个他所不认识的人表示偏私呢？\par

\Quote{但是祂应被显示给以色列，因此我用水来施洗。}\par

祂当时并不需要洗礼，那洗涤除了为所有人准备一条通往基督信仰的道路外，别无其他目的。因为他并没有说，\Quote{为叫我能洁净受洗者，}或\Quote{为叫我能把他们从罪中解救出来，}而是说，\Quote{为叫祂应被显示给以色列。}\Quote{为什么，请告诉我，他不能不用洗礼就宣讲并带领群众到祂那里呢？}但若如此，就决不会容易。因为如果宣讲没有洗礼，他们就不会都一起跑来；他们也不会通过比较而认识祂的优越性。因为群众聚来并不是为了听他的话，却是为了什么？为了\Quote{受洗，承认他们的罪。}但他们来了以后，就接受了关于基督的事以及祂洗礼的不同之处的教导。然而，即便是约翰的洗礼也比犹太人的更尊贵，因此众人都奔向它；但即便如此，它仍是不完善的。\par

他说：\Quote{那么你是如何认识祂的？}\Quote{藉着圣神的降临。}但是，为避免有人以为他像我们一样需要圣神，请听他如何消除这猜疑，表明圣神的降临只是为了宣告基督。因为他曾说：\Quote{我并不认识祂，}他接着说：\Quote{但那派遣我用水施洗的，对我说：\InnerQuote{你看见圣神降下并停留在谁身上，谁就是那用圣神施洗的。}}\BibleRef{若 1:33}\par

你是否看到，这圣神的工作就是指明基督？若翰的见证确实无可怀疑，但为了使它更加可信，他将它归诸天主和圣神。因为当若翰见证如此伟大奇妙、足以令所有听者震惊的事——即唯有祂承担了全世界的罪，且这恩赐的伟大足以支付如此大的赎价——之后，他证明了这个断言。而这证明就是祂是天主子，祂不需要洗礼，圣神降临的目的只是为了显示祂。因为若翰并没有能力赐予圣神，那些受他洗礼的人证明了这一点，他们说：\Quote{我们连有没有圣神都不知道。}\BibleRef{宗 19:2}事实上，基督不需要洗礼，无论是若翰的洗礼还是任何其他洗礼；相反，洗礼需要基督的能力。因为所欠缺的是那至高的祝福，即受洗者应被算为配得圣神；祂来到时便加上了这圣神的恩赐。\par

% structure: p0017 source=h2 target=subsection reason=relative-heading-rank; base=h2

\subsection{若望福音 1:32-34}

他多次重复\Quote{我不认识祂。}为什么，有何原因？祂按肉身是他的亲戚。天使说：\Quote{你看，你的亲戚依撒伯尔，她也怀了一个儿子。}\BibleRef{路 1:36}因此，免得他因亲戚关系而显得偏袒祂，他重复了\Quote{我不认识祂。}这发生得很有道理；因为他一直在荒野中度过，远离父家。\par

那么，如果他在圣神降临之前不认识祂，而且那时他第一次认出了祂，他怎能在洗礼前阻止祂，说：\Quote{我需要受祢的洗，祢反而到我这里来吗？}\BibleRef{玛 3:14}？这证明了他非常认识祂。然而，他在此前或长时间内并不认识祂，这是有充分理由的；因为祂童年时代发生的奇事，如贤士们的情况和其他类似事件，早已发生，而若翰当时非常年幼，并且中间过了许久，祂自然不为人所知。因为如果祂被认识，若翰就不会说：\Quote{为叫祂应被显示给以色列，因此我用水来施洗。}\par

3. 因此我们清楚得知，那些他们所说的属于基督童年的奇迹是假的，是一些人捏造并传播的。因为如果祂从幼年就开始行奇事，若翰不可能不认识祂，群众也不需要一位老师来显明祂。但现在他说，他为此而来，\Quote{为叫祂应被显示给以色列}；因此他又说：\Quote{我需要受祢的洗。}后来，在获得了对祂更确切的认识以后，他向群众宣告说：\Quote{这就是我曾说过的：在我以后来的那个人，成了在我以前的。}因为\Quote{那派遣我用水施洗的}，并为此目的派遣我，\Quote{为叫祂应被显示给以色列}，祂自己在圣神降临之前就已启示了祂。因此，甚至在祂来到之前，若翰就说：\Quote{在我以后来的，成了在我以前的。}他是在来到约但河并给众人施洗之前不认识祂的；但当祂即将受洗时，他就认识了祂；这是由于父向先知启示祂，以及圣神在祂受洗时向犹太人指明祂，为了他们的缘故圣神确实降临了。因为为了使若翰的见证不被轻视，他曾说祂是\Quote{在我以前的}、\Quote{祂要用圣神施洗}、\Quote{祂要审判世界}，圣父发出声音宣告圣子，圣神降临，将声音引向耶稣的头上。因为既然一个人施洗，另一个人受洗，圣神来纠正在场某些人可能形成的想法，以为那些话是针对若翰说的。所以，当他说\Quote{我不认识祂}时，他说的是以前的时候，并非临近祂受洗时。否则他怎么阻止祂，说\Quote{我需要受祢的洗}？他又怎能说这样的话关于祂呢？\par

\Quote{但是，}有人说，\Quote{那犹太人为何没有相信？因为并非只有若望看见圣神像鸽子一样降下。}这是因为，即使他们看见了，这样的事不仅需要肉体的眼睛，更需要理智的视界，以免人们以为整个是虚妄的幻觉。因为当看见祂行奇迹，亲手触摸病人和死者，使他们恢复生命和健康时，他们仍被恶意所醉，宣称与所见相反；那么仅仅借着圣神的降临，他们又怎能摆脱不信呢？有些人说，并非所有人都看见了，只有若望和他们中那些心性较好的人看见。因为即使能用肉眼看见圣神像鸽子一样降下，也并不必然所有人都能看见这情景。因为匝加利亚看见许多有形的事物，达尼尔和厄则克耳也是如此，没有人分享他们所见的；梅瑟也看见许多别人未曾见的事物；并不是所有门徒都享见山上显圣容，也不是所有人在复活时都同样看见祂。路加清楚表明这一点，他说祂显现给\InnerQuote{天主预先拣选的见证人。}\BibleRef{宗 10:41}\par

\Quote{我看见了，便作证：这就是天主子。}\par

他在哪里\Quote{作证说：这就是天主子}呢？他确实称祂为\Quote{羔羊}，并说祂要\Quote{以圣神施洗}，但没有一处称祂为\Quote{天主子}。然而其他福音作者并未记载祂在受洗后说了什么，他们省略了中间的时间，只提到基督在若翰被囚后所行的奇迹。由此我们可以合理推断，这些以及许多其他事都被省略了。我们的福音作者本人在叙述结束时也如此声明了。因为他们远未发明任何关于祂的伟大事物，反而那些似乎带来羞辱的事，他们都异口同声、极其精确地记录下来；你不会发现他们中任何人省略了这样的事；但奇迹中，有些一部分留给他人叙述，一部分却都被略过不提。\par

我说这些并非无的放矢，而是为了回应外邦人的无耻。因为这是他们热爱真理的充分证据，表明他们不为偏袒而说任何话。这样，你们也得以武装自己，与他们进行辩论。但要注意。奇怪的是，医生、鞋匠、织工，总之所有工匠，都能正确地为自己技艺争辩，而一个自称基督徒的人却不能为自己的信仰辩护；然而前者的忽视只带来财产的损失，后者的忽视却毁灭我们的灵魂。但我们的心态如此可怜，以至于我们全心关注前者，而轻视那些必要的、我们救恩基础的事，仿佛无足轻重。\par

4. 正因如此，外邦人不能迅速唾弃自己的错误。因为当他们在谎言中扎根，用尽一切手段掩盖自己观点的可耻时，我们这些真理的仆人却连口都张不开，他们怎能不谴责我们教义的极大软弱？他们怎能不怀疑我们的宗教是欺骗和愚妄？他们怎能不亵渎基督为骗子、为欺诈者，利用众人的愚昧来推进他的骗局？而这亵渎要归咎于我们，因为我们不愿在虔敬的论证中警醒，反而视之为多余，只关心地上的事物。凡是欣赏舞者、车夫或与野兽搏斗的人，无不竭尽全力，想尽办法不在任何有关他们的争论中落败；他们连篇累牍地颂扬，为自己的辩护针对批评者，向对手投掷无数嘲讽；然而当基督信仰的论据被提出时，所有人都垂头丧气，搔首弄姿，打哈欠，最终在轻蔑中退去。\par

这难道不应招致极重的愤怒吗？当基督在你们眼中显得比一个舞者更不体面时？因为你们为他们的所作所为编造了万般辩护，尽管那些事比任何事都更可耻；但对于基督的奇迹，尽管它们将世界吸引到祂面前，你们却丝毫不能忍受思想或关心。我们相信圣父、圣子、圣神、肉身的复活和永生。如果现在有外教人说：\Quote{这位圣父是什么？这位圣子是什么？这位圣神是什么？你们说有三位神，又如何指责我们有多神？}你们要说什么？你们要回答什么？你们如何抵挡这些论点的攻击？但如果当你们沉默时，不信者又提出另一个问题，问说：\Quote{复活究竟是怎么回事？我们将在这一身体内复活？还是在别的、与这不同的身体内？如果在这身体内，何以需要它被分解？}你们要回答什么？又如果他问：\Quote{基督为什么现在来，而不是在古代？祂现在才关心人，而过去的所有年日中祂轻视了我们吗？}或者他再问其他问题，比这些还多？因为我不可提出许多问题，而对它们的回答保持沉默，免得这样做，伤害你们中较纯朴的人。已经说的话足以唤醒你们的沉睡。那么，如果他们问这些问题，而你们全然不能听那些话，请问，当我们成了坐在黑暗中的人如此错误的缘由时，我们只受轻微的惩罚吗？我曾经希望，如果你们有足够的闲暇，将某个不洁的外教哲学家所写的反对我们的整本书，以及另一本较早的书摆在你们面前，这样至少可以唤醒你们，使你们脱离极度的懈怠。因为如果他们警醒着说出这些反对我们的话，我们若不甚至知道如何击退对我们的攻击，还能得到什么宽恕？我们被带出来是为了什么目的？你们没有听见宗徒说：\BibleQuote{要随时准备回答每一个问你们心中盼望之缘由的人}？\BibleRef{伯前 3:15}保禄也同样劝勉说：\BibleQuote{要让基督的话丰丰富富地住在你们内。}\BibleRef{哥 3:16}那些比懒惰的人更懈怠的人如何回答这个？\BibleQuote{纯朴的灵魂是有福的，行为纯朴的人行走稳妥。}\BibleRef{箴 10:8}因为这就是各种邪恶的起因：许多人甚至不会正确运用圣经的见证。所以在此处，作者并不是指（用\Quote{纯朴}一词）愚昧或一无所知的人，而是指没有邪恶、不作恶、有智慧的人。若非如此，说\BibleQuote{你们要机警如同蛇，纯朴如同鸽子}就是无用的了。\BibleRef{玛 10:16}但我为什么提及这些事呢？因为这篇讲道来得不太合时宜？因为除了已经提到的事，我们还有别的事不妥当，就是关乎我们的生活和言行。我们在各方面都是可怜又可笑的，总是准备好彼此挑剔，但对自己身上责怪和控告邻人的事却改正得很慢。因此我劝诫你们，现在至少要注意我们自己，不要停留在挑剔上（这不足以平息天主），而是要显出各方面最卓越的改变，好使我们在此生为光荣天主而生活，从而享受将来的光荣；愿我们众人因我们的主耶稣基督的恩宠与慈爱都能获得它，愿光荣归于祂，直到永远。阿们。\par

\SourceRecord{Translated by Charles Marriott.；Nicene and Post-Nicene Fathers, First Series；Vol. 14.；Edited by Philip Schaff.；Buffalo, NY: Christian Literature Publishing Co.,；1889.；<http://www.newadvent.org/fathers/240117.htm>.}

% structure: p0001 source=h1 target=section reason=page-title

\section{若望福音讲道集 第十八篇}

% structure: p0002 source=h2 target=subsection reason=relative-heading-rank; base=h2

\subsection{若望福音 1:35-37}

1. 人的本性在某种意义上是一种怠惰的事物，容易趋向丧亡，不是由于本性自身的构造，而是由于那出自自主选择的怠惰。因此它需要很多的提醒。为此，保禄写信给斐理伯人说：\BibleQuote{把同样的事写给你们，对我并不厌烦，对你们却是稳妥的。} \BibleRef{斐 3:1}\par

土地一旦接受了种子，立刻就结出果实，不需要第二次播种；但我们的灵魂却不是这样，一个人必须满足于多次播种并表现出许多谨慎之后，才能一次收获果实。因为首先，所说的话很难在头脑中安顿下来，因为土地非常坚硬，且被无数荆棘缠绕，还有许多设埋伏并把种子夺走的人；后来，当它被固定并扎根之后，仍然需要同样的关注，才能成熟，并且成熟后能够不受伤害，不被任何事情损害。因为就种子而言，当穗子完全形成并获得适当的力量后，它很容易忽略锈病、干旱和所有其他事物；但教义却不是这样；在它们的情况下，当所有工作圆满完成之后，常常会来一场暴风雨和洪水，要么由于不愉快环境的攻击，要么由于善于欺骗之人的阴谋，要么由于各种其他针对它们的诱惑，而将它们毁灭。\par

1. 我说这些话不是没有原因的，而是为了使你们听到若翰重复同样的话时，不至于谴责他空谈，也不至于认为他无礼或令人厌烦。他希望一次说话就被听见，但因为许多人由于沉睡没有注意起初所说的话，他就通过这第二次呼唤再次唤醒他们。现在请注意：他曾说：\BibleQuote{那在我以后来的，成了在我以前的}；以及\BibleQuote{我连解祂鞋带也不配}；以及\BibleQuote{祂要以圣神和火施洗}；以及他\BibleQuote{看见圣神像鸽子降下，停在祂身上}，并\BibleQuote{作证说，这一位是天主子}。没有人留意，没有人询问，也没有人说：\Quote{你为什么说这些事？为谁？什么原因？}他又曾说：\BibleQuote{请看天主的羔羊，除免世罪者}；但即便如此，他也未能触动他们的麻木不仁。因此，在这之后，他被迫再次重复同样的话，好像通过耕种使一些坚硬顽固的土壤变软，并通过他的言语作为犁，搅动已经硬结成块的头脑，以便把种子种得深。为此，他也不把讲论弄得很长；因为他只渴望一件事，就是把那些人引过来，使他们与基督联合。他知道，一旦他们接受了这话并信服了，以后就不再需要有人为祂作证了。事实也正是如此。因为，如果撒玛黎雅人在听了祂之后能对那妇人说：\BibleQuote{现在我们信，不是因为你的话，因为我们知道这一位确实是世界的救主，基督}，那么门徒们就会更快地被降服，事实也正是如此。因为圣史说：\BibleQuote{他找到自己的兄弟西满，对他说：我们找到了默西亚，——默西亚翻出来就是基督。}并且，我请你们注意这一点：当他说\BibleQuote{那在我以后来的，成了在我以前的}，以及\BibleQuote{我连解祂鞋带也不配}时，没有吸引任何人；但当他说到救恩计划，并把讲论降低到更谦卑的调子时，门徒们便跟随了祂。\par

我们不仅可以在门徒身上注意到这一点，而且许多人并不那么被论及天主的伟大崇高之辞所吸引，反而被谈到恩慈和仁爱、与听者救恩有关的事所吸引。他们听见\BibleQuote{祂除去世罪}，就立刻跑向祂。因为他们说：\Quote{如果有可能洗掉我们身上的控告，我们为什么拖延？这里有一位无须我们劳苦就能解救我们的。推迟接受这份恩赐岂不是极端的愚昧吗？}让那些是慕道者、并且把救恩拖延到最后一口气的人听吧。\par

圣史说：\BibleQuote{再，若翰站着，说：请看，天主的羔羊。}基督一言不发，祂的使者说了一切。新郎也是如此。他暂时对新娘什么也不说，只是沉默地站在那里，有人把他指给新娘，另一些人把新娘交到他手里；她只是出现，他并没有亲自带走她，而是从另一个把她交给他的人那里接受了她。当他这样接受了新娘之后，就如此安排她，使她不再记得那些把她许配给他的人。基督也是如此。祂来是要把教会与自己结合；祂什么也没说，只是来了。是祂的朋友若翰，把新娘的右手放在祂手里，当他通过自己的讲论把人的灵魂交到祂手里时。祂接受了他们之后，就如此安排他们，使他们不再回到那把他们托付给祂的若翰那里。\par

2. 在此我们不仅要注意这一点，还要注意另一件事。如在婚礼上，少女不是到新郎那里去，而是新郎急忙到少女那里，即使他是国王的儿子，即使他要娶一个贫穷卑微的人，甚至一个仆人，这里也是如此。人的本性没有上升，而是卑微贫穷的它，祂来到它那里；而当婚姻完成后，祂不再允许它留在这里，而是把它带回自己那里，迁到了祂父的家中。\par

\newline{}\Quote{那麼，若翰為什麼不把自己的門徒帶到一邊，私下與他們談論這些事，然後把他們交給基督，卻公開對他們和全體人民說：\Quote{請看，天主的羔羊}？}這是為了避免讓人覺得這是安排好的；因為如果他們在私下受他告誡之後，好像為了給他面子而離開他去跟隨基督，他們或許很快就會離開；但如今，他們接受了從一般教導而來的跟隨祂的責任，此後便成了祂堅定的門徒，因為他們跟隨祂不是為了討好老師，而是純粹為了自己的益處。\par

\newline{}眾先知和宗徒都宣講缺席的祂；先知在祂按肉身來臨之前，宗徒在祂被接升天之後；唯獨若翰宣講在場的祂。因此他自稱\OriginalTerm{新郎的朋友}{friend of the Bridegroom}\BibleRef{若 3:29}，因為只有他出席婚禮，他完成了所有事，他開始了這項工作。\Quote{他看見耶穌走著，就說：\Quote{請看，天主的羔羊。}}他不僅用聲音，也用眼睛為基督作證，表達讚嘆，歡喜和光榮。他暫時沒有對跟隨者說任何勸勉的話，只顯示對在場者的驚嘆和崇敬，並向眾人宣告祂來賜予的恩賜和潔淨的方式。因為\Quote{羔羊}同時宣告了這兩點。他沒有說\Quote{誰將除去}或\Quote{誰已除去}，而是說\Quote{誰除去世界的罪孽}；因為祂始終如此行。祂不僅在受苦時除去罪，從那時直到現在仍然除去，不是被重複釘死（因為祂為罪獻上了一次祭獻），而是藉著那一次祭獻不斷潔淨他們。因此，正如\Quote{聖言}顯示祂的超卓性，\Quote{聖子}顯示祂與其他相比的優越，同樣，\Quote{羔羊、基督、那位先知、真光、善牧}等等帶有冠詞的名稱，都標示著巨大的差異。因為有許多\Quote{羔羊}、\Quote{先知}、\Quote{基督}和\Quote{兒子}，但若翰將祂與這一切遠遠分開。他不僅藉著冠詞，還藉著加上\OriginalTerm{獨生子}{Only-Begotten}來確保這一點；因為祂與受造物毫無共同之處。\par

\newline{}若有人覺得在\Quote{第十時辰}談論這些事不合時宜（那是白天的時刻，因為他說：\Quote{約在第十時辰}\BibleRef{若 1:39}），在我看來那人就大錯特錯了。的確，對於許多人以及事奉肉體的人來說，宴席後的時段不適合處理緊要事情，因為他們的心被食物所壓；但這裡是一個甚至不分享普通食物的人，他在晚上就像我們在早晨一樣清醒（或者更清醒；因為我們晚上所留食物的殘渣常在我們內充滿想像，但他器皿中沒有裝載這些東西）；他很有理由在晚間談論這些事。此外，他停留在約但河邊的曠野，所有的人都懷著極大恐懼到他那裡受洗，那時很少關心今世的事物；正如他們與基督同住了三天，卻沒有吃什麼。\BibleRef{瑪 15:32}因為這是熱心的傳報者和謹慎的農夫所做的事，即在他看到播下的種子牢牢紮根之前不停止。\Quote{那麼，他為什麼不走遍猶太各處宣講基督，卻站在河邊等候祂來，好指認祂呢？}因為他希望這事由祂的工作來完成；同時他自己的目標只是讓人認識祂，說服一些人聽從永生。但祂把更偉大的見證，即工作的見證，留給祂自己，正如祂所說：\Quote{我不接受人的見證；但父交給我的工作，正是我所作的，為我作證。}\EditorialNote{參見第34、36節}請注意這有多麼有效；因為當他投入一點火花時，火焰立刻沖天。那些先前甚至不留意他話語的人，後來卻說：\Quote{若翰所說的一切都是真的。}\BibleRef{若 10:41}\par

\newline{}此外，如果他四處說這些事，所做的事便會顯得出於某種人性的動機，而宣講也充滿了嫌疑。\par

\newline{}\Quote{那兩個門徒聽見他說，就跟隨了耶穌。}\par

\newline{}然而若翰還有其他門徒，但他們不僅沒有\Quote{跟隨耶穌}，甚至對祂心懷嫉妒。\Quote{辣彼，}有一個人說，\Quote{從前同你在約但河外、你所見證的那位，現今施洗，眾人都到他那裡去了。}\BibleRef{若 3:26}他們又似乎向他提出指控：\Quote{為什麼我們禁食，你的門徒卻不禁食呢？}\BibleRef{瑪 9:14}但那些比其餘更優秀的門徒卻沒有這種感覺，他們聽見後立刻跟隨；跟隨，並非輕視自己的老師，而是因他而被最完全地說服，並提供最有力的證據，表明他們這樣做是出於對他理智的正確判斷。因為他們不是聽從他的建議而行，那樣可能顯得可疑；而是當他只預言將要發生的事，即\Quote{祂要用聖神［和火］施洗}時，他們就跟隨了。他們並沒有離棄自己的老師，而是渴望學習基督帶來的比若翰更多的事。請注意熱誠與謙遜的結合。他們沒有立刻上前向耶穌詢問必要和最重要的問題，也沒有急於在眾人面前公開與祂交談，而是私下進行；因為他們知道老師的話不是出於謙卑，而是出於真理。\par

% structure: p0015 source=h2 target=subsection reason=relative-heading-rank; base=h2

\subsection{若望福音 1:40}

那么，为什么他没有也将另一位门徒的名字显明呢？有人说，因为那跟随者就是作者本人；另有人说并非如此，而是因为他不是那些出类拔萃的门徒之一；因此不宜多提。因为如果作者连那七十二门徒的名字都没有提及，那么知道他的名字又有什么益处呢？圣保禄也做了同样的事。他说：\BibleQuote{我们曾打发同他一起的那位兄弟，他在福音上得称赞}（他曾在许多事上积极努力）。\BibleRef{格后 8:18}。此外，他提及安德肋是出于另一个原因。是什么原因呢？就是当你得知西满与他一同听见\BibleQuote{跟随我，我要使你们成为渔人的渔夫}\BibleRef{玛 4:19}，并没有因这奇异的应许而感到困惑，你就可以知道他的兄弟已经在他内奠定了信德的初基。\par

% structure: p0017 source=h2 target=subsection reason=relative-heading-rank; base=h2

\subsection{\BibleRef{若 1:38}}

由此我们学到，天主并不以祂的恩赐来阻止我们的意愿，而是当我们开始、当我们提供意愿时，祂就赐给我们许多救恩的机会。\BibleQuote{你们找什么？}这是怎么回事？祂洞察人心，居住在我们的意念中，难道祂需要问吗？祂确实问了，并非为了得到信息（这怎么可能呢？），而是借此问题使他们更亲近，赋予他们更大的勇气，并显示他们堪受祂的倾听；因为很可能他们会羞怯害怕，既因不为人知，又因从老师的见证中听到了有关祂的如此描述。因此，祂为了消除这一切，即他们的羞愧和恐惧，就询问他们，不让他们一言不发地一路走到屋子。然而，即使祂不问，结果也会一样：他们会继续跟随祂，跟着祂的脚步到达祂的住所。那么祂为何要问呢？是为了达到我所说的目的：安抚他们仍被羞愧和焦虑搅扰的心灵，并给予他们信心。\par

他们不仅通过跟随显示出热切的愿望，也通过提问显示出来：因为那时他们还未从祂那里学过或听过任何东西，就称祂为\BibleQuote{师傅}，仿佛将自己挤入门徒之中，并宣告他们跟随的原因，为能听到一些有益的事。也请注意他们的智慧。他们没有说\Quote{请教导我们您的教义，或其他我们需要知道的事}，而是说什么？\BibleQuote{你住在哪里？}因为，正如我之前所说，他们希望安静地向祂说些什么，并听祂说些什么，以便学习。因此，他们没有拖延此事，也没有说\Quote{我们明天一定来，听你公开讲道}，而是表明了他们对听祂讲话的巨大渴望，甚至时辰也没有使他们回头，因为太阳已经快要下山了（若望说：\BibleQuote{那时大约是第十时辰}）。因此，基督没有告诉他们祂住处的标志或位置，反而通过显示祂已接纳他们来吸引他们跟随。为此，祂没有对他们说\Quote{现在进屋子不是时候，明天如果你们愿意就可以听，现在回家吧}，而是像对待朋友和长久与祂在一起的人那样与他们交谈。\par

那么，祂在另一处怎么说\BibleQuote{人子没有枕头的地方}\BibleRef{路 9:58}，而这里祂说\BibleQuote{你们来看}\BibleRef{若 1:39}我住的地方呢？因为\BibleQuote{人子没有枕头的地方}这句话的意思是，祂没有自己的住所，并非祂不住在房子里。这也是比较的含义。圣史提到\BibleQuote{他们那一天与祂同住}，但没有说明原因，因为原因很明显：他们跟随基督，基督吸引他们，除了为了得到教诲之外，没有其他动机；而他们在短短一夜之中如此丰盛而热切地享受了这教诲，以致他们立刻前去捕捉他人。\par

4. 那么，让我们也由此学会，视一切事物次于聆听天主的圣言，不认为任何时刻不适合，即使一个人可能需要进入别人的家，作为一个不相识的人向大人物自我介绍，即使天色已晚，或任何时刻，也决不忽视这交易。让食物、洗澡、宴席和此生其余事物有它们指定的时间；但让天上哲学的道理没有单独的时间，让每个季节都属于它。因为保禄说：\BibleQuote{不论时机是否适宜，务必劝勉、斥责、鼓励}\BibleRef{弟后 4:2}；先知也说：\BibleQuote{他昼夜默思祂的法律}\BibleRef{咏 1:3}；梅瑟也曾命令犹太人时常这样做。因为今生的事物，我是说洗澡和宴席，即使有必要，但不断重复会使身体虚弱；然而灵魂的道理，时间越长，就越使接受它的灵魂强壮。但现在我们将所有的时间都分配给了琐事和无益的瞎扯，我们在早晨、下午、甚至中午和晚上都闲坐在一起，并为此指定了地点；但每周听两次或三次神圣的道理，我们就感到厌烦和完全饱足。原因何在？我们的灵魂处于糟糕的状态；它渴望并追求这些事物的能力已完全松懈。因此，它没有足够的胃口来渴求精神食粮。而这一点，加上其他方面，是不饥不渴的软弱的显著证据。如果这在身体上发生，是严重疾病的确切迹象，并产生虚弱，那么灵魂中更是如此。\par

\Quote{那么，}有人问，\Quote{我们如何能使这已如此衰弱松弛的（灵魂）恢复力量呢？该做什么、该说什么呢？}借着专心研读先知、宗徒、福音及其他神圣的言语，我们便会知道，以此滋养远胜于不洁的食物——因为我们不得不这样称呼那些不合时宜的空谈和集会。请告诉我，哪一样更好：是谈论市场、法庭或军营中的事，还是谈论天上的事，以及我们离世之后将要发生的事？哪一样更好：是谈论邻人及其事务、操心别人的事，还是探究天使的事以及与我们自身相关的事？邻人的事与你毫不相干，而天上的事却是你的。\newline{}\newline{}\Quote{但是，}有人说，\Quote{人只要一次开口，就能把这些话题全部说尽。}为什么你在那些无益而空洞的交谈中却不这样想？为何你耗费一生于此，却从未穷尽这些话题？我还没有提到比这更卑劣的事。以上是较好的人彼此交谈的内容；而那些冷漠粗心的人则在言谈中谈论戏子、舞者和战车御者，玷污人的耳朵，败坏人的灵魂，借着这些叙述将他们的本性驱入疯狂的放纵，并通过这样的谈话将各种邪恶引入自己的想象。因为舌头一说出舞者的名字，灵魂立刻就勾勒出他的容貌、头发、精致的衣着，以及他比任何人都更娇媚的姿态。另一个人则以另一种方式扇起火焰，在谈话中引入某个娼妓，描述她的话语、姿态、眼神、慵懒的目光、卷曲的发丝、光滑的面颊和涂色的眼睑。当我这样描述时，你们难道没有些许触动吗？但不要羞愧，也不要脸红，因为本性的必然要求如此，它根据所谈论的内容的倾向来支配灵魂。然而，如果当我说话时，你们站在教堂中，远离这些事物，尚且因听闻而有所触动，那么想想那些坐在剧院里、毫无畏惧、远离这神圣可畏的集会、并以极大的无耻看见并听见这些事物的人，他们又会如何呢？\newline{}\newline{}\Quote{既然如此，}也许某个漫不经心的人会说，\Quote{如果本性的必然如此支配灵魂，你为什么放过那种必然，而指责我们呢？}因为当人听闻这些事时变得软弱，这是本性的作用；但去听闻它们却不是本性的过错，而是出于自由的选择。正如一个碰触火的人必然会受伤，这是由于我们本性的软弱；然而本性并不因此驱使我们走向火和随之而来的伤害；这只能来自故意的悖逆。因此，我恳求你们，要除去并纠正这一过错，免得你们自愿坠入悬崖，把自己推进邪恶的坑中，自行奔向烈焰，以免使我们陷入为魔鬼所预备的火中。但愿我们都能既摆脱这火，也摆脱那火，进入\Person{亚巴郎}的怀抱中，藉着我们的主耶稣基督的恩宠和慈爱，愿光荣归于圣父和圣神，至于无穷之世。阿们。\par

\SourceRecord{Translated by Charles Marriott.；Nicene and Post-Nicene Fathers, First Series；Vol. 14.；Edited by Philip Schaff.；Buffalo, NY: Christian Literature Publishing Co.,；1889.；<http://www.newadvent.org/fathers/240118.htm>.}

% structure: p0001 source=h1 target=section reason=page-title

\section{若望福音讲道 第十九篇}

% structure: p0002 source=h2 target=subsection reason=relative-heading-rank; base=h2

\subsection{若望福音 1:41-42}

1. 天主起初造人时，并不许人独居，而是赐给他女人作助手，使他们同居共处，因为他知道从这种同伴关系中将产生极大的益处。即使女人没有善用这一恩惠，但若有人充分了解这事理的本质，就会看到，对智者而言，同居共处带来极大益处；不仅在于夫妻关系，就是弟兄们若这样行，也必享受益处。因此，先知曾说：\Quote{有什么比兄弟同居共处更美好、更愉快呢？}\BibleRef{咏 132:1} 保禄也劝勉我们不可忽略聚会。\BibleRef{希 10:25} 正是在这一点上，我们与禽兽有别；为此我们建造了城市、市场与房屋，为使我们彼此联合，不仅在于居住的地点，更在于爱德之链。因为我们的本性由造物主造来时并不完善，且不能自足，天主为我们的益处，制定了以此需要借着互相交往的帮助而得以弥补，使一人的不足得由另一人补充，从而缺损的本性得以自足；例如，虽被造为有死者，却可借着延续长久保持不朽。我本可更详尽地阐述此理，以表明那些真诚纯洁交往者之间所得益处；但眼下另有更紧迫之事，即我们之所以说这番话的缘故。\par

安德肋与耶稣同住，得知他的所作所为后，并未将宝藏据为己有，而是急忙跑向他的兄弟，将所领受的好事分给他。但若望为何没有说明基督与他们谈论了什么事？何以知道他们\Quote{与祂同住}正是为此？这一点我们前日已证明；但从今日所读经文同样可以得知。请注意安德肋对他的兄弟所说的话：\Quote{我们找到了默西亚（意即基督）。}您看，他在短时间内所学到的，显明了说服他们的导师的智慧，以及他们自己的热忱——他们早已渴慕这些事，自起初就期待。因为\Quote{我们找到了}这句话，出自一个为他来临而忍受产痛、仰望他从天降下、且因期待之事实现而喜乐洋溢的灵魂，急切地将好消息传给他人。这就是手足之情、天然友谊、真诚性情的一部分：渴望在属神之事上彼此援手。此外，请注意他说话时带着冠词；他说的不是\Quote{默西亚}，而是\Quote{那默西亚}；可见他们期待的是那一位基督，与别的不相干。看哪，我恳求您，伯多禄的心从起初就顺服而易于受教；他毫不迟疑地跑到耶稣那里；圣若望说：\Quote{他带他到耶稣那里。}但不要有人责怪他性情轻信，因为他兄弟可能已更精确详尽地告诉了他这些事；只是圣史们为了简练，常删去许多细节。此外，经文并没有直接说他\Quote{相信}，而是说他\Quote{带他到耶稣那里}，为将来把他交托给祂，以便从他那里学到一切；因为另一位门徒也与安德肋同来，并为此作了协助。洗者若翰曾说祂是\Quote{羔羊}，并且祂\Quote{以圣神施洗}，就让他们从祂那里学习关于此事的更明确教理，那么安德肋更会如此行，不认为自己足以宣告全部，而是带着如此热忱与喜乐，将他引到光明的源头，以致另一位毫无迟延。\par

% structure: p0005 source=h2 target=subsection reason=relative-heading-rank; base=h2

\subsection{若望福音 1:42}

2. 祂从这时起开始启示属于祂天主性的事，并借预言逐步展开。祂在纳塔乃耳和撒玛黎雅妇人身上也如此行。因为预言引人归信不逊于奇迹，且免于炫耀之嫌。奇迹在愚人中间可能被毁谤，\BibleRef{玛 12:24} 但关于预言从未有过这类说法。现在对纳塔乃耳和西满，祂用了这种教导方式；但对安德肋和斐理伯，祂却没有这样做。为什么？因为那两位（安德肋和斐理伯）已有洗者若翰的见证作为不小的预备，而斐理伯看到那些曾与耶稣同在的人后，也得到了信德的可信证据。\par

祂说：\Quote{你是西满，若纳的儿子。} 借着现在，将来得到保证；因为那说出伯多禄父亲名字的祂，也预知将来。这预言伴随着赞美；但目的不是奉承，而是预告将来之事。至少听祂对撒玛黎雅妇人的话，是如何以严厉的责备发出预言：祂说：\Quote{你曾有五个丈夫，你现在所有的不是你的丈夫。}\BibleRef{若 4:18} 同样，祂的圣父也极为重视预言，当祂反对偶像崇拜时，祂说：\Quote{让它们向你们宣告将要来临的事}\BibleRef{依 47:13}；又说：\Quote{我曾宣告，我施行了拯救，你们中间没有外邦的神。}\BibleRef{依 43:12}；祂在整个预言中都提出这一点。因为预言尤其是天主的作为，连魔鬼也无法模仿，尽管它们竭力尝试。因为奇迹可能有欺骗；但准确预言将来只属于那纯粹的本性。或者，如果魔鬼曾这样做，也是欺骗了较简单的人；因此它们的预言总是容易识破。\par

但是\Person{伯多禄}没有回答这些话；他那时还不清楚任何事情，仍在学习。请注意，连预言也没有完全陈述出来；因为\Person{耶稣}并没有说：\Quote{我会把你的名字改为\Person{伯多禄}，在这磐石上我要建立我的教会}，而是说：\Quote{你要被称为\OriginalTerm{刻法}{Cephas}}。前一种说法会表现出过大的权威和能力；因为基督并不立刻或一开始就显示他的全部能力，而是暂时以更谦逊的口吻说话；所以当他证明了祂的天主性后，才更权威地说：\BibleQuote{你是有福的，\Person{西满}，因为我的父启示了你}；又说：\BibleQuote{你是\Person{伯多禄}，在这磐石上我要建立我的教会。}\BibleRef{玛 16:17} 因此祂这样命名他，又称\Person{雅各伯}和他的兄弟为\Quote{雷霆之子}。\BibleRef{谷 3:17} 那么祂为什么这样做呢？为显示正是祂赐予了旧盟约，正是祂改变了名字，祂曾将\Person{亚巴郎}称为\Person{亚巴辣罕}，\Person{撒辣依}称为\Person{撒辣}，\Person{雅各伯}称为\Person{以色列}。祂给许多人甚至从出生时就指定名字，如\Person{依撒格}、\Person{三松}，以及\Person{依撒意亚}和\Person{欧瑟亚}中的人\BibleRef{依 8:3}；但给其他人则在父母命名后赐予，如我们提到的那些人，以及农的儿子\Person{若苏厄}。这也是古人的一种习俗，根据事物取名，事实上\Person{肋阿}也这样做了；这并非没有原因，而是为了使人们拥有这个名字来提醒他们天主的仁慈，使名字所传达的预言的永久记忆能在领受者耳中回响。同样，祂也早早就给\Person{若翰}取名，因为那些从幼年起就要发光的人的德行，从那时起就得到了名字；而对于那些在较晚时期才成为伟大的人，这个称号也是在后来才赐予的。\par

3. 但是那时他们各人领受不同的名字，我们现在却都有同一个名字，那是比任何名字都更大的名字，被称为\Quote{基督徒}、\Quote{天主的儿子}、祂的\Quote{朋友}和祂的\Quote{身体}。因为这个名称本身比所有那些其他名称更能激励我们，促使我们更热心于修德。我们被称为基督的人（因为保禄这样称呼我们），看看我们身份的超越，就不应做出与这名字的荣耀不相称的事。让我们记住并尊重这个称呼的崇高。\BibleRef{格前 3:23} 因为如果有人据说是某位著名将领的后裔，或是其他显赫人物的后裔，就骄傲地被称为这人或那人的儿子，并认为这名字是极大的荣耀，且尽全力避免因自己的懈怠而给所自称为后裔的人带来羞耻；那么，我们被称为名字的，不是一位将领，也不是地上任何一位王子，也不是天使、总领天使、色辣芬，而是这些本身的天主，难道我们不应当甚至付出自己的生命，以免羞辱那尊荣了我们的一位吗？你们不知道围绕君王的皇家持盾者和持矛者享有何等荣耀吗？同样，我们这些被相称得以亲近祂，而且比那些人更亲近，就如身体比他们更接近头一样，我说，让我们用尽一切方法成为基督的模仿者。\par

基督说什么？\BibleQuote{狐狸有洞，天空的飞鸟有巢，但是人子却没有枕头的地方。}\BibleRef{路 9:58} 如果我现在向你们要求这一点，大多数你们可能觉得痛苦和沉重；因此，因着你们的软弱，我不说这种成全，而是希望你们不要被财富钉住；如同我因着许多人的软弱，从要求极致的德行中稍作退让，我渴望你们也在恶的方面这样做，甚至更多。我不责备那些有房屋、土地、财富和仆人的人，而是希望他们以安全而适宜的方式拥有这些东西。什么是\Quote{适宜的方式}？作为主人，而不是作为奴隶；使他们统治它们，而不是被它们统治；使他们使用它们，而不是滥用它们。这就是为什么它们被称为\Quote{可使用之物}，我们可以将它们用于必要的服务，而不是囤积它们；这是仆人的职分，而不是主人的；保存财物是奴隶的事，而花费财物是主人和有权势者的事。你们获得财富不是为埋葬，而是为分施。如果天主希望财富被囤积，他就不会把它赐给人，而是让它们留在地上；但因为祂希望它们被使用，所以就允许我们拥有它们，好让我们彼此分享。如果我们将其据为己有，我们就不再是它们的主人。但如果你们想让它们增加，因而将之封存，即使在这种情况下，最好的办法也仍然是把它们分散出去，因为不花出去就没有收入，不支出就没有财富。即便在世俗事务中也可以看到这点。对商人如此，对农夫也是如此：一个付出其财富，另一个付出其种子；一个航海去分售货物，另一个整年劳作，播种和照料。但在这里，这些都不需要：无需装备船只，无需套牛，无需耕田，无需担忧不确定的天气，无需害怕冰雹；这里没有波浪也没有岩石；这趟航行和这次播种只需要一件事：把我们的财富抛出去；其余都将是农夫所做的，正如基督所说：\BibleQuote{我的父是农夫。}\BibleRef{若 15:1} 那么，难道不是荒谬的吗，当我们可以在不劳而获的地方获得一切时却懒惰懈怠，而在有许多辛劳、烦恼和忧虑，且希望不确定的地方，却表现出全部的热忱？我恳求你们，不要对我们的得救如此麻木，而是让我们离开更麻烦的任务，奔向那更容易且更有益的，好使我们也能获得将来的美善；因我们的主耶稣基督的恩宠和慈爱，愿光荣归于祂与圣父及赋予生命的圣神，从今世直到永远。阿们。\par

\SourceRecord{Translated by Charles Marriott.；Nicene and Post-Nicene Fathers, First Series；Vol. 14.；Edited by Philip Schaff.；Buffalo, NY: Christian Literature Publishing Co.,；1889.；<http://www.newadvent.org/fathers/240119.htm>.}

% structure: p0001 source=h1 target=section reason=page-title

\section{讲道集第二十篇 论若望福音}

% structure: p0002 source=h2 target=subsection reason=relative-heading-rank; base=h2

\subsection{若 1:43-44}

1. 箴言说：\Quote{对于一个审慎的思想家，每一件事都有其益处。} \BibleRef{箴 14:23}，基督更强调说：\Quote{寻求的，就会找到。}\BibleRef{玛 7:8}。因此，我不再诧异斐理伯为何跟随基督。安德肋听了若翰的话就信了，伯多禄也因安德肋而信；但斐理伯并非从任何人那里学到什么，而只听了基督说的\Quote{跟随我}，便立刻服从，没有退缩，甚至还向别人宣讲。因为他跑去见纳塔乃耳，对他说：\Quote{我们找到了梅瑟在法律中和先知们所记载的那一位。}你看见他多么有心，勤于默想梅瑟的著作并期待基督的来临吗？因为\Quote{我们找到了}这句话，总是属于那些正在寻求的人。\Quote{次日，耶稣往加里肋亚去。}在无人跟随祂之前，祂没有召叫任何人；祂这样做并非无故，而是出于自己的智慧。因为如果当没有人主动来跟随祂时，祂就亲自吸引他们，他们可能会退缩；但现在他们自己选择了，后来就坚定地留住了。祂召叫斐理伯，他是比较熟悉祂的人；因为斐理伯生于加里肋亚，长于加里肋亚，比别人更认识祂。收了门徒之后，祂接着去捕获其他人，吸引了斐理伯和纳塔乃耳。\newline{} 对于纳塔乃耳来说，这并不那么奇妙，因为耶稣的名声已经传遍了全叙利亚。\BibleRef{玛 4:24}但奇妙的是关于伯多禄、雅各伯和斐理伯，他们在未有神迹之前就信了，而且他们来自加里肋亚，那里\Quote{不出先知}，\Quote{什么好事也不能出}；因为加里肋亚人似乎比其他人更粗野、更迟钝；但基督在其中彰显了祂的大能，从一片不结果实的土地上挑选了祂最精选的门徒。因此，斐理伯可能看到了伯多禄和安德肋，并听到了若翰的话，就跟随了；也可能基督的声音在他身上起了作用；因为祂知道哪些人会有用。但福音作者简略了所有这些要点。斐理伯知道基督要来，却不知道这一位就是基督；我这样说，是因为他可能从伯多禄或若翰那里听到过。但若望也提到他的村庄，好让你知道\Quote{天主拣选了世上软弱的}。\BibleRef{格前 1:27}\par

% structure: p0004 source=h2 target=subsection reason=relative-heading-rank; base=h2

\subsection{若 1:45}

他说这话，是为了使他的宣讲可信——它必须建立在梅瑟和先知之外，并以此使听众感到羞愧。因为纳塔乃耳是一个严谨的人，以真理看待一切，正如基督所证实的，事件也表明这一点；因此斐理伯明智地把他引向梅瑟和先知，好让他接受所宣讲的那一位。他并不因称祂为\Quote{若瑟之子}而困扰；因为当时人们仍以为祂是若瑟的儿子。\Quote{斐理伯，从哪里看得出来这就是祂呢？你向我们提出什么证据呢？因为仅仅断言是不够的。你看见了什么征兆，什么奇迹？在这样的问题上没有原因地相信是有危险的。那么你有什么证据呢？}\Quote{安德肋所给你的证据也一样，}他回答说；因为他虽然无法展现他所找到的财富，也无法用言语描述他的宝藏，但当他发现时，就把他的兄弟带到了那里。斐理伯也是如此。关于这位如何是基督，以及先知们如何预先宣告了祂，他没有说；但他把纳塔乃耳带到耶稣跟前，知道只要他一旦尝到祂的言语和教导，就不会再离开。\par

% structure: p0006 source=h2 target=subsection reason=relative-heading-rank; base=h2

\subsection{若 1:46-47}

他称赞并认可那人，因为他说了\Quote{纳匝肋能出什么好事吗？}然而他本应受到责备。当然不是；因为这话并非出于不信，也不是该受责备，而是该受称赞。\Quote{这是为什么？如何呢？}因为纳塔乃耳比斐理伯更熟悉先知书。他从圣经中听到基督必须来自白冷，来自达味所在的村庄。犹太人普遍有这个信念，先知早已宣告说：\Quote{你白冷啊，你在犹大诸城中绝不是最小的，因为将有一位统治者从你出来，牧养我的以色列民。}\BibleRef{玛 2:6}因此当他听说祂是\Quote{从纳匝肋来的}时，他就困惑、怀疑，因为斐理伯的宣告与先知的预言不符。\par

但请观察，即使在他怀疑时，他的智慧和坦诚。他没有立刻说：\Quote{斐理伯，你欺骗我，你说谎，我不相信你，我不会去；我从先知那里学到基督必须来自白冷，而你说\InnerQuote{从纳匝肋来}；}他没有说这样的话；而是怎样呢？他亲自去见祂；他不承认基督是\Quote{从纳匝肋来的}，显示出他对圣经的精确性和不易受骗的性格；而他没有拒绝报信的人，则显示出他对基督来临的极大渴望。因为他心里想，斐理伯可能弄错了地方。\par

2. 且请留意，我请求你们，祂是如何委婉地拒绝，以提问的方式使之温和。因为祂没有说：\Quote{加里肋亚不出产任何好东西}；但祂怎么说？\BibleQuote{纳匝肋能出什么好事吗？}斐理伯也很明智；因为他并不显得困惑、恼怒或烦恼，而是坚持劝导那人，从最初宣讲时就向我们显示了宗徒应有的坚定。因此基督也说：\BibleQuote{看，这确是一个以色列人，在他内毫无诡诈。}所以，确有假以色列人；但这人不是；因为基督说，他的判断是公正的，他说话不偏袒，也不怀恶意。然而，犹太人被问及基督应在何处出生时，回答说：\BibleQuote{在白冷}\BibleRef{玛 2:5}，并引用证据说：\BibleQuote{而你，白冷，在犹大的首领中绝不是最小的}\BibleRef{米 5:2}。他们未看见祂之前作此见证，但看见祂后却因恶意隐藏了证词，说：\BibleQuote{至于这人，我们不知道祂从哪里来}\BibleRef{若 9:29}。纳塔乃耳却不如此，他始终持守起初的看法，认为祂不是\Quote{从纳匝肋来的}。\par

那么先知们为何称祂为纳匝肋人？是因为祂在那里长大并居住。祂省略不说：\Quote{我不是斐理伯所告诉你的\InnerQuote{从纳匝肋来的}，而是从白冷来的}，免得立即显得说法可疑；此外，即使祂被相信，也不足以证明祂就是基督。因为若祂不是基督，有什么阻止祂像其他出生在那里的人一样来自白冷呢？因此祂省略了这一点；但祂做了最能说服他的事，因为祂表明当他们在交谈时祂就在场。因为当纳塔乃耳说了：\par

% structure: p0011 source=h2 target=subsection reason=relative-heading-rank; base=h2

\subsection{若 1:48}

要留意一个坚定而稳重的人。当基督说：\BibleQuote{看，这确是一个以色列人，}他并未因这称赞而自满，未追逐公开的赞美，而是继续更准确地寻求和探寻，渴望确知某事。他仍像对一个人询问，但耶稣如同天主回答。因为祂说：\BibleQuote{我从起初就认识你}（指他的性格和坦诚，祂不是作为人从密切跟随认识，而是作为天主从起初就认识），\BibleQuote{但方才我在无花果树下看见了你}；当时那里只有斐理伯和纳塔乃耳，他们私下说了这一切。经上记载，祂从远处看见他，就说：\BibleQuote{看，这确是一个以色列人}；为表明在斐理伯走近之前，基督就说了这些话，使这见证不被怀疑。为此祂也提到了时间、地点和树；因为如果祂只说：\Quote{斐理伯到你那里之前，我看见了你}，祂可能会被怀疑是派了他，并说了什么不起的事；但现在，通过提到当斐理伯对他说话时他在哪里，以及树名和谈话的时间，祂表明祂的预知是无可置疑的。\par

祂不仅向他显示了预知，还用另一种方式教导了他。因为祂使他回忆起他们当时所说的话，例如：\BibleQuote{纳匝肋能出什么好事吗？}尤其是因为这个原因纳塔乃耳接受了祂，因为当他说了这些话时，祂并没有责备，反而称赞并认可了他。因此他确信这确实是基督，既因祂的预知，也因祂精确探知了他的感受，这乃是那位表明自己知道他想什么者的作为；此外，还因祂没有责备，反而在他似乎说了反对祂的话时称赞了他。于是祂说，斐理伯曾\Quote{叫}了他；但斐理伯对他说了什么或他对斐理伯说了什么，祂没有提，留给他的良心，不愿进一步责备他。\par

3. 那么，是否只有在\Quote{斐理伯叫他之前}祂才\Quote{看见}了他？难道祂以前没有用祂不眠的眼睛看见他吗？祂看见了他，无人能否认；但这是当时需要说的。纳塔乃耳做了什么呢？当他获得了祂预知的确凿证据时，他赶紧承认了祂，以先前的迟延显示了他的谨慎，以之后的赞同显示了他的公正。因为，福音作者说：\par

% structure: p0015 source=h2 target=subsection reason=relative-heading-rank; base=h2

\subsection{若 1:49}

你们看见他的灵魂如何立刻充满了极大的喜乐，并如何以言语拥抱耶稣吗？他说：\BibleQuote{祢就是那期待的、那寻求的一位。}你们看见他是如何惊讶、如何赞叹吗？他如何欣喜跳跃、手舞足蹈？\par

因此，我们也应当喜乐，因为我们蒙恩得以认识天主子；不仅要心中喜乐，也要以行动表现出来。那些喜乐的人该做什么呢？就是服从那已向他们启示的主；而服从的人必须做祂所愿意的一切。因为如果我们去做令祂愤怒的事，又怎能表明我们喜乐呢？难道你们没有看到，在我们家中，若有人接待他所爱的人，他是何等欢喜地忙碌，四处奔走，即使需要花尽所有，也毫不吝惜，只为取悦他的客人吗？但如果邀请的人不关心他的客人，不做那些能使客人舒适的事，即使他说一万次他因客人的到来而喜乐，也绝不会被客人相信。这很公平；因为喜乐应以行动证明。那么，既然基督已来到我们中间，就让我们表明我们的喜乐，不做任何可能激怒祂的事；让我们装饰祂所来到的居所，因为喜乐的人正是这样做的；让我们在祂面前献上祂渴望享用的餐食，因为欢庆节日的人正是这样做的。这是什么餐食呢？祂亲自说：\BibleQuote{我的食物就是承行派遣我者的旨意。} \BibleRef{若 4:34} 当祂饥饿时，让我们喂饱祂；当祂口渴时，让我们给祂喝：即使你只给祂一杯凉水，祂也接受；因为祂爱你，而对于爱的人来说，所爱之人的奉献，即使微小，也显得伟大。只是不要懒惰；即使你只投入两个铜钱，祂也不会拒绝，反而当作巨大的财富接纳。因为祂一无所缺，接受这些奉献并非因为祂需要，所以区别合理之处不在于礼物的数量，而在于奉献者的心意。只要表明你爱那来临的主，为了祂你竭尽全力，你因祂的来临而喜乐。看祂是如何对待你的。祂为你而来，为你舍弃了自己的生命，在这一切之后，祂甚至不拒绝恳求你。保禄说：\BibleQuote{我们是基督的使者，好像天主借着我们劝勉你们一样。} \BibleRef{格后 5:20} 有一个人说：\Quote{谁这样疯狂，竟不爱自己的主人呢？} 我也这样说，我知道我们中没有一个人会在言语或意念上否认这一点；但被爱的人希望爱不仅用言语，也要用行动来表达。因为口说爱，行为却不似爱人者，不仅在天主面前，即使在世人眼中也是可笑的。既然仅在言语上承认祂，而在行为上反对祂，不仅无益，反而有害于我们；我恳求你们，也以行动做出宣认；这样我们也能在那一天得到祂的宣认，那时祂将在祂的父面前承认那些在基督耶稣我们的主内堪当的人。愿光荣归于圣父及圣神，藉着祂并与祂一起，从今世直到永远，永无穷尽。阿们。\par

\SourceRecord{Translated by Charles Marriott.；Nicene and Post-Nicene Fathers, First Series；Vol. 14.；Edited by Philip Schaff.；Buffalo, NY: Christian Literature Publishing Co.,；1889.；<http://www.newadvent.org/fathers/240120.htm>.}

% structure: p0001 source=h1 target=section reason=page-title

\section{圣若望福音讲道 第二十一篇}

% structure: p0002 source=h2 target=subsection reason=relative-heading-rank; base=h2

\subsection{若望福音 1:49-50}

1. 亲爱的，我们需要许多谨慎、许多警醒，才能深入圣经的深奥之处。因为不可能以粗心或沉睡的方式发现其含义，而需要仔细搜寻和恳切祈祷，使我们得以稍微窥见神圣言语的秘密。例如今天，摆在我们面前的绝非小事，而是需要极大热忱和探究的问题。当纳塔乃耳说：\BibleQuote{你是天主子，} 基督回答说：\BibleQuote{因为我对你说，我看见你在无花果树下，你就信吗？你将看见比这些更大的事。}\par

那么从这段经文产生的问题是什么？就是这一点。伯多禄在如此多的奇迹和高深道理之后，宣认\BibleQuote{你是天主子}（\BibleRef{玛 16:16}），被称为\Quote{有福的}，因为是从圣父领受了启示；而纳塔乃耳虽在未见或未听奇迹和道理之前就说了同样的话，却没有得到这样的话语，反而好像他没有说出应有的那么多，被引导到更大的事。这是为什么呢？这是因为伯多禄和纳塔乃耳说了同样的言辞，但意图不同。伯多禄宣认祂是\BibleQuote{天主子}，但作为真天主；纳塔乃耳则作为纯粹的人。这从哪里显明呢？从他随后所说的话；因为说了\BibleQuote{你是天主子}之后，他又加上\BibleQuote{你是以色列君王。}但天主子不只是\Quote{以色列君王}，而是全世界的君王。\par

我所说的话不仅从这一点，也从后续的事上显明。因为基督没有对伯多禄再加什么，反而好像他的信德已经完备，说他要在这宣认之上建立教会；但在另一情况下，祂并未如此行，反而相反。因为好像他的宣认还缺少某些大的、更好的部分，祂就加上了以下的话。因为祂说了什么？\par

% structure: p0006 source=h2 target=subsection reason=relative-heading-rank; base=h2

\subsection{若望福音 1:51}

你看祂如何逐步将纳塔乃耳从地上引领，使他不再仅仅认为祂是一个人？因为一位有天使服事、天使在祂身上上下往来的人，怎能是人呢？为此祂说：\BibleQuote{你将看见比这些更大的事。}为证明这一点，祂提出了天使的服事。祂的大意是这样：\BibleQuote{纳塔乃耳，这事在你看来很大，你就因此承认我是以色列君王吗？那么，当你看见天使在\InnerQuote{我}身上上上下下时，你将要说什么？}借此说服他，也要承认祂是天使的主。因为天使这御前侍从在祂——君王的亲子身上上下往来，一次在受难之时，又一次在复活和升天之时，并且在这之前，当他们\BibleQuote{前来服事祂}（\BibleRef{玛 4:11}），当他们宣告祂诞生的喜讯，高唱\BibleQuote{光荣归于天主高天，平安于地}（\BibleRef{路 2:14}），当他们来到玛利亚、来到若瑟那里。\par

祂现在所做的，正如在许多事例中所做的；祂说出两个预言，对其中一个给予当下的证明，并用已应验的来证实那将要成就的。因为祂的话有的已经证明，如\BibleQuote{在斐理伯叫你之前，你在无花果树下，我看见了你}；另一些尚未应验，但也部分应验了，即天使的降下和上升，在受难、复活和升天时；祂以祂的话语在事前就使其可信。那借着先前的事认识祂权能的人，又听祂说未来的事，会更乐意接受这预言。\par

那么纳塔乃耳做什么呢？他对这话没有回答。因此基督在此停止与他的对话，让他私下思考所说的话；并避免一下子倾注所有，既然已经将种子播在好土里，便任其自由生长。祂在另一处也显示了这一点，祂说：\BibleQuote{天国好像一个人撒好种子，但他睡觉的时候，仇敌来把莠子撒在麦子里。}\par

% structure: p0010 source=h2 target=subsection reason=relative-heading-rank; base=h2

\subsection{若望福音 2:1-2}

我先前曾说，祂在加里肋亚最被熟知；因此他们邀请祂赴婚宴，祂就来了。因为祂不看自己的尊荣，而看我们的益处。祂不轻视\BibleQuote{取了奴仆的形体}（\BibleRef{斐 2:7}），更不会轻视参加仆人的婚宴；祂曾\BibleQuote{同税吏和罪人一起坐席}（\BibleRef{玛 9:13}），更不会拒绝与婚宴上的人同坐。显然，邀请祂的人对祂没有正确的判断，也没有把祂当作大人物邀请，而只是当作一位普通熟人；福音作者已暗示了这一点，他说：\BibleQuote{耶稣的母亲在那里，祂的兄弟们也在那里。}正如他们邀请她和祂的兄弟们，他们也邀请耶稣。\par

% structure: p0012 source=h2 target=subsection reason=relative-heading-rank; base=h2

\subsection{若望福音 2:3}

这里值得探究，祂的母亲从哪里想到对祂儿子有伟大想法；因为祂还未行过奇迹，因为福音作者说：\BibleQuote{这是耶稣在加里肋亚加纳所行的第一个奇迹。}（\BibleRef{若 2:11}）\par

2. 若有人说这不足以证明那是\\Quote{祂奇迹的开始}，因为仅仅增加了\\Quote{在加里肋亚的加纳}这一句，只允许它成为在那里行的第一个奇迹，但并非绝对的第一，因为祂或许在别处行了别的奇事，我们将以先前已经说过的来回答他。是怎样的回答呢？若翰的话：\\BibleQuote{我也不曾认识祂，但为叫祂显示给以色列，为此我来用水施洗。}如果祂在早年就行过奇迹，以色列人就不需要别人来宣告祂了。因为那来到人间，并藉着自己的奇迹如此被传扬的祂，不仅被犹太地的人所知，也被叙利亚和境外的人所知，并且这只是在三年之内完成的，甚至不需要这三年来自我显示\\BibleRef{玛 4:24}，因为立刻并从起初，祂的名声就传遍各处。祂，我在说，祂在短时间内因着众多奇迹如此闪耀，以致祂的名被所有人熟知，那么，如果祂小时候从幼年起就行奇迹，就更不可能被隐藏这么久。因为作为一个孩子所行的事会显得更奇特，并且会有两倍或三倍的时间，甚至更多。但实际上，祂在孩童时期除了路加所见证的那一件事之外，什么也没有做\\BibleRef{路 2:46}：十二岁时，祂坐着听经师们的教导，并因祂的提问令人钦佩。此外，这也符合情理：祂没有从幼年立刻开始行奇迹；因为人们会认为那是幻觉。因为当祂成年时，许多人尚且怀疑这一点，更何况如果祂在十分年幼时就显奇迹，他们会在恶意的毒怒中更早、不合时宜地将祂推向十字架；而救恩计划的事实本身也会被质疑。\par

\\Quote{怎么，}有人问，\\Quote{祂的母亲怎么会想到祂有什么伟大之处呢？}祂现在开始显示自己，并通过若翰的见证以及祂对门徒所说的话清楚地显明出来。在此之前，受孕本身及其所有相关情境已使她对这孩子抱有极高的看法；因为，路加说，\\Quote{她把一切关于这孩子的话都听见了，并存在心里。}\\Quote{那么为什么，}有人说，\\Quote{她以前不说这话呢？}因为，如我所说，现在终于到了祂开始显示自己的时候。在此之前，祂生活如同众人中的一个，因此祂的母亲没有信心对祂说这样的话；但当她听说若翰是为祂而来，并如此为祂作证，而且祂有了门徒，此后她便有了信心，叫了祂来，在缺酒时说：\\BibleQuote{他们没有酒了。}因她既想帮他们的忙，又想通过她的儿子使自己更为显眼；也许她也有一些人情的想法，如同祂的弟兄们说\\InnerQuote{祢将自己显示给世界}\\BibleRef{若 17:4}时那样，希望从祂的奇迹中获得声誉。因此祂回答得有些严厉，说：\par

% structure: p0016 source=h2 target=subsection reason=relative-heading-rank; base=h2

\subsection{\\BibleRef{若 2:4}}

为证明祂极其尊敬祂的母亲，请听路加如何叙述祂\\Quote{服从}祂的父母\\BibleRef{路 2:51}，以及我们这位圣史如何记述祂在受难时还挂念着她。因为当父母在属于天主的事上不造成阻碍或妨害时，我们有义务顺服他们，不这样做就有大危险；但当他们不合时宜地要求某事，并在任何属灵的事上造成阻碍时，服从便不安全了。因此，祂在这里如此回答，又在别处说：\\BibleQuote{谁是我的母亲？谁是我的弟兄？}\\BibleRef{玛 12:48}，因为他们那时对祂的看法还不正确；而她因生了祂，便按照其他母亲的习惯，要求在一切事上指挥祂，而她本应尊敬和崇拜祂。这就是祂当时那样回答的原因。试想：当所有民众无论尊卑都围在祂身边，群众聚精会神听祂讲道，祂的教导已经开始展开时，她却来到中间，把祂从劝言的工作中带走，与祂私下交谈，甚至不愿进入里面，只把祂拉到外面到自己身边。因此祂说：\\BibleQuote{谁是我的母亲？谁是我的弟兄？}这并不是侮辱生了祂的那一位（断乎不可！），而是为了给她最大的益处，不让她对祂存有卑下的看法。因为如果祂关心别人，并想尽一切办法在他们心中植入对祂适当的看法，那么对祂的母亲更应如此。既然很可能，如果这些由她的儿子直接对她讲，她那时也不容易接受，反而在一切事上因是祂的母亲而要求优越地位，因此祂对说话的人这样回答；否则，如果她期望祂总是以儿子的身份尊敬她，而不是以师傅的身份来到，祂就无法将她的思想从现在的卑微引导到将来的高举。\par

3. 祂正是出于这个动机，在此处说：\Quote{女人，这于我和妳有什么关系？}，此外还有一个同样迫切的原因。那是什么呢？就是不让祂的奇迹受到猜疑。请求应当来自那些需要的人，而不是祂的母亲。为何如此？因为若出于亲友的请求，即使事迹伟大，也常令旁观者反感；但若由有需要的人自己请求，奇迹便免于猜疑，赞美纯粹，益处巨大。假设一位高明的医生进入一间有许多病人的屋子，没有一位病人或他们的亲属向他说话，却只由他的母亲指示，他必会遭到患者的怀疑和厌恶，谁也不认为他能行什么大而惊人的事。因此，祂在那时责备她，说：\Quote{女人，这于我和妳有什么关系？}，教导她以后不可再如此行；因为虽然祂谨慎地孝敬母亲，但祂更关心她灵魂的得救，以及为大众的益处，祂正是为此取了血肉之躯。\newline{}\par

这些话并非粗鲁地对母亲说的，而是属于一项智慧的安排，为使她进入正确的心理状态，并使奇迹得到应有的尊荣。且不说别的，单单这话带有责备口气的表象，就足以表明祂非常尊敬她；因为祂的不悦正显示了对她的极大敬意；其方式我们将在下篇讲道中说明。请深思这一点，并且当你听到某个妇人说：\Quote{怀过祢的胎，及祢所吮吸过的乳房，是有福的}，而祂回答说：\Quote{可是，那听从我父旨意的，更是有福的}\BibleRef{路 11:27}，你就应想到，前面那些话也是出于同样的用意。因为那回答并非拒绝母亲，而是为表明她生了祂，若她本身不是极善良忠信，便毫无益处。如今，即使撇开她灵魂的卓越，玛利亚因基督由她所生也无所得益，那么，若我们自己远离了德行，即使有德行高尚的父亲、兄弟或儿女，对我们又能有何益处呢？达味说：\Quote{兄弟不能赎人，人岂能赎人？}\BibleRef{咏 48:7}。我们必须将得救的希望寄托于别的东西，而只在于我们在天主的恩宠之后所行的义德。因为如果单凭这一点就能奏效，那么它早就对犹太人有效了（因为基督按肉身是他们的亲属），也会对祂出生的城镇有效，也会对祂的兄弟们有效。然而，只要祂的兄弟们不照顾自己，亲属关系的尊荣就对他们毫无益处，他们反而与世人一同被定罪；只有当他们在自己的德行上发光时，才被认可；而那座城毁灭了，被焚烧了，从这关系上一无所得；祂按肉身的亲属也被屠杀，悲惨地灭亡，未能因与祂的关系而得救，因为他们没有德行的庇护。相反，宗徒们却比任何人都伟大，因为他们跟随了与祂建立关系的真实而卓越的道路，即通过服从的道路。由此我们学到，我们总是需要信德，以及光明闪耀的生命，因为只有这才能拯救我们。因为虽然祂的亲属一度到处受尊崇，被称为主的亲属，但如今我们连他们的名字也不知道，而宗徒们的生活和名字却处处被传颂。\newline{}\par

因此，我们不要因肉身的血统高贵而骄傲；即使我们有千万位著名的祖先，我们也应亲自努力超越他们的德行，因为我们知道，在将来的审判中，他人的努力对我们毫无帮助；反而这更成为严厉的定罪：既然生于义德的父母，又有家中的榜样，却仍不效法我们的教师。我现在说这话，是因为我看见许多外教人，当我们引领他们信德、劝他们成为基督徒时，他们便投靠自己的亲属、祖先和家族，说：\Quote{我所有的亲属、朋友和同伴都是虔信的基督徒。}这于你何干呢？你这可怜不幸的人！正是这一点特别要毁掉你，因为你没有尊重你周围那么多的人，并奔向真理。另外一些信徒，但生活懈怠，当被劝勉行善时，也作同样的辩护，说：\Quote{我的父亲、祖父和曾祖父都是非常虔诚善良的人。}但这必然最定罪你，因为你是这些人的后裔，却行事不配你所出的根。因为请听先知对犹太人所说的话：\Quote{以色列曾为娶妻而服役，又为娶妻而牧羊}\BibleRef{欧 12:12}；基督又说：\Quote{你们的父亲亚巴郎曾欢欣踊跃地希望看到我的日子，他看见了，且喜欢。}\BibleRef{若 8:56}。他们处处举出他们祖先的义行，不仅为了赞美他们，更是为了加重对后代的指控。知道了这些，让我们用尽一切办法好使我们能因自己的行为而得救，免得我们因虚假地依靠他人而欺骗自己，等到知道受骗时，这知识已毫无益处。达味说：\Quote{在坟墓里，谁会称谢祢？}\BibleRef{咏 6:5}。让我们在此悔改，好获得永恒的福乐。愿天主恩赐我们众人如此，赖我们的主耶稣基督的恩宠和慈爱，祂与圣父及圣神，永受光荣，于无穷世。阿们。\newline{}\par

\SourceRecord{Translated by Charles Marriott.；Nicene and Post-Nicene Fathers, First Series；Vol. 14.；Edited by Philip Schaff.；Buffalo, NY: Christian Literature Publishing Co.,；1889.；<http://www.newadvent.org/fathers/240121.htm>.}

% structure: p0001 source=h1 target=section reason=page-title

\section{论若望福音第二十二篇讲道}

% structure: p0002 source=h2 target=subsection reason=relative-heading-rank; base=h2

\subsection{\BibleRef{若 2:4}}

1. 宣讲圣言是有些劳苦的，保禄在说：\BibleQuote{那些善于管理的长老，应被认为配受加倍的尊敬，尤其是那些在言语和道理上劳苦的人。} \BibleRef{弟前 5:17} 然而，你们有能力使这劳苦变得轻松或沉重；因为如果你们拒绝我们的话，或者虽不拒绝却不在行为上实践，我们的劳苦就沉重，因为我们徒劳无益地劳作；而如果你们留心听并藉行为证明，我们甚至不会感到劳苦，因为劳作所结的果实不会让那劳苦的宏大显明出来。所以，如果你们要激发我们的热忱，而不是熄灭或削弱它，我恳求你们，向我们显示你们的果实，使我们看见麦浪翻滚的田地，因丰收的希望而得到支持，并数算你们的财富，就不会在这美好的交易上懈怠。\par

今天给我们提出的也是一个不小的难题。首先，当耶稣的母亲说：\BibleQuote{他们没有酒了}，基督回答说：\BibleQuote{女人，我与妳有什么相干？我的时辰还没有到。}然后，说了这话后，祂却按母亲所说的做了；这一行为需要像话语一样去探究。那么，在呼求那行奇迹者之后，让我们开始解释。\par

这些话不仅在这里使用，也在其他地方出现；因为同一福音作者说：\BibleQuote{他们不能下手拿祂，因为祂的时辰还没有到} \BibleRef{若 8:20}；又说：\BibleQuote{没有人下手拿祂，因为祂的时辰还没有到} \BibleRef{若 7:30}；又说：\BibleQuote{时辰到了，请光荣你的子。} \BibleRef{若 17:1} 那么这些话是什么意思？我把更多例子集中起来，为能对所有这些给出一个解释。那解释是什么？基督说\BibleQuote{我的时辰还没有到}，并不是作为受季节必要性或遵守\Quote{时辰}的约束；祂怎能如此呢？祂是季节的制造者、时间和年代的创造者。那么祂指的是什么呢？祂想要表明这一点：祂按合宜的时辰作一切事，并非一次全做；因为如果祂不在合宜的时辰做，而是把祂的出生、复活和审判的来临全部混在一起，就会造成一种混乱。请看：创造要发生，但不是一次全部；男人和女人要受造，但也不是一起；人类要被判处死亡，并要有复活，但两者之间的间隔要大；法律要赐下，但恩宠不与法律同来，每样事物都在其合宜的时间施行。现在，基督不是受时辰必要性的约束，而是建立了它们的秩序，因为祂是它们的创造者；因此祂在此说：\BibleQuote{我的时辰还没有到。} 祂的意思是，祂尚未向许多人显明，祂甚至没有祂的全部门徒；安德肋跟随了祂，接着是斐理伯，但没有其他人。而且，这些人中，甚至祂的母亲和祂的兄弟都不完全认识祂；因为经许多奇迹之后，福音作者论到祂的兄弟说：\BibleQuote{因为连祂的兄弟也不信祂。} \BibleRef{若 7:5} 那些在婚筵上的人也不认识祂，因为他们在需要时肯定会来恳求祂。因此祂说：\BibleQuote{我的时辰还没有到}，也就是说：\Quote{我还不被这些人认识，他们甚至不知道酒已用尽；让他们先意识到这一点。我不应该从妳这里得知；妳是我的母亲，使奇迹显得可疑。需要酒的人应该来求我，并不是我需要这样，而是为使他们能以完全的同意接受奇迹。因为知道自己需要的人，在得到帮助时会非常感激；但那些没有意识到自己需要的人，永远不会对恩惠有清楚明确的感觉。}\par

那么为什么祂在说了\BibleQuote{我的时辰还没有到}并拒绝了她之后，却做了祂母亲所愿的？主要是为了那些反对祂、以为祂受\Quote{时辰}支配的人，能有充足的证据证明祂不受任何时辰的支配；因为如果祂受时辰支配，怎能在适当的\Quote{时辰}到来之前就做祂所做的呢？其次，祂这样做是为了光荣祂的母亲，免得在如此多人面前似乎完全反驳和羞辱生育祂的那一位；并且，为了不被人认为祂缺乏能力，因为母亲把仆人们带到祂那里。\par

此外，甚至在祂对客纳罕妇人说\BibleQuote{拿儿女的饼扔给小狗吃是不对的} \BibleRef{玛 15:26} 之后，祂仍然给了饼，因为考虑到了她的坚持；虽然在第一次回答后，祂说\BibleQuote{我奉差遣，不过是到以色列家迷失的羊那里去}，但说了这话后，祂还是医治了那妇人的女儿。由此我们学到，即使我们不配，常常也能因坚持而使自己配得领受。正因如此，祂的母亲留下来，并公开把仆人们带到祂面前，好让请求由更多人提出；因此她又加上：\par

% structure: p0008 source=h2 target=subsection reason=relative-heading-rank; base=h2

\subsection{\BibleRef{若 2:5}}

因为她知道，祂的拒绝不是出于缺乏能力，而是出于谦卑，并为了避免似乎无缘无故匆忙行奇迹；因此她把仆人们带来。\par

% structure: p0010 source=h2 target=subsection reason=relative-heading-rank; base=h2

\subsection{\BibleRef{若 2:6-7}}

福音作者说\BibleQuote{照犹太人取洁的规矩}，不是没有理由的，而是为使任何不信的人都不怀疑，以为缸底留有酒渣，水倒入和它们混合后就酿成了极淡的酒。因此他说\BibleQuote{照犹太人取洁的规矩}，以表明那些缸从未盛过酒。因为\Place{巴勒斯坦}是一个缺水的地方，溪流和水泉并非到处都有，所以他们常常用水缸装满水，以便万一被玷污时，不必匆忙赶往河边，而是手边就有取洁的用水。\par

\Quote{为何祂不在他们装满水以前行这奇迹，那本应更为奇妙得多？因为将已有的物质改变成另一种性质是一回事，从无中创造物质则是另一回事。}后者固然更为奇妙，但在众人看来却不会如此可信。因此，祂常常故意减弱祂奇迹的伟大，以便更易于被接受。\par

\Quote{但为什么，}有人问，\Quote{祂不亲自产生那水，后来显示为酒，却吩咐仆人们拿来呢？}出于同样的理由；此外，祂也要使那些汲水的人作证，所发生的事绝非幻觉，因为若有人倾向于无耻，那些侍者会对他们说：\InnerQuote{我们汲了水，我们装满了器皿。}除了我们刚才提到的，祂以此推翻了那些兴起反对教会的异说。因为有些人说，世界的造物主是另一位，所见之物并非这另一位造物主的作品，而是某个其他敌对神明的作品；为了遏制这些人的疯狂，祂行的大多数奇迹都是在现成的物质上。因为，如果这些物质的创造者是与祂对立的，祂就不会利用别人的东西来彰显自己的权能。但现在，为了表明正是祂在葡萄树中使水变为酒，正是祂使雨水经由根部转化为酒，祂在婚宴上瞬间成就了植物中需要长久才能完成的事。当他们装满水缸时，祂说，\par

% structure: p0014 source=h2 target=subsection reason=relative-heading-rank; base=h2

\subsection{\BibleRef{若 2:8-10}}

\Quote{这是一群醉汉的聚会，品味者的感官被败坏了，不能品尝所造之物，也不能判断所行之事，以致他们不知道所造的是水还是酒：因为他们醉了，}有人宣称，\Quote{总管本人已凭他所说的话表明了这一点。}这是最荒谬的，然而就连这猜疑，圣史也消除了。因为他并没有说客人对此事发表了意见，而是说\Quote{筵席总管}，他是清醒的，尚未尝过任何东西。因为你们当然知道，那些受托管理此类宴会的人是最清醒的，因为他们唯一的职责就是安排一切井然有序、有条不紊；因此主召这样一个人清醒的感官来见证所发生的事。因为祂并没有说\Quote{给坐席的人倒出来，}而是说\Quote{拿到总管那里去。}\par

\BibleQuote{当筵席总管尝了那变成酒的水，并不知道是从哪里来的（但仆人们知道），总管便叫新郎来。}\Quote{为何他不叫仆人们呢？因为这样奇迹就会显露出来。}因为耶稣没有亲自显明所行的事，而是愿意祂奇迹的能力渐渐地被知晓。假设那时就提及此事，仆人们所说的话绝不会被相信，反而会被人视为疯狂，竟向当时在众人眼中不过是一个普通人的人作这样的见证；尽管他们凭经验确知事情的真实性（因为他们不可能不相信自己的双手），但仍不足以说服他人。所以祂没有向所有人显明，而是向最能理解所行之事的人显明，将更清楚的认识留待将来；因为其他奇迹显现后，此事也会变得可信。因此，当祂将要治愈王臣的儿子时，圣史已经表明此事已更为清楚地为人所知；因为王臣正是因为听说了这奇迹才来求祂，正如若望顺便指出的：\BibleQuote{耶稣来到加里肋亚的加纳，祂曾在那里使水变酒。} \BibleRef{若 4:46}而且不仅仅是酒，还是上好的酒。\par

3. 因为基督的奇迹工作正是如此，它们远比自然运作更完美、更优越。这也在其他事例中可见：当祂治愈身体上任何有病的肢体时，祂使之比健全的更好。\par

当时所造的是酒，且是上好的酒，不仅仆人可以作证，新郎和筵席总管也可作证；而那是基督所造的，汲水的人可以作证；因此，尽管奇迹当时没有显露，但最终不可能被默默忽略，因为祂为将来提供了如此众多而有力的证据。祂使水变酒，有仆人为证；所造的酒是好的，有筵席总管和新郎为证。\par

或许人们会期待新郎对此（总管的讲话）做出回应并说些什么，但圣史因急于处理更紧迫的事项，只略略提及这奇迹，便继续前进了。因为我们需要知道的，是基督使水变成酒，而且是好酒；但新郎对总管说了什么，他认为没有必要添加。许多奇迹起初有些隐晦，随着时间的推移，当那些从一开始就知道它们的人更确切地讲述时，就变得更为清晰了。\par

那时，耶稣便以水变酒；如今祂也不断改变我们软弱不定的意志。因为确实存在一些人，他们在任何方面都与水无异——如此冰冷、软弱、动摇。但让我们把这样性情的人带到主前，使祂将他们的意志改变为酒的品质，叫他们不再稀薄，而是有实体，成为自己和他人的喜乐之源。然而，这些冰冷的人是谁呢？他们是那些心思专注于现世转瞬即逝之事的人，他们不轻视今世的奢华，喜爱荣耀和权势：因为这一切都是流动之水，永不稳定，总是急剧地冲下陡坡。今日的富人明日成为穷人；前一天还带着传令官、腰带、战车和众多随从出现的人，第二天往往成为地牢的居民，不情愿地放弃所有排场为他人腾出空间。同样，贪食和放荡之人，在把自己撑到爆裂后，甚至无法将佳肴所供应的养分保留一日；当养分消散后，为了重新补充，他不得不摄入更多，这与奔流毫无区别。因为正如在急流中，第一批水过去后，另一批依次接替；同样在贪食中，当一餐被消化后，我们又要另一餐。地上事物之本性及命运正是如此：从不稳定，总是倾泻和匆匆流过；但在纵欲方面，不仅仅是流动和匆忙；还有许多其他困扰我们的事。凭借其猛烈冲刷，它磨损身体的力量，剥去灵魂的刚毅；河流最急的水流也不能如此轻易地侵蚀并崩塌其河岸，如同纵欲和放荡扫除我们健康的全部屏障；若你进入医生的家并询问他，你会发现几乎所有疾病的起因都源于此。因为节俭和简朴的餐桌是健康之母，因此医生如此命名；他们把不满足称作\Quote{健康}（因为对食物不满足就是健康），把节省的饮食称为\Quote{健康之母}。如果匮乏的状态是健康之母，显然饱足是疾病和虚弱的母亲，并且产生连医生也无能为力的发作。因为脚上的痛风、中风、视力模糊、手上疼痛、震颤、瘫痪发作、黄疸、缠绵和炎症热病，以及更多其他疾病（我们没有时间一一列举），它们不是节制和适量饮食的自然产物，而是贪食和饱足的自然产物。你若察看由此产生的灵魂的疾病，你会发现贪恋、懒惰、忧郁、迟钝、不洁和各种愚昧的感觉都源于此。因为在此类盛宴之后，纵欲者的灵魂并不比驴子更好，被这些（情欲）猛兽撕碎。还要我说明那些侍奉奢华的人有多少痛苦和不悦吗？我无法全部列举，但通过一个主要点，我能使一切清晰。在我说到的那种餐桌上，即丰盛的餐桌上，人们从未愉快地进食；因为不仅是健康和愉快的母亲是节制，而饱足不仅是疾病之源和根，也是不悦之源。因为哪里有饱足，哪里就不能有欲望；哪里没有欲望，哪里能有愉快呢？因此，我们会发现穷人不仅比富人更明理和健康，而且享受更大程度的愉快。当我们思考这一点时，让我们逃避醉酒和纵欲，不仅是餐桌上的，而且是一切显于现世事物中的；让我们以属于精神事物的愉快来换取它，并一如先知所说，在主内欢欣：\Quote{你要在主内喜乐，祂必满足你心中的愿望}\BibleRef{咏 36:4}；这样，我们就能在此世和来世享受美善，借我们的主耶稣基督的恩宠和慈爱，借着祂并与祂一同，归于父及圣神，光荣永无穷尽。阿们。\par

\SourceRecord{Translated by Charles Marriott.；Nicene and Post-Nicene Fathers, First Series；Vol. 14.；Edited by Philip Schaff.；Buffalo, NY: Christian Literature Publishing Co.,；1889.；<http://www.newadvent.org/fathers/240122.htm>.}

% structure: p0001 source=h1 target=section reason=page-title

\section{若望福音第二十三篇讲道}

% structure: p0002 source=h2 target=subsection reason=relative-heading-rank; base=h2

\subsection{若望福音 2:11}

1. 魔鬼的攻击频繁而猛烈，从四面围攻我们的救恩；因此我们必须警醒、清醒，处处防备牠的攻击，因为只要牠取得了一点点立足之地，就会为自己开辟一条宽阔的道路，逐步引入牠的所有兵力。如果我们对自己的救恩还有丝毫关心，就绝不可让牠在微小的事上接近我们，这样就可以在重要的事上预先遏制牠；因为如果牠如此急切地要摧毁我们的灵魂，而我们却没有付出同等的努力来保卫自己的救恩，那将是极端的愚蠢。\par

我这样说并非无故，而是因为我担心那匹狼如今正不为我们所见地站在羊圈中间，有些羊因自己的疏忽和牠的狡诈而离开羊群、不再听命，成为牠的猎物。倘若伤口可感知，或是身体受到打击，洞察牠的阴谋并不困难；但既然灵魂是看不见的，而受伤的恰恰是它，我们就需要极大的警醒，好让每个人审视自己；因为\BibleQuote{除了人内里的心神外，没有人知道人的事。}\BibleRef{格前 2:11}这话确实是向众人所说的，并作为普遍的药方供给有需要的人，但每一个听者都有责任摄取适合自己病症的。我不知道谁有病，谁健康。因此我用各种论点，并引入适合所有疾病的药方，时而谴责贪婪，时而论及奢侈，再论不洁，然后又写作赞美和劝勉爱德，以及依次论及其他德行。因为我害怕，如果我集中论及某一主题，我会在不知不觉中为你们治疗一种疾病，而你们却患有其他疾病。所以，如果这个会众只有一个人，我就不会觉得必须使我的讲论如此多样化；但既然在这么多人中间很可能有许多病症，我并非不合理地使我的教导多样化，因为当我的讲论涵盖你们所有人时，它必能达到目的。正因如此，圣经也是多方面的，谈论成千上万的事，因为它指向全人类的共同本性，而在如此众多的人群中，必有一切灵魂的情欲；虽然并非每人都有全部。让我们洁净这些，并如此聆听天主的话，以痛悔的心听取今天给我们诵读的。\par

那是甚么呢？\BibleQuote{耶稣在加里肋亚的加纳行了这个神迹的开始。}我前些日子告诉过你们，有些人说这不是开始。\Quote{若加上\InnerQuote{加里肋亚的加纳}又怎样？}一人说：\Quote{这表明在加纳祂所行的正是\InnerQuote{开始}。}但在这些事上，我不敢断定任何确切的事。我以前已经表明，祂在受洗后才开始行神迹，在受洗前没有行过任何神迹；但至于受洗后所行的神迹中，这个是第一个还是别的为第一个，在我看来不必断然肯定。\par

\Quote{并显示出了祂的光荣。}\par

\Quote{怎样？}一人问道：\Quote{并且以何种方式？因为只有仆役、管宴席的、新郎注意到所发生的事，而不是在场的多数人。}那么祂怎样\Quote{显示出了祂的光荣}呢？至少祂自己方面是显示了；如果当时在场的人没有全都听到神迹，之后也会听到，因为直到如今它被传颂，并未被忽略。所有人在同一天没有知道这事，从以下经文可以清楚看出，因为说了祂\Quote{显示出了祂的光荣}之后，福音作者补充道：\par

\Quote{祂的门徒就信了祂。}\par

祂的门徒，在此之前已经对祂怀着惊奇。你们看到吗？行神迹尤其必要在心中有正直的人在场的时候，因为这些人会更乐意相信，并更仔细地留意情况。\Quote{若没有神迹，祂怎能被人认识呢？}因为祂的教义和先知的能力足以使听者的灵魂感到惊奇，以致他们以正直的倾向留意祂所行的事，因为他们的心已经对祂怀有好感。因此，在许多其他地方，福音作者说，由于那里居民的邪恶，祂没有行神迹。\BibleRef{玛 12:38}\par

% structure: p0010 source=h2 target=subsection reason=relative-heading-rank; base=h2

\subsection{若望福音 2:12}

那么祂同\Quote{祂的母亲来到葛法翁}是为何故？因为祂在那里没有行神迹，那城的居民并非正直地对待祂，而是彻底败坏的。基督表明这一点时说：\BibleQuote{而你，葛法翁，已被高举到天上，将被打入地狱。}\BibleRef{路 10:15}那么祂为何去呢？我认为，是因为祂打算稍后上耶路撒冷去，所以祂就去葛法翁，以免带着祂的母亲和弟兄到处走动。因此，祂离开并停留了一小段时间以光荣祂的母亲后，又开始了祂的神迹，将生养祂的母亲送回了家。所以福音作者说：\Quote{过了不多几日}，\par

% structure: p0012 source=h2 target=subsection reason=relative-heading-rank; base=h2

\subsection{若望福音 2:13}

祂在逾越节前几天领受了圣洗。但上耶路撒冷后，祂做了什么事呢？一件充满权柄的事：祂把那些兑换银钱和卖鸽子、牛、羊的人，以及为此目的在那里逗留的人，都赶出了圣殿。\par

2. 另一位福音作者写道：当祂赶出他们时，祂说：不要使我父的殿成为\BibleQuote{贼窝}；但这一位说：\par

% structure: p0015 source=h2 target=subsection reason=relative-heading-rank; base=h2

\subsection{若望福音 2:16}

它们在此并不彼此矛盾，而是显示祂做了第二次，并且这两句话并非在同一场合说的，而是祂在传教开始时这样做了一次，又在祂临近受难时再次这样做。因此，（在后者场合）用了更强硬的措辞，称之为（成了）\BibleQuote{贼窝}；但在此处，在祂奇迹的开端，祂并没有这样做，而是用了更为温和的斥责；由此看来，这件事可能是第二次发生。\par

\Quote{何故，}有人说，\Quote{基督为何这样做，并对这些人使用如此严厉的手段，这种事在别处从未见祂做过，即使祂被侮辱、被辱骂，被他们称为\InnerQuote{撒玛黎雅人}和\InnerQuote{附魔的}？因为祂甚至不满足于只言片语，而是拿起了鞭子，把他们赶了出去。} 是的，但正是在别人得益时，犹太人控告并恼怒反对祂；当祂的斥责有可能激怒他们时，他们对祂却没有表现出这种态度，因为他们既没有控告也没有辱骂祂。他们说了什么？\par

% structure: p0018 source=h2 target=subsection reason=relative-heading-rank; base=h2

\subsection{若望福音 2:18}

你看到他们极度的恶意了吗？以及施于他人的善行（比责备）更激怒他们？\par

当时祂说，圣殿被他们弄成了\BibleQuote{贼窝}，表明他们所卖的是通过偷窃、抢劫和贪财得来的，并且是借着别人的灾难致富；另一次，祂说是\BibleQuote{买卖之所}，指斥他们无耻的交易。\Quote{但是祂为何这样做？} 因为祂即将在安息日治病，并做许多此类他们视为违反法律的事，为了使祂不至于显得好像是来作为一位与祂父亲对立的敌对之神，祂抓住这个机会来纠正他们的任何此类猜疑。因为一位对圣殿表现出如此热心的人，不可能反对那圣殿的主宰、在圣殿中被敬拜的那位。毫无疑问，即使是祂在律法下生活的早年，也足以表明祂对立法者的敬畏，并且祂不是来颁布相反的法律的；但由于那些年岁可能因时间流逝而被遗忘，因为祂在一个贫穷卑微的住所长大，不为众人所知，因此祂后来在众人面前这样做（因为许多人因节期临近而在场），并且冒着很大风险。因为祂不仅\Quote{把他们赶出去}，还\Quote{推翻了桌子}，\Quote{倒出了钱币}，借此让他们明白，那位为了圣殿的秩序而置身危险的人，绝不会轻视祂的主子。如果祂是出于伪善而这样做，祂只需规劝他们；但让自己置身危险是非常大胆的。因为将自己暴露于如此多商贩的怒气中，以斥责和击打激起一群最野蛮的小贩反对自己，这不是一个伪装者的行为，而是一个为了圣殿的秩序选择承受一切的人的行为。\par

因此，不仅通过祂的行为，也通过祂的话语，祂显示祂与圣父的一致；因为祂没有说\Quote{圣殿}，而是说\BibleQuote{我父的圣殿}。看，祂甚至称祂为\Quote{父}，他们并没有生气；他们认为祂是在泛泛而言：但当祂继续讲下去，更清楚地说话，以便在他们面前提出祂与圣父平等的观念时，那时他们就生气了。\par

他们说什么？\Quote{祢做这些事，给我们显什么征兆呢？} 唉，他们何等疯狂！难道他们需要先看到征兆才能停止恶行，将天主的圣殿从这种不名誉中解放出来吗？祂对圣殿如此热心，这难道不是祂卓越的最大征兆吗？事实上，善良好义的人正因此而被区分出来，因为经上说：\Quote{他们，}祂的门徒，它说，\par

% structure: p0023 source=h2 target=subsection reason=relative-heading-rank; base=h2

\subsection{若望福音 2:17}

但犹太人没有忆起这预言，反而说：\Quote{祢给我们显什么征兆呢？} \BibleRef{咏 68:9} ，他们既因自己可耻的交易被中断而悲伤，又期望借此阻止祂，还想要向祂挑战行奇迹，并挑剔祂所做的事。因此祂不会给他们征兆；之前，当他们来问祂时，祂给了他们同样的回答：\BibleQuote{邪恶淫乱的世代寻求征兆，除了约纳先知的征兆，再没有征兆赐给它。} \BibleRef{玛 16:4} 只是那时回答清楚，现在却更为含糊。祂这样做是因为他们极其麻木；因为那位在他们未求时就先阻止他们、并给予他们征兆的，绝不会在他们求的时候拒绝他们，除非祂看到他们的心思邪恶虚假，他们的意图险恶。试想这问题本身从一开始是多么邪恶。当他们本该为祂的诚恳和热心而鼓掌，本该因祂如此关心圣殿而惊讶时，他们却责备祂，声称买卖是合法的，除非有人向他们显示征兆，否则任何人停止他们的买卖都是不合法的。基督说什么？\par

% structure: p0025 source=h2 target=subsection reason=relative-heading-rank; base=h2

\subsection{若望福音 2:19}

祂说了许多这样的话，在当时听者听来不能理解，但对后来的人却是可以理解的。祂为什么这样做呢？为了当祂预言的应验得以实现时，祂可以被看出是从起初就预知将要发生的事；这预言正是如此。因为，福音作者说，\par

% structure: p0027 source=h2 target=subsection reason=relative-heading-rank; base=h2

\subsection{若望福音 2:22}

但在说这话的时候，犹太人困惑于它可能是什么意思，并设法探究，说：\par

% structure: p0029 source=h2 target=subsection reason=relative-heading-rank; base=h2

\subsection{若望福音 2:20}

他们说：\BibleQuote{四十六年}，指的是后来的建筑，因为前一座用了二十年时间建成。 \BibleRef{厄上 6:15}\par

3. 那么，为何祂不解除难题，说：\Quote{我不是指那座圣殿，而是指我的身体}？为何圣史在日后撰写福音时解释这句话，而耶稣当时却保持沉默？祂为何如此沉默？因为人们不会接受祂的话；如果连门徒都无法理解这话，何况群众呢？\Quote{当，}圣史说，\Quote{祂从死者中复活后，他们才想起，并相信了圣经和祂的话。}当时有两件事阻碍了他们：一是复活的事实，二是更重大的问题，即祂是否是内住的天主；这两件事祂都隐约地说到了，当祂说：\Quote{你们拆毁这座圣殿，我三天内要把它重建起来。}圣保禄指出这是祂天主性的一个不小证明，他写道：\Quote{按圣德的神性，藉着从死者中的复活，被宣示为具有权能的天主之子。}（\BibleRef{罗 1:4}）\par

但为何祂在那里、在这里、在每一处都以此作为记号，有时说：\Quote{当你们举起人子时，那时你们就知道我是那一位}（\BibleRef{若 8:28}）；在另一处说：\Quote{除约纳先知的记号外，不会给你们其他记号}（\BibleRef{玛 12:39}）；又在此处说：\Quote{三天内我要把它建起来}？因为最能表明祂不是普通人的，是祂能设立战胜死亡的凯旋纪念碑，如此迅速地废除其长期暴政，结束那艰苦的战争。因此，祂说：\Quote{那时你们就知道。}\Quote{那时。}何时？当我在复活后吸引（所有）世界归向我时，那时你们就知道我以天主、真天主之子的身份做这些事，以报复对我父的侮辱。\par

那么，为何祂没有说：\Quote{需要用\InnerQuote{记号}来制止恶行吗？}而是承诺给他们一个记号？因为若那样说，祂会更激怒他们；但这样反而使他们惊讶。然而他们对此没有回应，因为祂对他们所说的话似乎难以置信，以致他们甚至没有停留下来追问祂，而是认为不可能而放过了。但如果他们明智，尽管当时看来难以置信，但当祂行了许多奇迹时，他们就会前来追问祂，请求祂为他们解释难题；但因为他们是愚昧的，他们对所说的话一部分完全不予理会，另一部分则以恶意听闻。因此基督以隐晦的方式向他们说话。\par

问题仍然存在：\Quote{门徒们怎么会不知道祂必须从死者中复活呢？}这是因为他们尚未蒙赐圣神的恩赐；因此，尽管他们不断听到祂关于复活的讲论，却不理解，而是心中思忖这究竟是什么意思。因为那声称人能自己复活，且要以这样的方式复活的断言，是非常离奇和矛盾的。因此，当伯多禄对复活一无所知，说：\Quote{主，愿这事远离祢}（\BibleRef{玛 16:22}）时，受到了斥责。基督在事件发生前没有向他们清楚启示，恐怕他们从一开始就跌倒，因事情极不可能而对祂的话产生不信，也因为他们尚未清楚认识祂是谁。因为没有人能不信那些由事实大声宣告的事，而有些人可能会不信口传的事。因此，祂起初允许祂话语的意思被隐藏；但当祂藉着他们的经验证实了祂的话后，祂随后赐予他们理解祂话语的能力，以及圣神的恩赐，使他们能立刻接受一切。\Quote{祂，}耶稣说，\Quote{要叫你们想起一切。}（\BibleRef{若 14:26}）因为那些在一夜之间抛弃了对祂的一切尊重，逃跑并否认认识祂的人，若非享有圣神的丰富恩宠，是很难记住祂所做所说的一切。\par

\Quote{但是，}有人说，\Quote{如果他们要从圣神那里听见，为何还需要跟随基督，既然他们不会记住祂的话？}因为圣神并非教导他们，而是提醒他们基督先前所说的话；并且这对基督的荣耀不无好处，即他们被引向记起祂曾对他们所说的话。起初，乃是天主的恩赐，使圣神的恩宠如此丰富地降临在他们身上；但此后，乃是他们自己的德行使他们持守了恩赐。因为他们显示出光耀的生命、极大的智慧、辛勤的劳作，轻视今生，视世事如无物，超脱万有；如同一只轻翅的鹰，以他们的功绩高飞，直抵天堂，并以此拥有圣神不可言喻的恩宠。\par

那么，让我们效法他们，不要熄灭我们的灯，而是藉着施舍使它们保持明亮，因为如此这火的光才得以保存。让我们趁还在此时，把油收集到我们的器皿里，因为我们一旦去了那另一个地方，就不能再买油，也不能从别处得到，除非经由穷人之手。因此，让我们非常充足地收集油，如果我们至少渴望与新郎一同进去的话。但如果我们不这样做，就必须留在洞房之外，因为不可能，即使我们行了成千上万的其他善事，没有施舍就无法进入天国的门。那么，让我们非常充足地行施舍，好使我们能享受那些不可言喻的福乐；愿我们都能藉着我们的主耶稣基督的恩宠和慈爱，与祂一起，将光荣归于圣父及圣神，直到永远。阿们。\par

\SourceRecord{Translated by Charles Marriott.；Nicene and Post-Nicene Fathers, First Series；Vol. 14.；Edited by Philip Schaff.；Buffalo, NY: Christian Literature Publishing Co.,；1889.；<http://www.newadvent.org/fathers/240123.htm>.}

% structure: p0001 source=h1 target=section reason=page-title

\section{论\BibleBook{若望}{福音}讲道集第二十四篇}

% structure: p0002 source=h2 target=subsection reason=relative-heading-rank; base=h2

\subsection{\BibleRef{若 2:23}}

1. 那时的人中，有些人固守他们的错误，另一些人则把握了真理；而在后者中，有些人保持了真理不久又再次偏离。基督在比喻中将他们比作没有深种、根在土表的种子，并说他们很快就会枯萎。福音作者在这里向我们指出了这些人，说：\par

\BibleQuote{当祂在耶路撒冷过逾越节过节时，有许多人看见祂所行的奇迹，就信了祂。}\BibleRef{若 2:23}\par

% structure: p0005 source=h2 target=subsection reason=relative-heading-rank; base=h2

\subsection{\BibleRef{若 2:24}}

因为在他门徒中，更成全的是那些不单因奇迹而来，也因他的教导而来的人。比较粗俗的人被奇迹吸引，但更善于思考的人被他的预言和道理吸引；因此被他的教导所吸引的人，比被奇迹所吸引的人更坚定。基督也称他们为\Quote{有福的}，说：\BibleQuote{那些没有看见而相信的，才是有福的。}\BibleRef{若 20:29} 但这里提到的那些人并不是真正的门徒，下文显示了出来，因为它说：\BibleQuote{耶稣却不信任他们。}\BibleRef{若 2:24} 为什么呢？\BibleQuote{因为他知道一切。}\BibleRef{若 2:25}\par

% structure: p0007 source=h2 target=subsection reason=relative-heading-rank; base=h2

\subsection{\BibleRef{若 2:25}}

意思是这样的。\Quote{祂居住在人的心中，进入人的思想，不注意外在的言语；深知他们的热心只是暂时的，不信任他们如同完全的弟子，也不把一切道理托付给他们，好像他们已经成了坚定的信徒。}然而，知道人心深处唯独属于天主，\BibleQuote{祂一一铸造了人心}\BibleRef{咏 32:15}；因为撒罗满说：\BibleQuote{唯独祢知道人心}\BibleRef{列上 8:39}；因此祂不需要见证来了解祂受造物的思想，也不信任他们只因他们暂时的、肤浅的信德。人既不知道现在也不知道未来，常常毫无保留地告诉和托付给那些以诡诈接近他们、不久就会背离的人；但基督却非如此，因为祂深知他们一切的秘密思想。\par

如今也有许多这样的人，他们确实拥有信德的名号，却是不稳定的，容易被带走；因此现在基督也不将自己托付给他们，反而向他们隐藏许多事；正如我们不信任泛泛之交而只信任真正的朋友，基督也是如此。请听祂对门徒说的话：\BibleQuote{从此以后，我不再称你们为仆人，你们是我的朋友。}\EditorialNote{参阅若望福音第十五章14至15节} 这是为何？\BibleQuote{因为我父所听的一切，我都已告诉你们了。}因此，祂没有给那些求征兆的犹太人任何征兆，因为他们怀着试探的心求祂。的确，求征兆是试探者的做法，那时如此，现在亦然；因为甚至现在也有一些寻求征兆的人说：\Quote{为什么现今不再有奇迹发生呢？}如果你有信德，正如你应当有的，并且爱基督如你应当爱的，你就不需要征兆，它们是为不信的人所赐的。\Quote{那么，}有人问，\Quote{它们不是赐给犹太人了吗？}它们确实赐下了；如果有些时候他们求却未得到，那是因为他们求不是为了除去不信，而是为了更加证实他们的邪恶。\par

% structure: p0010 source=h2 target=subsection reason=relative-heading-rank; base=h2

\subsection{\BibleRef{若 3:1-2}}

此人也出现在福音的中部，为基督辩护；因为他曾说：\BibleQuote{我们的法律，不经审问是不定人的罪的。}\BibleRef{若 7:51} 犹太人愤怒地回答他说：\Quote{你且研读，从加里肋亚不会出先知。}在基督受难后，他又悉心料理主的身体埋葬之事：福音作者说：\BibleQuote{尼苛德摩也来了，他以前曾在夜里来见过主，带来了一百斤的没药和沉香。}\BibleRef{若 19:39} 即使那时，他对基督存有善意，但并非理应如此，也没有怀着适当的情感，因为他仍然被犹太人的软弱所纠缠。因此他夜里来，因为他害怕白天来。然而仁慈的天主没有因此拒绝或斥责他，也没有剥夺对他的教导，反而以极大的仁慈与他交谈，向他启示了非常高超的道理，虽然是以隐晦的方式，但毕竟启示了。因为他远较那些因邪恶而行的人更值得宽恕。那些人完全无可推诿；而他虽然应受谴责，却不至于同等程度。\Quote{那么，福音作者为什么关于他没有说这样的话呢？}他在另一处说过：\BibleQuote{官长中也有许多人信了祂，但因法利塞人的缘故，不敢明认，免得被逐出会堂。}\BibleRef{若 12:42} 但在这里，他通过提到他夜里来，暗示了全部。那么尼苛德摩说了什么呢？\par

\BibleQuote{辣彼，我们知道祢是从天主而来的师傅，因为祢所行的这些奇迹，除非天主与祢同在，是谁也不能行的。}\BibleRef{若 3:2}\par

2. 尼苛德摩仍停留在低处，对祂怀着人的想法，将祂当作先知谈论，从祂的奇迹中想不到任何伟大的事。他说：\Quote{我们知道祢是从天主而来的师傅。}\Quote{那么，你为什么夜里暗暗地来到那位讲论天主之事、从天主而来的祂面前呢？为什么你不公开与祂交谈呢？}但耶稣没有对他说这样的话，也没有斥责他；因为先知说：\BibleQuote{压伤的芦苇，祂不折断；将熄的灯芯，祂不吹灭；祂不争竞，也不喧嚷。}\BibleRef{依 42:2} 并且祂自己又说：\BibleQuote{我来不是为审判世界，乃是为拯救世界。}\BibleRef{若 12:47}\par

\BibleQuote{除非天主与祂同在，是谁也不能行这些奇迹的。}\BibleRef{若 3:2}\par

尼苛德摩在此仍像异端者一样说话，声称祂内在有一种运作的能力，且需要他人的帮助才能行祂所行的事。那么基督说什么呢？请注意祂无比的屈尊就卑。祂暂时忍住不说：\Quote{我不需要别人的帮助，而是以大能行一切，因为我是天主的真子，拥有与我的父相同的权能}，因为这对祂的听者来说太难以承受。因为我现在所说的，正是我一贯所说的：基督所渴望的，与其说是暂时彰显自己的尊威，不如说是说服众人，祂所做的没有一件与祂的父相违背。因此，在许多地方，祂在言语上似乎受到限制，但在行动上却并非如此。因为当祂行奇迹时，祂以权能做一切，说：\BibleQuote{我愿意，你洁净了吧！}（见：\BibleRef{玛 8:3}）\BibleQuote{塔里塔，起来！}（见：\BibleRef{谷 5:41}）\BibleQuote{伸出你的手来！}（见：\BibleRef{谷 3:5}）\BibleQuote{你的罪赦了！}（见：\BibleRef{玛 9:2}）\BibleQuote{不要作声，平靖了吧！}（见：\BibleRef{谷 4:39}）\BibleQuote{拿起你的床，回家去吧！}（见：\BibleRef{玛 9:6}）\BibleQuote{你这污秽的魔鬼，我命令你：从他身上出来！}（见：\BibleRef{谷 9:25}）\BibleQuote{照你的愿望成就吧！}（见：\BibleRef{玛 15:28}）\BibleQuote{若有人对你们说什么，你们就说：主要用它。}（见：\BibleRef{谷 11:3}）\BibleQuote{今天你就要与我一同在乐园里。}（见：\BibleRef{路 23:43}）\BibleQuote{你们一向听过对古人说：不可杀人；我却对你们说：凡向弟兄无故发怒的，就要受裁判。}（见：\BibleRef{玛 5:21}）\BibleQuote{来跟随我，我要使你们成为渔人的渔夫。}（见：\BibleRef{谷 1:17}）到处我们都能看出祂的权威极大；因为在祂的行动中，无人能指责祂所做的事。怎么可能呢？倘若祂的话没有应验，没有按祂的命令成就，任何人都可以说那是疯子的命令；但既然它们确实应验了，成就的事实就堵住了人的口，甚至违背他们的意愿。但至于祂的言论，他们常常在傲慢中指责祂癫狂。因此，现在在尼苛德摩的情况下，祂没有公开说什么，而是用隐晦的言辞将他从卑微的思想中提升，教导他：祂在自己内有足够的能力显奇迹，因为祂的父生了祂，使祂完美、全备，毫无瑕疵。\par

但让我们看看祂如何实现这一点。尼苛德摩说：\Quote{辣彼，我们知道祢是由天主而来的师傅，因为除非天主与祂同在，没有人能行祢所行的这些奇迹。}他以为自己对基督说了如此伟大的话时，自己说了什么了不起的事。那么基督说什么呢？为了表明他甚至还没有踏上正确知识的门槛，也没有站在门廊上，而是仍在宫殿外的某个地方游荡，无论是他，还是其他说类似话的人，并且他对独生子怀有这种看法时，甚至没有瞥见真正的知识，祂说什么呢？\par

% structure: p0017 source=h2 target=subsection reason=relative-heading-rank; base=h2

\subsection{若望福音 3:3}

也就是说，\Quote{除非你重生并接受正确的道理，你是在外面游荡，远离天国。}但祂说得并不如此直白。为了使这句话不那么难以承受，祂并没有直接针对他说，而是不指定对象地说：\Quote{人除非重生}；几乎是说：\Quote{你和其他任何对我有如此看法的人，都在天国之外。}如果祂不是出于建立此事的渴望而说话，祂的回答本来会与所说的话相符。但现在，如果这些话是对犹太人说的，他们会嘲笑祂并离去；但尼苛德摩在这里也表现出他对教导的渴望。这就是为什么在许多地方基督说得隐晦，因为祂希望激起听者提问，使他们更加专注。因为说得直白的话常常被听者忽略，而隐晦的话则使他更活跃、更热忱。现在祂所说的话，大致是这样的：\Quote{你若不重生，若不藉着重生之洗的圣神分施，就不能对我有正确的看法，因为你所有的看法不是属神的，而是属血肉的。}（参阅：\BibleRef{铎 3:5}）但祂并没有这样说，因为祂不愿使一个已经尽其所能带来他所拥有的、尽其所能说话的人感到困惑；祂在不知不觉中引导他提升到更大的知识，说：\Quote{人除非重生。}\Quote{重生}这个词，有些人理解为\Quote{从天而来}，另一些人理解为\Quote{从起初}。\Quote{不可能，}基督说，\Quote{不如此重生的人，不能看见天主的国}；在此指向自己，并宣告除了自然的视觉之外还有另一种视觉，我们需要另一双眼睛来看基督。听到这些话后……\par

% structure: p0019 source=h2 target=subsection reason=relative-heading-rank; base=h2

\subsection{若望福音 3:4}

你称祂为\Quote{老师}，你说祂是\Quote{从天主而来}，却又不接受祂的话，反而对你的老师使用一种表达极大困惑的说话方式？因为\Quote{如何}，是那些没有坚固信仰、仍属世俗之人的怀疑问题。因此撒辣在说\Quote{如何？}时笑了。许多其他问过这个问题的人，都从信仰中堕落了。\par

3. 异端者之所以仍持守其异端，是因为他们常这样探究，有人说：\Quote{祂是如何受生的？}另有人说：\Quote{祂是如何成了血肉的？}并将那\OriginalTerm{无限本体}{Infinite Essence}置于自身推理的软弱之下。我们既知此事，就当避免这种不合时宜的好奇，因为凡探究这些事的人，尚未学会\Quote{如何}，便已偏离了正信。为此，\Person{尼苛德摩}怀疑不解，询问此事如何可能（因为他明白所说的话是指自己而言），他困惑、头昏、不知所措，他是来到一个人面前，却听见超越人言的话语，是从来没有人听过的；一时之间，他因话语的崇高而振奋，但仍处于黑暗之中，摇摆不定，四处飘荡，不断从信仰中坠落。因此他坚持证明不可能，以激发祂更清楚的教导。\par

\Quote{人岂能再进入母腹而重生么？}他说。\par

你看见了吗？当人将属灵之事交于自己的推理时，说话是何等可笑，似乎是在玩笑，或如醉酒一般，当他窥探那超越天主美意所说的话，且不肯接受信德的服从时。\Person{尼苛德摩}听到了\OriginalTerm{属神重生}{spiritual Birth}，却不领会它是属神的，反而将这话拖到血肉的低贱中，使如此伟大崇高的道理依附于肉体的后果。于是他编造出无聊和无谓的难题。因此\Person{保禄}说：\Quote{属血气的人，不接受天主圣神的事。}\BibleRef{格前 2:14}但即便如此，他仍保持对\Person{基督}的尊敬，因为他没有嘲笑所说的话，而是认为不可能，便默然不语。当时有两个难题：这种重生和天国；因为无论是在犹太人中间，还是这种重生，都从未听说过天国的名号。但他停留在第一个难题上，这最使他心神不安。\par

我们既知此事，就不要凭推理探究有关天主的事，也不要以地上的后果来规范天上的事，更不要使它们屈从于本性的必然；却要怀着敬畏之心思想一切，相信圣经所说的话；因为忙碌好奇的人一无所获，而且除了找不到所寻求的，还要受极重的刑罚。你们听见了（圣父）生了（圣子）：要相信你们所听见的；但不要问\Quote{如何}，以致取消了受生；这样做是极其愚昧的。因为如果这个人，由于听见一种受生（不是那不可言喻的受生，而是因恩宠而来的受生）而没有领会其中的伟大，反而生出人性而属地的思想，因而陷于黑暗和怀疑，那么那些忙碌好奇地探究那\OriginalTerm{不可言喻的受生}{that ineffable Generation}和那\OriginalTerm{至可怕的受生}{that most awful Generation}（超越一切理性和理智）的人，该受怎样的刑罚呢？因为没有什么比人的推理更令人头昏，它的一切话语都是属地的，不能忍受从上而来的光照。属地推理充满了泥泞，因此我们需要天上的活水，好使泥泞沉淀，清澈部分上升，与天上的教训相融合；当我们献上诚实的灵魂和正直的生活时，这事就会发生。因为不仅不合时宜的好奇心，而且腐败的品行也能使理智昏暗；因此\Person{保禄}对格林多人说：\Quote{我曾给你们喝奶，不是吃饭，因为那时你们还不能吃，就是如今还是不能，因为你们仍是属血肉的；因为在你们中间有嫉妒、纷争和分裂，你们岂不是属血肉的么？}\BibleRef{格前 3:2}又在致希伯来人书及许多地方，可见\Person{保禄}断言这是邪说的原因；因为为情欲所困的灵魂不能看见任何伟大或尊贵的事，反而好像被一层薄膜遮蔽，遭受最痛苦的视力模糊。\par

那么，让我们洁净自己，点燃知识的明灯，不要将种子撒在荆棘中。荆棘是什么，你们知道，无需我们告诉你们；因为你们常听见\Person{基督}用这个名号称呼今生的忧虑和财富的欺骗。\BibleRef{玛 13:22}这是有理由的。因为荆棘不结果实，这些也是如此；荆棘刺伤那些触摸它们的人，这些情欲也是如此；荆棘容易被火点燃，被农夫憎恨，世俗之物也是如此；在荆棘中隐藏着野兽、蛇和蝎子，在财富的欺骗中也是如此。但让我们点燃圣神之火，好烧尽荆棘，赶走野兽，为农夫清理田地；清洁之后，让我们用圣神的活水浇灌，栽种那结果实的橄榄树，那最温和的树，常青的、光明的、滋养的、健康的。这一切品质，施舍都具备，它如同一个印记盖在拥有它的人身上。这植物，即使死亡来临，也不能使它枯萎，反而永远熠熠生辉，照亮心智，滋养灵魂的力量，使其力量更坚强。如果我们恒常拥有它，我们就能坦然无惧地见到新郎，进入婚筵；愿我们众人能达到那境界，赖我们的主耶稣基督的恩宠和慈爱，愿光荣归于祂及圣父和圣神，至于无穷之世。阿们。\par

\SourceRecord{Translated by Charles Marriott.；Nicene and Post-Nicene Fathers, First Series；Vol. 14.；Edited by Philip Schaff.；Buffalo, NY: Christian Literature Publishing Co.,；1889.；<http://www.newadvent.org/fathers/240124.htm>.}

% structure: p0001 source=h1 target=section reason=page-title

\section{若望福音讲道集 第二十五篇}

% structure: p0002 source=h2 target=subsection reason=relative-heading-rank; base=h2

\subsection{\BibleRef{若 3:5}}

每日到老师那里去的小孩子们，接受功课，背诵功课，从不停止这种获取，有时甚至昼夜不停，而他们是为会朽坏和短暂的东西被迫如此行。现在，我们并不要求你们这些成年人像你们要求你们的孩子那样辛劳；因为我们并非每天，而只是一周两天劝你们倾听我们的话语，并且只占用一天中的一小部分，好使你们的功课变得容易。为同一理由，我们也把圣经所写的分成小份给你们，使你们能容易地接受并储存在你们心灵的仓库里，并且如此努力记住这一切，使你们自己能亲自准确重述给别人，除非有人昏昏欲睡、迟钝，比一个小孩子更懒惰。\par

现在我们留心前面所说的后续。当尼苛德摩陷于错误，将基督的话曲解为肉身的出生，说一个老年人不可能重生时，请注意基督如何更清楚地启示这重生的方式，即使对属肉体的质问者仍有困难，却仍能将听者从其对此的低浅看法中提升起来。祂说了什么？\BibleQuote{我实实在在告诉你们：人除非由水和圣神而生，不能进天主的国。}祂所宣告的是：\Quote{你们说这是不可能的，我却说这是如此绝对可能，以致成为必要，并且若不如此，甚至不能得救。}因为天主使必要的事也成为极其容易的。那按照肉身而生的出生，是出于尘土，因此天堂对它是关闭的，因为尘土与天堂有什么相干呢？但那由圣神而生的，却轻易为我们敞开天上的穹苍。听着，你们这些尚未受光照的人，战栗、哀叹吧！这威胁是可怕的，这判决是可怕的。\BibleQuote{没有由水和圣神而生的人，不可能进天国}，因为他穿着死亡、诅咒、丧亡的衣服，尚未接受他主的标记，他是陌生人和外人，没有王室的暗号。\BibleQuote{除非一个人由水和圣神而生，不能进天国。}\par

即使如此，尼苛德摩仍不明白。再没有比将属灵之事交付推论更糟的了；正是这一点使他不能设想任何崇高伟大之事。这就是为什么我们被称为信友，好使我们抛下人类理性之软弱于下，升上信德的高峰，并将我们的大部分福分托付于她的教导；如果尼苛德摩这样做了，这件事就不会被他视为不可能。那么基督做了什么？为了将他从低下的想象中领出来，并表明祂所说的不是属于肉身的出生，祂说：\BibleQuote{除非一个人由水和圣神而生，不能进天国。}祂说这话，是愿意借着这威胁的恐惧吸引他归向信德，并说服他不要认为这事不可能，且努力使他离开对肉身出生的想象。\Quote{我指的是另一次出生，啊，尼苛德摩。你为什么将这话拖到地上？你为什么将这事置于自然的必然性之下？这次出生对这些产痛而言太高了；它与你们毫无共同之处；它确实被称为\InnerQuote{出生}，但只是在名称上有共同之处，实际上却不同。脱离那寻常熟悉的事物吧；我将一种不同的生育带入世界；我愿人以另一种方式被产生：我来是要带来一种新的受造方式。我用地和水造了人；但那被造的没有用处，器皿歪斜了；我不再用地和水造他们，而是用\InnerQuote{由水}和\InnerQuote{由圣神}。}\par

如果有人问：\Quote{\Emph{由水}是什么意思？}我也要问：由地是什么意思？那泥土是如何被分成不同部分的？原料是单一的（只是土），而由它制成的物品却是多种多样、各式各样的？骨骼、筋腱、动脉、静脉从何而来？薄膜、器官的管道、软骨、组织、肝、脾、心脏从何而来？皮肤、血、黏液、胆汁从何而来？这般巨大的能力从何而来？如此多样的颜色从何而来？这些不属于地或泥土。当土壤接收种子时，它如何使种子发芽，而肉体接收种子却消耗它们？土壤如何滋养投入其中的东西，而肉体却由这些东西滋养，却不滋养它们？例如，土地接收水，使之变成酒；肉体却经常接收酒，使之变成水。那么，既然根据前面所说，地的本性是与身体的本性相反的，我们又怎能清楚这些事物是由地形成的呢？我无法通过推理发现，我只凭信德接受。既然每日发生且我们亲手处理的事物需要信德，那么那些比这些更神秘、更属灵的事物岂不更需要信德吗？因为正如大地，虽无灵魂且静止不动，却因天主的意愿而获得能力，在其中行出如此奇迹；那么当圣神与那水同在时，所有那些如此奇异且超越理性的事，岂不更容易发生吗？\par

那么，不要因为你们看不见这些事就不信；你们看不见自己的灵魂，却仍相信你们有灵魂，且它是在肉体之外的另一个东西。\par

但基督没有用这个例子引领他，而是用了另一个例子；祂没有提出灵魂的例子（虽然灵魂是无形的），因为听者的心性还太迟钝。祂摆在他面前的是另一个例子，与有形物体的密度无关，却又未达到无形本性的高度，那就是风的活动。祂从水开始，水比地轻，但比空气密。正如起初地是质料，但一切出于塑造它者之手；同样，现在水是质料，而一切出于圣神的恩宠：那时\BibleQuote{人成了有生命的灵魂}\BibleRef{创 2:7}，现在他成了\BibleQuote{赋予生命的圣神}。但二者之间的差别极大。灵魂不能赋予生命给别人，只能赋予拥有它的人；圣神不仅活着，也赋予别人生命。例如，宗徒们甚至复活了死人。那时，人在创造完成后才被造；现在恰好相反，新人在新创造之前被造；他首先出生，然后世界被重新塑造。\BibleRef{格前 15:45} 正如起初祂完整地造了他，现在祂也完整地再造他。那时祂说\BibleQuote{让我们给他造一个助手}\BibleRef{创 2:18}，但在这里祂没有说这类话。那领受了圣神恩赐的人还需要什么别的助手呢？那属于基督身体的人还需要什么帮助呢？那时祂按天主的肖像造了人，现在祂使人自己与天主合一；那时祂命人统治鱼和野兽，现在祂将我们的初果高举于诸天之上；那时祂赐给人一个园子作为居所，现在祂为我们打开了天堂；那时人在第六天被造，世界几乎完成；但现在是在第一天，在起初，在光被造的时候。从这一切可以清楚看出，所成就的事属于另一个、更美好的生命，属于没有终结的状态。\newline{}\par

第一次创造，即亚当的创造，出于泥土；第二次，女人的创造，出于他的肋骨；第三次，亚伯的创造，出于种子；然而我们无法理解其中任何一个，也无法用推理证明这些情况，尽管它们属于最属地的性质；那么我们怎能说明那借着圣洗的不可见重生呢？这远比这些更为崇高，或为那奇异而奇妙的诞生要求论证呢？因为当那诞生发生时，连天使都站在一旁，却不能述说那奇妙运作的方式，他们只是站着，并非执行什么，而是观看所发生的事。圣父、圣子和圣神成就了一切。让我们相信天主的宣告；这比实际看见更可信。视觉常常出错，而天主的圣言不可能落空；让我们相信它。那使无变为有的话语，当它谈论事物的本性时，是可信的。那么它说什么呢？说所成就的是重生。如果有人问\Quote{如何}，用天主的宣告堵住他的口，那是最有力、最明确的证明。如果有人追问\Quote{为什么包括水}，让我们也反过来问：\Quote{起初在造人时为何用土？}因为天主不用土也能造人，这对每个人都是显然的。不要过分好奇。\newline{}\par

水是绝对必要、不可或缺的，你可以这样明白。有一次，圣神在用水之前已经降下，宗徒并未停留于此，而是仿佛水是必要而非多余的，看他说什么：\BibleQuote{这些人既已领受了圣神，如同我们一样，谁能禁止用水给他们施洗呢？}\BibleRef{宗 10:47}\newline{}\par

在圣洗中，我们与天主盟约的承诺得以实现——埋葬与死亡，复活与生命；这些同时发生。因为我们把头浸入水中时，旧人被埋葬在下面的坟墓里，永远沉没；然后我们抬起头时，新人代替他升起。我们容易浸入再抬起头，同样天主容易埋葬旧人、显出新人。这样做三次，好叫你明白圣父、圣子和圣神的能力成就了这一切。为了表明我们所说的不是推测，请听保禄说：\BibleQuote{我们借着圣洗与祂同葬于死亡}：又说：\BibleQuote{我们的旧人已与祂同钉十字架}：又说：\BibleQuote{我们已与祂同植于祂死亡的相似中。}\BibleRef{罗 6:4} 不但圣洗被称为\Quote{十字架}，十字架也被称为\Quote{圣洗}。基督说：\BibleQuote{我所受的圣洗，你们也将要受}\BibleRef{谷 10:39}；以及\BibleQuote{我有当受的圣洗}\BibleRef{路 12:50}（你们不知道）；因为我们容易浸入再抬起头，同样祂照自己的意愿容易死去并复活，或者更轻松，只是为了一奥秘的安排而延迟了三天。\newline{}\par

3. 那么，我们这些蒙受如此奥秘的人，要活出与恩赐相称的生命，即最卓越的品行；而你们尚未蒙恩宠的人，要尽力争取，好使我们成为一体，成为兄弟。因为只要我们在这方面分裂，纵然有人是父亲、儿子、兄弟或任何其他关系，他也不是真正的亲属，因为他断绝了从上而来的关系。被凡人家庭的纽带束缚，如果没有属灵的联系，又有何益处？地上的血亲关系有什么用，倘使我们在天上成为陌生人？因为慕道者对于信友是陌生人。他没有同一个头，没有同一个父亲，没有同一座城，也没有同样的食物、衣服、餐桌和房屋，一切都不同：前者的一切都在地上，后者的一切都在天上。一人以基督为君王，另一人以罪与魔鬼为君王；一人的食物是基督，另一人的食物是那会朽坏腐败的食物；一人以虫蛀之物为衣，另一人以天使之主为衣；一人的城是天堂，另一人的城是地上。既然我们没有共同之物，请告诉我，我们将在什么上相通？莫非我们出于同一产痛，出自同一母胎？这与那最完善的关系毫无关系。那么，让我们努力成为天上之城的居民。我们还要在边境徘徊多久，而本该收回我们古老的祖国？我们冒的不是寻常的危险；因为倘若发生（愿天主不许！）死亡突然来临，我们未受圣洗而离开人世，即使拥有万般德行，我们的命运也不外乎地狱、毒虫、不灭之火和不可解之锁链。但愿天主使所有听到这些话的人都不受那种刑罚！而这将实现，如果我们蒙受神圣奥秘的恩宠，并在那根基上建造金、银、宝石；因为这样我们离世后，就能在那地方显现为富有，因为我们不把财富留在这里，而是藉穷人之手将它们运到不可侵犯的宝库，当我们借给基督时。我们在那里有许多债务，不是金钱的，而是罪的；那么让我们把财富借给祂，好使我们获得罪的赦免；因为祂就是审判者。我们不可在这里祂饥饿时忽略祂，好使祂在那里永远喂养我们。在这里我们给祂穿衣，好使祂不使我们失去祂所赐的安全。如果我们在这里给祂喝，我们就不会像那富人一样说：\Quote{打发拉匝禄来，用指尖蘸点水，凉凉我的舌头。} 如果我们在这里接待祂到家中，祂就在那里为我们预备许多住所；如果我们到监狱去探望祂，祂也会解除我们的锁链；如果我们收留作客的祂，祂就不使我们成为天国的陌生人，反而将我们在天上之城中赐予一份；如果我们探望患病的祂，祂也会迅速解救我们的病弱。\par

那么，我们受领大恩，虽然只付出少许，仍要付出那少许，好获得那更大的。趁着还有时间，让我们播种，好能收割。当冬天临到我们，当海不再可航行，我们就不再是这交易的主人了。但冬天何时来临？当那伟大而显赫的日子临近时。那时我们将停止在这广阔的大海上航行，因为现世的生命就像这样。现在是播种的时候，那时是收获和得利的时候。如果一个人在播种时不撒种，而在收获时播种，他不仅一无所获，而且会显得可笑。但如果现在是播种时节，那么这就不是聚集的时候，而是分散的时候；让我们分散，好能聚集，不要现在聚集，免得失去我们的收成；因为，如我所说，这个季节邀请我们播种、花费、付出，而不是收集和积攒。那么，我们不要放弃这个机会，而要撒下充足的种子，不要吝惜我们的积蓄，好使我们能连本带利地收回，藉着我们主耶稣基督的恩宠和慈爱，愿光荣归于祂及圣父和圣神，至于无穷之世。阿们。\par

\SourceRecord{Translated by Charles Marriott.；Nicene and Post-Nicene Fathers, First Series；Vol. 14.；Edited by Philip Schaff.；Buffalo, NY: Christian Literature Publishing Co.,；1889.；<http://www.newadvent.org/fathers/240125.htm>.}

% structure: p0001 source=h1 target=section reason=page-title

\section{\WorkTitle{若望福音}讲道集第二十六篇}

% structure: p0002 source=h2 target=subsection reason=relative-heading-rank; base=h2

\subsection{若望福音 3:6}

天主的独生子使我们配享的奥秘是何等伟大！何等伟大，且是我们不配的，却是祂理当赐予的。若按我们的功德计算，我们不单不配领受这恩赐，还该受罚和惩罚；但祂因不看这些，不单免了我们的罚，还白白赐给我们比最初更光耀的生命，领我们进入另一个世界，使我们成为另一个受造物。\BibleQuote{若有人在基督里}，\Person{保禄}说，\BibleQuote{他就是一个新受造物。}\BibleRef{格后 5:17} 这\BibleQuote{新受造物}是什么样的？听基督亲自宣告：\BibleQuote{人若不从水和圣神而生，不能进天主的国。}乐园曾交托给我们，而我们显出连居住在那里也不配，祂却高举我们到天上。在起初的事上我们被证明不忠信，祂却把更大的事托付我们；我们连一颗树也不能节制，祂却为我们预备了天上的福乐；我们没有守住乐园的居所，祂却为我们开了天上门。\Person{保禄}说得对：\BibleQuote{啊，天主的智慧与知识的富藏是多么深奥！}\BibleRef{罗 11:33} 如今不再有母亲、产痛、睡眠、交合和身体的拥抱；从此我们本性的结构全由圣神和水在上界形成。水被使用，成为那重生者的出生；水对信徒如同子宫对胚胎；因为信徒是在水中被塑造和形成的。起初有话说：\BibleQuote{水要滋生有生命的爬虫}\BibleRef{创 1:20}；但从主进入约旦河的水流时起，水不再生出\BibleQuote{有生命的爬虫}，而是生出有理性和携带圣神的灵魂；关于太阳所说的话——他\BibleQuote{好像新郎出洞房}\BibleRef{咏 17:6}——我们现在倒可以这样说信徒，因为他们发出比太阳更光亮得多的光芒。在子宫中成形需要时间，但在水中却不然，一切都在一瞬间完成。此处的生命是可朽的，且起源于其他身体的腐朽；那将要出生的是缓慢的（因为这是身体的本性，需要时间才能达至完善），但属灵之事却不然。为什么呢？因为所造之物从一开始就是被造成完美的。\par

当尼苛德摩听见这些事仍感困惑时，请看基督如何向他部分揭开这奥秘的秘密，并使那暂时对他隐晦的事变得清楚。祂说：\BibleQuote{那由肉生的属于肉，那由圣神生的属于神。}祂引导他离开一切感官之物，不容他用血肉之眼徒然探究所启示的奥秘；\Quote{\BibleQuote{我们所说的不是肉，而是神，}尼苛德摩啊}（这句话暂时把他引向天上），\Quote{不要寻求任何属于感官的事；圣神绝不能为那些眼睛所见，也不要以为圣神产生肉身。}\Quote{那么怎样呢？}也许有人会问，\Quote{主的肉体是怎样产生的呢？}不是只由圣神，也是由肉，正如\Person{保禄}所宣告的：\BibleQuote{生于女人，生于法律之下}\BibleRef{迦 4:4}；因为圣神固然不是从无中塑造了祂（否则何需子宫？），而是从童贞女的肉中。我不能向你们解释究竟如何；但这确实发生了，为的是没有谁认为那诞生的与我们的本性相异。因为即使在这事发生后，仍有一些人不相信这样的诞生，倘若祂没有分享童贞女的肉，他们会陷于何等不敬呢！\par

\BibleQuote{由圣神生的属于神。}你看见圣神的尊严吗？祂似乎是执行天主的工作；因为上文说有些人\BibleQuote{是由天主生的}\BibleRef{若 1:13}，这里祂说圣神生了他们。\par

\BibleQuote{由圣神生的属于神。}祂的意思是这样：\BibleQuote{那由圣神生的，就是属神的。}因为祂在这里所说的出生不是按本质，而是按尊荣和恩宠。如果圣子也是这样生的，祂在哪些方面优于这样生的人？祂怎么是独生子呢？因为我也由天主而生，虽然不是出于祂的本质；如果祂也不是出于祂的本质，那么在这方面祂与我们有何区别？不，那样祂就会被发现低于圣神；因为这种出生是借着圣神的恩宠。难道祂需要圣神的帮助才能继续作为子吗？这与犹太教义有何不同？\par

基督于是说了\BibleQuote{由圣神生的属于神}之后，又见他困惑，就把话题引向一个来自感官的例子，说：\par

% structure: p0008 source=h2 target=subsection reason=relative-heading-rank; base=h2

\subsection{若望福音 3:7-8}

因为祂说\BibleQuote{不要惊奇}，表明他灵魂的困惑，并把他引向比身体更轻微的东西。祂已经引导他离开属肉之事，藉着说\BibleQuote{由圣神生的属于神}；但当尼苛德摩不明白\BibleQuote{由圣神生的属于神}是什么意思时，祂又把他带到另一个比喻，不是把他带到身体的密度，也不是纯粹谈论无形之事（因为若听见那些，他不能接受），而是找到一种介于有体和无体之间的东西，即风的运动，然后把他带到那近似的。祂论风说：\par

\BibleQuote{你听见风的响声，却不能辨别它从哪里来，往哪里去。}\par

虽然祂说：\Quote{风随意向哪里吹}，祂这样说并非因为风有选择的能力，而是说明其自然运动无法被阻止，并且带有力量。因为圣经惯于这样谈论无生命之物，如它说：\Quote{受造之物被屈服于虚空之下，并不是出于自愿。}\BibleRef{罗 8:20} 因此，\Quote{风随意向哪里吹}这句话，是表明它不能被约束，它散布各处，无人能阻止它来往，而是带着大能到处运行，无人能扭转它的暴力。\par

2. \Quote{你听见它的声音}（即它的沙沙声、响声），\Quote{却不晓得从哪里来，往哪里去；凡从圣神生的，也是如此。}\par

这里是整个事情的结论。祂说：\Quote{如果你不能解释你所听见和触摸到的风的运动与路径，为何要过分担忧天主圣神的运作呢？你连风的声音都听见了，却不明白它。\InnerQuote{风随意向哪里吹}这句话也用来建立护慰者的能力；因为如果无人能抓住风，它随意移动，那么自然的律法、肉体生育的界限或任何类似的东西，就更不能约束圣神的运作。}\par

\Quote{你听见它的声音}这句是指风而言，从这一点可以清楚看出：祂不会在和一个不信者、不认识圣神运作的人交谈时说\Quote{你听见它的声音}。既然风虽然发出声音却不可见，那么属灵之生的诞生也非肉眼可见；然而风是一个身体，尽管非常精细；因为任何感官对象都是身体。既然你不因看不见这个身体而抱怨，也不因此不信，那么当你听到\Quote{圣神}时，为何犹豫并要求如此精确的解释，而你面对一个身体时却不这样做？那么尼苛德摩如何呢？他仍然持守他低下的犹太见解，即使如此清晰的例子已经摆在他面前。因此当他再次怀疑地说：\par

% structure: p0015 source=h2 target=subsection reason=relative-heading-rank; base=h2

\subsection{\BibleRef{若 3:9-10}}

请注意，祂从未指责这人的邪恶，只指责他的软弱和单纯。有人可能会问：\Quote{这诞生与犹太人的事有什么共同之处？}但请问我，有什么不是共同的呢？因为最初被造的人、从肋骨形成的女人、不生育的妇人、以及藉着水所成就的事——我指的是厄里叟使铁器浮起的泉源、犹太人渡过的红海、天使搅动的池子、在约但河中被洁净的叙利亚人纳阿曼——这些都在事先以预像宣告了将要到来的诞生与净化。先知的话也暗示了这诞生的方式，如：\BibleQuote{一个世代要被报告给主，他们要来，他们要向主宣告祂的正义，给那将要出生的百姓，就是主所造的。}\BibleRef{咏 21:30} 以及：\BibleQuote{你的青春更新如鹰。}\BibleRef{咏 103:5} 以及：\BibleQuote{耶路撒冷，发光吧！看，你的君王来了！}\BibleRef{依 60:1} 以及：\BibleQuote{得赦其过者，是有福的。}\BibleRef{咏 31:1} 依撒格也是这诞生的预像。因为，尼苛德摩，请告诉我，他是如何诞生的？是按自然的律法吗？绝不是；他诞生的方式介于我们所说的自然之道之间：自然在于他是通过同房而生；另一面在于他不是由血气所生，而是由天主的旨意。我要指出，这些预像不仅提前宣告了这诞生，也宣告了由童贞女所生。因为没有人会轻易相信童贞女能生子，所以先有不生育的妇人，然后不仅是不能生育，而且是年迈的。女人由肋骨所造，确实比不育者怀孕更为奇妙；但因为是远古的事，所以给出了一个新的、新鲜的预像，即不育妇人的预像，为相信童贞女生子铺路。因此，耶稣为了提醒他这些事，说：\Quote{你是以色列的师傅，连这些事都不知道吗？}\par

% structure: p0017 source=h2 target=subsection reason=relative-heading-rank; base=h2

\subsection{\BibleRef{若 3:11}}

祂加上这话，用另一个论据使祂的话可信，并屈尊俯就对方的软弱。\par

3. 祂所说的\Quote{我们所说的，是我们知道的；我们所见证的，是我们见过的}是什么意思？因为对我们来说，视觉是最可靠的感官；如果我们想让人相信，我们会说我们亲眼看见了，而不是听说；因此基督按人的方式对他说话，也藉此使祂的话可信。这确实如此，并且祂没有想建立别的，也没有指肉眼的视觉，这一点很清楚：在说了\Quote{由肉生的属于肉；由神生的属于神}之后，祂加上\Quote{我们所说的，是我们知道的；我们所见证的，是我们见过的}。然而这（关于圣神的事）尚未发生；那么祂怎么说\Quote{我们见过的}呢？难道不是说明这知识是确切无误的吗？\par

\Quote{没有一个人接受我们的见证。}祂使用\Quote{我们知道}这个措辞，或用之于祂自己和圣父，或仅用于祂自己；而\Quote{没有人接受}这话，并非出于不悦，而是陈述一个事实：因为祂并没有说：\Quote{还有什么比你们这些不接受我们如此确切宣告的人更愚昧呢？}而是藉着祂的作为和言语，显示全然的温柔，没有说出任何类似的话；祂温和而慈祥地预言将要发生的事，如此也引导我们趋向全然的温柔，并教导我们，当我们与任何人交谈而未能说服他们时，不要气恼或变得粗暴；因为一个心浮气躁的人不可能达成目的，反而会使听他说话的人更加不信。因此，我们必须戒绝愤怒，使我们的言语在各方面都可信，不仅避免发怒，也要避免高声说话，因为高声是激情的燃料。\par

让我们捆住这匹马，好能制服骑手；让我们剪掉愤怒的翅膀，如此邪恶就不再升高。愤怒是一种强烈的激情，强烈且善于偷窃我们的灵魂；因此我们必须从各方面提防它的进入。我们能够驯服野兽，却忽略自己野蛮的心灵，这是奇怪的。愤怒是烈火，吞噬一切；它伤害身体，毁灭灵魂，使人容貌扭曲丑陋；如果愤怒的人能在愤怒时看见自己，他就不需要别的劝诫了，因为没有什么比愤怒的面容更令人不悦。愤怒是一种醉酒，甚至比醉酒更严重，比被魔鬼附身更可怜。但如果我们小心不要高声说话，就会发现这是通往节操行为的最佳途径。因此保禄要除去喧哗以及愤怒，他说：\Quote{你们应断绝一切愤怒和喧哗。}\BibleRef{弗 4:31} 让我们服从这位一切智慧的导师，当我们对仆人发怒时，让我们思想自己的过犯，并因他们的忍耐而羞愧。因为当你傲慢无礼，而你的仆人默默忍受你的侮辱，当你行为不端，他却如同智者一样，把这当作警告而不是别的。虽然他是你的仆人，但他仍是一个人，拥有不死的灵魂，并被共同的主赐予与你相同的恩赐。如果他——在更重要、更属灵的事上与我们平等——因着一些贫乏而微不足道的人性优越，如此温顺地承受我们的伤害，那么我们有什么可原谅的，有什么可辩解的？我们不能（或更确切地说，不愿）因敬畏天主而像他因敬畏我们那样有智慧。考虑到这一切，并记起自己的过犯和人类共同的本性，让我们时刻小心说话温和，好能心中谦卑，为我们的灵魂找到安息，无论是现在的，还是将来的；愿我们都能获得这安息，藉着我们的主耶稣基督的恩宠和慈爱，愿光荣藉着祂归于圣父及圣神，从今世直到永远。阿们。\par

\SourceRecord{Translated by Charles Marriott.；Nicene and Post-Nicene Fathers, First Series；Vol. 14.；Edited by Philip Schaff.；Buffalo, NY: Christian Literature Publishing Co.,；1889.；<http://www.newadvent.org/fathers/240126.htm>.}

% structure: p0001 source=h1 target=section reason=page-title

\section{若望福音讲道集 第二十七篇}

% structure: p0002 source=h2 target=subsection reason=relative-heading-rank; base=h2

\subsection{\BibleRef{若 3:12-13}}

一、我曾多次说过的话，现在要再说，并且不会停止再说。那是什么呢？就是：耶稣在将要论及崇高的道理时，常因听众的软弱而克制自己，不长久停留在配得上祂伟大身份的主题上，反而停留在那些含有屈尊俯就的主题上。因为崇高而伟大的事，只消说出来一次，就足以建立那样的特质，就我们所能听受而言；但除非不断说出更卑微的、更接近听众理解力的话，那些更崇高的事就不容易抓住一个属地的听众。因此，基督的话中，卑微的多于崇高的。然而，为免这又造成另一种伤害，即把门徒滞留在下界，祂并非不先说明祂为何说出这些话，就只将祂较低微的话摆在人前；事实上，祂在此处正是这样做的。因为当祂说了关于圣洗以及在地上因恩宠而生的重生之后，渴望让他们进入祂那自己的、神秘而不可理解的重生时，祂暂时悬置，不让他们进入，然后告诉他们祂不让他们进入的原因。那是什么呢？就是听众的迟钝和软弱。祂指着这一点又加了这些话：\BibleQuote{我对你们说过地上的事，而你们不相信；若我对你们说天上的事，你们又怎能相信呢？}因此，无论祂说什么平凡谦卑的话，我们都要归因于听众的软弱。\par

有人说，\Quote{地上的事}在这里是指风而言；即是说：\Quote{我若从地上的事给你们举例，你们尚且不信，那么你们怎能学习更崇高的事呢？}祂在此称圣洗为\Quote{地上的}事，你们不必惊奇，因为祂这样称呼它，或是因为它是在地上施行，或是因为这样称呼它是相对于祂自己那最可敬畏的重生而言的。因为虽然我们的这个重生是属天的，但相较于那出自圣父实体的真正重生，它仍是属地的。\par

祂并没有说：\Quote{你们没有明白}，而是说：\Quote{你们没有相信}；因为当一个人对那些可能通过理智理解的事物持不良态度，不愿轻易接受时，他可以被公正地指责为缺乏理解力；但当他不接受那些不能通过推理、只能通过信德理解的事物时，对他的指责就不再是缺乏理解力，而是不信。因此，祂引导他远离用推理来探究已说过的事，用指责他缺乏信德的方式更严厉地触动他。如果我们自己的重生必须通过信德接受，那么那些忙于用推理探究独生者受生的人，该当何罪？\par

但或许有人会问：\Quote{如果听众不信这些话，为什么还要说出来呢？}因为虽然\Quote{他们}不信，但后来的人会相信并从其中获益。因此，基督极其严厉地触动他，继续表明祂不仅知道这些事，还知道其他更多、远超过这些的事。祂藉着接下来的话宣告了这一点，当祂说：\BibleQuote{没有人上过天，除了那自天降下而仍在天上的人子。}\par

\Quote{这是怎样的后续呢？}有人问道。这是非常紧密的，与前面所说的完全一致。因为尼苛德摩曾说过：\BibleQuote{辣彼，我们知道祢是由天主而来的师傅}，正是在这一点上，祂纠正了他，几乎是说：\Quote{不要以为我像许多出自地上的先知那样是一位师傅，因为我（现在）是从天上来的。没有一位先知升到那里，而我却住在那里。}你们看见吗？即使那看似非常崇高的说法，也完全配不上祂的伟大？因为祂不仅在天上，而是在各处，充满万有；但祂仍然按照听者的软弱说话，渴望逐渐引领他上升。在此处，祂并未称肉体为\Quote{人子}，而是，可以这样说，从较低的本体来称呼祂整个的自我；事实上，这是祂的习惯，常从祂的天主性，也常从祂的人性来称呼祂的整个位格。\par

% structure: p0008 source=h2 target=subsection reason=relative-heading-rank; base=h2

\subsection{\BibleRef{若 3:14}}

这一点似乎也依赖于前面所说的，并且与它有非常紧密的联系。因为在说了圣洗带给人的极大恩惠之后，祂接着提到另一种恩惠，这是前者的原因，并且不比它逊色；即藉着十字架。如同保禄在论及格林多人时，也把这些恩惠放在一起，他说：\BibleQuote{保禄为你们被钉死了吗？或者你们是因保禄的名受洗的吗？}因为这两件事最显明祂那不可言喻的爱：祂既为仇敌受苦，又为仇敌死了之后，藉着圣洗白白地赐给他们完全的罪之赦免。\par

二、但为什么祂没有明说：\Quote{我将要被钉在十字架上}，而是把听众引向古代的预象呢？第一，为使你们学到古事与新知是相近的，彼此并不隔绝；其次，为使你们知道祂并非不情愿地走向祂的苦难；此外，除了这些理由，还为使你们学到，从这事实本身祂并不受任何损害，反而为许多人带来救恩。因为，为免有人说：\Quote{那相信一个被钉死的人怎么可能得救呢？当他自己被死亡所束缚时？}祂引导我们走向古代的故事。如今，如果犹太人藉着瞻仰铜蛇的像逃脱了死亡，那么相信被钉者的人，更有理由享有更大的恩惠。因为这事的发生，并非由于被钉者的软弱，也不是因为犹太人比祂强壮，而是因为\BibleQuote{天主爱了世界}，因此祂的活殿被钉在十字架上。\par

% structure: p0011 source=h2 target=subsection reason=relative-heading-rank; base=h2

\subsection{\BibleRef{若 3:15}}

你看到被钉十字架的起因和由此而来的救恩了吗？你看到预像与实体的关系了吗？在那里，犹太人逃脱了死亡，但只是暂时的；在这里，信徒逃脱了永恒的死亡。在那里，挂起的蛇治愈了蛇的咬伤；在这里，被钉的耶稣治愈了灵性毒龙所伤的创伤。在那里，肉眼观看的人得到痊愈；在这里，以理解之眼观看的人除去了一切罪恶。在那里，挂起的是铜制成的蛇像；在这里，是主的身体，由圣神所建造。在那里，蛇咬伤，蛇治愈；在这里，死亡毁灭，死亡拯救。但毁灭人的蛇有毒液，拯救人的蛇无毒；同样，在这里，杀害我们的死亡带有罪恶，如同蛇带毒液；但主的死亡完全没有罪恶，如同铜蛇没有毒液。因为\Person{伯多禄}说：\BibleQuote{祂没有犯过罪，口中也找不到诡诈。}\BibleRef{伯前 2:22}\Person{保禄}也宣告说：\BibleQuote{解除了率领者和掌权者的武装，公然羞辱他们，靠着十字架向他们夸胜。}\BibleRef{哥 2:16}正如一个高贵的冠军，将对手高举并摔下，使得胜利更加荣耀；同样，基督在全世界面前，将敌对权势摔下，并治愈了旷野中被击打的人，将他们从一切折磨他们的毒兽中救出，借着被钉在十字架上。但祂没有说\Quote{必须悬挂}，而是说\Quote{必须被举起}\BibleRef{宗 28:4}；因为祂用了这个似乎更温和的词，为了听众的缘故，也因为它符合预像。\par

% structure: p0013 source=h2 target=subsection reason=relative-heading-rank; base=h2

\subsection{若 3:16}

祂所说的话是这样：不要因我将被举起以使你得救而惊奇，因为这为父所喜悦，祂如此爱你们，竟将祂的儿子赐给奴隶，忘恩负义的奴隶。然而一个人甚至不会为朋友这样做，即使是义人也不轻易，正如\Person{保禄}所说：\BibleQuote{为义人死，是罕有的事。}\BibleRef{罗 5:7}现在他说话较长，因为是向信徒说话，但在这里基督说话简洁，因为是对\Person{尼苛德摩}讲话，但仍然更有意义，因为每个词都有深意。因为\Quote{如此爱}和另一句\Quote{天主爱世界}，祂显示出祂爱的大力量。两者之间的差距巨大无穷。祂，不朽的，无始的，无限威严；他们不过是尘土和灰烬，充满成千上万的罪恶，忘恩负义，时刻得罪祂；而祂\Quote{爱}了他们。同样，之后添加的话也很有意义，祂说\Quote{祂赐下了祂的独生子}，不是仆人，不是天使，不是总领天使。虽然如此，没有人会为自己的孩子如此焦虑，如同天主为祂忘恩负义的仆人那样。\par

然后祂将祂的苦难不是十分公开地放在他面前，而是较为隐晦地；但祂更清楚地加上苦难的好处，说：\BibleQuote{叫一切信祂的人，不致丧亡，反而获得永生。}因为当祂说了\Quote{必须被举起}并暗示死亡时，免得听众因这些话而沮丧，对祂生出一些纯粹属人的想法，以为祂的死亡是归于无有，请看祂如何纠正这一点，说那被赐下的是\BibleQuote{天主子}，是生命的源头，永生的源头。祂藉死亡为别人获得生命，自己岂会常在死亡中？因为如果信被钉者的人不致丧亡，被钉者自己就更不会丧亡。祂除去别人的缺乏，自己就更加远离缺乏；祂赐给别人生命，自己就更涌流生命。你看到处处都需要信德吗？因为祂称十字架为生命之泉；这为理性不易承认，如同现在外邦人嘲笑所证明。但那超越理性软弱的信德，却能轻易接受并持守它。天主\BibleQuote{如此爱了世界}是从哪里来的？没有别的源头，只出于祂的良善。\par

3. 让我们现在因祂的爱而羞愧，因祂慈爱的极至而羞惭，因为祂为了我们竟没有怜惜自己的\OriginalTerm{独生子}{Only-begotten Son}，而我们却为了自己的伤害而吝惜财富；祂为我们赐下了自己的儿子，但为了祂，我们却甚至不能轻看金钱，甚至不为我们自己。这些事怎能得到宽恕呢？如果我们看见一个人为我们受苦受死，我们会把他放在所有人之上，将他算作我们最亲密的朋友，把我们所有的一切都交给他，并认为那是他的而不是我们的，即使这样，我们也并不认为我们给了他应得的回报。但对着\Person{基督}，我们连这种程度的正确感觉都没有。祂为我们舍了性命，为我们流出了祂宝贵的血，而我们原本既非善意也非良善，但我们甚至不为自己的缘故流出我们的金钱，还忽略了那位为我们而死的主，当祂赤身露体又作客旅时；谁将救我们脱离那将来的刑罚呢？因为假设不是\Person{天主}惩罚，而是我们自己惩罚自己；我们岂不投票反对自己吗？我们岂不把自己判入地狱的火中，因为容许那为我们舍命的人因饥饿而憔悴吗？但我何必提及金钱呢？即使我们有千万条性命，难道不该都为祂舍弃吗？然而即便如此，我们仍不能与祂的恩惠相称。因为那首先施恩的人，明确证明了他的善意，但那受了恩惠的人，无论作什么回报，都是还债，而不是施恩；尤其是那首先行善的人，是施恩于他的仇敌。而回报的人，既向恩人奉献礼物，自己也收获其果实。但即便如此也不能感动我们；我们比任何人都更愚蠢，把金项链戴在我们的仆役、骡子和马匹身上，却忽略了我们的主，祂赤身露体，挨家挨户行走，常站在我们的门口，向我们伸出双手，而我们却常常以无情的眼光看待祂；然而这些正是祂为我们所受的。祂甘愿饥饿，为的是你能得饱足；祂赤身行走，为的是给你提供不朽衣袍的材料，然而即便如此，你也不肯放弃自己的一分。你的一些衣服被虫蛀了，另一些则成为你箱柜的负担，对拥有者来说是无用的麻烦，而那位赐给你这些以及你所拥有的一切的主，却赤身露体。\par

然而也许你并不把它们收藏在你的箱柜里，而是穿在身上，用它来装扮自己。你由此得到什么好处呢？是让街上的人看见你吗？那又怎样呢？他们不会佩服你这穿着这样衣服的人，而是佩服供应衣服给穷人的那个人；所以如果你渴望被人佩服，就通过给别人衣服，你反而会得到无穷的赞扬。那时，\Person{天主}和人都会赞美你；现在没有人赞美你，反而所有人都嫉妒你，看见你的身体穿着华服，而灵魂却被忽视。娼妓也有装饰，她们的衣服往往比平常更昂贵华丽；但灵魂的装饰只属于那些活在美德中的人。\par

这些事我不断地说，并且不会停止说它们，并不那么因为我关心穷人，而是因为我关心你们的灵魂。因为他们将会得到一些安慰，如果不是从你们这里，也会从别处；甚至如果他们得不到安慰，而是死于饥饿，对他们的伤害也不会很大。贫穷和饥饿的消耗曾怎样伤害了\Person{拉匝禄}呢？但是，如果你们得不到穷人的帮助，没有人能把你们从地狱中救出来；我们要对你们说那曾对那个不断受煎熬却得不到安慰的富翁说的话。愿\Person{天主}使没有人听见那些话，而是让所有人都能进入\Person{亚巴郎}的怀抱；借着我们的\Person{主耶稣基督}的恩宠和慈爱，愿光荣藉着祂并同着祂，归于\Person{圣父}及\Person{圣神}，直到永远。阿们。\par

\SourceRecord{Translated by Charles Marriott.；Nicene and Post-Nicene Fathers, First Series；Vol. 14.；Edited by Philip Schaff.；Buffalo, NY: Christian Literature Publishing Co.,；1889.；<http://www.newadvent.org/fathers/240127.htm>.}

% structure: p0001 source=h1 target=section reason=page-title

\section{\WorkTitle{若望福音第二十八篇讲道}}

% structure: p0002 source=h2 target=subsection reason=relative-heading-rank; base=h2

\subsection{\BibleRef{若 3:17}}

1. 许多较为粗心的人，利用天主的慈爱来增加自己罪过的严重和轻蔑的程度，这样说：\Quote{没有地狱，没有将来的惩罚，天主赦免我们所有的罪。}为堵住这些人的口，一位智者说：\BibleQuote{不要说：祂的仁慈广大，祂必会因我众多的罪而平息；因为仁慈和义怒都来自祂，祂的震怒停留在罪人身上} \BibleRef{德 5:6}；又说：\BibleQuote{祂的仁慈如何伟大，祂的惩戒也如何严厉} \BibleRef{德 16:12}。\Quote{那么，}有人说，\Quote{祂的慈爱在哪里，如果我们按自己的罪过受报应？}我们确实要\Quote{按自己的罪过受报应}，听听先知和\Person{保禄}的宣告：一个说：\BibleQuote{祢要按每人的行为报应他} \BibleRef{咏 61:12}；另一个说：\BibleQuote{祂要按每人的行为报应他} \BibleRef{罗 2:6}。然而我们仍可看到，即使如此，天主的慈爱仍是伟大的：祂将我们的存在分为两个时期——今生和来世，使前者成为考验的定命，后者成为加冕的地方，即使在这点上，祂也显示了极大的慈爱。\par

\Quote{如何，以何种方式？}因为当我们犯了许多严重的罪，从青年到老年没有停止以千万的恶行玷污我们的灵魂时，祂没有向我们要求任何一份账目，反而藉着重生的洗礼赐予我们罪的赦免，并白白赐给我们义德和圣化。\Quote{那么，}有人说，\Quote{如果一个人自幼就被认为堪当奥迹，之后却犯下千万的罪，该如何？}这样的人应受更重的惩罚。因为我们若在入门后犯罪，为同样的罪所受的惩罚是不同的。保禄如此宣告说：\BibleQuote{谁弃绝了梅瑟法律，凭两三个见证人，应处死不获宽恕；那么，你们想，那践踏了天主子，将那盟约的血视为不洁，又冒犯了恩宠之神的人，应受怎样更重的惩罚呢？} \BibleRef{希 10:28}这样的人理当受更重的惩罚。然而即使为他，天主也开启了悔改之门，并赐予他许多洗净过犯的方法，只要他愿意。你想，这些慈爱的明证是什么？藉恩宠赦罪，不惩罚那在恩宠后犯了罪而应受惩罚的人，反而给他一段时辰和指定的时间得以清理。基于这一切，基督对尼苛德摩说：\BibleQuote{天主没有派遣子来审判世界，而是为拯救世界} \BibleRef{若 3:17}。\par

因为基督有两次来临，一次已经过去，一次尚未来到；两次的目的不同。第一次的来临不是要审查我们的行为，而是要赦免；第二次的目的不是赦免，而是审问。因此，关于第一次祂说：\BibleQuote{我来不是为审判世界，而是为拯救世界} \BibleRef{若 3:17}；但关于第二次，祂说：\BibleQuote{当人子在祂父的光荣中降来时，祂要把绵羊放在右边，山羊放在左边} \BibleRef{玛 25:31}。这些要进入永生，那些要进入永罚。然而，祂的第一次来临，按公义的准则，本是为审判。为什么？因为在祂来临之前，有自然律、先知、还有成文的律法、教导、无数的许诺、奇迹的显示、惩罚、报复，以及许多其他可能使人归正的事物，因此祂本应要求为这一切交账；但因祂是仁慈的，祂暂时宽恕而不追究。因为若祂追究，所有人都将立即被投入毁灭。因为经上说：\BibleQuote{所有的人都犯了罪，亏欠了天主的荣耀} \BibleRef{罗 3:23}。你看到祂的慈爱那不可言喻的浩大吗？\par

% structure: p0006 source=h2 target=subsection reason=relative-heading-rank; base=h2

\subsection{\BibleRef{若 3:18}}

然而，如果祂\Quote{没有来审判世界}，那么\Quote{不信的人已经受了审判}是怎么回事，如果\Quote{审判}的时刻还没有到呢？祂的意思或是：不信而不悔改这一事实本身就是一种惩罚（因为处于光明之外，本身就包含极重的惩罚），或是预先宣布将要发生的事。正如谋杀犯，虽然还没有被法官的判决定罪，但从事物的本质上已被定罪，不信者也是如此。既然\Person{亚当}在吃树果的那一天就死了；因为判决是这样说的：\BibleQuote{你吃树果的那一天，你必定要死} \BibleRef{创 2:17}；但他却活了。那么他怎样\Quote{死}了呢？是因判决；因事物的本质；因为那使自己处于惩罚之下的人，就已经在承受它的刑罚，即使暂时未实际承受，但在判决上已是如此。\par

以免任何人听见\Quote{我没有来审判世界}就以为自己可以犯罪不受罚，从而变得更加粗心，基督用\Quote{已经受了审判}这句话来制止这种轻率；并且因为\Quote{审判}是将来的、尚未临近，祂就把报应的恐惧拉近，描述惩罚已经到来。这本身就是大慈爱的标记：祂不仅赐下自己的子，甚至还延迟审判的时候，使那些犯了罪和那些不信的人，能够有能力洗净自己的过犯。\par

\Quote{那信从子的，不受审判。}他所\Quote{信}的，不是过分好奇的人；他所\Quote{信}的，不是多管闲事的人。但如果他的生活不洁，行为邪恶呢？\Person{保禄}特别指出，这样的人根本不是真正的信徒：\BibleQuote{他们承认认识天主，但在行为上却否认祂} \BibleRef{铎 1:16}。但基督在这里说，这样的人在这一方面不受\Quote{审判}；因为他的行为固然要受更重的惩罚，但既然一度相信了，他就不会因不信而受罚。\par

2. 你看见祂如何以令人畏惧的事开始祂的讲论，又以同样的事结束吗？因为起初祂说：\BibleQuote{人若不从水和圣神而生，不能进天主的国}；这里又说：\BibleQuote{不信从子的，已经受了审判。}祂说：\Quote{不要以为迟延对罪人有什么益处，除非他悔改，因为那不信的，将与被定罪和受惩罚的人无异。}\par

% structure: p0011 source=h2 target=subsection reason=relative-heading-rank; base=h2

\subsection{若 3:19}

祂所说的是这样的意思：\Quote{他们受罚，是因为他们不愿离开黑暗，奔向光明。}于是祂继续剥夺他们将来的一切借口：\Quote{假若我来是为惩罚和追究他们的行为，他们或许可以说：\InnerQuote{这就是我们躲避祢的原因}，但如今我来是为将他们从黑暗中释放，领他们进入光明；那么，谁还会怜悯一个不愿从黑暗走向光明的人呢？他们既无法指控我们，又接受了无数的恩惠，却仍躲避我们。}祂在另一处也提出了这项指控，说：\BibleQuote{他们无故地恨我}\BibleRef{若 15:25}；又说：\BibleQuote{我若没有来向他们说话，他们就没有罪。}\BibleRef{若 15:22} 因为在光明未到时坐在黑暗中的人，或许可得宽恕；但光明来临后仍留在黑暗中，便证明了他自己的乖戾和好辩的性情。接着，因为祂的话对大多数人来说似乎难以置信（因为没有人会宁愿选择\BibleQuote{黑暗而不爱光明}），祂便指出了这种情感在他们心中的原因。那是什么呢？\par

% structure: p0013 source=h2 target=subsection reason=relative-heading-rank; base=h2

\subsection{若 3:19-20}

然而祂来不是为审判或查问，而是为宽恕和赦免过犯，并藉着信德赐予救恩。那么他们为何逃避呢？假若祂来并坐在审判席上，祂所说的或许合理；因为自知行为邪恶的人，通常会躲避审判官。但相反，那些犯罪的人甚至会奔向那施予宽恕者。因此，如果祂来是为宽恕，那些自知犯了许多过犯的人自然最应赶紧到祂面前；事实上，许多人的情况正是如此，因为连税吏和罪人也与\Person{耶稣}同席。那么祂所说的是指什么呢？祂是指那些始终愿意留在邪恶中的人。祂的确来了，为要赦免人从前的罪，并保护他们免于将来的罪；但既然有些人如此懈怠、如此无力于德行的劳苦，以至于愿意直到最后一口气都持守邪恶，永不停止，祂便在这里说话斥责他们。祂说：\Quote{因为基督徒的生活不仅要求正确的教义，还要求圣善的品行，他们害怕归向我们，因为他们不愿显示正义的生活。生活在异教中的人，无人指责，因为有了这样的神，以及像他们的神一样污秽可笑的仪式，他们所行的也符合教义；但属于真天主的人，若生活懈怠，便会受到所有人的责备和指控。真理甚至令它的仇敌如此钦佩。}请观察祂是如何精确地陈述祂的话。祂说：\Quote{不是\InnerQuote{行恶的人不来就光}，而是\InnerQuote{那始终行恶、始终愿在罪污中打滚的人，不愿服从我的律法，宁愿留在外面，无所畏惧地行淫，并做一切禁止的事。因为他若来到我这里，便如贼在光中暴露无遗，所以他躲避我的统治。}}例如，现今仍可听到许多异教徒说：\Quote{他们不能接受我们的信仰，因为他们不能放弃酗酒、淫乱和类似的放纵。}\par

有人说：\Quote{那么，难道没有作恶的基督徒和品行端正的异教徒吗？}我知道有作恶的基督徒，但我还不确知是否有品行端正的异教徒。请不要向我提及那些生性善良而有节制的人（这不是德行），而要告诉我那能够忍受强烈情欲冲击却保持节制的人。你找不到。因为如果天主之国的应许、地狱的威胁和如此多的准备都几乎不能使人持守德行，那么不相信这些的人更难追求德行。或者，如果有人假装这样做，那也只是为了炫耀；而为炫耀而行的人，在能逃避注意时，不会克制自己的恶欲。但为了不使人认为我们有争论之心，我们姑且承认异教徒中有品行端正的人；但这也不违背我的论点，因为我所说的是普遍情况，而非罕见事例。\par

请看祂如何以另一种方式剥夺他们的一切借口，当祂说\BibleQuote{光来到世上}时。祂说：\Quote{是他们自己寻找光吗？他们劳苦、努力了吗？光本身来到他们面前，即使如此他们也不肯奔向它。}如果有些基督徒生活邪恶，我可以说祂所说的不是指那些从一开始就是基督徒、并从祖先继承真宗教的人（虽然这些人也大多因邪恶生活而偏离了正确教义），但我仍认为祂现在所说的并非关于他们，而是关于异教徒和那些本应归向真信仰的犹太人。因为祂表明，没有一个生活在谬误中的人会选择归向真理，除非他事先为自己规划了正义的生活；也没有人会留在不信中，除非他先前选择了永远邪恶。\par

不要告诉我某人节制、不抢劫；这些事情本身并非德行。因为如果一个人拥有这些，却仍是虚荣的奴隶，并因畏惧朋友的交往而停留在他的错误中，这有什么益处呢？这不是正当的生活。声誉的奴隶并不比奸淫者罪过更少；相反，他比奸淫者做出更多更严重的行为。但请告诉我，有谁完全脱离了所有的情欲和不义，却仍留在外邦人中？你无法做到这一点；因为即使在那些夸耀大事、且据说已克服贪婪或贪饕的人中，大多数也是声誉的奴隶，而这是万恶之因。因此，犹太人也就继续作犹太人；为此基督斥责他们说：\BibleQuote{你们彼此寻求光荣，却不寻求从天主来的惟一光荣，怎能相信呢？}\BibleRef{若 5:44}\par

\Quote{请问，他为何没有对\Person{纳塔乃耳}谈论这些事，就是他曾为他作证真理的那一位，也没有将他的谈论延长？}因为连\Person{纳塔乃耳}也没有像\Person{尼苛德摩}那样带着热忱而来。\Person{尼苛德摩}把这当作他的工作，把别人用来休息的时间用作聆听的时间；但\Person{纳塔乃耳}是因别人的邀请而来的。然而，耶稣也并没有完全忽略他，因为他对他说：\BibleQuote{你要看见天开，天主的天使在人子身上上去下来。}\BibleRef{若 1:51}但耶稣对\Person{尼苛德摩}并没有这样说，而是与他谈论了救恩计划与永生，按照个人意志的状况，对每人做了不同的、合适的讲话。对\Person{纳塔乃耳}来说，因为他知道先知们的著作，也不那么胆怯，只需听到这些就够了；但\Person{尼苛德摩}仍然被恐惧所缠住，所以基督并没有清楚地向他启示整个事情，而是摇动他的心思，以便用恐惧驱除恐惧，宣布不信的人已被审判，不信来自邪恶的良心。由于他看重从人来的光荣甚于惩罚（福音作者说：\BibleQuote{许多} \BibleQuote{领袖中也有许多人信了他，但为了法利塞人不敢承认}\BibleRef{若 12:42}），基督就在这一点上触动他，说：\Quote{不信我的人，不可能有别的缘故，只因为他生活不洁。}随后他说：\BibleQuote{我是光}\BibleRef{若 8:12}，但在这里他说：\Quote{光来到了世上}；因为起初他说话有些隐晦，但后来更加清楚。即使如此，那人仍因顾及众人意见而被约束，因此不能按他所应当的那样勇敢地说话。\par

我们要逃避虚荣，因为这是一种比任何情欲都更专横的情欲。由此产生贪婪、爱财、仇恨、战争和争斗；因为想要超过自己所有的人，永远不能停止，而他之所以想要，没有别的原因，只是因为爱虚荣。请告诉我，为什么那么多人用一大群宦官、成群结队的奴隶和许多排场环绕自己？不是因为他们需要这些，而是为了让那些遇见他们的人成为这不合时宜的炫耀的见证。如果我们斩断了这个，我们就连同头一起杀死了邪恶的其他肢体，就没有什么能阻止我们在地上如同在天上生活。这个恶行不仅把它的俘虏推入邪恶，甚至与他们的德行共存，当它不能完全将我们从这些德行中逐出时，它仍在德行操练本身中给我们造成很大损害，迫使我们承受劳苦，却剥夺了果实。因为怀着这种目的而斋戒、祈祷、施舍的人，已得到了他的赏报。还有什么比这样的损失更可怜的呢？即人徒然无益地哀叹自己，成为可笑的对象，反而失去了天上的光荣？因为既追求两者的人不能两者兼得。的确，当我们不追求两者而只追求那从天而来的光荣时，两者是可能的；但贪求两者的人不能两者兼得。因此，如果我们希望得到光荣，让我们躲避人的光荣，只渴望那出自天主的；这样，我们就能得到两者。愿我们众人，因我们的主耶稣基督的恩宠和慈爱，偕同他，和他一起，归于圣父及圣神，光荣至于永远。阿们。\par

\SourceRecord{Translated by Charles Marriott.；Nicene and Post-Nicene Fathers, First Series；Vol. 14.；Edited by Philip Schaff.；Buffalo, NY: Christian Literature Publishing Co.,；1889.；<http://www.newadvent.org/fathers/240128.htm>.}

% structure: p0001 source=h1 target=section reason=page-title

\section{\WorkTitle{若望福音讲道集 29}}

% structure: p0002 source=h2 target=subsection reason=relative-heading-rank; base=h2

\subsection{\BibleRef{若 3:22}}

1. 没有什么比真理更清晰、更有力，正如没有什么比谎言更脆弱，即使它被千万层帷幕遮蔽。因为即便如此，它也容易被识破，容易消散。但真理毫无遮掩地展现在所有愿意观看其美丽的人面前；它不寻求隐藏，不惧怕危险，不因阴谋而颤抖，不渴求众人的赞誉，不受任何可朽之物的审判，反而超越一切，成为千万秘密阴谋的对象，却依然不可战胜，并以自身超越的力量守护着投奔它的人；它避开隐秘之处，将属于它的一切摆在众人面前。当基督与比拉多交谈时，祂宣告了这一点，说：\BibleQuote{我向来公开教导，在暗地里我没有说过什么。}\BibleRef{若 18:20}正如祂当时所说，现在祂也这样做了，因为，福音史家说：\BibleQuote{此后，耶稣和门徒出去，到了犹太地，在那里同他们住下，并施洗。}在节日期间，祂上到城市，在他们中间宣讲教义，并藉奇迹提供帮助；但节日过后，祂常常前往约但河，因为有许多人聚集在那里。因为祂总是选择最拥挤的地方，并非出于炫耀或虚荣，而是因为祂愿意向最多的人提供帮助。\par

然而福音史家后来又说，\BibleQuote{耶稣不亲自施洗，而是祂的门徒施洗}；由此可见，此处也是这个意思。为何耶稣不施洗？洗者若翰先前曾说过：\BibleQuote{祂要用圣神和火给你们施洗。}但那时祂尚未赐下圣神，因此祂没有施洗是合情合理的。但祂的门徒这样做，是因为他们愿意引导许多人归向这救恩的道理。\par

当耶稣的门徒施洗时，为何若翰没有停止施洗？为何他继续施洗，甚至直到他被投入监狱？\par

第 23 节。\BibleQuote{若翰也在艾农施洗}；并补充说，\BibleRef{若 3:23}\par

第 24 节。\BibleQuote{若翰还没有被下在监狱里}，\BibleRef{若 3:24}\par

因为说这些话是要表明，直到那时他都没有停止施洗。但他为何施洗到那时？因为如果他在门徒开始施洗时就停止，反倒会使耶稣的门徒显得更可敬。那么他为何还要施洗？那是为了不激发他的门徒产生更强烈的竞争，使他们更加好斗。因为，尽管他千万次宣告基督，将首位让给祂，使自己如此卑微，却仍无法说服他们投奔祂；如果他再添加这个举动，就会使他们更加敌对。正因如此，基督在若翰被移除后才更持续地开始宣讲。而且，我认为，若翰之死被允许，并且发生得非常迅速，是为了使群众的全部注意力转向基督，使他们不再对两人持有分歧的意见。\par

此外，即使在他施洗时，他也不断地劝诫他们，向他们展示耶稣的高贵和可畏。因为他给他们施洗，告诉他们的无非是要他们相信那在他以后来的那一位。那么，一个这样做的人，如果停止施洗，又怎能使基督的门徒显得可敬呢？相反，人们会以为他是出于嫉妒和激情而这样做。但继续宣讲提供了一个更强的证据；因为他不是为自己寻求荣耀，而是将听众送到基督那里，并与基督一同工作，不亚于基督自己的门徒，甚至比他们更多，因为他的见证是无可怀疑的，并且他在众人眼中的声望远比他们高得多。福音史家暗示了这一点，他说：\BibleQuote{全犹太和约但河一带的人都出来到他那里，受了他的洗。}\BibleRef{玛 3:5}即使门徒施洗时，仍有许多人不断地涌向他。\par

如果有人问：\Quote{门徒的洗礼在哪些方面比若翰的更好？}我们将回答：\Quote{毫无区别}；两者都没有圣神的恩赐，双方施洗的理由相同，就是引导受洗者归向基督。因为为了使他们不必总是跑来跑去召集那些应该相信的人，如同在西满的例子中他的兄弟所做，以及斐理伯对纳塔乃耳所做，他们设立了洗礼，以便藉此轻易地将所有人带到他们面前，并为将要来临的信仰预备道路。但这两次洗礼并无高低之分，这从下文可以看出。那是什么呢？\par

% structure: p0011 source=h2 target=subsection reason=relative-heading-rank; base=h2

\subsection{\BibleRef{若 3:25}}

因为若翰的门徒总是嫉妒基督的门徒和基督本人，当他们看到他们施洗时，就开始与那些受洗的人争论，好像他们的洗礼在某种程度上优于基督门徒的洗礼；他们抓住一个受洗的人，试图说服他这一点，但没有说服他。请听福音史家如何使我们明白，是他们攻击了他，而不是他发起了问题。他没有说\Quote{一个犹太人同他们辩论}，而是说\Quote{若翰的门徒同一位犹太人关于洁净的礼发生了争论}。\par

2. 请观察福音史家的温和无恶意。他没有用谩骂的方式说话，而是尽其所能缓和指责，仅仅说\Quote{发生了争论}；而后续的内容（他也以温和的方式记载了）清楚表明所说的是出于嫉妒。\par

% structure: p0014 source=h2 target=subsection reason=relative-heading-rank; base=h2

\subsection{\BibleRef{若 3:26}}

就是说，\Quote{你所施洗的那一位}；因为他们说\Quote{你所见证的那一位}时，暗示了这一点，就好像他们说\Quote{你所指出为显赫、使之著名的那一位，竟敢与你做同样的事。}然而他们不说\Quote{你所施洗的那一位}施洗（因为那样他们就不得不提及从天而来的声音和圣神的降临），而是说什么呢？\Quote{那曾与你在约旦河对岸、你所见证的那一位}，就是说\Quote{那居于门徒地位、与我们无异的这个人，竟自分离出去，并且施洗。}因为他们以为这样不但能使他嫉妒，还能断言他们自己的声望正在衰落。他们说\Quote{众人都到祂那里去了。}由此显然，他们并未在与那犹太人辩论中占上风；但他们说这些话是因为他们心意不完善，尚未摆脱竞争之心。那么\Person{若翰}做什么呢？他没有严厉责备他们，唯恐他们再次离开他，做出什么恶事。他说了什么？\par

% structure: p0016 source=h2 target=subsection reason=relative-heading-rank; base=h2

\subsection{\BibleRef{若 3:27}}

不要惊奇，若他论及基督时用低微的语气；要一下子完全教导那些被激情占据的人是不可能的。但他想要暂时以敬畏和恐惧击打他们，并表明他们与基督争战，就是与天主争战。这里他暗中肯定了\Person{加玛里耳}所宣称的那真理：\Quote{你们不能推翻它，免得你们甚至被发现是与天主作对。}\BibleRef{宗 5:39}因为说\Quote{没有人能领受什么，除非是天赐予他的}，无异于宣告他们在尝试不可能的事，因此会发现是与天主作对。\Quote{好吧，但是\Person{特乌达}和他的追随者不是\InnerQuote{从自己领受}吗？}他们确实领受了，但他们立刻就被驱散和毁灭了；基督所拥有的却并非如此。\par

借此他也温和地安慰他们，表明不是一个人，而是天主在尊荣上超越了他们；因此他们不必惊奇，若属于祂的事物是荣耀的，并且\Quote{众人都到祂那里去}：因为这是神圣事物的本性，并且是天主成就了这些事，因为从未有人有能力行如此的事。一切属人的事都容易看透、败坏，迅速消融和消亡；这些事却不然，因此不是属人的。注意，当他们说\Quote{你所见证的那一位}时，他把他们以为可以用来贬低基督的话转而攻击他们自己，并以显示耶稣的荣耀并非来自他的见证来使他们沉默；他说\Quote{人不能自己领受什么，除非是天赐予他的。}\Quote{如果你们多少坚持我的见证，并相信它是真的，要知道，借着那见证，你们应当优先的不是我，而是祂。因为我见证了什么呢？我请你们自己作证。}\par

% structure: p0019 source=h2 target=subsection reason=relative-heading-rank; base=h2

\subsection{\BibleRef{若 3:28}}

\Quote{那么，如果你们坚持我的见证（而且你们现在引用它，说\InnerQuote{你所见证的那一位}），祂非但不因接受我的见证而减损，反而因此被高举；此外，那见证不是我的，而是天主的。所以，如果我在你们眼中是可信的，我所说的还包括这一件：\InnerQuote{我是在祂前面被派遣的。}}你看他如何逐渐显示这声音是神圣的？因为他所说的是这一类的：\Quote{我是仆人，说的是那派遣我者的话，并非出于人的偏爱而奉承基督，而是事奉那派遣我的圣父。我并非作为礼物给出见证，而是我被派去说什么，就说什么。不要因此以为我伟大，因为这显示祂是伟大的。祂是万有的主。}他接下来继续宣告这一点，并说：\par

% structure: p0021 source=h2 target=subsection reason=relative-heading-rank; base=h2

\subsection{\BibleRef{若 3:29}}

\Quote{但是，那位曾说\InnerQuote{我连替祂解鞋带也不配}的，如今怎么称自己是祂的\InnerQuote{朋友}呢？}他说这话并非为了高抬自己，也不是夸口，而是为了表明他也大力促进这一点（即基督的高举），并且这些事的发生并非违背他的意愿或使他忧伤，而是他渴望并热切期望它们，并且他的一切行动都是专门为此而做的；他非常智慧地用了\Quote{朋友}一词来表明这一点。因为在婚姻之类的场合，新郎的仆人并不如他的\Quote{朋友}那样欢喜快乐。他并非出于证明地位平等的愿望（绝无此意），而只是出于极度的喜悦，并且更是为了迁就他们的软弱，才称自己为\Quote{朋友}。他之前已藉着说\Quote{我是在祂前面被派遣的}宣告了他的服务。为此，又因为他们以为他对所发生的事感到烦恼，他称自己为\Quote{新郎的朋友}，以表明他不但不烦恼，反而大大喜乐。\Quote{因为，}他说，\Quote{我来正是为了成就这事，我非但不为所做的事忧伤，反而若不成就，我就会大为忧伤。倘若新娘没有来到新郎那里，我便会忧伤，但现在不会，因为我的任务已经完成。当祂的仆人前进时，我们才获得尊荣；因为我们所渴望的已经成就，新娘认识新郎，你们也为此作证，当你们说：\InnerQuote{众人都到祂那里去了。}我切望这事，我为此做了一切；如今我看见它成就了，我就欢喜、快乐、踊跃。}\par

3. 但\Quote{那站着并听见祂的，因新郎的声音而大大喜乐}是什么意思？他将比喻的表述转移到当前的论题上；因为在提及新郎和新娘之后，他展示了新娘如何被领回家，即藉着\Quote{声音}和教导。因为教会就是这样与天主结合的；因此保禄说，\BibleQuote{信德来自听闻，听闻来自天主的言语。}\BibleRef{罗 10:17}\Quote{因这\InnerQuote{声音}，}他说，\Quote{我喜乐。}他不是无故地用了\Quote{站着}，而是要表明他的职务已经止息，他已将\Quote{新娘}交付给祂，并且将来必须站着听祂；他是一位仆人和管家；他的美好希望和喜乐如今已经成就。因此他说，\par

\Quote{因此我的喜乐得以满全。}\par

\Quote{应由我完成的工作已经结束，今后我不能再做什么了。}\par

% structure: p0026 source=h2 target=subsection reason=relative-heading-rank; base=h2

\subsection{若 3:30}

\Quote{我的如今已停滞，并从此终止，但祂的却增长；因为你们所惧怕的，将不只在现在，而是在将来会更多。而正是这一点尤其彰显了我的工作更为光辉；我为此而来，我欢喜的是祂的进展如此之大，而那些我所有行动所促成的事已经实现。}你们看见他是如何温和而极明智地缓和了他们的情绪，扑灭了他们的嫉妒，向他们表明他们是在尝试不可能的事——这是遏制邪恶的最佳方法吗？为此，这些事发生在若翰还在世且施洗的时候，是为让他的门徒有他作为基督优越性的见证，这样，如果他们不信，就无可推诿。因为若翰并非主动说这些话，也不是回答其他询问者，而是他们自己提出了问题，并听到了回答。如果他自发的说话，他们的信心就不如在听到他回答他们问题时所得的自责判断那样强烈；正如犹太人，他们从家中派人到他那里，听到了他们所做的，却仍不信，尤其使自己失去了借口。\par

那么我们由此学到什么？那就是对虚荣的疯狂欲望是万恶之源；这导致了嫉妒，而当嫉妒稍息时，又再次激起。因此他们来到耶稣面前，说：\BibleQuote{为什么你的门徒不禁食？}\BibleRef{玛 9:14}因此，亲爱的，让我们避免这种欲望；因为如果我们避免它，我们将逃脱地狱。因为这恶习尤其点燃地狱之火，处处扩展其统治，暴虐地占据每个年龄和每个阶层。它搅乱了教会，在国家事务中为害，推翻了家庭、城市、人民和民族。你何以惊奇？它甚至进入了旷野，并在那里显出了它的大能。因为那些彻底告别财富和世间一切繁华、不与任何人交谈、战胜了肉体上更霸道的欲望的人，这些人被虚荣心俘虏，常常丧失一切。由于这种虚荣心，一个辛勤劳作的人可能变得比一个根本没有劳作、反而犯了万种罪的人更糟；法利塞人比税吏更糟。然而，谴责这种虚荣心很容易（人人都同意这样做），但问题是如何战胜它。我们怎能做到？以荣誉对抗荣誉。因为正如我们因仰望另一种财富而轻视地上的财富，因思想比这一生命远为优越的那一生命而轻视这一生命，同样，当我们认识另一种远比它庄严的光荣——那才是真正的光荣时，我们就能唾弃这世界的虚荣。一种是虚浮空洞的，有名无实；但从天上来的那一种是真实的，并能得到天使、总领天使和总领天使之主的赞美，或者我应该说，也有人的赞美。现在，如果你注目于那个舞台，学习那里有什么冠冕，将自己转移到从那里来的掌声，那么地上的事物永远不能束缚你；它们来临时，你不以为大；它们离去时，你不追求。因为甚至在世上的宫殿中，站在国王周围的卫士中，没有人会忽略取悦那位戴着王冠、坐在宝座上的君王，而去操心乌鸦的叫声，或苍蝇和蚊虫在他周围飞翔嗡嗡的声音；而来自人的美誉并不比这些更好。因此，认识人事的虚妄，让我们把我们的一切聚集在不被破坏的库房里，让我们寻求那持守不动摇的光荣；愿我们众人都能获得，藉着我们的主耶稣基督的恩宠和慈爱，愿光荣藉着祂、并与祂一起归于圣父和圣神，从今直到永远，永世无穷。阿们。\par

\SourceRecord{Translated by Charles Marriott.；Nicene and Post-Nicene Fathers, First Series；Vol. 14.；Edited by Philip Schaff.；Buffalo, NY: Christian Literature Publishing Co.,；1889.；<http://www.newadvent.org/fathers/240129.htm>.}

% structure: p0001 source=h1 target=section reason=page-title

\section{若望福音讲道 第30篇}

% structure: p0002 source=h2 target=subsection reason=relative-heading-rank; base=h2

\subsection{\BibleRef{若 3:31}}

1. 贪求光荣是一件可怕的事，可怕且充满诸多邪恶；它是一根难以拔除的荆棘，一头难以驯服的多头野兽，武装起来对付喂养它的人；因为正如虫蛀蚀它所从生的木头，锈腐蚀它所从出的铁，蛾蛀蚀羊毛，虚荣也摧毁滋养它的灵魂；因此我们需要极大的勤勉来除掉这种情欲。请看，\Person{若翰}对那些受了这种影响的门徒使用了多么长久的劝导，却几乎不能安抚他们。因为他除了已提到的那些话外，还用另外的话来软化他们。这些话是什么呢？他说：\Quote{那从上面来的，} \Quote{是在万有之上；那从地上来的，是属于地，并且讲论地。} 既然你们以我的见证为大事，并因此说比我更可信，你们必须知道这一点：一位从天上来的，不可能因一位住在世上的人而使祂的可信度得到加强。\newline{}\par

而\Quote{在万有之上}是什么意思？这句话想向我们表明什么？那就是\Person{基督}一无所缺，祂自给自足，并且远超过万有；\Person{若翰}说自己是\Quote{属于地，并且讲论地}。这不是说他凭自己的心思讲话，而是如\Person{基督}所说：\Quote{我对你们说地上的事，你们尚且不信，} 祂如此称呼圣洗圣事，不是因为它是\Quote{地上的事}，而是因为祂讲话时拿它比作自己的\OriginalTerm{不可言喻的受生}{Ineffable Generation}；同样，\Person{若翰}说他讲论\Quote{地}，是将自己的教导与\Person{基督}的教导相比较。因为\Quote{讲论地}的意思无非是：\Quote{我的事与祂的事相比是渺小、卑微、贫乏的，且是属地的本性所能领受的。在祂内\InnerQuote{藏着一切智慧和知识的宝藏。}} \BibleRef{哥 2:5} 他说这不是出于人的推理，这一点是显而易见的。他说：\Quote{那从地上来的，是属于地。} 然而他里面并非一切都是属地的，而是较高的部分是属天的，因为他有灵魂，并分享了非属地的\Person{圣神}。那么他怎能说自己是\Quote{属地的}呢？你难道不明白他只是在说：\Quote{我是渺小、不足道的，行走在地上，生在地上；但\Person{基督}从上面来到我们中间。} 他用这一切方法平息了他们的情欲之后，便更公开地讲论\Person{基督}；因为在此之前，说出那些无法在听者理解中扎根的话是无用的；但当他拔除了荆棘之后，便大胆地撒下种子，说：\newline{}\par

% structure: p0005 source=h2 target=subsection reason=relative-heading-rank; base=h2

\subsection{\BibleRef{若 3:31-32}}

他说了关于祂伟大而崇高的话之后，又将讲论降低到更谦卑的层面。因为\Quote{祂所听见所看见的}这种说法更适合一个普通人。祂所知道的，不是通过看见或听见而学到的，而是祂的本性包含了全部，从\Person{圣父}的怀里完美地出来，无需任何人教导。因为祂说：\Quote{\Person{圣父}认识我，我也认识\Person{圣父}。} \BibleRef{若 10:13} 那么\Quote{祂讲祂所听见的，并见证祂所看见的}是什么意思呢？由于我们通过这些感官正确认识一切事物，并且当我们教导那些我们的眼睛所看见、耳朵所听到的事情时，我们认为可靠，因为在这种情况下我们不会编造或说谎，\Person{若翰}在此想要确立这一点，便说：\Quote{祂所听见所看见的}，即\Quote{凡出于祂的没有虚假，全都是真实的。} 同样，当我们对某事好奇追问时，常常问：\Quote{你听见了吗？} \Quote{你看见了吗？} 如果这一点得到证实，那见证就是无可置疑的；同样，当\Person{基督}自己说：\Quote{我所听见的，我就这样审判} \BibleRef{若 5:30}；\Quote{我从\Person{圣父}所听见的，我就讲论} \BibleRef{若 15:15}；\Quote{我们说我们所看见的} \BibleRef{若 3:11}；以及祂所说的其他类似话语，都不是要使我们以为祂是因受教于谁而说这些话（这样想是极其愚蠢的），而是为了不让无耻的犹太人怀疑所说的话。因为他们对祂还没有正确的看法，祂便不断归依于\Person{圣父}，并由此使祂的话可信。\newline{}\par

2. 你若怪祂归依于\Person{圣父}，而祂也常常求助于先知和圣经，这一点都不奇怪？正如祂说：\Quote{他们是给我作见证的。} \BibleRef{若 5:39} 难道我们会因此说祂比先知低，因为祂从他们那里引用见证吗？绝不可这样想。这是由于听者的软弱，祂才如此安排讲论，并说祂所讲的，是听自\Person{圣父}的，并非祂需要一位老师，而是为叫他们相信祂所说的没有虚假。\Person{若翰}的意思是这样的：\Quote{我愿意听祂所说的话，因为祂从上面来，从那里带来那些唯独生命正确知道的讯息；因为\InnerQuote{祂所看见所听见的}是宣告这一点的人的表达。}\newline{}\par

\Quote{没有人接受祂的见证。} 然而祂有门徒，并且许多人留心听祂的话。那么\Person{若翰}怎说\Quote{没有人}呢？他说\Quote{没有人}，意思是\Quote{很少人}；因为若他是指\Quote{完全没有一个人}，他怎能加上，\par

% structure: p0009 source=h2 target=subsection reason=relative-heading-rank; base=h2

\subsection{\BibleRef{若 3:33}}

在此，他触及自己的门徒，因为他们在一段时间内似乎不会坚定地相信。此后他们仍未相信祂，从后来所说的话中可以清楚看出；因为若翰即使身在狱中，还派他们到基督那里，好更使他们依附于祂。然而，他们即使在那时也几乎不相信，基督提及这点时曾说：\Quote{凡不因我而绊倒的，是有福的。}\BibleRef{玛 11:6}因此，他现在说：\Quote{没有人接受祂的见证，}为要坚定自己的门徒；几乎是说：\Quote{不要因为一时很少有人信祂，就认为祂的话是假的；因为\InnerQuote{祂所说的，是祂所看见的。}}此外，他说这话也是为了触动犹太人的麻木。这是福音作者在开始时对他们提出的控诉，说：\Quote{祂来到自己的地方，自己的人却没有接受祂。}\BibleRef{若 1:11}因为这不是对祂的责备，而是对那些不接受祂的人的指控。\par

\Quote{那接受祂见证的人，就以印证实天主是真实的。}这里他也以显示不信任祂的人不仅不信祂，也不信圣父，来恐吓他们；因此他接着说：\par

% structure: p0012 source=h2 target=subsection reason=relative-heading-rank; base=h2

\subsection{若望福音 3:34}

既然祂说祂的话，信与不信的人，就是信或不信天主。\Quote{以印证实}，即\Quote{声明了}。然后，为增加他们的恐惧，他说：\Quote{天主是真实的；}从而显示，没有人能在不信基督的同时，不使派遣祂的天主成为说谎者。因为既然祂所说的没有不是出自圣父的，但祂所说的都是圣父的，那么不听祂的，就是不听派遣祂的那一位。看，这些话语又如何使他们战栗。那时他们仍认为不听基督不算什么大事；因此他在不信者头上悬着如此大的危险，好让他们明白，谁不听基督，就是不听天主本身。然后他继续讲论，降低到他们软弱的程度，说：\par

\Quote{因为天主赐圣神没有限量。}\par

再次，如我所说，他将讲论降到较低层面，使其变化并适合当时听众的接受；否则，他无法提升他们并增加他们的恐惧。因为如果他关于耶稣本身说了任何伟大崇高的事，他们不会相信，甚至可能轻视祂。因此他将一切引向圣父，暂时像谈论人一样谈论基督。但他说\Quote{天主赐圣神没有限量}是什么意思呢？他要显示我们众人领受圣神的运作是有限量的（此处他以\Quote{圣神}指圣神的运作，因为运作是分施的），而基督却拥有其全部运作，没有限量，且是完整的。如果祂的运作没有限量，那么祂的本质更是无限。你也看到圣神是无限的吗？那么，那位领受了圣神全部运作、知道天主的事、说\Quote{我们所说的，是我们所听见的；我们所见证的，是我们所看见的}\BibleRef{若 3:11}的祂，怎能被怀疑呢？祂所说的没有不是\Quote{从天主}来的，没有不是\Quote{从圣神}的。他暂时没有说任何关于天主圣言的事，而是藉着（提及）圣父和圣神使他的全部教导可信。因为他们知道有一位天主，也知道有一位圣神（尽管他们对祂没有正确的认识），但他们不知道有一位圣子。正是为此，他不断求助于圣父和圣神，从而证实他的话。如果有人不考虑这个理由，而单独审视他的话，那么他的话会远远低于基督的尊严。基督并不因为祂有圣神的运作而配得他们的信心（因为祂不需从那里得到帮助），而是祂本身自足；只是洗者若翰暂时按照较简单的人的理解说话，渴望逐步将他们从低下的观念提升起来。\par

我说这些，是要我们不要粗心地放过圣经所载的内容，而要全面考虑讲者目的、听众的软弱，以及其中的许多其他要点。因为教师并非完全照自己的意愿说话，而一般按照他们软弱（听众）的情况所需。因此保禄说：\Quote{我从前不能把你们当作属神的人，只得把你们当作属血肉的人；我喂养你们的是奶，不是饭食。}\BibleRef{格前 3:1}他的意思是：\Quote{我本来愿意把你们当作属神的人说话，却不能}；不是因为他不能，而是因为他们不能那样听。同样，若翰愿意教导门徒一些伟大的事，但他们还不能承受接受，因此他大多停留在较卑微的层面上。\par

因此，我们必须仔细探究一切。因为圣经的话是我们的属灵武器；但我们若不知道如何装配这些武器，以及如何正确装备我们的学者，它们虽保持其本有的力量，却无法帮助那些接受它们的人。举例说，假设有坚固的铠甲、头盔、盾牌和长矛；有人拿起这些装备，把铠甲穿在脚上，头盔戴在眼睛上而不是头上，不把盾牌护在胸前，反而颠倒地绑在腿上——他能从这些武器中得到任何益处吗？他不是反而会受害吗？任何明眼人都知道必然如此。但这并非由于武器的薄弱，而是由于那不善使用它们之人的笨拙。同样，对于圣经，如果我们错乱了它们的次序，它们虽仍保持其本有的力量，却对我们毫无益处。虽然我常在私下和在公开场合向你们讲论这事，却毫无效果；反见你们终其一生被钉于今世的事务上，连做梦也未曾想到属灵的事。因此，我们的生命变得漫不经心，我们这些为真理奋斗的人力量微薄，成了希腊人、犹太人和异端者的笑柄。倘若你们在其他事上疏忽，而在此时此地显露出与别处同样的冷漠，那么你们的行为尚且不可原谅；但如今在今生的事务上，你们每个人，无论工匠或政客，都比刀剑更锋利，而在必要和属灵的事上，我们却比任何人都迟钝；把正事当作副业，而不把本该视为比任何事务更紧迫的事，甚至当作副业看待。你们岂不知圣经不仅为最初的人类而写，也为了我们吗？你们岂没听见\Person{保禄}说：\BibleQuote{它们是为我们这些末世的人所写的警戒；使我们因圣经的忍耐和安慰，常怀有望}？\BibleRef{格前 10:11} 我知道我是在白费唇舌，但我仍不会停止说话，因为这样我在天主面前可以自辩，即使没有人听我。那向留心的人说话的人，至少还有听众的顺服来鼓励他的言语；但那人不断说话却不被听从，仍不停止说话，也许比前者更配得尊荣，因为他尽己所能履行天主的旨意，即使没有一人听从。尽管如此，虽然因你们的悖逆，我们的赏报会更大，我们却宁愿它减少，而你们的得救增进，认为你们被称赞乃是巨大的赏报。我们现在说这些，并非要使我们的讲道令你们痛苦和沉重，而是要向你们显示因你们的冷漠我们所感到的忧伤。愿天主恩赐我们都从中解脱，使我们紧握属灵的热忱，获得天上的福乐，因我们的主\Person{耶稣基督}的恩宠和慈爱。愿光荣归于祂，及\Person{圣父}、\Person{圣神}，于无穷世。阿们。\par

\SourceRecord{Translated by Charles Marriott.；Nicene and Post-Nicene Fathers, First Series；Vol. 14.；Edited by Philip Schaff.；Buffalo, NY: Christian Literature Publishing Co.,；1889.；<http://www.newadvent.org/fathers/240130.htm>.}

% structure: p0001 source=h1 target=section reason=page-title

\section{若望福音讲道集第三十一篇}

% structure: p0002 source=h2 target=subsection reason=relative-heading-rank; base=h2

\subsection{\BibleRef{若 3:35-36}}

1. 谦卑的获益在一切事上显得多么巨大。正因如此，我们不是一下子从老师那里学全技艺，而是逐渐臻于完善；正因如此，我们缓慢地、一点一点地建造城市；正因如此，我们维持生命。如果这事在今生的事务中有如此大能，那么在属灵的事上，人可能会发现这种智慧的力量是多么巨大，这就不足为奇了。因为犹太人正是这样得以摆脱他们的偶像崇拜，被温和地、一点点地引导，起初既没有听到关于教义也没有听到关于生活的崇高之事。同样，在基督来临之后，当更崇高的教义时期到来之际，宗徒们起初没有宣讲任何崇高之事，就把所有人都归化了。基督在开始时对大多数人似乎也是这样说话的；现在若翰也是这样，他提及基督如同一个奇妙的人，并暗暗地引入高深的内容。\par

例如，当他开始讲论时，他这样说：\Quote{人不能领受任何东西除非是从天上赐给他的} \BibleRef{若 3:27}；然后加上一句崇高的声明，说：\Quote{从天上来的那一位在万有之上}，他又把话题降卑下来，并且除了许多其他事情之外，还说：\Quote{天主赐圣神没有限量}。接着他继续说：\Quote{父爱子，并把一切交给了祂手中}。之后，知道惩罚的力量是巨大的，且大多数人与其说是被好事的应许所引导，不如说是被可怕的威胁所推动，他就以这些话结束了他的讲论：\Quote{信从子的人，有永生；不信从子的人，不得见生命，天主的义怒却常在他身上}。这里他又把惩罚的叙述归于父，因为他说的是\Quote{子的义怒}（其实祂就是审判者），而是把父置于他们之上，希望以此更有效地恐吓他们。\par

\Quote{那么，只要信从子就够了，}有人说，\Quote{就能得永生吗？} 绝对不是。并且听基督亲自宣告这一点，说：\Quote{不是每一个对我说：主啊，主啊的人，都能进入天国} \BibleRef{玛 7:21}；而且亵渎圣神的罪本身就足以使人堕入地狱。但我何必只谈一部分教义呢？即使一个人对圣父、圣子、圣神的信仰正确，但如果他不度正直的生活，他的信德对他的救恩将毫无益处。因此，当祂说：\Quote{这是永生，就是他们认识祢是唯一的真天主} \BibleRef{若 17:3}时，我们不要以为这里所说的（认识）就足以使我们得救；我们还需要最严谨的生活和品行。因为祂在这里说了：\Quote{信从子的人，有永生}，并且在同一处说了更强硬的话（因为他编织他的讲论时不仅包括祝福，还包括它们的反面，他这样说：\Quote{不信从子的人，不得见生命，但天主的义怒常在他身上}）；然而，我们也不能仅凭这一点就断言信德就足以得救。福音许多地方给出的生活指示也证明了这一点。所以他没有说：\Quote{这本身是永生}，也没有说：\Quote{仅仅相信子的人就有永生}，而是通过这两种说法表明，信德确实包含生命，但如果正确的品行没有随之而来，就会有严厉的惩罚。他说：\Quote{常在他身上}，而不是\Quote{临到他身上}，意思是\Quote{永远不离开他}。免得你以为\Quote{不得见生命}是暂时的死亡，而要相信惩罚是持续的，他用这个说法表明惩罚一直停留在他身上。他用这些话，强迫他们归向基督。因此，他没有特别针对他们发出警告，而是使之普遍化，这是最能引导他们的方式。因为他没有说：\Quote{如果你们信}，和\Quote{如果你们不信}，而是使他的讲话具有一般性，使他的话免受怀疑。他这样做比基督做得更有力。因为基督说：\Quote{不信的人已经被定罪了}，但若翰说：\Quote{不得见生命，但天主的义怒常在他身上}。这是有充分理由的；因为一个人谈论自己和另一个人谈论他是不同的。他们会认为基督常常出于自爱谈论这些事，并且他是一个自夸者；但若翰完全没有嫌疑。如果后来基督也使用了更强硬的措辞，那是在他们开始对祂产生崇高看法的时候。\par

% structure: p0006 source=h2 target=subsection reason=relative-heading-rank; base=h2

\subsection{\BibleRef{若 4:1-3}}

耶稣本人并没有施洗，但那些传递消息的人，想要激起听者嫉妒，就这样报告了。那么，祂为什么\Quote{离开}呢？不是出于恐惧，而是要除去他们的恶意，并缓和他们的嫉妒。祂确实能够在他们来攻击祂时制止他们，但祂不愿继续这样做，免得降生奥迹被怀疑。因为如果祂常常被抓住又逃脱，那就会引起许多人怀疑；因此，祂多数时候是按照人的方式处理事情。正如祂希望人们相信祂是天主，同样也相信祂身为天主却取了肉身；因此即使在复活后，祂对门徒说：\Quote{你们摸摸我，看看，因为鬼神没有骨肉} \BibleRef{路 24:39}；同样，当伯多禄说：\Quote{主，这事绝不会临到祢身上} \BibleRef{玛 16:22}时，祂责备了伯多禄。由此可见，祂对这一事是多么关心。\par

2. 因为这是教会教义中不小的一部分；它是为我们成就的救恩的关键点；一切都借此得以成就并取得成效，因为正是这样，死亡的束缚被解开，罪恶被除去，咒诅被废除，万般恩惠被引入我们的生命。因此，祂特别希望这个降生奥迹被相信，因为它是对我们无数恩惠的根源和泉源。\par

然而，祂在这样对待祂的人性时，并没有让祂的天主性被遮蔽。因此，在祂离开后，祂再次使用了与先前相同的言语。因为祂前往加里肋亚并非单纯地离开，而是为了完成某些重要的事，即那些在撒玛黎雅人中间的事；祂也不是随意地处理这些事，而是以属于祂的智慧，并且为了不给犹太人留下任何无耻借口的余地。福音作者指出这一点时说道：\par

% structure: p0010 source=h2 target=subsection reason=relative-heading-rank; base=h2

\subsection{若 4:4}

这表明祂将此作为行程中的次要工作。宗徒们也是如此；因为正如他们受犹太人迫害时，便到了外邦人那里；同样，当犹太人驱逐祂时，基督便着手处理撒玛黎雅人，正如祂在叙利亚-腓尼基妇人的事上所做的一样。这样做是为了剥夺犹太人所有的辩护，使他们不能说：\Quote{祂离开了我们，到了未受割损的人那里。}因此，宗徒们为自己辩护说：\BibleQuote{天主的话本当先讲给你们听，但因你们断定自己不配得永生，看，我们就转向外邦人。}\BibleRef{宗 13:46}祂自己又说：\BibleQuote{我奉差遣，不过是到以色列家迷失的羊那里去。}\BibleRef{玛 15:24}又说：\BibleQuote{拿儿女的饼丢给狗吃是不合理的。}但当他们驱赶祂时，便为外邦人打开了门。然而，祂并不是明确地来到外邦人那里，而是顺道而来。那么，顺道而来时，\par

% structure: p0012 source=h2 target=subsection reason=relative-heading-rank; base=h2

\subsection{若 4:5-6}

为什么福音作者对地点如此精确？这是为了当你听见妇人说\Quote{我们的祖先雅各伯把这口井给了我们}时，你可以不觉得奇怪。因为这个地方正是肋未和西默盎因狄纳的事发怒，进行那次残忍屠杀的地方。并且，或许值得叙述一下撒玛黎雅人是由哪些人组成的，因为整个这一地区都称为撒玛黎雅。那么，他们从何得名呢？那座山因其主人而被称为\BibleQuote{索默尔}\BibleRef{列上 16:24}；正如依撒意亚所说：\BibleQuote{厄弗辣因的首都是索默隆}\BibleRef{依 7:9}；但居民并非被称为\Quote{撒玛黎雅人}，而是\Quote{以色列人}。但随着时间的推移，他们得罪了天主。在培卡黑在位期间，提革拉特-丕肋色尔上来，攻取了许多城，攻打厄拉，杀了他，将王国交给了曷舍亚。\BibleRef{列下 15:29}撒玛乃色来攻打他，夺取了其他城，使他们臣服并进贡。\BibleRef{列下 17:3}起初他顺服了，但后来背叛了亚述的统治，转而与雇士人结盟。亚述王得知此事，便向他们开战，摧毁了他们的城邑，不再容许那民族留在那里，因为他认为他们可能会叛变。\BibleRef{列下 17:4}但他将他们掳到巴比伦和玛代人那里，并从各处带来不同的民族，安置在撒玛黎雅，以便他的统治将来得以稳固，使自己的人民占据该地。此后，天主为了显示祂并非因软弱而放弃犹太人，而是因为那些被放弃的人的罪恶，便派遣狮子攻击那些外邦人，狮子蹂躏了他们整个民族。这些事报告给国王，他便派了一位司祭去向他们传授天主的法律。即便如此，他们仍未完全停止不敬，只是半心半意。但随着时间的推移，他们反倒抛弃了偶像，敬拜了天主。当事情发展到这种状态时，返回的犹太人从此对他们怀有嫉妒之情，视他们为陌生人和仇敌，并以山名称呼他们为\Quote{撒玛黎雅人}。因此，他们之间也产生了不小的敌对。撒玛黎雅人并不使用全部圣经，只接受梅瑟的著作，对先知们的作品则不甚重视。然而，他们却急于将自己挤入高贵的犹太血统，以亚巴郎自豪，称他为祖先，因为亚巴郎是加色丁人；他们也称雅各伯为父亲，因为他们是雅各伯的后裔。但犹太人憎恶他们，正如憎恶所有其他民族。因此，他们用这话来羞辱基督，说：\BibleQuote{你是撒玛黎雅人，并且附有魔鬼。}\BibleRef{若 8:48}正因为如此，在从耶路撒冷下耶里哥的人的比喻中，基督使那向他显示怜悯的人成为一个\Quote{撒玛黎雅人}\BibleRef{路 10:33}，一个在他们眼中被视为卑贱、可鄙、可憎的人。而在十个癞病人的事例中，祂因这个缘故称其中一人为\Quote{外方人}（因为他是一个\Quote{撒玛黎雅人}），并这样命令门徒说：\BibleQuote{不要走外邦人的路，也不要进入任何撒玛黎雅人的城。}\BibleRef{玛 10:5}\par

3. 福音作者提醒我们雅各伯，并不仅仅是为了描述这个地方，而是为了表明犹太人被弃绝早已发生。因为在他们祖先的时代，这些犹太人拥有这片土地，而不是撒玛黎雅人；而那些并不属于他们的祖先所获得的产业，当他们成为土地的拥有者时，却因自己的懈怠和过犯而失去了。由此可见，如果子孙不像他们的祖先，卓越的祖先也毫无益处。此外，那些外邦人仅仅经历狮子的试验之后，便立刻转而遵行犹太人的正确敬拜；而他们，在忍受了这样的打击之后，甚至仍未能恢复理智。\par

基督现在来到了这个地方，祂始终拒绝闲逸安逸的生活，而展现了辛劳活跃的生活。祂不骑任何牲畜，而是长途跋涉，甚至因旅途而疲惫。祂一直教导人应当亲自劳作，放弃多余之物，不有许多需求。不仅如此，祂如此渴望我们远离多余之物，以至于连许多必需之物也予以缩减。为此祂说：\Quote{狐狸有穴，天空的飞鸟有巢，人子却没有枕首的地方。}\BibleRef{玛 8:20}因此，祂大部分时间都在山上和旷野中度过，不仅白天如此，夜间也是如此。达味在预言时也说：\Quote{祂要在路旁喝溪水}\BibleRef{咏 109:7}，以此显示祂俭朴的生活方式。圣史在此处也指出了这一点。\par

% structure: p0016 source=h2 target=subsection reason=relative-heading-rank; base=h2

\subsection{若望福音 4:6-8}

由此我们得知祂旅途的勤勉，对食物的不关心，以及祂如何将其视为次要之事。门徒也受教导要自己培养同样的心态，因为他们没有携带路上所需的食物。另一位圣史也指出了这一点，说当祂向他们谈及\Quote{法利塞人的酵母}\BibleRef{玛 16:6}时，他们以为这是因为他们没有带饼；当他叙述他们掐麦穗吃\BibleRef{玛 12:1}，以及他说耶稣因饥饿来到无花果树前\BibleRef{玛 21:18}时，这一切无非是要教导我们轻视肚腹，不要以为服侍肚腹是必须要操心的。例如，在此处他们既没有带任何东西，也没有因为没带任何东西而在一天之初就为此操心，而是在其他人用餐时才买食物。不像我们，一从床上起来就先处理这事，请来厨师和管家，认真吩咐，然后再处理其他事务，将属世之事置于属灵之事之前，把本应视为次要的当作必要的。因此一切都混乱了。相反，我们应当重视所有属灵之事，完成这些之后，再处理其他事。\par

在此处，不仅显出祂的勤劳，也显出祂的无骄傲；不仅因祂疲倦，或因祂坐在路旁，更因祂独自一人，门徒与祂分离。然而，祂本有能力，若祂愿意，可以不打发他们全部离开，或在他们离开后有其他仆人伺候。但祂不肯，因为祂如此训练祂的门徒，将一切骄傲践踏在脚下。\par

\Quote{这有什么可惊奇的呢？}一人说，\Quote{如果他们欲望有节制，既然他们是渔夫和制帐棚的人？}是的！他们本是渔夫和制帐棚的人，但转瞬间已升到天堂的高处，并变得比世上所有君王更尊贵，被视为配做世界之主的同伴，并追随那位万众敬畏的主。你们也知道，特别是那些出身低微的人，一旦获得身份，就更容易陷入愚妄，因为他们完全不知如何承受突如其来的尊荣。因此，祂在他们现有的谦卑心态中克制他们，教导他们始终谦逊，从不要求有人服侍他们。\par

% structure: p0020 source=h3 target=subsubsection reason=relative-heading-rank; base=h2

\subsubsection{祂因此旅途疲惫，就这样坐在井旁}

你们看见祂坐下是因为疲倦吗？因为炎热吗？因为等候门徒吗？祂确实知道撒玛黎雅人那里将要发生的事，但祂并非为此主要而来；然而，虽非为此而来，祂也不应拒绝那位来到祂面前、表现出如此求知渴望的妇人。当祂来到犹太人那里时，他们却驱赶祂；当祂往别处去时，外邦人却将祂吸引到自己身边。那些人心生嫉妒，这些人却相信了祂。那些人生气，这些人却尊敬和朝拜祂。那么该怎样呢？难道祂要忽视这么多人的得救，打发走如此高贵的热情？这就有负于祂的仁爱。因此，祂以相称于祂的智慧安排了眼前的一切。祂坐下休息身体，并在泉边使其凉爽；因为那时正是正午，正如圣史所说：\Quote{约在第六时辰。}\par

祂\Quote{就这样}坐着。\Quote{这样}是什么意思？不是坐在宝座上，也不是坐在坐垫上，而是简单地、照祂本来的样子坐在地上。\par

% structure: p0023 source=h3 target=subsubsection reason=relative-heading-rank; base=h2

\subsubsection{有一个撒玛黎雅妇人来打水}

4. 观察他如何说明那妇人出来是为另一目的，从而以各种方式堵住了犹太人无耻的反对，免得有人说祂行事违背了自己的命令，即吩咐（祂的门徒）不要进入任何撒玛黎雅人的城，却与撒玛黎雅人交谈。\BibleRef{玛 10:5}因此圣史写下了，\par

% structure: p0025 source=h3 target=subsubsection reason=relative-heading-rank; base=h2

\subsubsection{因为门徒已经往城里买食物去了}

那妇人做了什么？当她听见\Quote{请给我水喝}，就非常聪明地把基督的话作为提问的契机，说：\par

% structure: p0027 source=h2 target=subsection reason=relative-heading-rank; base=h2

\subsection{若望福音 4:9}

她又从哪里以为祂是犹太人呢？或许是从祂的衣着，从祂的口音。请看，我请你们注意，这妇人是何等体贴。如果需要谨慎，那需要谨慎的是耶稣，而不是她。因为她没有说\Quote{撒玛黎雅人不与犹太人往来}，而是说\Quote{犹太人不接纳撒玛黎雅人}。然而，虽然她自己无可指责，但当她以为别人正陷入过错时，她仍不肯就此缄默，反而纠正了她认为违法的事。或许有人会问：耶稣为何向她讨水喝，而法律却不允许呢？如果回答说，是因为祂预先知道她不会给，那么正因如此，祂就不该讨。那我们能说什么呢？可以说，祂如今已不在乎拒绝这些规条；因为那促使他人废除这些规条者，更会自己越过它们。祂说：\BibleQuote{不是入口的污秽人，而是出口的。}\BibleRef{玛 15:11} 这番与妇人的对话，对犹太人将是一个不小的指控。因为祂常藉言语和行动吸引他们归向自己，他们却不肯听从；看她如何因一个简单的请求而停留。因为祂尚未着手办理这事，也未指出途径，但若有人到祂这里来，祂并不阻止。祂也曾对门徒如此说：\Quote{撒玛黎雅人的任何城市，你们不要进去。}祂没有说：\Quote{但当他们来到你们跟前时，你们要拒绝他们}；那将极不符合祂的仁爱。因此祂回答妇人说：\par

% structure: p0029 source=h2 target=subsection reason=relative-heading-rank; base=h2

\subsection{若 4:10}

首先，祂表明她堪当听闻，不应被忽视，然后祂启示了自己。因为她一旦知道祂是谁，就会立刻听从并留心祂；这是无人能对犹太人说的，因为他们虽然知道了，却不向祂求什么，也不愿受教于任何有益的事，反而侮辱祂、赶走祂。但妇人听了这些话，请看她是多么温和地回答：\par

% structure: p0031 source=h2 target=subsection reason=relative-heading-rank; base=h2

\subsection{若 4:11}

祂已将她从对祂的卑微看法中提升，从以为祂只是个凡人的想法中提升。因为她在这里称祂为\Quote{主}，并非无因，而是赋予祂崇高的尊荣。她说这些话是为尊敬祂，从之后的话可以明显看出，因为她没有嘲笑、没有讥讽，只是迟疑了一会儿。不要奇怪她没有立刻明白一切，因为尼苛德摩也是如此。他说：\Quote{这事怎能成就？}又说：\Quote{人已年老，怎能重生？}又说：\Quote{难道他能再入母腹而生吗？}但这个女人更加恭敬：\Quote{先生，你没有打水的器具，井又深，那活水从哪里来呢？}基督说了一件事，她却想象另一回事，只听见字面意思，尚未能形成任何高深的思想。然而，若她轻率发言，她可能会说：\Quote{你若真有那活水，就不会向我讨，而会为自己预备。你不过是夸口。}但她没有说这样的话；她十分温和地回答，无论是起初还是后来。起初她说：\Quote{你既是犹太人，怎么向我一个撒玛黎雅妇人要水喝？}她没有像对异族和敌人说话那样说：\Quote{我断不可给你，因为你是我们民族的仇敌和外人。}之后，当她听见祂说出伟大的话时——这是仇敌最恼怒的事——她并没有嘲笑或讥讽；但她说的是什么呢？\par

% structure: p0033 source=h2 target=subsection reason=relative-heading-rank; base=h2

\subsection{若 4:12}

请注意她如何将自己纳入犹太人的高贵世系。因为她的话大意如此：\Quote{雅各伯曾用这井，他没有更好的给我们。}她说这话，表明从（基督的）第一个回答开始，她已有了一个伟大崇高的想法；因为\Quote{他本人、他的儿女和他的牲畜都喝过这井的水}这句话，无非是暗示她意识到有一种更好的水，但从未找到，也未清楚认识。更清楚地解释她想要说的是什么，她话中的意思是：\Quote{你不能断言雅各伯给了我们这井，而他自己却用另一口井；因为他和他的儿女都喝过这口井的水，倘若他们有另一口更好的，就不会这样做了。如今，这井的水你不能给我，你也不可能有另一口更好的水，除非你承认你比雅各伯更大。那么，你所应许要赐给我们的水，是从哪里来的呢？}犹太人与祂交谈并不如此温和，尽管祂对他们讲了同样的话题，提到类似的水，但他们毫无所得；而当祂提到亚巴郎时，他们甚至企图用石头砸祂。这妇人却不是这样对待祂；她极其温和，在正午的炎热中，耐心地听和说一切，甚至没有想过犹太人很可能断言的话：\Quote{这人是疯了，神志不清：他把我拴在这水泉和井边，什么也不给我，只说大话}；不，她忍耐坚持，直到找到了她所寻求的。\par

5. 如果一位撒玛黎雅妇女如此热切地学习有益之事，虽然还不认识基督，却仍留在祂身边，那么，我们这些既认识祂，又不在井旁，不在旷野，不在正午，不在烈日之下，而是在清晨，在这样的屋顶下，享受荫凉和舒适，却无法忍受听任何道理，反而感到厌倦的人，将获得什么宽恕呢？那妇女并非如此；她被耶稣的话深深吸引，以致还叫别人来听。相反，犹太人非但不叫别人，反而阻碍和拦阻那些愿意到祂面前来的人，说：\BibleQuote{你看，有哪位首领相信了祂？但这些不认识法律的平民，是被诅咒的。}那么，让我们效法这位撒玛黎雅妇女，让我们与基督交谈。因为如今祂仍站在我们中间，藉着先知和门徒向我们说话；让我们聆听并服从。我们要徒然无益地生活多久呢？因为不做天主所喜悦的事，就是徒然生活，甚至不仅是徒然，而且是有害的；因为当我们把所赐的时间浪费在毫无益处的目的上时，我们将离开此生，因不合时宜的放纵而遭受最严厉的惩罚。因为一个人领了钱去做买卖，却把它吃光了，那托付给他的人岂能不找他算账？同样，一个人像我们这样虚度光阴，也绝不能逃脱惩罚。天主领我们进入今生，并赋予我们灵魂，并非为此，要我们只利用现在的时间，而是要以未来的生命为念，处理一切事务。只有无理性的动物才仅对今生有用；但我们拥有一个不死的灵魂，为使我们能用一切方法预备自己，进入那另一个生命。因为如果有人问马、驴、牛及其他这类动物的用处，我们会告诉他，无非是服务今生；但对我们却不能这样说；我们最好的状态是在离开之后；我们必须尽一切努力，好在那里发光，加入天使的歌唱团，永远侍立在君王面前，直至永无穷尽。因此，灵魂是不死的，身体也将不死，好使我们能享受永无穷尽的福乐。但是，当天上之物摆在你面前时，你仍被钉在地上，想想这是对施恩者多大的侮辱——祂向你展示天上的事物，而你却不重视它们，反而选择了地上的事物。因此，因被你轻视，祂用地狱来威胁你，为要你由此得知你使自己丧失了多么大的福乐。愿天主使无人经历那惩罚，而是因我们中悦基督，藉着我们的主耶稣基督的恩宠和慈爱，获得永远的福乐；愿光荣归于祂，及圣父和圣神，自今而始，以至永远，永世无尽。阿们。\par

\SourceRecord{Translated by Charles Marriott.；Nicene and Post-Nicene Fathers, First Series；Vol. 14.；Edited by Philip Schaff.；Buffalo, NY: Christian Literature Publishing Co.,；1889.；<http://www.newadvent.org/fathers/240131.htm>.}

% structure: p0001 source=h1 target=section reason=page-title

\section{《若望福音》第卅二篇讲道}

% structure: p0002 source=h2 target=subsection reason=relative-heading-rank; base=h2

\subsection{若 4:13-14}

1. 圣经有时将圣神的恩宠称为\Quote{火}，有时称为\Quote{水}，表明这些名称并非描述其本质，而是描述其运作；因为圣神是无形而单纯的，不可能由不同的实质组成。若翰宣告这一点，如此说：\BibleQuote{祂要用圣神与火洗你们}\BibleRef{玛 3:11}；另一方面，基督说：\BibleQuote{从他的腹中要流出活水的江河}\BibleRef{若 7:38}。\BibleQuote{但这话}，若望说，\BibleQuote{是指着祂要赐给的圣神说的。}同样，在与那妇人交谈时，祂称圣神为水；因为：\BibleQuote{凡喝了我所赐的水，永远不再渴。}同样，祂也用\Quote{火}的名称来称呼圣神，暗示恩宠的激励和温暖特性，及其摧毁过犯的能力；而用\Quote{水}来宣告它所施行的净化，以及它给那接受它的心灵带来的巨大舒爽。这是有充分理由的；因为它使甘愿的灵魂像一座花园，长满各种常青、结果累累的树木，不容它感到沮丧或遭受撒旦的阴谋，并扑灭恶者的一切火箭。\par

请观察基督的智慧，祂是如何温柔地引导那妇人；因为祂并没有一开始就说：\BibleQuote{你若知道是谁对你说：给我水喝}，而是在给了她称祂为\Quote{犹太人}的机会，并使她陷于这样称呼的罪责之后，驳斥这指控时祂才说：\BibleQuote{你若知道是谁对你说：给我水喝，你必早求祂}；通过祂伟大的应许迫使她提及圣祖，这样祂才让妇人看透，然后当她反驳说：\BibleQuote{你比我们的祖先雅各伯还大吗？}祂并没有说：\Quote{是的，我更大}（因为那似乎只是自夸，既然证据尚未显明），而是藉着祂所说的话达到了这一点。因为祂没有简单地说：\Quote{我要给你水}，而是先撇下雅各伯所给的水，高举祂自己所给的水，意在从所赐之物的性质显示两位施予者之间有多么大的差距和区别，以及祂自己胜过圣祖。祂说：\Quote{如果，你钦佩雅各伯，因他给了你这水，那么，若我给你远胜于此的水，你会说什么呢？你自己已首先承认我比雅各伯大，因为你反驳我说：\InnerQuote{你比雅各伯还大吗，竟应许给我更好的水？}你若接受了那水，必定会承认我更大。}你看到那妇人正直的判断吗？她根据事实，既对圣祖，也对基督，作出了决定？犹太人并没有这样行；他们即使见祂赶鬼，不仅不称祂比圣祖大，甚至说祂附有魔鬼。但这妇人却不如此，她按照基督所期望的，从祂的行为所显明的证据形成她的看法。因为祂以此为自己辩护，说：\BibleQuote{我若不行我父的工程，你们就不必信我；但若行了，你们纵然不信我，也当信这些工程。}\EditorialNote{若 10:37-38}这样妇人就被引向信仰了。\par

因此，当祂听见\Quote{你比我们的祖先雅各伯还大吗}之后，祂撇下雅各伯，论及那水说：\BibleQuote{凡喝这水的，还要再渴；但喝我所赐的水，永远不渴。}祂作比较，并非贬低前者，而是显示后者的卓越；因为祂没有说\Quote{这水是无用的}，也没有说\Quote{它是低劣可鄙的}，而是说出连本性也见证的话：\BibleQuote{凡喝这水的，还要再渴；但喝我所赐的水，永远不渴。}那妇人先前已听见\Quote{活水}\BibleRef{若 4:10}，却不明白其含义。因为那水被称为\Quote{活水}，是指从不断的泉源中涌出、永不枯竭的水，她以为就是指这样的水。因此，祂更清楚地指明这一点，如此说，并以比较确立了其优越性。那么祂说什么呢？\BibleQuote{喝我所赐的水，永远不渴。}这一点以及接下来所说的尤其显示了优越性，因为物质的水不具备这些特性。接下来是什么呢？\BibleQuote{要在他内成为涌到永生的水泉。}因为正如一个人腹内有泉，绝不会被渴所困，同样，拥有这水的人也不会。\par

那妇人立刻相信了，显得比尼苛德摩聪明得多，而且不仅更聪明，也更刚毅。因为尼苛德摩听见了成千上万这样的教导，既没有邀请别人来听，也没有自己公开说出来；但这妇人却表现出了宗徒的行为，向众人传福音，召唤他们到耶稣那里，并吸引整座城的人到祂面前。尼苛德摩听见后说：\BibleQuote{这事怎么能成就呢？}而当基督向他提出一个清楚的比喻，即\Quote{风}的比喻，他连这样也没有接受圣言。但这妇人却不然；起初她曾疑惑，但后来，不是通过任何系统的论证，而是以陈述的形式接受了圣言，便立刻急忙拥抱了它。因为当基督说：\BibleQuote{要在他内成为涌到永生的水泉}时，妇人立刻说：\par

% structure: p0007 source=h2 target=subsection reason=relative-heading-rank; base=h2

\subsection{若 4:15}

你看她是如何逐渐被引领至高深的道理？起初，她以为祂是某个违反法律的犹太人；当祂驳斥了那指控后（因为教导她如此之事的人不该被怀疑），她听了\Quote{活水}，就以为这话是指物质的水；后来得知这话是属神的，她便相信那水能解除口渴的必需，但仍不知道这水是什么；她仍疑惑，以为它确实超越物质，却不确切明白。但在此，她有了更清晰的认识，却尚未完全领悟全部（因为她说：\Quote{请给我这水，使我不渴，也不用来这里打水}），这时她把祂置于雅各之上。她说：\Quote{因为}她说：\Quote{我若从你得了那水，就不需要这井了。}你看她如何将祂置于族长之上？这是一颗公正判断灵魂的行为。她曾显示她对雅各有何等大的敬意，现在她看见一位比他更好者，并未被先入之见所阻。因此，这妇人既非轻易动怒（她并未轻率接受所说的话，她如此精确地追问，岂会如此？），也非不顺服或好争论，她的请求显示出这一点。然而，祂曾对犹太人说：\Quote{凡吃我肉的人永远不饿，信我的人永远不渴}\BibleRef{若 6:35}；但他们不仅不信，反而因祂而跌倒。这妇人没有这种感觉，她留下并恳求。祂对犹太人说：\Quote{信我的人永远不渴}；对妇人却没有这样说，而是更具体地说：\Quote{喝这水的人永远不渴。}因为那应许指的是属神且不可见的事。因此，祂用应许提升她的心智后，仍停留在与感官相关的表达上，因为她尚不能领悟属神事物的确切表达。因为若祂说：\Quote{你若信我，就不渴}，她不会理解祂的话，不知道对她说话的是谁，也不明白祂所说的是哪种渴。那么，为什么祂在犹太人面前不这样做呢？因为他们见过许多神迹，而她没见过任何神迹，只是先听了这些话。因此祂后来借着预言启示祂的能力，并没有直接引入责备，而是说了什么？\par

% structure: p0009 source=h2 target=subsection reason=relative-heading-rank; base=h2

\subsection{若望福音 4:16-19}

2. 啊，这妇人的智慧何其大！她多么温顺地接受责备！\Quote{她岂能不如此？}有人说。请告诉我，她为何该如此？祂不是也常常责备犹太人，且用比这更严厉的责备吗？（因为揭露心中的隐秘思想与显明私下行的事不同：前者只有天主知道，除了心中有那思想的人，无人知道；后者，所有参与其事的人都知道；）但受责备时，他们并没有耐心忍受。当祂说：\Quote{你们为什么想要杀我？}\BibleRef{若 7:19}，他们不仅没有像这妇人那样钦佩，反而讥笑侮辱祂；然而他们从其他神迹得到了证明，她只听到了这番话。但他们不仅不钦佩，反而侮辱祂，说：\Quote{你有魔鬼，谁想要杀你？}而她不仅不侮辱，反而钦佩、惊讶于祂，并以为祂是先知。的确，这责备对妇人的触动比对他们的更大；因为她的过犯是她自己的，他们的则是普遍的；我们通常被个别的事刺痛，而非普遍的事。此外，他们以为若能杀害基督，便得到极大的目标，而妇人所做的事众人都认为是邪恶的；然而她并不愤慨，反而惊讶赞叹。基督在纳塔乃耳的事上也做了同样的事。祂没有一开始就引入预言，说：\Quote{我看见你在无花果树下}，而是当纳塔乃耳说：\Quote{你从哪里认识我？}那时祂才引入这话。因为祂渴望从那些亲近祂的人身上取得祂神迹和预言的起始，好让他们因所行的事更受吸引，并避免虚荣的嫌疑。祂在此也这样做；因为若一开始就指责她\Quote{你没有丈夫}，会显得沉重而多余；但从她自身引出说话的理由，然后纠正所有这些点，则十分合理，并柔和了听者的性情。\par

\Quote{有人说：\InnerQuote{去，叫你的丈夫}这话有什么关联呢？}谈话涉及超越必死本性的恩赐和恩宠：妇人急切地寻求接受它。基督说：\Quote{叫你的丈夫}，表明他也必须分享这些事；但她急于接受（恩赐），并掩盖当时情况的羞耻，以为自己在与一个男人交谈，就说：\Quote{我没有丈夫。}基督听了这话，适时引入了责备，准确地提到了两点：祂数说了她所有的前夫，并责备她如今想隐藏的那一位。那么，这妇人做了什么？她没有恼怒，也没有离开祂逃跑，也不认为这是侮辱，反而更钦佩祂，并坚持得更久。她说：\Quote{我看您是一位先知。}注意她的谨慎；她没有立刻跑向祂，仍思考着祂，并惊叹于祂。因为\Quote{我看}意思是\Quote{在我看来，您是先知。}然后当她怀疑这事时，她不再问关于今生的事，不是关于身体的健康、产业或财富，而是立刻问关于教理。因为她说的是什么？\par

% structure: p0012 source=h2 target=subsection reason=relative-heading-rank; base=h2

\subsection{若望福音 4:20}

3. 你看见她如何在思想上变得更高超了吗？她曾担心因口渴而烦恼，如今却追问教义。那么基督做什么？祂并不解答这问题，（因为简单地回应人的话语并非祂所关心，因为那是无必要的，）却引领这妇人达到更高的境界，且与她谈论这些事情，直到她承认祂是一位先知，好使她此后能以充足的信念聆听祂的圣言；因为既已被说服于此，她再也不能对所要对她说的话产生怀疑。\par

此后，让我们感到羞愧和脸红。一个曾有过五位丈夫的撒玛黎雅妇人，竟对教义如此热切，以致一天中的时辰、她来此的另有目的，或任何其他事，都不能使她远离询问这些事；而我们不仅不询问教义，反而对一切教义的态度都是漫不经心和冷漠的。因此一切都荒废了。因为你们谁在家中曾拿起一本基督信仰书籍，翻阅其内容，并研究圣经？没有一个人能说他这样做，但在我们大多数人中，我们会看到的是赌博和骰子，而书籍却无处可寻，除了少数人之外。即便这少数人，他们的心态也与多数人相同；因为他们将书籍捆扎起来，总是收藏在盒子里，他们所有的关心都在于羊皮纸的精致和字母的美丽，而不是为了阅读。因为他们购买书籍并非为了从中获得利益和益处，而是费力于这些事以显示他们的财富和骄傲。这便是虚荣的过分。我没有听见任何人夸耀他知晓内容，而是夸耀他拥有一本用金字写成的书。告诉我，这有什么益处？圣经赐给我们不仅仅是为了这个，即我们能将它们保存在书中，而是为了我们能将它们铭刻在我们的心上。因为仅仅以文字遵守诫命这种拥有方式，属于犹太人的虚荣；但对我们来说，律法根本不是这样赐下的，而是在我们心的肉版上。我说这话，并非阻止你们购置圣经，相反，我规劝并恳切地祈求你们这样做，但我希望你们从那些书中将文字和意义传达到你们的理智中，这样当理智接受了书写的意义时，它就能被净化。因为如果魔鬼不敢接近一个放有福音书的房屋，那么更不会有任何恶神或任何有罪的本性会触及或进入一个怀有如同福音书所包含之情感的灵魂。因此，要圣化你们的灵魂，圣化你们的身体，常将这些话存于你们的心中及口中。因为如果污秽的言语玷污并招引魔鬼，那么显然，属灵的阅读圣化并吸引圣神的恩宠。圣经是神圣的符咒，让我们从它们中获取疗治我们自己和我们的灵魂情欲的良药。因为如果我们明白所读的是什么，我们就会以极大的乐意去听它。我常这样说，并且不会停止这样说。那些坐在市场旁边的人能说出车夫、舞者的名字、家族和城邦，以及每个人所拥有的能力种类，并精确说出那些马匹的好坏特性；而来到这里的人却对自己所做的事一无所知，甚至连圣书的总数都不知晓，这难道不奇怪吗？如果你们追求那些世俗之事是为了享乐，我要向你们证明，这里也有更大的快乐。告诉我，哪个更甘甜，哪个更奇妙？是看见一个人与另一个人角力，还是看见一个人与魔鬼较量，一个身体与一个无形之力搏斗，并且你们同族的人得胜？让我们注视这些角力，这些也是可敬且有裨益去效仿的，效仿它们，我们就能得到冠冕；而不是那些因竞争而使效仿者蒙羞的角力。如果你们观看那一种竞赛，你们是与魔鬼一同观看；另一种则是与天使、总领天使以及总领天使的主一同观看。现在告诉我，如果你们被允许与总督和君王同坐，观看并享受那景象，你们岂不认为那是极大的光荣？但在这里，当你们与天使之王一同作观众，当你们看见魔鬼被从背后抓住，竭力想要得胜却无能无力时，你们岂不跑去追求这样的景象吗？\Quote{这怎么可能呢？}有人说。如果你们手中拿着圣经，因为你们将在其中看到竞技场、漫长的赛跑、他的抓握和义人的技巧。因为观看这些事，你们也将学会自己如何角力，并避开魔鬼；外邦人的表演是魔鬼的聚会，而非人的剧院。为此我劝你们远离这些撒殚的集会；因为若进入偶像的殿宇是不合法的，那么进入撒殚的节庆更是如此。我不会停止说这些话并使你们厌烦，直到我看见一些改变；因为说这些话，如保禄所说，\Quote{于我并不为难，于你们却是安全的。}\BibleRef{斐 3:1} 因此不要因我的劝告而反感。如果有人应当反感，那是我常常说话却未被人听，而不是你们常常听却总是不服从。愿天主使你们不常受此谴责，而是从这羞耻中获得释放，堪当享受属灵景象及将来的光荣，藉着我们的主耶稣基督的恩宠与慈爱，愿光荣归于祂及圣父和圣神，直到永远。阿们。\par

\SourceRecord{Translated by Charles Marriott.；Nicene and Post-Nicene Fathers, First Series；Vol. 14.；Edited by Philip Schaff.；Buffalo, NY: Christian Literature Publishing Co.,；1889.；<http://www.newadvent.org/fathers/240132.htm>.}

% structure: p0001 source=h1 target=section reason=page-title

\section{若望福音讲道集 第三十三篇}

% structure: p0002 source=h2 target=subsection reason=relative-heading-rank; base=h2

\subsection{\BibleRef{若 4:21-22}}

亲爱的诸位，我们处处需要信德，信德是福佑之母，救恩之药；没有它，便不可能掌握任何一项伟大的道理。没有信德，人就像那些试图不用船横渡大海的人，他们靠游泳坚持一小段路程，手脚并用，但走得更远时，便迅速被波浪淹没：同样，那些在尚未学到任何东西之前就运用自己推理的人，也遭遇了船难；正如保禄所说：\Quote{有人关于信德，便遭遇了船难。} \BibleRef{弟前 1:19} 为避免我们也如此，让我们紧握那神圣的锚，基督正是借此引领那位撒玛黎雅妇人。因为她曾说过：\Quote{你们怎么说耶路撒冷是应当朝拜的地方呢？} 基督回答说：\Quote{女人，你相信我罢，时候将到，你们朝拜圣父，既不在这座山上，也不在耶路撒冷。} 祂向她启示了一个极其重要的道理，是祂未曾向尼苛德摩或纳塔乃耳提及的。她急于证明自己的特权比犹太人的更尊贵，并巧妙地根据祖传（的习俗）来辩论，但基督并未正面回答这个问题。因为当时谈论此事并说明为何祖先在这座山上朝拜，而犹太人在耶路撒冷朝拜，会分散注意力。因此祂对此保持沉默，剥夺了两地尊贵上的优先性，通过指出与将要赐予的相比，无论犹太人或撒玛黎雅人都没有拥有什么伟大的事物，来唤醒她的心灵；然后祂引入了区别。然而即便如此，祂仍宣告犹太人更为尊贵，并非偏爱一个地方胜过另一个，而是因其意向而给予他们优先。如同祂说：\Quote{关于朝拜的\InnerQuote{地方}，你们今后无需争论，但在\InnerQuote{方式}上，犹太人比你们撒玛黎雅人更有优势，因为\InnerQuote{你们}，祂说，\InnerQuote{朝拜你们所不认识的；我们却认识我们所朝拜的。}}\par

那么，撒玛黎雅人如何\Quote{不认识}他们所朝拜的呢？因为他们以为天主是地方性的、有偏私的；他们至少就是这样事奉祂的，因此他们派人去见\Person{波斯人}，报告说：\Quote{这地方的天主向我们发怒。}\par

\SourceRecord{Translated by Charles Marriott.；Nicene and Post-Nicene Fathers, First Series；Vol. 14.；Edited by Philip Schaff.；Buffalo, NY: Christian Literature Publishing Co.,；1889.；<http://www.newadvent.org/fathers/240133.htm>.}

% structure: p0001 source=h1 target=section reason=page-title

\section{\BibleBook{若望}{福音}讲道集 第三十四篇}

% structure: p0002 source=h2 target=subsection reason=relative-heading-rank; base=h2

\subsection{\BibleBook{若望}{福音} 4:28-29}

1. 我们需要极大的热情和被激起的奋勉，因为缺乏这些，就不可能获得应许给我们的福分。为表明这一点，基督一时说：\BibleQuote{除非人背起自己的十字架跟随我，不配是我的} \BibleRef{玛 10:38}；另一时说：\BibleQuote{我来是把火投到地上，我何等渴望它已经燃烧起来？} \BibleRef{路 12:49}；祂藉这两者都期望向我们描绘一个充满热力和火焰、预备好面对一切危险的门徒。这样的一个人就是这妇人。因为她被祂的话如此点燃，以至于她撇下了水罐和她来的目的，跑进城里，把所有人都带到耶稣跟前。\BibleQuote{来，}她说，\BibleQuote{看一个人，他把我所作的一切都告诉了我。}\par

请注意她的热忱和智慧。她来打水，当她遇见了真实的井之后，便轻视了物质的井；甚至藉这微小的事例教导我们，当我们聆听属灵之事时，要忽视现世的事物，不予以计较。因为宗徒们所做的，这妇人也按她的能力做了。他们蒙召时，撇下了网；她则自动自发，未经任何人命令，撇下了水罐，并因喜乐而展翅，履行了福音传道者的职分。她不像安德肋和斐理伯那样只叫一两个人，而是唤醒了整座城和百姓，将他们带到祂面前。\par

也请注意她说话何等谨慎；她没有说：\BibleQuote{来，看基督，}而是以同样基督网住她的屈尊俯就来吸引众人到祂那里；她说：\BibleQuote{来，看一个人，他把我所作的一切都告诉了我。}她不以说\BibleQuote{他把我所作的一切都告诉了我}为耻。然而她本可这样说：\Quote{来，看一个说预言的人}；但当灵魂被圣火点燃时，她便不看任何地上的事物，无论荣耀或羞辱，而只专注于那占据她的火焰。\par

\BibleQuote{莫非这就是基督？}请再次注意这妇人的大智慧；她既没有明确宣告事实，也没有保持沉默，因为她不愿凭自己的断言带他们来，而是想使他们透过听祂说话而分享这个意见；这使她的话更容易被他们接受。然而祂并没有把她的一生都告诉她，只是从她所说的话，她就确信祂知道其余的事。她也没有说：\Quote{来，相信，}而是说：\Quote{来，看}；这比另一种说法更温和，也更能吸引他们。你看到这妇人的智慧了吗？她知道，她确知，只要尝了那井，他们就会像她一样受影响。然而任何较为粗俗的人都会隐瞒耶稣给她的责备；但她却炫耀自己的生命，将它摆在众人面前，以吸引和俘获所有人。\par

% structure: p0007 source=h2 target=subsection reason=relative-heading-rank; base=h2

\subsection{\BibleBook{若望}{福音} 4:31}

\Quote{问，}在此是\Quote{恳求，}在他们本地语言中；因为他们看见祂旅途疲乏，又酷热难当，就恳求祂；因为他们关于食物的请求不是出于急躁，而是出于对老师的热爱？那么基督说了什么？\par

% structure: p0009 source=h2 target=subsection reason=relative-heading-rank; base=h2

\subsection{\BibleBook{若望}{福音} 4:32-33}

那么，你们为何惊奇，那妇人听见\Quote{水}时，仍以为指的是普通的水，而连门徒也处于同样的情况，至今不想到任何属灵的事，反而困惑呢？虽然他们仍表现出对老师惯有的谦逊和敬畏，彼此交谈，却不敢向祂提问。他们在其他场合也这样做，想要问祂，却没有问。那么基督说了什么？\par

% structure: p0011 source=h2 target=subsection reason=relative-heading-rank; base=h2

\subsection{\BibleBook{若望}{福音} 4:34}

祂在此称人的救恩为\Quote{食物}，表明祂多么热切地为我们预备；因为正如我们渴望食物，祂也渴望我们得救。请听，祂在各地如何不立刻揭示一切，而是先使听者困惑，以便他们开始寻求所说之话的意义，然后因困惑和困难，当所求出现时，就更愿意接受，并更留心倾听。因为祂为何不立刻说：\BibleQuote{我的食物就是承行我父的旨意？}（虽然这也不清楚，但比另一个清楚。）但祂说什么？\BibleQuote{我有食物吃，是你们不知道的}；因为祂，如我所说，首先想通过他们的不确定使他们更留心，并藉着这样的隐晦说法使他们习惯于听从祂的话。但什么是\Quote{父的旨意}？祂接下来谈到这一点并解释了。\par

% structure: p0013 source=h2 target=subsection reason=relative-heading-rank; base=h2

\subsection{\BibleBook{若望}{福音} 4:35}

看，祂又藉熟悉的话语引导他们思考更大的事；因为当祂说到\Quote{食物}时，祂所指的无非是那些将要来到祂面前的人的救恩；同样，\Quote{田地}和\Quote{收获}也指同一件事，就是预备好接受宣讲的众多灵魂；而祂所说的\Quote{眼睛}，既指肉身的眼睛，也指心灵的眼睛（因为他们此刻正看见撒玛黎雅人的群众前来）；祂称他们意志的准备为\BibleQuote{已经发白的田地}。因为正如谷穗变白、预备收割时一样，这些人，祂说，已预备好、适合得救。\par

为何祂不直接称它们为\Quote{田地}和\Quote{收割}，却明白地说：\Quote{那些人前来相信，并准备好接受圣言，因为他们已受过先知的训导；如今正结出果实}？祂使用这些比喻有何用意？因为祂不仅在此处如此，在整个福音中皆是如此；先知们也采用同样的方法，用隐喻的方式讲述许多事。那么，其原因何在？因为圣神的恩宠并非无故这样安排，而是有原因和理由的。有两个原因：其一，为使讲论更生动，更清晰地将所言之物呈现在我们眼前。因为心灵一旦把握住关于所论事物的熟悉形象，就更加活跃，仿佛在画中看到它们，从而更深刻地沉浸其中。这是原因之一；另一个原因是，为使陈述更悦耳，并使对所言之事的记忆更持久。因为断言对于普通听众的影响，不如借事物叙述和经验的再现那样能征服并引导他们。这一点，我们在此可从比喻中看到最明智的体现。\par

% structure: p0016 source=h2 target=subsection reason=relative-heading-rank; base=h2

\subsection{若望福音 4:36}

因为地上收割的果实无益于永生，而只属于暂时的生命；但属神的果实却归于那既无年龄亦无死亡者。你看见了吗？表述是属感觉的，思想却是属神的，并且祂亲自用这些话将属地的与属天的分别开来。因为当祂论及水时，祂将\Quote{喝这水的，永远不渴}作为天上之水的独特属性；此处祂也是如此，当祂说：\Quote{这果实被收取归入永生}。\par

\BibleQuote{好使那播种的和那收割的一同欢乐。}\par

谁是\Quote{那播种的}？谁是\Quote{那收割的}？先知们是播种者，但他们没有收割，宗徒们却是收割者。\Quote{然而他们并不因此被剥夺了劳苦的喜乐和报酬，反而与我们一同欢乐喜悦，虽然他们没有与我们一同收割。因为收割不像播种那样的工作。因此，我把你们留在这项工作上，其中辛劳较少而喜乐更大，而不是让你们去播种，因为播种有许多艰难和劳苦。收割时回报大，劳动却不那么大；甚至有很多便利。}祂用这些论证在此证明：\Quote{先知们的愿望是，所有人都要到我这里来。}法律也致力于此；为此他们播种，为能结出这果实。祂还显示祂也派遣了他们，并且新约与旧约之间有非常密切的联系，这一切祂都通过这个比喻同时达到了。祂还提及一个普遍流传的谚语。\par

% structure: p0020 source=h2 target=subsection reason=relative-heading-rank; base=h2

\subsection{若望福音 4:37}

众人常用这句话，当一方付出辛劳而另一方收获果实之时；祂说：\Quote{在这事例中，这句谚语尤其真实，因为先知们劳苦，而你们收获他们劳苦的果实。}祂没有说\Quote{报酬}（因为他们的巨大劳苦也并非没有报酬），而是说\Quote{果实}。达尼尔也如此行，因为他也提到一句谚语：\Quote{邪恶出于恶人}；达味在哀歌中也提到类似的谚语。因此祂预先说：\Quote{好使那播种的和那收割的一同欢乐。}因为祂即将宣告：\Quote{一人播种，另一人收割}，免得有人认为先知们被剥夺了报酬，祂便断言了一个奇怪且悖逆的事，这在感觉事物中从未发生，而只属于属神的事物。因为在感觉事物中，如果一人播种而另一人收割，他们不会\Quote{一同欢乐}，而是播种者忧伤，因为他们为别人劳苦，只有收割者单独欢乐。但在此并非如此，那些没有收割自己所播种的人，与收割者一同欢乐；由此清楚看出，他们也共享报酬。\par

% structure: p0022 source=h2 target=subsection reason=relative-heading-rank; base=h2

\subsection{若望福音 4:38}

藉此祂更加鼓励他们；因为当走遍全世界宣讲福音似乎是一件非常困难的事时，祂向他们表明，这甚至是最容易的。那非常困难的工作是另一项，需要巨大辛劳：播种，将未入门者的灵魂引向认识天主。但祂为何说出这些话？为的是当祂派遣他们去宣讲时，他们不会困惑，仿佛被派去完成一项艰难的任务。\Quote{因为先知们的那份工作更为困难，事实为我的话作证：你们来到了容易的工作上；因为如同在收割时，果实被轻易收集，一瞬间禾场上堆满了禾捆，它们不等待季节的轮回、冬天、春天和雨水，现在也是如此。事实大声宣告了这一点。}正当他说这些话时，撒玛黎雅人出来了，果实立刻被聚集起来。为此祂说：\Quote{你们举目向田观看，庄稼已经发白。}祂如此说了，事实便清晰了，言语藉着事件被证实（为真）。因为圣若望说：\par

% structure: p0024 source=h2 target=subsection reason=relative-heading-rank; base=h2

\subsection{若望福音 4:39}

他们看出，那妇人不会出于偏爱而赞美一位指责她罪过的人，也不会为了取悦别人而暴露自己的生活经历。\par

3. 那么，我们也应效法这位妇人，在自己犯罪时，不要以人为羞，而要敬畏那察看一切的天主，以及那将来惩罚不肯悔改者的天主。目前我们却反其道而行之，因为我们不敬畏那审判我们的主，反而畏惧那些不能伤害我们的人，并因来自他们的羞辱而战栗。因此，在我们所畏惧的事上，我们反而招致惩罚。因为现今只顾及人的指责，当主察看时却毫不羞耻地做非礼之事，且不愿悔改归正的人，在那一天将不仅在一两个人面前，而是在全世界面前成为耻辱的榜样。因为那里将有一大群观众，观看善行与恶行，这可以由绵羊与山羊的比喻教导你们，蒙福的保禄也说：\BibleQuote{因为我们众人必须出现在基督的审判台前，好使每人按他本身所行的，或善或恶，领取应得的报应。}\BibleRef{格后 5:10} 又说：\BibleQuote{祂要揭发黑暗中的隐密事。}\BibleRef{格前 4:5} 你行了或想了什么恶事，并向人隐瞒吗？但你却不能向天主隐瞒。然而你毫不关心；人的眼睛才是你所畏惧的。那么思想吧，在那一天你甚至无法逃避人的眼光；因为所有一切将像一幅画一样摆在我们眼前，以致每个人都将自我定罪。这从富人的例子便可以看出，我说的是\Person{拉匝禄}，那个被他忽视的穷人，他看见他站在眼前，并且他过去常常厌恶的手指，如今他恳求它能够成为他的安慰。因此，我劝告你们，即使没有人看见我们的行为，我们每人也要进入自己的良心，把理性立为审判官，带出自己的过犯，如果他不希望它们在那可怕的一天被公开揭露，就让他现在治愈自己的伤口，给它们敷上悔改的药方。因为一个身负万伤的人，完全能够痊愈地离开。因为祂说：\BibleQuote{如果你们宽恕人的过犯，你们的天父也必宽恕你们的过犯。}\BibleRef{玛 6:14} 因为正如在圣洗中所埋葬的罪不再显现，同样，如果我们愿意悔改，这些罪也将消失。悔改就是不再重犯；因为重蹈覆辙的人，就像狗呕吐又转回去吃，又像俗语所说把羊毛梳到火里，把水倒进破桶里。因此，我们应当在行动和思想上脱离我们所敢做的事，并脱离之后，应该用与我们的罪相反的药方来医治伤口。例如：你是否贪婪吝啬？戒除强取，以施舍医治伤口。你是否犯奸淫？戒除奸淫，以贞洁医治伤口。你是否诋毁兄弟并伤害他？停止指责，以仁爱医治伤口。让我们在每件犯罪的事上都这样做，不要漫不经心地放过我们的罪，因为以后、以后会有一个清算的季节。因此\Person{保禄}也说：\BibleQuote{主已经近了，应当一无挂虑。}\BibleRef{斐 4:5} 但也许我们必须加上与此相反的话：\Quote{主已经近了，应当谨慎。}因为那些生活在苦难、劳苦和试炼中的人，确实应该听到\Quote{一无挂虑}；但那些以强取豪夺为生、或生活奢华、并将要交出沉重账目的人，按理应该听到的不是这句话，而是另一句：\Quote{主已经近了，应当谨慎。}因为如今离终结已没有多少时日，世界正加快走向它的终点；战争宣告了这一点，苦难宣告了这一点，地震宣告了这一点，冷淡的爱也宣告了这一点。正如身体在最后喘息濒临死亡时，会引来万种痛苦；又如一所房屋将要倒塌时，屋顶和墙壁往往事先掉落许多部分；同样，世界的末日已近在咫尺，因此万种灾祸到处散布。如果那时主已经\Quote{近了}，那么如今祂更是\Quote{近了}。如果三百年前，当这些话被说出时，\Person{保禄}称那时为\Quote{时期的满全}，那么他更会称现在为如此。但也许正因如此，有些人不相信，然而他们正因此更应相信。因为人啊，你怎么知道终结不是\Quote{近了}，这些话要很快应验呢？因为正如我们说一年的末了，不是指最后一天，而是指最后一个月，尽管它有三十天；同样，如果我在许多年中称四百年为\Quote{末了}，也不算错；所以保禄那时预先称它为末了。那么让我们调整自己，以敬畏天主为乐；因为如果我们在此生活而不敬畏祂，祂的来临将突然袭击我们，那时我们既不谨慎，也不期待祂。正如基督所说：\BibleQuote{就如\Person{诺厄}和\Person{罗特}的日子一样，人子显现的日子也要这样。}\BibleRef{玛 24:37} \Person{保禄}也宣告说：\BibleQuote{几时人正说：\InnerQuote{平安无事}，危险就突然来到他们身上，如同孕妇的疼痛一样。}\BibleRef{得前 5:3} 所谓\Quote{如同孕妇的疼痛}是什么意思？孕妇常常在嬉戏、吃饭、洗澡或逛街时，对即将发生的事情毫无预感，突然被疼痛抓住。既然我们的情况类似，让我们时刻准备着，因为我们不会永远听到这些事，也不会永远有能力去做这些事。\Person{达味}说：\BibleQuote{在阴府里，谁能称颂祢？}\BibleRef{咏 6:5} 让我们在此悔改，好使我们在将来之日能得天主仁慈，获享充足的宽恕；愿我们众人藉着我们的主耶稣基督的恩宠和慈爱获得这一切，愿光荣与权能归于祂，从今世直到永远。阿们。\par

\SourceRecord{Translated by Charles Marriott.；Nicene and Post-Nicene Fathers, First Series；Vol. 14.；Edited by Philip Schaff.；Buffalo, NY: Christian Literature Publishing Co.,；1889.；<http://www.newadvent.org/fathers/240134.htm>.}

% structure: p0001 source=h1 target=section reason=page-title

\section{\WorkTitle{若望福音讲道 第三十五篇}}

% structure: p0002 source=h2 target=subsection reason=relative-heading-rank; base=h2

\subsection{若望福音 4:40-43}

没有比嫉妒与恶意更坏的了，没有比虚荣更有害的了；它惯于毁坏无数美善。因此，那些在知识上胜过撒玛黎雅人、且素常亲近先知的犹太人，竟因此显得不如他们。因为这些人甚至因那妇人的见证而相信，且未见任何奇迹，便前来恳求基督与他们同住；而犹太人虽目睹祂的奇事，非但不留祂在他们中间，反而驱逐祂，用尽一切方法将祂从他们的土地上赶出去，尽管祂的来临本是为他们的缘故。犹太人驱逐了祂，这些人却恳求祂与他们同住。那么，请告诉我，岂不更合宜的是，祂应接纳那些求告祂、恳求祂的人，而不是等候那些设计谋害祂、拒绝祂的人？而祂却将自己不给予那些爱祂、渴望留住祂的人吗？这实在不配祂的慈爱关怀；因此祂接纳了他们，并与他们同住了两天。他们渴望祂始终留在他们中间（因为福音作者已藉这段话表明：\BibleQuote{他们恳求祂与他们在同住}），但祂不曾这样做，只与他们同住了两天；在这些日子里，更多的人信了祂。然而，这些人本不可能相信，因为他们未见任何奇迹，且对犹太人怀有敌意；但因他们诚实地判断祂的话，这并未阻碍他们，反而使他们获得了超越障碍的认知，并竞相更加尊敬祂。因为福音作者说：\BibleQuote{他们对那妇人说：现在我们信，不是因为你的话；因为我们亲自听见了，并知道这确实是基督，世界的救主。}学生们超过了他们的女教师。他们有充分理由因相信祂、迎接祂而谴责犹太人。犹太人——为了他们的缘故，祂设计了整个计划——不断企图用石头砸祂；而这些人，当祂甚至无意来到他们那里时，却将祂吸引到自己身边。那些犹太人，即使见了奇迹，仍不改正；这些人，虽无奇迹，却对祂显示了极大的信德，并以此夸耀——他们不需奇迹就相信；而犹太人却不断要求奇迹并试探祂。\par

处处都需要诚实的灵魂；如果真理抓住这样的人，便轻易征服他；如果未征服他，那不是因为真理的虚弱，而是由于灵魂本身缺乏坦诚。正如太阳，当它遇到清澈的眼睛时，轻易地照亮它们；若未照亮，那是因它们自身的软弱，而非太阳的不足。\par

那么，听这些人说什么：\BibleQuote{我们知道这确实是基督，世界的救主。}你看见他们怎样立刻明白祂将把世界吸引到祂身边，祂来是为妥善安排我们共同的救恩，祂无意将祂的关怀局限于犹太人，而是要在各处撒播祂的圣言吗？犹太人并不这样，他们谋求建立自己的义，而不服从天主的义；这些人却承认众人都应受惩罚，与宗徒一同宣告：\BibleQuote{众人都犯了罪，缺少了天主的荣耀，因祂的恩宠自由地成义。}\BibleRef{罗 3:23}因为他们称祂为\BibleQuote{世界的救主}，就表明祂是面向一个失落的世界，并且祂不只是普通的救主，而是极其大能的一位。因为曾有先知和天使前来\Quote{拯救}，但这个人说：这是真正的救主，祂赐予真正的救恩，而非暂时的救恩。这出于纯正的信德。他们在两方面都令人钦佩：因为他们相信了，并且是在未见奇迹的情况下相信的（基督也称他们为\Quote{有福的}，说：\BibleQuote{那些没有看见而相信的人，是有福的。}\BibleRef{若 20:29}），而且他们是真诚地相信。虽然他们听见那妇人含糊地说：\Quote{这不是基督吗？}但他们并没有说：\Quote{我们也怀疑}或\Quote{我们以为}，而是说：\Quote{我们知道}，且不只是\Quote{我们知道}，而是\Quote{我们知道这确实是世界的救主。}他们承认基督不是众人中的一位，而是真正的\Quote{救主}。然而他们看见谁得救了？他们只是听了祂的话，便说出如同目睹许多伟大奇事时所说的话。为何福音作者不告诉我们这些话语，以及祂如何精彩地演讲？为使你明白，他们略过许多重要的事，却藉着结果将一切启示给我们。因为祂藉着自己的话说服了一整个民族、一整个城市。当听众未被说服时，写作者就不得不提及所说的话，免得有人因听者的迟钝而对发言者作出不利的判断。\par

\BibleQuote{两天后，祂离开那里，往加里肋亚去了。}\par

% structure: p0007 source=h2 target=subsection reason=relative-heading-rank; base=h2

\subsection{若望福音 4:44}

为何加上这句话？因为祂不是往葛法翁去，而是往加里肋亚去，又从那里往加纳。为避免你追问为何祂不与自己的百姓同住，却与撒玛黎雅人同住，福音作者说明了原因：他们不理会祂；为此祂没有去那里，以免他们的罪罚更大。因为我想，在此处祂称葛法翁为\Quote{祂的本乡}。为表明祂在那里未受尊敬，请听祂说：\BibleQuote{而你，葛法翁，已升到天上，却要被推下地狱。}\BibleRef{玛 11:23}祂称它为\Quote{祂自己的地方}，因为在那里祂宣讲了救恩工程的圣言，并且更常居住在那里。\Quote{那么，}有人说，\Quote{我们岂不见许多人在亲属中受尊敬吗？}首先，这类判断不可依据罕见的例子；其次，即使有些人在本乡受尊敬，他们在异乡也会受更多尊敬，因为熟识常使人容易受轻视。\par

% structure: p0009 source=h2 target=subsection reason=relative-heading-rank; base=h2

\subsection{若望福音 4:45}

你看，这些被人如此恶评的人，却最能来就祂。因为有一个说：\Quote{从纳匝肋还能出什么好事吗？}（\BibleRef{若 1:46}），另一个说：\Quote{你且查考，便知从加里肋亚不会出先知。}（\BibleRef{若 7:52}）他们说这些话是侮辱祂，因为许多人以为祂是纳匝肋人，他们还辱骂祂是撒玛黎雅人；有人说：\Quote{你是撒玛黎雅人，并且附有魔鬼。}（\BibleRef{若 8:48}）然而，看，撒玛黎雅人和加里肋亚人都信了，使犹太人蒙羞，而且撒玛黎雅人比加里肋亚人更佳，因为前者藉妇人的话就接待了祂，后者是在见了祂所行的奇迹之后才信的。\par

% structure: p0011 source=h2 target=subsection reason=relative-heading-rank; base=h2

\subsection{\BibleRef{若 4:46}}

圣史提起这奇迹，是为了彰显撒玛黎雅人的赞美。加纳的人因祂在耶路撒冷和那地方所行的奇迹而接待了祂；但撒玛黎雅人却不同，他们仅因祂的教导就接待了祂。\par

祂来到加纳，圣史已经说了，但没有补充祂来的原因。祂来到加里肋亚，是由于犹太人的嫉妒；但为何去加纳呢？起初祂是被邀请去赴婚宴；但如今呢？我想，是为了以祂的临在坚固那因祂的奇迹而植入的信德，并且通过自动前来、离开自己的家乡、偏爱他们，更加吸引他们归向祂。\par

\Quote{有一个王臣，他的儿子在葛法翁病了。}\par

% structure: p0015 source=h2 target=subsection reason=relative-heading-rank; base=h2

\subsection{\BibleRef{若 4:47}}

这人当然出身王族，或因官职而拥有某些尊严，带有\Quote{贵族}的称号。有些人认为这就是玛窦所提的那个人（\BibleRef{玛 8:5}），但无论从尊严还是信德来看，他都显为不同的人。那人即使基督愿意去，他也恳求祂留下；这人却当祂没有如此提议时，就拉祂到自己家中。那人说：\Quote{我不配你进我的屋顶下}；这人却恳求祂说：\Quote{请你趁我的孩子还没有死就下去。}在那事例中，祂从山上下来，进入葛法翁；但这里，祂从撒玛黎雅来，没有进入葛法翁，而是进入加纳，这人便遇见了祂。那人的仆人患了瘫痪，这人的儿子患了热症。\par

\Quote{他来求祂医治他的儿子，因为他快要死了。}基督怎么说？\par

% structure: p0018 source=h2 target=subsection reason=relative-heading-rank; base=h2

\subsection{\BibleRef{若 4:48}}

然而，他来求祂，这本身就是信德的标记。而且，之后圣史为他作证，说当耶稣说：\Quote{去吧，你的儿子活了}时，他信了祂的话，便去了。那么，祂在这里说的是什么呢？或者祂用这话来赞许撒玛黎雅人，因为他们不靠奇迹就信了；或者是为了触动那被认为是祂本城、且这人也属于的葛法翁。此外，路加福音中另一个人说：\Quote{主，我信}，又加上\Quote{帮助我的不信吧。}因此，如果这王臣也信了，却信得不完全或正确，这从他询问\Quote{孩子是什么时候退的热}可以看出，因为他想知道热病是自动退的，还是因基督的命令。因此，当他知道是\Quote{昨天第七时辰}时，那时\Quote{他自己和他全家就都信了}。\par

你看见吗？他是因了仆人的话，而不是因基督的话而信的。因此祂责备他来见祂时和说话时的心理状态（这样祂也更能引导他走向信德），因为他在奇迹之前信得不坚定。他来求祂并不稀奇，因为父母出于极大的爱心，不仅会去求他们信赖的医生，也会与不信赖的人交谈，希望不遗漏任何办法。事实上，他来时并无坚定目标，这从以下事实可见：当基督来到加里肋亚时，他才见到祂；而如果他坚定地信祂，他的孩子将死时，他绝不会犹豫去犹太。或者如果他害怕，那也是不可容忍的。\par

请看，这些话本身也显明此人的软弱。当他在基督责备了他的心理状态之后，本应对祂产生某种伟大的想法，即使他先前没有，请听他是如何拖在地上。\par

% structure: p0022 source=h2 target=subsection reason=relative-heading-rank; base=h2

\subsection{\BibleRef{若 4:49}}

好像祂不能使他死后复活，好像祂不知道孩子的状况似的。为此基督责备他，触动他的良心，以表明祂行奇迹主要是为了灵魂的好处。因为这里祂医治了心灵有病的父亲，不亚于儿子，为说服我们注意祂，不是因祂的奇迹，而是因祂的教导。因为奇迹不是为信者，而是为不信者和粗俗之人。\par

3. 那时，由于激动，王臣没有太留心这些话，或只留心那些关于他儿子的话；但他之后会想起所说的话，并从中获得最大的益处。事实正是如此。\par

然而，为什么在百夫长的事上祂主动提出要去，而在这里虽被邀请却不去呢？因为在前者，信德已成熟，因此祂同意去，让我们学习那人的善心；但这里王臣尚未成全。因此，当他不断催逼祂说：\Quote{下去}，并且还不清楚即使祂不在场也能医治时，祂便表明这对祂是可能的，为使他从耶稣的不去中，获得百夫长自己所拥有的那种认识。因此，当祂说：\Quote{除非你们看见神迹和奇事，你们总是不信}时，祂的意思是：\Quote{你们还没有正确的信德，仍然把我看作是一位先知。}因此，为了启示自己，并表明他即使没有奇迹也该信，祂说了对斐理伯也说过的话：\Quote{你信父在我内，我也在父内吗？或不然，你们因那些事而信吧。}（\EditorialNote{参阅：若 14:10-11}）\par

% structure: p0026 source=h2 target=subsection reason=relative-heading-rank; base=h2

\subsection{\BibleRef{若 4:51-53}}

你们看到这奇迹是何等明显吗？这孩子得到解脱并非简单平常，而是突然之间，这样所发生的事就被视为不是自然的结果，而是基督的作为。因为当他已来到死亡的门前时，正如他的父亲所说：\BibleQuote{趁我的孩子未死之前，下来吧}；他立刻就摆脱了疾病。这一事实也惊动了仆人们，因为他们或许来迎接主人，不仅是为了给他带来好消息，也认为耶稣的到来如今已是多余（因为他们知道主人已往那里去了），所以他们甚至在路上就遇见了他。这人从恐惧中解脱出来，从此逃入信德，渴望表明所发生的事是他此行的结果，并且从此他力求不使自己显得白费力气；于是他仔细查问了一切，\BibleQuote{他自己和全家都信了。}因为从此证据确凿无疑。因为那些没有在场、没有听见基督说话、也不知道时辰的人，当他们从主人口中听到某时某刻时，就有了祂大能不可辩驳的证明。因此他们也信了。\par

我们由此受到什么教导呢？不要等待奇迹，也不要寻求天主大能的保证。我看到许多人现在变得更加热心，当他们因孩子受苦或妻子患病而得到任何安慰时，但他们即使得不到，也应当坚持同样地感恩、光荣天主。因为这是正直仆人的本分，也是那些对主人怀有应有深情和爱戴之人的本分，不仅在被宽恕时，也在受责打时，都要投奔祂。因为这些也是天主慈爱眷顾的效果；经上说：\BibleQuote{主所爱的，祂必管教，又鞭打凡所收纳的儿子。} \BibleRef{希 12:6}因此，当一个人只在安逸时期事奉祂时，他就没有证明出伟大的爱，也没有纯洁地爱基督。我何必提健康、丰富、贫穷或疾病呢？即便你听到地狱或任何其他可怕的事，也不可停止赞美你的主人，而应因你对祂的爱忍受并做一切事。因为这是正直仆人和坚定不移灵魂的本分；凡如此存心的人，必能轻易忍受现世，获得将来的美善，并在天主面前享有极大的信心；愿我们众人赖我们的主耶稣基督的恩宠和慈爱都能获得，愿光荣归于祂，偕同圣父及圣神，自今至永远，及世之世。阿们。\par

\SourceRecord{Translated by Charles Marriott.；Nicene and Post-Nicene Fathers, First Series；Vol. 14.；Edited by Philip Schaff.；Buffalo, NY: Christian Literature Publishing Co.,；1889.；<http://www.newadvent.org/fathers/240135.htm>.}

% structure: p0001 source=h1 target=section reason=page-title

\section{第36篇讲道　论若望福音}

% structure: p0002 source=h2 target=subsection reason=relative-heading-rank; base=h2

\subsection{\BibleRef{若 4:54-5:1}}

如同在金矿中，善于辨认的人不会忽略哪怕是最细微的矿脉，认为它能产出巨大的财富；同样，在圣经中，若忽略了一笔一画，就必有损失，我们必须探究一切。因为这一切都是圣神所启示的，其中没有任何无用之物。\par

例如，请看福音作者在此处所说的：\BibleQuote{这又是耶稣从犹太往加里肋亚去以后行的第二次神迹。}他甚至特意加上了\Quote{第二次}这个词，原因何在？是为了更加彰显撒玛黎雅人的可赞美之处，表明即使行了第二次神迹，看见它的人仍不及那些一个神迹也未看见的人。\par

\BibleQuote{这些事以后，到了犹太人的一个节期。}这是什么\Quote{节期}？我想是五旬节。\BibleQuote{耶稣便上耶路撒冷去。}祂常在节期时进城，一方面是为与他们一同过节，另一方面是为吸引那些纯朴的群众；因为在这些日子里，尤其如此，那些较为单纯的人比平时更容易聚集。\par

% structure: p0006 source=h2 target=subsection reason=relative-heading-rank; base=h2

\subsection{\BibleRef{若 5:2-3}}

这是一种怎样的医治？它为我们预示了什么奥秘？因为这些事情不是随意或无目的地记载下来的，而是如同预象和预表，概要地显示出将来的事，以免那极其奇异之事突然出现时，在众人中损害信德的力量。那么，它们所预示的是什么呢？一种圣洗将要被赐予，它具有极大的能力，是最伟大的恩赐，能洗净一切罪过，使死人变为活人。这些事正是由那水池以及许多其他环境预先描绘出来的，如同图画一般。首先，它提供了一种水，可以洗净我们身体上的污秽，以及那些不是真正污秽而看似污秽的东西，例如因触摸死者、患麻风病以及其他类似原因所导致的不洁；在旧约之下，我们可以看到许多因水而成就的事。不过，现在让我们继续讨论当前的话题。首先，正如我先前所说，祂藉着水除去了我们身体的污秽，后来又除去了各种疾病。因为天主愿意使我们更接近对圣洗的信德，便不再只医治污秽，也医治疾病。那些在时间上更接近实体的预象，无论是关于圣洗、受难还是其他，都比更古老的那些更为清晰；正如王子身边的侍卫比前面的更为华丽一样，预象也是如此。\BibleQuote{有一位天使从天降下，搅动了那水，}使之具有治愈的能力，为叫犹太人可以知道，天使之主更能医治灵魂的疾病。然而，在这里，治愈的并非单纯是水的本性（否则便会经常发生），而是水与天使的作为相结合；同样，在我们身上，并非仅仅是水在起作用，而是当它接受了圣神的恩宠时，才除去我们一切的罪过。在这水池周围，\BibleQuote{躺着许多患病的人，有瞎眼的、瘸腿的、枯干的，等候水动，}但那时，疾病成了想要得医治之人的障碍；现在，每个人都有能力接近，因为搅动水的不是天使，而是天使之主亲自成就一切。病人如今不能说：\Quote{我没有一个人}；不能说：\Quote{当我来的时候，别人先我而下。}即使全世界的人都来，恩宠也不会耗尽，能力也不会枯竭，反而与先前一样强大。\par

首先，正如我先前所说，祂藉着水除去了我们身体的污秽，后来又除去了各种疾病。因为天主愿意使我们更接近对圣洗的信德，便不再只医治污秽，也医治疾病。那些在时间上更接近实体的预象，无论是关于圣洗、受难还是其他，都比更古老的那些更为清晰；正如王子身边的侍卫比前面的更为华丽一样，预象也是如此。\BibleQuote{有一位天使从天降下，搅动了那水，}使之具有治愈的能力，为叫犹太人可以知道，天使之主更能医治灵魂的疾病。然而，在这里，治愈的并非单纯是水的本性（否则便会经常发生），而是水与天使的作为相结合；同样，在我们身上，并非仅仅是水在起作用，而是当它接受了圣神的恩宠时，才除去我们一切的罪过。在这水池周围，\BibleQuote{躺着许多患病的人，有瞎眼的、瘸腿的、枯干的，等候水动，}但那时，疾病成了想要得医治之人的障碍；现在，每个人都有能力接近，因为搅动水的不是天使，而是天使之主亲自成就一切。病人如今不能说：\Quote{我没有一个人}；不能说：\Quote{当我来的时候，别人先我而下。}即使全世界的人都来，恩宠也不会耗尽，能力也不会枯竭，反而与先前一样强大。正如太阳的光芒每日照耀，却不会枯竭，其光明也不会因如此丰盛的供应而减弱；同样，圣神的能力也完全不因享受它的人数众多而减少。行这神迹的目的，是要人学会：藉着水可以医治身体的疾病，并且经过长时间的操练，便更容易相信它也能医治灵魂的疾病。\par

但是，耶稣为什么撇下其余的人，只来到那个患病三十八年的人面前？祂为什么问他：\BibleQuote{你愿意痊愈吗？}并非祂需要知道，而是为显明那人的恒心，并叫我们知道，祂正是因此才撇下别人来到他那里。那人说了什么？他说：\BibleQuote{主，是的，但是当水动的时候，没有人把我放到水池里，而当我来的时候，别人先我而下。}\par

正是为了叫我们得知这些情况，耶稣才问：\BibleQuote{你愿意痊愈吗？}祂并没有说：\Quote{你愿意我医治你吗？}（因为那时那人还没有对祂形成高超的观念），而是说\Quote{你愿意痊愈吗？}瘫痪者的恒心令人惊奇，他患病三十八年，每年都希望摆脱疾病，继续守候，没有退却。如果他不是非常恒心，难道过去的时间，即使不是将来的，还不足够使他离开那地方吗？请想一想，我请求你，既然水动的时间不确定，其他患病的人必定会如何地警觉。瘸腿的可能注意到，但瞎子如何看见？或许他们从喧哗声中得知了。\par

2. 那么，亲爱的，让我们羞愧，让我们羞愧，为我们过度的懈怠而叹息。那人等候了\Quote{三十八年}，没有得到他所渴望的，却没有退去。他之所以没有成功，并非由于自己的疏忽，而是因为受到别人压制和暴力，即便如此，他也没有变得灰心。而我们呢，若我们坚持十天为某事祈祷而没有得到，后来就懈怠得不愿再付出同样的热忱。我们对人尚且等待那么久，争战、受苦、从事奴仆般的工作，常常最终期望落空；但对我们确实会从祂那里获得超过我们劳苦的赏报的主（因为宗徒说：\BibleQuote{望德不叫人蒙羞}\BibleRef{罗 5:5}），我们却不愿以应有的勤勉等待祂。这该受何等惩罚！因为即使我们从祂那里得不到什么，难道我们不应该认为与祂不断交谈本身就是万福的缘由吗？\Quote{但是不断的祈祷是件辛苦的事。}哪一件属于德行的事不辛苦呢？\Quote{事实上，}有人说，\Quote{这一点充满了极大的困难：快乐与恶习相连，劳苦与德行相连。}我想，许多人把这一点当作一个问题。那么原因是什么呢？天主起初赐给我们无忧无虑、免于劳苦的生活。我们没有正确使用这份恩赐，反而因无所事事而堕落，并被逐出乐园。因此，祂使我们将来的生活成为劳苦的，仿佛向人类说明理由，说：\Quote{我起初允许你们过享乐的生活，但你们因自由而变得更糟，所以我命令从此将劳苦和汗水加在你们身上。}即使这劳苦也没有约束我们，祂随后赐给我们一条包含许多诫命的律法，将它加在我们身上，就像放在不驯服的马上的嚼子和缰绳，以抑制它的跳跃，正如驯马者所做的那样。这就是为什么生活是劳苦的，因为不劳苦往往会成为我们的毁灭。因为我们的本性不能忍受无所事事，而是容易转向邪恶。让我们假设，那节制的人以及正确履行其他美德的人不需要劳苦，他们在睡眠中做一切事，那么我们又会如何利用我们的安逸呢？难道不是用于骄傲和自夸吗？\Quote{但是为什么，}有人说，\Quote{极大的快乐被附着于恶习，极大的劳苦和辛劳被附着于德行？}如果事情不是困难的，你又能有什么感谢，又为什么能得到赏报呢？即便现在，我可以向你们指出许多人，他们天生厌恶与女人交往，视与她们交谈为不洁而避开；那么，我们要称这些人为贞洁的吗？要为他们加冕，宣告他们为胜利者吗？决不。贞洁是自我约束，是战胜那些争斗的快乐，正如在战争中，当战斗激烈时的战利品最为光荣，而不是当无人向我们举手时。许多人天生没有激情；我们要称他们为好脾气吗？绝不是。因此，主在提到三种阉人的状况后，让其中两种不加冕，而让一种进入天国。\BibleRef{玛 19:12}有人说：\Quote{但是邪恶有什么需要呢？}我也这么说。\Quote{那么，是什么使邪恶存在呢？}除了我们故意的疏忽，还有什么呢？\Quote{但是，}有人说，\Quote{应该只有好人。}那么，好人的特质是什么？是警醒和持守，还是睡觉和打鼾？\Quote{为什么，}有人说，\Quote{一个人不劳苦而行为正确似乎不好呢？}你的话适合牲畜或贪食者，或那些以自己的肚腹为神的人。为了证明这些话是愚昧的，请回答我这个问题。假设有一个国王和一个将军，当国王睡觉或醉酒时，将军忍受艰苦并树立了一座记功碑，你认为胜利属于谁？谁会享受所做的事的快乐？你看到灵魂更特别倾向于那些她为之劳苦的事物吗？因此，天主将劳苦与德行结合起来，希望使我们依恋她。为此，我们钦佩德行，尽管我们自己行为不正，而我们谴责恶习，即使它非常愉快。如果你说：\Quote{为什么我们不更钦佩那些天生善良的人，胜过那些因选择而善良的人呢？}我们回答：因为偏爱劳苦的人胜过不劳苦的人，是公正的。因为我们为什么劳苦？是因为你没有以节制承受不劳苦。不仅如此，如果仔细探究，懈怠也以其他方式常常毁掉我们，并给我们带来很多麻烦。如果你愿意，让我们把一个人关起来，只喂养和溺爱他，不允许他行走，也不带他出去工作，只让他享受餐桌和床铺，持续奢侈；还有什么比这样的生活更悲惨呢？\Quote{但是，}有人说，\Quote{工作是一回事，劳苦是另一回事。}是的，但那时人能够不劳苦而工作。\Quote{这可能吗？}他说。是的，这是可能的；天主甚至愿意如此，但你们却不承受。因此，祂把你放在园中工作，指定了职业，但并未与劳苦结合。因为如果人起初就劳苦，天主后来就不会以劳苦作为惩罚。因为工作是可能的，而不会疲倦，如同天使那样。为证明他们工作，请听大卫所说：\BibleQuote{你们这些大有能力、遵行祂话语的。}\BibleRef{咏 103:20}现在力量的缺乏导致许多劳苦，但那时并非如此。因为有人说：\BibleQuote{那进入祂安息的，就歇了自己的工，正如天主歇了祂的工。}\BibleRef{希 4:10}这里不是指懒惰，而是指停止劳苦。因为天主甚至现在仍在工作，正如基督所说：\BibleQuote{我父作工直到现在，我也作工。}\BibleRef{若 5:17}因此，我劝勉你们，放下一切懈怠，为德行而发热心。因为邪恶的快乐是短暂的，痛苦是持久的；相反，德行的喜乐不老旧，劳苦只是暂时的。德行甚至在分派冠冕之前就鼓舞她的工人，以希望喂养他；恶习甚至在报复之前就惩罚为她工作的人，折磨和恐吓他的良心，使它易于想象一切（邪恶）。这些难道不比任何劳苦、任何辛劳更糟吗？如果事情并非如此，如果还有快乐，还有什么比那种快乐更无价值的呢？因为它一出现就飞走了，在尚未抓住之前就枯萎和逃逸，无论你说的是美貌的快乐、奢侈的快乐，还是财富的快乐，因为它们每日都在不断消逝。此外，当（除了这快乐）还有惩罚和报复时，还有什么比追求它的人更可悲的呢？既然如此，让我们为德行忍受一切，这样我们就会享受真正的快乐，借着我们的主耶稣基督的恩宠和慈爱，愿荣耀归于祂，及父及圣神，从今时直到永远，世世无尽。阿们。\par

\SourceRecord{Translated by Charles Marriott.；Nicene and Post-Nicene Fathers, First Series；Vol. 14.；Edited by Philip Schaff.；Buffalo, NY: Christian Literature Publishing Co.,；1889.；<http://www.newadvent.org/fathers/240136.htm>.}

% structure: p0001 source=h1 target=section reason=page-title

\section{讲道集第37篇 论\WorkTitle{若望福音}}

% structure: p0002 source=h2 target=subsection reason=relative-heading-rank; base=h2

\subsection{\BibleBook{若望}{福音} 5:6-7}

1. 神圣圣经的益处极大，来自圣经的援助也全然充足。\Person{保禄}曾这样宣告说：\BibleQuote{先前所写的一切，都是为教训我们而写的，因为我们这末世的人，借着圣经的忍耐和安慰，可以怀有希望。}\BibleRef{罗 15:4} 因为天主的话语是各类药物的宝库，无论需要扑灭骄傲、抑制欲望、践踏贪财、轻视痛苦、鼓舞信心、获得忍耐，都能从中找到丰富的资源。因为那些长期与贫穷搏斗、或困于重病的人，读了我们面前的这段经文，谁能不获得极大的安慰呢？这个人瘫痪了三十八年，每年都看见别人得释放，而自己却被疾病捆绑，即使如此，他也没有退却而失望；其实，不仅过去的沮丧，而且未来的绝望，也足以使他崩溃。现在且听他说什么，并学习他受苦之深。因为当基督说：\BibleQuote{你愿意痊愈吗？} \BibleQuote{主，是的，}他说，\BibleQuote{但我没有人，当水动的时候，把我放进池子里。} 还有什么比这些话更可怜的呢？还有什么比这些处境更悲惨的呢？你看见一颗因长期疾病而破碎的心吗？你看见一切暴力都被制伏了吗？他没有说出亵渎的话，也没有像我们常见许多人在逆境中那样，咒骂自己的日子，也没有对问题发怒，也没有说：\Quote{你来是嘲弄我们吗？你问我是否愿意痊愈吗？} 而是温和地、极其柔和地回答：\BibleQuote{主，是的}；然而他并不知道问话的是谁，也不知道祂要治愈他，但他仍然温和地陈述一切情况，并没有进一步要求，就好像在向一位医生说话，只想讲述自己的痛苦。也许他希冀基督能够在这件事上帮助他，把他放进水里，并想用这些话引起祂的注意。那么耶稣说什么呢？\par

% structure: p0004 source=h2 target=subsection reason=relative-heading-rank; base=h2

\subsection{\BibleBook{若望}{福音} 5:8}

现在有些人认为这就是\BibleBook{玛窦}{福音}中那个\BibleQuote{躺在床上}\BibleRef{玛 9:2}的人；但并非如此，在许多方面都很明显。首先，因为他需要有人为他出面。那个人有许多人照顾和抬他，而这个人一个也没有；因此他说：\BibleQuote{我没有人。} 其次，从回答的方式来看：那一个人没有说一句话，而这个人讲述了他全部的情况。第三，从季节和时间来看：这个人是在一个庆节，在安息日被治愈的，那一个人则是在另一天。地方也不同：一个是在房屋里治愈的，另一个是在池边。治愈的方式也不同：在那里基督说：\BibleQuote{你的罪赦了}，但在这里祂先坚固身体，然后照顾灵魂。在那里有罪的赦免（因为祂说：\BibleQuote{你的罪赦了}），但在这里有警告和威胁，以坚固此人将来：\BibleQuote{不要再犯罪，免得遭遇更坏的事。}\BibleRef{若 5:14} 犹太人的指责也不同：在这里他们反对耶稣在安息日工作，在那里他们指责祂亵渎。\par

现在，我请你们思想天主的无限智慧。祂没有立刻扶起那人，而是先通过问话使他熟悉，为即将来临的信德铺路；祂不仅扶起他，还命令他\BibleQuote{拿起床来}，以证实所行的奇迹，使没有人认为所发生的是幻觉或做戏。因为除非他的四肢已经坚固而完全连接，否则他不能拿起他的床。基督常这样做，有效地堵住那些想要傲慢的人的口。在饼的事上也是如此，为使人不能断言那些人只是吃饱了，所行的是幻觉，祂使饼有许多剩余。同样，对洁净的癞病人，祂说：\BibleQuote{你去，把你自己给司祭看}\BibleRef{玛 8:4}；同时提供了洁净的确切证据，并堵住了那些断言祂立法反对天主之人的无耻之口。在酒的事上，祂也照样做了；因为祂不仅把酒指给他们看，还叫人把它送到管宴席的人那里，为使一个对所发生之事一无所知的人，通过他的承认，为祂作一个无偏见的见证；所以福音作者说，管宴席的人\BibleQuote{不知道从哪里来的}，以此显示他的见证公正。在另一处，当祂复活死人时，祂说：\BibleQuote{你们给他吃的吧}；提供了真正复活的证据，并以此甚至说服愚蒙的人，祂不是骗子，不是搞幻觉的，而是为了整个人类本性的救恩而来的。\par

2. 但为什么耶稣没有要求这个人的信德，像在其他人身上那样，说：\BibleQuote{你信我能做这事吗？} 这是因为那人还不清楚知道祂是谁；并非在行奇迹之前，而是在行奇迹之后，才看到祂这样做。对于那些曾看见祂向别人施展权能的人，合理地对他说这话；而对于那些还不知道祂是谁、后来要通过标记认识祂的人，是在奇迹之后才要求信德。所以\Person{玛窦}没有把基督说这句话放在祂奇迹的开端，而是在祂治愈许多人之后，只对着那两位盲人。\par

然而，请看这瘫子的信德是何等样。当他听见说：\BibleQuote{拿起你的床，行走吧！}，他并没有讥笑，也没有说：\Quote{这是什么意思？天使下来搅动那水，只医治一个人，而你，一个凡人，竟想凭一句话和命令就能做比天使更大的事？这不过是虚空、夸口、嘲弄。}但他没有说也没有想任何类似的话，而是立刻听见就起来了，痊愈了，没有违抗那发命令的；因为他立刻痊愈了，就\BibleQuote{拿起他的床，行走了。}随后发生的事更令人钦佩。起初，当没有人搅扰他时，他相信，这并不那么奇妙；但后来，当犹太人满怀疯狂，从四面八方逼迫他，控告他，围困他，说：\BibleQuote{不许你拿起你的床！}时，他却不理会他们的疯狂，反而在众人中极其勇敢地宣告了他的恩主，并使他们的无耻舌头沉默——我说，这乃是极大的勇气之举。因为当犹太人起来攻击他，并以责骂和傲慢的态度对他说（\Quote{这是不合法的……}）：\par

% structure: p0009 source=h2 target=subsection reason=relative-heading-rank; base=h2

\subsection{\BibleRef{若 5:10-11}}

请听他所说的话（\BibleQuote{那使我痊愈的……}），他几乎是在说：\Quote{你们真是愚蠢而疯狂，竟吩咐我不要以那解救了我长久沉重疾病的人为我的老师，也不要听从他所命令的任何事。}如果他选择不公正地行事，他可能会说不同的话，比如这样：\Quote{我这样做不是出于自己的意愿，而是出于另一个人的命令；如果这件事该受责备，就责备那下命令的人，我会放下床。}他也可能隐瞒这治愈，因为他深知他们如此烦恼，不是因为违反安息日，而是因为他的病弱被治愈了。然而他既没有隐瞒这事，也没有那样说，更没有请求原谅，反而大声宣认并宣扬这恩惠。瘫子就是这样做的；但看一看他们行事何等不公。因为他们没有说：\BibleQuote{谁使你痊愈了？}——在这点上他们沉默，却不断提出表面的过犯。\par

% structure: p0011 source=h2 target=subsection reason=relative-heading-rank; base=h2

\subsection{\BibleRef{若 5:12-13}}

耶稣为何隐藏自己？首先，为的是在祂不在场时，那人的见证无可怀疑，因为如今感到自己痊愈的人乃是这恩惠的可信证人。其次，为的是不致使犹太人的愤怒更加激化，因为见到自己所嫉妒的人，在存恶意者心中常会点燃不小的火苗。因此祂退开了，让这行为本身在他们中间为自己辩护，以免祂亲自为自己说什么，而是让被治愈的人以及控告者去说。就连这些后者也暂时为这奇迹作证，因为他们没有说：\BibleQuote{你为什么吩咐在安息日做这些事？}而是说：\BibleQuote{你为什么在安息日做这些事？}这并非因着过犯而不悦，而是因着瘫子得复原而嫉妒。然而，就其是人的工作而言，瘫子所做的更算是一项工作，因为另一项只是一句话和一个字。这里祂命令别人违反安息日，但在别处祂自己也这样做，和了泥，抹在人的眼睛上\EditorialNote{参\BibleRef{若 9}}；然而祂做这些事并非违反法律，而是超越法律。这一点我们以后将谈论。因为当犹太人因安息日而控告祂时，祂并非总是用同样的话为自己辩护，这一点我们必须仔细观察。\par

3. 但是，让我们想一想，嫉妒是多么大的邪恶，它如何使灵魂的眼睛失明，以致危及那被嫉妒所占有之人的救恩。因为如同疯子常常把剑刺向自己的身体，心怀恶意的人也只盯着一件事，即伤害他们所嫉妒的人，而不关心自己的救恩。这些人比野兽更坏；野兽在饥饿时，或先被我们激怒时，才会武装起来对付我们；但这些人受到恩惠时，常常以怨报德，好像被亏待了一样。他们比野兽更坏，如同魔鬼，或许比魔鬼更坏；因为魔鬼虽然对我们怀有不断的敌意，但并不会谋害自己的同类（因此耶稣以此使犹太人哑口无言，当他们说祂是靠着贝耳则步驱魔时），但这些人既不尊重共同的本性，也不顾惜自己。因为他们在激怒所嫉妒的人之前，先激怒了自己的灵魂，使灵魂充满各种烦恼和沮丧，无益而徒劳。人啊，你为什么要因邻人的兴旺而忧伤呢？我们应该因自己所受的苦难而忧伤，而不是因为看到别人享有好名声。因此，这种罪没有任何借口。奸淫者可以推脱他的情欲，窃贼可以推脱他的贫穷，杀人者可以推脱他的愤怒——这些借口虽然牵强无理，但他们仍可提出。但是，告诉我，你会提出什么理由呢？除了极端的邪恶之外，根本没有任何理由。如果我们被命令爱仇人，那么当我们恨自己的朋友时，该受怎样的惩罚？如果爱那些爱我们的人的人，其境况并不比外邦人好，那么伤害那些没有亏待他的人，还有什么借口、什么开脱呢？请听保禄所说的：\BibleQuote{即使我将身体焚烧，却没有爱，对我毫无益处} \BibleRef{格前 13:3}；现在，每个人都清楚，哪里有嫉妒和恶意，哪里就没有爱德。这种情感比奸淫和通奸更坏，因为后者只影响行奸淫的人，而嫉妒的暴政却倾覆了整个教会，毁灭了全世界。嫉妒是谋杀之母。因着嫉妒，加音杀了亚伯；因着嫉妒，厄撒乌（几乎杀了）雅各伯，他的兄弟们也（几乎杀了）若瑟；因着嫉妒，魔鬼（毁灭了）全人类。你确实现在没有杀人，但你做了许多比杀人更坏的事：渴望你的兄弟行为不当，四面设下陷阱，瘫痪他在德行方面的努力，因他取悦世界之主而悲伤。然而，你并不是与你的兄弟争战，而是与他所服务的那位争战；当你把你的荣耀置于祂的荣耀之上时，你是在侮辱祂。而事实上最糟糕的是，这罪似乎无关紧要，但实际上比任何其他罪都更严重；因为即使你施行仁慈、守夜并禁食，如果你嫉妒你的兄弟，你就比任何人都更受诅咒。从以下情况也可以清楚看出：曾有一个格林多人犯奸淫，但他被指控后不久就恢复了成义；加音嫉妒亚伯，却没有得到医治，虽然天主亲自不断安抚伤口，他反而更加痛苦和动荡，并被催促去杀人。因此，这种情欲比那种（奸淫）更坏，除非我们留心，否则它不易被治愈。那么，让我们无论如何将它连根拔起，考虑到这一点：当我们因嫉妒他人的祝福而得罪天主时，当我们与他们一同喜乐时，我们就悦乐祂，并使我们自己成为为义人所存留的好事的参与者。因此，保禄劝勉我们说：\BibleQuote{应与喜乐的一同喜乐，与哭泣的一同哭泣} \BibleRef{罗 12:15}，这样我们就能在两方面获得极大的益处。\par

那么，既然即使我们不劳苦，因着与劳苦者一同喜乐，我们也成为他冠冕的分享者，就让我们抛开一切嫉妒，将爱德植根于我们的灵魂中，这样，通过称赞那些悦乐天主的弟兄们，我们就能获得现在和将来的福乐，借着我们的主耶稣基督的恩宠和慈爱，愿光荣归于祂，并与祂一起，归于圣父和圣神，从今世直到永远，无穷世之世。阿们。\par

\SourceRecord{Translated by Charles Marriott.；Nicene and Post-Nicene Fathers, First Series；Vol. 14.；Edited by Philip Schaff.；Buffalo, NY: Christian Literature Publishing Co.,；1889.；<http://www.newadvent.org/fathers/240137.htm>.}

% structure: p0001 source=h1 target=section reason=page-title

\section{若望福音讲道 第三十八篇}

% structure: p0002 source=h2 target=subsection reason=relative-heading-rank; base=h2

\subsection{若望福音 5:14}

1. 罪是一件可畏的事，是可畏的，灵魂的毁灭；其恶行往往因过盛而泛滥，甚至攻击人的身体。因为大多数时候，当灵魂患病时，我们感觉不到痛苦；但身体即使受一点小伤，我们也会竭力摆脱它的病弱，因为我们察觉到了病弱。因此，天主常常为灵魂的过犯而惩罚身体，以便通过较低部分的鞭笞，使较优部分也能获得某种医治。同样，在格林多人中，保禄也恢复了那犯奸淫的人，藉肉身的毁灭来抑制灵魂的疾病，将刀施加于身体，如此抑制了那邪恶（\BibleRef{格前 5:5}）；如同一位优秀的医生，对水肿或脾脏病症使用外部烧灼，当它们拒绝接受内服药物时。基督对瘫子也这样做；祂在说话时表明：\BibleQuote{看，你已经痊愈了；不要再犯罪，免得遭遇更坏的事。}\par

那么我们从中学到什么？首先，他的疾病是由他的罪产生的；其次，应相信地狱之火的道理；第三，惩罚是长久的，甚至无穷的。现在那些说\Quote{我在一小时之内杀了人，我在片刻之间犯了奸淫，难道我要受永罚吗？}的人在哪里？因为看，这个人犯罪的时间并不如他受苦的年数那样长，他在惩罚的长度中度过了整个一生；罪过不是按时间，而是按过犯的性质受审判。除此之外，我们还可以看到另一件事：即使我们因先前的罪受了严厉的惩罚，如果后来我们重蹈覆辙，我们将受更严厉的惩罚。这理应是合理的；因为一个人即使受惩罚也不能改善，那么他就被视为麻木不仁和藐视者，之后会遭受更重的惩罚。过失本身应当足以使失足者警惕并更清醒，但当加上惩罚仍不能生效时，他自然需要更痛苦的折磨。如果即使在这个世界，当我们在受罚后仍犯同样的罪时，我们受到的惩罚比之前更严厉，那么当我们犯罪后根本没有受到惩罚时，岂不更当极其惧怕和战栗，如同将要遭受某种不可挽回的事？\Quote{那么，}有人说，\Quote{难道不是所有的人都这样受罚吗？}因为我们看见许多坏人身体健壮、精力充沛、并享受极大的幸福。但让我们不要自信，我们要最为他们哀悼，因为他们在此生未受任何苦，反而助长了来世更严厉的报复。正如保禄所说：\BibleQuote{但我们受审的时候，是被主惩戒，免得我们和世界一同被定罪}（\BibleRef{格前 11:32}）；因为此世的惩罚是警告，来世的惩罚是报复。\par

\Quote{那么，}有人说，\Quote{所有的疾病都是从罪来的吗？}不是全部，但大部分是的；有些是从不同的放荡生活而来，因为贪食、放纵和懒惰会产生类似的痛苦。但我们应遵守的一条规则是：感恩地承受每一次打击；因为它们是为我们的罪而加给的，就像在列王纪中我们看到一人患痛风（\BibleRef{列上 15:23}）；它们也是为了使我们成为可嘉许的，正如主对约伯说：\BibleQuote{你以为我向你说话，只是为了显你为义吗？}（\BibleRef{约 60:8}）\par

但是为什么在这些瘫子的案例中，基督提起了他们的罪？因为祂也对\BibleBook{玛窦}{福音}中躺在床上的人说：\BibleQuote{孩子，放心罢！你的罪赦了}（\BibleRef{玛 9:2}）；又对这一个人说：\BibleQuote{看，你已经痊愈了；不要再犯罪。}我知道有些人毁谤这个瘫子，断言他是控告基督的人，因此这段话是对他说的；那么对于\BibleBook{玛窦}{福音}中的另一个人，他听到了几乎相同的话，我们该说什么呢？因为基督也对他说：\BibleQuote{你的罪赦了。}由此清楚可见，这个人也不是因为他们所声称的原因得此话语。从下文我们可以看得更清楚；因为福音作者说：\Quote{后来耶稣在殿里遇见他}，这是他极大虔诚的表示；因为他没有去市场或散步，也没有沉溺于奢侈和安逸，而是留在圣殿里，尽管即将承受如此猛烈的攻击并受到那里所有人的骚扰。然而这些事没有一件劝他离开圣殿。此外，基督遇见他之后，即使祂已与犹太人交谈过，也没有暗示这类事。因为如果祂想指控他，祂会对他说：\Quote{你又在犯从前的罪吗？你难道没有因你的痊愈变得更好吗？}但祂没有说这类话，只是为将来而保全他。\par

2. 那么，为什么当祂治愈了跛子和残废者时，祂没有在任何情况下提及类似的事？我认为这些（瘫子的）疾病是由罪恶的行为产生的，而那些人的疾病来自自然的虚弱。或者如果不是这样，那么藉由这些人以及对它们所说的话，祂也对其他人说了话。因为既然这种疾病比其他任何疾病更严重，祂藉由较严重的纠正较轻的。而且，正如祂在治愈另一个人时，命令他把光荣归于天主，将这一劝告不仅针对他，而且通过他针对所有人，同样祂也对这些，并通过这些对全人类，当面给予劝告和建议。除此之外，我们还可以说，耶稣察觉到他的灵魂有极大的忍耐，便对他说话，如同对一位能接受祂命令的人，藉由恩惠和对未来祸患的恐惧来保持他的健康。\par

请注意祂如何不夸耀自己。祂没有说：\Quote{看，我使你痊愈了}，而是说：\Quote{你已经痊愈了；不要再犯罪了}。又不说：\Quote{免得我惩罚你}，而是说：\Quote{免得你遭遇更坏的事}。这两种说法都未以第一人称，并且表明治愈出于恩宠而非功绩。因为祂并未向那人宣告，他是在受了应得的惩罚之后才获得释放，而是因着慈爱而痊愈。若非如此，祂便会说：\Quote{看，你已为你的罪受了足够的惩罚，以后要坚定不移}。但如今祂没有这样说，而是如何说的呢？\Quote{看，你已经痊愈了；不要再犯罪了}。让我们不断向自己重复这些话；若我们在受惩戒之后获得释放，愿每个人对自己说：\Quote{看，你已经痊愈了；不要再犯罪了}。但若我们虽继续行同样的事却未受惩罚，就让我们以宗徒的那句话作为我们的警诫：\Quote{天主的仁慈引导（我们）悔改，但（我们）因着顽固和不悔的心，（为自己）积蓄愤怒}\BibleRef{罗 2:4}。\par

祂不仅通过强壮病人的身体，也以另一种方式向那人提供了祂天主性的大证据；因为祂说\Quote{不要再犯罪了}，便表明祂知道那人从前所犯的一切过犯；借此祂会使那人相信关于未来的事。\par

% structure: p0010 source=h2 target=subsection reason=relative-heading-rank; base=h2

\subsection{若望福音 5:15}

再留意那人如何保持同样的正确心态。他没有说：\Quote{这就是那说\InnerQuote{拿起你的床}的人}，但当他们不断提出这表面上的指控时，他也不断提出辩护，再次宣告他的治愈者，并设法吸引和引导他人归向祂。因为他并非那样无情，以致在如此恩惠和指控之后背叛他的恩人，并怀着恶意说出他所说的话。倘若他是野兽，倘若他不像人而是石头，那恩惠和恐惧也足以约束他，因为他内怀威胁，生怕遭受\Quote{更坏的事}，既已得到他医师之权能的最大保证。此外，若他想要毁谤祂，他就不会提及自己的痊愈，而会提出并反对祂违反安息日。但事实并非如此，绝非如此；这话是出于极大的勇气和坦率；他宣扬他的恩人不亚于那个盲人。因为那盲人说了什么？\Quote{祂和了泥，抹在我的眼睛上}\BibleRef{若 9:6}；同样，我们现在所说的这个人说：\Quote{是耶稣使我痊愈了。}\par

% structure: p0012 source=h2 target=subsection reason=relative-heading-rank; base=h2

\subsection{若望福音 5:16}

基督怎么说呢？\par

% structure: p0014 source=h2 target=subsection reason=relative-heading-rank; base=h2

\subsection{若望福音 5:17}

当需要为门徒辩护时，祂引用了他们的同仆达味，说：\Quote{你们没有念过达味饿了的时候做了什么吗？}\BibleRef{玛 12:2}但当需要为祂自己辩护时，祂便转向圣父，以两种方式表明祂的平等：特别称天主为祂的父，并做与祂相同的事。\Quote{为何祂没有提那在耶里哥所发生的事呢？}因为祂要提升他们脱离尘世，使他们不再只视祂为人，而是视为天主，视为应当立法的那一位；因为若非祂是真正的圣子并同体，这辩护就比指控更糟。倘若一位总督更改了王法，当被问及时如此辩解，说：\Quote{是的，因为王也废除了法律}，他无法逃脱，反而会加重指控。但在这里，既然尊严相等，辩护便在最稳固的基础上得以圆满。\Quote{从你们赦免天主的那些指控中}，祂说，\Quote{也赦免我吧。}因此祂先说\Quote{我的父}，为要说服他们即便不情愿，也因着对祂明确宣称的儿子身份的敬意，而允许祂同样的事。\par

若有人说：\Quote{圣父在第七日停了一切工，又怎能\InnerQuote{工作}呢？}让他了解祂\Quote{工作}的方式。那么祂工作方式是怎样的？祂眷顾、维系一切受造之物。因此，当你看见太阳升起、月亮运行、湖泊、泉源、河流、雨水、种子和我们肉身及无理性生物中的自然程序，以及构成这宇宙的一切其余事物时，就学会圣父不断的工作。\Quote{因为祂叫太阳上升，照恶人也照善人；降雨给义人也给不义的人。}\BibleRef{玛 5:45}又说：\Quote{天主这样装饰田野的草，今天还在，明天就被投入炉中}\BibleRef{玛 6:30}；论到飞鸟祂说：\Quote{你们的天父养活它们。}\par

3. 在那里，祂在安息日只以言语行了一切，没有增添什么，却以殿中所行的事和他们的实践驳斥了他们的指控。但在这里，祂命令做一件事，即拿起床（从奇迹看，这并非什么大事，但借之祂表明了一点：明显违反安息日），祂将论述引向更高的事，更愿借圣父之尊严威慑他们，引导他们上升到更高思想。因此，当祂论及安息日时，祂并不只作为人或只作为天主辩护，而是有时用这种方式，有时用另一种方式；因为祂要说服他们相信救恩计划的屈尊和祂天主性的尊严。所以现在祂以天主的身分辩护；因为若祂总是只以人的方式与他们交谈，他们就会停留在同样卑下的状态。为免如此，祂引出圣父。然而受造物本身在安息日也\Quote{工作}（因为太阳奔跑，河流流动，泉源涌出，妇女生育），但为使你学习祂不属于受造物，祂不说：\Quote{是的，我工作，因为受造物工作}，而是说：\Quote{是的，我工作，因为我的父工作。}\par

% structure: p0018 source=h2 target=subsection reason=relative-heading-rank; base=h2

\subsection{若望福音 5:18}

祂不仅以言语，更以行动证实了这一点；因为祂不仅以言词，更以更频繁的行动宣告了这一点。为何如此？因为他们可能会反驳祂的话语，指控祂傲慢；但当他们看到祂行动的真实性由结果证明，祂的能力由作为彰显之后，他们便无可反驳了。\par

然而，那些不肯以正确心思领受这些话语的人断言，\Quote{基督并没有使自己与天主平等，而是犹太人如此猜测。}那么，让我们从头回顾一下所说的话。告诉我，犹太人迫害祂，还是没迫害？人人都清楚他们迫害了。他们是为这个还是为别的而迫害？同样公认是为这个。祂到底有没有违反安息日？祂确实违反了，无人能否认。祂称天主为祂的父，还是没有这样称呼？这也是真的。那么其余的也以同样的方式随之而来；因为称天主为祂的父、违反安息日、并因此前特别是后者而受犹太人迫害，这不属于虚假的想象，而是属于实际事实；因此使自己与天主平等，正是同一意义的宣告。\par

这一点从祂先前所说的话中可以看得更清楚，因为\BibleQuote{我父工作，我也工作}，这是宣告自己与天主平等的话。因为在这些话中，祂没有指出任何区别。祂没有说\Quote{祂工作，我服务}，而是\Quote{正如祂工作，我也工作}，并宣告了完全平等。但如果祂不希望确立这一点，而犹太人却毫无理由地这样认为，祂就不会容许他们的心思被欺骗，而会纠正这一点。此外，福音作者也不会对此缄口不言，而会明确地说，犹太人如此猜测，但耶稣并没有使自己与天主平等。正如在另一处，祂做了同样的事，当祂察觉到某话以一种方式说出，却被以另一种方式理解时；例如，基督说\BibleQuote{你们拆毁这座圣殿，三天之内我要把它重建起来}\BibleRef{若 2:19}，指的是祂的身体。但犹太人不明白这一点，以为这话是指犹太人的圣殿，便说\BibleQuote{这圣殿建造了四十六年，你能在三天内把它重建起来吗？}既然祂说了一件事，而他们想象了另一件（因为祂说的是祂的身体，而他们以为这话是指他们的圣殿），福音作者对此加以评论，或者更确切地说，纠正他们的想象，接着说道\Quote{祂却以祂的身体为圣殿。}所以，在这里，如果基督没有使自己与天主平等，没有希望确立这一点，而犹太人却想象祂如此行，作者在这里也会纠正他们的猜测，并且会说\Quote{犹太人以为祂使自己与天主平等，但祂其实并非谈论平等。}这不仅在这一处、由这位福音作者如此行，而且在另一处，另一位福音作者也被看见做同样的事。因为当基督警告门徒说\BibleQuote{你们要提防法利塞人和撒杜塞人的酵母}\BibleRef{玛 16:6}，他们彼此议论说\Quote{这是因为我们没有带饼吧}，而祂说了一件事，称他们的教导为\Quote{酵母}，但门徒想象了另一件，以为这话是关于饼的；这时不是福音作者纠正他们，而是基督亲自这样说\Quote{你们为什么不明白，我不是对你们论饼呢？}但在这里，并无此类情形。\par

\Quote{但是，}有人说，为了消除这种想法，基督接着说道，\par

% structure: p0023 source=h2 target=subsection reason=relative-heading-rank; base=h2

\subsection{\BibleBook{若望}{福音} 5:19}

人啊！祂所做的恰恰相反。因为祂说这话不是要消除，而是要证实祂的平等。但请留心，因为这不是一个普通的问题。\Quote{凭自己}这一表达在圣经中多处出现，涉及基督和圣神，我们必须领会这个表达的力量，以免陷入极大的错误；因为如果按照显而易见的理解方式单独取其字面意思，请思考将产生多么荒谬的后果。祂不是说自己能做某些事而不能做另一些事，而是普遍地说：\par

4. \BibleQuote{子凭自己不能做什么。}那么我问我的对手：\Quote{子凭自己不能做什么，请告诉我？}如果他回答说\Quote{祂不能做什么}，我们将说，祂凭自己行了最大的善事。正如保禄大声宣告说：\BibleQuote{祂虽具有天主的形体，却不以自己与天主同等为应当把持的，反而使自己空虚，取了奴仆的形体。}\BibleRef{斐 2:6} 而基督在另一处又说：\BibleQuote{我有权舍弃我的生命，也有权再取回它来}；并且\BibleQuote{没有人能夺去我的生命，而是我自愿舍弃的。}\BibleRef{若 10:18} 你看见祂对生命和死亡拥有权柄，并且凭自己完成了如此伟大的救恩工程吗？我为何论基督呢？就连我们这些比什么都卑微的人，也凭自己做了许多事：我们凭自己选择罪恶，凭自己追求德行，如果我们不凭自己、没有能力这样做，那么做错了就不会下地狱，做对了也不会享天国。\par

那么，\BibleQuote{凭自己不能做什么}是什么意思？意思是祂不能做任何与圣父相悖的事，不能做任何与圣父无关、相异的事，这尤其宣告了平等和完全的一致。\par

可是祂为甚么不说\Quote{祂不作相悖的事}，却说\Quote{祂不能作}呢？这是要借此再次显示平等的不变与精确，因为这种说法并非归咎于祂软弱，反而显明祂的大能；正如\Person{保禄}在另一处论圣父时说：\BibleQuote{藉着两件不可更改的事，在这事上，天主不可能说谎}\BibleRef{希 6:18}；又说：\BibleQuote{如果我们否认祂——祂仍忠实，因为祂不能否认自己}\BibleRef{弟后 2:12}。实际上，\Quote{不可能}这个说法不是宣告软弱，而是能力——不可言喻的能力。祂的意思是这样的：那本性不容许这类事情。正如我们也说\Quote{天主不可能作恶}，并非归咎于祂任何软弱，而是承认祂那不可言喻的能力；同样，当祂说\BibleQuote{我由自己不能作什么}\BibleRef{若 5:30}时，祂的意思是：\Quote{不可能，本性不容许我作任何与圣父相悖的事}。为使你明白这确实是祂所说的，让我们看下文，看看\Person{基督}是否同意我们所说的，还是同意你们所说的。你们说，这句话否定了祂的能力和祂固有的权柄，显示祂的能力是软弱的；但我却说，这证明了祂的平等、祂与圣父不变的相似，以及一切都是凭同一旨意、能力和权柄完成的。那么，我们亲自询问基督，看祂随后所说的话是按你们的假设还是我们的假设来解释这些词句。祂究竟说了什么？\par

\BibleQuote{凡圣父所作的，圣子也照样作。}\par

你看见祂如何从根本上推翻了你们的论断，而证实了我们所说的吗？因为，如果基督不由自己作什么，那么圣父也不会由自己作什么，因为基督凡事都照样效法祂。若非如此，就会引出另一个荒谬的结论。因为祂没有说\Quote{凡祂看见圣父所作的，祂就作}，而是说\Quote{除非祂看见圣父作什么，祂就不作}，将祂的话扩展到一切时间；那么，按你们的说法，祂将不断地学习同样的事。你看见这思想何其崇高，甚至这表达上的极度谦卑也迫使最无耻和最不情愿的人避免卑下的想法，就是那些与祂尊位不相称的想法吗？因为谁会如此可怜可悲，竟主张圣子日复一日学习祂必须作的事呢？那\BibleQuote{祢是同样的，祢的岁月不会穷尽}\BibleRef{咏 102:27}，以及\BibleQuote{万物是藉着祂造的，凡受造的没有一样不是藉着祂造的}\BibleRef{若 1:3}，又怎能成立呢？如果圣父作了某些事，而圣子看见并模仿祂？你看见从前面所说的和后面所说的，都证明祂独立的权能吗？如果祂引用某些谦卑的说法，不必惊奇；因为既然他们听见祂崇高的言论后逼迫祂，视祂为天主的仇敌，所以祂在言辞上稍加降低，再次将论述引向更崇高的教义，然后又转向较低的，变换祂的教导，以便即使对不愿意听的人也易于接受。请看，在说了\Quote{我父工作，我也工作}，并宣称自己与天主平等之后，祂补充说：\Quote{圣子由自己不能作什么，除非看见圣父作。}然后又以更高的语调说：\Quote{凡圣父所作的，圣子也照样作。}然后又以较低的语调说：\par

% structure: p0030 source=h2 target=subsection reason=relative-heading-rank; base=h2

\subsection{若 5:20}

你看见这谦卑何等之大吗？这也合理；因为之前我说过的、将不断重复的，现在我再重申：当祂说任何低微或谦卑的话时，祂会过度使用，以便言辞的贫乏本身能说服甚至不愿意接受的人以虔敬的理解来领悟这些概念。因为若非如此，请看从词句本身推断所断言的内容何等荒谬。因为当祂说\BibleQuote{且要显示比这些更大的工程给祂看}时，祂将被发现尚未学会许多事，甚至不能这样论及宗徒们；因为他们一旦领受了圣神的恩宠，眨眼间就知道并能作一切需要知道和有能力作的事，而基督将被发现尚未学会许多祂需要知道的事。还有什么比这更荒谬的呢？\par

那么祂的意思是什么呢？因为祂曾治愈了瘫子，且将要复活死人，所以祂这样说，几乎等于说：\Quote{你们奇怪我治愈了瘫子吗？你们将看见比这些更大的事。}但祂并没有这样说，而是以某种更谦卑的方式继续，为要平息他们的疯狂。为使你明白\Quote{显示}不是绝对用语，请再听下文。\par

% structure: p0033 source=h2 target=subsection reason=relative-heading-rank; base=h2

\subsection{若 5:21}

然而，\BibleQuote{祂不能自己做什麼} 與 \BibleQuote{祂所愿意的人} 相對：因为如果祂赐生命给 \BibleQuote{祂所愿意的人}，祂就能做些事情 \BibleQuote{靠自己}，（因为 \BibleQuote{愿意} 表示能力，）但是，如果祂 \BibleQuote{不能自己做什麼}，那么祂就不能 \BibleQuote{赐生命给祂所愿意的人}。因为 \BibleQuote{如同圣父复活死人} 这句话显示能力上的不变相似，而 \BibleQuote{祂所愿意的人} 则显示权柄的平等。所以，你是否看到，\BibleQuote{不能自己做什麼} 是那一位的表达，不是否认祂自己的权柄，而是宣告祂的能力与旨意（与圣父的）的不变相似？按此意义也要理解 \BibleQuote{将显示给他} 这话；因为在另一处祂说：\BibleQuote{我要在末日使他复活}。\BibleRef{若 6:40} 并且，为显示祂这样做不是从高处接受内在能力，祂说：\BibleQuote{我是复活，我是生命}。\BibleRef{若 11:25} 为使你们不至于断言祂只复活祂所愿意的死人并赐生命给他们，而其他事却不如此行，祂预先防止了每一种此类异议，说道：\BibleQuote{凡圣父所做的，圣子也照样做}，从而宣告祂行圣父所行的一切事，并且如同圣父行它们；无论你说的是复活死人，还是塑造身体，还是赦免罪过，或任何其他事情，祂都像生祂的圣父一样工作。\par

但是，漠视自己救恩的人，毫不留意这些事；爱慕首位是何等大的邪恶。这一直是异端之母，这坚定了外邦人的不虔。因为天主愿意祂那不可见的事，借着这世界的受造得以被人领悟，\BibleRef{罗 1:20}，但他们离弃这些，拒绝由这种教导方式而来，为自己开辟了另一条路，因此被逐出真理。而犹太人不信，是因为他们彼此领受光荣，而不寻求那来自天主的光荣。但是，亲爱的诸位，让我们极其恳切地避免这疾病；因为即使我们有成千上万的美德，这虚荣的瘟疫也足以使它们全部化为乌有。\BibleRef{若 5:44} 因此，如果我们渴望赞美，就让我们寻求那来自天主的赞美，因为人的赞美，无论何种，一旦出现就消逝了，或者即使不消逝，也不能带给我们任何益处，而且常常来自腐坏的判断。那来自人的光荣有什么可羡慕的呢？那是年轻舞女、淫荡妇女、贪婪和强夺的人所享受的。但蒙天主认可的人，不是与这些人一同蒙认可，而是与那些圣人们，即先知和宗徒，他们彰显了天使般的生活。如果我们感到任何渴望，想要带领众多人跟随我们，或被他们注目，就让我们独自考虑这事，我们就会发现它完全毫无价值。最后，如果你喜爱人群，就吸引天使的军队到你身边，成为魔鬼所惧怕的，那么你就不会关心必朽的事物，反而会将一切辉煌之物踩在脚下，如同泥和土；并且会清楚看到，没有什么比追求光荣更使灵魂蒙羞；因为，渴望这荣光的人不可能过被钉在十字架上的生活，正如另一方面，把这种虚荣踩在脚下的人也不可能不制伏其他大部分情欲；因为那制伏了这虚荣的人，将胜过嫉妒、贪婪和一切严重的疾病。\Quote{怎样}，有人说，\Quote{我们才能胜过它？} 如果我们仰望那来自天上的另一种光荣，而这类虚荣竭力将我们从那光荣中驱逐出去。因为那天上的光荣既使我们在今生受尊敬，又随我们进入来世的生命，并解救我们脱离一切肉身的奴役——我们现在最悲惨地服事这奴役，完全将自己委身于大地和地上的事物。因为如果你进入市场，进入房屋，进入街道，进入兵营，进入旅馆、酒馆、船只、岛屿、宫殿、法庭、议会厅，你在各处都会发现对现世和此生之事的焦虑，每个人都在为这些事劳苦，无论去或来，旅行或居家，航海，耕种田地，在田野，在城市，总而言之，一切人。那么，当我们住在天主的土地上却不关心天主的事，当被命令要远离世俗之事时，我们却是天堂的陌生人和地上之公民，我们还有什么得救的希望呢？有什么比这麻木更糟糕呢？当我们每天听到审判和天国时，却效法诺厄时代的人和索多玛人，等待通过实际经历来学习一切？然而，所有那些事记载下来正是为此目的：若有人不信将要来临的事，他可以从已经发生的事，得到未来事情的确定证据。因此，考虑到这些事，无论是过去的还是未来的，让我们至少从这艰难的奴役中稍作喘息，也为我们自己的灵魂作些打算，好使我们获得现世和来世的福分；借着我们的主耶稣基督的恩宠和慈爱，愿光荣归于祂，偕同圣父和圣神，从今世直到永远，永无穷尽。阿们。\par

\SourceRecord{Translated by Charles Marriott.；Nicene and Post-Nicene Fathers, First Series；Vol. 14.；Edited by Philip Schaff.；Buffalo, NY: Christian Literature Publishing Co.,；1889.；<http://www.newadvent.org/fathers/240138.htm>.}

% structure: p0001 source=h1 target=section reason=page-title

\section{\WorkTitle{若望福音 第39篇讲道}}

% structure: p0002 source=h2 target=subsection reason=relative-heading-rank; base=h2

\subsection{若 5:23-24}

亲爱的，我们应当在一切事上极其谨慎，因为我们将要为自己的言行交账并接受严格的审查。我们的关注不止于现世，因为另一种生命状态将在今世之后接纳我们，我们将被带到可畏的审判台前。\BibleQuote{因为我们众人必须出现在基督的审判台前，好使每人藉着身体所行的事，或善或恶，都按他所行的受报。}\BibleRef{格后 5:10} 愿我们时常铭记这个审判台，以便能时刻持守德行；因为正如那从自己灵魂中驱逐了那日的人，就像脱缰之马冲向悬崖，（因为\BibleQuote{他的道路常是污秽的}\BibleRef{咏 9:5}），然后圣咏作者又加上原因：\Quote{他把你的审判从眼前远远抛开}；）所以，那常存此敬畏的人，必行事谨慎。\BibleQuote{你要记得，}有人说道，\BibleQuote{你的终末，你就永远不会犯罪。}\BibleRef{德 7:40} 因为那现在赦免我们罪过的主，将来却在审判中犯罪；那为我们而死的主，将来要再次显现以审判全人类。宗徒说：\BibleQuote{凡是期待祂的，}\BibleQuote{祂将第二次显现，与罪无关，而为救恩。}\BibleRef{希 9:28} 因此，在此处祂也说：\Quote{父不审判任何人，但将一切审判权交给了子，为使众人都尊敬子，如同尊敬父一样。}\par

有人问：\Quote{那么，我们也要称祂为父吗？} 断乎不可！祂使用\Quote{子}这个词，是要我们尊敬祂，祂仍然是子，如同我们尊敬父一样；但称祂为\Quote{父}的人，并不像尊敬父那样尊敬子，反而混淆了一切。再者，人往往因惩罚而非因恩惠而被吸引，因此祂说了如此严厉的话，甚至藉着恐惧吸引他们尊敬祂。当祂说\Quote{一切}时，意思是祂有权惩罚和尊荣，并随己意而行。使用\Quote{给予}一词，是免得你以为祂不是受生的，因而以为有两个父。因为凡是父所有的，子也都有，祂是受生的，且仍然是子。并且，为使你明白\Quote{给予}等同于\Quote{生了}，请听另一处所说：\BibleQuote{就如父在自己内有生命，同样也给了子在自己内有生命。}\BibleRef{若 5:26} 有人问：\Quote{那么，祂是先生了子，然后才赐予祂生命吗？因为给予的对象已是存在之物。难道祂生下来时是没有生命的吗？} 连魔鬼也不会如此想象，因为这是极其愚昧且不敬的。所以，正如\Quote{赐予生命}就是\Quote{生了祂，祂就是生命}，同样，\Quote{赐予审判}就是\Quote{生了祂，祂将成为审判者。}\par

为使你听到祂以父为根源，不臆想本质的差异或尊荣的逊色，祂来审判你，藉此证明祂的平等。因为那有权随心所欲惩罚和尊荣的，与父有同等的权能。既然不是这样，如果祂是在受生之后才得到尊荣，那么祂后来是如何受到尊荣的？通过何种进步方式达到了如此高度，以致接受并被任命于这种尊位？你们难道不羞耻如此无耻地将这些属肉体的、卑劣的想象应用到那不接纳任何增加的纯洁本性上吗？\par

有人问：\Quote{那么，基督为何如此说话？} 为使祂的话被欣然接受，并为崇高的言论铺路；因此祂将这些与那些混合，将那些与这些混合。请看祂如何做；因为从一开始就看出这一点是好的。祂说：\BibleQuote{我父工作，我也工作，}\BibleRef{若 5:17} 藉此宣告他们的平等和同等的尊荣。但他们\Quote{想要杀祂。}祂又怎样呢？祂确实降低了讲话的方式，并用同样的含义说：\Quote{子不能自己做任何事。}然后祂又将讲论提高到高处，说：\Quote{凡父所做的，子也同样做。}然后又回到低处：\Quote{因为父爱子，并将自己所作的一切显示给祂，还要将比这些更大的事显示给祂。}然后又升高：\Quote{就如父复活死者并赐予他们生命，同样子也随己意赐生命。}此后祂又将高低结合：\Quote{因为父不审判任何人，而将一切审判交给了子，}然后又升高：\Quote{为使众人都尊敬子，如同尊敬父一样。}你看见祂如何变化讲论，用高低不同的词语和表达编织起来，使当时的人可以接受，并使后来的人不受损害，从高处部分获得对其余部分的正确认识吗？因为如果不是这样，如果这些话不是出于迁就，那么为何要加上那些崇高的表达？因为一个有权谈论自己伟大之事的人，当他说一些卑微的话时，他有合理的借口，即他这样做是出于某种策略；但一个应当卑微自谦的人若说出伟大的话，他凭什么说出超越自己本性的话呢？这完全没有理由，而是极度的不敬。\par

2. 因此，我们能够为这些谦卑的言辞找到一个理由，一个充分且合宜于天主的理由，即祂的屈尊俯就，祂教导我们要节制，以及由此为我们成就的救恩。为说明这一点，祂自己在另一处说：\Quote{我说这些话，是为叫你们得救。}因为当祂放下自己的见证，转向\Person{若翰}的见证时（这是一件与祂的尊贵不相称的事），祂给出了如此谦卑言语的理由，并说：\Quote{我说这些话，是为叫你们得救。}而你们这些断言祂与生祂的那一位没有同等权威和能力的人，当你们听到祂说出那些宣告祂与圣父同等的权威、能力和光荣的话语时，你们能说什么呢？那么，如果祂像你们所断言的那样非常低等，祂又如何要求同等的尊荣呢？祂甚至不止于此，还继续说：\par

\Quote{不尊敬圣子的，就是不尊敬派遣祂来的圣父。}你们看，圣子的尊荣如何与圣父相连？\Quote{那有什么了不起？}有人说。\Quote{我们在宗徒身上也看到同样的事；基督说：\InnerQuote{谁接待你们，就是接待我。}}\BibleRef{玛 10:40} 但在那处，祂这样说，是因为祂将仆人的事当作自己的事；而在这里，却是因为性体和光荣是同一的（与圣父）。因此，关于宗徒并没有说\Quote{他们应尊敬}，而祂说得对：\Quote{不尊敬圣子的，就是不尊敬圣父。}因为如果有两个国王，一个受侮辱，另一个也受侮辱，尤其当受侮辱的是儿子时，更是如此。即使他的一个士兵被虐待，他也受侮辱；但方式不同，那是在别人身上，而这里却是在自己身上。所以祂预先说：\Quote{他们应尊敬圣子，如同尊敬圣父一样}，这样当祂说\Quote{不尊敬圣子的，就是不尊敬圣父}时，你们就能明白这尊敬是相同的。因为祂不单单说\Quote{不尊敬圣子的}，而是说\Quote{不按我所说的那样尊敬祂的}\Quote{就是不尊敬圣父。}\par

\Quote{那派遣者和那被派遣者怎么可能是同一性体？}有人说。你们又落到属血肉的推论中，并且不明白这一切话除了让我们知道祂是原因，不要陷入\OriginalTerm{撒伯流}{Sabellius}的错误，并以此医治犹太人的软弱，使祂不被视为天主的敌人之外，别无目的；因为他们曾说：\Quote{这个人不是出于天主}\BibleRef{若 9:16}，\Quote{这个人不是从天主来的。}现在，为消除这种怀疑，高深的言论不如谦卑的言语贡献大，因此祂不断到处说祂是\Quote{被派遣的}；不是要你们以为那表达是对祂伟大的任何贬低，而是为堵住他们的口。也是为此原因，祂不断转向父，并提及自己高贵的出身。因为如果祂完全按照祂的尊贵来说话，犹太人就不会接受祂的话，因为仅仅几句这样的表达，他们就迫害祂，并且多次用石头砸祂；而如果完全针对他们，祂只使用谦卑的表达，将来许多人可能会受害。所以祂混合并融合祂的教导，既用这些谦卑的言语堵住犹太人的口，如我所说，也用与祂尊贵相称的表达，从有识之士心中消除任何对祂所说的话的低贱看法，并证明这样的看法根本不适用于祂。\par

\Quote{被派遣}这个表达表示位置的改变——但天主是无所不在的。那么祂为什么说祂是\Quote{被派遣}的呢？祂是以属人的方式说话，表明祂与父同心。至少祂塑造了随后的言语，渴望实现这一点。\par

% structure: p0011 source=h2 target=subsection reason=relative-heading-rank; base=h2

\subsection{若望福音 5:24}

你们看见祂如何不断提出同一件事，为医治那种怀疑感，既在此处，也在后续通过恐惧和祝福的应许消除他们对祂的嫉妒，然后又再次在言语上大大屈尊吗？因为祂没有说：\Quote{那听我的话并相信我的}，因为他们必定会认为那是骄傲，是多余的辞藻；因为，如果经过很长时间和无数奇迹之后，当祂以这种方式说话时，他们还怀疑，那时他们会更加如此。正是为此，在后来的时期他们才对祂说：\Quote{亚巴郎死了，先知们也死了，祢怎么说：若有人遵守我的话，他永远尝不到死味呢？}\BibleRef{若 8:52} 因此，为避免他们在这里也发怒，请看祂说什么：\Quote{那听我的话，并相信派遣我来者的，就有永生。}这大大有助于使祂的讲论被接受，当他们得知，那些听祂的人也会相信父；因为他们在欣然接受这一点后，就会更容易接受其余部分。因此，正是这种谦卑的说话方式有助于并引向更高的事；因为说过\Quote{有永生}之后，祂接着又说：\par

\Quote{并且不受审判，而是已从死亡进入生命。}\par

通过这两点，祂使祂的讲论被接受：第一，因为是信靠父；第二，因为信徒享受许多祝福。而\Quote{不受审判}的意思是\Quote{不受惩罚}，因为祂所说的死亡不是\Quote{这里的}死亡，而是永远的死亡，同样，那\Quote{生命}也是不朽的生命。\par

% structure: p0015 source=h2 target=subsection reason=relative-heading-rank; base=h2

\subsection{若望福音 5:25}

说了这些话，祂也谈到了以行为作证明。因为当祂说：\BibleQuote{就如圣父唤起死者并使他们生，同样圣子也随意使他人。}为了不让这事被认为是空夸和骄傲，祂以事迹提供证明，说：\BibleQuote{时候将到。}然后，使你不要认为时间漫长，祂加上：\BibleQuote{而且现在就是，死者要听见天主子的声音，听见的人就活了。}你看到这里祂绝对而不可言喻的权柄了吗？因为如同在复活时一样，祂说，现在也是这样。此外，当我们听见祂命令我们的声音时，我们就复活了；因为宗徒说：\BibleQuote{在天主的命令下，死者将复活。}\Quote{也许有人会问：从何而见这些话并非空夸？}从祂所加的\BibleQuote{现在就是}；因为如果祂的应许只指向未来的某个时候，他们就会怀疑祂的话，但现在祂给他们提供了一个证明：\Quote{当我还在你们中间时，这事就要发生。}如果祂没有权能，祂就不会应许那个时候，免得因应许而招致更大的嘲笑。然后祂又加上一个证明祂论断的论证，说：\par

% structure: p0017 source=h2 target=subsection reason=relative-heading-rank; base=h2

\subsection{若 5:26}

3. 你看到这宣告了完全的相似，除了一点：一位是圣父，另一位是圣子？因为\Quote{赐给}这个说法仅仅引入了这个区别，但宣告其余的都平等且完全相同。由此可知，圣子像圣父一样有权柄和能力行一切事，祂的权力并非来自别处，因为祂\Quote{有生命}如同圣父一样。因此，紧接着就加上了下一句，使我们由此也能理解另一句。这是什么？\par

% structure: p0019 source=h2 target=subsection reason=relative-heading-rank; base=h2

\subsection{若 5:27}

为何祂不断提到\Quote{复活}和\Quote{审判}？因为祂说：\BibleQuote{就如圣父唤起死者并使他们生，同样圣子也随意使他人}；又说：\BibleQuote{圣父不审判任何人，但将一切审判交给了圣子}；又说：\BibleQuote{如圣父在自己内有生命，同样也赐给圣子在自己内有生命}；又说：\BibleQuote{听见天主子声音的人将活}；而这里又说：\BibleQuote{赐给祂施行审判的权柄。}为何祂不断说到这些？我说的是\Quote{审判}、\Quote{生命}和\Quote{复活}？因为这些主题最能吸引即使倔强的听者。因为一个相信自己将复活并向基督交账的人，即使没有看到其他征兆，但一旦承认了这一点，必定会奔向祂以平息祂的审判。\par

% structure: p0021 source=h2 target=subsection reason=relative-heading-rank; base=h2

\subsection{若 5:28}

\Person{撒摩沙塔的保禄}不是这样解释；他是怎样解释的呢？\Quote{赐给祂施行审判的权柄，\Emph{因为}祂是人子。}如此阅读的经文结果不一致，因为祂得到审判不是\Quote{因为}祂是人（否则，什么能阻止所有人都成为法官？），而是因为祂是那不可言说之本质的圣子，所以祂是审判者。因此我们必须读作：\Quote{祂是人子，你们不要惊奇这事。}因为当祂所说的话在听者看来不一致，他们认为祂不过是一个凡人，而祂的话却超越凡人甚至天使，唯独属于天主时，为了解决这个异议，祂加上：\par

% structure: p0023 source=h2 target=subsection reason=relative-heading-rank; base=h2

\subsection{若 5:28-29}

为何祂不说：\Quote{不要惊奇祂是人子，因为祂也是天主子}，反而提到\Quote{复活}？祂确实在前面提到了这一点，说：\BibleQuote{将要听见天主子的声音。}如果这里祂对此保持沉默，不必惊奇；因为在提及一个专属天主的工作后，祂让听者从中推导出祂是天主，是天主子。因为如果祂不断自己断言这点，当时会触犯他们，但通过奇迹的论证证明，祂的教义就不那么沉重。所以那些构造三段论的人，在提出前提并充分证明问题后，常常不自己下结论，而是为了使听者更公平地接受，并使自己的胜利更明显，让对手自己宣告判决，这样当对手做出有利于他们的决定时，旁观者会更同意他们。因此，当祂提到\Person{拉匝禄}复活时，祂没有谈到审判（因为拉匝禄不是为了这个复活）；但当祂泛泛而谈时，祂也加上：\BibleQuote{行善的要去复活得生命，作恶的要去复活受审判。}同样，\Person{若望}也通过谈论审判来引导他的听者，说：\BibleQuote{不信从圣子的，不得见生命，天主的义怒常在他身上}（\BibleRef{若 3:36}）。同样，祂也引导\Person{尼苛德摩}：\BibleQuote{信从圣子的，不被审判；不信的，已被审判了}（\BibleRef{若 3:18}）。所以这里祂提到审判台和恶行后的惩罚。因为祂前面说过：\BibleQuote{听我的话而信那派遣我来的，不被审判}，免得有人以为这足以得救，祂还增加了人生活的结果，宣称：\BibleQuote{行善的要去复活得生命，作恶的要去复活受审判。}既然祂说过全世界都要向祂交账，所有人都因祂的声音而复活，这是新奇而陌生的事，甚至现在许多似乎相信的人也不信，更不用说当时的犹太人了，请听祂如何证明这一点，再次迁就听者的软弱。\par

% structure: p0025 source=h2 target=subsection reason=relative-heading-rank; base=h2

\subsection{若 5:30}

虽然祂刚刚以治愈瘫子证明了复活（这并非微不足道的证据），因此祂也在做了几乎等于复活的事之前并未提及复活。祂在使身体康复后暗示了审判，说：\BibleQuote{看，你已经痊愈了，不要再犯罪，免得你遭遇更糟的事}；然而祂仍预先宣告了拉匝禄的复活和世界的复活。当祂谈及这两者时——拉匝禄的复活几乎即刻发生，而全世界的复活则在遥远的将来——祂以瘫子以及时间的临近来证实前者，说：\BibleQuote{时候已到，并且就是现在}；又以拉匝禄的复活来证实后者，以已经发生的事使尚未发生的事呈现在他们眼前。我们处处可见祂这样做，提出两个或三个预言，并总是以过去的事证实将来的事。\par

4. 然而，在说了和做了这么多之后，因他们仍然十分软弱，祂并不满足，而是用其他言辞平息他们好争辩的性情，说：\BibleQuote{我由我自己什么也不能作；我所判断的，是由于我所听见的，而我的判断是正义的，因为我不寻求我的意愿，而寻求那派遣我者的意愿。} 因为祂似乎说出了一些与先知们不同的奇怪主张（因为先知们说审判全地，即人类，的是天主；达味处处高声宣扬这一真理：\BibleQuote{祂要以正义审判万民}，和\BibleQuote{天主是正义的审判者，坚强而忍耐} \BibleRef{咏 95:10}；所有先知和梅瑟也是如此；但基督说：\BibleQuote{圣父不审判任何人，但将一切审判交给了圣子}；这话足以使听到的犹太人困惑，进而怀疑基督是天主的仇敌），祂在这里极大程度地俯就他们的言语，并按照他们软弱的程度，为要根除这种有害的意见，说：\BibleQuote{我由我自己什么也不能作}；也就是说：\Quote{你们不会看见我做出或说出任何离奇、不同或圣父所不愿的事。} 而且祂先前已宣布祂是\BibleQuote{人子}，而因他们当时以为祂是一个人，所以祂在此也这样表达。正如先前祂说：\BibleQuote{我们说我认识的，我们见过的事}；以及若翰说：\BibleQuote{祂所见过的事，祂就作证，却没有人接受祂的见证} \BibleRef{若 3:32}；这两种说法都是关于精确的知识，不仅仅是关于听见和看见；所以在此处当祂说到\Quote{听见}时，无非是宣告祂不可能愿意任何事情，除非圣父所愿意的。然而祂没有说得如此明确（因为他们还不会立即接受这样断言的话）；那怎么说呢？以非常俯就并适合一个普通人的方式：\BibleQuote{我所判断的，是由于我所听见的}。祂在此处再次使用这些词，不是指\Quote{教导}（因为祂没有说\Quote{如我被教导}，而是\Quote{如我听见}）；也不是像祂需要聆听（因为祂不仅不需要被教导，甚至不需要聆听）；而是为了宣告祂与圣父决定的一致与同一，就好像祂说：\Quote{我如此判断，如同圣父自己判断一样。} 接着祂又说：\BibleQuote{并且我知道我的判断是正义的，因为我不寻求我的意愿，而寻求那派遣我者的意愿。} 你说什么呢？祢有一个与圣父不同的意愿吗？然而在另一处祂说：\BibleQuote{我与祢是一}（这里说的是意愿与同心），\BibleQuote{愿他们也合而为一在我们内} \BibleRef{若 17:21}；也就是说：\Quote{在关于我们的信仰中。} 你看见那些看似最谦卑的言语其实隐藏着高深的含义吗？因为祂所暗示的是：并非圣父的意愿是一，祂自己的是另一个；而是：\Quote{如同在一个心智内有一个意志，我的意愿和圣父的意愿也是如此。}\par

不要惊讶祂宣称如此紧密的结合；因为\Person{保禄}也曾用这个比喻论及圣神：\Quote{什么人知道人的事，除非是那人内里的心神？同样，天主的事，除了天主圣神外，没有人知道。} 所以基督的意思不外乎是：\Quote{我没有一个与圣父的旨意不同且分离的旨意；若祂愿意什么，那么我也愿意；若我愿意，那么祂也愿意。因此，既然无人能反对圣父审判，同样也无人能反对我，因为每一个的判决都出自同一个心志。} 如果祂说这些话更像是作为一个人，不要惊讶，因为他们仍把祂视为一个凡人。因此，对于这类经文，不仅需要探究字句的意义，还需要考虑听者的猜疑，并视这些话为针对那猜疑而说的。否则，许多困难就会随之而来。例如，祂说：\Quote{我不寻求我自己的旨意}；按此，祂的旨意就是不同的（与圣父），是不完美的，不，不仅不完美，甚至是有害的。\Quote{因为如果它是有益的，如果它符合圣父的旨意，你为什么不寻求它？} 人们有理由这样说，因为他们有许多与圣父所喜悦的相反的旨意；但祢，祢在一切事上都像圣父，为何说这话？因为连一个成全并被钉十字架的\Quote{人}也不会说这话。如果\Person{保禄}如此将自己与天主的旨意相融合，以至说：\Quote{我活着，但不再是我，而是基督在我内生活} \BibleRef{迦 2:20}，那么万有之主怎么说：\Quote{我不寻求我自己的旨意，而是那派遣我来的圣父的旨意}，好像那旨意是不同的？那么祂的意思是什么？祂将祂的言论应用，好像情况是一个凡人，并使祂的语言适应听者的猜疑。因为当祂通过之前所说的证明了自己的话，部分像天主，部分像凡人，祂又作为一个人努力确立同样的事，并说：\Quote{我的审判是正义的。} 从哪里看出呢？\Quote{因为我不寻求我自己的旨意，而是那派遣我来的圣父的旨意。} \Quote{因为如同在人当中，那摆脱自私的人不能被公正地指控作出了不公正的判决，同样，你们现在也不能指控我。那想要建立自己事的人，或许会被许多人怀疑有意败坏公正；但那不看自己事的人，有什么理由不公正地判决呢？现在将这推理应用到我的情况。若我说我不是由圣父派遣的，若我没有将所行事的荣耀归于祂，你们中有些人或许会怀疑，我为了自己得荣耀而说了不真实的话；但如果我将所做的事归功于别人，你们从何处、因何故还有理由怀疑我的话？} 你看见祂如何证实自己的话，并断言\Quote{祂的审判是正义的}，以任何一个普通人可能用来为自己辩护的论证？你看见我常说的话是清楚可见的吗？那是什么？那就是，这些表达的极度谦卑最能说服明智的人，不是立即接受这些话然后（陷于低俗思想），而是努力使它们达到其意义的高度；这谦卑也轻易地抬起了那些曾匍匐在地上的人。\par

现在，牢记这一切，我劝你们，不要漫不经心地忽视基督的话，而要仔细探究这一切，处处思考所说之事的理由；不要以为无知和单纯就足以成为我们的借口，因为祂命令我们不仅要\BibleQuote{纯朴}，还要\BibleQuote{机警}。\BibleRef{玛 10:16}因此，让我们在教义和生活的正当行为上，以纯朴实践智慧；让我们在此审判自己，免得将来与世人一同被定罪；让我们对待同仆，如同我们希望主对待我们一样：因为我们说：\BibleQuote{宽免我们的罪债，犹如我们也宽免得罪我们的人。}\BibleRef{玛 6:12}我知道受打击的灵魂不会温顺地忍受，但如果我们考虑到这样做是善待自己而非那惹恼我们的人，我们很快就会放下愤怒的毒液；因为那没有宽恕得罪他的同伴一百个德纳的人，并非亏负了他的同伴，而是亏负了自己，使自己再度对先前已被宽免的一万塔冷通负有责任。\BibleRef{玛 18:30}因此，当我们不宽恕别人时，我们就不宽恕自己。所以，我们不仅要对天主说：\Quote{不要记念我们的过犯}，每个人也要对自己说：\Quote{让我们不要记念同仆对我们的过犯}。因为你先对自己的罪恶下了判决，然后天主才审判；你制定了关于宽恕和惩罚的法律，你宣布了在这些事上的决定，因此天主是否记念，取决于你。为此，保禄劝诫我们：\BibleQuote{如果某人与某人之间有怨尤，要彼此宽恕}\BibleRef{哥 3:13}，不仅宽恕，而且不留任何痕迹。因为基督不仅没有宣扬我们的过犯，也没有使我们这些犯罪者想起它们，也没有说：\Quote{你在这样那样的事上犯罪了}，而是豁免并涂抹了欠据，不算我们的过犯，正如保禄所宣告的。\BibleRef{哥 2:14}让我们也这样做；从我们的心中擦去所有对我们的冒犯；如果那得罪我们的人曾对我们做过任何好事，我们只记念那好事；但若有什么令人痛苦和难以忍受的事，我们就把它丢出去并涂抹掉，不留一丝痕迹。如果他没有对我们做过任何好事，那么当我们宽恕时，我们所获得的赏报和功德就更大。其他人通过守夜、睡在地上、千万种刻苦来洗去他们的罪恶，而你的方式更容易，就是不记念冤屈，就能使你的所有过犯消失。那么，你为什么像疯狂和狂暴的人那样，把剑刺向自己，把自己从未来的生命中放逐，而你应该用尽一切方法去获得它？因为如果现世生命如此可爱，那未来的生命——痛苦、悲伤和哀叹都已逃离——又该如何说呢？在那里不需要害怕死亡，也不要想像那些美好事物会有尽头。有福的，三倍有福的，甚至更多倍有福的，是那些享受那有福安息的人；而可怜的，三倍可怜的，甚至万倍可怜的，是那些把自己从那福乐中驱逐出去的人。有人说：\Quote{那么，是什么使我们享受那生命呢？}听审判者亲自与一个年轻人谈论此事。当那年轻人说：\BibleQuote{我该做什么才能获得永生？}\BibleRef{玛 19:16}基督在向他重复了其他诫命之后，以爱近人结束。也许像那富翁一样，我的听众中有些人会说：\Quote{我们也遵守了这些，因为我们没有抢劫、没有杀人、没有奸淫}；然而你肯定不能说，你爱了近人如你所应爱的。因为如果有人嫉妒或诽谤了他人，如果有人没有帮助受伤害的人，或者没有分施他的财物，那么他就没有爱那人。现在基督不仅命令这个，还有别的。那是什么呢？他说：\BibleQuote{你去变卖你所有的，施舍给穷人，然后来跟随我}\BibleRef{玛 19:21}，称在行动上效法祂为\Quote{跟随}祂。我们从中学习到什么？首先，没有这一切的人不能在那安息中达到首要的地位。因为年轻人说：\BibleQuote{这一切我都遵守了}，基督好像为了使他完全蒙悦纳还缺一件大事，就回答说：\BibleQuote{你若愿意成为成全的，去变卖你所有的，施舍给穷人，然后来跟随我。}首先我们可学到这一点；其次，基督责备这人虚妄的自夸；因为一个人生活在如此富裕中，却不关心其他穷困的人，他怎能爱近人呢？因此连在这件事上他也没有说实话。但让我们做这两件事：让我们渴望倒空我们的财物，以购买天堂。因为如果人们为了世俗的荣誉常常耗尽全部财产，这种荣誉本应停留在此世，甚至在此世也不能长久伴随我们（因为很多人甚至在死前就被剥夺了高位，其他人因此而丧命，然而尽管知道这些，他们仍为它花费一切）；如果他们为这种荣誉做了这么多，那么我们若为了那持久且不能被夺去的荣誉而不愿放弃一点点，也不愿供给别人那些我们不久就会留在此世的东西，那还有什么比我们更可怜的呢？这该是多么疯狂啊，当我们有能力自愿给予别人，从而带走那些我们即使不情愿也会被剥夺的东西，却拒绝这样做？如果一个人被领向死亡，而给他机会放弃一切财产以获得自由，我们会认为他蒙受了恩惠；而我们这些被领向深渊的人，当我们被允许放弃一半就能自由时，却宁愿受罚，无益地保留不属于我们的东西，以致连属于我们的也失去？我们还有什么借口，什么求恕的理据？我们有如此容易的道路通向生命，却冲下悬崖，行走在无益的道路上，剥夺自己今世和来世的一切，而我们本可安全地享受两者？如果我们以前没有这样做，至少现在停下来；恢复理智，妥善处理当下的事，以便我们容易获得将来的事，借着我们的主耶稣基督的恩宠和慈爱，愿光荣归于祂和圣父及圣神，至于无穷之世。阿们。\par

\SourceRecord{Translated by Charles Marriott.；Nicene and Post-Nicene Fathers, First Series；Vol. 14.；Edited by Philip Schaff.；Buffalo, NY: Christian Literature Publishing Co.,；1889.；<http://www.newadvent.org/fathers/240139.htm>.}

% structure: p0001 source=h1 target=section reason=page-title

\section{\WorkTitle{若望福音讲道集} 第四十篇}

% structure: p0002 source=h2 target=subsection reason=relative-heading-rank; base=h2

\subsection{\BibleBook{若望}{福音} 5:31-32}

1. 若有人不谙技艺而试图开矿，他必得不到金子，反倒胡乱混杂一切，徒劳无益且有害：同样，那些不理解圣经的方法，不探究其特点和规律，而是粗心大意、千篇一律地通读经书的人，会将金子与泥土混在一起，永远发现不了其中隐藏的宝藏。我这样说，是因为眼前的这段经文含有许多金子，并非显而易见，而是被重重晦暗所覆盖，因此我们必须通过挖掘和提炼，才能获得正确的意义。因为谁不会一听到基督说\BibleQuote{我若为自己作证，我的证词便不真实}就感到困扰呢？祂似乎常常为自己作证。例如，祂与撒玛黎雅妇人谈话时说\Quote{我是那与你说话的}；同样，对盲人说\BibleQuote{是祂在与你讲话}\BibleRef{若 9:37}；并斥责犹太人说\BibleQuote{你们说，我因说了\InnerQuote{我是天主子}而亵渎}\BibleRef{若 10:36}。此外，在许多其他地方祂也如此做。如果所有这些声明都是假的，我们还有什么得救的希望？当真理本身宣称\BibleQuote{我的证词不真实}时，我们到哪里去找真理呢？这似乎还不是唯一的矛盾；还有另一个同样突出的矛盾。祂接着说\BibleQuote{虽然我为自己作证，但我的证词是真实的}\BibleRef{若 8:14}。那么，请告诉我，我该接受哪一个，又认为哪一个虚假呢？如果我们就这样简单地从上下文摘取这些说法，而不仔细考虑说话的对象、原因以及其他类似情况，那么它们都将成为虚假。因为如果祂的证词\BibleQuote{不真实}，那么这句断言也不真实——不仅是第二句，第一句也是如此。我们需要极大的警觉，更确切地说，需要天主的恩宠，以免停留在字面上；因为异端者正是如此犯错，他们不探究说话者的意图和听者的心态。如果我们在这些之外不加上时间、地点以及听众的看法等其他要点，就会有许多荒谬的后果。\par

那么究竟是什么意思呢？犹太人即将反驳祂说\BibleQuote{你如果为自己作证，你的证词便不真实}\BibleRef{若 8:13}。因此祂预先说了这些话；似乎是说\Quote{你们一定会对我说：我们不相信你，因为为自己作证的人在人群中不容易被信任。}所以\Quote{不真实}这句话不应绝对理解，而是指向他们的猜疑，似乎是说\Quote{对你们而言，它不真实}。祂说这话并非着眼于自己的尊威，而是指向他们内心的隐秘想法。当祂说\BibleQuote{我的证词不真实}时，祂是在斥责他们对祂的看法，以及他们即将提出的反对意见；但当祂说\BibleQuote{虽然我为自己作证，但我的证词是真实的}\BibleRef{若 8:14}时，祂声明了事情本身的性质，即作为天主，他们理当相信祂，即使祂为自己说话。因为祂曾谈论过死人的复活、审判，以及那信祂的不被定罪，反而进入生命；祂将坐在审判位上向所有人索账，祂拥有与圣父同等的权柄和能力。由于祂将要再次用其他方式证明这些事，所以祂先提出他们的反对意见。祂说\BibleQuote{我告诉过你们：\InnerQuote{正如圣父使死人复活并赐予生命，同样，圣子也随意赐予生命}；我告诉过你们：\InnerQuote{圣父不审判任何人，而是将一切审判交给了圣子}；我告诉过你们：人必须\InnerQuote{尊敬圣子如同尊敬圣父一样}；我告诉过你们：\InnerQuote{那不尊敬圣子的，也不尊敬圣父}；我告诉过你们：\InnerQuote{那听到我的话并相信的人，不会见到死亡，而已经出死入生}\BibleRef{若 5:24}；我的声音将使死人复活，一些现在，一些将来；我将向所有人追讨他们的过犯，我将公义地审判，并报答那些正直行走的人。}现在，既然所有这些都只是断言，所断言的事又重要，并且尚未向犹太人提供清楚证明，只有一个模糊的暗示，所以祂在将要确立自己的断言之前，先提出他们的反对意见，说了一些这样的话（即使不是完全相同的字句）\Quote{也许你会说：你断言这一切，但你并不是可信的证人，因为你在为自己作证。}因此，祂首先通过陈述他们会说的话来压制他们好辩的精神，表明祂知道他们内心的秘密，并以此作为祂能力的第一证明。提出反对意见后，祂提供了其他清楚且无可争辩的证明，为祂的话提出了三个证人：祂所行的作为、圣父的见证，以及若翰的宣讲。而祂首先提出若翰这个不那么重要的证人。因为祂说\BibleQuote{另有为我作证的一位，我知道他的证词是真实的}，随后祂加上：\par

% structure: p0005 source=h2 target=subsection reason=relative-heading-rank; base=h2

\subsection{\BibleBook{若望}{福音} 5:33}

然而，如果祢的证词不真实，祢怎能说\BibleQuote{我知道若翰的见证是真的，他已为真理作证}呢？你（人啊）岂不看到，由此清楚可知，\BibleQuote{我的证词不真实}这句话，是对他们内心的隐秘想法说的？\par

2. \Quote{那么，}有人说，\Quote{如果若翰的见证有所偏颇。}为避免犹太人可能这样断言，请看祂如何消除这种疑虑。因为祂并没有说：\Quote{若翰为我作证，}而是说：\Quote{你们曾差人到若翰那里去，你们若不是认为他可信，就不会差人去。}不仅如此，他们差人去不是问关于基督的事，而是问他本人，而他们既在关于他自己的事上认为他可信，就更会在关于别人的事上认为他可信。因为，可以这样说，我们所有人的天性都是不太相信那些谈论自己的人，而更相信那些谈论别人的人；然而他们却认为他如此可信，以至于连关于他自己的事也不需要别的见证。因为被差去的人并没有说：\Quote{你关于基督说什么？}而是说：\Quote{你是谁？你关于你自己说什么？}他们对这个人怀有如此巨大的钦佩。现在基督提到这一切，说：\Quote{你们曾差人到若翰那里。}为此，福音作者不仅记述了他们去差人，而且还确切地说明了被差的人是司祭和法利塞人，不是普通或卑贱的人，也不是可能被贿赂或欺骗的人，而是能准确理解他所说的话的人。\par

% structure: p0008 source=h2 target=subsection reason=relative-heading-rank; base=h2

\subsection{若 5:34}

\Quote{那么，你为什么引述若翰的见证？} 他的见证不是\Quote{人的见证，}因为，他说：\Quote{那派遣我用水施洗的，祂对我说：} \BibleRef{若 1:33} 所以若翰的见证是天主的见证；因为他从祂那里学习，才说出他所说的。但为了没有谁会问：\Quote{从哪里显明他从天主学习了？} 并停在这里，祂仍针对他们的思想说话，大量地使他们哑口无言。因为许多人未必知道这些事；他们迄今一直听从若翰，如同听从一位谈论自己的人，因此基督说：\Quote{我不接受人的见证。} 而且为免犹太人问：\Quote{如果你不接受人的见证，并不藉此加强自己，为什么你又引述这个人的见证？} 请看祂如何藉着接下来的话纠正这个矛盾。因为说了\Quote{我不接受人的见证}之后，祂又加上了：\par

\Quote{但我这些话，是为叫你们得救。}\par

祂所说的是这样的：\Quote{我，作为天主，并不需要若翰那属于人的见证，然而因为你们更留意他，认为他比任何人都可信，跑向他如同一位先知（因为全城的人都涌向约旦河），并且即使我行了奇迹，你们仍不相信我，因此我提醒你们关于他的那个见证。}\par

% structure: p0012 source=h2 target=subsection reason=relative-heading-rank; base=h2

\subsection{若 5:35}

为避免他们反驳：\Quote{如果他真的说了话而我们没有领受他，那又如何？}祂表明他们确实接受了若翰的话：因为他们派去的不是普通人，而是司祭和法利塞人，并且愿意欢乐；他们如此钦佩这个人，同时对他的话无话可说。但\Quote{一时}这个说法，是祂注意到他们的轻浮，以及他们很快离开了他这一事实。\par

% structure: p0014 source=h2 target=subsection reason=relative-heading-rank; base=h2

\subsection{若 5:36}

\Quote{因为如果你们愿意按照事实的（自然）结果接受信仰，我本可以藉我的作为比藉他的言语更使你们信服。但既然你们不愿意，我就引你们到若翰那里，不是因为我需要他的见证，而是因为我做这一切都是\InnerQuote{为使你们得救}。因为我比若翰有更大的见证，即来自我的作为的见证；但我不仅仅考虑如何藉可靠的证据使你们接受我，还考虑如何藉着你们所认识和钦佩的人（使你们接受）。}然后祂看了他们一眼，说他们在若翰的光中一时欢乐，祂宣布他们的热心只是暂时的、不确定的。\par

祂称若翰为一盏灯，表明他自己没有光，而是因圣神的恩宠；但导致祂与若翰绝对区别的情况，即祂是义德的太阳，祂尚未提出；但仅仅暗示这一点，祂就尖锐地触动了他们，表明从使他们轻视若翰的同一性情，他们也不能相信基督。因为他们只是暂时钦佩那个他们确实钦佩的人，而这个人，如果他们不如此行，本来很快就会牵着他们的手到耶稣那里。既然证明了他们完全不配得宽恕，祂继续说到：\Quote{我有比若翰更大的见证。}\Quote{是什么？}就是来自祂作为的见证。\par

\Quote{因为那些事工，}祂说，\Quote{圣父交给我要我完成的，就是我正在做的这些事工，为我作证：是圣父派遣了我。}\par

藉此祂提醒他们已痊愈的瘫子，以及许多其他事。也许他们中有人可能断言这些话只是空谈，并且是因为若翰对祂的友谊而说的（虽然即使对若翰——一个完全实践智慧的人，并因此受他们钦佩的人——他们也无权这样说），但是那些事工即使在他们中最狂的人中间也不能引起这种怀疑；因此祂添加了这第二个见证，说：\Quote{圣父交给我要我完成的那些事工，就是我正在做的这些事工，为我作证：是圣父派遣了我。}\par

3. 在此处，祂也应对了关于违反安息日的指控。因为那些人争论说：\Quote{既因他不守安息日，他怎能是从天主来的呢？}\BibleRef{若 9:16}，因此祂说：\Quote{我父所赐给我的。}然而实际上，祂是以绝对的权能行事，但为了极其充分地显示祂不做任何与圣父相悖的事，故此祂用了表示极为卑微的语句。既然祂为何不说：\Quote{父所赐给我的工程，证明我与父同等。}？因为这两项真理都应从工程中得知，即祂未尝做任何相悖的事，且祂与生祂的那位同等；这一点祂在别处证实，祂说：\Quote{你们纵然不信我，应当信这些工程，为叫你们知道并相信我在父内，父也在我内。}\BibleRef{若 10:38}因此，从两方面看，工程都为祂作证，即祂与父同等，且祂未尝做任何与祂相悖的事。那么祂为何不如此说，反而略去较大的而提出这个呢？因为确立这一项是祂的首要目标。因为虽然让世人相信祂是从天主来的，远不及让世人相信祂与天主同等那么重大（因为前者也属于先知们，但后者则从未有过），祂仍然在次要之点上煞费苦心，因祂知道，一旦这一点被承认，另一项随后就容易被接受。所以祂绝不提证词中较重要的部分，而提出其次要的功能，好让他们借此也接受另一项。完成此事后，祂接着说：\par

% structure: p0020 source=h2 target=subsection reason=relative-heading-rank; base=h2

\subsection{若 5:37}

祂在何处为祂\Quote{作证}？在约旦河：\Quote{这是我的爱子，我所喜悦的。}\BibleRef{玛 3:16}；你们要听祂！然而即使这还需要证明。若翰的证词那时是清楚的，因为他们自己曾派人到他那里，无法否认。来自奇迹的证词也同样清楚，因为他们看见了这些奇迹的施行，并且从被治愈的人那里听见了，也相信了；他们也由此提出了控告。因此还需要为圣父的证词提供证明。接着，为达成此目的，祂接着说：\par

\Quote{你们从来没有听过祂的声音}：\par

那么梅瑟如何说呢？\Quote{上主发言，梅瑟回答}\BibleRef{出 19:19}；以及达味：\Quote{他听见了他不认识的言语}\BibleRef{咏 80:5}；以及梅瑟再次：\Quote{有哪一个民族曾\InnerQuote{听过天主的声音？}呢？}\BibleRef{申 4:33}\par

\Quote{也没有见过祂的形状。}\par

然而依撒意亚、耶肋米亚和厄则克尔，以及许多其他人，都被说曾见过祂。那么基督现在说的是什么呢？祂逐步引导他们到一个哲学道理，显示在天主那里既没有声音也没有形状，而是祂高于这样的形体和声音。因为正如当祂说：\Quote{你们没有听过祂的声音}时，祂并非说天主确实发出一种声音，而是一种无法听见的声音；同样，当祂说：\Quote{也没有见过祂的形状}时，祂并非说天主有一种形状，尽管是不可见的，而是说这两样都不属于天主。为了使他们不致说：\Quote{你是个夸口者，天主只向梅瑟说了话}；（至少他们确实说了：\Quote{我们知道天主曾与梅瑟交谈；至于这个人，我们不知道祂是从哪里来的}\BibleRef{若 9:29}）为此缘故祂那样说话，以显示天主那里没有声音也没有形状。\Quote{但是为什么，}祂说，\Quote{我提这些事？你们不仅\InnerQuote{既没有听过祂的声音也没有见过祂的形状}，而且你们甚至不能肯定你们最夸耀且你们都最确信的事，即你们领受了并遵行了祂的诫命。}因此祂接着说：\par

% structure: p0026 source=h2 target=subsection reason=relative-heading-rank; base=h2

\subsection{若 5:38}

即典章、诫命、法律和先知。因为即使天主规定了这些，它们仍不在你们那里，因为你们不信我。因为，如果圣经处处说必须留意我，而你们却不信，那么很明显祂的话已从你们那里收回。因此祂再次说：\par

\Quote{因为祂所派遣的那位，你们不信。}\par

然后，为使他们不致争论：\Quote{如果我们没有听过祂的声音，祂如何为你们作证呢？}祂说：\par

% structure: p0030 source=h2 target=subsection reason=relative-heading-rank; base=h2

\subsection{若 5:39}

因为圣父藉着这些（圣经）给出了祂的证词。祂确实也在约旦河和山上给出了证词，但基督没有引出那些声音；或许这样做祂会被不信；因为其中一个（山上的）他们没有听见，另一个（约旦河的）他们听见了却不在意。为此缘故，祂将他们引到圣经，显示圣父的证词来自圣经，祂首先移除了他们习惯夸耀的旧依据，即他们曾看见过天主或听过祂的声音。因为他们很可能不信祂的声音，并自行想象在西乃山上所发生的事，在首先纠正了他们在这几点上的疑虑，并表明所发生的事是一种迁就之后，祂才将他们引到圣经的证词。\par

4. 从这些事，我们也当在与异端争斗时武装和坚固自己。因为\BibleQuote{凡圣经都是天主所默示的，于教训、督责、使人归正、教导人学义都是有益的，叫属天主的人得以完全，预备行各样的善事}\BibleRef{弟后 3:16}；不是叫他有一些而没有其他的，因为这样的人并不是\Quote{完全}。因为请告诉我，若有人不住地祈祷，却不慷慨施舍，有什么益处？或者他慷慨施舍，却贪婪或凶暴？或者他不贪婪也不凶暴，却（慷慨）为了在人前炫耀，并博得旁观者的称赞？或者他施舍时严谨且合乎天主的喜悦，却因这事而自高自大？或者他谦卑并恒常禁食，却贪婪、贪财、钉在地上，将那万恶之源引入自己的灵魂？因为贪财是万恶之根。让我们因此战栗于这行为，逃避这罪；这罪使世界成为荒芜，这罪使一切陷入混乱，这罪诱惑我们离开基督最蒙福的服事。主说：\Quote{你们不能事奉天主而又事奉\OriginalTerm{玛门}{mammon}。}因为玛门发出与基督的命令相矛盾的命令。一个说：\Quote{施舍给那需要的人}；另一个说：\Quote{掠夺穷人的财物。}基督说：\Quote{宽恕那亏负你们的人}；另一个说：\Quote{为那没有亏负你们的人设下网罗。}基督说：\Quote{要慈悲、仁慈}；玛门说：\Quote{要残忍、残酷，视穷人的眼泪为无物}；为要使审判者在那一日严厉地对待我们。因为那时我们的一切行为都要出现在我们眼前，那些曾被我们伤害和剥夺的人，将断绝我们的一切托辞。既然拉匝禄没有从富翁那里受到任何伤害，只是没有享受他的任何好处，却在那时站出来作为严厉的控告者，不让他获得任何宽恕，那么，请告诉我，那些不但没有从自己的财物中施舍，还夺取别人的财物、拆毁孤儿房屋的人，将有什么借口呢？如果那些没有给饥饿的基督食物的人，已经将这样的火引到自己头上，那么那些抢夺根本不属于自己的东西、编织无数诉讼、不义地攫取所有人财产的人，将享受什么安慰呢？因此，让我们赶出这欲望；我们若能思想那些在我们之前行不义、贪婪而已经逝去的人，就必能赶出它。难道不是别人享受他们的财富和劳苦，而他们却躺在惩罚、报应和无法忍受的灾祸中吗？这岂不就是极端的愚妄——使自己劳累和烦恼，活着时竭力受苦，离世后却要遭受无法忍受的惩罚和报应，而本来我们可以在现世享乐（因为没有比施舍的意识更带来快乐的了），离世后则可以免除一切灾祸，获得千万祝福？因为正如邪恶惯于惩罚那些追随它的人，甚至在他们（到达）坑穴之前；同样，德行甚至在（天国的恩赐）之前，就为在此实践它的人提供快乐，使他们生活在美好的盼望和不断的喜乐中。因此，为要获得这一切，无论是今世还是来生，让我们持守善工，这样我们必将得到未来的冠冕；愿我们众人藉着我们的主耶稣基督的恩宠和慈爱，达于此境，愿光荣藉着祂并同祂，归于圣父及圣神，现在及永远，世世无尽。阿们。\par

\SourceRecord{Translated by Charles Marriott.；Nicene and Post-Nicene Fathers, First Series；Vol. 14.；Edited by Philip Schaff.；Buffalo, NY: Christian Literature Publishing Co.,；1889.；<http://www.newadvent.org/fathers/240140.htm>.}

% structure: p0001 source=h1 target=section reason=page-title

\section{若望福音讲道集 第四十一篇}

% structure: p0002 source=h2 target=subsection reason=relative-heading-rank; base=h2

\subsection{若望福音 5:39-40}

\BibleQuote{你们查考圣经，因为你们以为在其中得永生；正是这些为我作证；但你们却不愿到我这里来得[永生]生命。}\par

亲爱的，我们要高度重视属灵之事，不要认为随便追求它们就足以得救。因为如果在此世的事物中，一个人若以漫不经心、偶然的方式处理，便不能获得多大益处，那么对于属灵之事更是如此，因为它们需要更大的注意。因此，当基督将犹太人引向圣经时，祂并不是叫他们只是阅读，而是要仔细、审慎地查考；因为祂不是说：\Quote{阅读圣经}，而是说：\Quote{查考圣经}。因为关于祂的言论需要极大的注意（它们从一开始就被隐藏起来，以利于当时的人），祂现在命令他们仔细挖掘，以便能够发现深埋之下的东西。这些言论不在表面，也没有公开显露，而是像深藏的宝藏一样。凡寻找隐藏之物的人，若不精心、费力地寻找，就永远找不到他所寻的目标。因此祂说：\Quote{查考圣经，因为你们以为在其中得永生。}祂不是说：\Quote{你们有}，而是说\Quote{你们以为}，表明他们从圣经中并没有得到什么伟大或崇高的东西，因为他们期望仅凭阅读就得救，而不加上信德。因此祂所说的是这样：\Quote{你们不羡慕圣经吗？你们不以为它们是所有生命的根源吗？我现在借此确认我的主张，因为它们正是为我作证的，但你们却不愿到我这里来，为得永生。}因此祂有充分理由说\Quote{你们以为}，因为他们不愿服从，却仅仅以肤浅的阅读自夸。然后，免得因祂极其温柔的关怀，祂在他们中间招致虚荣的怀疑，又因祂渴望被他们相信，而被视为寻求自己的事（因为祂提醒他们若翰的话、天主的见证、以及祂自己的作为，并尽一切所能吸引他们到祂那里，并应许他们\Quote{生命}）；因此，我说，很可能许多人怀疑祂说这些话是出于好名的欲望，请听祂所说的话：\par

% structure: p0005 source=h2 target=subsection reason=relative-heading-rank; base=h2

\subsection{若望福音 5:41}

这就是说：\Quote{我不需要它}：祂说：\Quote{我的本性不是这样的，不需要来自人的荣耀，因为如果太阳不能从烛光中得到增加，那么我更不需要来自人的荣耀。}有人问：\Quote{那么，如果你不需要，为什么说这些话呢？}\Quote{为使你们得救。}这是祂在上面明确肯定的，同样祂在这里也暗示了，通过说\Quote{使你们得生命。}此外，祂还提出了另一个理由：\par

% structure: p0007 source=h2 target=subsection reason=relative-heading-rank; base=h2

\subsection{若望福音 5:42}

因为当他们假装爱天主而迫害祂时，因为祂使自己与天主平等，祂知道他们不会相信祂，免得有人问：\Quote{你为什么说这些话？}祂说：\Quote{我说这些话，是要指证你们，即你们迫害我不是出于爱天主，如果祂借着作为和圣经为我作证的话。因为正如先前你们认为我是天主的敌人而赶走我，现在既然我已经宣告了这些事，如果你们真正爱天主，你们就应该赶快到我这里来。但你们不爱祂。因此我说了这些话，以表明你们充满了过度的骄傲，你们是在虚妄地夸耀，并掩饰你们自己的嫉妒。}\par

% structure: p0009 source=h2 target=subsection reason=relative-heading-rank; base=h2

\subsection{若望福音 5:43}

2. 你们看见祂处处宣告祂被\Quote{派遣}，审判权由圣父交给祂，祂自己什么也不能做，为的是要杜绝他们不公的所有借口吗？但是祂在这里说的是谁要\Quote{以自己的名}而来呢？祂这里暗示的是假基督，并提出了他们不公的无可辩驳的证据。\Quote{因为如果你们因爱天主而迫害我，那么在假基督的情况下，这件事更应发生。因为他既不会说自己被圣父派遣，也不会说自己按祂的旨意而来，而是凡事相反，像暴君一样夺取不属于他的东西，并宣称他就是万有之上的天主，正如保禄所说：\BibleQuote{他高举自己超越一切称为神或受人敬拜者，甚至显示自己就是神。}\BibleRef{得后 2:4}这就是\InnerQuote{以自己的名而来}。我却不这样，而是以我父的名而来。}他们没有接受那位说自己受天主派遣的那位，这足以证明他们不爱天主；但现在从与此相反的事实，即他们将要接受假基督，祂显示出他们的无耻。因为当他们不接受那位宣称自己被天主派遣的那位，却要敬拜一位不认识天主并宣称自己是万有之上的天主的一位，显然他们的迫害出于恶意和对天主的仇恨。因此，祂为自己的话提出了两个理由：首先是较仁慈的：\Quote{为使你们得救}，以及\Quote{为使你们得生命}；当他们嘲笑祂时，祂提出了另一个更惊人的理由，表明即使听众不信，天主仍然习惯于做祂自己的事。现在保禄论及假基督时曾预言说：\BibleQuote{天主将给他们一种强大的迷惑，以致他们全都被定罪，因为他们不信真理，反倒喜爱不义。}\BibleRef{得后 2:11}基督没有说：\Quote{祂必来}，而是说\Quote{若祂来}，出于对听众的温柔；因为他们的顽固尚未完全。祂对祂来的原因保持沉默；但保禄，对于能理解的人，已特别暗示了。因为正是他剥夺了他们一切的借口。\par

基督于是也指出了他们不信的原因，说……\par

% structure: p0012 source=h2 target=subsection reason=relative-heading-rank; base=h2

\subsection{若望福音 5:44}

故此，祂再一次表明，他们关注的并非天主的事，而是藉此托辞满足私情，他们远非为祂的光荣而行此事，反倒将来自人的光荣置于来自祂的光荣之上。那么，他们既如此轻看那种光荣，以致偏爱来自人的光荣，又如何可能对那光荣怀有敌意呢？\par

在告诉他们他们没有天主的爱，并以祂身上所发生的事、以及将来在假基督身上发生的事证明了这一点，又表明他们已失去一切借口之后，祂接着将梅瑟带出来作为他们的控告者，继续说道：\par

% structure: p0015 source=h2 target=subsection reason=relative-heading-rank; base=h2

\subsection{\BibleRef{若 5:45-47}}

祂所说的话大意如下：\Quote{是梅瑟，而不是我，因你们待我的行为而受了更大的侮辱，因为你们不信他，甚于不信我。}请看，祂如何从各方面剥夺了他们所有的借口。\Quote{你们说，当你们迫害我时，你们爱天主；我已经表明，你们这样做是出于对祂的仇恨。你们说我违反安息日、废除法律；我也摆脱了这种毁谤。你们声称，因你们胆敢对我所做的事，你们相信梅瑟；相反，我表明这正是最不信梅瑟；因为我远非反对法律，那位控告你们的人，正是给你们法律的这位。}正如祂论及圣经时所说的，\Quote{你们以为在其中得永生}，照样论及梅瑟时祂也说，\Quote{你们所信赖的}，处处以他们自己的武器战胜他们。\par

有人说：\Quote{从何看出梅瑟会控告我们，而你不是自夸呢？你和梅瑟有何关系？你违反了他命令我们遵守的安息日；他怎能控告我们呢？又怎能看出我们会相信另一个以自己名而来的人呢？所有这些主张你都没有证据。}其实这一切在上述经文中都已得到证明。\Quote{因为}（基督会回答）\Quote{既然从这些作为、若翰的呼声和圣父的见证，都公认我是从天主而来的，那么显然梅瑟会控告犹太人。}因为他怎么说？\Quote{若有一人行奇迹，引导你们归向天主，并真实预言未来之事，你们必须欣然听从。}基督行了这一切。祂确实行了奇迹，祂引导众人归向天主，并且（如此祂）使祂的预言得以应验。\par

\Quote{但从何看出他们会相信另一个呢？}从他们仇恨基督看出，因为那些偏离按天主旨意而来者的人，显然会接纳天主的仇敌。若祂现在提出梅瑟，尽管祂说过\Quote{我不接受人的证据}，也不必惊奇，因为祂不是指引他们到梅瑟那里，而是到天主的圣经。然而，由于圣经对他们的威吓较小，祂将论述转向这个人本身（梅瑟），将他们的立法者作为控告者摆在他们面前，使威吓更为有力；祂驳斥了他们每个主张。请注意：他们说他们因爱天主而迫害祂，祂表明他们其实是因恨天主而这样做；他们说他们坚守梅瑟，祂表明他们这样做是因为不信梅瑟。因为若他们热心于法律，就该接纳那成全法律者；若他们爱天主，就该相信那引领他们归向祂的那一位；若他们相信梅瑟，就该敬拜梅瑟所预言的那一位。\Quote{但是}（基督说）\Quote{若梅瑟在我来临之前已不被相信，那么，我这个由他预告的人被你们赶走，也不足为奇。}正如祂已从他们对待祂的举止中表明，那些赞赏若翰的人（实际上）轻看他，同样现在祂表明，那些自以为相信梅瑟的人其实不信他，并将他们以为可用来为自己辩护的一切反击到他们自己头上。\Quote{我远非}，祂说，\Quote{使你们离开法律，我反而召你们的立法者亲自作你们的控告者。}\par

祂曾说过圣经为祂作证，但祂并未加添何处作证，意在使他们更加敬畏，并促使他们去查考，迫使他们不得不询问。因为若祂不经他们询问就轻易告诉，他们会拒绝这见证；但现在，若他们对祂的话稍加留意，他们首先需要询问，并从祂那里得知那见证是什么。为此祂更多用断言和威吓，而非仅用证据，以便即使如此，祂也能藉祂所说的恐惧把他们争取过来；但他们仍然沉默。邪恶就是这样；无论人说什么或做什么，它都不被触动，而是保持其特有的毒液。\par

为此，我们必须从灵魂中驱除一切邪恶，不再图谋任何欺诈；因为有人曾说：\BibleQuote{天主使弯曲的道路临于乖戾的人} \BibleRef{箴 21:8}；又说：\BibleQuote{训诫的圣神必远避欺诈，远离无知的思念。} \BibleRef{智 1:5} 因为没有什么比邪恶更使人愚昧；当一个人奸诈、不公、忘恩（这些都是邪恶的不同形式），当他在未受委屈时使别人受苦，当他编造欺诈时，他岂不表现出极端的愚昧吗？同样，没有什么比德行更使人明智；它使人感恩、公正、慈悲、温和、谦逊、坦率；它常是其他一切福惠之母。还有什么比如此性情的人更有理解力呢？因为德行本身就是明智的泉源和根基，正如一切邪恶始于愚昧。因为傲慢和愤怒的人因缺乏智慧而成为各自情欲的猎物；因此先知说：\BibleQuote{我身上没有一块完好的肌肉：我的创伤流脓发臭，皆因我的愚昧。} \BibleRef{咏 38:3}：表明一切罪恶始于愚昧；所以敬畏天主的德行之人比任何人更有理解力；因此一位智者曾说：\BibleQuote{敬畏上主是智慧的开端。} \BibleRef{箴 1:7} 如果敬畏天主即是拥有智慧，而恶人没有这样的敬畏，那么他就被剥夺了真正的智慧——而被剥夺了真正的智慧，他比任何人都更愚昧。然而许多人却钦佩恶人能为非作歹，不知道他们理应视其为最可怜的人，那些以为在伤害别人的人，其实是把剑刺向自己——这是极端的愚昧行为：一个人打击自己，却甚至不知道自己在这样做，反而以为在伤害别人，其实是在杀死自己。因此\Person{保禄}知道我们击打别人时是在杀害自己，便说：\BibleQuote{你们为什么不情愿受欺？为什么不情愿吃亏？} \BibleRef{格前 6:7} 因为不受委屈在于不委屈任何人，同样不受虐待在于不虐待他人；尽管这种说法在许多人看来，在那些不愿学习真正智慧的人看来，可能像谜语一样。知道了这一点，我们不要称那些受伤害或受侮辱的人为可怜或为他们悲哀，而要为那些施加伤害的人悲哀；他们才是受害最深的人，他们使天主与他们为敌，开启了无数控告者的口，在今生得恶名，并为来世招致严厉的惩罚。而那些被他们冤枉并高贵地忍受了一切的人，却得天主的眷顾，所有人都同情、赞美、接待他们。这样的人在今生将享有极好的声誉，作为真正智慧最强有力的榜样，并在来世分享永恒的福乐；愿我们众人赖我们的主耶稣基督的恩宠和慈爱都能达到，愿光荣归于他及圣父和圣神，从现在直到永远，世世无尽。阿们。\par

\SourceRecord{Translated by Charles Marriott.；Nicene and Post-Nicene Fathers, First Series；Vol. 14.；Edited by Philip Schaff.；Buffalo, NY: Christian Literature Publishing Co.,；1889.；<http://www.newadvent.org/fathers/240141.htm>.}

% structure: p0001 source=h1 target=section reason=page-title

\section{若望福音讲道第四十二篇}

% structure: p0002 source=h2 target=subsection reason=relative-heading-rank; base=h2

\subsection{若望福音 6:1-4}

1. 亲爱的，我们不要与强暴的人争斗，而要学习，当这样做对我们的德性无伤时，就退避他们的恶计；因为如此，他们的一切蛮横就会受到抑制。正如箭镞落在坚硬、牢固而有抵抗力的物体上时，会以极大的力量反弹向投掷者，但当投掷的猛力没有遇到阻力时，它很快就会变弱而停止；傲慢的人也是如此：当我们与他们相争时，他们变得更凶猛，但当我们让步退却时，就轻易地平息了他们的一切狂怒。因此，主知道法利塞人听说\Quote{耶稣所收的门徒和所施的洗比若翰还多}时，就去了加里肋亚，为熄灭他们的嫉妒，并通过退隐来软化这些传闻可能激起的愤怒。当祂第二次前往加里肋亚时，祂没有来到先前同样的地方；因为祂没有去加纳，而是到了\Quote{海的对岸}，有一大群人跟随祂，看见\Quote{祂所行的奇迹}。什么奇迹？为什么他没有具体提及？因为这位福音作者最希望将他的书的大部分用于（基督的）讲论和训诲。请看，例如，整整一年，甚至现在在逾越节庆期，他在奇跡方面只给了我们不多信息，仅仅提到祂治愈了瘫痪病人和贵族之子。因为他并不急于一一列举它们（那是不可能的），而是从许多伟大的奇迹中记录少数几个。\par

第2节。\Quote{有一大群人跟随祂，看见祂所行的奇迹。} 这里所讲述的表明一种不甚明智的心态；因为当他们享受了这样的教导后，仍更被奇迹所吸引，这是更属粗俗状态的标志。因为经上说：\Quote{奇迹不是为信者，而是为不信者。} \BibleRef{若 6:2} 玛窦所描述的人们行为并非如此，而是怎样？他说，所有人都\Quote{因祂的教训而惊异，因为祂教训他们，像有权威的人。} \BibleRef{玛 7:28}\par

\Quote{为什么祂现在上了山，并和祂的门徒坐在那里？} 因为即将发生的奇迹。只有门徒和祂一起上去，这是对没有跟随祂的人群的指责。然而祂上山并非仅为此，也是为了教导我们要时常从日常生活的喧嚣和混乱中休息。因为独处是适于研究智慧的事。祂也时常独自上山，整夜祈祷，教导我们，最亲近天主的人必须摆脱一切纷扰，寻求没有混乱的时间和地点。\par

% structure: p0006 source=h2 target=subsection reason=relative-heading-rank; base=h2

\subsection{若望福音 6:4}

\Quote{那么，有人说，祂为什么不上节去，却当众人都赶往耶路撒冷时，自己去了加里肋亚，而且不单是祂自己，还带着门徒，然后从那里往葛法翁去？} 因为从此祂就悄悄地废除法律，藉着犹太人的恶行来行事。\par

% structure: p0008 source=h2 target=subsection reason=relative-heading-rank; base=h2

\subsection{若望福音 6:5}

这表明祂并非在任何时候都无所事事地与门徒坐着，而是或许在仔细地与他们交谈，使他们留心并转向祂，这特别标志出祂对他们的温柔关怀，以及祂对待他们的谦卑和屈尊。因为他们与祂同坐，或许彼此相望；然后祂举目看见许多群众向祂走来。其他福音作者说，门徒来求祂，请祂不要叫他们空腹散去，而圣若望说，这个问题是由基督向斐理伯提出的。在我看来，两种说法都是真实的，但并非发生在同一时间，前一个叙述先于另一个，所以两者完全不同。\par

那么，祂为什么问\Quote{斐理伯}？祂知道哪个门徒最需要教导；因为这就是后来那说\Quote{请将父显示给我们，我们就满足了}的人\BibleRef{若 14:8}，因此耶稣事先使他进入适当的状态。因为如果奇迹只是简单地发生，其奇妙之处似乎不会如此巨大；但现在祂先强迫他承认当前的匮乏，好使他了解事情的状况后，能更准确地认识即将发生的奇迹的伟大。因此祂说……\par

% structure: p0011 source=h2 target=subsection reason=relative-heading-rank; base=h2

\subsection{若望福音 6:5}

在旧约中，祂也对梅瑟说话，因为祂先问他\Quote{你手里是什么？}然后才施行征兆。因为那些出乎意料并突然发生的事，常会使人忘记之前的事，因此祂先使他承认当前的情况，好使当他感到惊异时，之后不能驱散他所承认的记忆，从而通过比较得知奇迹的伟大。事实上，此处正是如此；因为斐理伯被问时回答说：\par

% structure: p0013 source=h2 target=subsection reason=relative-heading-rank; base=h2

\subsection{若望福音 6:7}

2. \Quote{为试探他}是什么意思？难道祂不知道他会说什么？我们不能肯定这一点。那么，这句话的意思是什么？我们可以从旧约中领悟。因为那里也说：\Quote{这些事以后，天主试探亚巴郎，对他说：带你心爱的独生子去} \BibleRef{创 22:1}；然而，在那里也不显示祂说这话时是在等待试探的结果，即亚巴郎是否服从（祂怎能呢？祂在一切事物存在之前就知道一切！）。但两处的话都是照人的方式说的。因为正如（圣咏作者）说祂\Quote{窥察人心}，他不是指出于无知的窥察，而是指确切的知识；同样，当福音作者说祂试探（斐理伯）时，仅指祂确切知道。或许还可以说另一件事：正如祂曾使亚巴郎更得称赞，同样祂也使这人这样，通过这问题使他确切认识奇迹。因此，福音作者为了不使你停留在措辞的软弱上，从而对所说的话形成不当的看法，加上说：\Quote{祂自己原知道要做什么。}\par

此外，我们必须注意这一点：每当有错误的怀疑时，作者会立刻非常仔细地加以纠正。正如在此处，为使听众不致产生任何此类怀疑，他加了纠正的话说：\Quote{因为祂自己知道祂要做什么}；同样在另一处，当祂说\Quote{犹太人迫害祂，因为祂不但破坏了安息日，而且还说天主是自己的圣父，使自己与天主平等}时，若没有基督自身藉其工作所确认的断言，他也会在此处加上这个纠正。因为即使在基督自己所说的话语中，福音作者也谨慎地不让任何人产生怀疑，更何况在他人论及祂时，若他察觉到关于祂的不当意见盛行，他必定会仔细审视。但他没有这样做，因为他知道这是祂的意思和不可动摇的定意。因此，在说了\Quote{使自己与天主平等}之后，他未曾使用任何此类纠正；因为所谈论的事不是他们错误的想象，而是祂自己的断言，且由祂的工作所确认。于是，斐理伯被问后，\par

第8、9节。\Quote{西满的兄弟安德肋说：\InnerQuote{这里有一个孩子，他有五个大麦饼和两条小鱼：但这么多人中，这些算什么呢？}}\par

安德肋比斐理伯更高尚，但尚未达到一切。然而我不认为他说话没有目的，而是因为他曾听过先知们的奇迹，以及厄里叟如何用饼行了一个神迹\BibleRef{列下 4:43}；为此他上升到一定高度，却未能达到顶端。\par

那么，我们这些沉溺于奢华的人，让我们学习那些伟大而可敬的人的食物是什么；并在质与量上，让我们注视并效法他们餐桌的俭朴。\par

接下来的话也表现出极大的软弱。因为在说了\Quote{有五个大麦饼}之后，他加上\Quote{但这么多人中，这些算什么呢？}他以为奇迹的施行者会从少中变少，从多中变多。但事实并非如此，因为祂从多中或少中使饼生出同样容易，因为祂不需要任何质料。但为使受造之物不显得与祂的智慧无关（正如后来的诽谤者和那些染上马西昂病的人所说），祂使用了受造之物本身作为祂奇事的基础。\par

当两位门徒都承认自己无能为力时，祂便行了奇迹；因为这样他们获益更大，首先承认了事情的困难，以便当事情发生时，他们能明白天主的德能。并且因为将要行的奇迹也曾由先知们施行过，尽管程度不同，而且祂要先感谢后才行此事，免得他们因祂方面有任何软弱的怀疑，请注意，祂如何藉其行事的方式完全提升他们的思想并显示区别。因为当饼尚未出现时，为使你们学习，那不存在的在祂眼中如同存在（正如保禄所说：\Quote{呼召那不存在的如同存在的}\BibleRef{罗 4:17}），祂命令他们仿佛宴席已预备好，立即坐下，以此激发祂门徒的心智。因为他们已从询问中获益，便立即服从，没有困惑，也没有说：\Quote{这是怎么回事？祢为何叫我们坐下，面前却一无所有？}这些人先前如此不信，以至于说：\Quote{我们从哪里买饼呢？}现在甚至在看见奇迹之前就开始相信，以至于他们乐意让群众坐下。\par

3. 但是为什么当祂要复活瘫子时，祂没有祈祷，也没有在复活死人或平定大海时祈祷，而在此处却对饼做了这事呢？这是为了表明，当我们开始用餐时，应当感谢天主。而且祂尤其在一个较小的事上这样做，为使你们明白祂这样做并非因为祂有需要；因为如果祂有需要，在更大的事上祂更会这样做；但祂是以自己的权柄行那些大事，显然祂在较小的事上这样做是出于迁就。此外，有许多群众在场，有必要使他们相信祂是按照天主的旨意而来的。因此，当祂在没有见证者面前行奇迹时，祂不展示任何这类举动；但当祂在许多人在场时行奇迹，为了说服他们祂不是天主的敌人，不是生祂者的敌对者，祂藉感谢除去怀疑。\par

\Quote{祂便分给坐下的人，他们都吃饱了。}\par

你们看见仆人与主人之间有多大的差距吗？他们按度量领受恩宠，相应地行奇迹；但天主以自由的权能行事，一切都极丰沛地成就。\par

\BibleRef{若 6:12} . \Quote{祂对门徒说：\InnerQuote{把剩余的零碎收拾起来。}——他们便收拾起来，装满了十二筐。}\par

这并不是多余的表现，而是为使此事不被人视为纯属幻像；因此祂从已经存在的物质中创造。\Quote{但祂为何不将饼交给群众拿着，而只交给祂的门徒呢？}因为祂最渴望教导这些将要成为世界教师的人。群众从奇迹中尚未获得多大的益处（至少他们立刻忘记了这一个，而求问另一个），而门徒们则会获得非同寻常的益处。此外，所发生的事也是对提篮子的犹大的一种不寻常的谴责。这些事是为教训他们而做的，这从后来祂提醒他们时所说的话可以清楚看出：\Quote{你们还不明白——你们收起了多少筐子？}\BibleRef{玛 16:9} 同样，碎块的筐子数目与门徒的数目相等；后来，当他们受了教训，就不再收起那么多，而只有\Quote{七个筐子}。\BibleRef{玛 15:37} 我不仅惊叹于所造饼的数量，而且惊叹于剩余量的精确性：祂使超出的部分不多不少，正合祂所愿意的量；祂预见了他们要消耗多少；这标志着无可言喻的能力。碎块然后证实了这事，显示了这两点：所发生的不是幻像，而且这些碎块正是来自喂饱众人的饼。至于鱼，此时它们是从已有的鱼中产生的，但后来，在复活之后，它们不是从已有的物质中产生的。\Quote{为什么呢？}为叫你们明白，即使现在祂使用物质，并非出于必需，也不是需要任何基础（来工作），而是为堵住异端者的口。\par

\Quote{群众说：这真是那位先知。}\par

唉，贪食的过度！祂曾行了成千上万比这更令人赞叹的事，但只有在他们吃饱时，他们才作了这样的宣认。然而，由此可知他们期待一位卓越的先知；因为那些人对若翰说过：\Quote{你是那位先知吗？}而这些人则说：\Quote{这真是那位先知。}\par

第15节。\Quote{耶稣知道他们要来强迫祂为王，就又退到山上去了。}\BibleRef{若 6:15}\par

奇妙！贪食的暴政何等巨大，人心的反复无常何等巨大！他们不再维护法律，不再关心安息日的触犯，不再为天主热心；当他们的肚腹被填满时，这一切考虑都被抛在一边。在他们眼中祂是先知，他们将要选祂为王。但基督逃走了。\Quote{为什么呢？}为教导我们轻视世上的尊荣，并表明祂在地上毫无所需。因为那选择一切卑微事物——母亲、房屋、城市、养育和衣着——的主，不会后来因世上的事而变得尊贵。祂从天上的事物是荣耀而伟大的：天使、一颗星、祂的圣父高声说话、圣神作证、先知从远处宣告祂；地上的事物却全是卑微的，为叫祂的能力更显明。祂也来教导我们轻视世界上的事，不被今生的荣华所惊异或震惊，反倒嘲笑它们，并渴慕将来的事。因为谁若赞叹现实的事，必不会赞叹天上的事。因此祂也对比拉多说：\Quote{我的王国不属于这世界}\BibleRef{若 18:36}，免得后来祂似乎只是利用属人的恐怖或权势来说服人。既然如此，为何先知说：\Quote{看，你的君王到你这里来，谦逊地骑在驴上}？\BibleRef{匝 9:9} 祂所说的是那天上的王国，不是这地上的；为此基督说：\Quote{我不接受人的光荣。}\BibleRef{若 5:41}\par

我们务要学习，亲爱的，轻视并不要贪求来自人的荣耀；因为我们已获得了极大的荣耀，与之相比，那另一种实在是一种侮辱、讥讽和嘲弄。就如今世的财富与那财富相比是贫穷，这生命若无那生命便是死亡，（因为\BibleQuote{让死人埋葬他们的死人}\BibleRef{玛 8:28}），同样，这荣耀与那荣耀相比便是羞耻和嘲弄。所以我们不要追求它。若那赐予荣耀的人比影子和梦还微不足道，那荣耀本身更是如此。\BibleQuote{人的光荣有如草上的花}\BibleRef{伯前 1:24}；有什么比草上的花更卑贱呢？就算这光荣是永久的，它又能对灵魂有什么益处呢？毫无益处。相反，它大大地损害我们，使我们成为奴隶，是比用金钱买来的更糟糕的奴隶，服从的不是一个主人，而是两个、三个、一万个，都发着不同的命令。做一个自由人，不受人的奴役，只服从天主的统治，这比做奴隶好多少呢？简言之，你若渴望光荣，就渴望吧，但要渴慕那不朽的光荣，因为那光荣是在更光荣的舞台上展示，并带来更大的益处。因为这里的人命令你花钱取悦他们，但基督相反，祂因你所给与的，赐给你百倍的回报，还加上永生。那么这两者中哪个更好：是在地上受赞美，还是在天上？是由人，还是由天主？是损失，还是得益？是戴冠一天，还是永世无穷？要施舍给那有需要的人，但不要施舍给舞者，免得你损失钱财又毁灭他的灵魂。因为你是他走向丧亡的原因，由于不合时宜的慷慨。假如舞台上的人知道他们的行当无利可图，他们早就停止操练了；但当他们看到你喝彩、追随、为他们花费和耗费钱财时，即使他们本无意继续，也会被获利的欲望所束缚。如果他们知道没有人会称赞他们所行的，他们就会因无利而很快停止劳作；但当他们看到这行为被许多人赞赏时，他人的赞美就成了他们的诱饵。那么，我们停止这种无益的开支吧，让我们学习该在谁身上和何时花费。我恳求你们，不要在两方面都惹怒天主：从不应取的地方收取，又在不应散的地方散施。你们的行为岂不配得愤怒吗？当你忽略穷人，却给妓女施舍时，即使你从正当的收入中支付，难道雇佣罪恶的代价、把荣耀给予本应惩罚的地方，不也是对你的控诉吗？更何况你是通过抢夺孤儿、冤屈寡妇来滋养你的不洁！想想看，对那些胆敢如此行事的人，怎样的烈火已为他们预备！听听保禄所说的：\BibleQuote{他们不仅自己做这些事，还赞同别人去做}\BibleRef{罗 1:32}。\par

或许我们刺痛了你，然而即使我们不刺痛你，那些犯罪而不悔改的人仍有实际的惩罚在等待。那么，用言语取悦那些将被现实惩罚的人，又有什么益处呢？你喜爱一个舞者，你赞美他、钦佩他吗？那么你比他更坏；他的贫穷给了他一个借口，尽管不是合理的，但你连这个辩护也被剥夺了。若我问他：\Quote{你为什么离开其他的技艺，来从事这受诅咒和不洁的行当？}他会回答说：\Quote{因为我可以用很少的劳力获得很大的利润。}但若我问你，为什么你钦佩一个在污秽中消磨时光、为害众人的人，你就不能拿出同样的借口，只能低头羞愧脸红。现在，当我们叫你来解释时，你无话可说；当那可怕且严厉的审判来临时，我们还要为思想和行为以及一切做解释，那时我们怎能站立？我们以怎样的眼睛瞻仰我们的审判者？我们说什么？我们作什么辩护？我们提出什么合理或不合理的借口？我们能否列举花费？消遣？或者藉着他的技艺我们毁灭的他人？我们无话可说，而必须受那无终止、无限制的刑罚。为了不让这事发生，让我们从此谨守一切，好使我们带着美好的希望离世，获得永福。愿我们众人赖我们的主耶稣基督的恩宠和慈爱，得以到达这永福；藉着祂并同祂，偕同圣父及圣神，享受光荣，从现在直到永远，万世无疆。阿们。\par

\SourceRecord{Translated by Charles Marriott.；Nicene and Post-Nicene Fathers, First Series；Vol. 14.；Edited by Philip Schaff.；Buffalo, NY: Christian Literature Publishing Co.,；1889.；<http://www.newadvent.org/fathers/240142.htm>.}

% structure: p0001 source=h1 target=section reason=page-title

\section{若望福音讲道集第四十三篇}

% structure: p0002 source=h2 target=subsection reason=relative-heading-rank; base=h2

\subsection{若望福音 6:16-18}

基督不仅当祂肉身临在时，也为远离时的门徒谋福；因祂有丰富的手段和技巧，藉相反的行动达到同一目的。例如，看祂在此所行的。祂离开门徒，上山去；门徒到了晚上，就下到海边。他们等候祂直到傍晚，期待祂会来到他们那里；但到了晚上，他们再也忍不住不去寻找他们的主；他们心中充满如此大的爱。他们没有说：\Quote{现在已是傍晚，黑夜临到我们，我们要往哪里去？地方危险，时间不安全}；但因渴望所驱使，他们上了船。因为福音作者并非无端地记载了时间，而是借此显示他们爱情的热烈。\par

那么为什么基督让他们离开，而不显现自己？再者，为什么祂独自在海面上行走显现？前者教导他们被祂舍弃是多么大的恶，并使他们的渴望更加强烈；后者再次显示祂的权能。因为在教导中，他们并没有与群众一同听到所有的事，同样在奇迹中，他们也没有与众人一起看到一切，因为那些将要接受管理世界的职责的人，需要比其余的人多有一些。有人问：\Quote{他们单独看见了什么样的奇迹？} 山上显圣容；这次在海面上，以及复活后的那些奇迹，既多又重要。从这些我推测还有别的奇迹。他们来到葛法翁，没有任何确定的消息，但期待在那里或在中途找到祂；福音作者暗示这一点，说：\Quote{那时天已经黑了，耶稣还没有来到他们那里。}\par

\Quote{因着大风吹刮，海就翻腾起来。} 他们做了什么？他们不安，因为有多种原因迫使他们如此。他们因时间而害怕，因为天黑了；因风暴而害怕，因为海翻腾了；因地点而害怕，因为他们不靠近陆地；但是，\par

% structure: p0006 source=h2 target=subsection reason=relative-heading-rank; base=h2

\subsection{若望福音 6:19}

最后，因着这事的神奇，因为——\par

\Quote{他们看见祂在海面上行走。} 当他们非常不安时，\par

% structure: p0009 source=h2 target=subsection reason=relative-heading-rank; base=h2

\subsection{若望福音 6:20}

那么祂为何显现？为要显示正是祂要使风暴平息。因为福音作者已显示这一点，说：\par

% structure: p0011 source=h2 target=subsection reason=relative-heading-rank; base=h2

\subsection{若望福音 6:21}

祂不仅给了他们安全的航行，而且还有顺风。\par

祂不在群众面前显示在海面上行走，因为那奇迹对他们软弱的本性来说过于伟大。即便是门徒，祂也没有长时间这样做，而是显现后立刻退去。在我看来，这与玛窦福音第十四章中的奇迹不同；从许多理由可以清楚看出那是不同的。因为祂常行同样的奇迹，为使观看者不仅极其惊讶，而且也以大的信德领受它们。\par

\BibleQuote{是我，不要害怕。} 当祂说话时，祂从他们灵魂中驱逐了恐惧。但在另一次并非如此；因此伯多禄说：\BibleQuote{主，如果是祢，就请祢命令我来到祢那里。}\BibleRef{玛 14:28} 那么，为什么那时他们没有立刻相信这个，而现在却被说服了？因为那时风暴继续颠簸小船，但现在祂一说话，风平浪静了。或者如果不是这个原因，那就是我之前提到的另一个原因：祂常行同样的奇迹，藉第一个使第二个更容易被接受。但为什么祂没有上船？因为祂要使奇迹更大，更公开地向他们显示祂的天主性，并要向他们表明，当祂先前感谢时，祂并非因为需要帮助而这样做，而是为了迁就他们。祂允许风暴起，使他们永远寻求祂；祂平息风暴，为使他们认识祂的权能；祂没有上船，为要使奇迹更大。\par

% structure: p0015 source=h2 target=subsection reason=relative-heading-rank; base=h2

\subsection{若望福音 6:22}

为什么若望如此精确？为什么他不说群众在第二天过后就离开了？他想要教导我们另外的事，即耶稣允许群众以隐秘的方式，即使不是公开地，猜想所发生的事。因为，他说：\Quote{他们看见那里没有别的船，只有那一只，并且耶稣没有同祂的门徒上船。}\par

% structure: p0017 source=h2 target=subsection reason=relative-heading-rank; base=h2

\subsection{若望福音 6:24}

那么他们还能猜想什么，除非祂是步行渡过海而到达那里？因为不可能说祂是从别的船过去的。因为福音作者说：\Quote{只有一只，祂的门徒上了那船。} 尽管如此，当他们因如此大的奇迹而来见祂时，他们没有问祂如何过来，如何到达那里，也没有寻求明白如此大的征兆。但他们说什么？\par

% structure: p0019 source=h2 target=subsection reason=relative-heading-rank; base=h2

\subsection{若望福音 6:25}

除非有人断定，他们在此用\Quote{何时}一词表示\Quote{如何}的意思。但在此也值得注意他们冲动的反复无常。凡说\Quote{这位真是那位先知}，凡急于\Quote{要强逼祂立祂为王}的人，如今当他们找到祂时，却不再有此图谋；他们丢弃了惊奇，不再因祂先前的作为而钦佩祂。他们寻找祂，渴望再次享受像第一次那样的筵席。\par

梅瑟领导下的犹太人渡过了红海，但那件事与这事大不相同。祂以祈祷并以仆人的身份做了一切，而基督则以绝对权能行事。那里当南风吹起时，海水退去，使他们能在干地上过去，但这里的奇迹更大。\BibleRef{出 14:21} 因为海保留其本性，却在其表面承载了它的主，从而证实了圣经所说的：\BibleQuote{祂行走在海面上，如同走在铺石路上。}\BibleRef{约 9:8}\par

这是有理由的：当祂将要进入固执而不驯服的\Place{葛法翁}时，祂行了增饼的奇迹，因为祂不仅愿意藉城内所发生的事，也愿意藉城外所行的奇迹来软化那城的不驯服。难道众多群众如此热切地涌向那座城，还不足以软化甚至石头吗？然而他们却没有这样的感受，反而又渴望身体的食物；为此他们受到\Person{耶稣}的责备。\par

那么，诸位亲爱的，我们既然知道这些事，就应为感官的事物感谢天主，但更为属神的事物感谢；因为这是祂的旨意，祂正是为了后者才赐予前者，仿佛藉着这些事物引领那些较为不完美的人，并给予他们准备性的教导，因为他们还贪恋世俗。但是当这些人接受了这些世俗事物，便安于其中，那时他们就会受到斥责和责备。例如那个瘫痪的人，基督首先愿意赐予属神的事物，但当时在场的人却无法忍受；因为当祂说\BibleQuote{你的罪赦了}\BibleRef{玛 9:2}时，他们喊道\Quote{这人亵渎了。}让我们不要这样受影响，而应更重视那些（属神）的事物。为什么呢？因为当属神的事物与我们同在时，缺少肉身的事物不会带来伤害；但当属神的事物不在时，还有什么希望、什么安慰能留给我们呢？因此，我们应该时常呼求天主，为这些事向祂恳求。基督也正是这样教导我们祈祷的；因为如果我们展开那篇祈祷文，就会发现其中没有一样属于血肉，全都是属神的，甚至连那一小部分看似与感官有关的内容，也因方式而成为属神的。因为吩咐我们只求我们的\Quote{连续不断的}食粮，即我们的\Quote{每日}食粮，就标志着一个属神而真正明智的心智。并且看看那之前的祷词：\Quote{愿你的名被尊为圣，愿你的国来临，愿你的旨意奉行在人间，如同在天上}；然后，在提及那暂时的需要之后，祂迅速离开它，再次将我们带回属神的教导，说：\Quote{宽免我们的罪债，如同我们也宽免得罪我们的人}。祂在这祈祷文中从未放入财富、光荣或权柄，而是全部有助于灵魂的益处；没有一样是地上的，样样都是天上的。如果我们被命令远离今世的事物，那么向天主祈求那些即使拥有也被祂命令我们抛弃（为了使我们免除挂虑）并且祂命令我们不要为之操劳的事物，我们怎能不感到可怜和悲惨呢？这就是\Quote{重复无用的言词}；也正是我们的祈祷毫无效果的原因。\Quote{那么，}有人说，\Quote{恶人怎能致富，不义和不洁之人、劫掠者和贪得无厌者又怎能富足？}并非出于天主的赐予；（断乎不可！）而是由于劫掠和索取超过他们应得的。\Quote{天主为何允许他们？}正如祂允许那个富翁，将他留作更严厉的惩罚。\BibleRef{路 16:25} 请听（\Person{亚巴郎}）对他所说的话：\BibleQuote{孩子，你活着时已享尽了你的福，而\Person{拉匝禄}同样受尽了苦；现在他在这里受安慰，而你却受苦。}\BibleRef{路 16:25}因此，为了使我们也不致听到那声音，由于我们生活安逸懒惰，为自己积聚了许多罪恶，让我们选择真正的财富和正确的智慧，好能获得那所应许的美物；愿我们众人全靠我们的主\Person{耶稣基督}的恩宠和慈爱，到达那境界；祂和圣父及圣神，同享光荣，从现在直到永远。阿们。\par

\SourceRecord{Translated by Charles Marriott.；Nicene and Post-Nicene Fathers, First Series；Vol. 14.；Edited by Philip Schaff.；Buffalo, NY: Christian Literature Publishing Co.,；1889.；<http://www.newadvent.org/fathers/240143.htm>.}

% structure: p0001 source=h1 target=section reason=page-title

\section{论若望福音第44篇讲道}

% structure: p0002 source=h2 target=subsection reason=relative-heading-rank; base=h2

\subsection{若望福音 6:26-27}

\BibleQuote{耶稣回答他们说：我实实在在告诉你们：你们找我，并不是因为看见了神迹，而是因为吃饼吃饱了。不要为那可朽坏的食粮劳碌，而要为那存留到永生的食粮劳碌。}\par

1. 温和与温柔并不总是有用，有时教师需要更尖锐的语言。因为如果门徒迟钝粗俗，那么为了深深触动他的迟钝，我们必须用刺棒唤醒他。天主子在现在以及许多其他情况下都这样做了。因为当群众找到耶稣，并奉承祂说：\Quote{老师，祢什么时候到这里来的？}时，为了表明祂不渴望人的尊荣，而只关注一件事，即他们的救恩，祂严厉地回答他们，不仅想以这种方式纠正他们，而且通过揭示和暴露他们的思想。因为祂说什么？\BibleQuote{我实实在在告诉你们（正面地和肯定地说），你们找我，并不是因为看见了神迹，而是因为吃饼吃饱了。}祂用这些话责备和斥责他们，但并不突然或猛烈，而是非常克制。因为祂不说：\Quote{噢，你们这些贪吃的人和肚腹的奴隶！我行了那么多神迹，你们却从来没有跟随我，或对我的作为感到惊奇。}而是温和而温柔地以这种方式说：\BibleQuote{你们找我，并不是因为看见了神迹，而是因为吃饼吃饱了。}不仅指过去，也指现在的神迹。\BibleQuote{使你们惊讶的（祂说），不是饼的神迹，而是吃饱了。}祂说这话不是猜测，他们立刻显明，因这个缘故他们第二次来，如同要享受与先前同样的食物。所以他们又说：\Quote{我们的祖先在旷野吃了玛纳。}再次把他们引到肉身的食物上，这是对他们主要的控告和指责。但祂不停止于责备，也加上教导，说：\BibleQuote{不要为那可朽坏的食粮劳碌，而要为那存留到永生的食粮劳碌。}\par

\BibleQuote{人子要赐给你们的，因为天主圣父已印证了祂。}\par

祂所说的是这样的：\Quote{不要看重属地的，而要看重属神的食物。}但有些想要懒惰度日的人滥用这话，好像基督要完全废除工作，因此现在适合对他们说些什么。因为他们可以说是诽谤所有基督徒，使基督教因懒惰而受嘲笑。首先，我们必须提及保禄的那句话。他说什么呢？\Quote{你们应记住主耶稣的话，他说过：\BibleQuote{施予比领受更为有福。}} \BibleRef{宗 20:35}那么，没有的人怎能给予呢？那么耶稣对玛尔大怎么说？\BibleQuote{你为许多事操心忙碌，但只有一件事是必要的，玛利亚选择了更好的一份}\BibleRef{路 10:41}？又：\BibleQuote{不要为明天忧虑}\BibleRef{玛 6:34}。因为现在必须解决所有这些疑问，不仅为了阻止那些懒惰的人，也为了使天主的圣言不显得自相矛盾。\par

保禄在另一处说：\BibleQuote{我们恳求你们，弟兄们，你们要更加进步，立志过安静的生活，做自己的事，对外人行事端正。}\BibleRef{得前 4:10}；又说：\BibleQuote{那偷窃的，不要再偷，却要劳苦，亲手做正经事，好使他能周济有需要的人。}\BibleRef{弗 4:28}。在这里，宗徒不仅命令\Quote{工作}，而且工作如此积极和勤劳，以至能有多余的给别人。在另一处，同一位又说：\BibleQuote{这双手供应了我的需要，也供应了同我在一起的人的需要。}\BibleRef{宗 20:34}。他在给格林多人写信时说：\BibleQuote{那么我的赏报是什么呢？就是我在传福音时，能使人免费得享福音。}\BibleRef{格前 9:18}。在格林多城，他与阿桂拉和普黎史拉同住，\BibleQuote{并工作，因为他们的职业是制帐篷的。}\BibleRef{宗 18:3}。\par

这些段落显示出在字句上更为明确的对照；因此我们现在必须提出解答。那么我们的回答应当是什么呢？那就是：\Quote{不要忧虑}的意思并非\Quote{不工作}，而是\Quote{不被钉在今世的事物上}；也就是说，不为明天的安逸操心，而视之为多余。因为一个人可以不工作，却（仍然）为明天积蓄财宝；另一个人可以工作，却一无所虑；因为忧虑和工作不是同一回事；人工作，并不是依靠自己的工作，而是\Quote{为了能够分施给有需要的人。}而且，对玛尔大所说的话也不是指工作和劳作，而是指这一点：我们有责任知道合宜的时机，不把用来聆听的时间花费在属肉体的事上。因此基督说这些话并非敦促她\Quote{懒惰}，而是要她专注于聆听。祂说：\Quote{我来，是为教导你们必需的事，而你却为饭食焦虑。你愿意接待我，为我预备丰盛的筵席吗？请预备另一种款待：通过乐意聆听，并效法你妹妹对教义的渴慕。}祂说这话并不是禁止她的款待（绝无此意！怎么可能呢？），而是表明她在聆听的时节不应忙碌于其他事务。因为说：\Quote{不要为那必朽坏的食物劳作}，并不是暗示我们应当懒惰；（事实上，这尤其就是\Quote{必朽坏的食物}，因为懒惰常教导一切邪恶；）而是说我们应当工作，并分施。这就是永不朽坏的食物；但如果有人懒惰、贪食、关心奢侈，那么他就是在为\Quote{那必朽坏的食物}劳作。同样，如果有人通过他的劳作喂养基督、给祂喝、给祂穿，那么有谁会如此愚昧和疯狂，说这样的人为那必朽坏的食物劳作呢？因为为此有将来天国的应许和那些美善之物。这食物永远长存。但那时，因为群众不关心污秽，也不寻求学习是谁行了这些事，凭着什么权能，只渴望一件事，就是不劳而获地填饱肚子；基督有充分的理由称这样的食物为\Quote{必朽坏的食物}。祂说：\Quote{我喂养了你们的身体，为叫你们此后寻求那另一份长存的食物，那滋养灵魂的食物；但你们又追逐那属地上的事。所以你们不明白，我引导你们不是到这不完全的食物，而是到那赐予不是暂时而是永生、不是滋养身体而是滋养灵魂的食物。}然后，当祂说了如此多关于自己的伟大话语，并说祂要赐下这食物时，为了使所说的话不成为他们的绊脚石，为了使祂的话可信，祂将此供应归于圣父。因为说了\Quote{人子必要赐给你们的}之后，祂接着说：\Quote{因为天主圣父已印证了祂}，也就是说，\Quote{派遣祂为此目的，使祂能把食物带给我们。}这话也有另一种解释；因为在另一处基督说：\Quote{那听我话的，就证实天主是真实的}\BibleRef{若 3:33}，也就是说，\Quote{不容置疑地显明了。}在我看来，这话似乎也暗示了此处，\Quote{因为圣父已印证了}不是别的，正是\Quote{已宣告了}、\Quote{已藉着祂的见证启示了。}事实上，祂自己也宣告了，但因为祂是对犹太人说话，所以提出了圣父的见证。\par

2. 那么，亲爱的诸位，让我们学习向天主祈求那些我们应当向祂祈求的事。因为那些属于今生的事，无论它们如何发展，都不能伤害我们；因为如果我们富有，我们只在此世享受奢华；如果我们陷入贫穷，我们也不会遭受什么可怕的事。因为今生无论荣华还是痛苦，在令人沮丧或欢乐方面都没有多大力量；它们是可轻视的，并且很快就消逝。因此它们被有理由地称为\Quote{一条路}，因为它们经过，并且按其本性不会长久存留；但将来的事，无论是惩罚的还是天国的，都永远长存。那么，让我们在这些事上殷勤努力，避免前者而选择后者。因为这世界的奢华有什么益处呢？今天在，明天就不在了；今天是一朵鲜艳的花，明天是散落的尘土；今天是一团燃烧的火，明天是冒烟的灰烬。但属灵的事却不是这样：它们永远闪耀、兴盛，并且日益光辉。那财富永不朽坏，永不离去，永不停止，不带来忧虑、嫉妒或责备，不毁坏身体，不腐蚀灵魂，没有恶意，不堆积仇恨；而这一切都伴随着另一种财富。那尊荣不使人陷入愚妄，不使人自高自大，永不停止，永不暗淡。再者，天上的安息和喜乐持续不断，永远稳固不朽，人找不到它的尽头和界限。那么让我们渴望这生命吧，因为如果我们这样行，就不会看重现在的事物，反而藐视和嘲笑它们的一切；即使有人邀请我们进入王宫，我们若怀有这盼望，也不会选择去行；然而世上似乎没有什么比这样的许可更接近幸福；但对于那些被对天国的爱所占据的人，甚至连这也显得微小、卑贱，不值一提。任何有终结的事都不应太被渴慕；一切终止的东西，今天在，明天就不在，即使它非常大，也显得非常微小和可轻视。那么，让我们不要依附于那些转瞬即逝、溜走和离去的事物，而要依附于那些持久和不动摇的事物。愿我们众人都能达到这境界，藉着我们主耶稣基督的恩宠和慈爱，藉着祂和与祂，圣父及圣神，得享光荣，从今世直到永远，世世无尽。阿们。\par

\SourceRecord{Translated by Charles Marriott.；Nicene and Post-Nicene Fathers, First Series；Vol. 14.；Edited by Philip Schaff.；Buffalo, NY: Christian Literature Publishing Co.,；1889.；<http://www.newadvent.org/fathers/240144.htm>.}

% structure: p0001 source=h1 target=section reason=page-title

\section{若望福音讲道 45}

% structure: p0002 source=h2 target=subsection reason=relative-heading-rank; base=h2

\subsection{若 6:28-30}

1. 没有比贪饕更糟糕、更可耻的了；它使心智迟钝，使灵魂属血肉；它蒙蔽人，不容人看清。请看，犹太人就是如此；因为他们贪图口腹，完全沉溺于世俗之事，毫无属灵的思想，所以虽然基督用无数言论引导他们，既尖锐又宽容，他们仍然不站起来，而继续匍匐于下。请想：祂对他们说：\BibleQuote{你们找我，并不是因为见了奇迹，而是因为吃了饼而吃饱了}；祂用责备触动了他们，指示他们应当寻求什么食物，说：\BibleQuote{不要为那必坏的食物劳力}；祂将奖赏放在他们面前，说：\BibleQuote{要为那存留到永生的食物劳力}；然后又对可能有的反对提出了补救，宣告自己是圣父所派遣的。\par

他们怎么做了呢？好像什么也没有听见，他们说：\BibleQuote{我们该做什么，才能做天主的工作？} 他们这样说，并不是为了学习并履行它们（如后文所示），而是为了诱导祂再次供给他们食物，并渴望说服祂满足他们。那么基督说什么呢？\BibleQuote{天主的工作就是：你们要信祂所派遣的那位。} 于是他们问：\BibleQuote{你显什么奇迹给我们看，好让我们看见而信你呢？}\par

% structure: p0005 source=h2 target=subsection reason=relative-heading-rank; base=h2

\subsection{若 6:31}

没有比这些人更无知、更无理的了！奇迹还在他们手中，就好像没有发生过一样，他们这样说：\Quote{你显什么奇迹？} 说了这话，他们甚至不让祂选择奇迹的权利，以为要迫使祂只显一个像他们祖先时代所行的那样的奇迹；因此他们说：\Quote{我们的祖先在旷野吃过玛纳}，以为这样能激起祂行一个为他们提供肉体营养的奇迹。否则，为什么他们不提古代的其他奇迹，尽管那些时期发生了许多奇迹，无论是在埃及、在海边还是在旷野，却只提玛纳呢？岂不是因为他们被肚腹的暴政所支配，极其渴望那一个吗？你们这些人，在看见祂的奇迹时称祂为先知，并试图立祂为王，怎么现在就像什么奇迹都没有发生过一样，变成了忘恩负义、居心不良的人，求一个奇迹，说出适合寄生虫或饿狗的话呢？难道现在玛纳对你们来说就显得神奇吗？你们的灵魂现在并不干渴。\par

也要注意他们的伪善。他们没有说：\Quote{梅瑟行了这奇迹，你行什么呢？} 以为这会惹恼祂；但有一段时间，他们因为期待食物，就以极大的尊敬对祂说话。所以他们既没有说：\Quote{天主行了这事，你行什么呢？} 以免显得使祂与天主同等；也没有提梅瑟，以免显得贬低祂，而是以一种中间的形式提出问题：\Quote{我们的祖先在旷野吃过玛纳}。祂本来可以回答说：\Quote{我现在行了比梅瑟更大的奇迹，不需要棍杖，不需要祈祷，全都靠我自己；如果你们回想玛纳，看，我给了你们饼。} 但这不是说这些话的时候；祂热切渴望的一件事，就是引导他们得到属灵的食物。请看祂的无限智慧和祂回答的方式。\par

% structure: p0008 source=h2 target=subsection reason=relative-heading-rank; base=h2

\subsection{若 6:32}

为什么祂没有说：\Quote{那不是梅瑟赐给你们的，而是我}，却把天主放在梅瑟的位置，把自己放在玛纳的位置？因为听众的软弱很大。从下文可以看出。因为即使祂这样说了，也没有吸引他们的注意，尽管祂起初说过：\BibleQuote{你们找我，并不是因为见了奇迹，而是因为吃了饼而吃饱了。}\BibleRef{若 6:26} 现在因为他们寻求这些（属血肉的）事，祂本想用随后的话纠正他们，但即使这样他们也没有停止。当祂向撒玛黎雅妇人许诺给她\Quote{水}时，祂没有提及圣父。祂说什么？\BibleQuote{如果你知道向你要水喝的是谁，你或许会求祂，祂必会给你活水}\BibleRef{若 4:10}；又说：\BibleQuote{我所赐的水}。祂没有把她指向圣父。但在这里祂提到了圣父，好让你明白撒玛黎雅妇人的信德何等大，而犹太人的软弱何等大。\par

那么，玛纳不是从天降下的吗？怎么又说它是从天降下的呢？这正如圣经说的：\BibleQuote{天上的鸟}\BibleRef{咏 8:8}；又说：\BibleQuote{上主从天上打雷}\BibleRef{咏 17:13}。祂称那另一种食物为\Quote{真正的食粮}，不是因为玛纳的奇迹是假的，而是因为它是预像，并非真理本身。但在提到梅瑟时，祂没有将自己与梅瑟相比，因为犹太人尚未认为祂高于梅瑟，他们仍对梅瑟怀有更高的敬意。所以，在说\Quote{不是梅瑟赐给}之后，祂没有接着说\Quote{是我赐给}，而是说圣父赐给，不是梅瑟。他们听见这话，回答说：\Quote{请把这食粮赐给我们吃吧}；因为他们仍然以为那是物质的东西，仍然期望满足自己的食欲，所以急忙跑到祂跟前。基督做什么呢？祂一步一步地引导他们，说：\par

% structure: p0011 source=h2 target=subsection reason=relative-heading-rank; base=h2

\subsection{若 6:33}

祂说：不是单为犹太人，而是为整个\Quote{世界}；不仅是食物，而是\Quote{生命}，另一种不同的\Quote{生命}。祂称它为\Quote{生命}，因为他们都死在罪中。但他们仍然俯首向下，说：\par

% structure: p0013 source=h2 target=subsection reason=relative-heading-rank; base=h2

\subsection{若 6:34}

然后祂为了责备他们，因为当他们以为食物是物质时就跑向祂，但得知是属灵的种类时却不然，便说：\par

% structure: p0015 source=h2 target=subsection reason=relative-heading-rank; base=h2

\subsection{若 6:35–36}

2. 因此，若翰也预先喊道：\Quote{祂所说的，祂是知道的；祂所见证的，是祂见过的，然而没有人接受祂的见证。}\BibleRef{若 3:32}；而基督自己又说：\Quote{我们所说的，是我们知道的；我们见证的，是我们见过的，}\BibleRef{若 3:11}，\Quote{而你们却不相信。}祂这样做，是为劝阻他们，并向他们表明，这事并不使祂困扰，祂不求荣耀，祂并非不知道他们心中的秘密，也不知道现在和将来的事。\par

\Quote{我是生命的食粮。}现在祂开始将奥秘托付给他们。首先祂论述祂的神性，说道：\Quote{我是生命的食粮。}因为这话不是论及祂的身体（关于身体，祂在末尾说：\Quote{我所要赐给的食粮，就是我的肉}），而是指祂的神性。因为那通过天主圣言而成的，是食粮，正如这饼通过降临于其上的圣神，成了天上的食粮。这里祂没有像先前讲话那样使用见证，因为祂有增饼奇迹为祂作证，而犹太人自己也曾一度假装相信祂；先前他们反对并控告祂。这就是为何在这里祂宣告自己。但那些人，由于期望享受一场肉性的盛宴，直到他们放弃希望时才受到搅扰。然而基督并不因此沉默，而是说了许多斥责的话。因为那些一边吃一边称祂为先知的人，在这里却感到冒犯，称祂为木匠的儿子；但吃饼的时候却不是这样——那时他们说：\Quote{祂真是那位先知}，并想要立祂为王。现在他们似乎因祂宣称\Quote{从天上降下来}而愤慨，但事实上，引起他们愤慨的并非这个，而是想到他们将不会享受物质的宴席。假如他们真的愤慨，他们本该询问并探究祂如何是\Quote{生命的食粮}，如何\Quote{从天上降下来}；但他们没有这样做，反而窃窃私语。而从另一件事可以清楚看出，冒犯他们的不是这个。当祂说\Quote{我父把食粮赐给你们}时，他们并没有喊\Quote{求祂赐下}，而是说什么？\Quote{请把那个食粮赐给我们}；然而祂并没有说\Quote{我赐下}，而是说\Quote{我父赐下}；尽管如此，他们因对食物的渴望，认为祂值得信任以供应食物。那么，那些认为祂在赐予食物上值得信任的人，怎么后来听到\Quote{父赐下}时反而感到冒犯呢？原因何在？这是因为他们听到自己将不再吃，就又不信了，并以\Quote{这话难懂}为不信的借口。因此祂说：\Quote{你们看见了我，还是不信}\BibleRef{若 5:39}；这部分是指祂的奇迹，部分是指圣经的见证；\Quote{因为为我作证的，}祂说，\Quote{就是这些经书}\EditorialNote{参 5:43-44}；以及\Quote{我因我父的名而来，你们却不接纳我}；以及\Quote{你们既然彼此求光荣，怎能相信呢？}\par

% structure: p0018 source=h2 target=subsection reason=relative-heading-rank; base=h2

\subsection{若望福音 6:37}

注意祂如何为那些得救的人做一切事；因此祂加上这话，为显得不是随便说说、毫无目的地讲这些话。但祂说的是什么？\Quote{凡父赐给我的人，必到我这里来}\BibleRef{若 6:37}，以及\Quote{在末日我要叫他复活}\BibleRef{若 6:40}？为何祂提到普遍的复活（就连不敬虔的人也有分），好像这是信祂之人的特殊恩赐？因为祂不是单纯地讲复活，而是讲一种特殊的复活。因为祂先说了\Quote{我不丢弃他，也不失落其中的任何}，然后才讲复活。因为在复活中，有些人被丢在外面（\Quote{把他丢在外面的黑暗中}\BibleRef{玛 22:13}），有些人被灭亡（\Quote{更当怕那能连灵魂和身体一并毁灭在地狱里的}\BibleRef{玛 10:28}）。而\Quote{我赐给他们永生}\BibleRef{若 10:28}这话也表明了这一点；因为\Quote{作恶的，复活起来受审；行善的，复活起来得生}\BibleRef{若 5:29}。因此，这里所指的，正是那得永生的复活。但祂说\Quote{凡父赐给我的人，必到我这里来}是什么意思？祂触到了他们的不信，指出凡不信祂的人就是违背了父的旨意。因此祂并非直截了当地说，而是以隐蔽的方式，这是祂一贯的做法，为要表明不信的人是与父不和，而不仅仅与祂不和。因为如果这是父的旨意，如果祂为此而来，为拯救世人，那么不信的人就是违背了父的旨意。\Quote{因此，}祂说，\Quote{凡父引导的人，没有什么能阻止他到我这里来}；另一处又说：\Quote{除非父吸引，谁也不能到我这里来}\BibleRef{若 6:44}。保禄说，祂将他们交付于父：\Quote{那时，基督将王国交付给父}\BibleRef{格前 15:24}。正如父给予时并非先使自己贫乏，同样，圣子交付时也不排除自己。祂被说成是交付我们，因为通过祂我们得以接近（父）。\par

3. \Quote{借着谁}一词也应用于父，就如宗徒所说：\Quote{你们为祂所召，是共享祂圣子的}\BibleRef{格前 1:9}；以及\Quote{借着父的旨意}。又说：\Quote{西满巴尔多禄茂，你是有福的，因为不是血肉启示了你}\BibleRef{玛 16:17}。祂在此暗示的是这样的事：\Quote{信德并不是寻常的事，而是需要来自上界的推动}；祂在整个讲论中都确立了这一点，表明这信德需要一种高贵的灵魂，并且是受天主牵引的。\par

\Quote{凡圣父所赐给我的，和凡祂所吸引的人，都必到我这里来；若没有人能到我这里来，除非是从上头赐给他的，那么圣父所不赐给的人，就无可指责或控告了。} 这些不过是空洞的言语和借口。因为我们也需要本身的自主选择，因为我们是否愿意受教是一个选择的问题，我们是否愿意相信也是如此。在此处，藉着\Quote{圣父所赐给我的人}，祂所要表达的不过是\Quote{信我并不是一件平常的事，也不是出于人的理智，而是需要来自上头的启示，以及一个井然有序的灵魂来接受那启示。} 而\Quote{到我这里来的，必得救}，意思是他将受到极大的眷顾。\Quote{因为为了这些，}祂说，\Quote{我来了，并取了肉身，取了奴仆的形体。} 然后祂接着说：\par

% structure: p0022 source=h2 target=subsection reason=relative-heading-rank; base=h2

\subsection{\BibleRef{若 6:38}}

祢说什么？怎么，祢的旨意是一个，而祂的又是另一个吗？为了没有人怀疑这一点，祂藉着下面的话解释，说：\par

% structure: p0024 source=h2 target=subsection reason=relative-heading-rank; base=h2

\subsection{\BibleRef{若 6:40}}

那么，这不是祢的旨意吗？祢怎么说：\Quote{我来要把火投在地上，我多么渴望它已经点燃了}？\BibleRef{路 12:49} 因为如果祢也愿意这样，那么很明显，祢的旨意与圣父的旨意是同一个。在另一处祂也说：\BibleQuote{因为父怎样叫死人复活，赐他们生命，子也照样随祂的意愿赐生命。}\BibleRef{若 5:21} 但是圣父的旨意是什么？难道不是不让其中任何一个丧亡吗？祢也愿意这样。\BibleRef{玛 18:14} 因此，一者的旨意与另一者的旨意没有不同。同样，在另一处，祂被看见更坚固地确立祂与圣父的平等，说：\Quote{我与我的圣父\InnerQuote{要来，并要在他那里作我们的住所。}}\BibleRef{若 14:23} 那么祂所说的是这样：\Quote{我来不是要行任何不同于圣父所愿意的事，我没有自己单独的旨意与圣父的不同，因为凡圣父所有的都是我的，凡我所有的也都是圣父的。} 既然圣父与子的一切都是共通的，祂有理由说：\Quote{不是要行我自己的旨意。} 但在这里祂不是这样说的，而是把这一点留到最后。因为，如我所说，祂暂时隐藏和遮盖了崇高的事，并想证明：即使祂说\Quote{这是我的旨意}，他们也会轻视祂。因此祂说\Quote{我合作于那旨意}，想要这样更使他们震惊；好像祂说过：\Quote{你们以为如何？你们因不信而激怒我吗？不，你们激怒了我的圣父。} \BibleQuote{因为这是派遣我来的那位的旨意：凡祂赐给我的，我应不失落其中之一。}\BibleRef{若 6:39} 这里祂显示祂不需要他们的服务，祂来不是为了自己的利益，而是为了他们的救恩；也不是为了从他们那里得到荣耀。这祂确实在先前的话中宣告过，说：\BibleQuote{我不接受从人来的荣耀}\BibleRef{若 5:41}；又说：\BibleQuote{我说这些话，是要你们得救。}\BibleRef{若 5:34} 因为祂处处努力说服他们，祂来是为了他们的救恩。祂说，祂使圣父得荣耀，为的是不使他们怀疑祂。而祂之所以这样说，祂藉着下面的话更清楚地启示了。因为祂说：\BibleQuote{那寻求自己旨意的，是寻求自己的荣耀；但那寻求那派遣祂者的荣耀的，是真实的，并且祂里面没有不义。}\BibleRef{若 7:18} \BibleQuote{这就是父的旨意：凡看见子并相信祂的，得享永生。}\BibleRef{若 6:40}\par

\Quote{在末日我要使他复活。} 为什么祂不断谈论复活？是为了使人不单凭现世事物判断天主的眷顾；使他们即使在此世未享成果，也不因此沮丧，而是等待将来的事；并使他们不因现时罪未受罚而轻视祂，而是盼望另一个生命。\par

如今那些人一无所获，但让我们努力获得益处，经由使复活不断在我们耳边回响；如果我们想要贪婪、或偷窃、或做任何错事，让我们立刻将那日子放在心上，让我们想象审判座，因为这样的反思比任何嚼环更能遏制邪恶冲动。让我们不断对别人和自己说：\Quote{有复活，还有一个可怕的审判座在等候我们。}如果看到任何人傲慢且因今世的好事而自大，让我们也对他说同样的话，并向他表明所有那些事都停留在此；如果我们看到另一个人悲伤且不耐烦，让我们也对他说同样的话，并向他指出他的悲伤将有终结；如果我们看到一个人粗心大意且放荡，让我们对他念这同样的咒语，并表明他必须为他的粗心交账。这句话比任何其他疗法更能治愈我们的灵魂。因为有复活，并且那复活就在我们门前，不是遥远，也不是远距离。\BibleQuote{因为还有一点点时候，那要来的就来，并不迟延。}\BibleRef{希 10:37}又说：\BibleQuote{我们众人必须出现在基督的审判座前}\BibleRef{格后 5:10}；也就是说，坏的和好的，前者在众人面前蒙羞，后者在众人面前更显光荣。因为正如在此审判的人公开惩罚恶人、尊荣好人，在那里也是如此，使前者更加羞愧，后者更加显赫的光荣。让我们每天想象这些事。如果我们常常回味这些，今世的事就无法刺痛我们。\BibleQuote{因为所见的是暂时的，所不见的是永远的。}\BibleRef{格后 4:18}让我们不断对自己和别人说：\Quote{有复活，有审判，有对我们行为的审查}；让许多认为有命运存在的人也重复这话，他们立刻就会从他们疾病的黑烂中被解救出来；因为如果有复活和审判，就没有命运，即使他们提出千万个论据，并极力证明它。但我羞于教导基督徒关于复活的事：因为需要学习有复活，并且没有坚定地相信这个世界的事情不是由命运、没有计划、随意运行的人，不能是基督徒。因此，我劝勉并恳求你们，我们洁净自己脱离一切邪恶，尽我们所能在那一天获得宽恕和赦免。\par

也许有人会说：\Quote{何时才有终结？何时才有复活？看，已经过了这么长时间，这类事一件也没有发生？}但它必定会发生，你尽管放心。因为洪水以前的人也是如此说，嘲笑\Person{诺厄}，但洪水来了，冲走了所有那些不信的人，却保全了相信他的人。\Person{罗特}时代的人也没有料到天主的那一击，直到那些闪电和雷霆降下，将他们全部毁灭。无论是这些人，还是\Person{诺厄}时代的人，对将要发生的事都没有任何预兆，但当他们全都尽情享受，饮酒狂欢，疯狂醉倒时，这些难以忍受的灾祸就临到了他们身上。复活也是如此；没有任何预兆，而是在我们处于好时光的时候。因此，\Person{保禄}说：\BibleQuote{因为当他们说：\InnerQuote{平安稳妥}的时候，毁灭忽然临到他们，如同产痛临到怀胎的妇人；他们绝不能逃脱。}\BibleRef{得前 5:3}天主这样安排，使我们常常奋斗，甚至在平安时也不自信。你说什么？你难道不期望有复活和审判吗？魔鬼也承认这些，你竟不羞愧？他们说：\Quote{你来到，}他们说，\Quote{要在时候以前折磨我们吗？}\BibleRef{玛 8:29}如今他们说着将有\Quote{折磨}，就知道审判、清算和惩罚。因此，我们除了胆敢作恶之外，不要再因不信复活的话而激怒天主。因为正如在其他事上基督是我们的元始，在这事上也是如此；为此祂被称为\BibleQuote{死者中的首生者。}\BibleRef{哥 1:18}如果没有复活，祂怎能称为\Quote{首生者，}而\Quote{死者}中没有一个要跟随祂呢？如果没有复活，天主的公义如何保存，当这么多恶人亨通，这么多好人受苦并死于苦难之中？如果没有复活，这些人各得其所应得的在哪里？没有一个生活正直的人不信复活，反而他们每天祈祷并重复那句圣言：\BibleQuote{愿祢的国来临。}那么，不信复活的是谁呢？是那些走不圣洁之路、过不洁生活的人；正如先知所说：\BibleQuote{他的道路时常污秽；你的判断远在他眼前。}\BibleRef{咏 9:5}因为一个人不可能不信复活而过纯洁的生活；因为那些自觉无罪的人既谈论它，又渴望它，又相信它，为要得到他们的报答。那么，我们不要再激怒祂，而要听祂说：\BibleQuote{要害怕那能把身体和灵魂都投入地狱的祂}\BibleRef{玛 10:28}；好藉那恐惧我们得以改善，并被从那丧亡中解救出来，堪当获得天国。但愿我们众人，赖我们的主耶稣基督的恩宠和慈爱，都能达到天国；藉着祂，偕同祂，圣父和圣神，永受光荣，自今至于永远，及世之世。阿们。\par

\SourceRecord{Translated by Charles Marriott.；Nicene and Post-Nicene Fathers, First Series；Vol. 14.；Edited by Philip Schaff.；Buffalo, NY: Christian Literature Publishing Co.,；1889.；<http://www.newadvent.org/fathers/240145.htm>.}

% structure: p0001 source=h1 target=section reason=page-title

\section{若望福音讲道集 第46篇}

% structure: p0002 source=h2 target=subsection reason=relative-heading-rank; base=h2

\subsection{若望福音 6:41-42}

1. \BibleQuote{他们的天主是肚腹，并以自己的羞辱为光荣}\BibleRef{斐 3:19}，保禄在致斐理伯人书中论到某些人如是说。现在犹太人显然具有这种性格，从前面的事和他们来对基督说的话都可以看出。因为当祂赐给他们饼，填饱他们的肚腹时，他们说祂是先知，并想立祂为王；但当祂教导他们关于灵性之粮、关于永生，带领他们远离感官之物，对他们谈论复活，提升他们的思想到更高的事理时，在他们最该钦佩的时候，他们却抱怨而退却。然而，如果祂是他们先前所宣称的那位先知，声称梅瑟所说\BibleQuote{主你的天主要从你们弟兄中给你们兴起一位像我一样的先知，你们应听从祂}\BibleRef{申 18:15}的那位，那么当祂说\Quote{我从天降下}时，他们本该听从祂；但他们却不听从，反而抱怨。他们仍然尊敬祂，因为饼的奇迹刚发生，所以他们没有公开反驳祂，而是通过抱怨表达不悦，因为祂没有给他们想要的饭食。他们抱怨说：\Quote{这不是若瑟的儿子吗？}由此明显看出，他们尚不知道祂奇妙的、令人惊叹的出生。所以他们仍然说祂是若瑟的儿子，而祂没有被责备；祂也没有对他们说：\Quote{我不是若瑟的儿子}；不是因为祂是他的儿子，而是因为他们还无法承受那奇妙降生的奥迹。如果他们不能承受被告知祂按肉身的出生，那么他们就更不能承受那从上面来的、不可言说的出生。如果祂没有向他们启示那较低的事，那么祂更不会将那更高的事托付给他们。虽然这一点大大冒犯了他们，即祂出生于一个卑微平凡的父，但祂仍然没有向他们启示真理，免得在消除一个冒犯的原因时，又制造另一个。那么，当他们抱怨时，祂说了什么呢？\par

% structure: p0004 source=h2 target=subsection reason=relative-heading-rank; base=h2

\subsection{若望福音 6:44}

摩尼派抓住这些话，说\Quote{任何事都不在我们自己的能力之内}；然而，这句话表明我们是我们意志的主人。有人会说：\Quote{人若到祂那里去，还需要什么吸引呢？}但这话并没有取消我们的自由意志，而是表明我们极需帮助。祂所指的不是一个不情愿的来者，而是一个得到许多帮助的人。然后祂也显示了祂吸引的方式；为避免人们又对天主形成任何物质性的观念，祂接着说：\par

% structure: p0006 source=h2 target=subsection reason=relative-heading-rank; base=h2

\subsection{若望福音 6:46}

有人问：\Quote{那么，圣父怎样吸引人呢？}那位先知古时早已解释，当祂预告并说：\par

% structure: p0008 source=h2 target=subsection reason=relative-heading-rank; base=h2

\subsection{若望福音 6:45}

你看到信德的尊贵了吗？这并非由人，也不是藉着人，而是由天主亲自教导他们。为使这说法可信，祂引用了他们的先知的话。\Quote{如果\InnerQuote{众人都要蒙天主教诲}，那么为什么有些人不会相信呢？}因为这话是针对多数人说的。此外，预言并非绝对指所有人，而是指所有愿意的人。因为教师准备好向所有人传授他所拥有的，并向所有人倾注他的教训。\par

% structure: p0010 source=h2 target=subsection reason=relative-heading-rank; base=h2

\subsection{若望福音 6:44}

这里，圣子的权威并非微小，既然圣父引领，祂就举起。祂没有将祂的工作与圣父的工作区分开（那怎么可能呢？），而是显示能力的平等。因此，正如在另一处说了\Quote{那派遣我的圣父为我作证}之后，为避免他们对这话过于好奇，祂就引用了圣经；同样在这里，为避免他们产生类似的猜疑，祂引用了先知的话，祂不断到处引用他们，以表明祂与圣父并非对立。\par

有人问：\Quote{那么，那些在祂以前的人又怎样呢？难道他们不是受天主教导的吗？为什么这话特别应用在这里呢？}因为古时，他们是通过人手学习天主的事，但如今是通过天主的独生子和圣神。然后祂接着说：\Quote{不是有人看见过圣父，惟独那出于天主的}，这里祂用此说法不是指原因，而是指存在的样式。如果祂是从前一种意义说的，那么我们都是\Quote{出于天主的}。那么圣子的独特本性又在哪里呢？有人问：\Quote{为什么祂没有更清晰地说明这一点呢？}因为他们的软弱。因为当祂说\Quote{我从天降下}时，他们已如此反感，如果祂加上这个，他们会怎样呢？\par

% structure: p0013 source=h2 target=subsection reason=relative-heading-rank; base=h2

\subsection{若望福音 6:48}

祂称自己为\Quote{生命之粮}，因为祂维持我们的生命，无论是现在的还是将来的，并且说：\Quote{谁吃了这粮，将永远活着。}这里\Quote{粮}可以指祂救恩的教义和对于祂的信德，也可以指祂的圣体；因为两者都加强灵魂。但在另一处祂说：\Quote{谁若听从我的话，永远见不到死亡。}\BibleRef{若 8:51}那时他们起了反感；但在这里他们也许没有这样的感觉，因为由于所行的饼的奇迹，他们仍然尊敬祂。\par

2. 请注意祂如何通过让他们听到每种食物的结果，来区分祂的粮和玛纳。为表明玛纳并没有给他们带来不寻常的益处，祂接着说：\par

% structure: p0016 source=h2 target=subsection reason=relative-heading-rank; base=h2

\subsection{若望福音 6:49}

然后祂确立了一件最可能说服他们的事，即他们被认为比他们的祖先（指梅瑟时代那些可敬的人）更配得更大的事，所以说了吃过玛纳的人都死了之后，祂接着说：\par

第51节。\BibleQuote{谁吃这粮，将永远活着。}\BibleRef{若 6:51}\par

祂提到\Quote{在旷野}并非无缘故，而是为了指出玛纳的供应并未持续很长时间，也没有随他们进入应许之地。但这\Quote{粮}却不同。\par

\Quote{我所赐的粮就是我的肉，我必要为世界的生命而赐下。}\par

这里有人可能会合理地问，这些话语既不能建立也不能获益，反而对已经建立的人造成伤害，这怎么是一个合适的时候呢？因为福音史家说：\Quote{从此，祂的许多门徒退去了}，说：\Quote{这话生硬，谁能听得下去呢？}\BibleRef{若 6:60}；因为这些事情本可以只托付给门徒，正如玛窦告诉我们的，祂单独与他们谈论。\BibleRef{谷 4:34}那么，我们该说什么呢？这些话有什么益处呢？益处和必要性是大的。因为他们逼迫祂，求肉身的食物，提醒祂先祖们日子的食物，并提到玛纳为大事，为要给他们显示，那些事不过是预像和影子，而事情的实体现在与他们同在，祂就提到属灵的食物。\Quote{但是，}有人说，\Quote{祂本应该说：你们的祖先在旷野吃过玛纳，但我给了你们饼。}但两个奇迹之间的间隔很大，后者似乎比前者逊色，因为玛纳是从天上降下的，而这饼的奇迹却是在地上行的。因此，当他们寻求\Quote{从天上降下的}食物时，祂不断地告诉他们：\Quote{我从天上降下。}若有人问祂为什么引入关于\OriginalTerm{奥迹}{Mysteries}的谈论，我们要回答说，这是谈论这些事情非常合适的时候；因为话语的含混总是激发听者，使他更专注。他们本不应被绊倒，反而应该询问和探究。但现在他们退去了。如果他们相信祂是先知，就该相信祂的话，所以绊倒是由于他们自己的愚昧，而不是话语的任何困难。且看祂如何一步步引导他们到自己这里来。在这里祂说是祂自己给予，不是父；\Quote{我所要赐给的饼，就是我的肉，是为世界的生命而赐给的。}\par

\Quote{但是，}有人说，\Quote{这道理对他们来说是陌生而不寻常的。}然而若望在更早的时候已经藉称呼祂为\Quote{羔羊}暗示了这一点。\BibleRef{若 1:29}\Quote{尽管如此，他们仍不知道。}我知道他们不知道；不，连门徒也不明白。因为如果他们尚不清楚复活，因此不知道\Quote{拆毁这座圣殿，三天内我要把它重建起来}\BibleRef{若 2:19}可能是什么意思，那么对这些话就更会一无所知了。因为这些话不如那些清楚。既然先知曾从死人中复活，他们知道，即使圣经没有如此清楚地谈论这一点，但从来没有一个人宣称有人吃过肉。然而他们仍然服从，跟随祂，并承认祂有永生的话。因为这是门徒的本分，不是过分好奇地探究老师的断言，而是听从并服从他，等待合适的时间来解答任何难题。\Quote{那么，}有人说，\Quote{为什么事情相反，这些人\InnerQuote{退去了}？}这是由于他们的愚昧。因为当对\Quote{如何}的追问进来时，不信也随之进来。尼苛德摩也是如此困惑，说：\Quote{人怎能再入母腹？}他们也是这样慌乱，说：\par

第52节。\Quote{这人怎能把自己的肉给我们吃呢？}\BibleRef{若 6:52}\par

如果你要问\Quote{如何}，为什么在饼的事上你不问这个，祂如何把五个饼扩展到如此大的数量？因为那时他们只想到吃饱，而不是看奇迹。\Quote{但是，}有人说，\Quote{他们当时的经验教导了他们。}那么，因着那经验，这些话本该容易被接受。为此祂预先行了那个奇异的奇迹，为使他们因之受教，不再不信祂以后所说的一切。\par

3. 那些人当时从所说的话中并未收获果实，但我们却在实际体验中得到了益处。因此，有必要理解这\OriginalTerm{奥迹}{Mysteries}的奇妙：它是什么，为何赐下，以及这行动的利益何在。我们成为一个身体，\BibleQuote{是祂肉和祂骨的肢体}\BibleRef{弗 5:30}。让已入门的人追随我所说的。为使我们不仅藉着爱，更在实际中成为这样，让我们融和于那肉中。这是藉着祂白白赐给我们的食物而成就的，祂愿意显示祂对我们的爱。因此祂将自己与我们混合；祂将祂的身体与我们的揉和，使我们成为一个整体，如同身体与头相连。因为这属于强烈相爱的人；例如，\Person{约伯}谈及他的仆人们，他们极其爱他，以致渴望与他血肉相连。因为他们说，为显示他们所感受到的强烈爱情：\BibleQuote{谁能让我们饱尝他的肉？}\BibleRef{约 31:31}。因此，基督也做了这事，为领我们进入更亲密的友谊，并显示祂对我们的爱；祂赐给那些渴望祂的人不仅看见祂，甚至触摸祂、吃祂、以牙齿咬祂的肉、拥抱祂，并满足他们所有的爱。让我们从那祭台如喷火的狮子般回来，成为魔鬼的恐怖；默想我们的元首，以及祂向我们显示的爱。父母常常将子女托付给别人喂养；\Quote{但我不这样}，祂说，\Quote{我用我自己的肉喂养你们，愿意你们全都生来高贵，并给你们将来美好的希望。因为那在这里把自己给你们的，在将来更会这样做。我甘愿成为你们的兄弟，为了你们我分享了血肉，反过来我又把那使我成为你们亲属的血肉给了你们。}这血使我们心中常有我们君王的形象，产生说不出的美丽，不容我们灵魂的高贵消逝，持续浇灌它，滋养它。由我们食物所生的血并非立即成为血，而是别的东西；但这血却不是如此，而是直接浇灌我们的灵魂，并在其中运行某种大能。这血，若正确领受，驱赶魔鬼，使他们远离我们，同时吸引天使和天使之主到我们这里。因为无论何处看见主的血，魔鬼就逃跑，天使们聚集。这血倾流，洗净了整个世界；有福的\Person{圣保禄}在\WorkTitle{希伯来书}中论及它时说了许多智慧的话。这血洁净了内室和至圣所。如果它的\OriginalTerm{预像}{type}在希伯来人的圣殿中，以及在埃及当中涂在门框上，具有如此大的力量，那么这实体更是如此。这血祝圣了金祭台；没有它，大司祭不敢进入内室。这血祝圣了司祭，这血在预像中洁净了罪恶。但如果它在预像中有如此力量，如果死亡在阴影前如此战栗，请告诉我，它岂能不畏惧实体本身吗？这血是我们灵魂的救恩，灵魂藉此被洗净，藉此成为美丽，藉此被点燃，这血使我们的悟性比火更明亮，使我们的灵魂比黄金更闪耀；这血倾流，使天堂变得可及。\par

4. 教会的\OriginalTerm{奥迹}{Mysteries}确实可畏，祭台也确实可畏。一道泉源从乐园涌出，发出物质的河流；而由这祭台涌出一股泉源，发出属灵的河流。在这泉源旁边栽种的，不是不结果子的柳树，而是直达高天的树木，常结时鲜而不朽的果实。若有人被热浪灼伤，让他来到这泉源旁，冷却他的灼热。因为它止息干旱，安慰一切被烧毁的，不是被太阳，而是被火箭。因为它始于高天，其源头就在那里，它的水也从那里流出。那\OriginalTerm{护慰者}{Comforter}所发出的泉源有很多支流，圣子是中保，不是拿着镐头开路，而是开启我们的心智。这泉源是光明的泉源，喷出真理的光芒。在它旁边站着高天万军，观看其支流的美丽，因为他们更清楚地领悟所陈列之事的权能，以及不可接近的闪耀。因为如同当金被熔化时，若有人（倘若可能）将手或舌浸入其中，它们立即变为金色；同样，但这里所陈列的以更大的程度运作于灵魂。这河比火更猛烈地沸腾，却不燃烧，只浸染它所触及的。这血自古以来就在义人的祭台和祭献中被预象，这是世界的代价，基督以此为自己购得了教会，以此完全装饰了她。因为如同一个人买仆役时给他们金子，而当他想装饰他们时也用金子；同样，基督用祂的血购买了我們，又用祂的血装饰了我们。分享这血的人与天使、总领天使和高天万军同列，身披基督自己的王袍，佩带圣神的武器。不，我还没说什么伟大的事：他们穿戴着王本身。\par

既然这是一件伟大而奇妙的事，你若以纯洁之心走近它，便是为得救恩而走近；但若怀着邪恶的良心，便是为受罚和报复而走近。\Quote{因为}，它说，\Quote{那不相称地吃喝}主的，\Quote{便是在吃喝自己的审判}\BibleRef{格前 11:29}；因为若那些玷污王袍的人与撕破王袍的人受同样的惩罚，那么那些以不洁的思想领受圣体的人，与那些用钉子撕裂圣体的人受同样的惩罚，就不足为奇了。至少请注意，\Person{保禄}宣布了何等可怕的惩罚，当他说：\Quote{那轻视\Person{梅瑟}法律的，凭两三个见证人，就不得怜恤而死；你们想，那践踏了天主子，并将那藉以成圣的盟约之血视为不洁之物的人，应受怎样更重的惩罚呢？}\BibleRef{希 1:28}那么，亲爱的，我们这些享有如此恩赐的人，要谨慎自己；若我们想说任何可耻的话，或发现自己被愤怒或任何类似的情欲所驱使，让我们想想我们曾被视为配得何等事，我们曾分享了何等伟大的圣神，这思想必会平息我们无理的情欲。我们还要被钉在现世的事物上多久呢？要到何时我们才醒悟呢？我们要忽略自己的救恩多久呢？让我们记住基督曾视我们为配得何等事，让我们感恩，让我们光荣祂，不仅借着我们的信德，也借着我们的行为，好使我们获得将来的福乐，借着我们的主耶稣基督的恩宠和慈爱，愿光荣归于祂偕同圣父及圣神，从今世直到永远。阿们。\par

\SourceRecord{Translated by Charles Marriott.；Nicene and Post-Nicene Fathers, First Series；Vol. 14.；Edited by Philip Schaff.；Buffalo, NY: Christian Literature Publishing Co.,；1889.；<http://www.newadvent.org/fathers/240146.htm>.}

% structure: p0001 source=h1 target=section reason=page-title

\section{《若望福音》讲道集第四十七篇}

% structure: p0002 source=h2 target=subsection reason=relative-heading-rank; base=h2

\subsection{\BibleRef{若 6:53-54}}

1. 当我们谈论属神之事时，愿我们灵魂中没有任何世俗之物，没有任何属地的思想，愿这一切思虑都退避、被驱散，让我们完全专注于聆听天主圣言。因为若在国王驾临时，一切纷扰都被驱散，更何况圣神与我们说话时，我们更需极大的宁静、极大的敬畏。今日所讲的是值得敬畏的。请听这是如何的。\BibleQuote{我实实在在告诉你们：你们若不吃人子的肉，不喝人子的血，在你们内便没有生命。} 既然犹太人先前断言这是不可能的，祂便表明不但不是不可能，而且是绝对必要的。因此祂接着说：\BibleQuote{谁吃我的肉，并喝我的血，必得永生。}\par

\BibleQuote{在末日我要使他复活。} 因为祂曾说过：\BibleQuote{谁吃这食粮，将永远不死} \BibleRef{若 6:50}，而这话很可能成为他们的障碍（正如他们先前所说：\BibleQuote{亚巴郎死了，先知们也死了，你怎么说：他决尝不到死味呢？} \BibleRef{若 8:52}），所以祂提出复活来解决这个问题，并表明（那吃祂的人）在末日不会死。祂不断谈论奥迹的主题，表明这行为的必要性，且必须无论如何要实行。\par

% structure: p0005 source=h2 target=subsection reason=relative-heading-rank; base=h2

\subsection{\BibleRef{若 6:55}}

祂说的是什么意思？祂或者要表明这是真正的食粮，能拯救灵魂；或者要证实所说的，使他们不以为这话只是谜语或比喻，而要知道必须吃这身体。然后祂说：\par

% structure: p0007 source=h2 target=subsection reason=relative-heading-rank; base=h2

\subsection{\BibleRef{若 6:56}}

祂说这话，是表明这样的人与祂结合。接下来似乎不连贯，除非我们探究其意义；因为有人说，在说了\BibleQuote{谁吃我的肉，住在我的内}之后，再加上这样的话，是什么关联呢？\par

% structure: p0009 source=h2 target=subsection reason=relative-heading-rank; base=h2

\subsection{\BibleRef{若 6:57}}

然而这些话完全和谐一致。因为祂常提到\Quote{永生}，为证明这一点，祂引入了\BibleQuote{住在我内}这句话；因为\Quote{如果他住在我内，而我活着，显然他也将活着。}然后祂说：\BibleQuote{就如那生活的父派遣了我}。这是比较和相似的说法，其意思是：\BibleQuote{正如父活着，我也活着}。 为免你以为祂是非受生的，祂立即补充说\BibleQuote{因父而活}，这并不是说祂需要某种运作的能力才能活着，因为祂先前已说过，为消除这种怀疑：\BibleQuote{就如父在自己内有生命，同样也赐给子在自己内有生命}；如果祂需要另一位的运作，那么要么父没有赐给祂如此，那么断言便是假的；要么父已如此赐给了，那么祂就不再需要另一位来支持祂。那么\Quote{因父而活}是什么意思？祂这里只是暗示原因，所说的是这样的：\BibleQuote{正如父活着，我也活着，那吃我的，也要因我而活着。} 祂所说的\Quote{生命}不仅是生命，而且是卓越的生命；因为祂不是单纯说生命，而是指那荣耀、不可言喻的生命，从这一点可以清楚看出。因为所有的人\Quote{活着}，即使是不信者、未领受奥迹者、不吃那肉的人。你看到这些话不是指这个生命，而是指另一个生命吗？祂所说的是这样的：\Quote{那吃我肉的人，当他死时，不会灭亡，也不会受惩罚}；祂说的不是普通的复活（因为所有人都同样复活），而是指特殊的、光荣的复活，那有赏报的复活。\par

% structure: p0011 source=h2 target=subsection reason=relative-heading-rank; base=h2

\subsection{\BibleRef{若 6:58}}

祂不断讲述同一要点，为将其铭刻在听众的理解中（因为关于这些点的教导是一种最终的教导），并巩固复活和永生的道理。因此祂提到复活，既然祂应许永生，表明那生命现在不是，而是在复活之后。\Quote{从何得知这些事呢？}有人说。\Quote{从圣经}；祂处处将犹太人引向圣经，命他们从圣经学习这些事。祂说\BibleQuote{给世界带来生命的}，以此激起他们的嫉妒心，使别人因嫉妒而享受恩赐，他们自己也不得落后。祂又不断提及玛纳，表明（它和祂的食粮之间的）差别，并引导他们走向信德；因为若祂能不用收割、谷粒或其他通常之物而供养他们的生命四十年，那么现在祂更能如此，因为祂是为了更大的目的而来的。而且，若那些只是预象，而人尚且能毫不费力地收集从天而降的东西，那么现在更是如此，因为在不死和享受真生命方面有极大的差别。祂多次说到\Quote{生命}，这是恰当的，因为这是人所渴望的，没有什么比不死更令他们喜悦。因为在旧约之下，应许的是长寿和许多日子，但现在不仅是长度，而是没有终结的生命。祂同时要表明，祂现在撤销了由罪带来的惩罚，废除了那定人死的判决，不仅要带来生命，而且是永生，与先前相反。\par

% structure: p0013 source=h2 target=subsection reason=relative-heading-rank; base=h2

\subsection{\BibleRef{若 6:59}}

2. 这是祂行了许多奇迹的地方，因此祂尤其应当在这里被听从。但为什么祂在会堂和圣殿里教导？一方面是因为祂想要吸引最多的人，另一方面是因为祂要表明祂与父并不对立。\par

% structure: p0015 source=h2 target=subsection reason=relative-heading-rank; base=h2

\subsection{\BibleRef{若 6:60}}

什么意思是\Quote{严厉}？严厉、劳苦、麻烦。然而祂并没有说这样的话，因祂不是谈论生活方式，而是谈论教义，不断地阐述那在祂内的信德。那么，\Quote{这话是严厉的}是什么意思呢？是因为它应许生命和复活吗？是因为祂说祂从天降下？还是因为不食祂肉的人不可能得救？告诉我，这些事是\Quote{严厉}的吗？谁能说是呢？那么，\Quote{严厉}是什么意思？意思是\Quote{难以接受}，\Quote{超越他们的软弱}，\Quote{带有很大的恐惧}。因为他们认为祂所说的话过于祂的真实身份，超过祂自己。所以他们说，\par

\Quote{谁能听呢？}\par

也许是他们为自己的离开找借口，因为他们即将离去。\par

% structure: p0019 source=h2 target=subsection reason=relative-heading-rank; base=h2

\subsection{\BibleRef{若 6:61-62}}

这也同\Person{纳塔乃耳}所做的一样，说：\Quote{因为我对你说：我看见你在无花果树下，你就信了吗？你要看见比这更大的事。}\BibleRef{若 1:50} 又对\Person{尼苛德摩}说：\Quote{没有人上过天，除了那自天降下而仍在天上的人子。}\BibleRef{若 3:13} 那么，祂是难上加难吗？不（这远非祂所愿），而是藉着教义之伟大和众多，祂想要赢得他们。因为如果有人简单地说：\Quote{我从天降下，} 却不再加添什么，他反而更容易使他们反感；但那说过\Quote{我的肉是世界的生命}，说过\Quote{就如生活的圣父派遣了我，我因圣父而生活}，又说\Quote{我从天降下}的祂，却解决了困难。因为一个人若只说自己一件伟大的事，也许会被怀疑是伪装，但连续说如此多的事，就消除了所有怀疑。祂所说所做的一切，都是为了引导他们离开\Person{若瑟}是祂父亲的念头。祂说这些，并非要加固这个绊脚石，而是除去它。因为凡认为祂是\Person{若瑟}之子的，不能接受祂的话，而那相信祂从天降下，并且要升到那里的，更容易听从祂的话：同时祂也提出了另一种解释，说：\par

% structure: p0021 source=h2 target=subsection reason=relative-heading-rank; base=h2

\subsection{\BibleRef{若 6:63}}

祂的意思是：\Quote{你们必须属神地听关于我的事，因为那属血肉地听的人，不得益处，也得不到任何好处。} 询问祂如何从天降下、认为祂是\Person{若瑟}的儿子、问\Quote{祂怎能把自己的肉给我们吃呢？}这些都是属血肉的。他们本该奥秘地、属神地理解这事。\Quote{但是，}有人说，\Quote{他们怎能明白\InnerQuote{吃肉}的意思呢？} 那时，他们理当等待适当的时候并询问，而不是离弃祂。\par

\Quote{我对你们所说的话，就是神，就是生命。}\par

那就是说，它们是神圣的、属神的，没有属血肉的成分，不受物理后果法则的约束，而是免于任何这种必然性，甚至超越为这世界所定的法则，并且还有另一种不同的意义。正如在此处祂说\Quote{神}，而不是\Quote{属神的}；同样，当祂说到\Quote{肉}时，祂的意思不是\Quote{属血肉的事物}，而是\Quote{以属血肉的方式听}，同时暗指他们，因为他们总是渴求属血肉的事物，而他们本应渴求属神的事物。因为若一个人以属血肉的方式接受它们，他就毫无益处。\Quote{那么，祂的肉不是肉吗？}确实是的。\Quote{那么，祂怎么说肉毫无益处呢？}祂不是指自己的肉（绝不！），而是指那些以属血肉的方式接受祂话的人。但是，什么是\Quote{以属血肉的方式理解}？它仅仅是看我们眼前的，而不想象任何超出的事。这就是属血肉的理解。但我们不可这样凭外表判断，而要用内在的眼睛察看一切奥秘。这就是以属神的方式看。那不吃祂的肉、不喝祂的血的人，在他内没有生命。那么，如果离了它我们不能活，\Quote{肉毫无益处}又如何說呢？你看见吗？\Quote{肉毫无益处}这句话不是指着祂自己的肉说的，而是指着属血肉的听？\par

% structure: p0025 source=h2 target=subsection reason=relative-heading-rank; base=h2

\subsection{\BibleRef{若 6:64}}

再次，按照祂的习惯，祂藉着预言将要发生的事，并表明祂这样说并非出于从他们得荣耀的愿望，而是因为祂关心他们，来加重祂话语的分量。当祂说\Quote{有些人}时，祂将门徒排除在外。因为起初祂说：\Quote{你们既看见了我，却不相信}\BibleRef{若 6:36}；但在这里说：\Quote{你们中间有些人不相信。}\par

因为祂\Quote{从开始就知道谁是不信的，谁要出卖祂。}\par

% structure: p0028 source=h2 target=subsection reason=relative-heading-rank; base=h2

\subsection{\BibleRef{若 6:65}}

3. 这里，福音作者向我们暗示了救恩工程的自愿性和祂对邪恶的忍耐。而且，\Quote{从开始}并非无端放在这里，而是为让你意识到祂从一开始就预知，并且早在话语说出之前，不是在人埋怨之后，也不是在他们跌倒之后，祂就知道那叛徒，而是之前，这是天主性的一个属性。然后祂又说：\Quote{除非从我的圣父得到赐予}；从而说服他们视天主为祂的圣父，而非\Person{若瑟}，并向他们表明相信祂不是一件寻常的事。好像祂说：\Quote{不信者不会扰乱我，不会麻烦我，不会使我惊讶。我早在他们受造之前就知道，我知道圣父赐予谁相信；} 并且，当你听到\Quote{祂赐予了}时，不要联想到一种任意的分配，而是若有人使自己配得上领受这恩赐，他就领受了。\par

% structure: p0030 source=h2 target=subsection reason=relative-heading-rank; base=h2

\subsection{\BibleRef{若 6:66}}

福音作者说得对，不是说他们\Quote{离开}，而是说他们\Quote{退回去}；表明他们切断了自身与任何德性长进的联系，并且藉着分离，他们失去了从前所有的信德。但十二位宗徒并非如此；因此祂对他们说：\par

% structure: p0032 source=h2 target=subsection reason=relative-heading-rank; base=h2

\subsection{\BibleRef{若 6:67}}

再次显示祂不需要他们的服务与伺候，并向他们证明祂带领他们同行并非为此。因为祂若甚至对他们也使用这样的言辞，又怎能需要他们呢？但祂为何不称赞他们？为何不嘉许他们？既是为了保持教师应有的尊严，也是要向他们表明，他们更应被这种对待方式所吸引。因为若祂称赞了他们，他们可能自以为是在为祂效劳，而怀有属人的情感；但通过表明祂不需要他们的陪伴，祂反而更牢牢地留住了他们。且注意祂说话何等谨慎。祂没有说\Quote{你们离开吧}（这会使他们远离祂），而是问他们说：\Quote{你们也愿意离开吗？}这话消除了任何强迫或压力，表明祂不愿他们因羞耻感而依附于祂，而愿他们怀有感恩之心。祂没有公开指责，而是温和地暗示他们，彰显了在这种情况下真正智慧的做法。但我们感受不同；这理所当然，因为我们凡事紧抓自己的体面，因此认为跟随者离开会降低我们的地位。但祂既不奉承也不排斥他们，只是问他们一个问题。这不是轻视他们的举动，而是不愿他们被强迫和压力所束缚：因为在这种条件下留下，无异于离开。那么伯多禄说什么呢？\par

% structure: p0034 source=h2 target=subsection reason=relative-heading-rank; base=h2

\subsection{若望福音 6:68-69}

你看见吗？那使人跌倒的并非言语，而是听众的疏忽、懈怠与心术不正。因为即使祂没有说这些话，他们也会跌倒，并且不会停止为肉身的食物焦虑，永远钉在地上。此外，门徒们与其他人同时听见，却发表了相反的意见，说：\Quote{我们去投奔谁呢？}这句话显出了极大的爱戴，因为它表明他们的老师比任何事物都更宝贵，比父亲、母亲或任何财产都更宝贵，并且若他们离开祂，便无处可逃。接着，免得人以为他说\Quote{我们去投奔谁}是因为无人接纳他们，他立刻加上：\Quote{你永生的话。}因为犹太人属血肉地、以人的理智来听，而门徒们却属神地，并将一切交托于信德。因此基督说：\Quote{我对你们所说的话就是神}——意思是：\Quote{不要以为我言语的教导受物质后果的法则或受造之物的必然性所支配。属神的事物并非如此，也不屈从于地上的律法。}保禄也宣告了这一点，说：\Quote{你心里不要说：谁要升到天上去？（就是领基督下来）或：谁要下到深渊中去？（就是将基督从死者中领上来）}\BibleRef{罗 10:6}\par

\Quote{你永生的话。}这些人已承认复活以及那时将要来临的一切分配。注意这位充满手足之爱与深情的人如何代表整个团体回答。因为他没有说\Quote{我知道}，而是说\Quote{我们知道}。或者，注意他如何直接沿用老师的话，不似犹太人那样说。犹太人说：\Quote{这是若瑟的儿子}，而他说：\Quote{你是基督，永生天主的儿子}，以及\Quote{你永生的话}；他或许曾听祂说：\Quote{信从我的，必永远不死，我在末日要叫他复活。}因为他回忆原话，表明他保留了一切所听过的。那么基督做了什么？祂没有称赞或赞赏伯多禄，尽管别处祂曾如此；但祂说什么呢？\par

% structure: p0037 source=h2 target=subsection reason=relative-heading-rank; base=h2

\subsection{若望福音 6:70}

因为既然伯多禄说\Quote{我们信}，耶稣便将犹达斯排除在团体之外。在别处，伯多禄并未提及门徒们；但当基督说\Quote{你们说我是谁？}时，他回答说：\Quote{你是基督，永生天主之子}\BibleRef{玛 16:15}；但在这里，既然他说\Quote{我们信}，基督有理不将犹达斯纳入那个团体。而且祂远在事件发生之前很久就这样做了，以遏制叛徒的邪恶，祂知道这毫无成效，但仍尽了自己的本分。\par

4. 且注意祂的智慧。祂没有公开叛徒，但也不让他隐藏起来；一方面使他不会全然失去羞耻心而变得更悖逆，另一方面使他不会以为自己被察觉而无所畏惧地行恶。因此祂逐渐对他提出更明显的责备。首先，祂将他同样列在众人之中，说：\Quote{你们中有人不信}（因为福音作者宣告祂将叛徒也算在内，说：\Quote{因为祂从起初就知道谁不信，谁要出卖祂}）；但当他仍保持原样时，祂对他提出更严厉的斥责：\Quote{你们中有一个是\OriginalTerm{魔鬼}{devil}}，却使恐惧普及所有人，愿意隐藏他。此处值得探讨：为何门徒们这时什么也不说，后来却害怕、疑虑，彼此相看，并问：\Quote{主，是我吗？}\BibleRef{玛 26:22}，那时伯多禄示意若望去查明叛徒是谁，通过询问老师谁是叛徒。原因是什么？伯多禄尚未听见\Quote{撒殚，退到我后面去！}因此他毫无惧怕；但当他被斥责后，尽管他出于深情说话，却未得赞许，反被称作\Quote{\OriginalTerm{撒殚}{Satan}}，之后他听见\Quote{你们中有一人要出卖我}时自然害怕。况且，祂甚至现在也没有说\Quote{你们中有一人要出卖我}，而是说\Quote{你们中有一个是魔鬼}；因此他们不明白所说话的意思，以为祂只是在指责他们的邪恶。\par

但是祂为什么说：\BibleQuote{我拣选了你们十二人，而你们中有一个是魔鬼}？这是为了表明祂的教导完全不讲谄媚。因为为了使他们不致以为祂会奉承他们——因为当所有人都离弃祂时，只有他们留了下来，并藉\Person{伯多禄}宣认祂是基督——祂便引导他们远离这种猜疑。祂所说的意思如下：\Quote{没有什么能使我羞于责备恶人；不要以为因为你们留下来了，我就选择奉承你们，或因为你们跟随了我，我就不责备恶人。因为连另一种情况也不会使我羞愧，而这种情况比这更易使一位教师羞愧。因为留下的人证明了他的爱戴，而被教师拣选后又遭弃绝的人，在愚昧人眼中会给教师带来愚拙的名声。然而，这也不能使我停止责备。}至少直到如今，外邦人仍冷酷而愚蠢地将此作为反对基督的理由。因为天主不惯于用强迫和强制作成善人，祂的拣选和选择对蒙召者也不是强制的，而是说服性的。为使你们知道这召叫并不强制，请想一想有多少蒙召的人竟丧亡了，如此可见，得救或丧亡也在于我们自己的意志。\par

5. 因此，听到这些事，我们要常保持清醒和警醒。因为如果那位曾被算在圣洁团体中，曾享有如此大恩赐，曾施行奇迹——因为他也与其他被派遣去复活死者和治愈癞病人的人一起——当他被贪婪的可怕疾病抓住并出卖了他的主时，无论是恩惠、恩赐、与基督同在、服侍祂、洗脚、同席、还是掌管钱囊，都无助于他，反而助长了他的惩罚，那么我们也当恐惧，免得我们因贪婪而效法\Person{犹达斯}。你不曾出卖基督。但是当你忽视因饥饿而消瘦、因寒冷而濒死的穷人时，那人就给你招致同样的定罪。当我们不配地参与圣事时，我们与杀害基督者一同丧亡。当我们掠夺、压迫比我们弱小的人时，我们将招致最严厉的惩罚。这合情合理；因为对现世之物的贪恋还要占据我们多久呢？它们本是多馀而无益的。财富在于多馀，其中毫无益处。我们要被虚妄之物钉住多久呢？我们要多久才不望向天空、保持清醒、不以这些易逝的尘世之物为满足，不从经验中学习它们的无价值呢？让我们想想那些在我们之前曾富有的人；所有那些岂不都是梦？岂不都是影儿、花朵？岂不都是流水？故事和传说？某人曾经富有，而今他的财富在哪里？它已逝去、消失，但因财富而犯的罪仍留在他身上，以及因罪而来的惩罚。是的，诚然，如果没有惩罚，如果没有天国摆在面前，我们也应当尊重同族同源的人，尊重那些与我们有同样感情的人。但如今我们喂养狗，我们中许多人喂养野驴、熊和各样的野兽，却毫不关心因饥饿而垂死的人；与我们有异的物比我们的亲族更受重视，我们自己的家人反而不如那些非我族类的受造物。\par

为自己建造华丽的房屋、拥有许多仆人、躺卧并凝视镀金的天花板，这难道是好事吗？那么，诚然，这是多馀而无益的。因为有别的建筑，远比这些更明亮、更庄严；我们应使眼睛喜悦地观看这些，因为无人阻止我们。你想看最美丽的屋顶吗？在傍晚仰望布满星辰的天空。有人说：\Quote{但这屋顶不属于我。}然而事实上，这个屋顶比那个更属于你。因为它为你而造，且是你与你的弟兄们共有的；那个不属于你，而属于在你死后继承它的人。一个可能为你带来最大的益处，以其美丽引导你归向它的创造者；另一个则带来最大的害处，在审判之日成为你最有力的控告者，因为它覆盖着黄金，而基督甚至连必要的衣服都没有。我恳求你们，不要陷于这样的愚昧，不要追求那些飞逝的事物，而逃避那些持久的事物；不要背叛我们自己的救恩，而要坚守对将来之事的希望；年长者，因为确知我们余生不多；年轻人，也深信所剩无几。因为那日子要像夜间的盗贼一样来到。知道这事，让妻子劝勉丈夫，丈夫规劝妻子；让我们教导青年男女，众人彼此训诲，不关心现世之事，而渴慕将来之事，使我们也能获得它们；借着我们的主耶稣基督的恩宠和慈爱，愿光荣归于圣父及圣神，从现在直到永远，世世无尽。阿们。\par

\SourceRecord{Translated by Charles Marriott.；Nicene and Post-Nicene Fathers, First Series；Vol. 14.；Edited by Philip Schaff.；Buffalo, NY: Christian Literature Publishing Co.,；1889.；<http://www.newadvent.org/fathers/240147.htm>.}

% structure: p0001 source=h1 target=section reason=page-title

\section{若望福音讲道集 48}

% structure: p0002 source=h2 target=subsection reason=relative-heading-rank; base=h2

\subsection{若望福音 7:1-2}

一、嫉妒与恶意没有比这更坏的了；通过它们，死亡进入了世界。因为当魔鬼看见人受尊荣时，它不能忍受他的兴旺，而用尽一切手段来毁灭他。\BibleRef{智 2:24} 从同一根源，随处都可以看到结出同样的果实。因此亚伯尔被杀害；达味及许多其他义人也几乎如此；由此犹太人也成为了杀害基督者。福音作者宣告这事说：\BibleQuote{此后，耶稣在加里肋亚行走，因为祂不能在犹太行走，因为犹太人想要杀死祂。} 你说什么，真福若望？难道祂没有能力，那能成就祂所愿意的一切者？祂说：\BibleQuote{你们找谁？}\BibleRef{若 18:6} 就使他们倒退？祂临在却不被看见\BibleRef{若 21:4}，难道祂没有能力吗？那么后来祂怎么来到他们中间，在圣殿正中，在庆节当中，当集会时，当渴望谋杀的人在场时，说出那些更激怒他们的话呢？是的，至少人们对此感到惊奇，说：\BibleQuote{这不是他们想要杀死的那个人吗？看，祂公开讲话，而他们什么也不对祂说。}\EditorialNote{25–26节} 这些谜是什么意思？不要这样说！福音作者并非如此说以让人觉得他是在说谜语，而是要表明祂同时显示天主性和人性的证据。因为当他说\Quote{祂没有能力}时，他是将祂作为人来说，做了许多符合人的方式的事；但当他谈到祂站在他们中间而他们并未捉拿祂时，他向我们显示了天主性的能力，（作为人祂逃避，作为天主祂显现，）而在两种情况下他都说出了真理。置身于那些图谋反对祂的人中间而不被捉拿，显示了祂无可比拟、不可抗拒的本性；退让则巩固并印证了救恩计划，好让撒摩沙塔的保禄、马西昂以及那些受他们病态影响的人无话可说。由此祂便堵住了他们所有人的口。\par

\BibleQuote{此后，是犹太人的帐棚节。} \Quote{此后}一词，仅表示作者在此处作了简略叙述，避开了很长一段时间，这从以下情况可以清楚看出。当基督坐在山上时，他说那是逾越节；而这里作者提到\Quote{帐棚节}，在这五个月内，除了饼的奇迹和给那些吃饼的人的讲道之外，他并没有叙述或告诉我们任何其他事。然而祂并未停止行奇迹和谈话，无论是在白天、晚上，还是经常在夜间；至少，祂这样带领祂的门徒，正如所有福音作者告诉我们的。那么他们为什么忽略那段时间呢？因为不可能完全记下所有的事，而且，因为他们急于提到那些曾引起犹太人挑剔或反驳的事情。有很多类似这里被省略的情况；因为祂复活死人、医治病人并受人钦佩，他们经常记录；但当他们有什么不寻常的事要讲，当他们必须描述任何看来是针对祂的指控时，他们才记下这些事；例如现在这件事：\Quote{祂的兄弟们不相信祂。} 因为这样的事并非带来少许怀疑，而且值得赞赏他们热爱真理的性情，他们不耻于叙述那些似乎给他们的老师带来耻辱的事，甚至比其它事情更急于报告这些事。例如，作者跳过许多征兆、奇迹和讲道，突然跳到这件事上。\par

% structure: p0005 source=h2 target=subsection reason=relative-heading-rank; base=h2

\subsection{若望福音 7:3-5}

二、这里有什么不信呢？他们劝祂行奇迹。这是大事；因为他们的言语、傲慢和不合时宜的直言都来自不信。因为他们认为，由于他们的亲属关系，他们可以大胆地对祂说话。他们的请求表面上是朋友的，但话语里带着极大的恶意。因为在此处他们指责祂胆怯和虚荣；因为说\Quote{没有人暗地里做事}是责备祂胆怯，并怀疑祂所做的事并非真正做到的；而加上\Quote{祂寻求被人知道}则是控告祂虚荣。但请注意，我请你们看看基督的能力。说这些话的人中，有一位成了耶路撒冷的第一位主教，真福雅各伯，保禄说：\Quote{我没有看见别的宗徒，只看见主的兄弟雅各伯}\BibleRef{迦 1:19}；犹达据说也是一个令人钦佩的人。然而这些人也在加纳，当酒变成时，但那时他们并未受益。那么他们从何处有这么大的不信呢？来自他们邪恶的心和嫉妒；因为亲属间的优越感常被那些未同样提升的人所嫉妒。他们在这里所称的门徒是谁？是跟随祂的群众，不是十二人。那么基督说什么呢？看祂如何温良地回答；祂没有说：\Quote{你们是谁，竟这样劝诫教导我？}而是说：\par

% structure: p0007 source=h2 target=subsection reason=relative-heading-rank; base=h2

\subsection{若望福音 7:6}

祂在此似乎向我暗示了与祂所表达的不同的事；也许他们在嫉妒中曾设计把祂交给犹太人；而向他们指出这一点，祂说：\BibleQuote{我的时候还没有到，} 那就是说：\Quote{十字架与死亡的时候，那么你们为何在时候未到之前急于杀害我？}\par

\BibleQuote{但你们的时候总是准备好的。}\par

好像祂说：\Quote{虽然你们常与犹太人同在，他们不会杀你们这些与他们愿望相同的人；但他们会立刻想要杀我。因此，你们时常可以安全地与他们在一起，但我的时候是十字架的时机临近时，是我必须死的时候。} 因为这就是祂的意思，祂通过后续的事表明了。\par

% structure: p0011 source=h2 target=subsection reason=relative-heading-rank; base=h2

\subsection{若望福音 7:7}

\Quote{就是说，因为我斥责并责备它，所以它恨我。} 由此让我们学习控制我们的愤怒，不要屈服于不配的激情，即使给我们劝告的是卑劣的人。因为若基督温和地容忍不信者劝告祂，虽然他们的劝告不当且不怀好意，那么，我们不过是尘土和灰烬，却对劝告我们的人恼怒，并认为我们受到了不配的待遇，即使这样做的人也许只比我们卑微一点，我们又将得到什么宽恕呢？请注意在此祂如何以全部温柔来驳斥他们的指控；因为当他们说：\Quote{祢将自己显给世人看}，祂回答说：\Quote{世人不能恨你们，却是恨我}，从而消除了他们的指控。祂说：\Quote{我远非寻求从人来的荣耀，以致我不断地责备他们，而且我知道这样做会招致对我的仇恨并为我预备死亡。} 有人问：\Quote{祂在何处责备了人？} 祂何时停止过这样做？祂不是说：\BibleQuote{不要想我要在圣父面前告你们；有一位告你们的，就是梅瑟。} \BibleRef{若 5:45} 又说：\BibleQuote{我知道你们没有天主的爱在你们里面} \BibleRef{若 5:42}；以及\BibleQuote{你们彼此领受光荣，却不寻求从唯一天主来的光荣，怎能相信呢？} \BibleRef{若 5:44} 你看见祂怎样处处表明，不是违反安息日，而是公开的责备，导致了人们对祂的仇恨吗？\par

祂为何打发他们去过节，说：\par

% structure: p0014 source=h2 target=subsection reason=relative-heading-rank; base=h2

\subsection{\BibleRef{若 7:8}}

为要表明祂说这些话，并非因为需要他们，或渴望得到他们的奉承，而是允许他们去作属于犹太人的事。有人说：\Quote{那么，祂说过\InnerQuote{我不上去}之后，怎么又上去了呢？} 祂不是一次地说\Quote{我不上去}，而是说\Quote{现在不上去}，即\Quote{不与你们一同上去}。\par

\Quote{因为我的时候还没有满。}\par

然而祂即将在即将到来的逾越节被钉十字架。\Quote{那么祂怎么也不上去呢？因为如果祂因时候未到而不上去，祂根本就不该上去。} 但祂上去不是为了受苦，而是为了教训他们。\Quote{但为何秘密地？因为祂若公开上去，既可以置身他们中间，又可以约束他们的暴躁冲动，正如祂常做的那样。} 这是因为祂不愿继续这样做。倘若祂公开上去，又再次使他们眼瞎，祂就会更大程度地彰显祂的天主性，这在当时并不适宜，祂反而更隐藏它。既然他们认为祂留下是因胆怯，祂便向他们显示相反的情况，即出于信心和一项安排，并且祂预先知道祂受苦的时候，当那时刻终于临近，祂会极其渴望上耶路撒冷去。我认为，祂说\BibleQuote{你们上去吧}，意思是\Quote{不要以为我强迫你们违背自己的意愿与我同住}，而补充说\BibleQuote{我的时候还没有完全到}，乃是宣告必须行神迹和讲道，以便更多群众相信，门徒因看见主的勇敢和苦难而更加坚定。\par

3. 让我们从所说的这些话中，学习祂的仁慈和温良；\BibleQuote{你们跟我学习，因为我是良善心谦的}\BibleRef{玛 11:29}；并让我们抛弃一切怨恨。若有人自高反对我们，让我们自卑；若有人傲慢，让我们服侍他；若有人以讥讽和戏弄咬啮我们，让我们不要被战胜，免得在自卫中毁灭自己。因为愤怒是一头野兽，一头凶猛而暴躁的野兽。让我们用从圣经中汲取的柔和话对自己说：\Quote{你是尘土和灰烬。} \BibleQuote{尘土和灰烬，你为什么骄傲？}\BibleRef{德 10:9}，以及\BibleQuote{愤怒的权势必成为他的毁灭}\BibleRef{德 1:22}：和\BibleQuote{愤怒的人不相称}\BibleRef{箴 11:25}；因为没有什么比因愤怒而涨红的面容更可耻、更丑陋的了。正如你搅动泥沼会发出恶臭，同样，当灵魂被激情扰乱时，便会有极大的不雅和不快。\Quote{但，}有人说，\Quote{我受不了敌人的侮辱。}为什么？请告诉我。如果控告属实，那么你甚至在受辱之前就应该心中刺痛，并感谢你的敌人斥责你；如果是假的，就轻视它。他骂你贫穷，嘲笑他；他骂你出身低贱和愚蠢，就为他哀悼；因为\BibleQuote{凡向弟兄说\InnerQuote{傻子}的，难免地狱的火。}\BibleRef{玛 5:22}因此，每当有人侮辱你时，想想他所受的惩罚；这样你不但不会生气，反而会为他流泪。因为没有人对发烧或发炎的人生气，而是怜悯和哭泣；愤怒的灵魂正是如此。不，即使你想要复仇，保持沉默，你就给了你的敌人致命一击；而如果你以辱骂还辱骂，你就点燃了烈火。\Quote{但，}有人说，\Quote{若我们保持沉默，旁观者会指责我们软弱。}不，他们不会指责你的软弱，反而会因你的智慧而钦佩你。此外，如果你因侮辱而刺痛，你就变得无礼；并且受了刺痛，就会迫使人们认为对你的指责是真实的。告诉我，为什么一个富人被说成贫穷时会发笑？岂不是因为他意识到自己并不贫穷吗？因此，如果我们嘲笑侮辱，我们就提供了最强有力的证据，表明我们并不意识到所指控的过错。此外，我们还要害怕向人交账到何时？我们还要轻视我们共同的主，并被钉在肉体上到何时？\BibleQuote{因为你们中间有嫉妒、纷争，这岂不是属乎肉体吗？}\BibleRef{格前 3:3}让我们成为属神的人，并勒住这可怖的野兽。愤怒与疯狂无异，它是一时的魔鬼，或者说，比拥有魔鬼更糟糕；因为拥有魔鬼的人或许可原谅，但愤怒的人应受万种惩罚，自愿投身毁灭的深渊，在未来的地狱之前，就已经在此世受惩罚，给灵魂的理性带来整夜整日不安的骚动和永不平静的愤怒风暴。因此，为了使我们摆脱此世的惩罚和来世的报应，让我们抛弃这种激情，展现出全部的柔和与温良，以便在此世和天国中获得灵魂的安息。愿我们众人藉着我们的主耶稣基督的恩宠和慈爱，达至这一境界。藉着祂、偕同祂，圣父和圣神，获享光荣，自今至永远，世世无尽。阿们。\par

\SourceRecord{Translated by Charles Marriott.；Nicene and Post-Nicene Fathers, First Series；Vol. 14.；Edited by Philip Schaff.；Buffalo, NY: Christian Literature Publishing Co.,；1889.；<http://www.newadvent.org/fathers/240148.htm>.}

% structure: p0001 source=h1 target=section reason=page-title

\section{\WorkTitle{若望福音讲道集 第四十九篇}}

% structure: p0002 source=h2 target=subsection reason=relative-heading-rank; base=h2

\subsection{\BibleRef{若 7:9-10}}

1. 基督按人的方式行事，不仅是为确立降生成人，也是为教导我们修德。因为如果祂完全以天主的方式行事，当我们遇上不如意之事时，怎能知道该如何做呢？例如，就在此地，犹太人要杀祂，祂却走到他们中间，平息了骚动。但若祂一直这样行事，我们这些无力如此行、又陷入同样处境的人，怎能知道应如何处理？是立刻送命，还是用些计谋以使圣言得以传播？因此，由于我们无力行此事，无法知道在仇敌中间该如何做，为此，祂亲自教导我们。因为圣史说：耶稣\Quote{说了这些话，仍留在加里肋亚；但祂的弟兄们上去过节以后，祂也上去了，但不是公开去的，而是暗中去的。}\Quote{祂的弟兄们上去以后}这句话表明祂不愿与他们同去。因此祂留在原地，没有现身，尽管他们似乎在催促祂。但祂为何一向公开说话，现在却\Quote{暗中}去呢？作者不说\Quote{秘密地}，而说\Quote{暗中}。因为，如我所说，祂似乎是在教导我们如何处理事情。此外，在对方情绪激动、桀骜不驯时来到他们中间，与在节期结束后再来，情况也是不同的。\par

% structure: p0004 source=h2 target=subsection reason=relative-heading-rank; base=h2

\subsection{\BibleRef{若 7:11}}

这些节期中的善行实在精彩！他们急于杀人，甚至在节期中也要捉拿祂。至少在另一处，他们这样说：\Quote{你们以为祂不来过节吗？}\BibleRef{若 11:56}；而此处他们说：\Quote{祂在哪里？}由于他们过度的仇恨和敌意，甚至不愿呼祂的名。他们对节期的敬畏何其大，谨慎何其大！他们竟想借着节期来陷害祂！\par

% structure: p0006 source=h2 target=subsection reason=relative-heading-rank; base=h2

\subsection{\BibleRef{若 7:12}}

我想，他们因行奇迹的地点而恼怒，大为激怒和害怕，与其说是因先前的事而愤怒，不如说是怕祂再次行类似的事。但一切都与他们的意愿相反，他们反而无意中使祂更加引人注目。\par

\Quote{有的说：祂是好人；有的却说：不然，祂迷惑群众。}\par

我想第一种意见是多数人的看法，第二种是首领和司铎的看法。因为诽谤祂正合他们的恶意和邪恶。他们说：\Quote{祂迷惑群众。}请问，怎么迷惑？是假装行奇迹而不真正行奇迹吗？但经验证明恰恰相反。\par

% structure: p0010 source=h2 target=subsection reason=relative-heading-rank; base=h2

\subsection{\BibleRef{若 7:13}}

你岂不见处处都是当权者败坏，被统治的人判断正确，却没有应有的勇气，而群众尤其缺乏这勇气吗？\par

% structure: p0012 source=h2 target=subsection reason=relative-heading-rank; base=h2

\subsection{\BibleRef{若 7:14}}

借着迟延，祂使他们更加留心；因为那些在头几天寻找祂、说\Quote{祂在哪里？}的人，当突然看见祂时，请看他们怎样靠近，好像要拥挤祂讲话，无论是说祂是好人的人，还是说祂不是的人；前者为受益而钦佩祂，后者为捉拿和拘留祂。有一方因教导和教义而不明其意，便说\Quote{祂迷惑群众}；另一方因奇迹而说\Quote{祂是好人}。因此，当祂平息了他们的怒气后，就这样来到他们中间，好让他们能从容听祂的话，不再有激情堵住他们的耳朵。祂教导了什么，圣史没有告诉我们；只说祂奇妙地教导，并赢得了他们，使他们归向祂。祂的话语就有这样的力量。那些曾说\Quote{祂迷惑群众}的人也改变了看法，\Quote{且感到惊奇}。因此他们也说：\par

% structure: p0014 source=h2 target=subsection reason=relative-heading-rank; base=h2

\subsection{\BibleRef{若 7:15}}

你注意到圣史在这里也显示他们的惊奇充满恶意吗？因为他不说他们赞叹教导，或接受话语，只说他们\Quote{惊奇}。也就是说，他们陷入惊愕状态，疑惑地说：\Quote{这个人从哪里得到这些？}他们正应从这困难中得知祂身上毫无纯属人性的事。但因他们不愿承认这事，只停留在惊奇上，请听祂说什么。\par

% structure: p0016 source=h2 target=subsection reason=relative-heading-rank; base=h2

\subsection{\BibleRef{若 7:16}}

祂再次回答他们心中的思念，将他们引向圣父，以此堵住他们的口。\par

% structure: p0018 source=h2 target=subsection reason=relative-heading-rank; base=h2

\subsection{\BibleRef{若 7:17}}

祂所说的是：\Quote{从你们自己身上除去那无故对我产生的恶意、愤怒、嫉妒和仇恨，那么没有什么能阻止你们知道我的话确实出自天主。因为目前这些事给你们投下黑暗，毁灭了正确判断的光明；如果你们除去它们，就不再是这种情况了。}不过祂没有（明白）这样说（因为那样会使他们极其困惑），而是通过说\Quote{谁若愿意承行祂的旨意，就会知道这教导是否出于天主，或者我是否凭自己说话}来暗示这一切；也就是说，\Quote{我是否说了什么不同、异常、违反天主的事。}因为\Quote{凭自己}总带有这个意思，即\Quote{我除了祂认为好的事以外，什么也不说；凡圣父所愿意的，我也愿意。}\par

\Quote{谁若愿意承行祂的旨意，就会知道这教导。}\par

\Quote{谁若愿意承行祂的旨意}是什么意思？\Quote{谁若热爱符合美德的生活，就会知道这些话的力量。} \Quote{谁若愿意留意预言，看我说的是否符合它们。}\par

2. 但祂的教导如何是祂的，又不是祂的？因为祂没有说：\Quote{这教导不是我的}；而是先说了\Quote{是我的}，将其视为己有，然后才加上\Quote{不是我的}。那么，同一件事怎能既是\Quote{祂的}又不是\Quote{祂的}呢？是\Quote{祂的}，因为祂说话不像一个受教者；不是\Quote{祂的}，因为那是圣父的教导。那么，祂怎么说：\Quote{凡圣父所有的，都是我的，我也属于祂}？\BibleRef{若 17:10}\Quote{因为如果教导是圣父的，就不是你的，那么另一断言便是假的，因为照此它应该是你的。}但\Quote{不是我的}提供了强有力的证据，证明祂的教导与圣父的教导是同一的；好像祂说：\Quote{它毫无不同，仿佛属于另一位。因为虽然我的人位不同，但我的言行如此，不被人认为我所言所行与圣父相悖，反而正是圣父所言所行。}然后祂加上另一个无可辩驳的论据，引用了某种纯粹属人的事物，并用他们熟悉的事例教导他们。那是什么呢？\par

% structure: p0023 source=h2 target=subsection reason=relative-heading-rank; base=h2

\subsection{若望福音 7:18}

即：\Quote{凡想要确立自己教导的，只为自己能享荣耀。现在，如果我不愿享荣耀，我何必要确立自己的教导呢？那凭自己说话的，即那说出独特或与众不同的道理的，便是为此以建立自己的荣耀；但如果我寻求那派遣我者的荣耀，我为何要选择教导别的事呢？}你看见了吗？祂在那里也说祂\Quote{凭自己什么都不做}，是有原因的。\EditorialNote{参见 7:19 和 8:28}是什么？是为了使他们相信祂不寻求众人的尊荣。因此，当祂的话谦卑时，祂说：\Quote{我寻求圣父的荣耀}，处处想要说服他们：祂自己不爱荣耀。祂使用谦卑的话有许多原因，例如：以免被视为非受生或与天主对立，因祂披戴肉身，因听众的软弱，为教导人谦虚，不说夸大自己的话；而说崇高的话只有一个原因：祂本性的伟大。而当祂说：\Quote{在亚巴郎以前，我就是}\BibleRef{若 8:58}时，他们便感到冒犯；那么，他们若一直听崇高的话语，又会如何呢？\par

% structure: p0025 source=h2 target=subsection reason=relative-heading-rank; base=h2

\subsection{若望福音 7:19}

有人说：\Quote{这有什么联系？这与先前所说的有什么关系？}犹太人向祂提出两项指控：一是祂违反安息日；二是祂称天主为祂的父，使自己与天主平等。而这一点并非他们的想象，而是祂明确判断，且祂说话不像众人那样，而是有一种特别而独特的意义，从这情况便可清楚。许多人常称天主为父，如：\Quote{我们岂不都有一个父？岂不是一个天主创造了我们？}\BibleRef{拉 2:10}但百姓并不因此与天主平等，因此听众并未感到冒犯。于是，当犹太人说：\Quote{这人不是从天主来的}时，祂常治好他们，并为违反安息日辩护；同样，如果他们认为祂的话是按照他们的想象，而非按照祂的意图，祂就会纠正他们并说：\Quote{你们为何以为我与天主平等？我不平等}；然而祂并未说任何此类的话，反而在后续中宣告自己是平等的。因为\Quote{父怎样使死人复活，并赐予生命，子也照样}\BibleRef{若 5:21}；以及\Quote{为叫众人尊敬子，如同尊敬父一样}；\Quote{父所做的事，子也照样做}；这些都建立了祂的平等。此外，关于法律，祂说：\Quote{不要以为我来是废除法律或先知。}\BibleRef{玛 5:17}这样，祂知道如何消除他们心中的恶意猜疑；但在这里，祂不仅不消除，反而证实了他们对祂平等的猜疑。因此，当他们在别处说：\Quote{你使自己成为天主}时，祂并未消除他们的猜疑，反而证实了它，说：\Quote{为叫你们知道人子在地上有权赦免罪过，就对瘫子说：拿起你的床，走吧。}\BibleRef{玛 9:6}于是，祂首先致力于使自己与天主平等，表明自己不是天主的仇敌，而是与祂说同样的话，教导同样的事，然后祂着手处理违反安息日的问题，说：\Quote{梅瑟不是给了你们法律吗？你们中却没有一人遵守法律？}好像祂说：\Quote{法律规定：不可杀人；但你们杀人，却指责我违反法律。}但祂为何说\Quote{你们中却没有一人}？因为他们都想杀祂。祂说：\Quote{即便我确实违反了法律，那是在救人，而你们却是为了邪恶违反法律。即使我的行为是违反，但那是为了拯救，我不该被你们这些在最大事上违反的人审判。因为你们的行为是在推翻整个法律。}然后祂进一步强调，尽管祂先前已对他们说过许多话，但那时祂以更崇高的方式说话，更适合祂的尊严，而此刻祂说话更谦卑。为何？因为祂不愿不断激怒他们。当时他们的怒气变得强烈，并走向谋杀。因此祂继续以两种方式抑制他们：责备他们的恶行，并说：\Quote{你们为什么图谋杀害我？}以及谦卑地自称\Quote{一个把真理告诉你们的人}\BibleRef{若 8:40}，并显出心怀杀意的人不配审判他人。且看基督问题的谦卑，以及他们回答的傲慢。\par

% structure: p0027 source=h2 target=subsection reason=relative-heading-rank; base=h2

\subsection{若望福音 7:20}

3. 这是愤怒与怒气的话，是一个因意外的责备而变得肆无忌惮、并在他们以为时候未到之前就陷入慌乱的灵魂的话。因为正如一群强盗得意于他们的阴谋，当他们想使所谋害的人放松警惕时，便通过沉默来达成目的；这些人也是如此。但祂并没有为此责备他们，以免使他们更加无耻，而是再次围绕安息日进行辩护，从法律出发与他们辩论。请注意祂是多么谨慎。祂说：\Quote{难怪你们不顺从我，你们不顺从你们以为遵守、并且认为是梅瑟所赐给你们的那法律。因此，你们不留意我的话，也并非新鲜事。}因为他们曾说：\Quote{天主对梅瑟说过话，至于这个人，我们不知道他从哪里来}\BibleRef{若 9:29}，祂表明他们也侮辱了梅瑟以及祂自己，因为梅瑟给了他们法律，而他们却不遵守。\par

% structure: p0029 source=h2 target=subsection reason=relative-heading-rank; base=h2

\subsection{若望福音 7:21}

请注意祂在需要为自己辩护、并使祂的辩护成为对他们的控告时，是如何论证的。因为关于所行的事，祂并未引入圣父的身份，而是祂自己的身份：\Quote{我作了一件事。}祂要表明，不作这件事就是违反法律，并且有许多事比法律更具权威，而且\Quote{梅瑟}曾接受一个违反法律的命令，这个命令比法律更具权威。因为\Quote{割礼}比安息日更具权威，然而割礼并非来自法律，而是来自\Quote{祖先}。\Quote{但我}，祂说，\Quote{作了比割礼更具权威、更好的一件事。}然后祂没有提及法律的命令——例如，司祭们在安息日渎圣，正如祂已经说过的——而是更广泛地谈论。\Quote{你们惊讶}\BibleRef{玛 12:5} 的意思是\Quote{你们困惑}、\Quote{你们不安}。因为如果法律是永久的，割礼就不会比它更具权威。祂没有说：\Quote{我作了一件比割礼更大的事}，而是通过说话大量地驳斥他们：\par

% structure: p0031 source=h2 target=subsection reason=relative-heading-rank; base=h2

\subsection{若望福音 7:23}

\Quote{你们看见法律在一个人违反它的时候最被确立吗？你们看见违反安息日就是遵守法律吗？如果安息日不被违反，法律就必须被违反？所以我也确立了法律。}祂没有说：\Quote{你们向我发怒，因为我作了一件比割礼更大的事}，而只是提及所行的事，留给他们判断：整个健康是否不比割礼更为必要？\Quote{法律}，祂说，\Quote{被违反，为叫人得到一个对健康毫无贡献的记号；你们却对为了使人脱离如此沉重的疾病而违反它感到愤怒和愤慨吗？}\par

% structure: p0033 source=h2 target=subsection reason=relative-heading-rank; base=h2

\subsection{若望福音 7:24}

\Quote{按外貌}是什么意思？\Quote{不要因为梅瑟在你们当中享有最大的尊荣，就按你们对人的评估来下判断，而要按事情的本性来判断；因为这才是判断正确。为什么你们当中没有人责备梅瑟？为什么没有人不顺从他，当他命令照一个从外部引入法律的诫命而违反安息日时？他允许一个诫命比他自己的法律更具权威；一个不是由法律引入、而是从外部引入的诫命，这尤其奇妙；而你们这些不是立法者的人，却过分地维护法律、保卫它。然而，命令法律被一个不属于法律的诫命所违反的梅瑟，比你们更值得信任。}于是祂说（我使）\Quote{整个人（健康）}，表明割礼也是\Quote{部分的}健康。割礼带来的健康是什么？经上说：\Quote{凡未受割礼的男子，都应从民间铲除。}\BibleRef{创 17:14}\Quote{但我扶起了一个不是局部受苦、而是完全败坏的人。}所以，\Quote{不要按外貌判断。}\par

我们要确信，这话不仅是对当时的人说的，也是对我们说的，叫我们在任何事上都不歪曲正义，而是为正义尽一切努力；无论一个人是穷人还是富人，我们都不看情面，只查察事情。经上说：\Quote{不可在穷人争讼的事上偏袒穷人。}\BibleRef{出 23:3} 这是什么意思？\Quote{不要心软，也不要屈服}，经上说，\Quote{如果那作恶的是个穷人。}既然你们不可偏袒穷人，何况富人？这话我不但对你们作审判官的说，也对所有的人说，叫他们任何地方都不歪曲正义，而在各处保持正义的纯洁。经上说：\Quote{上主喜爱正义}；并且\Quote{喜爱不义的，是憎恨自己的灵魂。}\BibleRef{咏 11:7} 我恳求，我们不要憎恨自己的灵魂，也不要喜爱不义。因为它在今世的利益微乎其微，甚至毫无益处，而对来世却带来巨大的损害。或者不如说，即使在此世，我们也无法享受它；因为当我们生活安逸，良心却有愧，这难道不是报应和惩罚吗？让我们就喜爱正义，永不偏离那条法律。因为如果我们不达至美德就离开今世，我们将从今世获得什么果实呢？到那时有什么会帮助我们？是友谊，还是亲属，或这个或那个人的青睐？我说什么？这个或那个人的青睐？即使我们有诺厄、约伯或达尼尔为父亲，若我们被自己的行为所出卖，这都无济于事。我们只需要一样东西，就是灵魂的美德。这将能保你安全度过、将你从永火中救出，这将护送你进入天国。愿我们众人，藉着我们的主耶稣基督的恩宠和慈爱，得以到达那里；藉着祂、偕同祂，圣父和圣神享受光荣，从现在直到永远，世世无穷。阿们。\par

\SourceRecord{Translated by Charles Marriott.；Nicene and Post-Nicene Fathers, First Series；Vol. 14.；Edited by Philip Schaff.；Buffalo, NY: Christian Literature Publishing Co.,；1889.；<http://www.newadvent.org/fathers/240149.htm>.}

% structure: p0001 source=h1 target=section reason=page-title

\section{若望福音讲道集第五十篇}

% structure: p0002 source=h2 target=subsection reason=relative-heading-rank; base=h2

\subsection{若望福音 7:25-27}

1. 圣经中没有一事是毫无理由地写下的，因为圣经是由圣神所启示的；因此，我们应当精确地查考每一点。因为从一句表述中，就可能找出某段经文的全部含义，就像我们眼前的情形一样。\Quote{耶路撒冷人中有的说：这个人不就是他们想要杀死的吗？然而，他还公开讲话，他们却不对他说什么。}那么，为什么加上了\InnerQuote{耶路撒冷人}呢？福音作者借此表明，那些最享受祂大能奇迹的人，比任何人都更可怜；他们亲眼见过祂天主性的最大明证，却将一切都托付于他们腐败领袖的判断。因为那些狂怒、决意谋害、四处寻找祂的人，当他们把祂掌握在手中时，却突然安静下来，这难道不是祂天主性的一大明证吗？谁能做到这事？谁这样扑灭了他们绝对的狂怒？然而在如此明证之后，请看这些人的愚昧和疯狂。\Quote{这个人不就是他们想要杀死的吗？}请看他们如何控告自己；经上说：\Quote{他们想要杀死祂，却不对祂说什么。}他们不仅不对祂说什么，甚至当祂\Quote{公开讲话}时也如此。因为一个公开讲话、完全自由的人，本应更激怒他们；但他们却什么也不做。\Quote{难道领袖们真的知道这就是基督吗？}\Quote{你们想怎样？你们给出什么意见？}经上说：相反。因此他们说：\Quote{我们知道这个人是从哪里来的。}何等的恶意，何等的矛盾！他们甚至不听从领袖的意见，反而提出另一套，既乖僻又配得上他们自己的愚昧：\Quote{我们知道祂是从哪里来的。}\par

\Quote{但当基督来时，没有人知道祂是从哪里来的。}\BibleRef{玛 2:4}\par

\Quote{然而你们的领袖被问到时回答说，祂应生在伯利恒。}另有些人又说：\Quote{天主曾对梅瑟说话；至于这个人，我们不知道祂是从哪里来的。}\BibleRef{若 9:29}\Quote{我们知道祂是从哪里来的，}又\Quote{我们不知道祂是从哪里来的}；请看这些醉言醉语。再次：\Quote{基督难道是出自加里肋亚吗？}\BibleRef{若 7:41}祂不\Quote{是出自伯利恒城}吗？你看见了吗？这是疯子的决断。\Quote{我们知道，}又\Quote{我们不知道}；\Quote{基督出自伯利恒}；\Quote{当基督来时，没有人知道祂是从哪里来的。}还有什么比这矛盾更明显的？因为他们只看一件事：就是不信。那么基督的回答是什么呢？\par

% structure: p0006 source=h2 target=subsection reason=relative-heading-rank; base=h2

\subsection{若望福音 7:28}

2. 又，\Quote{你们若认识我，也就认识我的父。}\BibleRef{若 8:19}那么，祂怎么说他们既\Quote{知道祂}，\Quote{知道祂是从哪里来的}，又\Quote{既不认识祂，也不认识父}呢？祂并不矛盾（断无此念！），而是与祂自己非常一致。因为祂谈到一种不同的认识，当祂说\Quote{你们不知道}时；正如祂说：\Quote{厄里的儿子是恶子，他们不认识上主}\BibleRef{撒上 2:12}；又说：\Quote{以色列不认识我。}\BibleRef{依 1:3}同样，保禄说：\Quote{他们声称认识天主，但在行为上却否认祂。}\BibleRef{铎 1:16}因此，是可能\Quote{知道}却又\Quote{不知道}的。这就是祂所说的：\Quote{你们若认识我，就知道我是天主子。}因为这里的\Quote{我从哪里来}并非指地方。这从下文可以清楚看出：\Quote{我不是由我自己而来，但那派遣我来的，是真实的，你们却不认识祂。}这里指的是他们在行为上显示的无知。[如保禄所说：\Quote{他们声称认识天主，但在行为上却否认祂。}]因为他们的过错不仅来自无知，也来自邪恶和恶意的意志；因为即使他们知道这事，他们却选择不知道。但这里有什么关联呢？祂如何责备他们，却用了他们自己的话？因为当他们说：\Quote{我们知道这个人是从哪里来的}，祂接着说：\Quote{你们也知道我。}难道他们说过\Quote{我们不知道祂}吗？不，他们说了\Quote{我们知道祂}。但是（请注意），他们说\Quote{我们知道祂是从哪里来的}，无非是表明祂是\Quote{出自地上}，是\Quote{木匠的儿子}；但祂却引导他们上升至天上，说：\Quote{你们知道我从哪里来}，即不是从你们所设想的地方，而是从派遣我来的那一位那里（祂派遣了我）。因为说\Quote{我不是由我自己而来}，是向他们暗示，他们知道祂是由父所派遣的，尽管他们不揭露这一点。因此，祂以双重方式责备他们：首先，他们私下说的话，祂公开说出来，以使他们羞愧；之后，祂又揭示了他们心中所想的。好像祂说：\Quote{我不是微贱之人，也不是无缘无故而来的人，而是那派遣我来的，是真实的，你们却不认识祂。}\Quote{那派遣我来的，是真实的}是什么意思？\Quote{如果祂是真实的，祂派遣我是为了真理；如果祂是真实的，那么被派遣的也应该是真实的。}祂也用另一种方式证明了这一点，以他们自己的话驳倒他们。因为他们曾说过：\Quote{当基督来时，没有人知道祂是从哪里来的}。祂由此证明祂自己就是基督。他们用\Quote{没有人知道}这句话，是指某种特定地方的区别；但祂却从同样的话表明祂是基督，因为祂是从父而来的；而且祂处处见证只有祂认识父，说：\Quote{不是有人见过父，惟独那从父来的，祂见过父。}\BibleRef{若 6:46}祂的话激怒了他们；因为对他们说\Quote{你们不认识祂}，并责备他们明知却假装无知，足以刺痛和骚扰他们。\par

% structure: p0008 source=h2 target=subsection reason=relative-heading-rank; base=h2

\subsection{若望福音 7:30}

你岂看见他们被无形地约束，他们的怒气被遏制了吗？但为何不说祂无形地约束了他们，却说\Quote{因为祂的时候还没有到}？圣史想要以更人性、更谦卑的方式说话，使基督也被视为人。因为基督处处谈论崇高之事，所以他便穿插此类言辞。当基督说\Quote{我从祂而来}，祂不是像先知那样学习，而是如同看见祂、并与祂同在。\par

% structure: p0010 source=h2 target=subsection reason=relative-heading-rank; base=h2

\subsection{\BibleRef{若 7:29}}

你岂看见祂不断试图证明\Quote{我不是凭自己来的}和\Quote{那派遣我的是真实的}，力求不被视为天主的仇敌？且看祂话语谦逊的益处何等大：因为经上说，此后有许多人说：\par

% structure: p0012 source=h2 target=subsection reason=relative-heading-rank; base=h2

\subsection{\BibleRef{若 7:31}}

有多少奇迹呢？实际上有三个：酒、瘫子、贵族之子；圣史也只记载了这些。由此显而易见，正如我常说的，作者们略过了大部分奇迹，只向我们讲述那些因之而领袖们虐待祂的奇迹。\Quote{于是他们想要捉拿祂}，并杀害祂。谁\Quote{想要}？不是群众，他们无意统治，也不致被恶毒俘获；而是司祭们。因为群众中有人说：\Quote{基督来的时候，祂会行更多的奇迹吗？}然而这也不是纯全的信德，而似乎是一群杂乱之人的想法；因为说\Quote{当祂来的时候}，并不是那些坚信祂是基督之人的表达。我们可以这样理解这些话，或者认为它们是群众聚集时说的。他们说：\Quote{既然我们的首领竭力证明这人不是基督，让我们假设他不是基督；那么基督会比他更好吗？}因为，我常重复，粗俗之人不是被道理，也不是被宣讲所引导，而是被奇迹所引导。\par

% structure: p0014 source=h2 target=subsection reason=relative-heading-rank; base=h2

\subsection{\BibleRef{若 7:32}}

你岂看见违反安息日只是借口？而最刺痛他们的是这怨言？因为此时，虽然祂的言行无可指摘，他们却因群众而想要捉拿祂。他们自己不敢行动，怀疑有危险，便派仆从去。唉！他们的专横和疯狂，或者更确切说，他们的愚昧。自己多次尝试未果后，竟将此事交给仆从，仅仅为满足怒气。然而祂在羊池说过许多话\EditorialNote{若望福音第5章}，他们并未如此行；他们确实寻找机会，但并未尝试；而此处，当群众将要投奔祂时，他们再也无法忍受。那么基督说什么呢？\par

% structure: p0016 source=h2 target=subsection reason=relative-heading-rank; base=h2

\subsection{\BibleRef{若 7:33}}

祂有能力使听众降服和畏惧，却说出了充满谦逊的话。仿佛祂说：\Quote{你们为何急于迫害和杀死我？稍等片刻，即使你们渴望留住我，我也不会忍受。}为使无人（如他们所想的）以为\Quote{我还有不多的时候与你们同在}指普通的死亡，为使无人这样想，或以为祂死后一无所成，祂又加上：\par

% structure: p0018 source=h2 target=subsection reason=relative-heading-rank; base=h2

\subsection{\BibleRef{若 7:34}}

倘若祂要继续留在死亡中，他们或许会到祂那里去，因为我们都往那地方去。因此祂的话使较纯朴的群众感到不安，使胆大者恐惧，使较聪明的人渴望听祂，因为时间不多了，他们不能永远享受这教训。祂不仅说\Quote{我在这里}，而是说\Quote{我与你们同在}，即\Quote{虽然你们迫害，虽然你们驱赶我，但在不多的时候，我不会停止为你们的益处行事，讲论并推荐那些关乎你们救恩的事。}\par

% structure: p0020 source=h2 target=subsection reason=relative-heading-rank; base=h2

\subsection{\BibleRef{若 7:33}}

这足以使他们恐惧并陷入痛苦。因为祂也宣告他们将需要祂……\par

% structure: p0022 source=h2 target=subsection reason=relative-heading-rank; base=h2

\subsection{\BibleRef{若 7:34}}

3. 那么犹太人何时\Quote{寻找祂}呢？路加说妇女们为祂哀恸，可能还有许多其他人，无论是在当时还是在城被攻陷时，都记起了基督和祂的奇迹，并渴望祂的同在。\BibleRef{路 23:49}祂加上这一切，意在吸引他们。因为时间短暂，祂离开后他们会后悔地渴望祂，且那时找不到祂，这一切足以说服他们到祂那里去。倘若祂的同在不会令他们后悔地渴望，祂的话在他们看来便无甚伟大；倘若渴望却又能找到祂，也不会搅扰他们。再者，倘若祂要与他们长久同在，他们也会懈怠。但现在祂从各方面迫使他们畏惧。而\Quote{我往那派遣我者那里去}，是宣告他们的阴谋不会伤害祂，祂的受难是自愿的。因此祂现在发出两个预言：不多时祂要离去，而他们不会到祂那里去；这并非人智慧所能及——预言自己的死亡。例如听听达味说：\BibleQuote{上主，求祢使我知晓我的结局，我的岁月有何数目，好使我知道我尚有何时光。}\BibleRef{咏 38:4}没有任何人知道这事；由一事证实另一事。我想祂是暗中对仆从们说话，将话指向他们，以此特别吸引他们，向他们显示祂知道他们到来的原因。仿佛祂说：\Quote{稍等片刻，我就要离去。}\par

% structure: p0024 source=h2 target=subsection reason=relative-heading-rank; base=h2

\subsection{\BibleRef{若 7:35}}

然而那些希望摆脱祂、竭尽所能不愿见祂的人，本不应问这个问题，而应说：\Quote{我们很高兴，何时离去？}但他们略显触动，并以愚蠢的猜疑彼此相问：\Quote{祂将往哪里去？}\par

\Quote{祂要到散居于外邦人中去吗？}\par

什么是\Quote{外邦人的分散}？犹太人这样称呼其他民族，因为他们到处分散，无所畏惧地彼此混杂。这耻辱他们自己后来也承受了，因为他们也成了\Quote{分散}。因为在古时，他们整个民族聚集在一处，除了在巴勒斯坦，你找不到一个犹太人；因此他们称外邦人为\Quote{分散}，以此指责他们，并对自己夸耀。那么，\Quote{我所去的地方，你们不能去}是什么意思？因为那时所有民族都和他们有来往，到处都是犹太人。所以，若他指的是外邦人，他就不会说：\Quote{你们不能去的地方。}在说了\Quote{他要往外邦人的分散那里去吗？}之后，他们并没有加\Quote{并毁灭}，而是加\Quote{并教训他们}。他们的怒气已如此平息，并且相信了他的话；因为他们若不相信，就不会彼此查问这话是什么意思。\par

这些话确实是对犹太人说的，但我也害怕它们也适用于我们，即\Quote{他在哪里}我们\Quote{不能去}，因为我们的生命充满了罪恶。因为关于门徒，他说：\BibleQuote{我愿意他们也在那里，同我在一起}\BibleRef{若 17:24}，但关于我们自己，我担心相反的话会被说出：\Quote{我在哪里，你们不能去。}因为当我们行事违反诫命时，怎能到那地方去呢？即使在今生，若一个士兵对国王行事不配，他不能见到国王，反而被剥夺权柄，遭受最严厉的惩罚；因此，若我们偷窃、贪恋，若我们伤害或打击他人，若我们不施行仁慈的工作，我们就不能到那里去，反而会遭受童女们所遭遇的。因为她所在的地方，她们不能进去，反而退去，她们的灯熄灭了，即恩宠离开了她们。因为如果我们愿意，我们可以增加我们由圣神的恩宠直接领受的那火焰的光亮；但若我们不愿这样做，我们就会失去它，当它熄灭时，我们的灵魂中就只剩下黑暗；因为如同灯燃烧时光线强烈，熄灭时则只有阴暗。因此宗徒说：\BibleQuote{不要熄灭圣神。}\BibleRef{得前 5:19}。它被熄灭，是因为没有油，因为有猛烈的大风，因为被压抑和限制（因为火就是这样被熄灭的），它被世俗的忧虑压抑，被邪恶的欲望熄灭。除了我们提到的原因，没有什么比不人道、残忍和掠夺更能熄灭它。因为当我们除了没有油之外，还往上浇冷水（贪恋之心就是这样，它使那些我们亏负的人因沮丧而冰冷），它从哪里重新点燃呢？因此，我们将带着灰尘和灰烬离去，并有许多烟雾证明我们有过灯却熄灭了它们；因为哪里有烟，哪里必然有过火被熄灭。愿我们中间没有人听见那句话：\BibleQuote{我不认识你们。}\BibleRef{玛 25:12}。我们又从哪里听见这句话，除非从不理会穷人，视而不见？若我们不肯认识饥渴的基督，他也不会在我们恳求他的仁慈时认识我们。这是公义的；因为那忽略受苦者、不给自己的东西的人，怎能寻求得到不属于他的东西呢？因此，我恳求你们，让我们做一切事，设法使油不缺乏，而是让我们整理我们的灯，与新郎一同进入洞房。愿我们所有人都能达到这一点，藉着我们的主耶稣基督的恩宠和慈爱，藉着他并同着他，圣父和圣神，获得荣耀，现在以至永远，世世无穷。阿们。\par

\SourceRecord{Translated by Charles Marriott.；Nicene and Post-Nicene Fathers, First Series；Vol. 14.；Edited by Philip Schaff.；Buffalo, NY: Christian Literature Publishing Co.,；1889.；<http://www.newadvent.org/fathers/240150.htm>.}

% structure: p0001 source=h1 target=section reason=page-title

\section{若望福音讲道 51}

% structure: p0002 source=h2 target=subsection reason=relative-heading-rank; base=h2

\subsection{若 7:37-38}

1. 凡前来听受神圣宣讲并注意信德的人，必须表现出干渴之人对水的渴望，并在自己心中燃起同样的渴望；这样他们才能极其仔细地牢记所说的话。因为正如口渴的人拿起碗来，急切地喝干然后停止一样，那些听受天主的圣言的人，如果他们渴慕地领受，就永远不会厌倦，直到完全喝尽。因为为了表明人应当常常饥渴，经上说：\BibleQuote{饥渴慕义的人是有福的}\BibleRef{玛 5:6}；而这里基督说：\BibleQuote{谁若渴，到我这里来喝罢！}祂所说的大意是：\Quote{我不用必要和强制拉任何人到我这里来；但若有人有大热忱，若有人被渴望所燃烧，我就召叫他。}\par

但是圣史为什么特别指出是\BibleQuote{末后的日子，即那大日子}呢？因为头一天和最后一天都是\Quote{大}的，而中间的日子他们倒是在享乐中度过的。那么祂为什么说\BibleQuote{在末后的日子}呢？因为那一天他们都聚集在一起。因为在第一天，祂没有来，也向祂的弟兄说明了原因；第二、第三天祂也没有讲这类话，免得祂的话落空，因为听众快要放纵自己。但在最后一天，当他们正要回家时，祂赐给他们救恩的供应，并大声呼喊，一方面向我们显示祂的勇敢，另一方面也因群众众多。为了表明祂说的不是物质饮料，祂加上说：\BibleQuote{谁信从我，就如经上所说，从他的腹中要流出活水的江河}。在这里祂以\Quote{腹}指心，正如在另一处经上说：\BibleQuote{你的法律在我腹中}\BibleRef{咏 39:10}。但是经上哪里说过\BibleQuote{从他的腹中要流出活水的江河}呢？没有。那么\BibleQuote{谁信从我，就如经上所说}是什么意思呢？这里我们必须加一个句号，使得\BibleQuote{江河要从他的腹中流出}成为基督的断言。因为许多人说：\Quote{这是基督}；又说：\Quote{基督来时，祂会行更多的奇迹吗？}祂表明应该具有正确的认识，不要从奇迹而得信服，更重要的是从圣经。事实上，许多人即使亲眼看见祂行奇事，也不接受祂为基督，反而总说：\Quote{圣经不是说过基督来自达味的后裔吗？}他们经常纠缠于此。于是，祂为了表明祂不回避来自圣经的证明，再次将他们引向圣经。祂先前说过：\BibleQuote{你们查考经典}\BibleRef{若 5:39}；又说：\BibleQuote{先知书上记载：他们都要受天主的教导}\BibleRef{若 6:45}；以及\BibleQuote{梅瑟控告你们}\BibleRef{若 5:45}；而这里说：\BibleQuote{就如经上所说，江河要从他的腹中流出}，这暗示恩宠的宽广与丰盛。正如在另一处祂说：\BibleQuote{涌到永生的水泉}\BibleRef{若 4:14}，这就是说，\Quote{他将拥有极大的恩宠}；在别处祂称之为\Quote{永生}，但这里称为\Quote{活水}。祂称那\Quote{活}的，就是永远运行的；因为圣神的恩宠，一旦进入心智并稳固下来，比任何泉源涌流得更旺，永不枯竭，永不空乏，永不止息。因此，为了同时表示其不竭的供应和无尽的运作，祂称之为\Quote{水泉}和\Quote{江河}，不是一条江河，而是无数条；在前一种情形中，祂用\Quote{涌}一词来表达其丰盈。如果一个人思量斯德望的智慧、伯多禄的口才、保禄的激昂，什么都不能阻挡他们，不能挫败他们，不是群众的愤怒，不是暴君的兴起，不是魔鬼的诡计，不是每日的死亡，而是像带着巨大响声奔腾的江河，他们把一切都席卷而去，那么就能清楚地明白这话的意思了。\par

% structure: p0005 source=h2 target=subsection reason=relative-heading-rank; base=h2

\subsection{若 7:39}

2. 那么先知如何预言并行了那万千奇事呢？因为宗徒们驱魔不是靠圣神，而是靠从祂领受的权能；正如祂自己所说：\BibleQuote{我若靠贝耳则步驱魔，你们的子弟凭什么驱魔呢？}\BibleRef{玛 12:27} 祂说这话，是表明在十字架受难之前，并非所有人都靠圣神驱魔，而是有些人靠从祂领受的权能。因此，当祂要派遣他们时，祂说：\BibleQuote{领受圣神吧}\BibleRef{若 20:22}；又说：\BibleQuote{圣神降临在他们身上}\BibleRef{宗 19:6}，然后他们就行奇事。但当祂派遣他们时，圣经没有说\BibleQuote{祂赐给他们圣神}，而是赐给他们\Quote{权能}，说：\BibleQuote{洁净癞病人，驱魔，复活死人，你们白白地领受了，也要白白地施与。}\BibleRef{玛 10:1} 但在先知们的情况下，所有人都承认那恩赐是圣神的恩赐。但自那话所说的日子起，这恩宠就受限、离去并从地上绝迹了：\BibleQuote{你们的房屋要被撇下，成为荒场}\BibleRef{玛 23:38}；而且那日之前它的匮乏就已经开始了，因为他们中间不再有先知，恩宠也不临到他们的圣物。既然圣神曾被扣留，但将来要丰沛地倾注，而这赋予的开端是在十字架受难之后，不仅在于其丰沛，也在于恩赐的更大提升（因为那恩赐更奇妙，如经上所说：\BibleQuote{你们不知道你们是什么精神}\BibleRef{路 9:55}；又说：\BibleQuote{因为你们所领受的，不是奴隶的圣神，而是义子的圣神}\BibleRef{罗 8:15}；而古时的人自己拥有圣神，却不分施给别人，宗徒们却用圣神充满成千上万的人），既然，我说，他们要领受这恩赐，但它尚未赐下，所以祂加上了：\BibleQuote{圣神还没有赐下}。既然主论及这恩宠，福音作者说了：\BibleQuote{因为圣神还没有赐下}，也就是说：\Quote{还没有赐下}。\par

\BibleQuote{因为耶稣还没有受光荣}\par

称十字架为\Quote{光荣}。因为既然我们是仇敌，犯了罪，亏缺了天主的恩赐，又是憎恨天主的，并且恩宠是我们修和的证明，而恩赐不赐给受憎恨的人，只赐给朋友和蒙悦纳的人；所以有必要先为我们献上祭献，好消除那在我们肉体中（对天主）的仇恨，使我们成为天主的友人，从而领受这恩赐。因为如果这事是关于对亚巴郎的恩许而做的，那么更关于恩宠。保禄也说明了这一点，说：\BibleQuote{如果承继人出于法律，信德就落了空——因为法律激起忿怒。}\BibleRef{罗 4:14} 他所说的是这样的：天主恩许将土地赐给亚巴郎和他的后裔；但他的后裔不配这恩许，凭自己的行为不能蒙天主悦纳。因此，信德来了，这是一项容易的举动，好吸住恩宠，不使恩许落空。经上又说：\par

\BibleQuote{因此，它是出于信德，为的使其出于恩宠，好使恩许得以稳固。}\BibleRef{罗 4:16} 因此它是出于恩宠，因为凭他们自己的劳苦，他们未能成功。\par

但为什么说了\Quote{按经上所载}之后，祂没有加上见证呢？因为他们的心思败坏；因为，\par

% structure: p0011 source=h2 target=subsection reason=relative-heading-rank; base=h2

\subsection{若望福音 7:40-42}

其他人说：\BibleQuote{当基督来时，没有人知道祂从哪里来}\BibleRef{若 7:27}；于是意见分歧，正如在混乱的人群中可预期的；因为他们不是专心地听祂的话，也不是为学习。因此祂不回答他们；但他们又说：\Quote{基督莫非来自加里肋亚吗？}祂曾称赞纳塔乃耳是\Quote{一个真以色列人}，而纳塔乃耳以更有力和显着的方式说过：\BibleQuote{纳匝肋能出什么好的吗？}\BibleRef{若 1:46} 但这些人，以及那些对尼苛德摩说：\BibleQuote{你去查考，看看从加里肋亚不会出先知的}\BibleRef{若 7:52} 的人，说这话并非为学习，而只是推翻关于基督的见解。纳塔乃耳说这话，是因为他是一个爱真理的人，并精确知道所有古代历史；但他们只注意一件事，就是消除祂是基督的见解，因此祂没有向他们启示什么。因为那些甚至自相矛盾，一时说\Quote{没有人知道祂从哪里来}，另一时又说\Quote{从白冷}的人，显然即使被告知，也会反对祂。因为就算他们不知道祂出生的地方，即祂来自白冷，因为祂居住在纳匝肋（但这不能承认，因为祂不是在那里出生的），难道他们也不知道祂的族系，即祂是\Quote{达味的家和后裔}吗？那么他们怎么说：\BibleQuote{基督不是出自达味的后裔吗？}\BibleRef{若 7:42} 因为他们想用那个问题来隐藏这个事实，他们所说的一切都带着恶意。他们为什么不来到祂面前说：\Quote{我们既在其它方面钦佩祢，祢又叫我们按经书相信祢，那么请告诉我们，经书怎么说基督必须来自白冷，而祢却来自加里肋亚呢？}但他们没有说任何这类话，一切都是出于恶意。为了表明他们说话不是追问，也不是愿意学习，福音作者立刻加上说：\par

% structure: p0013 source=h2 target=subsection reason=relative-heading-rank; base=h2

\subsection{若望福音 7:44}

这一点，即使没有别的，足以使他们感到愧疚，但他们不觉得，正如先知所说：\BibleQuote{他们虽被劈开，心却未受刺痛。}\BibleRef{咏 34:15}\par

3. 这就是恶意！它不会向任何事物让步，它只看一件事，就是要毁掉它所谋算的人。但圣经说什么？\Quote{凡为邻人掘坑的，必自陷其中。}\BibleRef{箴 26:27}当时的情况正是如此。因为他们想要杀害祂，以为这样就能阻止祂的宣讲；结果却适得其反。因为宣讲因基督的恩宠而兴盛，而属于他们的一切都熄灭和消亡了；他们失去了他们的国家、自由、安全、敬拜，被剥夺了一切繁荣，成了奴隶和俘虏。\par

既然如此，我们切勿谋害他人，要知道这样做是磨刀对准自己，给自己造成更深的创伤。有人得罪了你，你想报复他吗？不要报复；这样你才能得到报复；但如果你报复，你就得不到报复。不要以为这是谜语，而是真实的话。\Quote{如何，以何种方式？}因为如果你不报复他，你就使天主成为他的敌人；但如果你报复，就不再如此。\BibleQuote{复仇在我，我必酬报，主说。}\BibleRef{罗 12:19}因为如果我们有仆人，他们彼此争吵，不将判断和惩罚留给我们，而是擅自处理；即使他们来找我们千万次，我们不仅不会为他们报仇，反而会生他们的气，说：\Quote{你这个逃奴，你这个挨鞭子的，你本应把一切都交给我们，既然你抢在我们前面报复了自己，就别再来烦扰我们了。}何况天主，祂命令我们把一切都交托给祂，岂不更是如此？因为当我们要求仆人有如此的智慧和服从，却不愿在我们的主面前顺服，而我们却希望家仆顺服我们，这岂不荒谬？我说这些是因为你们随时准备互相施加惩罚。真正的智者甚至不应该这样做，而应宽恕和赦免冒犯，即使没有那伟大报酬——得到赦免作为回报。因为，告诉我，如果你谴责一个犯罪的人，你为什么自己也犯罪，陷入同样的过失？他侮辱了你？不要再侮辱回去，否则你侮辱了自己。他打了你？不要再打回去，因为这样你们之间就没有区别。他使你烦恼？不要再使他烦恼，因为毫无益处，你反而会与那些委屈你的人平起平坐。因此，如果你以温良和柔和忍耐，你就能责备你的敌人，使他羞愧，使他厌倦愤怒。没有人以恶治恶，而是以善治恶。这些智慧准则出自一些外邦人；既然在愚昧的外邦人中间尚有如此智慧，让我们羞愧，不要显得不及他们。他们中许多人受伤害而忍耐；许多人被恶意控告而不为自己辩护；被谋害而施以恩惠。并且我们不得不担忧，怕他们中有的人在生活中被证明比我们更伟大，从而使我们的惩罚更严厉。因为当我们这些分享了圣神、期待天国、为天上之事追求智慧、（不）畏惧地狱、被命令成为天使、享受奥迹的人；当我们达不到他们所达到的德行时，我们还能有什么赦免呢？如果我们必须超越犹太人（因为，\Quote{\BibleQuote{除非你们的义德超过经师和法利塞人的义德，你们决进不了天国}\BibleRef{玛 5:20}}），那么更要超越外邦人；如果超越法利塞人，那就更要超越不信者。既然如果我们不超越犹太人的义德，天国就向我们关闭，那么当我们证明自己比外邦人更差时，我们将如何能进入呢？那么让我们除去一切苦毒、愤恨和恼怒。讲说\BibleQuote{同样的事，对我并不为难，对你们却是稳妥的}\BibleRef{斐 3:1}。因为医生也常常使用同样的疗法，我们也不会停止向你们耳中讲说同样的事，提醒、教导、劝勉，因为世俗事物的喧嚣是巨大的，它使我们健忘，我们需要不断的教导。那么，为了使我们在此聚集不徒劳无益，让我们展示由行为证明的成果，使我们能获得将来的美物，借着我们的主耶稣基督的恩宠和慈爱，祂和祂与圣父及圣神，是荣耀，自今至永远，永世无穷。阿们。\par

\SourceRecord{Translated by Charles Marriott.；Nicene and Post-Nicene Fathers, First Series；Vol. 14.；Edited by Philip Schaff.；Buffalo, NY: Christian Literature Publishing Co.,；1889.；<http://www.newadvent.org/fathers/240151.htm>.}

% structure: p0001 source=h1 target=section reason=page-title

\section{若望福音讲道集 第五十二篇}

% structure: p0002 source=h2 target=subsection reason=relative-heading-rank; base=h2

\subsection{\BibleRef{若 7:45-46}}

1. 如果我们的行事不乖悖，真理便再清晰简明不过了；反之，如果我们行事乖悖，那就再困难不过了。因为你看，那些经师和法利塞人，他们自以为比他人更聪明，常常与基督同在，为的是图谋反对祂，目睹祂的奇迹，阅读圣经，却毫无获益，反而受到了损害；而那些差役，他们并没有这些优越条件，却被一次简短的讲演所折服；他们出发去捆绑祂，自己反而被惊讶所捆绑。我们不仅要惊叹他们的理解力——他们不需要神迹，单凭教训就被吸引（因为他们没有说：\Quote{从来没有一个人行过这样的奇迹}，而是说：\Quote{从来没有一个人讲话像这人那样讲话}）——我们还要惊叹他们的胆量，竟敢对派遣他们的人、对法利塞人、对祂的仇敌、对那些为发泄仇恨而行事的人说这样的话。圣史说：\Quote{差役回来了，法利塞人对他们说：\InnerQuote{你们为什么没有把祂带来？}}。他们\Quote{回来}，比起留在那里是更伟大的行为，因为留在那里他们可以摆脱这些人的烦扰；但现在他们却成了基督智慧的传讯者，并更加彰显了他们的胆量。他们不说：\Quote{我们无法通过人群，因为人们把祂当作先知听从}，而是说：\Quote{从来没有一个人像这人那样讲话}。然而他们本可以这样说，但他们却表明了自己正直的情感。因为这句话不仅出于钦佩祂的人，也是在责备他们的主人——因为他们派遣他们去捆绑那本应聆听的人。而且他们也没有听过一篇讲道，只是一篇简短的讲道；因为当心平气和时，不需要冗长的论证。真理就是这样。那么法利塞人说什么呢？当他们本应心中刺痛时，反而向差役回敬一个指控，说：\par

% structure: p0004 source=h2 target=subsection reason=relative-heading-rank; base=h2

\subsection{\BibleRef{若 7:47}}

他们对差役说话还算和气，没有严厉表达，担心他们完全离开自己，但尽管如此，他们仍然表现出愤怒的迹象，说话有所保留。因为当他们本应问祂说了什么，并对这些话感到惊奇时，他们却没有这样做（知道他们可能会被吸引），而是从一个非常愚蠢的理由与他们辩论；\par

% structure: p0006 source=h2 target=subsection reason=relative-heading-rank; base=h2

\subsection{\BibleRef{若 7:48}}

那么，请问，你们这是以此指控基督，而不是指控那些不信的人吗？\par

% structure: p0008 source=h2 target=subsection reason=relative-heading-rank; base=h2

\subsection{\BibleRef{若 7:49}}

那么，指控落在你们身上更重，因为百姓相信了，而你们却不相信。他们行事如同知道法律的人；那么他们怎么是被诅咒的呢？倒是你们是被诅咒的，因为你们不遵守法律，而不是他们，他们遵守法律。并且，凭着不信者的见证来毁谤他们不信的人，是不正确的，因为这是一种不公义的行事方式。因为你们也不信天主，正如保禄所说的：\BibleQuote{即使有人不信，又有什么关系呢？难道他们的不信能使天主的信实失效吗？绝对不可！} \BibleRef{罗 3:3}。因为先知们也曾责备他们说：\BibleQuote{索多玛的统治者，请听}；以及\BibleQuote{你们的统治不法} \BibleRef{依 1:10}；并且又说：\BibleQuote{难道你们不应知道审判吗？} \BibleRef{米 3:1}。他们在各处都猛烈地攻击他们。那么怎样？有人会为此责怪天主吗？绝不可这样想。这责怪是他们的。除了你们不遵守法律之外，还有什么别的证据能让人证明你们不知道法律呢？因为他们曾说过：\Quote{难道有统治者信了祂吗？}和\Quote{这些不知道法律的人}，尼苛德摩公平地责备他们说：\par

% structure: p0010 source=h2 target=subsection reason=relative-heading-rank; base=h2

\subsection{\BibleRef{若 7:51}}

他指出他们既不知道法律，也不行法律；因为如果那条法律命令不先听取一个人的意见就不能杀人，而他们在听取之前就急于行事，那么他们就是法律的违反者。因为他们说过：\Quote{难道有统治者信了祂吗？} \BibleRef{若 7:50}，因此圣史告诉我们尼苛德摩是\Quote{他们中的一个}，以表明连统治者都信了祂；虽然他们尚未表现出应有的胆量，但他们仍然渐渐依附于基督。请注意他多么谨慎地责备他们；他没有说：\Quote{你们想要杀死祂，在毫无证据的情况下就判定这人为骗子}；而是用更温和的方式说话，阻止他们过度的暴力和不假思索、嗜杀的性情。因此他将话题转向法律，说：\Quote{除非先仔细听取祂，并知道祂所做的事。}因此，要求的不是单纯的\Quote{听取}，而是\Quote{仔细听取}。因为\Quote{知道祂所做的事}的意思是\Quote{祂的意图}、\Quote{出于什么原因}、\Quote{为了什么目的}、\Quote{是为了颠覆秩序还是作为仇敌}。因此，他们因祂曾说过\Quote{法律不审判任何人}而感到困惑，于是他们既非激烈、也非容忍地回应祂。因为请告诉我，在祂说了\Quote{法律不审判任何人}之后，他们怎么会接着说：\par

% structure: p0012 source=h2 target=subsection reason=relative-heading-rank; base=h2

\subsection{\BibleRef{若 7:52}}

2. 当他们本应表明没有未经审判就派人去召祂，或者允许祂说话是不合适的，他们却用粗暴和愤怒的方式来回应。\par

\Quote{你们去查考吧！因为从\Place{加里肋亚}从来没有出过先知。}\par

那人说了什么呢？是基督是先知吗？不；他说，祂不应未经审判就被杀害；但他们却傲慢地回答，如同对一个对圣经一无所知的人；好像有人说：\Quote{去学吧}，这就是\Quote{去查考吧}的意思。那么基督做了什么？既然他们不断强调\Place{加里肋亚}和\Quote{那位先知}，为了将所有人从这种错误猜疑中解放出来，并表明祂不是先知中的一位，而是世界的主宰，祂说：\par

% structure: p0016 source=h2 target=subsection reason=relative-heading-rank; base=h2

\subsection{\BibleRef{若 8:12}}

不是\Quote{从\Place{加里肋亚}}，也不是从\Place{巴勒斯坦}，也不是从\Place{犹太}。那么犹太人说什么呢？\par

% structure: p0018 source=h2 target=subsection reason=relative-heading-rank; base=h2

\subsection{\BibleRef{若 8:13}}

哀哉！因他们的愚昧，祂不断引他们到圣经，如今他们却说：\Quote{祢为自己作证。}祂作了什么证？\Quote{我是世界的光。}这是一句伟大的话，确实伟大，但并未使他们大为惊讶，因为祂没有现在使自己与父平等，也没有声称祂是圣子，也没有说祂是天主，而只是暂时称自己为\Quote{光。}他们确实也想反驳这一点，然而这比说以下的话要伟大得多：\par

\Quote{跟随我的，决不在黑暗中行走。}\par

祂在灵性的意义上使用了\Quote{光}和\Quote{黑暗}这两个词，意指\Quote{不陷于错误。}在此处祂引述尼苛德摩，并把他作为曾非常勇敢发言的人引出来，又称赞那些也这样做的仆人。因为\Quote{大声呼喊}是渴望使他们也听见之人的行动。同时祂暗示那些暗中策划背叛的人，他们既在黑暗又在错误中，但他们不能胜过光。祂又提醒尼苛德摩祂从前说过的话：\BibleQuote{凡作恶的，恨光，也不来就光，怕他的行为被责备。}\BibleRef{若 3:20}因为他们曾断言没有一位首领信祂，所以祂说：\Quote{凡作恶的，不来就光，}以表明他们不来不是因为光弱，而是因为他们自己的邪恶意愿。\par

\Quote{他们回答祂说：祢为自己作证吗？}\par

祂怎样回答呢？\par

% structure: p0024 source=h2 target=subsection reason=relative-heading-rank; base=h2

\subsection{若 8:14}

祂从前说的话，这些人提出好像特别肯定一样。那么基督怎样呢？为了驳斥这一点，并表明祂用那些话是为了适应他们和他们的猜疑（他们以为祂只是一个普通人），祂说：\Quote{虽然我为自己作证，我的见证是真的，因为我知道我从哪里来。}这是什么意思？\Quote{我属于天主，我是天主，是圣子，而天主自己本身就是自己忠实的见证，但你们不认识祂；你们甘愿错误，明知假装不知，只说你们凭纯粹人的想象所说的话，选择不理解任何超越所见的事。}\par

% structure: p0026 source=h2 target=subsection reason=relative-heading-rank; base=h2

\subsection{若 8:15}

正如随从肉身生活就是生活得糟糕，同样随从肉身判断就是判断不义。\Quote{但我不判断任何人。}\par

% structure: p0028 source=h2 target=subsection reason=relative-heading-rank; base=h2

\subsection{若 8:16}

祂所说的是这样的：\Quote{你们判断不义。}\Quote{如果，}某人说，\Quote{我们判断不义，为什么祢不责备我们？为什么祢不惩罚我们？为什么祢不谴责我们？}\Quote{因为，}祂说，\Quote{我不是为此而来的。}这就是\Quote{我不判断任何人；然而如果我判断，我的判断是真的}的意思。\Quote{因为如果我愿意判断，你们就会在定罪之列。我说这话，并不是判断你们。我告诉你们我说这话，不是判断你们，并不是我不自信如果我判断你们，我本会定你们的罪；因为如果我判断了你们，我必然公义地谴责你们。但现在还不是审判的时候。}祂也暗示了将来的审判，说：\par

\Quote{我不是独自一人，而是我与那派遣我的圣父在一起。}\par

在此祂暗示，不是祂独自谴责他们，而是圣父也谴责。然后祂通过引导他们到祂自己的见证，隐藏了这一点。\par

% structure: p0032 source=h2 target=subsection reason=relative-heading-rank; base=h2

\subsection{若 8:17}

3. 异端者在这里会说什么？他们会说：\Quote{如果我们简单地接受祂所说的，祂比人好在哪裏？因为这条规则是针对人而定的，因为没有人独自可信。但在天主的情况下，怎能容忍这样的说法？那么，\InnerQuote{两个}这个词如何使用？是因为他们是两个，还是因为他们是人所以是两个？如果是因为他们是两个，为什么祂不求助於若翰，说：\InnerQuote{我为自己作证，若翰也为我作证}？为什么不去找天使？为什么不去找先知？因为祂可以找到一万个其他见证。}但祂想要表明的不仅是存在两位，而且是他们具有同一本体。\par

% structure: p0034 source=h2 target=subsection reason=relative-heading-rank; base=h2

\subsection{若 8:19}

由于他们明知却装作不知，好像试探祂，祂甚至不屑回答他们。因此从此以后祂更加清晰、更加大胆地讲话；从祂的奇迹和跟随祂的人的教导中引出见证，并因十字架临近而引出见证。因为，\BibleQuote{我知道，}祂说，\BibleQuote{我从哪里来。} 这话对他们影响不大，但加上\BibleQuote{以及我往哪里去，}会更令他们恐惧，因为祂不会停留在死亡中。但祂为什么不说\BibleQuote{我知道我是天主，}而说\BibleQuote{我知道我从哪里来}呢？祂总是将卑微的话语与崇高的话语混合，甚至遮掩这些话语。因为在说了\BibleQuote{我为自己作证，}并证明了这一点之后，祂降到了更谦卑的语调。好像祂说：\BibleQuote{我知道我从谁那里被派遣，并向谁而去。} 因为这样，当他们听到祂是从祂那里被派遣，并将回到祂那里时，他们就无话可说了。\BibleQuote{我不能说，}祂说，\BibleQuote{任何谎言，我从那里来，也往那里去，到真实的天主那里。} 但你们不认识天主，所以按肉身判断。因为如果听到这么多确凿的迹象和证据，你们仍然说\BibleQuote{你的见证不真实，}如果你们认为梅瑟可信，无论他对别人说的还是对自己说的，但对基督却不这样，这就是按肉身判断。 \BibleQuote{但我不判断任何人。}祂也确实说过：\BibleQuote{圣父不判断任何人。}\BibleRef{若 5:22} 那么，祂如何在这里宣告说，\BibleQuote{如果我判断，我的判断是真的，因为我不是独自一人}呢？祂再次回应他们的思想。\BibleQuote{我的判断就是圣父的判断。圣父判断时，不会作出与我不同的判断，我也不会作出与圣父不同的判断。} 为什么祂提到圣父？因为他们不会认为圣子值得相信，除非祂得到圣父的见证。此外，这话甚至不成立。因为对于人来说，当两个人对涉及另一个人的事情作证时，他们的见证是真的（这是两个人作证）；但如果一个人为自己作证，那么就不再是两个人了。你是否看到祂说这些话只是为了表明祂是同体的，祂不需要别的见证，并且丝毫不逊于圣父？至少要看到祂的独立性；\par

% structure: p0036 source=h2 target=subsection reason=relative-heading-rank; base=h2

\subsection{若 8:18}

如果祂是次等的实体，祂就不会提出这一点。但现在，为了不让你以为圣父被包括进来以凑足数目（两个），请注意祂的能力与（圣父的）毫无差异。一个人作证时，他自己是可信的，而不是他自己需要见证，并且那是在涉及另一个人的事上；但在涉及他自己的事上，当他需要他人的见证时，他并不可信。但在此情形中，完全相反。因为祂虽然为自己作证，并说别人为祂作证，却宣称自己是可信的，各方面都显示出祂的独立性。因为为什么，当祂说了\BibleQuote{我不是独自一人，而是我与派遣我的圣父在一起，}和\BibleQuote{两个人的见证是真的，}之后，祂没有保持沉默，反而补充说\BibleQuote{我是为自己作证的那一位}呢？这显然是为了显示祂的独立性。这里祂显示了祂的同等尊荣，并且他们自称认识圣父却并不认识祂，这毫无益处。祂说，这种（无知）的原因是他们不愿意认识祂。因此祂告诉他们，不认识祂就不可能认识圣父，即使如此，祂仍可能引导他们认识祂。因为既然他们离开祂，甚至寻求认识圣父，祂说：\BibleQuote{你们若不认识我，就不能认识圣父。}\BibleRef{若 8:19} 所以，那些亵渎圣子的人，不仅亵渎圣子，也亵渎了生了祂的那一位。\par

4. 让我们避免这事，并光荣圣子。倘若祂不是同一本性，就不会如此说话。因为若祂只是教导，而实体不同，人可能不认识祂，却认识圣父；再者，认识祂的人，也不会完全认识圣父；因为认识一个人，并不就是认识一位天使。\Quote{是的，}有人回答说，\Quote{认识受造物的人就认识天主。}绝非如此。许多人，或者我该说，所有人都认识受造物（因为他们看见它），但他们却不认识天主。那么，让我们光荣天主圣子，不仅用（言语的）光荣，也要用行为的光荣。因为前者没有后者便毫无价值。\Quote{你看，}圣保禄说，\Quote{你称为犹太人，依赖法律，并夸耀于天主——那么，你教导别人，却不教导自己吗？你夸耀于法律，却因触犯法律而羞辱天主吗？}\BibleRef{罗 2:17} 我们也当谨慎，免得我们以信仰的正统自夸，却因没有活出与信仰相称的生活而羞辱天主，致使祂受亵渎。因为祂愿意基督徒作世界的导师、酵母、盐和光。那光是什么？就是照耀且有黑暗不在其中的生命。光不为自己有用，酵母和盐也是如此，而是对他人显示其效用，因此我们被要求行善，不仅为自己，也为他人。因为盐若不成盐，就不是盐。此外，还有一点是明显的：如果我们为义，别人也必然会如此；但只要我们不为义，我们就不能帮助别人。在我们中间不要有任何愚蠢或愚妄的事；属世的事正是如此，此生的挂虑正是如此。因此童女被称为愚拙的，因她们忙于愚蠢、属世的事，在这里聚集财物，却没有在应当的地方积蓄财宝。恐怕我们也是如此，恐怕我们也穿着污秽的衣服离开，到那人人衣服明亮发光的地方去。因为没有什么比罪恶更污秽、更不洁的了。因此先知宣告其本质时大声说：\Quote{我的伤口发臭流脓。}\BibleRef{咏 37:5} 若你愿意充分认识罪恶是何等恶臭，要在它做成之后思想它；当你脱离了欲望，当那火不再烦扰你时，那时你便会看见罪恶是什么。当你平静时思想忿怒，当你不觉着它时思想贪婪。没有什么比掠夺和贪婪更可耻、更可咒骂的了。我们不断这样说，不是要烦扰你们，而是要获得一些伟大而奇妙的益处。因为那一次听后没有正确行事的人，或许第二次听后就会做；而第二次忽略的人，第三次后可能会改正。愿天主恩赐我们，使我们摆脱一切邪恶，拥有基督的馨香；因为光荣归于祂，偕同圣父及圣神，从今世直到永远。阿们。\par

\SourceRecord{Translated by Charles Marriott.；Nicene and Post-Nicene Fathers, First Series；Vol. 14.；Edited by Philip Schaff.；Buffalo, NY: Christian Literature Publishing Co.,；1889.；<http://www.newadvent.org/fathers/240152.htm>.}

% structure: p0001 source=h1 target=section reason=page-title

\section{论\BibleBook{若望}{福音}讲道集第五十三篇}

% structure: p0002 source=h2 target=subsection reason=relative-heading-rank; base=h2

\subsection{若 8:20}

1. 犹太人何等愚妄！逾越节前他们寻找祂，后来在他们中间找到了祂，屡次试图亲手或借他人之手捉拿祂却不能；即便如此，他们仍未被祂的大能所震慑，反而执意行恶，不肯罢休。因为经上记载，他们不断尝试：\BibleQuote{这些话是耶稣在圣殿里，在银库旁教导人时说的；没有人下手拿祂。}祂在圣殿里，以教师的身份说话，这更易激起他们的愤怒，祂所说的那些话正是刺伤他们的，他们控告祂使自己与圣父平等。因为\BibleQuote{两个人的见证是真的}证实了这一点。然而经上仍说：\BibleQuote{祂说了这些话}，在圣殿里，以教师的身份，\BibleQuote{没有人下手拿祂，因为祂的时候还没有到}；这就是说，祂被钉十字架的时刻尚未到来。由此可见，那件事的发生并非出于他们的能力，而是出于祂的安排，因为他们早已渴望，却不能得逞，即便那时，若非祂同意，他们也无法下手。\par

% structure: p0004 source=h2 target=subsection reason=relative-heading-rank; base=h2

\subsection{若 8:21}

为什么祂不断说这话？为要羞愧并震慑他们的心灵；请看这话如何在他们心中引起恐惧。尽管他们想要杀死祂以摆脱祂，他们却仍问\Quote{祂往哪里去}，从这事他们想象了何等大事。祂也想向他们显明另一件事：这事不会因他们的强力而成就；但祂事先以比喻向他们显示，并已借着这些话预言了复活。\par

% structure: p0006 source=h2 target=subsection reason=relative-heading-rank; base=h2

\subsection{若 8:22}

那么基督做什么？为消除他们的猜疑，并表明这样的行为是罪，祂说：\par

% structure: p0008 source=h2 target=subsection reason=relative-heading-rank; base=h2

\subsection{若 8:23}

祂所说的，意思是：你们想象这样的事并不奇怪，你们是属肉性的人，没有属神的思想；但我不会做任何这类的事，因为：\par

\BibleQuote{我是从上而来的；你们是出于这世界的。}\par

祂在此又谈论他们属世和属肉性的想象，由此可知，\BibleQuote{我不是出于这世界}并非指祂没有取了肉身，而是指祂远离他们的邪恶。因为祂甚至说，祂的门徒\BibleQuote{不属于世界}\BibleRef{若 15:19}，然而他们是有肉身的。因此，正如保禄所说：\BibleQuote{你们不在肉性内}\BibleRef{罗 8:9}，并非指他们无形体，同样，基督说祂的门徒\BibleQuote{不属于世界}，无非是见证他们属天的智慧。\par

% structure: p0012 source=h2 target=subsection reason=relative-heading-rank; base=h2

\subsection{若 8:24}

因为祂来是要除去世人的罪，而人若不藉着圣洗，无法脱去那罪，所以不信的人必然带着旧人离去；因为凡不愿藉着信德杀死并埋葬那旧人的，必死在旧人中，并要往那地方去，受他先前罪过的惩罚。因此祂说：\BibleQuote{不信的人已被定罪}\BibleRef{若 3:18}；这不仅仅是因为他不相信，而是因为他离开此地时仍带着他先前的罪过。\par

% structure: p0014 source=h2 target=subsection reason=relative-heading-rank; base=h2

\subsection{若 8:25}

何等愚妄！经过如此长的时间、如此多的征兆和教导，他们竟问：\BibleQuote{你是谁？}基督怎样回答？\par

\BibleQuote{我从起初就告诉你们的。}\par

祂所说的意思是：\Quote{你们根本不配听我的话，更不配知道我是谁，因为你们所做的一切都是试探我，不留意我的任何话语。这一切我如今都可以证明你们的不是。}因为这就是以下经文的意思：\par

% structure: p0018 source=h2 target=subsection reason=relative-heading-rank; base=h2

\subsection{若 8:26}

\Quote{我不但可以证明你们有罪，而且可以惩罚你们；但派遣我来的那位，即圣父，不愿意这样。因为我来不是为审判世界，乃是为拯救世界，因为天主派遣子来不是为审判世界，乃是为叫世界藉着祂得救。\BibleRef{若 3:17}既然祂为此差遣我，而祂是真实的，那么我有理由如今不审判任何人。但我所说这些话，是为你们的得救，并非为你们的定罪。}\par

% structure: p0020 source=h2 target=subsection reason=relative-heading-rank; base=h2

\subsection{若 8:27}

何等愚妄！祂不停地谈论圣父，他们却不认识祂。然后，在行了许多征兆、教导了他们之后，仍不能吸引他们归向自己，祂便对他们讲论十字架，说：\BibleQuote{当你们高举人子以后，你们便知道我就是那一位，我凭自己不做任何事；但我所说的这些，是圣父教导我的。那派遣我来的与我同在；圣父没有留下我独自一人，因为我常做那些祂所喜悦的事。\BibleRef{若 8:28-29}}\par

% structure: p0022 source=h2 target=subsection reason=relative-heading-rank; base=h2

\subsection{若 8:28-29}

2. 祂表明祂所说\Quote{我从起初就对你们说的}这句话是对的。他们何其不留意祂的话。\Quote{当你们高举了人子以后。}\Quote{你们难道不以为那时你们必能摆脱我、杀害我吗？但我告诉你们，那时你们必因奇迹、复活和（耶路撒冷的）毁灭而最清楚地知道我是谁。}因为这一切都足以彰显祂的能力。祂没有说\Quote{那时你们将知道我是谁}；因为祂说，\Quote{当你们看见我从死亡中丝毫未受损害时，你们就知道我是谁，即基督、天主子，治理万有，并与祂不相反对。}因此祂加上：\Quote{我凭自己什么都不说。}因为你们将知道我的能力以及我与圣父的一致。因为\Quote{我凭自己什么都不说}这句话表明祂的本体与（圣父的）并无不同，并且祂所说的没有超出圣父心意之外的。\Quote{因为当你们被赶出敬拜之地，甚至不再被允许像从前一样侍奉祂时，那时你们就知道祂这样做是为报复我，并因恼恨那些不肯听我的人而发怒。}祂仿佛说：\Quote{倘若我是天主的仇敌和陌路人，祂就不会向你们激起这样的愤怒。}这也正是依撒意亚所宣告的：\Quote{祂将恶人葬于其坟}\BibleRef{依 53:9}；以及达味：\Quote{那时祂要在震怒中对他们说话}\BibleRef{咏 2:5}；还有基督自己：\Quote{看，你们的房屋将荒凉}\BibleRef{玛 23:38}。祂的比喻也宣告同样的事，当祂说：\Quote{那葡萄园的主人要怎样处置那些园户呢？他要凶恶地消灭那些恶人。}\BibleRef{玛 21:40}你看，祂处处这样说，是因为祂还未被相信吗？但倘若祂要消灭他们，正如祂必会做的（因为经上说：\Quote{把那不愿意我作王的人带来，杀掉他们}），为何祂说这事不是祂做的，而是祂的父做的呢？祂是迁就他们的软弱，同时也尊荣那位生祂的。因此祂没有说：\Quote{我使你们的房屋荒凉}，而是说它\Quote{被撇下了}，祂用了无人称的说法。但祂说：\Quote{我多次愿意聚集你的子女……你们却不愿意}，然后加上\Quote{被撇下了}，这表明是祂成就了这荒凉。\Quote{因为}，祂告诉他们，\Quote{当你们蒙恩、疾病得医治时，你们不肯认识我，你们将要因受惩罚而知道我是谁。}\par

\Quote{圣父与我同在。}为使不把\Quote{那派遣我的}视为低微的标记，祂说：\Quote{与我同在}；前者属于救恩安排，后者属于天主性。\par

\Quote{祂没有留下我独自一人，}因为我常做那些讨祂喜悦的事。\par

祂再一次将祂的言论降到更谦卑的调子，不断反对他们所宣称的，即祂不是出自天主，且不守安息日。对此祂回答说：\Quote{我常做那些讨祂喜悦的事}，表明连安息日被破坏也讨祂喜悦。例如，就在十字架苦刑前，祂说：\Quote{你想我不能请求我的父吗？}\BibleRef{玛 26:53}然而仅仅说\Quote{你们找谁？}\EditorialNote{若 18:4、6}祂就使他们退后跌倒。那么为什么祂不说：\Quote{你想我不能消灭你们吗？}而祂已经用事实证明了这事？祂迁就他们的软弱。因为祂费了很大力气表明自己没有做任何违反圣父的事。因此祂说话更像人的方式；正如\Quote{祂没有留下我独自一人}所说的，同样\Quote{我常做那些讨祂喜悦的事}也是这么说。\par

% structure: p0027 source=h2 target=subsection reason=relative-heading-rank; base=h2

\subsection{若 8:30}

当祂把言论降到谦卑的调子时，许多人相信了祂。你还问为什么祂说话谦卑？然而福音使徒清楚地暗示了这一点，当他说：\Quote{祂说这些话时，许多人信了祂。}这几乎是向我们大声宣告：\Quote{哦，听者，你若听到任何卑微的言辞不要困惑，因为那些经过如此高深教导后仍不信祂来自圣父的人，理当听到更谦卑的话，以便他们相信。}这也是对那些将要谦卑说的话的辩解。他们那时信了，却并非如所当信的那样，而是漫不经心、偶然似的，因言语的谦卑而喜悦和舒快。因为他们的信德不完全，福音使徒通过此后他们的言论表明出来，他们再次侮辱祂。而且使徒也通过说：\par

% structure: p0029 source=h2 target=subsection reason=relative-heading-rank; base=h2

\subsection{若 8:31}

表明他们还未接受祂的教导，只是留意祂的话。因此祂说话更严厉。先前祂只说：\Quote{你们将找我}\BibleRef{若 7:34}，但现在祂加上更严厉的：\Quote{你们将死在你们的罪中。}\BibleRef{若 8:21}并且祂表明如何：\Quote{因为你们不能事后到那地方求我。}\par

\Quote{这些话是我对世界说的。}这些话表明祂现在要往外邦人那里去。但因他们还不知道祂是论及圣父对他们说话，祂又提到祂，而福音使徒已道出了这些卑微言辞的理由。\par

3. 如果我们现在这样仔细而不疏忽地查考圣经，我们就能获得救恩；如果我们持续研读圣经，我们就能学到正确的教义和完美的生活。因为即使一个人非常顽固、倔强、骄傲，在其他时候一无所获，至少此时他也能获得果实，并受益，即使这益处不足以让他察觉，他仍会得到它。因为如果一个人经过或坐在膏油店旁，即使不情愿也会沾上香气，那么来到教堂的人更是如此。因为正如懒惰生于懒惰，勤勉也生于劳作。即使你充满一万个罪，即使你不洁，也不要回避在此停留。\Quote{所以，}有人说，\Quote{我听了却不做吗？}这并非小益：认为自己可怜并非无用，这种恐惧并非无益，这种畏惧并非不合时宜。只要你叹息说\Quote{我听了却不做}，你终有一天也会付诸实践。因为与天主交谈、听天主说话的人不可能不获益。当我们想拿起圣经时，我们会立刻端正自己并洗手。你看到在阅读之前就有如此敬意吗？如果我们严谨进行，我们将获得巨大益处。因为若非为使灵魂保持敬畏，我们不会洗手；女人若未戴面纱会立刻戴上，显示内在的敬畏；男人若戴帽会脱帽。你看到外在行为如何宣告内在敬畏吗？此外，坐着聆听的人常常叹息，并谴责他目前的生活。\par

让我们，亲爱的，留意圣经；即使其他部分不被如此重视，至少让福音成为我们认真关注的对象，让我们把它们拿在手中。因为当你打开书卷时，你会立刻看到基督的名字，并听到有人说：\Quote{耶稣基督的诞生是这样的：祂的母亲\Person{玛利亚}许配给\Person{若瑟}后，因圣神怀孕。}\BibleRef{玛 1:18}听到这的人会立刻渴望贞洁，惊叹于降生，摆脱世俗之事。当你看到童贞女被圣神眷顾，天使与她谈话时，这不是小事。而这只是在表面；如果你坚持读到最后，你会厌弃一切与此生相关的事，嘲笑所有世俗之物。如果你是富人，当你听到她是木匠之妻、出身卑微，却成为你主的母亲时，你会轻视财富。如果你是穷人，当你听到世界的创造者不以最卑微的居所为耻时，你不会因贫穷而羞愧。想到这些，你不会抢劫，不会贪恋，不会夺取他人的财物，反而会喜爱贫穷，鄙视财富。如果是这样，你将消除一切罪恶。再者，当你看到祂躺在马槽里，你不会急于为你的孩子穿上金色衣服，或用银子镶嵌你妻子的床榻。如果你不关心这些，你也不会做因它们引起的贪婪和掠夺之事。你还能获得许多我无法一一列举的好处，但经历过的人会知道。因此，我劝你们既得到圣经，又连同圣经一起持守它们所阐述的情感，并将它们写在你们心中。犹太人因不留意而被吩咐将经卷挂在手上；但我们甚至不把经卷拿在手上，而是放在家里，而我们应该把它们刻在心中。这样洁净我们目前的生活，我们必获得将来的福分；愿我们众人赖我主\Person{耶稣基督}的恩宠和慈爱，到达那里；借着祂，偕同祂，圣父和圣神获得光荣，于今世，于永世，直到永远。阿们。\par

\SourceRecord{Translated by Charles Marriott.；Nicene and Post-Nicene Fathers, First Series；Vol. 14.；Edited by Philip Schaff.；Buffalo, NY: Christian Literature Publishing Co.,；1889.；<http://www.newadvent.org/fathers/240153.htm>.}

% structure: p0001 source=h1 target=section reason=page-title

\section{《若望福音》第五十四篇讲道}

% structure: p0002 source=h2 target=subsection reason=relative-heading-rank; base=h2

\subsection{\BibleRef{若 8:31-32}}

1. 亲爱的诸位，我们的处境需要极大的坚忍；而坚忍产生于教义深深扎根之时。因为正如没有风能借其冲击将橡树连根拔起，它向下扎根于地底深处并牢牢紧握在那里；同样，那被敬畏天主钉住的灵魂，无人能将其倾覆。因为被钉住比被扎根更为坚固。因此先知祈祷说：\BibleQuote{求祢用敬畏钉住我的肉}\BibleRef{咏 118:120}；\Quote{求祢这样固定我、结合我，如同钉入我内的一根钉子。}因为这类人难以被俘获，相反的那类人则是现成的猎物，容易被推倒。正如当时犹太人的情况；他们在听了并相信之后，又转离了正路。因此基督渴望加深他们的信德，使之不流于表面，就用更锐利的话语深挖他们的灵魂。因为信徒的本分是甚至忍受责备，但他们立刻发怒。但祂是怎样做的呢？祂首先告诉他们：\BibleQuote{如果你们存留在我的话内，你们就真是我的门徒；并且真理将使你们自由。}这几乎是说：\Quote{我即将作一个深的切割，但你们不要动摇}；或者更确切地说，通过这些话语祂平息了他们想象力的骄傲。\BibleQuote{将使你们自由}：从什么中，告诉我？从你们的罪中。那么那些自夸者说什么呢？\par

% structure: p0004 source=h2 target=subsection reason=relative-heading-rank; base=h2

\subsection{\BibleRef{若 8:33}}

立刻他们的想象沉下去了，这是因为他们因世俗之事而烦乱。\BibleQuote{如果你们存留在我的话内，}这话是那宣告祂内心所存、且知道他们确实信了但没有存留的一位说的。祂应许了一件大事，即他们将成为祂的门徒。因为此前有些人离开了祂，祂指着他们说：\BibleQuote{如果你们存留，}因为他们也听了、信了，却因不能存留而离开了。\BibleQuote{因为祂的门徒中有许多人退去了，不再和祂同行。}\BibleRef{若 6:66}\par

\BibleQuote{你们将认识真理，}即\BibleQuote{将认识我，因为我是真理。一切犹太人的事物都是预像，但你们将从我认识真理，它将使你们从你们的罪中获得自由。}正如祂对那些人说的：\BibleQuote{你们将死在你们的罪中，}同样祂对这些说：\BibleQuote{将使你们自由。}祂没有说：\BibleQuote{我将从奴役中拯救你们，}这一点祂让他们自己去推想。那么他们说什么呢？\par

\Quote{我们是亚巴郎的后裔，从来没有给谁做过奴隶。}而如果他们必须被激怒，本该对祂话的前一部分，即祂所说的\BibleQuote{你们将认识真理}感到恼怒，并且会回答说：\Quote{什么！我们现在不已经认识真理了吗？难道法律和我们的知识是谎言吗？}但他们毫不在意这些事，他们为世俗的事忧伤，而这些就是他们对奴隶制的看法。当然即使在现在，仍有许多人对无关紧要的事、对这类奴役感到羞耻，却对罪的奴役毫无羞耻感，他们宁愿十万次被称作后一种奴役的仆人，也不愿一次被称作前一种的仆人。这些人就是这样，他们甚至不知道还有其他奴役，他们说：\Quote{你们称亚巴郎的后裔、那些高贵出生的人为奴仆，他们因此本不该被称为奴仆？因为，有人说，我们从来没有给谁做过奴隶。}这就是犹太人的夸口。\Quote{我们是亚巴郎的后裔，}\Quote{我们是以色列人。}他们从不提自己正义的行为。因此若翰向他们呼喊说：\BibleQuote{不要想自己有亚巴郎为父。}\BibleRef{玛 3:9}为什么基督没有反驳他们呢？因为他们曾多次给埃及人、巴比伦人和许多其他人做奴隶。因为祂的话不是为了为自己赢得荣耀，而是为了他们的救恩，为了他们的益处，祂正朝着这个目标迈进。因为祂本可以提那四百年，本可以提那七十年，本可以提民长时期奴隶的年份，有时二十年，有时两年，有时七年；祂本可以说他们从未停止做奴隶。但祂不愿意表明他们是人的奴隶，而是表明他们是罪的奴隶，这是最悲惨的奴役，只有天主才能释放；因为赦免罪不属于其他人。而这一点他们也是承认的。既然他们承认这是天主的作为，祂就把他们引到这一点，并说：\par

% structure: p0008 source=h2 target=subsection reason=relative-heading-rank; base=h2

\subsection{\BibleRef{若 8:34}}

表明祂所说的自由正是脱离这种服务的自由。\par

% structure: p0010 source=h2 target=subsection reason=relative-heading-rank; base=h2

\subsection{\BibleRef{若 8:35}}

祂也温和地由此贬低了法律的事，暗示先前的时代。因为免得他们跑回法律说：\Quote{我们有梅瑟所吩咐的祭献，它们能拯救我们，}祂加上了这些话，否则这话有什么关联呢？因为\BibleQuote{众人都犯了罪，缺乏天主的荣耀，因祂的恩宠白白成义}\BibleRef{罗 3:23}，连司祭本身也是如此。因此保禄也论到司祭说：\BibleQuote{他理当为人民也为自己的罪献祭，因为他自身也为弱点所困。}\BibleRef{希 5:3}而祂所说的\BibleQuote{奴隶不在家里永住}也表明了这一点。这里祂也显示了祂与父同等的尊荣，以及奴隶与自由人的区别。因为比喻有这个含义，即\Quote{奴隶没有权能}，这就是\Quote{不住}的含义。\par

2. 但为何在论及罪时祂提到了一座\Quote{房屋}？是为显示：如主人管理自己的房屋，祂也管理一切。至于\Quote{不常存}的意思是\Quote{没有能力施恩，因不是房屋的主人}；但圣子是房屋的主人。因为这就是\Quote{永远常存}，借用人间事物作比喻。免得他们说：\Quote{祢是谁？}\Quote{一切都是我的（祂说），因为我是子，住在我父的家里}，用\Quote{房屋}之名指祂的权能。正如在另一处祂称天国为父的家：\Quote{在我父的家里有许多住处。}\BibleRef{希 14:2} 由于谈论的是自由与束缚，祂合宜地使用这个比喻，告诉他们，他们没有能力释放人。\par

% structure: p0013 source=h2 target=subsection reason=relative-heading-rank; base=h2

\subsection{若望福音 8:36}

你见到圣子与圣父的\OriginalTerm{同体}{consubstantiality}吗？祂如何宣告祂拥有与圣父同等的权能？\Quote{若子使你们自由，以后无人反对，你们便有稳固的自由。}因为\Quote{使成义的是天主，谁能定罪呢？}\BibleRef{罗 8:33} 在此祂显示自己无罪，并暗示那种只徒具其名的自由；这种自由连人也给予，但那超性的自由惟独天主赐予。于是祂说服他们不以这种奴役为耻，而应以罪的奴役为耻。为了显示他们之所以是奴隶，是因为拒绝了那自由，祂更以这话表明他们是奴隶：\par

\Quote{你们必会真正地自由。}\par

这是宣告那自由并非真实的话。然后，免得他们说\Quote{我们没有罪}（因为他们很可能这样说），请看祂如何使他们陷入这项控诉。祂不逐一指摘他们的一生，只提出他们正在着手且渴望去做的事，说：\par

% structure: p0017 source=h2 target=subsection reason=relative-heading-rank; base=h2

\subsection{若望福音 8:37}

祂温柔而逐步地将他们从那种亲属关系中驱逐出去，教训他们不可因之骄傲。因为正如自由与束缚取决于人的行为，亲属关系也是如此。祂没有直接说\Quote{你们不是亚巴郎的子孙，你们是杀害义人的人}；但有一段时间祂甚至顺着他们说：\Quote{我知道你们是亚巴郎的子孙。}然而，这并非问题的关键，在余下的谈话中祂用了更大的力度。因为我们大致上可以观察到，当祂要行任何大事时，在行过之后，祂说话更为大胆，仿佛来自祂工作的见证堵住了人的口。\Quote{但你们却图谋杀害我。}\Quote{那又如何，}有人说，\Quote{如果他们出于公义而行？}但事实并非如此；因此祂也说明原因：\par

\Quote{因为我的话在你们内没有容身之处。}\par

\Quote{那么，}有人说，\Quote{他们怎么又相信了祂呢？}正如我先前所说，他们又改变了。因此祂严厉地触动他们：\Quote{如果你们夸耀亚巴郎的亲属关系，就该显示出他的生活。}祂没有说\Quote{你们不接受我的话}，而是说\Quote{我的话在你们内没有容身之处}，以此宣告祂教义的崇高。然而，即便如此，他们也不该杀害祂，反而应当尊敬并服侍祂，以便学习。\Quote{但是，}有人说，\Quote{如果这些话是祢凭自己说的呢？}为此祂补充道：\par

% structure: p0021 source=h2 target=subsection reason=relative-heading-rank; base=h2

\subsection{若望福音 8:38}

\Quote{正如，}祂说，\Quote{我既凭言语也凭真理显明圣父，你们也凭你们的作为显明你们的父。因为我不但拥有同等的\OriginalTerm{本体}{Substance}，还拥有与圣父同等的\OriginalTerm{真理}{Truth}。}\par

% structure: p0023 source=h2 target=subsection reason=relative-heading-rank; base=h2

\subsection{若望福音 8:39-40}

祂在此反复处理他们的杀人意图，并提及亚巴郎。祂这样做，是希望转移他们对这种亲属关系的注意，除去他们过度的自夸，并说服他们不再将救恩的希望寄托于亚巴郎和依本性的亲属关系，而是寄托于那依意愿的。因为阻碍他们归向基督的，正是他们认为那种亲属关系足以使他们得救。但祂所说的\Quote{真理}是什么？就是祂与圣父平等。正是为此，犹太人才图谋杀害祂；祂说：\par

\Quote{你们图谋杀害我，因为我将我从我父所听见的真理告诉了你们。}\par

为显示这些事与圣父并不相悖，祂再次转向圣父。他们对祂说：\par

% structure: p0027 source=h2 target=subsection reason=relative-heading-rank; base=h2

\subsection{若望福音 8:41}

3. \Quote{你们说什么？你们有天主为父，你们竟指责基督宣称这一点吗？}你看见祂是以一种特殊的方式说天主是祂的父吗？因此，当祂将他们逐出与亚巴郎的亲属关系后，他们无言以对，竟敢做更大的事，转向天主。然而，祂也将他们从这一尊荣中驱逐，说：\par

% structure: p0029 source=h2 target=subsection reason=relative-heading-rank; base=h2

\subsection{若望福音 8:42-44}

祂已将他们逐出与亚巴郎的亲属关系，当他们敢做更大的事时，祂便再加一击，告诉他们不但不是亚巴郎的子女，甚至还是魔鬼的\OriginalTerm{子女}{children of the devil}，并给予足以抵消他们无耻的打击；祂也没有让这话孤立无援，而是以事实证明。\Quote{因为，}祂说，\Quote{谋杀属于魔鬼的邪恶。}祂不仅说\Quote{你们做他的工作}，还说\Quote{你们行他的私欲}，显示双方都执着于谋杀，而嫉妒是原因。因为魔鬼毁灭了亚当，并非因有任何控告他的理由，只是出于嫉妒。祂在此也暗示这一点。\par

\Quote{并且不站在\OriginalTerm{真理}{truth}中。}即，在正道上。因为他们不断控告祂不是从天主来的，祂告诉他们这也出于那里。因为魔鬼是\OriginalTerm{谎言}{lie}之父，当他说\Quote{你们吃的那一天，你们的眼就会开了}\BibleRef{创 3:5}时，他首先使用了它。因为人使用谎言并非作为本有之物，而是作为违反本性的外来之物，但他却是作为本有之物。\par

% structure: p0032 source=h2 target=subsection reason=relative-heading-rank; base=h2

\subsection{若望福音 8:45}

这是什么后果？\Quote{你们对我没有指控，却想要杀我。因为你们是真理的仇敌，所以你们迫害我。若不是这个原因，你们就会提出你们的指控了。}因此祂加上说，\par

% structure: p0034 source=h2 target=subsection reason=relative-heading-rank; base=h2

\subsection{\BibleRef{若 8:46}}

然后他们说：\Quote{我们不是从淫乱生的。}事实上他们中有许多人是从淫乱生的，因为他们进行了不相宜的结合。但祂并没有定他们的罪，而是转而论述另一个论点。因为当祂以所有这些标记证明他们不是出于天主，而是出于魔鬼时（因为杀人出于魔鬼，说谎也出于魔鬼，而你们两者都做），然后祂指出爱是出于天主的标记。\Quote{你们为什么不明白我的话？}因为他们总是怀疑，说：\Quote{祂说的\InnerQuote{我去的地方你们不能去}是什么意思？}因此祂告诉他们：\Quote{你们不明白我的话，因为你们没有天主的话。这临到你们，是因为你们的理解力是卑下的，因为我的事对你们来说太过高深了。}但假如他们不能明白呢？这里\Quote{不能}的意思是不愿意；因为\Quote{你们训练自己成为卑劣的，不想象任何伟大的事物。}因为他们说他们迫害祂是出于对天主的热情，为此祂处处努力表明迫害祂是那些憎恨天主的行为，而反之，爱祂是那些认识天主的人的行为。\par

\Quote{我们有一位父，就是天主。}以此他们夸耀自己，夸耀他们的身份而非他们的义行。\Quote{因此你们的不信并不证明我是天主的仇敌，但你们的不信是你们不认识天主的标记。原因是你们愿意说谎并做魔鬼的作为。但这出自灵魂的卑下；（正如宗徒所说：\InnerQuote{因为在你们中间有嫉妒和纷争，你们岂不是属血气的吗？}\BibleRef{格前 3:3}）又为什么你们不能呢？因为你们愿意行你们父亲的私欲，你们渴望、你们雄心勃勃（去行它们）。}你们看见\Quote{你们不能}表明意志的缺乏吗？因为\Quote{亚巴郎没有这样做。}\Quote{他的作为是什么？温柔、良善、服从。但你们却站在相反的一边，变得刚硬残忍。}\par

但他们怎么会想到投靠天主呢？祂已证明他们不配亚巴郎；因此为了逃避这一指控，他们攀得更高。因为当祂指责他们杀人时，他们如此说，把这当作一种借口，仿佛他们在为天主报仇。因此祂表明这件事本身就是反对天主的行为。而这个\Quote{我出来了}表明祂是从那里来的。祂说\Quote{我出来了}，暗示祂来到我们中间。但既然他们可能会对祂说：\Quote{你说一些奇怪和新鲜的事}，祂告诉他们祂是从天主来的。\Quote{因此你们有充分的理由不听它们，因为你们是出于魔鬼。凭什么理由你们要杀我？你们对我有什么指控？如果没有，为什么你们不信我？}因此，既然已经通过他们的谎言和杀人证明他们出于魔鬼，祂也表明他们与亚巴郎和天主疏远，既因为他们恨一位没有行过任何不义的人，又因为他们不愿听祂的话；并且祂从各方面证明祂并非反对天主，他们不信并非因为这一点，而是因为他们是天主的陌生人。因为当一位没有犯过罪、说自己从天主来并为天主所派遣、说真理并如此说话以致挑战所有人去验证，之后却不被相信，显然他们不信是因为他们属血气的。因为罪确实会、是的确实会贬低灵魂。因此经上说：\Quote{你们成了迟钝于听的人。}\BibleRef{希 5:11}\par

4. 因此，我劝勉你们，务要竭力使我们的生命成为正义的，使我们的心灵得以洁净，以致没有任何污秽阻碍我们；为自己点燃知识的明灯，不要在荆棘中撒种。因为一个不知道贪婪是恶事的人，怎能知道更大的善事呢？一个不戒避这些世俗事物的人，怎能持守那属天的事物呢？夺取是好的，但不是夺取那会朽坏的东西，而是夺取天国。经上说：\Quote{猛力的人}夺取了它。\BibleRef{玛 11:12} 那么，我们不可能通过怠惰，而只能通过热忱来获得它。但\Quote{猛力的人}是什么意思？需要很大的力量（因为路是狭窄的），需要一颗青春而高贵的灵魂。抢夺者总想超过一切人，他们不顾定罪、控告或惩罚，只专注于一件事，即获取他们想要夺取的东西，他们超越一切在前面的人。那么，我们就夺取天国吧，因为在这里，夺取不是过错，反而是赞美，过错在于不夺取。在这里，我们的财富并非来自他人的损失。让我们赶快夺取吧。如果情欲扰乱我们，如果淫欲扰乱我们，我们就对本性施行暴力，让自己变得更温良，稍微劳动一下，以便永远安息。不要夺取金子，而要夺取那使金子显得如泥土般的财富。请告诉我，如果铅和金摆在面前，你会拿哪一个？显然你会拿金。那么，在夺取者受惩罚的地方，你选择更有价值的；而在夺取者受光荣的地方，你却放弃更有价值的吗？如果两种情况都有惩罚，你岂不宁愿追求后者？但在这种情况下，根本没有惩罚，反而有福乐。有人说：\Quote{怎样夺取呢？}丢掉你手中已有的东西；因为你只要抓住它们，就不能夺取别的。请想想，一个人两手抓满银子，他能趁保留银子时去夺取金子吗？除非他先丢掉银子，空出手来。因为夺取的人必须束紧腰带，以免被拖住。现在也有敌对的力量冲下来抢夺我们，但我们要逃避它们，逃避它们，不要拖带任何可让它们抓住的东西，让我们砍断绳索，脱去地上的事物。绸缎衣服有什么用？我们还要把这戏弄展演多久？我们还要埋藏金子多久？我本不想不停地讲这些事，但你们不容许我，不断给我提供机会和理由。可是现在至少让我们停止吧，这样，我们以生活去教导他人后，就能得到所应许的美善事物，因我们的主耶稣基督的恩宠与慈爱，祂和圣父及圣神同享光荣，从现在直到永远，永世无尽。阿们。\par

\SourceRecord{Translated by Charles Marriott.；Nicene and Post-Nicene Fathers, First Series；Vol. 14.；Edited by Philip Schaff.；Buffalo, NY: Christian Literature Publishing Co.,；1889.；<http://www.newadvent.org/fathers/240154.htm>.}

% structure: p0001 source=h1 target=section reason=page-title

\section{第55篇讲道：论\WorkTitle{若望福音}}

% structure: p0002 source=h2 target=subsection reason=relative-heading-rank; base=h2

\subsection{若 8:48-49}

1. 邪恶是一种无耻而放肆的东西，当它本该隐藏时，反而更加猛烈。正如犹太人的情形。当他们本该因所听闻的话而感到刺痛，钦佩这些话语的大胆和决定性时，他们反而侮辱祂，称祂为撒玛黎雅人，说祂附有魔鬼，并问道：\Quote{我们说你是个撒玛黎雅人，并附有魔鬼，岂不正对吗？}因为每当祂说出崇高之言时，在愚昧人看来便是疯狂。然而之前福音作者从未说过他们称祂为\Quote{撒玛黎雅人}；但从这句话看来，他们很可能常常这样断言。\par

有人说：\Quote{你附有魔鬼。}谁附有魔鬼呢？是那尊敬天主的，还是那侮辱尊敬祂的人的呢？那么极其温良柔和的基督说什么呢？\Quote{我没有魔鬼，而是尊敬那派遣我来的。}在需要教导他们，挫败他们过度的傲慢，教训他们不要因亚巴郎而骄傲的地方，祂是严厉的；但在需要祂自己被侮辱而忍受时，祂用了极大的温和。当他们说：\Quote{我们有天主和亚巴郎作我们的父亲}时，祂严厉地责备他们；但当他们称祂为附魔者时，祂谦卑地回应，以此教导我们：要报复对天主的侮辱，却要忽略对我们自身的侮辱。\par

% structure: p0005 source=h2 target=subsection reason=relative-heading-rank; base=h2

\subsection{若 8:50}

祂说：\Quote{这些事，}祂说，\Quote{我说过，是为表明你们身为凶手，不该称天主为你们的父亲；所以我因尊敬祂而说这些话，为祂的缘故我听到这些辱骂，为祂的缘故你们羞辱我。然而我不在乎这种无礼；你们要为自己所说的话向祂交账，因祂的缘故我现在听到这些事。因为\InnerQuote{我不寻求自己的光荣}。因此我放弃惩罚你们，转而劝诫你们，并劝你们如此行事，好使你们不仅免于惩罚，而且获得永生。}\par

% structure: p0007 source=h2 target=subsection reason=relative-heading-rank; base=h2

\subsection{若 8:51}

这里祂不仅论及信德，也论及纯洁的生活。上面祂说：\Quote{必得永生}，但这里说：\Quote{不见死亡}。\BibleRef{若 6:40}同时祂暗示他们，他们不能对祂做什么，因为那遵守祂话的人不会死，更何况祂自己呢？至少他们如此理解，就对祂说：\par

% structure: p0009 source=h2 target=subsection reason=relative-heading-rank; base=h2

\subsection{若 8:52}

即是说：\Quote{凡听了天主之言的都死了，听了你的话的人岂能不死吗？}\par

% structure: p0011 source=h2 target=subsection reason=relative-heading-rank; base=h2

\subsection{若 8:53}

他们的虚荣心啊！他们又转而寻求与祂的血缘关系。其实他们本应说：\Quote{你比天主还大吗？或是听你的话的人比亚巴郎还大吗？}但他们没有这样说，因为他们认为祂甚至比亚巴郎还小。因此，起初祂表明他们是凶手，从而使他们离开这种血缘关系；但当他们固执己见时，祂用另一种方式处理，表明他们徒劳无功。关于\Quote{死亡}，祂没有对他们说什么，也没有启示或告诉他们祂指的是哪种死亡，但与此同时祂要他们相信祂比亚巴郎大，甚至借此使他们羞愧。祂说：\Quote{当然，}祂说，\Quote{若我是一个普通人，没有做任何错事，我就不该死；但当我宣讲真理，没有罪，由天主派遣，又比亚巴郎更大时，你们岂不疯狂？你们企图杀我，岂不徒劳？}那么他们如何回答？\Quote{现在我们确知你附有魔鬼。}撒玛黎雅妇人却不是这样说的。她没有对祂说：\Quote{你附有魔鬼}；只说：\Quote{你比我们的父亲雅各伯还大吗？}\BibleRef{若 4:12}因为这些人傲慢可咒，而她却渴望学习；因此她怀疑并适当地谦虚回答，称祂为\Quote{主}。那应许更伟大之事且值得信赖的人，本不该被侮辱，反而应受钦佩；然而这些人却说祂附有魔鬼。撒玛黎雅妇人的那些话是出于怀疑；这些人的话却是不信和乖谬的。\Quote{你比我们的父亲亚巴郎还大吗？}所以这（祂所说的话）使祂比亚巴郎大。\Quote{因此，当你们看见祂被举起时，你们将承认祂更大。}为此祂说：\BibleQuote{当你们举起人子时，你们将知道我是那一位。}\BibleRef{若 8:28}请注意祂的智慧。祂先使他们脱离亚巴郎的血缘，然后表明祂比亚巴郎大，以致祂也被视为比众先知大得多。的确，因为他们不断地称祂为先知，所以祂说：\BibleQuote{我的话在你们中没有位置。}\BibleRef{若 8:37}在另一处祂宣称祂能使死人复活，但这里祂说：\Quote{相信的人永不见死亡，}这比不让信徒被死亡束缚大得多。因此犹太人更加愤怒。那么他们说什么呢？\par

\Quote{你把自己当作什么人？}\par

这也是以侮辱的方式说的。他们中一人说：\Quote{你自以为了不起。}于是基督回答：\par

% structure: p0015 source=h2 target=subsection reason=relative-heading-rank; base=h2

\subsection{若 8:54}

2. 异端者对这里怎么说？他们听见这问题：\Quote{你比我们的父亲亚巴郎还大吗？}却不敢对他们说：\Quote{是的，我更大，}而是以隐晦的方式这样做了。那么怎样呢？祂的光荣是\Quote{虚无}？就他们而言，是虚无。正如祂说：\BibleQuote{我的见证不真实}\BibleRef{若 5:31}，是就他们会形成的看法而言，祂在这里也是这样说。\par

\Quote{有一位尊敬我的。}\par

祂为何不说\Quote{那派遣我的父}，像祂以前那样，而说：\par

\Quote{就是你们称为你们天主的那一位。}\par

% structure: p0020 source=h2 target=subsection reason=relative-heading-rank; base=h2

\subsection{若 8:55}

因为祂要表明，他们不但不认识祂的父，也不认识天主。\par

\Quote{但我认识祂。}\par

\Quote{所以，说\InnerQuote{我认识祂}并非自夸，而说\InnerQuote{我不认识祂}则是谎言；但你们说你们认识祂，是撒谎；正如你们说你们认识祂时撒谎，同样，我若说我不认识祂，也是撒谎。}\par

\\Quote{若我光荣我自己。}既然他们说：\\Quote{你把自己当作谁？}他回答说：\\Quote{若我使自己（任何事物），我的光荣算不了什么。那么正如我确切认识天主，你们却不认识天主。}又如亚巴郎的情形，他没有完全否认他们的话，而是说：\\Quote{我知道你们是亚巴郎的子孙}，以使对他们的指控更重；同样在这里，他也没有完全否认，而是什么？\\Quote{你们所说的那位。}通过允许他们夸口，他增加了对他们的指控力度。那么你们又怎会\\Quote{不认识天主}呢？\\Quote{因为你们侮辱那位说和做一切好让他受光荣的，即使他是从他那里被派遣的。}这个断言没有见证支持，但接下来的内容将确立它。\par

\\Quote{并且我遵守他的话。}\par

在这里他们本可以反驳他，如果他们至少有什么可说，因为这是他被天主派遣的最有力证据。\par

% structure: p0027 source=h2 target=subsection reason=relative-heading-rank; base=h2

\subsection{若望福音 8:56}

祂再次表明，如果他们为\\Person{亚巴郎}所喜乐的事而悲伤，他们就与\\Person{亚巴郎}的种族无关。\\Quote{我的日子}，在我看来似乎是指被钉十字架的日子，\\Person{亚巴郎}曾通过献祭公羊和\\Person{依撒格}来预表。他们如何回答？\par

% structure: p0029 source=h2 target=subsection reason=relative-heading-rank; base=h2

\subsection{若望福音 8:57}

因此我们得出结论，基督将近四十岁。\par

% structure: p0031 source=h2 target=subsection reason=relative-heading-rank; base=h2

\subsection{若望福音 8:58-59}

你看见他如何证明自己比\\Person{亚巴郎}更大吗？因为那因看见他的日子而喜乐并热切渴望的人，显然是因为那日子是为益处的，且属于一位比自己更大者。因为他们曾说：\\Quote{木匠的儿子} \\BibleRef{玛 13:55}，并且对他没有更高的想法，他就逐步引导他们对他有一个崇高的认识。因此，当他们听到\\Quote{你们不认识天主}这句话时，他们并不悲伤；但当他们听到\\Quote{在\\Person{亚巴郎}存在以前，我就是}时，仿佛他们血统的高贵被贬低，他们就暴怒，并要拿石头砸他。\par

\\Quote{他看见我的日子，就欢欣。}他表明，他并非不情愿地走向他的受难，因为他赞美那因十字架而欢喜的人。因为这是世界的救恩。但他们却拿石头砸他；他们如此准备谋杀，而且这是他们自愿做的，没有经过查问。\par

但为什么他没有说\\Quote{在\\Person{亚巴郎}存在以前，我曾是}，而说\\Quote{我是}？因为圣父使用这一表达\\Quote{我是}，基督也同样使用；因为它表示不受时间限制的永恒存在。因此，这话在他们看来好像是亵渎的。现在，如果他们连与\\Person{亚巴郎}的比较都受不了，尽管这只是一件小事，那么如果他不断地使自己与圣父相等，他们难道会停止拿石头砸他吗？\par

此后，他又像一个人一样逃走，并隐藏自己，已经向他们提供了充分的教导；在完成他的工作后，他走出圣殿，离开去医治瞎子，通过他的行动证明他是在\\Person{亚巴郎}之前。但也许有人会说：\\Quote{为什么他不使他们瘫痪？这样他们就会相信。}他医治了瘫子，但他们不信；不仅如此，他行了无数奇迹；在受难时，他将他们击倒在地，并使他们的眼睛昏暗，但他们仍然不信；那么如果他使他们瘫痪，他们又怎么会相信呢？没有什么比一个陷于绝望而刚硬的灵魂更糟糕了；虽然它看见征兆和奇迹，它仍然坚持着同样的无耻。因此，\\Person{法郎}遭受了无数打击，只在受惩罚时才清醒，并且一直保持这种性格直到他生命的最后一天，追赶那些他曾释放的人。所以\\Person{保禄}不断说：\\Quote{免得你们中间有人被罪恶的欺骗所硬化。} \\BibleRef{希 3:13} 因为正如身体上的茧子一旦形成就变得死亡，没有感觉；同样，灵魂当它被许多情欲占据时，对德行就变得死亡；无论你向它施加什么，它都得不到感知，无论你警告惩罚还是其他，它都继续麻木。\par

3. 因此，我恳求你们，趁我们还抱有救恩的希望，趁我们还能回头，要竭尽全力这样做。因为那些变得麻木的人，此后便陷入绝望舵手的盲目状态，任凭船只随风漂流，自己不再加以任何协助。同样，嫉妒的人只关注一件事，就是满足自己的欲望，即使他可能会受罚甚至被杀，仍完全被那情欲所控制；同样，放纵和贪婪的人也是如此。但如果情欲的统治力如此之大，那么德行的统治力就更大了；如果为了情欲我们不惜一死，那么为了德行更当如此；如果他们（罪人）不顾自己的性命，那么我们在救恩的事业上更应如此。因为我们还有什么可说的呢？如果那些丧亡的人为他们的丧亡如此积极，而我们对自身的救恩却连同等的积极性都表现不出来，反而不断因嫉妒而消耗自己？没有什么比嫉妒更糟糕了；它毁掉别人，也毁掉自己。嫉妒的眼睛因忧伤而憔悴，他活在持续的死亡中，他把所有从未伤害过他的人视为敌人。他为天主受尊荣而悲伤，却为魔鬼所喜乐的事而喜乐。有人受人尊荣吗？那不是真正的尊荣，不要嫉妒他。但如果他受到天主的尊荣呢？你要努力效法他。你不愿意吗？那又何必连自己也毁掉呢？何必抛弃你所拥有的呢？你既不能像他一样，也不能获得任何益处？那你为何还要为自己取恶，而本该与他一同喜乐，好使即使你无法分享他的辛劳，也能因与他一同喜乐而获益？因为有时就连意愿也能成就极大的善。至少厄则克耳说，摩阿布人因对以色列子民幸灾乐祸而受罚，而另外一些人因哀悼邻人的不幸而得救。\BibleRef{则 25:8} 如果为别人的灾祸哀哭尚且有安慰，那么为别人的尊荣喜乐就更有福了。他指责摩阿布人因对以色列子民幸灾乐祸，但其实是天主惩罚了他们；但即使天主惩罚人，他也不愿我们因受罚者而喜乐。因为他并不愿惩罚他们。如果我们必须同情受罚者，那么就更不应嫉妒受尊荣者。例如，科辣黑和达堂与其同党一同灭亡，使他们所嫉妒的人更加光彩，却将自己交于惩罚。因为嫉妒是剧毒的野兽，是不洁的野兽，是无可宽恕的蓄意恶行，是无以辩解的邪恶，是一切邪恶的根源和母亲。因此，让我们将它连根拔除，好使我们在此世脱离邪恶，并得获未来的恩赐；赖我们的主耶稣基督的恩宠和慈爱，愿光荣归于祂，并偕同祂，及圣父与圣神，于今世，于永世，于无穷世。阿们。\par

\SourceRecord{Translated by Charles Marriott.；Nicene and Post-Nicene Fathers, First Series；Vol. 14.；Edited by Philip Schaff.；Buffalo, NY: Christian Literature Publishing Co.,；1889.；<http://www.newadvent.org/fathers/240155.htm>.}

% structure: p0001 source=h1 target=section reason=page-title

\section{\WorkTitle{若望福音第56篇讲道}}

% structure: p0002 source=h2 target=subsection reason=relative-heading-rank; base=h2

\subsection{若望福音 9:1-2}

\Quote{耶稣经过时，看见一个生来瞎眼的人。祂的门徒问祂说：\InnerQuote{老师，是谁犯了罪？是他，还是他的父母，使他生来瞎眼？}}\par

1. \Quote{耶稣经过时，看见一个生来瞎眼的人。} 祂满怀对人的慈爱，关心我们的救恩，并渴望堵住愚昧人的口，尽管无人留意，祂仍尽己之责。先知知道此事，便说：\BibleQuote{好叫祢在发言时显为公义，在受审判时显为清白。} \BibleRef{咏 50:4} 因此，当此处他们不接受祂崇高的言论，反而说祂有魔鬼，并企图杀害祂时，祂离开圣殿，医治了瞎眼的人，以祂的缺席缓和他们的怒气，并通过行奇迹软化他们的硬心与残酷，巩固自己的断言。祂所行的奇迹非同寻常，是当时首次发生的。被治好的人说：\BibleQuote{自创世以来，从未听说有人开了生来瞎眼之人的眼睛。} \BibleRef{若 9:32} 也许有人开过瞎子的眼，但从未有过生来瞎眼的。而且，祂离开圣殿后，特意去行这工，这很清楚：是祂看见了瞎眼的人，而非瞎眼的人来找祂；祂如此专注地看他，连门徒都察觉了。至少由此，他们来问祂：因为看见祂专注地看着那人，他们问祂说：\Quote{是谁犯了罪？是他，还是他的父母？} 这问题有误，因为他在出生前怎能犯罪？若他父母犯了罪，他又怎会受到惩罚？他们从何提出这样的问题？先前，当祂治好瘫痪者时，祂说：\BibleQuote{看，你已经痊愈了，不要再犯罪。} \BibleRef{若 5:14} 因此他们明白了那人因罪瘫痪，便说：\Quote{好吧，那人是因罪瘫痪；但关于这人，祢要说什么？他犯了罪吗？这不可能，因为他生来瞎眼。他的父母犯了罪吗？也不能这样说，因为子不代父受罚。} 因此，当我们看见一个孩子受虐待时，我们呼喊：\Quote{这能说什么？这孩子做了什么？} 这不是在问问题，而是困惑；门徒这样说，与其说是请教，不如说是困惑。那么基督怎么说呢？\par

% structure: p0005 source=h2 target=subsection reason=relative-heading-rank; base=h2

\subsection{若望福音 9:3}

祂这样说，并非要免除他们的罪，因为祂并非简单地说：\Quote{既不是这人犯了罪，也不是他的父母。} 而是加上：\Quote{他生来瞎眼，是为叫天主子在他身上受光荣。} \Quote{因为这人确实犯了罪，他的父母也犯了罪，但他的瞎眼并非由此而来。} 祂这样说，并非意指虽然这人确实并非如此，但其他人瞎眼是由于父母犯罪的缘故，因为一人犯罪，另一人不能被惩罚。若我们承认这点，也必须承认他在出生前就犯了罪。因此，当祂说：\Quote{也不是这人犯了罪} 时，祂并非说人可以从出生起就犯罪并因此受罚；同样，当祂说：\Quote{也不是他的父母} 时，祂并非说人可以因父母之故受罚。祂通过厄则克耳的口移除了这一假设：\BibleQuote{上主说：我指着我永生起誓，你们在以色列中必不再引用这俗语：\InnerQuote{父亲吃了酸葡萄，儿子的牙齿酸软。}} \BibleRef{则 18:3} 梅瑟也说：\BibleQuote{父亲不可为儿子死，儿子也不可为父亲死。} \BibleRef{申 24:16} 至于某位国王，圣经说他正因这理由没有做某事，遵守了梅瑟的法律。但若有人辩驳：\Quote{那么，为何经上说：\BibleQuote{父亲犯罪，惩罚他们的儿子，直到第三、第四代}？} \BibleRef{申 5:9} 我们应该回答：这话并非普遍适用，而是针对那些出离埃及的人。其含义是这样的：\Quote{这些人出离埃及后，在神迹奇事后，变得比未见这些事的祖先更坏，他们将受罚。} 经上说：\Quote{同那些人所受的相同，因为他们犯了同样的罪。} 若有人留意这经文，就更能确知这话是对那些人说的。那么，他为何生来瞎眼呢？\par

祂说：\Quote{为叫天主的荣耀得以彰显。}\par

这里又有一个困难：若没有此人的受苦，天主的荣耀就不可能彰显。当然，经上并未说不可能，因为那是可能的，而是说：\Quote{为叫他身上得以彰显。} 有人说：\Quote{什么？他为了天主的荣耀而受冤屈吗？} 告诉我，有什么冤屈？若天主根本不愿创造他，又如何？但我断言，他甚至从瞎眼中受益，因为他恢复了内心的眼目。犹太人凭借他们的眼目得了什么益处？他们受了更重的惩罚，即使看见仍是瞎眼的。这人又因瞎眼受了什么损害？他借此恢复了视力。因此，现世的祸并非真祸，福也并非真福；只有罪才是恶，而瞎眼并非恶。那将这人从无有带入存在者，也有能力让他保持原状。\par

2. 但有些人说，这个连接丝毫不表示原因，而是指向奇迹的结果；正如祂所说的：\BibleQuote{我为审判来到这世界，叫看不见的可以看见，看得见的反而成为瞎子}\BibleRef{若 9:39}；然而祂来并非为此，使看得见的成为瞎子。又如保禄说：\BibleQuote{因为认识天主为他们是很明显的事，原来天主已将自己显示给他们了，使他们无可推诿}\BibleRef{罗 1:19}；但天主向他们显示，并非为此使他们无可推诿，而是使他们可以得着推诿。又如另一处说：\BibleQuote{法律本是后加的，为叫过犯增多}\BibleRef{罗 5:20}；然而法律不是为此而赐下，而是为抑制罪恶。你岂不见到处连接都是指向结果吗？就如同一位卓越的建筑师建造房屋的一部分，其余部分则留下未完工，为向那不信的人借此残余证明他是整体的作者；同样，天主也像建造一座朽坏的房屋般连结并完成我们的身体，医治枯干的手，使瘫痪的肢体有力，使瘸子直行，洁净癞病人，使病人痊愈，使残疾人健全，从死亡中召回死者，睁开闭着的眼睛，或添上本来没有的眼睛；这一切，都是由于我们本性软弱而产生的缺陷，祂通过矫正而显示了自己的大能。\par

但当祂说：\Quote{为彰显天主的荣耀}时，祂是指自己，而非圣父；祂的荣耀已彰显出来。因为既然他们曾听说天主用尘土造人，所以基督也用泥土造眼睛。如果说：\Quote{我就是那个取尘土造人的那一位}，在听者看来似乎难以接受；但这由实际行动显明时，便不再成为障碍了。因此，祂用泥土和唾沫混合，显明了隐藏的荣耀；因为被认为创造万物的建筑师，这绝非小事。\par

此后，其余部分也相继而来；由部分证明了整体，因为对更大事物的信仰也证实了较小的事物。因为人比任何受造物更尊贵，而在我们的肢体中，眼睛最为尊贵。这就是为什么祂不是以普通方式，而是以祂所行的方式塑造了眼睛。虽然那肢体体积小，却比身体任何部分更为必要。保禄证实了这一点，他说：\BibleQuote{假若耳朵说：\InnerQuote{因为我不是眼，所以不属于身体。}它就不属于身体吗？}\BibleRef{格前 12:16}。因为我们内的一切都是天主智慧的彰显，但眼睛更是如此；它引导整个身体，美化整体，装饰面容，是众肢体的光。太阳在世界中是什么，眼睛在身体中也是什么；熄灭太阳，你就毁灭并搅乱一切；熄灭眼睛，脚、手、灵魂便毫无用处。当眼睛失效时，连知识也消失，因为借着它我们认识天主。\BibleQuote{其实，自从天主创世以来，祂那看不见的美善，即祂永远的大能和神性，都可凭祂所造的万物，辨认洞察出来}\BibleRef{罗 1:20}。因此，眼睛不仅为身体，而且超越身体为灵魂成为光。为此，它如同在堡垒中安设，获得更高的位置，并统辖其他感官。于是基督就塑造了眼睛。\par

为使你们不以为祂工作时需要物质，并使你们学到祂根本不需要泥土（因为那位从无中造出更大存在的，更可无需物质而造出此物），我说，为使你们学到祂这样做并非出于必要，而是为显示祂起初就是造物主，当祂抹上泥土后说：\Quote{去洗，好使你们知道我无需泥土来创造眼睛，但我的荣耀可借此彰显出来。}\par

% structure: p0013 source=h2 target=subsection reason=relative-heading-rank; base=h2

\subsection{\BibleRef{若 9:4}}

即是说：\Quote{我必须彰显自己，并做那显示我与圣父作相同的事}；不是\Quote{相似的}事，而是\Quote{相同的}，这说法标志着更大的不变性，且用于那些毫厘不差者。那么，此后谁还能面对祂，当他看见祂与圣父有同样的大能时？因为祂不仅形成或开了眼睛，还赐予视力，这证明祂也吹入灵魂。因为如果灵魂不运作，即使眼睛完美，也绝不能看见任何东西；因此祂赐予来自灵魂的能力，也赐予肢体——具备一切，包括动脉、神经、静脉，以及构成身体的一切。\par

\BibleQuote{我必须趁着白日做工。}\BibleRef{若 9:4}\par

\Quote{趁着白日，趁着人仍能信从我，趁着今生尚存，我必须工作。}\par

\BibleQuote{黑夜来到}\BibleRef{若 9:4}——即未来——\BibleQuote{那时没有人能工作}\BibleRef{若 9:4}\par

他说：\Quote{我不能工作}并不是说：\Quote{没有人能工作}：也就是说，不再有信德、劳苦或悔改。因为为表明祂称信德为\Quote{工作}，当他们对祂说：\Quote{我们该做什么，才能做天主的工作呢？}\BibleRef{若 6:28}，祂回答：\Quote{这是天主的工作，就是信祂所派遣的那一位。} 那么，将来的人怎能不做这工作呢？因为那里没有信德，所有人都自愿或不自愿地屈服。免得有人说祂这样做是出于荣誉的渴望，祂表明祂所做的一切都是为了怜悯那些只能\Quote{在这里}相信、却不能再\Quote{在那里}获得任何好处的人。为此，尽管盲人没有来到祂面前，祂仍然做了祂所做的：因为那个人是值得治愈的，如果他看见了，他会相信并来到基督面前；如果他听人说祂在场，他也不会疏忽，这从后来他的勇气和信德中可以看出。因为他很可能会在心里想，并说：\Quote{这是什么？祂和了泥，抹在我的眼睛上，对我说：\InnerQuote{去，洗吧；}难道祂不能先治好我，然后再打发我到\Place{西罗亚}去吗？我常和其他许多人一起在那里洗，却没有得到任何好处；如果祂有能力，祂会当着我的面治好我。} 正如\Person{纳阿曼}论及\Person{厄里叟}时所说的；因为他也被命令去\Place{约旦河}洗，却不相信，尽管当时\Person{厄里叟}的名声远扬。\BibleRef{列下 5:11} 但盲人既没有不信，也没有反驳，也没有在心里想：\Quote{这是什么？祂应该抹泥吗？这反而是让人更加盲目：谁曾这样恢复视力呢？} 但他没有这样的推理。你看到他坚定的信德和热忱吗？\par

\Quote{黑夜来临。}接着祂表明，即使在受难之后，祂仍会关心不敬虔的人，并引领许多人归向自己。因为\Quote{还是白天。}但之后，祂将他们完全剪除，并宣告这一点，祂说：\par

% structure: p0020 source=h2 target=subsection reason=relative-heading-rank; base=h2

\subsection{\BibleBook{若望}{福音} 9:5}

3. 正如祂也对其他人说：\Quote{趁着有光，相信吧。}\BibleRef{若 12:36} 那么，为什么保禄称此生命为\Quote{黑夜}，而称那生命为\Quote{白日}呢？并不是反对基督，而是说同样的事，即使不是言语上，也是意义上；因为他也说：\Quote{黑夜已深，白日将近。}\BibleRef{罗 13:12} 他称现今为\Quote{黑夜}，是因为那些坐在黑暗中的人，或者因为他将其与将来的白日相比较；基督称将来为\Quote{黑夜}，因为那里罪没有能力运作；但保禄称现世生命为黑夜，因为那些继续生活在邪恶和不信中的人处于黑暗。于是他向信友说：\Quote{黑夜已深，白日将近，}因为他们应享受那光；他称旧的生命为黑夜。他说：\Quote{让我们脱去黑暗的行为。}你看到他对他们说这是\Quote{黑夜}吗？因此他说：\Quote{让我们行事为人正如在白日，}好使我们享受那光。因为如果此光如此美好，想想那光将是怎样的；正如阳光比烛光更明亮，同样，那光远胜于此光。基督这样说，表示\Quote{太阳将变暗。}因为那光亮的极盛，连太阳都看不见。\par

如果现在为了在这里拥有光线充足、空气流通的房屋，我们花费巨资，建造和劳碌，那么想想我们应该如何付出我们的身体本身，好使我们在天上——那里有不可言喻的光——能建起荣耀的房屋。在这里，关于地界和墙壁有争吵和纷争，但那里将不会有任何这类事情，没有嫉妒，没有恶意，没有人会为划定地界与我们争辩。这住所我们也必定要离开，但那住所却永远与我们同在；这住所会因时间而朽坏，遭受无数的损害，但那住所却会永不衰老地存留；这住所穷人无法建造，但那另一个住所却可以用两个小钱建造，如同那位寡妇所做的。因此，我悲痛欲绝，因为那么多恩惠摆在我们面前，我们却懒惰，藐视它们；我们竭尽全力要在这里拥有辉煌的房屋，但对于如何在天上获得哪怕一小块安息之地，我们却毫不关心，不去思想。请告诉我，你愿意你的住所在这里哪里？在旷野里，还是在一个较小的城市中？我想不是；而是在某个最皇家、最宏伟的城市里，那里商业更繁荣，辉煌更盛大。但我要带领你进入这样一座城，它的建造者和创造者是\Person{天主}；在那里我劝你奠基和建造，花费更少（劳力更少）。那座房屋是由穷人的双手建造的，它确实是\Quote{建筑}，就像这里建造的建筑物是极度愚蠢的产物一样。因为如果有人带你到波斯境内，去看那里的景象然后返回，然后命令你在那里建造房屋，难道你不会认为他极其愚蠢，竟命令你在不合时宜的时候花费吗？那么，当你这样对待这片你不久就要离开的大地时，你岂不是在做同样的事吗？\Quote{但我要把它留给我的孩子，}有人说。然而他们也要在你不久之后离开它；甚至常常在你之前；他们的继承者也一样。即使在这里，你看到你的继承者没有保留他们的财产，这对你来说是件愁苦的事；但那里你无需担心任何这类事情；财产是稳固的，属于你、你的孩子和他们的后代，只要他们效仿同样的良善。那建造是由\Person{基督}着手进行，建造那房屋的人无需指定管家，无需操心，无需焦虑；因为当\Person{天主}承担了这工程，还需要操心什么？祂把一切聚集起来，建起房屋。这并非唯一奇妙的事，而且祂建造的方式也令你喜悦，甚至超过你所喜悦的，超出你的愿望，因为祂是最卓越的艺术家，极其关心你的益处。如果你贫穷，渴望建造这房屋，它不会给你带来嫉妒，不会引发对你的恶意，因为那些知道如何嫉妒的人看不见它，只有知道为你的福乐而欢喜的天使看见它；没有人能侵犯它，因为附近没有一个人患有这类情欲。你在那里有圣人作邻居：\Person{伯多禄}和\Person{保禄}以及他们的同伴，所有先知，殉道者，众多天使和总领天使。因此，为了这一切，让我们把我们的财物施舍给穷人，好使我们获得那些居所；愿我们都能通过我们的\Person{主耶稣基督}的恩宠和慈爱获得这些居所，愿光荣藉着祂并同祂一起归于\Person{圣父}和\Person{圣神}，从今世直到永远。阿们。\par

\SourceRecord{Translated by Charles Marriott.；Nicene and Post-Nicene Fathers, First Series；Vol. 14.；Edited by Philip Schaff.；Buffalo, NY: Christian Literature Publishing Co.,；1889.；<http://www.newadvent.org/fathers/240156.htm>.}

% structure: p0001 source=h1 target=section reason=page-title

\section{《若望福音》讲道集 第五十七篇}

% structure: p0002 source=h2 target=subsection reason=relative-heading-rank; base=h2

\subsection{若望福音 9:6-7}

1. 凡想从阅读中获得益处的人，连经文中的极小部分也不可放过；正因如此，我们奉劝要\Quote{搜查}圣经，因为大部分经文虽然初看容易，但深处却隐藏着许多意义。请看目前的情形。\Quote{祂说了这些话}，经上说，\Quote{就吐唾沫在地上。}什么话？\Quote{为叫天主的光荣彰显出来}，以及\Quote{我必须做那派遣我者的工作}。因为福音作者并非无缘无故向我们提及祂的话，又加上\Quote{祂吐唾沫}，而是要表明祂以行动证实了祂的话。祂为什么不用水而用唾沫和泥呢？祂正要打发那人往西罗亚去；因此，为使人不将任何功劳归于那泉水，而叫你明白从祂口中发出的德能，既塑造了那人的眼睛，又开启了它们，祂\Quote{吐唾沫在地上}；至少福音作者在说\Quote{用唾沫和泥}时指明了这一点。然后，为使成功的果效不显得是出于泥土，祂吩咐他去洗。但祂为什么不立刻完成这事，却打发他往西罗亚去呢？为叫你学习这盲人的信德，也使犹太人的固执被堵住：因为很可能他们都会看见他离开，眼睛上还涂着泥；由于这事的奇异，他会吸引所有人，无论认识或不认识他的人，都会仔细观察他。又因为认出复明的盲人并不容易，祂先藉着路途的长度使许多人成为见证人，并藉着景象的奇异使他们成为精确的观察者，好叫他们更加留意，再也不能说：\Quote{是他，不是他。}此外，打发他到西罗亚去，祂愿意证明祂并不疏远法律和旧（盟约），此后也无需担心西罗亚会得荣耀，因为许多常在那里洗眼的人并未得到这样的益处；因为在那里也是基督的德能成就了一切。为此，福音作者为我们补充了名字的解释；说了\Quote{在西罗亚}之后，他加上,\par

\Quote{即是\InnerQuote{被派遣者}。}\par

为叫你明白，那里也是基督医治了他。正如保禄说：\BibleQuote{他们饮了那随从他们的属灵盘石，那盘石就是基督。}\BibleRef{格前 10:4} 正如基督是属灵的盘石，祂也是属灵的西罗亚。在我看来，水的突然涌入也暗示了一项不可言喻的奥秘。那是什么？就是祂显现的出人意料的本性，超越一切期望。\par

但请看这盲人的心志，在一切事上顺服。他并没有说：\Quote{如果真是泥或唾沫给了我眼睛，那还需要西罗亚做什么？或者如果需要西罗亚，那还需要泥做什么？祂为什么给我涂抹？又为什么吩咐我洗？}他没有存这样的念头，他只预备了一件事：在一切事上服从发令者，所行的一切都不使他反感。若有人问：\Quote{那么，当他除去泥时，他是怎样复明的呢？}他从我们这里听不到别的答案，只能说我们不知道那方式。我们不知道有什么稀奇，因为连福音作者也不知道，连那被治好的人自己也不知道？他所知道的是所发生的事，但方式他不能理解。所以当人问他时，他说：\Quote{祂把泥放在我的眼睛上，我洗了，就看见了}；但这事是怎样发生的，即使他们问千万次，他也说不出来。\par

% structure: p0007 source=h2 target=subsection reason=relative-heading-rank; base=h2

\subsection{若望福音 9:8-9}

所成就之事的奇异甚至使他们不信，尽管已经安排了那么多，使他们不致不信。他们说：\Quote{这不是那坐着讨饭的人吗？}哦，天主的慈爱！祂降卑到何等程度，以极大的慈爱甚至医治乞丐，从而堵住犹太人的口，因为祂不以显赫、杰出、掌权的人，而以无名之辈作为同一眷顾的对象。因为祂来是为众人的救恩。\par

在那瘫子身上发生的事，也发生在这人身上：二者都不知道谁治好他们。这是由于基督的退隐；因为耶稣在治好之后经常退隐，为叫一切怀疑从奇迹中除去。因为那些不认识祂的人怎能奉承祂，或参与策划所行的事呢？这人也不是那些四处走动的人，而是那些坐在圣殿门口的人。当众人都对他疑惑时，他说什么？\par

\Quote{我就是。}\par

他不以先前的瞎眼为耻，也不惧怕百姓的怒气，他不拒绝显露自己，为能宣告他的恩主。\par

% structure: p0012 source=h2 target=subsection reason=relative-heading-rank; base=h2

\subsection{若望福音 9:10-11}

你说什么？\Quote{一个人}行这样的事？那时他对祂还没有什么崇高的认识。\par

\Quote{那名叫耶稣的人，和了泥，抹在我的眼睛上。}\par

2. 请看他是多么诚实。他没有说祂用什么和的泥，因为他没有提及他所不知道的事；他并未看见祂吐唾沫在地上，但他从感觉和触摸知道祂涂抹了泥。\par

\Quote{并对我说：你往西罗亚池里去洗。}\par

这事也是他的耳朵为他作证的。但他如何认出祂的声音？从祂与门徒的谈话。他述说了这一切，并藉着工作得了见证，但医治的方式他不能说出。如果连感觉和触摸的事都需要信德，那么在看不见的事上就更需要了。\par

% structure: p0018 source=h2 target=subsection reason=relative-heading-rank; base=h2

\subsection{若望福音 9:12}

他们说：\Quote{祂在哪里？}已经心中怀着杀害祂的意图。但请注意基督的谦逊，祂没有继续与那些被治好的人在一起；因为祂既不愿收获光荣，也不愿吸引群众，也不愿炫耀自己。也请注意那盲人是多么诚实地回答一切问题。犹太人想要找到基督，把祂带到司祭们那里，但没有找到祂，就把那盲人带到法利塞人那里，如同带到那些会严厉审问他的人面前。因此，福音作者指出，那一天是\Quote{安息日}\BibleRef{若 9:14}，为了指出他们的恶念，以及他们寻找祂的理由，仿佛他们找到了把柄，能借着似乎是违反法律的行为来贬低这奇迹。这从他们一见到他就只说\Quote{祂怎样开了你的眼睛？}就很清楚。也请注意他们说话的方式；他们不说\Quote{你怎样看见了？}而是说\Quote{祂怎样开了你的眼睛？}这样给了他一个诽谤耶稣的借口，因为祂做了工。但他简短地回答他们，如同对待已经听见的人；因为他没有提祂的名字，也没有说\Quote{祂对我说：\InnerQuote{去洗洗}}，就立刻说：\par

% structure: p0020 source=h2 target=subsection reason=relative-heading-rank; base=h2

\subsection{若望福音 9:15}

因为诽谤如今变得严重了，犹太人曾说：\Quote{看，耶稣在安息日做的是怎样的工作，祂用泥抹！}但请观察，我恳求你，那盲人是多么不为所动。当受审问时，他在那些没有危险的人面前说话，说实话并不是什么大不了的事，但令人惊奇的是，如今当他处于更大的恐惧之中时，他既没有否认也没有反驳他先前所说的话。那么法利塞人做了什么呢？或者，其他人也做了什么？他们把他带到法利塞人那里，以为他会否认；但相反，他们所不愿的事发生了，他们得到了更确切的消息。而他们在神迹的事上处处都得忍受这一点，这一点我们将在后面更清楚地说明。法利塞人说了什么？\par

% structure: p0022 source=h2 target=subsection reason=relative-heading-rank; base=h2

\subsection{若望福音 9:16}

你们难道没有看见他们被神迹引动了吗？请听他们现在说的话，这些人先前曾派人去抓祂。即使并非所有人都如此（因为身为统治者，因虚荣而陷入不信），但大部分统治者仍然相信了祂，只是没有承认祂。至于群众，很容易被忽视，因为他们在会堂中无足轻重，但统治者更为显眼，因此更难大胆直言；有些人被统治欲所束缚，另一些人被怯懦和对许多人的恐惧所束缚。因此祂也说：\Quote{你们彼此互相求光荣，而不求那独一天主的光荣，怎能信我呢？}\BibleRef{若 5:44} 而这些不义地想要杀害祂的人说，他们是属于天主的，但那治好盲人的不可能是属于天主的，因为祂不守安息日；另一些人则反对说，一个罪人不能行这样的神迹。那些心怀恶意地保持沉默的人，只提那表面的过犯；因为他们不说：\Quote{祂在安息日治病}，而是说：\Quote{祂不守安息日}。而另一些人则软弱地回答，因为他们本应证明安息日没有被违背，却只依赖神迹；这也情有可原，因为他们仍然认为祂是一个人。如果不是这样，他们本可以辩解说，祂是安息日的主，是祂自己所设立的，但那时他们还没有这种看法。总之，没有人敢公开说出他所想说的话，或以一种断言的方式，而只是以一种怀疑的方式；有些人是因为没有说话的胆量，另一些人则是因为统治欲。\par

\Quote{因此他们中间起了纷争。}这纷争首先在民众中开始，然后也在统治者中开始，有些人说：\Quote{祂是一个好人}；另一些人说：\Quote{不，祂迷惑了群众。}\BibleRef{若 7:12} 你们难道没有看见统治者比群众更缺乏理解力，因为他们分争得比群众晚？而分争之后，当他们看到法利塞人逼迫他们时，他们并没有表现出任何高尚的情感。因为如果他们完全与他们分离，他们很快就会知道真理。因为在分离中，是可以做得很好的。因此祂自己也曾说：\Quote{我来不是为带和平，而是带刀剑。}\BibleRef{玛 10:34} 因为有一种邪恶的和谐，也有一种善意的分歧。那些建造巴贝耳塔的人\BibleRef{创 11:4}同心合意，却害了自己；而这些人又被迫分离，虽然不情愿，却对他们有益。同样，科辣黑和他的同伴同心作恶，因此他们被分离是为好；犹达斯与犹太人同心作恶。所以分争可能是好的，同心可能是恶的。因此经上说：\Quote{若是你的右眼使你跌倒，就把它挖出来；若是你的脚使你跌倒，就把它砍下来。}\BibleRef{玛 5:29} 如果我们必须与一个接合不好的肢体分离，难道不是更应该与那些因恶而与我们联合的朋友分离吗？所以同心并非总是好事，正如分争并非总是坏事。\par

3. 我说这些事，是为叫我们躲避恶人，而追随善人；因为论到我们的肢体，若有一部分腐烂而无法医治，我们便将它割去，恐怕其余的身体染上同样的疾病，我们这样做并非不关心那部分，而是为保全其余部分，那么，对那些与我们同流合污的人，我们岂不更当如此？如果我们能救治他们而不致自身受损，就应当用尽一切方法去做；但他们若执迷不悟，且可能伤害我们，就必须将他们割除并弃绝。因为这样他们往往反而获益。因此保禄也劝勉说：\BibleQuote{你们应将那恶人从你们中间除去}；并且，\BibleQuote{叫行这事的人从你们中间赶出去。}\BibleRef{格前 5:13}与恶人交往，是一件可怕的事，实在可怕！瘟疫或疥癣传染给接触患者的人，远不如恶人的邪恶传得快。因为\BibleQuote{邪恶的交往败坏了良好的品行。}\BibleRef{格前 15:33}先知又说：\BibleQuote{你们要从他们中间出来，并要分开。}\BibleRef{依 52:11}因此，谁也别让恶人作自己的朋友。倘若我们有不良的儿子，我们公开不认他们，不顾天性及其规律，不顾它所加给我们的约束，那么，对于恶毒的同伴和相识，岂不更应当远离？因为即使我们未受他们的伤害，总难免遭受恶名，因为外人不深究我们的生活，却从我们的同伴来判断我们。这番劝告我特别对青年男女说。经上说：\BibleQuote{要预备光明的事，不但在主面前，也在众人面前。}\BibleRef{罗 12:17}我们当用尽一切方法，免得使近人跌倒。因为一种生活，即使非常正直，若使他人跌倒，就失去了一切。但正直的生活如何能使人跌倒呢？当与不正派的人交往给它带来恶名时，就是这样；因为当我们自信而交结恶人时，即使自己未受害，却使他人跌倒。我对男人、女人和少女说这些话，让他们凭良心仔细看出从这事产生出多少罪恶。或许我，或许任何更成全的人，都不怀疑什么恶事；但较单纯的弟兄却因你们的成全而受害；你们也当为他的软弱操心。即使他未受害，希腊人却受害了。保禄命我们\BibleQuote{不要成为绊脚石，无论是犹太人、希腊人，或是天主的教会。}\BibleRef{格前 10:32}（我对童女不怀疑什么恶事，因为我爱童贞，而\BibleQuote{爱不作任何恶事}\BibleRef{格前 13:5}；我非常钦佩那种生活状态，我甚至不能对它有任何不当的思想。）我们如何说服那些外人？因为我们也必须为他们预作考虑。让我们这样安排有关我们自己的事，使任何不信者都不能找到正当的控诉理由来攻击我们。因为正如那些表现正直生活的人光荣天主，那些行事相反的人却使祂受亵渎。愿我们当中没有这样的人；但愿我们的善行如此照耀，使我们在天之父得受光荣，而我们也能享受由祂而来的尊荣。愿我们众人藉着我们的主耶稣基督的恩宠和慈爱，都得达到这目标。祂和圣父及圣神，永受光荣，至于无穷之世。亚孟。\par

\SourceRecord{Translated by Charles Marriott.；Nicene and Post-Nicene Fathers, First Series；Vol. 14.；Edited by Philip Schaff.；Buffalo, NY: Christian Literature Publishing Co.,；1889.；<http://www.newadvent.org/fathers/240157.htm>.}

% structure: p0001 source=h1 target=section reason=page-title

\section{\WorkTitle{若望福音 第五十八篇讲道}}

% structure: p0002 source=h2 target=subsection reason=relative-heading-rank; base=h2

\subsection{若望福音 9:17-18}

我们必须以极大的精确性研读圣经，不得漫不经心或草率行事，以免陷入困境。因为即使在此处，也有人可能合理地提出问题：当他们断言\Quote{这个人不是出于天主，因为祂不守安息日}之后，如今他们对那人说\Quote{你说祂怎样，因为祂开了你的眼睛？}而不是说\Quote{你说祂怎样，因为祂犯了安息日？}但如今他们提出了辩护的根据，而不是控告的根据。那么，我们该如何回答呢？这些人并不是那些说\Quote{这个人不是出于天主}的人，而是那些从他们中分离出来的人，他们也说\Quote{一个罪人不能行这样的奇迹}。因为他们为了更有效地堵住对手的口，以免显得他们偏袒基督，便把那曾亲身领受祂权能证明的人带上来，质问他。现在请看这个穷人的智慧，他比他们所有人都更有智慧地说话。首先他说：\Quote{祂是一位先知}；并且没有因那些悖逆的犹太人反对祂所说的\Quote{这个人怎能是出于天主的，因为祂不守安息日？}而退缩，反而回答他们说：\Quote{祂是一位先知。}\par

\Quote{他们不相信他曾是盲人而得到了视力，直到他们叫来了他的父母。}\par

请观察他们试图掩盖和抹去这个奇迹的多种方式。但这是真理的本性：正是通过那些似乎被人攻击的手段，它变得更加强大，它通过被掩盖的手段而发光。因为如果这些事没有发生，许多人可能会怀疑这个奇迹；但现在，仿佛他们想要揭露真相，他们用尽一切手段，而且如果他们做了所有为基督辩护的事，他们就不会采取别的行动。因为首先他们试图借此方式（医治）摧毁祂，说：\Quote{祂怎样开了你的眼睛？}即\Quote{是否使用了某种巫术？}在另一处，当他们无法对祂提出指控时，他们也试图侮辱医治的方式，说：\Quote{祂驱魔，除非藉贝耳则步。}\BibleRef{玛 12:24}在这里，当他们无话可说时，又诉诸时间（医治），说：\Quote{祂犯了安息日}；又说：\Quote{祂是个罪人。}然而，祂问你们这些想要杀祂、随时准备抓住祂行为的人，最直白地说：\Quote{你们中间谁能指证我有罪？}\BibleRef{若 8:46}却没有人说话，也没有人说：\Quote{你亵渎，因为你使自己无罪。}如果他们有能力这样说，他们就不会保持沉默。因为那些因听见祂说在亚巴郎之前就有了而要砸死祂的人，说祂不是出于天主，他们这些谋杀者却自夸出于天主，但说这位行了如此奇迹、在施行医治后不是出于天主的人，因为祂不守安息日，如果他们哪怕有丝毫指控的阴影，也绝不会放过它。即使他们称祂为罪人，因为祂似乎犯了安息日，这个指控也被证明是站不住脚的，因为那些与他们同流合污的人谴责了他们的极度冷漠和心灵狭隘。因此，处处受困后，他们后来转而采取更无耻和厚颜的手段。是什么呢？经上说，他们\Quote{不相信他曾是盲人而得到了视力}。那么，他们怎么又指控基督不守安息日呢？显然，他们原本是相信的。但为什么你们不注意那许多人？不理会那些认识他的邻居？如我说过的，虚假处处自取败亡，正是通过那些似乎困扰真理的手段，使真理显得更加光明。现在的情况正是如此。为了使没有人说他的邻居和见过他的人没有精确地说出来，而是凭着相似猜测，他们便带来他的父母，通过他们，尽管不情愿，却成功证明了所发生的事是真实的，因为父母最了解自己的孩子。当他们无法恐吓那人本人，却见他无所畏惧地宣扬他的恩人时，他们想通过他的父母来损害这个奇迹。请看他们问话的恶意。因为经上怎么说？把他们放在中间，使他们陷入困境，便用极大的严厉和愤怒进行质问，\par

% structure: p0006 source=h2 target=subsection reason=relative-heading-rank; base=h2

\subsection{若望福音 9:19}

仿佛他们是在欺骗，为基督策划。哦，你们这些可咒骂的，彻底可咒骂的！哪个父亲会选择编造这样的谎言来反对自己的孩子？因为他们几乎就是说：\Quote{你们声称他是瞎子，不仅如此，还把这事到处传扬。}\par

\Quote{那么，他如今怎么看见了？}\par

2. 哦，愚蠢！有人说：\Quote{你们的诡计和机巧。}因为通过这两件事，他们试图引导父母否认：使用\Quote{你们所说的那人}和\Quote{那么他如今怎么看见了？}这些话。现在，当被问及三个问题：他是否是他们的儿子，他是否曾是盲人，以及他如何获得视力，父母只承认了其中两个，却没有补充第三个。这是为了真理的缘故而发生，为的是除了那个被治好的、也是可信的人之外，没有别人承认这件事。父母又如何会偏袒（基督）呢？他们连所知道的有一部分也因害怕犹太人而不敢说。他们说什么？\par

% structure: p0010 source=h2 target=subsection reason=relative-heading-rank; base=h2

\subsection{若望福音 9:20-21}

通过使他有信用，他们为自己辩解；他们说：\Quote{他不是孩子，也不是无能，而是能够为自己作证。}\par

% structure: p0012 source=h2 target=subsection reason=relative-heading-rank; base=h2

\subsection{若望福音 9:22}

注意圣史如何再次引述他们的见解和想法。我这样说，是因为他们先前所说的那句话，当他们说：\Quote{祂使自己与天主平等}。\BibleRef{若 5:18} 因为如果那也是犹太人的看法，而非基督的判断，他就会加上一句说：\Quote{那是一个犹太人的看法}。因此，当父母把他们指到那被治好的人那里时，他们又第二次叫他来，并没有公然无耻地说：\Quote{否认基督治好了你}，却想假借虔敬的外表来达到这个目的。\par

% structure: p0014 source=h2 target=subsection reason=relative-heading-rank; base=h2

\subsection{\BibleRef{若 9:24}}

因为如果对父母说：\Quote{否认他是你的儿子，他是生来瞎眼的}，会显得非常可笑。再者，对他本人这样说，显然是厚颜无耻。因此他们并不这样说，而是用另一种方式处理此事，说：\Quote{归光荣于天主}，就是说，\Quote{承认这人什么事也没有做}。\par

\Quote{我们知道这人是罪人。}\par

\Quote{那么，为什么当祂说：\InnerQuote{你们中间谁能指证我有罪？}\BibleRef{若 8:46}时，你们没有定祂的罪？你们从哪里知道祂是罪人？} 在他们说了\Quote{归光荣于天主}之后，那人没有回答，基督遇见他并称赞他，并没有责备他说：\Quote{你为什么没有归光荣于天主？}但祂说什么？\Quote{你信天主子吗？}\BibleRef{若 9:35}，好叫你们知道这就是\Quote{归光荣于天主}。如果祂在尊荣上与父不平等，这就不是光荣；但既然尊敬子的也尊敬父，那盲者没有被责备是理所当然的。当法利塞人指望父母会反驳并否认奇迹时，他们对那人本人什么也没说，但看到这样毫无结果，他们又回到他那里，说：\Quote{这人是罪人。}\par

% structure: p0018 source=h2 target=subsection reason=relative-heading-rank; base=h2

\subsection{\BibleRef{若 9:25}}

那位盲者难道没有害怕吗？断乎没有。那么，他先前说\Quote{祂是先知}\BibleRef{若 9:17}，现在却说\Quote{祂是不是罪人，我不知道}？他这样说，并不是出于这样的心态，也不是自己相信了这件事，而是希望借着事实的见证来辩明祂，免受他们的指控，并非凭自己的声明，而是使辩护可信，让善行本身的见证来定他们的罪。因为如果在那盲者说了许多话之后，说\Quote{这人若不是正义的，就不能做这样的奇迹}\BibleRef{若 9:33}，他们如此恼怒以至于回答说：\Quote{你完全生在罪中，竟敢教训我们吗？}，那么，如果他一开始就这样说，他们会说什么呢？他们会做出什么事来呢？\Quote{祂是不是罪人，我不知道}；好像他在说：\Quote{我不为这人说什么，目前我不作任何声明，但我确实知道并要肯定，如果祂是罪人，祂就不能做这样的事。}这样，他使自己免于嫌疑，他的见证也无懈可击，因为他不是出于偏袒，而是根据事实作证。因此，当他们既不能推翻也不能移除所发生的事时，他们又回到原来的计谋，琐碎地询问治愈的方式，就像那些四处搜寻眼前的猎物，却无法伤害它的人，急促地四处奔走；他们重提那人先前的话，想通过不断的提问来使之失效，并说：\par

% structure: p0020 source=h2 target=subsection reason=relative-heading-rank; base=h2

\subsection{\BibleRef{若 9:26}}

他的回答是什么？他既已战胜并打倒他们，就不再谦卑地对他们说话。只要事情还需要查问和论证，他就谨慎地说，同时提供证据；但当他得胜并获得辉煌的胜利时，他就鼓起勇气，践踏他们。他说了什么？\par

% structure: p0022 source=h2 target=subsection reason=relative-heading-rank; base=h2

\subsection{\BibleRef{若 9:27}}

你看见一个乞丐对经师和法利塞人的胆量吗？真理如此坚强，谎言如此软弱。真理，即使抓住平凡人，也能使他们显得荣耀；谎言，即使与强者同在，也显出他们的软弱。他所说的是这样：\Quote{你们不留意我的话，所以我将不再说话或不断回答你们，你们无缘无故地问我，不是想听以学习，而是要侮辱我的话。}\par

\Quote{你们也愿意做祂的门徒吗？}\par

现在他已将自己列于门徒的行列，因为\Quote{你们也愿意？}这句话是一个宣称自己是门徒的人说的。然后他大大地嘲笑和激怒他们。因为他知道这深深打击了他们，所以他这样说，想极其严厉地责备他们；这是一个勇敢的灵魂的行为，它高飞在天上，藐视他们的疯狂，指出这尊严的伟大，他对此非常自信，并表明他们侮辱了他这个值得钦佩的人，但他并不把这侮辱当作自己的，反而将他们认为的耻辱视为荣耀。\par

% structure: p0026 source=h2 target=subsection reason=relative-heading-rank; base=h2

\subsection{\BibleRef{若 9:28}}

\Quote{但这不可能。你们既不是摩西的，也不是这人的；因为如果你们是摩西的，你们也会成为这人的。}因此基督先前对他们说，因为他们不断使用这些话：\Quote{如果你们信摩西，你们也必信我，因为他为我作了记载。}\BibleRef{若 5:46}\par

% structure: p0028 source=h2 target=subsection reason=relative-heading-rank; base=h2

\subsection{\BibleRef{若 9:29}}

凭谁的话，谁的见证？有人说：\Quote{我们祖先的。}难道他不比你们的祖先更可信吗？祂以奇迹证实祂来自天主，祂所说的来自天上。他们没有说：\Quote{我们听说天主对摩西说了话}，而是说：\Quote{我们知道。}犹太人啊，你们把听来的当作知道的，却把亲眼看见的当作不如听来的可靠？但你们没有看见，只是听见；而祂，你们没有听见，却看见了。那么，那盲者说什么呢？\par

% structure: p0030 source=h2 target=subsection reason=relative-heading-rank; base=h2

\subsection{\BibleRef{若 9:30}}

\Quote{一个在你们中间并非卓越、高贵或显赫的人，竟能做这样的事；因此全面清楚地表明祂是天主，不需要人的帮助。}\par

% structure: p0032 source=h2 target=subsection reason=relative-heading-rank; base=h2

\subsection{\BibleRef{若 9:31}}

既然他们最先说\Quote{一个罪人怎能行这样的奇迹呢？}\newline{}\BibleRef{若 9:16}，如今他也引述他们的判断，提醒他们自己的话。他说：\Quote{这个意见，}\newline{}\Quote{是我和你们共同的。现在就坚持它吧。}请看他的智慧。他从各方面翻转这个奇迹，\newline{}因为他们无法否定它，并由此得出推论。你看见他起初说\Quote{他是不是罪人，我不知道}吗？\newline{}不是怀疑（愿天主禁止！）而是知道他不是罪人。至少现在，当他有机会时，看他怎样辩护：\newline{}\Quote{我们知道天主不听罪人：}\par

\Quote{但若有人是崇拜天主的，并遵行祂的旨意。}\par

这里他不仅为他洗脱了罪，还宣称他非常中悦天主，并完全遵行祂的旨意。因为既然他们自称是崇拜天主的人，他就加上\Quote{并遵行祂的旨意}；\Quote{因为，}他说，\Quote{认识天主是不够的：人还必须遵行祂的旨意。}然后他夸大所行的事，说：\par

% structure: p0036 source=h2 target=subsection reason=relative-heading-rank; base=h2

\subsection{若 9:32}

\Quote{如果你们现在承认天主不听罪人，而这个人行了一个奇迹，一个从来没有人行过的奇迹，那么显然他在德行上超越了一切，他的力量大于人所具有的。}那么他们说什么呢？\par

% structure: p0038 source=h2 target=subsection reason=relative-heading-rank; base=h2

\subsection{若 9:34}

只要他们指望他会否认基督，他们就认为他可信，再三召唤他。如果你认为他不可信，为什么又第二次召唤并审问他呢？但是当他毫无畏惧地说出真理时，他们本该最钦佩他，却定了他的罪。但是\Quote{你全然生在罪孽中}是什么意思？他们在这里毫不留情地责备他的瞎眼本身，好像说：\Quote{你从幼年就生在罪中}；暗示他因此生来瞎眼；这是不合逻辑的。至少在这方面，基督安慰他说：\Quote{我为审判来到这世界，叫那看不见的看见，叫那看见的变成瞎眼。}\newline{}\BibleRef{若 9:39}\par

\Quote{你全然生在罪孽中，还要教训我们吗？}什么，这个人说过什么？他提出了自己的意见吗？他不是提出了一个共同的判断，说\Quote{我们知道天主不听罪人}吗？他不是引用了你们自己的话吗？\par

\Quote{他们就把他赶出去了。}\par

你看见这真理的传报者吗？贫穷如何没有阻碍他的真智慧？你看见他从起初就忍受了怎样的斥责、怎样的苦难，如何用言语和行动作证吗？\par

4. 这些事被记录下来，是为了我们也能效法。如果那个瞎子、乞丐，甚至没有见过他，就在被基督鼓励之前，立刻显出这样的勇气，对抗整个充满杀意、被附身、疯狂的人民，他们想借他的声音来定基督的罪；如果他既不屈服也不退缩，而是最勇敢地堵住他们的口，宁愿被赶出去也不出卖真理；那么，我们这些生活在信德中这么久的人，见过信德所行的无数奇迹，领受比他更大的恩惠，恢复了内在眼睛的视力，看见了不可言喻的奥迹，并被召叫到这样尊荣的地位——我们应该如何向那些企图控告、说反对基督徒的话的人表现出全部的勇气，堵住他们的口，而不是不费力气地默许呢？如果我们勇敢，留心圣经，并认真地听，我们就能做到这一点。因为如果有人经常到这里来，即使在家不读经，只要注意这里所说的，一年就足以使他精通圣经；因为我们不是今天读一种圣经，明天读另一种，而是始终反复读同样的。然而，许多人的可悲心态是，在读了这么多之后，他们连经卷的名称都不知道，也不羞愧，不畏惧，如此漫不经心地进入一个可以听到天主言语的地方。可是，如果有一个弹竖琴的、跳舞的或演戏的召唤全城，他们都急切地跑去，为那召唤而感激他，并花上半天时间只注意他一个人；但当天主借着先知和宗徒向我们说话时，我们打哈欠，搔痒，昏昏欲睡。夏天，似乎太热，我们就跑到市场；冬天，雨和泥也成了障碍，我们就待在家里；然而在赛马场，即使没有屋顶遮雨，大多数人在大雨倾盆、风把水吹到脸上的时候，却像疯子一样站着，不顾寒冷、潮湿、泥泞和路途遥远，什么也不能把他们留在家中或阻止他们去那里。可是在这里，头顶有屋顶，温暖宜人，他们反而退缩，不聚拢来；何况这收益是自己的灵魂。告诉我，这如何能忍受？因此，我们在那些事上比任何人都更在行，而在必要的事上却比孩子更无知。如果有人称你为车夫或舞者，你说你受了侮辱，想尽办法洗刷这个侮辱；但如果他拉你去看表演，你却不退缩，那你不愿听到名字的技艺，你几乎每次都去追求。可是在你应当既有行为又有名字、既成为又被称为基督徒的地方，你甚至不知道那行为是什么。还有什么比这愚昧更糟的呢？我一直想对你们说这些话，但我怕我会白白地招致怨恨。因为我看不仅年轻人疯狂，老年人也疯狂；关于他们，我尤其感到羞耻，当我看见一个白发可敬的老人，玷污他那白发，并拉着一个孩子跟他走时。还有比这更可笑的吗？还有比这更无耻的吗？父亲教孩子做不体面的事。\par

5. 这些话刺痛人吗？这正是我所希望的：愿你们因这些话而受苦，好能摆脱因行为而来的羞耻。因为有些人比这些更冷漠，甚至不以所谈论的事为耻，反而为行为编织长篇辩护。你若问他们亚毛斯或敖巴狄雅是谁，或先知与宗徒的数目是多少，他们连嘴都张不开；但讲到马匹和战车，他们比诡辩家或修辞学家更巧妙地编造借口。这一切之后，他们却说：\Quote{有什么害处？有什么损失？} 这正是我所悲叹的：你们竟不知这行为是损失，也不感觉其罪恶。天主赐给你们一段指定的生命时期为事奉祂，而你们却虚度光阴，漫无目的，一事无成，还问：\Quote{有什么损失？} 若你们为无益的事花了一点钱，便称之为损失；但当你们将整日的时间耗费在魔鬼的炫耀上时，难道以为没有做错吗？你们应将全部生命用于恳求与祈祷，而你们却将生命与财物轻率地、有损地浪费在喧嚷、骚动、可耻的言语、争斗、不合时宜的欢乐与诡诈的行为上，之后还要问：\Quote{有什么损失？} 岂不知除了时间，什么都不应挥霍。金子，你若花掉了，还能再得；但时间若失去，便难以挽回。在此生中，我们所得的时间很少；若我们不适当使用，当我们离开前往\Quote{那里}时，将说什么？请告诉我，若你吩咐你的一个儿子学一门技艺，而他一直待在家里，或甚至在其他地方消磨时间，那师父岂不要拒绝他？他岂不对你说：\Quote{你与我立了约，定了时间；如今你的儿子若不与我共度这段时间，而去别处，我怎能把他作为学徒交给你？} 我们也当这样说。因为天主也要对我们说：\Quote{我给了你时间来学习这门虔敬的技艺，为何你们愚蠢而无益地浪费了那时间？为何你们既不常到师父那里，也不留意他的话？} 为表明虔敬是一门技艺，请听先知的话：\BibleQuote{孩子们，来，听我的话，我要教导你们敬畏上主。}\BibleRef{咏33:11} 又说：\BibleQuote{上主，你所教训的人，你所教导于他法律的人，是有福的。}\BibleRef{咏93:12} 因此，当你们虚度了这时间，还有什么借口？有人会说：\Quote{为何祂只给我们很短的时间？} 哦，愚昧与忘恩！你本应为此最当感谢祂：祂缩短了你的劳苦，减轻了你的辛劳，并使其余的时间长久而永恒——你竟为此抱怨，心怀不满吗？\par

但我不知我们怎会把话题引到这里，并讲了这么久；因此现在必须缩短。因为这也是我们可怜的一部分：在此若讲道太长，我们都变得漫不经心；而在那里，他们从午间开始，直到火炬与灯光下才散去。然而，为免我们总是责骂，现在我们恳求并央求你们：赐予我们和你们自己这份恩惠；摆脱一切其他事务，让我们专注于此。这样，我们将从你们那里得到喜乐与欢欣，因你们而获得光荣，并为这些劳苦而得到酬报；而你们将收获一切赏报，因为你们从前如此疯狂地迷恋舞台，如今却因敬畏天主并因我们的劝勉，从那疾病中挣脱出来，挣断锁链，奔向天主。不仅\Quote{那里}你们将得到赏报，\Quote{这里}也将享受纯正的喜乐。德性就是这样：除了在天上赐予我们冠冕，甚至在此世也使生活甘饴。因此，让我们被所说的话说服，好获得今世与来世的福乐，藉着我们的主耶稣基督的恩宠与慈爱，祂和圣父、圣神同受光荣，现在、永远，直到无穷之世。阿们。\par

\SourceRecord{Translated by Charles Marriott.；Nicene and Post-Nicene Fathers, First Series；Vol. 14.；Edited by Philip Schaff.；Buffalo, NY: Christian Literature Publishing Co.,；1889.；<http://www.newadvent.org/fathers/240158.htm>.}

% structure: p0001 source=h1 target=section reason=page-title

\section{若望福音第五十九篇讲道}

% structure: p0002 source=h2 target=subsection reason=relative-heading-rank; base=h2

\subsection{若望福音 9:34-36}

1. 凡为真理和宣认基督的缘故而忍受任何可怕之事并受到侮辱的，这些人尤其受到尊荣。因为正如那为祂的缘故失去自己财产的人，正是最得着它们的人；正如那恨恶自己生命的人，正是最爱它的人；同样，那受侮辱的人，正是最受尊荣的人。正如在瞎子身上所发生的。犹太人把他从圣殿里赶出去，圣殿之主却找到了他；他与那有害的团体分离，遇见了救恩的泉源；他被那些侮辱基督的人所侮辱，却被天使之主所尊荣。真理的奖赏就是这样。因此我们，若我们在此世舍弃我们的财产，就在来世找到自信；若在此世我们周济受苦的人，就在天堂得安息；若我们为天主的缘故受侮辱，无论在此世还是来世，我们都受尊荣。\par

当他们把他从圣殿里赶出去以后，耶稣找到了他。圣史表明，祂是特意来会他的。请注意祂如何酬报他，以那最首要的福分。因为祂使那素不认识祂的人认识了自己，并将他列入祂自己的门徒之列。也请注意圣史如何精确地描述当时的情形；因为当基督说：\Quote{你信天主子吗？}那人回答说：\Quote{主，是谁？}因为他那时还不认识祂，虽然已被治好；因为他在来到他的恩主面前之前是瞎眼的，而治愈之后，他又被那些狗折磨。因此，祂如同赛场上的一位裁判，接见了那位曾辛勤努力并获得冠冕的斗士。祂说了什么？\Quote{你信天主子吗？}这是什么意思？在与犹太人争论了那么多之后，说了那么多话之后，祂问他：\Quote{你信吗？}祂这样说并非出于不知，而是想要使自己被认识，并表明祂温和地看重此人的信德。祂说：\Quote{这群众多的人侮辱了我，但我毫不理会；我关心一件事，就是你应当信。因为一个承行天主旨意的人，胜过一万个犯法者。}\Quote{你信天主子吗？}如同祂曾亲临并赞许他所说的话，祂提出这个问题；首先，祂使他心中渴望自己。因为祂没有直接说：\Quote{信，}而是以询问的方式。那人怎样说？\Quote{主，是谁，好让我信祂？}这话出自一个渴望并询问的灵魂。他不知道他为之说了那么多话辩护的祂是谁，好叫你们学习他的爱真理之心。因为他还没有见过祂。\par

% structure: p0005 source=h2 target=subsection reason=relative-heading-rank; base=h2

\subsection{若望福音 9:37}

祂没有说：\Quote{我就是，}而是以一种中间而保留的方式说：\Quote{你已经见过祂了。}这仍然不确定；因此祂更清楚地补充说：\Quote{和你说话的就是祂。}\par

% structure: p0007 source=h2 target=subsection reason=relative-heading-rank; base=h2

\subsection{若望福音 9:38}

祂没有说：\Quote{我就是那治好你的，那吩咐你：去，在史罗亚洗的；}而是对这一切保持沉默，祂说：\Quote{你信天主子吗？}然后那人显出极大的热诚，立刻朝拜；这在被治好的人中是少有的；例如，那些癞病人和一些其他人；以此行为宣告了祂的神性大能。免得有人认为他所说的话只是空谈，他加上了行动。当他朝拜时，基督说：\par

% structure: p0009 source=h2 target=subsection reason=relative-heading-rank; base=h2

\subsection{若望福音 9:39}

保禄也这样说：\Quote{那么，我们可说什么呢？那未追求义的外邦人获得了义，就是出于信德由耶稣而来的义；但以色列追求义的法律，却未能达到义的法律。}\BibleRef{罗 9:30} 藉着说：\Quote{我为审判来到这世界，}祂既使那人在信德上更坚固，也唤起了跟随祂的人；因为法利塞人正跟着祂。而\Quote{为审判，}祂是就更大的惩罚而言；表明那些宣判祂的人，自己已经接受了判决；那些定祂为罪人的人，正是被定罪的人。在此处祂谈到两种复明和两种盲目：一种是感性的，另一种是属神的。\par

% structure: p0011 source=h2 target=subsection reason=relative-heading-rank; base=h2

\subsection{若望福音 9:40}

如同在另一处他们说：\Quote{我们从未做过任何人的奴隶，}以及\Quote{我们不是由淫乱生的}\EditorialNote{见第八章33、41节}；如今他们只盯住物质事物，并以此为羞耻。然后为表明他们盲目比看见更好，祂说：\par

% structure: p0013 source=h2 target=subsection reason=relative-heading-rank; base=h2

\subsection{若望福音 9:41}

既然他们把这灾难视为可耻之事，祂便将这转回他们自己头上，告诉他们：\Quote{这件事实在会使你们的惩罚更可容忍}；从各方面削除他们属人的思想，引导他们达到崇高而奇妙的认识。\par

\Quote{但如今你们说：我们看见。}\par

如同在另一处祂说：\Quote{你们曾说祂是你们的天主}\BibleRef{若 8:54}；同样在此处：\Quote{如今你们说你们看见，但你们却没有看见。}祂表明他们引以为荣的大事，反而给他们带来惩罚。祂也安慰那生来瞎眼的人，有关他先前的残缺，然后论及他们的盲目。因为祂把整篇谈话导向这个目的，免得他们说：\Quote{我们并非因盲目而拒绝到你这里来，而是我们避开你，视你为欺骗者。}\par

2. 圣史并非无端地提到，那些与祂在一起的法利塞人听见了这些话，就说：\Quote{我们也盲目吗？}而是为提醒你们，这些人就是首先离开祂，然后拿石头打祂的人；因为他们只是表面上跟随祂，且容易改变为相反的意见。那么祂如何证明自己不是欺骗者，而是牧人呢？藉着列出牧人与欺骗者、掠夺者的区别标志，并由此给他们机会去查究此事的真理。首先，祂指出谁是欺骗者和掠夺者，从圣经中称他为这样，并说：\par

% structure: p0018 source=h2 target=subsection reason=relative-heading-rank; base=h2

\subsection{若望福音 10:1}

请注意强盗的特征：首先，他不公开进入；其次，不按照圣经，因为这就是\Quote{不从门进入}。这里祂也指那些在祂之前和之后的人，假基督和伪基督，犹达斯和特乌达，以及其他一切同类的人。祂有充分理由称圣经为\Quote{门}，因为圣经引领我们到天主那里，向我们开启天主的认识，祂使羊成为羊，保护它们，不允许狼跟在它们后面进来。因为圣经，像一道可靠的门，阻挡了异端者的通道，使我们处于安全的状态，关于我们所渴望的一切，不让我们迷失；如果我们不打开它，我们将不容易被敌人征服。通过它，我们可以认识一切，无论是牧人还是非牧人。但什么是\Quote{进入羊栈}？它指的是羊和照顾羊的人。因为那不使用圣经，却\Quote{从别处爬进去}，即为自己开辟另一条不寻常的道路的人，\Quote{就是贼}。你从这里也看到基督与父一致，因为祂提出圣经？为此祂也对犹太人说：\Quote{查考圣经}（见：\BibleRef{若 5:39}），并引证梅瑟，称他和所有先知为证人，因为\Quote{所有听从先知的人，}祂说，\Quote{都要到我这里来}；以及\Quote{如果你们信了梅瑟，就会信我。}但这里祂用比喻说了同样的事。祂说\Quote{从别处爬进去}，暗示了经师，因为他们把人的诫命当作教义来教导，并违犯了法律（见：\BibleRef{玛 15:9}）；祂责备他们，说：\Quote{你们中间没有一个人遵行法律。}（见：\BibleRef{若 7:19}）祂说得很好：\Quote{爬进去}，而不是\Quote{进入}，因为爬是贼想要翻越墙的行为，做一切事都带着危险。你看到祂如何描绘强盗了吗？现在观察牧人的特征。是什么呢？\par

% structure: p0020 source=h2 target=subsection reason=relative-heading-rank; base=h2

\subsection{若望福音 10:2-4}

3. 祂列出了牧人和恶人的特征；现在让我们看祂如何将它们与接下来的内容配合。祂说：\Quote{守门者就给他开门}，祂继续使用比喻，使话语更有力。但如果你愿意逐字考察这比喻，没有什么妨碍你假设梅瑟是守门者，因为上主的圣言交托给了他。\Quote{羊听他的声音，他按着名字叫自己的羊。}因为那些人到处说祂是骗子，并用自己的不信来证明这一点，说：\Quote{有哪个官长信了他？}（见：\BibleRef{若 8:48}）祂表明，他们不应该因为那些人的不信而称祂为掠夺者和骗子，而是因为他们自己不听祂，甚至被排除在羊的行列之外。因为如果牧人的部分是通过通常的门进入，而祂是通过这门进入，那么所有跟随祂的人可能是羊，但那些自我分离的人，不损害牧人的声誉，而是将自己从羊的亲属中驱逐出去。如果祂进一步说祂是\Quote{门}，我们不必再次困扰，因为祂也称自己为\Quote{牧人}和\Quote{羊}，并以不同的方式宣布祂的安排。因此，当祂带我们到父那里时，祂称自己为\Quote{门}；当祂照顾我们时，祂称自己为\Quote{牧人}；并且为了不让你以为带我们到父那里是祂唯一的职务，祂称自己为牧人。\Quote{羊听他的声音，他按着名字叫自己的羊，并把他们领出来，在他们前面走。}牧人确实做相反的事，因为他们跟在后面；但祂为了表明祂将带领所有人走向真理，做了不同的事；正如当祂派遣羊时，祂不是派遣他们避开狼的道路，而是\Quote{在狼群中}（见：\BibleRef{玛 10:16}）。因为这种牧羊的方式比我们的更奇妙。在我看来，祂也暗示了那个瞎子，因为祂也\Quote{叫}了他，并把他\Quote{领出}犹太人中间，那人听见了\Quote{祂的声音}，并\Quote{认出}了它。\par

% structure: p0022 source=h2 target=subsection reason=relative-heading-rank; base=h2

\subsection{若望福音 10:5}

当然，这里祂说的是特乌达和犹达斯（因为\Quote{所有信了他们的人都四散了}——\BibleRef{宗 5:36}），\newline{}或是在那之后会欺骗人的假基督。为了不让人说祂是其中之一，祂以多种方式将自己与他们区别开来。第一个区别是祂从圣经而来的教导；因为祂借此引领人们到祂那里，而那些人并不用圣经吸引人跟随他们。第二个区别是羊的服从；因为所有的人都信了祂，不仅祂活着的时候，而且在祂死了之后；而那些人，人们立刻离开了他们。此外，我们还可以提到第三个区别，并非不重要。那些人都是作为反叛者行事，引起叛乱，而祂则远离这种嫌疑，以至于当他们要立祂为王时，祂逃走了；当他们问：\Quote{给凯撒纳税可以吗？}祂吩咐他们缴纳，并且自己缴纳了两德拉克马（见：\BibleRef{玛 17:27}）。除此之外，祂来是为了拯救羊，\Quote{为使他们获得生命，且获得更丰富的生命}（见：\BibleRef{若 10:10}），而那些人甚至剥夺了他们今世的生命。那些人背叛了托付给他们的人并逃跑了，而祂勇敢地忍受，甚至舍弃了自己的生命。那些人是不情愿、被迫、想要逃脱而遭受他们所遭受的，而祂是自愿、选择忍受一切。\par

% structure: p0024 source=h2 target=subsection reason=relative-heading-rank; base=h2

\subsection{若望福音 10:6}

祂为什么说得隐晦？因为祂要使听众更加专注；当祂达到了这个目的，祂就消除了隐晦，说：\par

% structure: p0026 source=h2 target=subsection reason=relative-heading-rank; base=h2

\subsection{若望福音 10:9}

好似祂曾说过：\Quote{必得安全与稳妥}，（但祂在此以\Quote{草场}指祂对羊群的养育和喂养，以及祂的权能与主权，）即\Quote{必得留在内，无人能将他赶出}。这便发生在宗徒们身上，他们安然而入又安然出，如同已成为全世界的主人，无人能赶逐他们。\par

% structure: p0028 source=h2 target=subsection reason=relative-heading-rank; base=h2

\subsection{若望福音 10:8}

祂在此并非指先知（如异端者所断言），因为凡信基督者，也听从并信服他们；而是指特乌达、犹大及其他煽动叛乱者。况且，祂说\BibleQuote{羊却不听他们}是称赞他们；我们从不曾见祂赞许那些拒绝听从先知的人，反而严厉责备并控告他们；由此显然\Quote{不听}是指那些叛乱首领。\par

% structure: p0030 source=h2 target=subsection reason=relative-heading-rank; base=h2

\subsection{若望福音 10:10}

这便在他们（其跟随者）全被杀灭时发生了。\par

\BibleQuote{但我来，是为叫他们获得生命，且获得更丰富的生命。}\par

那么，请问\Quote{更丰富}比生命还多的是什么？是天国。但祂尚未明言，而停留于\Quote{生命}这名号上，因这是他们所熟知的。\par

% structure: p0034 source=h2 target=subsection reason=relative-heading-rank; base=h2

\subsection{若望福音 10:11}

此处祂接着讲论受难，显示这应为世界的救恩，祂并非不情愿地走向它。然后祂再次提及牧人与佣工的特质。\par

\BibleQuote{因为牧人舍命。}\par

% structure: p0037 source=h2 target=subsection reason=relative-heading-rank; base=h2

\subsection{若望福音 10:12}

在此祂宣称自己与父同为主人，若祂是牧人，羊群便是祂的。你们看见祂在比喻中言说更为高超，而意义隐晦，不给听众留下明显把柄？那么，这佣工做了什么？他\Quote{看见狼来，就撇下羊，狼就来，把它们驱散}。那假教师正是如此，而祂相反。因当祂被捉拿时，祂说：\BibleQuote{让这些人去罢，为应验那话}\EditorialNote{若 18:8-9}，使祂没有一个失落。此处我们也可猜测暗指属灵的狼；因为基督也未容许它去攫取羊群。但它不仅是狼，也是狮子。经上说：\BibleQuote{因为你们的仇敌魔鬼，如同咆哮的狮子巡游}\BibleRef{伯前 5:8}。它也是蛇和龙；因为\BibleQuote{你们践踏蛇和蝎子}\BibleRef{路 10:19}。\par

4. 因此，我恳求你们，让我们留在这位牧人之下放牧；我们若听从他，若听从他的声音，若不相随陌生人，我们便会留下。祂的声音是什么？\BibleQuote{神贫的人是有福的，心地纯洁的人是有福的，怜悯人的人是有福的。}\BibleRef{玛 5:3} 我们若这样行，就必留在这牧人之下，狼就不能进来；即使它来攻击我们，也必自受亏损。因为我们有一位如此爱我们的牧人，他甚至为我们舍弃了自己的生命。因此，当他既有大能又爱我们时，还有什么能阻碍我们得救呢？没有，除非我们自己背叛他。我们怎能背叛呢？听他说：\BibleQuote{你们不能事奉两个主人：天主和\OriginalTerm{玛门}{mammon}。}\BibleRef{玛 6:24} 如果我们事奉天主，就不会屈服于玛门的暴政。实在说，比任何暴政更苦的事就是贪图财富；因为贪财不能带来快乐，反而带来忧虑、嫉妒、阴谋、仇恨、诬告，以及千万种对德行的障碍：懒惰、放荡、贪婪、酗酒；这些使自由人成为奴隶，而且比用金钱买来的奴隶更糟，不是作人的奴隶，而是作最痛苦的情欲和灵魂疾病的奴隶。这样的人敢做许多天主和人都厌恶的事，唯恐有人夺去他的统治。啊，痛苦的奴役，魔鬼般的暴政！因为最严重的是，当我们陷于这些恶中时，我们竟喜欢并拥抱自己的锁链，居住在充满黑暗的监狱里，拒绝来到光明中，反而把恶牢牢拴在自己身上，以我们的疾病为乐。因此我们无法获得自由，比那些在矿坑中工作的人更糟，他们虽忍受劳苦和磨难，却不能享受果实。而实际上，比一切更糟的是，如果有人想要解救我们脱离这痛苦的囚禁，我们却不允许，甚至恼怒和不悦；在这方面，我们并不比疯子好，甚至比他们更可怜，因为我们甚至不愿意从疯狂中被解救出来。什么？人啊，你被带到世上是为这个吗？你被造成人，是为在矿坑中劳作，聚集金子吗？天主不是为此按他的肖像造了你，而是为让你中悦他，为获得将来的事物，为加入天使的唱诗班。你为什么现在将自己从这样的关系中放逐，把自己投入极致的羞辱和卑贱中呢？那与你经历同样生产之痛（我指属灵的生产之痛）的人，正因饥饿而垂危，而你却胀得饱满：你的兄弟赤身露体，你却给衣服再添衣服，为虫子堆积这些衣服。若你把它们穿在穷人身上，那将好得多；这样它们就不会毁坏，会使你免去一切忧虑，并为你赢得来世的生命。如果你不愿意它们被虫蛀，就送给穷人吧，这些人知道如何好好抖落这些衣服。基督的身体比柜子更珍贵、更安全，因为他不仅保护衣服安然无恙，不仅保全它们不被消耗，甚至使它们更加明亮。柜子往往连同衣服一起被偷，使你遭受极大损失，但基督的身体这个安全之处，连死亡也不能伤害。与他在一起，我们不需要门、锁、警觉的仆人或任何其他安全保障，因为我们的财产免受一切诡诈攻击，且被看守起来，如同我们可以想象天上存放的东西那样；因为那个地方是邪恶无法进入的。我们不断向你们说这些事，你们听了却不信服。原因是我们灵魂卑贱，只盯着大地，匍匐在地上。或者说，但愿天主不许我定你们全都有罪，好像所有人都患了不治之症。因为即使那些沉醉于财富的人掩耳不听我的话，那些生活在贫穷中的人却能明白我所说的。\Quote{但是什么，}有人说，\Quote{这与穷人有什么关系？因为他们没有金子，也没有这样的衣服。} 不，但他们有面包和凉水，他们有\OriginalTerm{两文小钱}{two obols}，他们有脚可以探访病人，他们有舌头和言语可以安慰卧床者，他们有房屋和住处可以接待客旅。我们不向穷人要求若干金币，这些我们从富人要求。但若有人贫穷，来到别人的门口，我们的主并不会羞于收下一个\OriginalTerm{奥波尔}{obol}，反而会说，他从这施予者那里接受的比那些投入许多的人更多。现在站在这里的多少人曾渴望生在那个时候，当基督带着肉身在地上行走时，能与他交谈，与他同席？看，如今也可以这样做，我们可以邀请他吃饭，与他一同坐席，而且更有益处。因为当时与他同席的人中，许多甚至丧亡了，如犹达斯及跟他一样的人；但现在每个邀请他到家中，与他共享餐桌和屋顶的人，都将享受极大的祝福。它说：\BibleQuote{你们蒙我父赐福的，来承受从创世以来为你们预备的国度吧。因为我饿了，你们给我吃的；我渴了，你们给我喝的；我作客，你们收留了我；我患病，你们看顾了我；我在监里，你们来探望了我。}\BibleRef{玛 25:34} 为使我们可以听到这些话，让我们给赤身者衣服穿，收留客旅，给饥饿者食物，给口渴者喝，探访病人，看望囚犯，好使我们能坦然无惧，获得罪过的赦免，并分享那些超越言语和思想的良好事物。愿我们都能获得这一切，借着我们的主耶稣基督的恩宠和慈爱，愿光荣和大能归于他，直到永远。阿们。\par

\SourceRecord{Translated by Charles Marriott.；Nicene and Post-Nicene Fathers, First Series；Vol. 14.；Edited by Philip Schaff.；Buffalo, NY: Christian Literature Publishing Co.,；1889.；<http://www.newadvent.org/fathers/240159.htm>.}

% structure: p0001 source=h1 target=section reason=page-title

\section{若望福音讲道集第六十篇}

% structure: p0002 source=h2 target=subsection reason=relative-heading-rank; base=h2

\subsection{若 10:14-15}

1. 亲爱的，治理教会是一件大事，一件大事；需要如同基督所说的智慧和勇气，即人为羊舍命，永不撇下羊群或使之赤身露体；要勇敢地对抗豺狼。因为牧人与佣工的区别就在于此：一个总是顾及自身的安全，不关心羊群；另一个总是寻求羊群的益处，忽略自己。因此，在列出牧人的标志之后，基督提出了两种破坏者：一种是贼，杀死并偷窃；另一种是不做这些事，但当这些事发生时却不在意、不阻止。第一种是指\Person{特乌达}及其同伙；第二种是指犹太人的教师，他们既不关心也不顾及托付给他们的羊群。为此，从前\Person{厄则克耳}斥责他们，说：\BibleQuote{祸哉，以色列的牧者！牧者岂是牧养自己？牧者岂不是牧养羊群？}\BibleRef{则 34:2}但他们反其道而行，这是最恶劣的邪恶，也是一切罪恶的根源。因此经上说：\BibleQuote{他们没有召回迷路的，没有寻找失落的，没有包扎受伤的，没有医治患病的，因为他们牧养自己，而非羊群。}\BibleRef{则 34:4}\Person{保禄}也在另一处宣告说：\BibleQuote{众人只顾自己的事，不顾耶稣基督的事。}\BibleRef{斐 2:21}又说：\BibleQuote{人不要寻求自己的利益，而要求别人的利益。}\BibleRef{格前 10:24}基督从这两方面将自己区分开来：从那些来破坏的，祂说：\BibleQuote{我来，是为使他们获得生命，且获得更丰富的生命。}\BibleRef{若 10:10}从那些不关心羊群被狼叼走的，祂从不离弃羊群，甚至舍命为羊，使羊不丧亡。因为当众人想要杀祂时，祂既不改换教导，也不背叛信祂的人，而是坚持到底，选择死亡。因此祂不断地说：\BibleQuote{我是善牧。}然后，因为祂的话似乎没有见证支持，（因为\BibleQuote{我舍弃我的生命}这句话不久后就应验了，但\BibleQuote{使他们获得生命，且获得更丰富的生命}要在他们离开此世后的来生才应验），祂做什么呢？祂用一个证明另一个：通过祂的会死生命（证明）祂赐予不朽生命。正如\Person{保禄}也说：\BibleQuote{如果我们还在为仇敌的时候，因着祂圣子的死，与天主和好了，那么，在和好之后，我们更要因祂的生命得救了。}\BibleRef{罗 5:10}又在另一处说：\BibleQuote{祂没有怜惜自己的儿子，反而为我们众人把祂交出，岂不也把万有与祂一同赐给我们吗？}\BibleRef{罗 8:32}\par

但是，为何他们现在不拿以前加给祂的指控来攻击祂？那时他们说：\BibleQuote{你为自己作证，你的证据是不可信的。}\BibleRef{若 8:13}因为祂曾多次堵住他们的口，并且祂的勇气因奇迹而增加。然后，因为祂在上文说\BibleQuote{羊听他的声音，跟从他}，恐怕有人说：\BibleQuote{那么，那些不信的人又怎么样呢？}请听祂补充说：\BibleQuote{我认识我的羊，我的羊也认识我。}正如\Person{保禄}声明说：\BibleQuote{天主并没有弃绝祂预先所知道的百姓}\BibleRef{罗 11:2}；\Person{梅瑟}说：\BibleQuote{主认识属于祂的人}\BibleRef{弟后 2:19}；祂说：\BibleQuote{那些}，即\BibleQuote{祂所预知的人。}然后，以免你认为认识的程度相等，请听祂如何通过补充来纠正：\BibleQuote{我认识我的羊，我的羊也认识我。}但这认识并不相等。\BibleQuote{哪里相等？}在圣父和我的情况下，因为在那里：\BibleQuote{圣父认识我，我也认识圣父。}\BibleQuote{我认识圣父正如圣父认识我一样。}如果祂不想证明这一点，祂为何要加上那句话？因为祂常将自己列于众人之中，所以为了避免有人以为祂像人一样认识，祂补充说：\BibleQuote{圣父认识我，我也认识圣父。}\BibleQuote{我认识圣父正如圣父认识我一样。}因此祂说：\BibleQuote{除了父，没有人认识子；除了子，也没有人认识父}\BibleRef{路 10:22}，谈到一种独特的认识，是其他任何受造物无法拥有的。\par

2. \BibleQuote{我舍弃我的生命。}祂不断这样说，以表明祂不是欺骗者。同样，宗徒想要表明他是真教师，并与假宗徒辩论时，也以他的危险和死亡来建立权威，说：\BibleQuote{多受鞭打，屡次濒于死亡。}\BibleRef{格后 11:23}因为说\BibleQuote{我是光}和\BibleQuote{我是生命}，在愚笨的人看来似乎是傲慢；但说\BibleQuote{我愿意死}，却没有存有任何恶意或嫉妒。因此，他们不对祂说：\BibleQuote{你为自己作证，你的证据是不可信的}，因为这话表现出对祂们的深切关怀，如果祂确实愿意为那些要砸死祂的人舍己。为此，祂也适时引进了外邦人的话题。\par

% structure: p0006 source=h2 target=subsection reason=relative-heading-rank; base=h2

\subsection{若 10:16}

请再注意，这里所用的\Quote{必须}一词，并不表示必然性，而是宣告某件事必定会发生。如同祂曾说：\Quote{你们为何惊奇这些羊要跟随我，我的羊要听我的声音？当你们看见其他人也跟随我、听我的声音时，你们会更加惊讶。}当你们听祂说\Quote{不属于这羊栈的}时，不要惊慌；因为这区别只与法律有关，正如\Person{保禄}所说：\BibleQuote{割损算不了什么，不割损也算不了什么。}\BibleRef{迦 5:6}\par

祂说：\BibleQuote{我也必须领他们来。}祂显示这些和那些都分散混杂，没有牧人，因为善牧还没有来。然后祂预告他们将来要合一，即：\par

\BibleQuote{将成为一个羊栈。}\par

同样的事，保禄也宣告说：\BibleQuote{为在祂内将二者造成一个新人。}\BibleRef{弗 2:15}\par

% structure: p0011 source=h2 target=subsection reason=relative-heading-rank; base=h2

\subsection{若望福音 10:17}

如果我们的主因我们而被爱，因为祂为我们而死，还有什么比这句话更充满人情味呢？那么如何？请告诉我，难道祂在先前不被爱吗？难道圣父现在才开始爱祂，而我们成了祂爱的原因？你看见祂如何俯就吗？但祂在这里想要证明什么？因为他们说祂与圣父相异，是个骗子，来是要毁坏和消灭，祂告诉他们：\Quote{如果别无其他，这一点会说服我爱你们，即你们如此被圣父所爱，以致我也被祂所爱，因为我为你们而死。}此外，祂还想要证明另一点，即祂并非不情愿地从事这一行动（因为若是不情愿，所做的事怎么能引起爱呢？），而且这一点特别为圣父所知。如果祂以人的方式说话，不必惊奇，因为我们多次提过其原因，再说同样的话就多余且令人不悦了。\par

\Quote{我舍掉我的性命，目的是为了再取回它。}\par

% structure: p0014 source=h2 target=subsection reason=relative-heading-rank; base=h2

\subsection{若望福音 10:18}

因为他们屡次商议要杀祂，祂告诉他们：\Quote{除非我愿意，你们的劳作是徒劳的。}并且，祂以第一点证明第二点，即借着死亡证明复活。因为这是奇特而奇妙的事。两者都以新的方式发生，超乎寻常的惯例。但我们应仔细留意祂所说的话：\Quote{我有权舍掉我的性命。}谁没有权舍掉自己的性命呢？因为凡愿意的人都有能力自杀。但祂并非这样说，而是怎样呢？\Quote{我拥有如此舍掉性命的权能，以致没有人能违背我的意愿这样做。}而这并非人所有的权能；因为除了自杀外，我们没有办法舍掉性命。并且，若我们落入谋害我们的人手中，他们有能力杀我们，我们就不再能自由选择是否舍掉性命，而是他们甚至违背我们的意愿将其夺去。但基督并非如此，即使别人谋害祂，祂也有权不舍掉性命。因此祂说了\Quote{没有人能从我夺去它}，又补充说：\Quote{我有权舍掉我的性命}，即\Quote{只有我能决定是否舍掉它}，这件事不取决于我们，因为还有许多人能将其从我们夺去。祂起初并非这样说（因为这话似乎不可信），但当祂得到了事实的见证，且当众人屡次谋害祂却无法抓住祂时（因为祂无数次从他们手中逃脱），祂才说：\Quote{没有人能从我夺去它。}但如果这是真的，那么另一点随之而来，即祂是自愿走向死亡的。如果这也是真的，那么下一点也确定无疑：祂能按自己的意愿\Quote{再取回它}。因为如果死亡是人无法做到的更伟大的事，那么关于其他事就不要再怀疑了。既然唯独祂能舍弃自己的性命，这就表明祂能运用同样的权能再取回它。你看见了吗？祂如何从第一点证明第二点，并借着祂的死亡表明祂的复活是无可争议的？\par

\Quote{我从圣父接受了这条诫命。}\par

这是什么诫命？是为世界而死。那么，祂是先等候听见，然后选择，并且需要学习它吗？有理智的人谁会主张这点？但在此之前，当祂说：\Quote{因此我的圣父爱我}，祂表明最初的行动是自愿的，并且消除了对圣父的一切反对的猜疑；同样，当祂说从圣父接受了诫命时，除了声明：\Quote{我所做的这事在祂看来是好的}以外，别无他意。如此，当他们要杀死祂时，他们不会以为自己是杀死了一个被圣父抛弃和放弃的人，也不会用他们所用的那些辱骂来责备祂：\BibleQuote{他救了别人，却不能救自己}；以及，\BibleQuote{如果你是天主子，就从十字架上下来}\BibleRef{玛 27:42}；然而，祂不下来的真正原因正是祂是圣子。\par

3. 那么，以免你听见\Quote{我从圣父接受了诫命}时，以为这成就并不属于祂，祂事先说了\Quote{善牧为羊舍掉自己的性命}；以此表明羊是祂的，所发生的一切都是祂的成就，并且祂不需要任何诫命。因为若祂需要诫命，怎能说\Quote{我自愿舍掉它}？因为自愿舍掉的人不需要诫命。祂也指出了祂这样做的原因。那原因是什么？即祂是牧人，而且是善牧。善牧不需要任何人唤醒他尽本分；如果人在此事上如此，天主更是如此。因此，保禄说：\BibleQuote{祂空虚了自己。}\BibleRef{斐 2:7}所以，这里所说的\Quote{诫命}不是别的，只是显示祂与圣父的一致；而若祂以如此谦卑和人性化的方式说话，原因是听者的软弱。\par

% structure: p0019 source=h2 target=subsection reason=relative-heading-rank; base=h2

\subsection{若望福音 10:19}

因为祂的话超越凡人所言，不属寻常，他们便说祂有魔鬼，这样称祂已是第四次。因为他们先前说过：\Quote{你有魔鬼，谁想要杀你？}\BibleRef{若 7:20}；又说：\Quote{我们说得不错吗：你是撒玛黎雅人，且有魔鬼？}\BibleRef{若 8:48}；此处又说：\Quote{祂有魔鬼且疯了，为什么你们听祂？}或者不如说，祂不是第四次听到这话，而是屡次听到。因为问\Quote{我们说得不错吗：你有魔鬼？}，表示他们不是说过两次或三次，而是许多次。\Quote{另有人说：这不是有魔鬼者的话；魔鬼岂能开启瞎子的眼睛？}因为他们无法用言语使对手沉默，现在便从祂的事迹提出证据。当然，其言也不像有魔鬼者；但若你们不被言语说服，就因事迹而羞愧吧。因为若它们不是有魔鬼者的行为，且超越人性，那么很明显它们源自某一种神圣的能力。你看到这论证了吗？因为它们超越凡人是明显的，因犹太人曾说\Quote{祂有魔鬼}；而祂没有魔鬼，祂以所为来证明。\par

那么基督做了什么？祂对这些事没有回答。在此之前祂曾回答说：\Quote{我没有魔鬼}；但现在却不这样；因为祂已用行动提供了证据，便保持沉默。因为那些人也不配得到回答，他们因着那些本应令他们钦佩并认祂为天主的事迹，却说祂着了魔。当他们彼此反对、相互驳斥时，祂哪里还需要进一步的驳斥？因此祂沉默，温和地忍受一切。这不仅为此原因，也为了教导我们所有温和与忍耐。\par

4. 现在让我们效法祂。\newline{}因为祂不仅当时沉默不语，甚至还再次来到他们中间，受询问时予以回答，并显示有关祂预知的事；尽管被那些从祂领受无数恩惠的人（不止一次两次，而是多次）称为\Quote{附魔者}和\Quote{狂人}，祂不仅没有报复，反而继续恩待他们。\newline{}我说恩待吗？祂甚至为他们舍命，且在受钉时在圣父面前为他们求情。\newline{}因此，让我们也效法祂，因为作基督的门徒，就是温和与良善。\newline{}但这种温和从何而来？如果我们不断数算自己的罪过，如果我们哀恸，如果我们哭泣；因为一个常与如此多忧伤相处的灵魂，必不忍发怒或被激怒。\newline{}因为哪里有哀恸，就不可能有愤怒；哪里有悲伤，一切怒气便消散；哪里有心灵的破碎，便没有挑衅。\newline{}因为心灵被忧伤鞭笞时，便无暇被激怒，反而会痛苦呻吟，更加哀哭。\newline{}我知道许多人听到这些会发笑，但我不会停止为那些发笑者哀恸。\newline{}因为现在是哀恸、哭号与悲叹的时刻，因为我们言语行为犯了许多罪，地狱等候那些犯下这些过犯的人，那火河沸腾咆哮，以及被逐出天国——这是最痛苦的事。\newline{}当这些事被警告时，请告诉我，你还发笑、傲慢吗？当你的主发怒警告时，你竟漠不关心，难道不怕因此为自己点燃更旺的火炉吗？\newline{}你难道没有听见祂每天呼喊的话？\Quote{\BibleQuote{我饿了，你们没有给我吃的；我渴了，你们没有给我喝的；你们去吧，到那为魔鬼和他的使者预备的永火里去。}\EditorialNote{玛 25}}\newline{}祂每天都这样警告。\newline{}有人说：\Quote{但我给过祂食物。}\newline{}何时？多少天？十天或二十天？但祂所要的不仅仅限于这段时间，而是你在世上的全部岁月。\newline{}因为那些童女也有油，却不足以拯救自己；她们也点灯，却被关在婚宴厅外。\newline{}这也是合理的，因为新郎到来之前，灯已经熄灭了。\newline{}为此我们需要大量的油，和丰富的慈爱。\newline{}至少听听先知说的：\Quote{\BibleQuote{天主，求祢按祢的大仁慈怜悯我。}\BibleRef{咏 50:1}}\newline{}因此，我们也必须按祂向我们施与的大仁慈，去怜悯我们的邻人。\newline{}因为我们怎样对待同仆，也必发现主怎样对待我们。\newline{}什么样的\Quote{仁慈}算得上\Quote{大}？就是当我们不是从盈余中，而是从匮乏中给予。\newline{}但如果连从盈余中都不肯给，我们还有什么希望？我们怎能从那些灾祸中解脱？我们哪里能逃避并找到救恩？\newline{}因为那些童女经历了那么多、那么大的劳苦之后，仍得不到任何安慰；当我们听见审判者亲口对我们说那些可怕的话、责备我们时，谁能为我们辩护？因为祂说：\Quote{\BibleQuote{我饿了，你们没有给我吃的；因为我实在告诉你们，凡你们没有给这些最小中的一个做的，就是没有给我做。}}祂不仅指祂的门徒，也不仅指那些过苦修生活的人，而是指每一个有信德的人。\newline{}因为这样的人，即使是奴隶，或是那些在市场上讨饭的，只要他信天主，按理就应享有我们全部的善意。\newline{}如果我们忽视这样赤身露体或饥饿的人，我们就会听到那些话。\newline{}这是有理由的。\newline{}祂向我们要求了什么困难或痛苦的事？岂不都是最轻省容易的事？\newline{}祂不是说：\Quote{\BibleQuote{我病了，你们没有治好我}}，而是说：\Quote{\BibleQuote{你们没有来看顾我。}}\newline{}祂不是说：\Quote{\BibleQuote{我在监里，你们没有释放我}}，而是说：\Quote{\BibleQuote{你们没有到我这里来。}}\newline{}因此，命令越容易，违背者的惩罚就越大。\newline{}请告诉我，有什么比走出去进入监牢更容易？有什么比这更愉快？\newline{}因为当你看见有些人被捆绑，有些人满身污秽，有些人头发不剪、衣衫褴褛，有些人饥饿将死、像狗一样跑到你脚边，有些人肋旁深深凹陷，有些人戴着锁链从市场回来，整天乞讨却连必需的食物都得不到，晚上还要向监管者提供那邪恶野蛮的服务；即使你像石头一样硬，一定会变得更加仁慈；即使你过着放荡安逸的生活，也一定会变得更加明智，因为你在他人的不幸中观察到人事的本质；你必定会对那可怕的日子及其各种惩罚产生概念。\newline{}反复思考这些事，你一定能摒弃愤怒、享乐和对世俗之物的贪爱，使你的灵魂比最平静的港湾更宁静；你也会思考那审判台，思想如果人间尚有如此多的先见、秩序、恐惧和警告，那么在天主那里岂不更多？\newline{}\Quote{\BibleQuote{因为权柄没有不是从天主来的。}\BibleRef{罗 13:1}}\newline{}因此，祂容许统治者如此安排这些事，祂自己更会如此行。\par

5. 诚然，若无此畏惧，一切都将丧失，因为尽管有如此惩罚悬于头顶，仍有许多人投靠邪恶。你若明智地观察这些事，便会更乐意施舍，收获远超过那些从剧院出来之人的喜悦。因为他们离开那里时，被欲望点燃并燃烧。眼见那些女子在舞台上徘徊，从她们身上受到千万创伤，他们将如翻腾的大海一样不得安宁，因为面容、姿态、言辞、步态以及其余一切的形像都立在眼前，围攻他们的灵魂。但从监狱出来的人却不会遭受此类事情，他们将享有极大的平静与安宁。因为因看见囚犯而生的痛悔扑灭了所有那火焰。若一个妓女或荡妇遇见从囚犯中出来的人，她无法加害于他。因为他在那段时间变得仿佛无法塑形，因此不会被她容貌的网罗所捕获，因为那放荡的容貌之前，将会有审判的畏惧置于他眼前。因此，那尝遍各种奢华的人说：\Quote{往遭丧的家去，胜过往宴乐的家去。}\BibleRef{训 7:2} 因此\Quote{在这里}你将显出极大的智慧，\Quote{在那里}你将听到那些值得万千祝福的话语。那么，让我们不要忽略这样的实践和操练。因为我们即使无法给他们带来食物，也无法用金钱帮助他们，但藉着言语，我们能安慰他们，振作沮丧的精神，并透过与那些把他们投入监狱的人交谈，使看守变得更仁慈，从而在许多其他方面帮助他们，我们必定会成就或大或小的善事。但若你说那里的人既非有地位，亦非善良温顺，而是杀人犯、盗墓者、扒手、奸淫者、放纵者、充满诸多恶行，那么这又向我显示了一个常去那里的紧迫理由。因为我们并未奉命只怜悯好人而惩罚恶人，而是要向众人表示这种仁爱。经上说：\Quote{你们应当像你们的天父一样，因为他使太阳上升，光照恶人，也光照善人；降雨给义人，也给不义的人。}\BibleRef{玛 5:45} 因此，不要刻薄地指责他人的过错，也不要作严厉的审判官，而要温和而仁慈。因为如果我们没有犯奸淫、盗墓或扒窃，我们仍有其他过犯，应受无限的惩罚。或许我们曾称弟兄为\Quote{傻子}，这为我们准备了地狱之火；我们曾以不贞洁的目光注视妇女，这构成了完全的奸淫；而比这一切更严重的是，我们不配地领受圣体圣事，这使我们干犯基督的体血。因此，不要苛刻地追究他人的行为，而要省察自己的状况，这样我们就会停止这种不人道和残忍。此外，可以说，我们在那里会找到许多好人，常常有堪比全城价值的人。因为甚至那关押若瑟的监狱里也有许多恶人，但那位义人看管他们所有人，并且他的真实身份被隐藏了；因为他价值等同于全埃及地，却仍住在监狱中，里面无人认识他。同样，如今也很有可能有许多良善而有德的人，虽然他们不为所有人所见，而你照顾这样的人所付出的努力，会为整个群体带来回报。或者，若没有这样的人，即便如此，你的赏报也是大的；因为你的主不仅与义人来往而避开不洁的人，而是仁慈地接待了客纳罕妇人和撒玛黎雅妇人——那可憎而污秽的；另外还有一个妓女，犹太人为此责难他，他却接待并医治了她，并容许自己被那污秽之人的眼泪洗脚，教导我们要俯就那些在罪中的人，因为这是最大的仁爱。你说什么？强盗和盗墓者住在监狱里？那么请告诉我，住在城里的所有人都是义人吗？难道不是有许多甚至比他们更坏，以更大的无耻进行抢劫吗？因为前者，若没有其他借口，至少以独处和黑暗为遮羞布，并暗中行事；但后者却撕下面具，明目张胆地行恶，粗暴、贪婪、贪得无厌。很难找到一个远离不义的人。\par

6. 若我们不抢夺金子或若干亩土地，但在较小的事情上，我们却能以欺骗和抢劫达到同样的目的，并在我们能做到的地方达至同样结果。因为当我们立约，或买或卖时，我们争论并尽力压低价格，竭尽全力使这成为现实，这行动难道不是抢劫？难道不是偷窃和贪得无厌？不要对我说你没有夺走房屋或奴隶，因为不义不是以夺取之物的多少来判断，而是以抢劫者的意图来判断。因为\Quote{义}和\Quote{不义}在大事和小事上具有同样的效力；我称那割开钱包取金子的为扒手，也把那从市场上的人那里购买东西时压低价格、扣除应得部分的人称为扒手；那打破墙壁偷窃东西的是入室盗贼，但那种败坏正义并从邻人处取走任何东西的人，同样是入室盗贼。\par

那么，我们不要忽略自己的过犯，而去判断他人的；也不要在应当施行仁爱的时候，去追查别人的邪恶；反要思想我们自己从前是怎样的，如今就存温柔和恩慈。我们从前是怎样的呢？请听\Person{保禄}说：\BibleQuote{因为我们从前也是无知、悖逆、受迷惑，服侍各样的私欲和宴乐，常存恶毒和嫉妒之心，彼此可憎可恨。} \BibleRef{铎 3:3} 又说：\BibleQuote{我们生来就是义怒之子。} \BibleRef{弗 2:3} 但是，天主看见我们仿佛被囚禁在监牢里，被沉重的锁链捆绑，比铁链更沉重，祂并没有以我们为耻，反而前来进入监牢，尽管我们该受千万种惩罚，祂却将我们从那里领出，带我们进入王国，使我们比诸天更光荣，为使我们也能按自己的力量照样去行。因为当祂对门徒说：\BibleQuote{我，你们的主和师傅，既然洗了你们的脚，你们也当彼此洗脚；因为我给你们立了榜样，叫你们照着我向你们所做的去做。} \BibleRef{若 13:14} 祂立此法不仅是为洗脚，也是为祂向我们显示的一切其他行为。住在监牢里的难道是个杀人犯吗？然而，我们行善给祂不可懈怠。难道是个盗墓者或奸淫者吗？我们当怜悯他的不幸，而不是他的邪恶。但如前所述，那里常会发现一个比万人更有价值的人；如果你不断前往探视囚犯，就不会错失如此伟大的奖赏。因为正如\Person{亚巴郎}因接待普通客人而曾遇见天使，同样，我们若将此作为常业，也会遇见伟大的人物。容我说句奇怪的话：接待伟人的人，不如接待困苦可怜的人值得赞美。因为前者在自己生活中不乏受到善待的机会，但后者被众人唾弃和抛弃，只有一处港湾——恩人的怜悯；所以这尤其是纯粹的仁爱。此外，那关注备受敬仰和显赫之人的人，往往是为了在人前炫耀；而那照料卑微绝望之人的人，则只因天主的命令而行。所以，我们若设宴，就奉命款待瘸腿和残疾的；若行慈悲之事，就奉命向最卑微和最小的人施行。因为经上说：\BibleQuote{你们既做在我这弟兄中一个最小的身上，就是做在我身上了。} \BibleRef{玛 25:45} 因此，既然知道那里积存的宝藏，让我们不断进入，以此为业，并把我们对剧院的热情转向那里。如果没有什么可捐助的，就献上言语的安慰。因为天主不仅酬报那供食的，也酬报那进入的。当你进入，唤起那战兢恐惧的灵魂，劝勉、援助、应许帮助、教导其真智慧，你必从中获得不小的赏报。因为如果你在监狱外这样说话，许多人甚至会嘲笑，因为他们被过度的奢华所放纵；但那些在逆境中的人，心灵谦卑，会温柔地听从你的话，赞美它们，并成为更好的人。因为甚至当\Person{保禄}宣讲时，犹太人也常常讥笑他，但囚犯们却极其安静地聆听。因为没有什么比灾难、诱惑和苦难的压力更能使灵魂适合天上的智慧。想到这一切，以及我们若常与囚犯来往，将对他们和我们自己成就多少益处，让我们把从前花在市场和无益事务上的时间，用在监狱里；这样我们既能赢得他们，也能使自己喜乐，并因使天主得荣耀，而获得永恒的福乐，借着我们的主\Person{耶稣基督}的恩宠和仁爱。愿光荣归于圣父及圣神，从今世直到永远。阿们。\par

\SourceRecord{Translated by Charles Marriott.；Nicene and Post-Nicene Fathers, First Series；Vol. 14.；Edited by Philip Schaff.；Buffalo, NY: Christian Literature Publishing Co.,；1889.；<http://www.newadvent.org/fathers/240160.htm>.}

% structure: p0001 source=h1 target=section reason=page-title

\section{若望福音讲道集第六十一篇}

% structure: p0002 source=h2 target=subsection reason=relative-heading-rank; base=h2

\subsection{若望福音 10:22-24}

1. 每一种德行都是好事，但最美好的乃是温和与良善。这使我们成为人，这使我们与野兽有别，这使我们能与天使媲美。因此，基督不断地花费许多言语谈论这一德行，命令我们要温和、良善。祂不仅用言语教诲，也以身作则：一时受掌掴而忍受，一时受辱骂和谋害，却又来到那些谋害祂的人那里。因为那些曾称祂为附魔者、撒玛黎雅人，并多次想杀死祂、向祂投石的人，如今却围住祂，问祂说：\Quote{祢是基督吗？}即便如此，在如此众多且巨大的谋害之后，祂并没有拒绝他们，而是以极大的温和回答了他们。\par

但我们有必要从头考察整个段落。\par

圣经说：\Quote{那时在耶路撒冷举行重建节，正是冬天。}这一节日是盛大而民族的。因为他们以极大的热忱庆祝圣殿在从波斯长期被掳归来后重建的日子。基督也参加了这个节日，因为从此以后祂一直住在犹大地，因为受难近了。\Quote{犹太人围住祂，说：\InnerQuote{祢使我们的心悬疑不定，要到几时呢？}}\par

\Quote{祢若是基督，就明明地告诉我们。}\par

祂并没有回答：\Quote{你们问我什么呢？你们多次称我为附魔者、狂人、撒玛黎雅人，视我为天主的仇敌和骗子，你们刚才还说：\InnerQuote{祢为自己作证，祢的见证不真实}；那么你们又怎能来寻求并想从我这里学习，而你们却拒绝我的见证呢？}但祂没有说任何这类话，尽管祂知道他们发问的意图是恶意的。因为他们围住祂说：\Quote{祢使我们的心悬疑不定，要到几时呢？}这似乎出于某种渴望和学习的意愿，但他们发问的意图却是败坏和诡诈的。因为既然祂的作为不容他们诽谤和侮辱，而他们却可以通过在祂的话语中找出另外的含义来攻击祂的言论，所以他们不断地提出问题，想通过祂的言论来使祂沉默；当他们在祂的作为上找不到缺点时，就希望在祂的话语中找到把柄。因此他们说：\Quote{告诉我们}；然而祂已经多次告诉过他们。因为祂对撒玛黎雅妇人说：\Quote{同你谈话的我就是。}（\BibleRef{若 4:26}）对瞎子说：\Quote{你已看见祂，同你谈话的就是祂。}（\BibleRef{若 9:37}）祂也曾告诉过他们，即使不是同样的措辞，至少也用其他话语。事实上，如果他们明智，并且渴望正确询问，他们本应通过言语承认祂，因为通过作为祂已经多次证明了所讨论的事。但现在请看他们乖谬和好辩的性情。当祂讲话并以言语教导他们时，他们说：\Quote{祢显什么神迹给我们看呢？}（\BibleRef{若 6:30}）但当祂通过作为给他们证明时，他们对祂说：\Quote{祢是基督吗？明明地告诉我们}；当作为大声宣告时，他们寻求话语；而当话语教导时，他们又转向作为，总是自相矛盾。但他们并非为学习而询问，结局证明了这一点。因为他们认为祂如此值得信赖，以至于接受祂对自己的见证，而当祂说了几句话后，他们立刻拿起石头要砸祂；可见他们围住并逼迫祂是出于恶意。\par

提问的方式也充满了极大的仇恨。\Quote{明明地告诉我们：祢是基督吗？}然而祂所有的话都是公开说的，常在他们的节日中出现，没有在暗中说什么；但他们提出欺诈的话：\Quote{祢使我们的心悬疑不定，要到几时呢？}为了引诱祂出来，然后再找到把柄攻击祂。因为每次他们提问都不是为了学习，而是为了挑剔祂的话，这不仅从这一段，也许多其他段落可以看出来。因为他们来问祂：\Quote{纳税给凯撒，可以不可以？}（\BibleRef{玛 22:17}），他们谈论休妻的事（\BibleRef{玛 19:3}），他们询问那个据说有过七个丈夫的妇人（\BibleRef{玛 22:23}），都被证实他们提出问题不是出于学习的愿望，而是出于恶意。但那里祂斥责他们说：\Quote{你们为什么试探我，假善人？}表明祂知道他们隐秘的心思，而在这里祂没有说这类话；教导我们不要总是斥责那些谋害我们的人，而是要以温和良善忍受许多事。\par

既然当作为大声宣告时，寻求话语的见证是愚昧的，请听祂如何回答他们，同时暗示他们这些询问是多余的，不是为了学习，并且表明祂发出比言语更清晰的声音，即通过作为。\par

% structure: p0010 source=h2 target=subsection reason=relative-heading-rank; base=h2

\subsection{若望福音 10:25}

2. 这是他们中比较温和的人常对彼此说的话：\Quote{一个罪人不能行这样的奇迹。}又说：\Quote{魔鬼不能开瞎子的眼睛。}以及：\Quote{除非天主与他同在，没有人能行这样的奇迹。}（\BibleRef{若 3:2}）看到祂所行的奇迹，他们说：\Quote{这不是基督吧？}另有人说：\Quote{基督来的时候，祂所行的奇迹，难道会比这人所行的更多吗？}（\BibleRef{若 7:31}）而那些当时愿意信祂的人说：\Quote{祢显什么神迹给我们看，好叫我们看见而信祢呢？}（\BibleRef{若 6:30}）当那些没有被如此伟大作为说服的人，假装会被一句空话说服时，祂斥责他们的邪恶，说：\Quote{如果你们不信我的作为，怎能信我的话语呢？所以你们的询问是多余的。}\par

% structure: p0012 source=h2 target=subsection reason=relative-heading-rank; base=h2

\subsection{若望福音 10:26}

\Quote{因为就我而言，我已经完成了一位牧者所应做的一切；如果你们不跟随我，不是因为我不是牧者，而是因为你们不是我的羊。}\par

% structure: p0014 source=h2 target=subsection reason=relative-heading-rank; base=h2

\subsection{若望福音 10:27-30}

请注意祂在拒绝时如何激励他们跟随祂。\Quote{你们不听我，}祂说，\Quote{因为你们不是羊；但那跟随的，才是属于羊群的。}祂这样说，是为使他们努力成为羊。然后，通过提及他们将获得什么，祂使这些人产生嫉妒，从而激励他们，并使他们渴望这样的事物。\par

\Quote{那么，是因为父的能力没有人能夺去他们，而你没有力量，太软弱以致不能保护他们吗？}绝不。为了让你明白\Quote{父赐给我}这句话是为他们的缘故而说的，免得他们再次称祂为天主的仇敌，因此，在断言\Quote{没有人能从我的手中夺去他们}之后，祂继续表明祂的手和父的手是同一的。因为如果不是这样，祂很自然会说：\Quote{赐给我羊的父比万有都大，没有人能从我的手中夺去他们。}但祂没有这样说，而是说：\Quote{从我父的手中。}然后，免得你以为祂确实软弱，而羊因父的能力才安全，祂补充说：\Quote{我与父原是一体。}好像祂说：\Quote{我并非因为父的缘故而说没有人能夺去他们，好像我太软弱不能保护羊。因为我和父是一体。}这里祂指的是能力，因为祂的全部讲论都是关于这个；如果能力相同，显然本质也是相同的。当犹太人用各种手段，图谋并把人赶出会堂时，祂告诉他们，他们的一切计谋都是无用和枉然的；\Quote{因为羊在我父的手中}；正如先知所说：\Quote{在我手上我刻画了你的墙壁。}\BibleRef{依 49:16}然后，为了表明手是同一的，祂有时说那是祂自己的，有时说是父的。但当你听到\Quote{手}这个词时，不要理解为任何物质性的东西，而是能力、权威。再者，如果是因为父赐给祂能力，所以没有人能夺去羊，那么后面所说的\Quote{我与父原是一体}就是多余的。因为如果祂次于父，这将是一个非常大胆的说法，因为它只宣告了能力的平等；犹太人意识到了这一点，就拿起石头要砸祂。\BibleRef{若 10:31}即使如此，祂也没有消除这个意见和猜疑；不过，如果他们的猜疑是错误的，祂应该纠正他们，并说：\Quote{你们为什么做这些事？我这样说并不是要证明我的能力和父的相等}；但祂却做了正好相反的事，证实了他们的猜疑，并加以强调，而且是在他们激怒的时候。因为祂没有为所说的话作任何辩解，好像那些话说错了，反而责备他们没有对祂持有正确的意见。因为他们说：\par

% structure: p0017 source=h2 target=subsection reason=relative-heading-rank; base=h2

\subsection{\BibleBook{若望}{福音} 10:33\-36}

祂所说的是这样的：\Quote{如果那些因恩宠而领受这尊荣的人，称自己为神并不被指责，那么那本性拥有这尊荣的，如何配受责备呢？}然而祂没有这样说，而是在稍后才加以证明，先是在祂的讲论中稍微放松和让步，并说：\Quote{父所祝圣并派遣来的。}当祂缓和了他们的怒气后，祂便提出明确的断言。为了让祂的话被接受，祂暂时用较低微的语气说话，但随后祂提高了声音，说：\par

% structure: p0019 source=h2 target=subsection reason=relative-heading-rank; base=h2

\subsection{\BibleBook{若望}{福音} 10:37\-38}

你看到祂如何证明我所说的：祂在什么方面也不次于父，而是在各方面与祂相等吗？因为既然无法看见祂的本质，祂便从工作的相等和相同来提供能力不变的证明。那么，请告诉我，我们应该相信什么？\par

3. \Quote{我在父内，父也在我内。}\par

\Quote{因为我无非是父所是的，但仍然是子；祂无非是我所是的，但仍然是父。如果有人认识我，他就认识父；如果人认识父，他也认识了子。}如果能力是次等的，那么关于认识的说法也是假的，因为不可能通过另一种本质或能力来认识一种本质或能力。\par

% structure: p0023 source=h2 target=subsection reason=relative-heading-rank; base=h2

\subsection{\BibleBook{若望}{福音} 10:39\-41}

当祂说了任何伟大崇高的话后，祂很快便退下，向他们的怒气让步，好让愤怒因祂的离开而减退和止息。那时祂就是这样做的。但为什么福音作者提到那个地方呢？为的是让你知道，祂去那里是为了提醒他们若翰在那里所做所说的事，以及他的证词；至少当他们到达那里时，他们立刻记起了若翰。因此他们也说：\Quote{若翰确实没有行过神迹}，因为除非那地方使他们想起洗者若翰，并且他们记起他的证词，他们怎么会加上这句话呢？注意他们如何形成无可辩驳的三段论。\Quote{若翰确实没有行过神迹}，但\Quote{这人却行神迹}，有人说；\Quote{因此可见祂的优越性。如果人们相信那没有行神迹的人，那么他们更应该相信这人。}然后，因为是若翰作了证，免得他没有行神迹似乎证明他不配作证人，他们又加上：\Quote{但他若没有行神迹，他仍关于这人说了所有真话}；不再藉着若翰证明基督可信，而是藉着基督所行的事证明若翰可信。\par

% structure: p0025 source=h2 target=subsection reason=relative-heading-rank; base=h2

\subsection{\BibleBook{若望}{福音} 10:42}

有许多事吸引他们。他们记得若翰所说的话，称基督为\Quote{比他更强}、\Quote{光}、\Quote{生命}、\Quote{真理}等等。他们记得从天上来的声音，以及以鸽子形象显现的圣神，将祂指示给众人；此外，他们还回想起奇迹所提供的明证，注目于这些，他们今后就坚定了。\Quote{因为，}有人说，\Quote{如果我们应该相信若翰，就更应该相信这人；如果相信没有行奇迹的他，就更应该相信这人，他除了若翰的见证，还有来自奇迹的凭据。}你们看到吗？住在此处、摆脱恶人的同在，使他们获益如此之大？因此耶稣不断地引导并使他们远离那些人的陪伴；正如祂在旧盟约下似乎也这样做过，在旷野、远离埃及人之处，全面地塑造和管理犹太人。\par

祂现在也劝告我们这样做，吩咐我们避开公共场所、喧哗和扰乱，在内室平安地祈祷。因为免于混乱的船顺风航行，与世俗事务分离的灵魂安息在港湾。因此，妇女应该比男人更有真智慧，因为她们多半被束缚在家中。例如，雅各伯是个纯朴的人，因为他住在家里，免于公共生活的喧嚣；因为圣经如此记载，它说：\Quote{住在家中。}\BibleRef{创 25:27}\Quote{但是，}某个妇女说，\Quote{即使在家中也有很大的混乱。}是的，当你愿意这样时，你给自己带来一大堆烦恼。因为那在市场和法庭间消磨时间的男人，被外在的烦恼如同波浪般淹没；但那坐在家中如同在真智慧学校里的妇女，在内里收敛心神，将能致力于祈祷、读经和其他天上的智慧。如同那些住在旷野的人无人打扰，她持续在内里也能享受持久的平静。即使她有时需要外出，也绝无混乱的理由。因为妇女离家外出的必要场合，或是为了到这里来，或是身体需要在浴池中洁净；但大部分时间她坐在家中，既能自己真正聪明，又能在丈夫激动时接待他、安抚他、使他镇定，减轻他思想的过度和激烈，这样再送他出去，他已卸下从市场收集的一切恶习，并带着在家中学习到的一切善。因为没有什么，没有什么比一位虔诚而明智的妇女更有力量，能使男人归于正轨，并随心塑造他的灵魂。因为他不会像容忍他的伴侣劝告他那样容忍朋友、师长或统治者，因为这劝告甚至带有某种愉悦，因为给出劝告的她深受爱戴。我可以讲述许多刚硬不顺从的男人就这样被软化了。因为那与他同席、同床、同拥抱、同言语、同秘密、同出入以及许多其他事物的女子，她完全交付于他，与他结合，如同身体与头相连，如果她恰巧谨慎而和谐，她将在管理丈夫上超越并胜过所有其他人。\par

4. 因此，我劝勉妇女以此为己任，并给予合宜的劝告。因为她们既大有能力行善，也同样有能力行恶。一个女人毁灭了\Person{阿贝沙隆}，一个女人毁灭了\Person{阿默农}，一个女人几乎毁灭了\Person{约伯}，一个女人从屠杀中救出了\Person{纳巴耳}。妇女保全了整个民族；因为\Person{德波辣}和\Person{友弟德}展现了堪与男子媲美的成功；成千上万的妇女也是如此。因此\Person{保禄}说：\BibleQuote{因为妻子，你哪里知道是否能救你的丈夫？} \BibleRef{格前 7:16} 在那个时代，我们看到\Person{培黎息}、\Person{玛利亚}和\Person{普黎史拉}参与宗徒们的劳苦\EditorialNote{罗马书第十六章}；我们也必须效法她们，不仅用言语，更用行动，规正与我们同住的人。但我们如何用行动教导他呢？当他看到你没有不良倾向，不贪图花费或装饰，不要求过度的金钱供给，而满足于你所拥有的，那时他才会容忍你劝告他。但如果你言语明智，而行动却相反，他会指责你为极其愚蠢的谈论。但当你不仅用言语，也以你的行为给他教导时，他就会更容易接纳你并听从你；正如你不贪求黄金、珍珠或昂贵的衣服，而代之以谦逊、节制、温和；当你自己表现出这些德行，并在他身上要求它们。因为如果你必须做些什么来取悦你的丈夫，你应该装饰你的灵魂，而不是装饰从而糟蹋你的身体。你戴在身上的黄金不会使你对他如此可爱和令人渴望，如同谦逊和善待他，以及为你的伴侣赴死的准备；这些事最能征服男人。事实上，衣着的华丽甚至会惹他不悦，因为它束缚他的财力，给他带来许多开销和忧虑；但我所提到的那些事却会使丈夫紧紧依恋妻子；因为善良、友谊和爱不引起忧虑，不产生开销，反而恰恰相反。外在的装饰因使用而变得乏味，但灵魂的装饰却日益绽放，燃起更强烈的火焰。所以，如果你愿意取悦你的丈夫，就用谦逊、虔诚和家务管理来装饰你的灵魂。这些事既更能征服他，也永不消失。年龄不摧毁这装饰，疾病不损耗它。身体的装饰常被时间破坏，被疾病和许多其他事物损耗，但涉及灵魂的则超越这一切。那种装饰引起嫉妒，燃起猜疑，但这种装饰纯洁无病，远离一切虚荣。这样，家事就会更容易，你的收入也没有麻烦，当黄金不是戴在你的身上或缠绕你的手臂，而是用于必要的用途，如仆人的饮食、子女的必要照顾和其他有益的目的。但如果不是这样，如果（妻子的）脸上挂满装饰，而（丈夫的）心里充满焦虑，有什么益处，有什么好处呢？一方忧伤，就不让另一方奇妙的美被看见。因为你知道，你知道，即使一个人看见所有女人中最美的，当他的灵魂忧伤时，他也不能在看见时感到愉悦，因为为了感到愉悦，一个人必须首先欢喜和高兴。并且当他所有的黄金都堆积起来装饰一个女人的身体，而他的住所中却有痛苦时，她的伴侣就不能有快乐。所以，如果我们希望取悦我们的丈夫，就让我们给他们快乐；我们若除去我们的装饰和华丽服饰，就会给他们快乐。因为所有这一切在结婚之时看似提供一些快乐，但随后就随时间消褪了。因为如果当天空如此美丽，太阳——你无法说出任何物体与之相比——如此明亮，我们却因习惯而较少欣赏它们，我们又怎能欣赏一个用花哨饰品装扮的身体呢？我说这些，是希望你们用\Person{保禄}所吩咐的健康装饰来装饰自己；\BibleQuote{不以黄金、珍珠或昂贵的服装，而以善行——这合乎自称虔敬的妇女。} \BibleRef{弟前 2:9} 但你是否希望取悦陌生人，得他们的赞美？那么，这肯定不是一个谦逊妇女的愿望。然而，如果你希望如此，通过按我所说的去做，你也会使陌生人非常爱你，并赞美你的谦逊。因为对于装饰自己身体的妇女，没有有德和清醒的人会赞美，只有放纵和淫荡的人；不，甚至这些人也不会赞美她，反而会说她的坏话，因为他们被围绕她的放荡所煽动；但另一种妇女，所有的人都赞同，无论是这一类人还是那一类人，因为他们从她那里没有受到伤害，反而受到天上智慧的教导。从人那里，她将得到极大的赞美，并在天主那里得到极大的赏报。那么，让我们努力追求这样的装饰，使我们在此能够无忧无虑地生活，并获得未来的福乐；愿我们众人赖我们的主耶稣基督的恩宠和慈爱获得这一切，愿光荣归于他，直到永远。阿们。\par

\SourceRecord{Translated by Charles Marriott.；Nicene and Post-Nicene Fathers, First Series；Vol. 14.；Edited by Philip Schaff.；Buffalo, NY: Christian Literature Publishing Co.,；1889.；<http://www.newadvent.org/fathers/240161.htm>.}

% structure: p0001 source=h1 target=section reason=page-title

\section{《若望福音》讲道集第六十二篇}

% structure: p0002 source=h2 target=subsection reason=relative-heading-rank; base=h2

\subsection{若望福音 11:1-2}

1. 许多人看见那些中悦天主的人遭受任何可怕的事，例如陷入疾病、贫穷或其他类似的苦难，就会感到困惑，不知道承受这些事正是特别亲近天主的人所应有的；因为\Person{拉匝禄}也是基督的朋友之一，并且患病了。至少那些派去的人说：\Quote{请看，你所爱的人病了。}但让我们从头来思考这段经文。经上说：\Quote{有一个人，}经上说，\Quote{病了，就是\Place{伯达尼}的\Person{拉匝禄}。}作者并非无缘无故地提到\Person{拉匝禄}是哪里人，而是出于一个他日后会告诉我们的理由。现在让我们专注于眼前的经文。他也为了我们的益处告诉我们\Person{拉匝禄}的姐妹是谁，并且进一步指出\Person{玛利亚}比另一位更特别，接着说：\Quote{就是那用香液抹了主的\Person{玛利亚}。}这里有些人疑惑说：\Quote{主怎能容忍一个女人做这事呢？}首先必须明白，这不是\BibleBook{玛窦}{福音}中所提到的那位妓女（\BibleRef{玛 26:7}），或\BibleBook{路加}{福音}中的那位（\BibleRef{路 7:37}），而是另一个人；那些妓女充满许多恶习，但她却是庄重而真诚的；因为她对款待基督表现出了她的热心。福音作者也想表明，姐妹二人也爱祂，但祂仍允许\Person{拉匝禄}死去。但他们为什么不像百夫长和贵族那样，留下患病的兄弟，自己来见基督，而是派人去呢？他们非常信赖基督，对祂有深厚的亲切感。此外，她们是软弱的妇女，被忧愁压垮；后来她们表明，她们这样做并不是轻视祂。显然，这位\Person{玛利亚}不是那个妓女。有人问：\Quote{但是为什么，}有人问，\Quote{基督接纳那个妓女呢？}那是为了除掉她的罪孽，为了显示祂的慈爱，为了让你知道没有任何疾病能胜过祂的良善。所以不要只看祂接纳了她这一点，也要考虑祂如何改变了她。但（回到正题），为什么福音作者向我们叙述这段历史呢？或者说，他通过这句话想向我们表明什么呢？\par

% structure: p0004 source=h2 target=subsection reason=relative-heading-rank; base=h2

\subsection{若望福音 11:5}

就是叫我们永远不要因为好人，即那些亲近天主的人患病而感到不满或烦恼。\par

% structure: p0006 source=h2 target=subsection reason=relative-heading-rank; base=h2

\subsection{若望福音 11:3}

他们想要激发基督的怜悯，因为他们仍然把祂当作一个人来对待。这从他们的话中可以明显看出：\Quote{你若在这里，他就不会死，}以及他们不说\Quote{请看，\Person{拉匝禄}病了，}而说\Quote{请看，你所爱的人病了。}那么基督说了什么呢？\par

% structure: p0008 source=h2 target=subsection reason=relative-heading-rank; base=h2

\subsection{若望福音 11:4}

请注意祂如何再次断言祂的荣耀与圣父的荣耀是同一的；因为祂先说了\Quote{属于天主，}然后加上\Quote{为叫圣子受到光荣。}\par

\Quote{这病不至于死。}因为祂打算在原地再停留两天，所以暂时用这个回答打发走了报信的人。因此我们必须佩服\Person{拉匝禄}的姐妹，她们听说这病\Quote{不至于死，}却看见他死了，尽管结果完全相反，她们却没有感到困惑。但即便如此，她们还是来到祂面前，并不认为祂说了谎话。\par

这段经文中的\Quote{那}一词所指的不是原因，而是结果；这病由其他原因引起，但祂却用它来为天主的光荣。\par

% structure: p0012 source=h2 target=subsection reason=relative-heading-rank; base=h2

\subsection{若望福音 11:6}

祂为什么停留？是为了让\Person{拉匝禄}咽下最后一口气，并被埋葬；为使人无法断言祂是在他还未死时恢复了他，声称这只是昏睡、晕厥、发作，而不是死亡。正因如此，祂停留了那么久，直到腐烂开始，他们便说：\Quote{他现在已经臭了。}\par

% structure: p0014 source=h2 target=subsection reason=relative-heading-rank; base=h2

\subsection{若望福音 11:7}

为什么，祂在其他地方从未预先告诉门徒祂要去哪里，却在这里告诉他们？他们当时非常害怕，既然他们处于这样的状态，祂就预先警告他们，免得突然的事使他们不安。那么门徒们说什么呢？\par

% structure: p0016 source=h2 target=subsection reason=relative-heading-rank; base=h2

\subsection{若望福音 11:8}

所以他们也为祂害怕，但多半是为自己害怕；因为他们尚未成全。于是\Person{多默}恐惧战栗地说：\Quote{我们也去，好同祂一起死。}\BibleRef{若 11:16}，因为\Person{多默}比其他人更软弱、更不信。但请看耶稣如何通过祂的话鼓励他们。\par

% structure: p0018 source=h2 target=subsection reason=relative-heading-rank; base=h2

\subsection{若望福音 11:9}

祂要么这样说：\Quote{凡自己意识到没有邪恶的人，必不会受任何可怕的事；只有作恶的才会受苦，所以我们不必害怕，因为我们没有做任何该死的事；}或者这样：\Quote{凡\InnerQuote{看见这世界的光}的人，是安全的；如果看见这世界之光是安全的，那么那与我同在的，若不离开我，更是安全。}用这些话鼓励他们之后，祂接着说，他们去那里的原因很急迫，并向他们表明他们不是要去耶路撒冷，而是去\Place{伯达尼}。\par

% structure: p0020 source=h2 target=subsection reason=relative-heading-rank; base=h2

\subsection{若望福音 11:11-12}

也就是说，\Quote{我去不是为了像以前那样与犹太人辩论争执，而是为了唤醒我们的朋友。}\par

% structure: p0022 source=h2 target=subsection reason=relative-heading-rank; base=h2

\subsection{若望福音 11:12}

他们这样说并非没有理由，而是想阻止去那里。其中一人问：\Quote{你说他睡着了？那就没有急迫的理由去了。}然而正因如此，祂先前说：\Quote{我们的朋友，}以表明去那里是必要的。因此，当他们的态度有些勉强时，祂说：\par

% structure: p0024 source=h2 target=subsection reason=relative-heading-rank; base=h2

\subsection{若望福音 11:14}

2. 祂先前的话，是想要证明祂不爱自夸；但由于他们不明白，祂补充说：\Quote{他死了。}\par

% structure: p0026 source=h2 target=subsection reason=relative-heading-rank; base=h2

\subsection{若望福音 11:15}

为什么\Quote{为了你们}？\Quote{因为我预先告诉了你们他的死亡，当时我不在那里，并且当我再次复活他时，就不会有欺骗的嫌疑。}你们看见门徒们的心性还不完美，对祂的权能认识得不够充分吗？这是由于插入的恐怖事物扰乱了他们的灵魂所造成的。当祂说\Quote{他睡了}时，祂又加上\Quote{我去唤醒他}；但当祂说\Quote{他死了}时，祂没有加上\Quote{我去复活他}；因为祂不愿以言语预言祂要以行动确实成就的事，处处教导我们不要虚荣，也不要在没有理由的情况下许诺。如果祂对待百夫长的召请时也是如此（因为祂说\BibleQuote{我去治好他}，\BibleRef{玛 8:7}），那是为显示百夫长的信德才这样说。如果有人问：\Quote{门徒们怎么会想到睡觉？他们怎么会不明白从祂的话说\InnerQuote{我去唤醒他？}是指死亡？因为他们如果指望祂走十五斯塔狄亚去唤醒他，那是愚蠢的。}我们会回答，他们认为这是一个隐晦的说法，如同祂经常对他们说的那样。\par

现在他们都害怕犹太人的攻击，但\Person{多默}超过其他人；因此他也说，\par

% structure: p0029 source=h2 target=subsection reason=relative-heading-rank; base=h2

\subsection{若望福音 11:16}

有人说他（指\Person{多默}）自己愿意死；但事实并非如此；这句话更像是胆怯的表现。然而他并没有受到责备，因为基督还支持他的软弱，但后来他变得比所有人都坚强，不可战胜。因为奇妙之处在于：我们看到一个在受难前如此软弱的人，在受难后，并在他相信复活后，变得比任何人都热忱。基督的权能何其伟大！那个不敢与基督一起前往\Place{伯达尼}的人，正是他在没有见到基督的情况下几乎跑遍了有人居住的世界，居住在充满谋杀、并想要杀害他的民族中间。\par

但是如果\Place{伯达尼}相距\Quote{十五弗隆}（即两英里），\Person{拉匝禄}怎么会\Quote{死了四天}？\Person{耶稣}停留了两天，在那两天之前的一天有人来报信（正是在同一天\Person{拉匝禄}死了），然后在第四天祂抵达了。祂等待被召请，因此未受邀请而来，以免有人怀疑所发生的事；那些祂所爱的妇女也没有亲自来，而是派了别人来。\par

% structure: p0032 source=h2 target=subsection reason=relative-heading-rank; base=h2

\subsection{若望福音 11:18}

他并非无故提及此事，而是想告诉我们伯达尼离得近，因此很可能有许多人在那里。所以他宣告这一点后接着说，\par

% structure: p0034 source=h2 target=subsection reason=relative-heading-rank; base=h2

\subsection{若望福音 11:19}

但是他们既然约定，若有人承认基督，就要被逐出会堂，又怎能安慰基督所爱的妇女呢？或是因为灾祸的严重性，或是他们敬重她们出身高贵，又或者这些来的人并非恶类，他们中至少有许多人相信了。福音作者提到这些情况，是为了证明\Person{拉匝禄}确实死了。\par

3. 但是，为什么（\Person{玛尔大}）去见基督时，没有带她的妹妹一起？她希望私下遇见他，告诉他发生的事。但当他使她怀有美好的希望时，她就去叫\Person{玛利亚}，\Person{玛利亚}在悲伤最甚时来见他。你看到她的爱是多么炽热吗？这就是他曾说\Quote{玛利亚选择了更好的一份}的那位\Person{玛利亚}。\BibleRef{路 10:42} \Quote{那么，}有人说，\Quote{\Person{玛尔大}怎么显得更热忱？}她并非更热忱，而是因为另一位尚未得知，因为\Person{玛尔大}是较软弱的一个。她即使从基督那里听到了这样的事，却仍然用卑下的语气说：\Quote{现在他必发臭，因为他已死了四天。} \BibleRef{若 11:39} 但是\Person{玛利亚}，虽然她什么也没听到，却说出这样的？而是立刻相信了，说：\par

% structure: p0037 source=h2 target=subsection reason=relative-heading-rank; base=h2

\subsection{若望福音 11:21}

请看这些妇女的天上智慧有多么伟大，尽管她们的领悟力薄弱。因为她们见到基督时，没有像我们在悲伤中见到熟人时那样，哀哭号啕，而是立即尊敬她们的老师。所以这两位姊妹都相信基督，但不是以正确的方式；因为他们还不确定知道祂是天主，也不知道祂凭自己的能力和权柄做这些事；在这两点上，祂都教导了她们。因为他们对前者的无知表现在说：\Quote{你若在这里，我们的兄弟就不会死；}对后者的无知表现在说：\par

% structure: p0039 source=h2 target=subsection reason=relative-heading-rank; base=h2

\subsection{若望福音 11:22}

好像他们是在谈论某个有德行的、被认可的人。但请看基督说什么；\par

% structure: p0041 source=h2 target=subsection reason=relative-heading-rank; base=h2

\subsection{若望福音 11:23}

祂至此驳斥了先前的话：\Quote{凡你求的}；因为祂没有说\Quote{我求}，而是说什么？\Quote{你的兄弟必要复活。}如果说\Quote{女人，你仍往下看，我不需要别人的帮助，而是自己做一切，}就会使她悲伤，成为她路上的绊脚石；但说\Quote{他必要复活}，则是选择了一种折中的说话方式。并且借着随后的话，祂暗示了我所提到的要点；因为当\Person{玛尔大}说：\par

% structure: p0043 source=h2 target=subsection reason=relative-heading-rank; base=h2

\subsection{若望福音 11:24}

% structure: p0044 source=h2 target=subsection reason=relative-heading-rank; base=h2

\subsection{若望福音 11:25}

表明祂不需要别人帮助，因为祂自己就是生命；因为如果祂需要别人，祂怎能是\Quote{复活与生命}？然而祂没有明确说明，只是暗示。但当她又说\Quote{凡你求的}时，祂回答：\par

\Quote{信我的人，即使死了，仍要活着。}\par

表明祂是美善的赐予者，我们必须向祂求。\par

% structure: p0048 source=h2 target=subsection reason=relative-heading-rank; base=h2

\subsection{若望福音 11:26}

注意看祂如何引导她的心神向上；因为所求的不仅是复活拉匝禄，更必须使她和与她同在的人都学习复活。因此，在复活死者之前，祂先用言语教导天上的智慧。但若祂是\Quote{复活}和\Quote{生命}，便不受地方限制，而是处处临在，知道如何医治。所以，如果他们像百夫长那样说：\Quote{祢只要说一句话，我的仆人就会痊愈}\BibleRef{玛 8:8}，祂也会这样做；但他们邀请祂到他们那里，恳求祂前来，祂便俯就，为将他们从对祂形成的卑微看法中提升起来，并来到那地方。然而在俯就的同时，祂也显示即使不在场，也有医治的能力。为此，祂也延迟，因为若非尸体先有气味，慈悲便不会像赐予时那样明显。但这妇人如何知道将有复活？他们曾听到基督多次谈论复活，但她现在仍渴望看见祂。且看她仍如何囿于下层；因为听见了\Quote{我是复活和生命}，她甚至也没有说\Quote{请使他复活}，但，\par

% structure: p0050 source=h2 target=subsection reason=relative-heading-rank; base=h2

\subsection{若 11:27}

基督的回答是什么？\Quote{凡信我的，即使死了，仍要活着}（此处是指人人共有的死亡。）\Quote{凡活着而信我的，必永远不死}\BibleRef{若 11:26}，意指那另一种死亡。\Quote{既然我是复活和生命，你勿忧愁，即使你的兄弟已经死了，但要相信，因为这不是死亡。}祂暂且安慰她所发生的事，并给她希望的曙光，说\Quote{他必要复活}，以及\Quote{我是复活}；并且复活后，即使再死，也不会受害，因此无需害怕这死亡。祂所说的是这样的：\Quote{这人没有死，你也不会死。} \Quote{你信这话吗？}她说：\Quote{我相信你是基督、天主子。}\par

\Quote{那要来到世界上的。}\par

这妇人似乎并不明白这话；她意识到这是某种大事，但并未领会全部含义，因此被问及一事时，却回答另一事。不过至少暂时，她得到了这一点好处：她缓和了悲伤；基督的话语就有这样的力量。为此，玛尔大先出来，玛利亚随后。因为她们对师傅的爱慕不允许她们强烈感受到眼前的忧伤；所以这些妇人的心灵确实是既明智又充满爱心的。\par

4. 但在我们这个时代，在其它罪恶之中，有一种弊端在我们妇女中非常普遍：她们在哀悼和哭号中大做文章，裸露手臂，撕扯头发，在脸颊上划出沟痕。她们这样做，有些出于悲伤，有些出于炫耀和攀比，有些出于放荡；她们裸露手臂，而且是在男人面前。女人啊，你为何如此？你以不雅的方式暴露自己，告诉我，你这基督的肢体，在集市中间，当男人在场时？你拔自己的头发，撕裂衣服，大声哭号，加入舞蹈，全程保持着与酒神节妇女相似的模样，难道你不觉得自己在冒犯天主吗？这是何等疯狂？外教人岂不嘲笑？他们岂不认为我们的教义是神话？他们会说：\Quote{没有复活——基督徒的教义是嘲弄、诡计和捏造。因为她们的妇女哀哭，仿佛死后一无所有；她们毫不留意铭刻在书卷上的话；所有那些话都是虚构，而这些妇女正表明它们确实是虚构。因为，倘若她们相信死了的人并非死了，而是迁往更好的生命，她们就不会哀悼他如同他不再存在，就不会这样击打自己，就不会说出这样充满不信的话：\InnerQuote{我再也见不到你了，我再也不能得回你了。}她们的整个宗教是个神话，如果连这至善之物都被她们如此全盘不信，那么她们所尊崇的其他事物就更不用说了。}外教人并非如此女子气，他们中许多人实践了天上的智慧；一位妇女听到她的儿子在战斗中阵亡，立刻问道：\Quote{城中的情况如何？}另一位真正睿智的人，在戴冠时听说他的儿子为国捐躯，便摘下花冠，问是哪个儿子；当他知道是哪一个后，随即又戴上花冠。许多人甚至献出他们的儿女为他们的邪神作祭品；斯巴达妇女鼓励她们的儿子要么从战场平安带回盾牌，要么就死在盾牌上被抬回来。因此，我惭愧的是，外教人在这些事上显出真智慧，而我们却行为不当。那些对复活一无所知的人，装作知道的样子；而那些知道的人，却装作不知道的样子。常常有许多人因为在人前害羞而做了那些他们本不因天主而做的事。因为上等阶层的妇女既不撕扯头发，也不裸露手臂；这件事本身正是对她们最严重的指控，不是因为她们不裸露身体，而是因为她们这样做不是出于虔诚，而是为了不被人认为有失体面。难道她们的羞耻心比悲伤更强，而对天主的敬畏却不更强吗？这难道不应受最严厉的谴责吗？富家妇女因财富而做的事，穷人也该因敬畏天主而去做；但如今恰恰相反：富人因虚荣而行得明智，穷人因心灵狭隘而行得不妥。还有什么比这种反常更糟糕呢？我们做一切事都是为人，都是为地上的事。这些人说出的话充满了疯狂和许多讥讽。主的确说过：\BibleQuote{哀恸的人是有福的} \BibleRef{玛 5:4}，这是指为自己的罪而哀恸的人；没有人哀哭那种哀恸，也没有人关心迷失的灵魂；而这另一种哀哭我们并未受命去做，我们却去做了。\Quote{什么？}有人问道，\Quote{人可能不哭泣吗？}不，我并非禁止哭泣，而是禁止击打自己、无节制的哭泣。我既非野蛮，亦非残忍。我知道我们的本性寻求并思念朋友和日常伴侣；它不能不悲伤。基督也如此表明，因为祂曾为拉匝禄哭泣。所以你也当如此；哭泣，但要温柔，要端庄，要怀着对天主的敬畏。如果你这样哭泣，你这样做并非不相信复活，而是因为不能忍受分离。因为即使对于那些离开我们、远赴他乡的人，我们也会哭泣，但我们这样做并非出于绝望。\par

5. 所以，你也要哀哭，好像送一个人到异地。我说这些话，不是定下行动规则，而是迁就（人性的软弱）。因为如果死者是个罪人，在许多事上得罪了天主，就该哀哭（更确切地说，不只是哀哭，因为对他毫无助益，而要尽力通过施舍和奉献为他谋求一些安慰；）但也该为此欢喜，即他的恶行被截断了。如果他是个义人，又该欢喜，因为他的所有现已处于安稳之中，免于将来的不确定性；如果是年轻人，因为他已迅速脱离人生常见的邪恶；如果是老年人，因为他已饱享那被认为可欲之物后离世。但你忽略这些事，却使唤你的婢女作哀悼者，仿佛你是在尊敬死者，其实这是极大的不敬。因为对死者的尊敬不是哀嚎和痛哭，而是赞美诗、圣咏和卓越的生活。善人离世时，虽有天使陪伴，即使无人靠近他的遗体；但恶人，即使全城为他送葬，也毫无益处。你愿意尊敬已逝的人吗？用另一种方式尊敬他：通过施舍、善行和公益服务。许多哀悼有何益处？我还听到另一件可悲的事：许多妇女以悲泣吸引情人，因哀哭的热切而获得对丈夫有情的名声。魔鬼的意图！撒殚的发明！我们不过是尘土和灰烬，不过是血肉之躯，要到几时呢？我们举目望天，思想神性的事。当我们做这些事时，我们怎能责备外教人，怎能劝勉他们呢？我们怎能就复活与他们辩论？怎能论及其他天上智慧呢？我们自己怎能没有恐惧地生活？你不知道悲伤带来死亡吗？因为悲伤使灵魂的视觉部分暗淡，不仅阻碍它觉察任何当觉察之事，还给它带来大害。因此，一种方式我们得罪天主，既无益于自己也无益于已逝之人；另一种方式我们中悦天主，并在人中获得尊敬。如果我们不自行消沉，祂很快会挪去我们沮丧的残余；如果我们不满，祂容许我们陷于悲伤。如果我们感恩，我们就不会沮丧。\Quote{但是，}有人说，\Quote{失去儿子、女儿或妻子，怎么可能不悲伤呢？}我回答不是\Quote{不悲伤}，而是\Quote{不要过分悲伤}。因为如果我们考虑是天主带走了，我们所拥有的丈夫或儿子是有死的，我们就会很快得到安慰。不满是那些寻求高于本性之事者的行为。你生而为人，是有死的；那么你为什么悲伤于本性之事的发生呢？你悲伤于靠饮食维生吗？你寻求不靠此生活吗？对待死亡也是如此，作为有死之人，不要现在就寻求不朽。这事已一次而永远地规定了。所以不要悲伤，不要扮演哀悼者，而要服从加于众人的律法。要为你们的罪悲伤；这是好的哀悼，这是最高的智慧。让我们因此持续为这原因哀悼，好使我们获得那里的喜乐，藉着我们的主耶稣基督的恩宠和慈爱，愿光荣归于祂，直到永远。阿们。\par

\SourceRecord{Translated by Charles Marriott.；Nicene and Post-Nicene Fathers, First Series；Vol. 14.；Edited by Philip Schaff.；Buffalo, NY: Christian Literature Publishing Co.,；1889.；<http://www.newadvent.org/fathers/240162.htm>.}

% structure: p0001 source=h1 target=section reason=page-title

\section{\WorkTitle{若望福音讲道集第六十三篇}}

% structure: p0002 source=h2 target=subsection reason=relative-heading-rank; base=h2

\subsection{\BibleRef{若 11:30-31}}

\Emph{\BibleQuote{耶稣还没有进村庄，仍在玛尔大迎接祂的地方。那些同她在一起的犹太人，}以及下文。}\par

一、哲学是一件伟大的恩赐，我指的是我们所有的哲学。因为外教人所拥有的不过是空谈和寓言，而且这些寓言中也没有任何真正的智慧；因为他们所做的一切都是为了名声。真正的智慧乃是一大恩赐，甚至在此世也给我们回报。因为那轻视财富的人，立即从中获益，摆脱了多余且无益的忧虑；那践踏荣耀的人，立即得到报酬，不为任何人所奴役，享有真正的自由；那渴望天上事物的人，在此世就得到补偿，视现世之事为无物，轻易超越一切忧伤。请看，例如这位妇人是如何因实践真正的智慧，甚至在此世就得了赏报。因为当众人都坐在她身边，陪她哀哭时，她没有等待老师到她那里去，也没有坚持她似乎应得的礼遇，也没有被悲伤所束缚（因为，除了其他的可怜之处，哀伤的妇人还有这种毛病：她们希望因自己的遭遇而受人看重），但她完全没有这样受影响；她一听见，就立刻来到祂那里。\BibleQuote{耶稣还没有进村庄。}祂稍微慢慢地前行，为的是不显得祂急于行神迹，而是由他们恳求。至少，福音作者这样说\BibleQuote{急忙起身}，要么是暗示这点，要么是表明她跑得很快，以致在基督到达之前就迎见了祂。她不是独自前来，而是带着屋里的犹太人跟在后面。她的妹妹很明智地悄悄叫她，以免打扰聚集的人，也没有说明原因；因为肯定有许多人会退回去，但此刻好像她要去哭泣一样，所有人都跟着她。通过这些，再次证明拉匝禄已经死了。\par

% structure: p0005 source=h2 target=subsection reason=relative-heading-rank; base=h2

\subsection{\BibleRef{若 11:32}}

她比她的妹妹更热切。她不顾众人，也不顾及人们对祂的猜疑，因为那里有许多祂的仇敌，他们说：\BibleQuote{这个开了瞎子眼睛的人，难道不能使这个人也不死么？}（\BibleRef{若 11:37}）但在她主人面前，她抛下一切世俗之事，只专注于一件事，即那主人的荣耀。她说了什么呢？\par

\BibleQuote{主，祢若早在这里，我的兄弟就不会死了。}\par

基督做了什么？祂此刻不与她交谈，也没有对她说像对她妹妹说过的话（因为有许多群众在场，这不是说这种话的时候）；祂只是适度行动，屈尊俯就；为证明祂的人性，祂默默流泪，暂时推迟了神迹。因为那个神迹是伟大的，是祂不常行的，并且许多人要因此相信；为了不因在群众不在场时施行神迹而成为他们的绊脚石，以至于他们因这神迹的伟大而一无所获，为了不失去猎物，祂以屈尊俯就吸引了众多见证人，并显出了祂人性的证据。祂流泪、心神烦乱；因为悲伤常会激动情绪。然后祂斥责这些情绪（因为\BibleQuote{在心神内叹息}意思是\Quote{抑制祂的烦恼}），就问：\par

% structure: p0009 source=h2 target=subsection reason=relative-heading-rank; base=h2

\subsection{\BibleRef{若 11:34}}

这样，问题就不会伴随着哀哭。但祂为什么问？因为祂不想自己强行（施行神迹），而是要从他们那里了解一切，并在他们的请求下做一切，以使神迹免于任何嫌疑。\par

\BibleQuote{他们对祂说：主，请来看。}\par

% structure: p0012 source=h2 target=subsection reason=relative-heading-rank; base=h2

\subsection{\BibleRef{若 11:35}}

你看见祂还没有显出任何复活迹象，而且去的时候不像要复活拉匝禄，倒像是要去哭泣？因为犹太人表明，在他们看来，祂似乎是去哀悼，而不是去复活他；至少他们说：\par

% structure: p0014 source=h2 target=subsection reason=relative-heading-rank; base=h2

\subsection{\BibleRef{若 11:36-37}}

即使在灾难中，他们也没有放宽他们的恶行。然而祂所要做的是一件更奇妙得多的事；因为赶走已来临并已取胜的死亡，比阻止它的来临要伟大得多。因此他们用那些本应令他们惊叹祂能力的事来诽谤祂。他们暂时承认祂开了瞎子的眼睛，而他们本应因那个神迹钦佩祂，却借此案给那个神迹抹黑，仿佛它根本没有发生过。不仅从这里看出他们全然败坏，还因为当祂尚未到来，也未显出任何行动时，他们不等事情结束就抢先控告祂。你看到他们的判断是多么败坏吗？\par

二、然后祂来到坟墓前，再次斥责自己的情绪。为什么福音作者在好几处仔细提到\BibleQuote{祂流泪}和\BibleQuote{祂叹息}？为叫你明白祂确实取了我们的本性。因为当这位福音作者比其他人更显著地宣讲关于基督的伟大事物时，在关乎身体的事上，他在这里也说得比他们更为谦卑。例如，关于祂的死，他丝毫没有说这类的话；其他福音作者宣告祂极其忧伤、痛苦挣扎；但若望相反却说祂甚至把差役打倒在地。所以，他在这里以提及祂的忧伤弥补了那里的省略。当说到祂的死时，基督说：\BibleQuote{我有权舍去我的生命}（\BibleRef{若 10:18}），然后祂没有说任何谦卑的话；因此，在受难时，他们赋予祂许多人性的事，以显示救恩计划的真实性。玛窦以痛苦挣扎、忧伤、战栗和汗珠证明这点；但若望以祂的忧伤证明。因为如果祂不是我们的本性，祂不会一次又一次地被悲伤所控制。耶稣做了什么？祂没有就他们的指控进行辩护；因为何必用言语使那些即将被事实封口的人缄默呢？这是一种较不烦扰且更适于羞辱他们的方式。\par

% structure: p0017 source=h2 target=subsection reason=relative-heading-rank; base=h2

\subsection{\BibleRef{若 11:39}}

为什么祂不在远处就呼唤拉匝禄，将他放在他们眼前？或者说，为什么祂不让石头还压在坟墓上时就使他复活？因为祂既能以声音驱动尸体，并使之再次获得生命，那么祂用同样的声音也能轻易挪开石头；祂曾以声音使那被殓布捆绑缠绕的人行走，岂不更能挪开石头？那么祂为什么不这样做呢？是为了使他们成为奇迹的见证人，免得他们像对待瞎子那样说：\Quote{是他}，\Quote{不是他}。因为他们的手和来到坟墓前证实了确实是那个人。如果他们不来，他们可能会以为自己看见的是幻象，或是另一个人代替了那个人。但现在，来到现场、搬开石头、吩咐他们解开那用殓布包裹的死者身上的束缚；抬他出坟墓的朋友们从殓布认出是他；他的姊妹没有留在后面；其中一位说：\Quote{他现在已经臭了，因为已经四天了}——这一切，足以堵住心怀恶意之人的口，因为他们成了奇迹的见证人。为此，祂吩咐他们搬开墓口的石头，以表明祂使人复活。为此，祂也问：\Quote{你们把他安放在哪里？}使那些说\Quote{你来看看}并带领祂的人，不能说是祂复活了另一个人；使他们的声音和手（声音说\Quote{你来看看}，手搬开石头并解开殓布）与他们的眼睛和耳朵（耳朵听见祂的声音，眼睛看见拉匝禄出来），以及他们的嗅觉（因玛尔大说：\Quote{他现在已经臭了，因为已经四天了}）一同作证。\par

因此我有理由说，那妇人根本没有理解基督的话：\Quote{他虽然死了，仍然要活着。}至少请看，她说话时好像由于时间相隔，这事是不可能的。确实，复活一个死了四天、已经腐烂的尸体是件奇事。耶稣对门徒说：\Quote{为使人子得到光荣}，是指祂自己；但对那妇人说：\Quote{你将看到天主的光荣}，是指圣父。你看见听者的软弱是话语差异的原因吗？因此祂提醒她曾对她说的话，几乎是在责备她，好像她忘记了。然而祂此时不愿使观众困惑，因此祂说：\par

% structure: p0020 source=h2 target=subsection reason=relative-heading-rank; base=h2

\subsection{\BibleRef{若 11:40}}

3. 信德确实是一项巨大的祝福，巨大，并且使那些以善生正确持守它的人变得伟大。藉着它，人们得以奉祂的名行天主的事。基督说得很好：\Quote{如果你们有信德，你们向这座山说：移动，它就会移动}\BibleRef{玛 17:20}；又说：\Quote{信我的人，我所做的事，他也要做，并且要做比这更大的事。}\BibleRef{若 14:12}祂说\InnerQuote{更大的}是什么意思？是指门徒们此后所行的事。因为连伯多禄的影子也复活了一个死人；这样，基督的能力更被显扬了。因为祂活着行奇迹固然奇妙，但祂死了以后，别人竟能因祂的名行比祂所行的更大的事，这更是无可争辩的复活证据；即使所有人都看见了那复活，也不会如此被相信。因为人们可能会说那是个幻象，但一个看见只因祂的名就成就了比祂与人交往时更大的奇迹的人，除非极其愚蠢，否则不能不信。那么，信德是一项巨大的祝福，当它出自热烈的感情、伟大的爱和热忱的灵魂时；它使我们真正有智慧，它隐藏我们人性的卑微，它撇下推理，默想天上的事；或者更确切地说，人的智慧无法发现的事，它却能充分领悟并成功。让我们因此紧握信德，不要将与我们相关的事交托给推理。因为请告诉我，希腊人为什么什么也没能发现？他们不是通晓了外邦人的一切智慧吗？为什么他们不能胜过渔夫、制帐棚者和无知的人？难道不是因为前者将一切托付给论证，而后者托付给信德吗？因此后者战胜了柏拉图、毕达哥拉斯，简言之，战胜了所有误入歧途的人；他们超越了那些一生耗在占星术、几何学、数学、算术上，并精通各种学问的人，正如真正的、实在的哲学家超越那些天生愚昧和失去理智的人。因为请看，这些人断言灵魂是不朽的，或者更确切地说，他们不仅断言，还说服了别人。相反，希腊人起初不知道灵魂是什么东西，当他们发现了并将它与身体区分后，又陷入同样的情形：一个说它是无形的，另一个说它是有形的，并与身体一同消散。关于天，一个说它有生命且是神，但渔夫们教导并说服众人说它是天主的作品和设计。希腊人使用推理并不奇怪，但那些看似信徒的人\Quote{他们}竟被发现是属血肉的，这才值得悲叹。正因如此，他们误入歧途：有人说他们认识天主如同天主认识自己——连任何希腊人都不敢如此断言；有人说天主不能无感而受生，甚至不让祂有任何超越人的地方；又有人说正义的生活和严谨的交往毫无用处。但现在不是驳斥这些的时候。\newline{}4. 然而，如果生活败坏，正确的信德也毫无用处，这是基督和保禄都宣告的，他们更关心后一部分；基督教导说：\Quote{不是每一个对我说\InnerQuote{主啊，主啊}的人，都能进入天国}\BibleRef{玛 7:21}；又说：\Quote{在那一天，许多人要对我说：\InnerQuote{主啊，我们不是因祢的名说过预言吗？}那时我要向他们声明说：\InnerQuote{我从来不认识你们；这些作恶的人，离开我吧！}}\BibleRef{玛 22:23}（因为不慎自守的人，即使有正确的信德，也容易滑入邪恶）；保禄在致希伯来人书中这样劝勉说：\Quote{你们要跟众人追求和平与圣德，没有圣德谁也不能看见主。}\BibleRef{希 12:14}此处\InnerQuote{圣德}指贞洁，所以每个人都应满足于自己的妻子，不与别的女人发生关系；因为不如此是不可能的，他必定灭亡，即便有万般善行，因为奸淫不能进入天国。或者更确切地说，这从此不再是奸淫，而是通奸；因为女人若与男人结合，又与别的男人结合，就是犯了通奸；同样，男人若与女人结合，又有别的女人，就是犯了通奸。这样的人不能承受天国，而要堕入火坑。请听基督关于这些人所说的话：\Quote{他们的虫子不死，火也不灭。}\BibleRef{谷 9:44}因为一个有了妻子和妻子的安慰之后，又对另一个女人无耻行事的人，是没有宽恕的；因为这从此是纵欲。如果许多人在斋戒或祈祷的时候甚至戒绝自己的妻子，那么一个不满足于自己的妻子、反而与别的女人苟合的人，为自己积聚了多大的火！如果一个人休弃并赶出自己的妻子后，不许可与别的女人结合（因为这是通奸），那么一个妻子还在家中却领进另一个女人的人，犯了多大的恶！因此，不要让任何人容留这疾病在他的灵魂中；要连根拔除。他并不怎么伤害妻子，而是伤害自己。因为这罪如此严重和不可宽恕，以致如果一个女人没有丈夫（一个偶像崇拜者）的同意就离开他，天主会惩罚她；但如果她离开一个淫乱者，就不这样。你们看出这是多大的恶吗？它说：\Quote{如果某个信主的妇人有不信的丈夫，而丈夫情愿与她同住，她不可离弃他。}\BibleRef{格前 7:13}关于妓女却不是这样；而是怎样？\Quote{凡休妻的，除非因奸淫的缘故，就是使她犯奸淫。}\BibleRef{玛 5:32}因为若结合使两人成为一体，那么与妓女结合的人必然成为与她一体。那么，那端庄的妇人，身为基督的肢体，怎能接纳这样的人？或者她怎能将妓女的肢体与自己结合？请注意一种（奸淫）比另一种（偶像崇拜）的严重性。与不信者同住的妇人并不污秽；（因为\Quote{不信的丈夫因妻子而成圣}\BibleRef{格前 6:15}）但对妓女却不是这样；而是怎样？\Quote{那我怎能将基督的肢体作为妓女的肢体呢？}\BibleRef{格前 6:15}在前者，圣化仍然保留，即使不信者与妻子同住也不被除去；但在后者，圣化离去。奸淫是一个可怕的、可怕的事，它导致永远的惩罚；即使在此世，它也带来万般灾祸。如此犯罪的人被迫过着焦虑和劳苦的生活；他并不比那些受惩罚的人更好，他怀着恐惧和极大的战栗潜入别人的家，怀疑所有的人，无论是奴隶还是自由人。因此我劝你们从这疾病中得释放，如果你们不听从，就不要踏上神圣的门槛。长了疥疮、充满疾病的羊不能与健康的羊群在一起；我们必须把它们赶出羊圈，直到它们除去疾病。我们成了基督的肢体；我恳求你们，不要让我们成为妓女的肢体。这地方不是妓院，而是教会；如果你们有妓女的肢体，就不要站在教会里，免得侮辱这地方。即使没有地狱，没有惩罚，但在那些婚约、婚礼火炬、合法的床、生儿育女、夫妇交往之后，你怎么能忍受与另一个结合？你怎么不羞愧、不脸红？你不知道吗？许多人在自己的妻子死后领进另一个女人到家里，都会受到许多人的责备？但这行为没有刑罚；而你妻子还活着，你却领进另一个。这是何等的邪淫！要听关于这样的人所说的话：\Quote{他们的虫子不死，火也不灭。}\BibleRef{谷 9:44}要战栗于这威胁，畏惧这惩罚。此处的快乐不如那里的惩罚大，但愿没有人成为这惩罚的承受者，而是藉着操练圣德得以看见基督，获得所应许的美善，愿我们众人藉着我们的主耶稣基督的恩宠和慈爱，与圣父和圣神，获得这一切，永世无尽。亚孟。\par

\SourceRecord{Translated by Charles Marriott.；Nicene and Post-Nicene Fathers, First Series；Vol. 14.；Edited by Philip Schaff.；Buffalo, NY: Christian Literature Publishing Co.,；1889.；<http://www.newadvent.org/fathers/240163.htm>.}

% structure: p0001 source=h1 target=section reason=page-title

\section{若望福音讲道 第六十四篇}

% structure: p0002 source=h2 target=subsection reason=relative-heading-rank; base=h2

\subsection{\BibleRef{若 11:41-42}}

1. 我常说的，现在还要说：基督所顾及的，并非自己的荣耀，而是我们的救恩；不是如何说出一些崇高的话，而是如何说出能吸引我们归向祂的话。因此，祂崇高而有力的言论很少，并且也是隐藏的，但谦卑卑微的言论却很多，且遍布祂的谈话中。因为既然人们更被这些话语所吸引，祂便持续使用它们；祂一方面不普遍地讲论这些（崇高的话），免得后来的人受损，另一方面也不完全隐藏它们，免得当时的人反感。因为那些从心志卑下达到成全的人，将能从即使一句崇高的道理中洞察全体，但那些始终心志卑下的人，若不常听到这些卑微的话，就根本不会到祂这里来。事实上，即使听了这么多这样的话语，他们仍不坚定，反而用石头砸祂、迫害祂，想杀祂，称祂为亵渎者。当祂使自己与天主相等时，他们说：\BibleQuote{这人说了亵渎的话。}\BibleRef{玛 9:3}当祂说：\BibleQuote{你的罪赦了。}\BibleRef{若 10:20}他们竟又称祂为附魔的人。所以当祂说，听祂话的人比死亡更强，或说：\BibleQuote{我在父内，父在我内。}\BibleRef{若 8:51}他们便离弃祂；而当祂说祂从天降下时，他们便反感。\EditorialNote{参见若6:33、60}如果连这些言论（尽管不常讲）都使他们不能忍受，那么，假使祂的谈论总是崇高，总是这种结构，他们还能留心听祂吗？因此当祂说：\BibleQuote{父怎样命令了我，我就怎样讲。}\BibleRef{若 14:31}和\BibleQuote{我不是凭自己来的。}\BibleRef{若 7:28}时，他们便相信了。他们当时信了，从福音作者另外指明这一点可以清楚，他说：\BibleQuote{祂说这些话的时候，许多人信了祂。}\BibleRef{若 5:30}既然卑微的话吸引人信德，崇高的话却吓跑他们，那么不一眼看出如何盘算那些卑微言论的唯一原因，即它们是为了听众的缘故而说，岂不是极端的愚昧吗？因为在另一处，当祂想说一些崇高的事时，祂没有说，并附加了这个理由说：\BibleQuote{免得我们得罪他们，往海里去投钩。}\BibleRef{玛 17:27}祂在这里也是如此；因为在说了\Quote{我知道祢常常俯听我}之后，祂接着说：\BibleQuote{但为了站在周围的群众，我说了这话，为叫他们相信。}\BibleRef{若 12:30}这些话是我们的吗？这是人的猜想吗？那么当一个人不肯被所写的说服，即他们因崇高的事反感，他却听见基督说祂用卑微的方式说话，免得他们反感，此后，他怎能还怀疑那些平庸的话是出于祂的本性，而非祂的屈尊就卑呢？同样，在另一处，当有声音从天降下时，祂说：\BibleQuote{这声音不是为我的缘故，而是为你们的缘故。}\BibleRef{若 12:30}被高举的人可以允许说自己卑微的事，但卑微的人对自己说出任何伟大或崇高的事是不合法的。因为前者出自屈尊就卑，其原因是听者的软弱；或者更确切地说，是由于引导他们谦卑，以及祂穿上血肉，教导听众不要对自己说任何大事，以及祂被视为天主的敌人，不被相信是从天主来的，祂被怀疑违背法律，以及听众以恶眼看祂、对祂怀有恶意，因为祂说与天主平等。但一个卑微的人对自己说任何大事，既没有合理的原因，也没有不合理的原因；只能是愚蠢、无耻和不可饶恕的胆大妄为。那么，基督为什么说话卑微，而祂本是那不可言说、伟大的实体呢？由于上述的原因，也为了祂不被视为非受生者；因为保禄似乎也害怕类似的事；所以他接着说：\BibleQuote{除了那使万有屈服于祂的。}\BibleRef{格前 15:27}连想这样的事也是不虔敬的。因为如果祂小于那生祂的，且是不同实体的，却被认为相等，难道祂不会使用一切手段使这事不被认为吗？但现在祂却相反，说：\BibleQuote{我若不做那差我来者的工，你们就不要信我。}\BibleRef{若 10:37}的确，祂说\BibleQuote{我在父内，父在我内}\BibleRef{若 14:10}，向我们暗示了平等。如果祂是较低的，本应十分强烈地驳斥这种意见，并且完全不该说\BibleQuote{我在父内，父在我内}\BibleRef{若 10:30}，或\BibleQuote{我们是一体}\BibleRef{若 10:30}，或\BibleQuote{看见我的，就是看见了父。}\BibleRef{若 14:9}同样，当祂论及权力时，祂说：\BibleQuote{我与父原是一体。}\BibleRef{若 10:30}当祂论及权威时，祂又说：\BibleQuote{因为父怎样唤醒死者并使他们活，子也照样随祂的意愿使人活。}\BibleRef{若 5:21}如果祂是不同实体的，就不可能做这些；或者即使可能，也不该如此说，免得他们怀疑实体是同一个。因为为了使他们不以为祂是天主的仇敌，祂甚至经常说出不合祂身份的话，那么那时祂更应当这样做；但现在，祂说：\BibleQuote{为使众人尊敬子如同尊敬父一样。}\BibleRef{若 5:23}祂说：\BibleQuote{祂所做的工，我也做。}\BibleRef{若 5:19}祂说祂是\Quote{复活、生命、世界的光}\EditorialNote{参见若11:25；若8:12}，这些是使自己与那生了祂的相等、并证实他们所持怀疑的表达。你看见祂如何做出这番讲话和辩护，以表明祂没有违反法律，并且不仅没有消除，甚至证实了祂与父平等的意见吗？所以当他们说：\BibleQuote{你说亵渎的话，因为你把自己当作天主}\BibleRef{若 10:33}时，祂从工作的相等确立了这一点。\par

2. 我为什么说圣子做了这事，而那没有取肉身的圣父也做同样的事呢？因为祂也为了听众的救恩，忍受了许多关于祂的卑微说法。因为，\BibleQuote{亚当，你在哪里？}\BibleRef{创 3:9}，以及\BibleQuote{让我知道他们是否完全按照那呼喊做了}\BibleRef{创 18:21}；和\BibleQuote{现在我知道你敬畏天主}\BibleRef{创 22:12}；和\BibleQuote{如果他们愿意听}\BibleRef{则 3:11}；和\BibleQuote{如果他们愿意明白}\BibleRef{申 5:29}；和\Quote{谁会给这百姓这样的心？}以及这句话，\BibleQuote{主啊，在诸神中没有像祢的}\BibleRef{咏 79:29}；这些以及其他许多在旧约中类似的句子，如果有人把它们挑出来，他会发现它们不相称于天主的尊威。在阿哈布的事上记载着，\BibleQuote{谁为我引诱阿哈布？}\BibleRef{编下 18:19}。而且祂不断将自己与外邦人的神祇比较，所有这些都不相称于天主。然而，在另一方面，它们被变得相称于祂，因为祂是如此仁慈，为了我们的救恩，祂不在乎那些符合祂尊严的表达。的确，成为人是不相称于祂的，取了仆人的形体、说谦卑的话、穿上谦卑的衣服，若看祂的尊严就不相称，但若考虑祂慈爱那说不尽的丰富，就是相称的。祂话语谦卑还有另一个原因。那是什么？就是他们认识并承认圣父，却不认识祂。因此祂不断转向那被他们承认的圣父，因为祂自己那时尚未被认为值得信赖；这不是因为祂自己有任何低劣，而是因为听众的愚昧和软弱。为此祂祈祷说：\Quote{父啊，我感谢祢，因为祢俯听了我。}因为如果祂愿意赐予谁生命，就赐予生命，和圣父一样，为什么祂要呼求圣父呢？\par

但现在是从头开始讲解这段经文的时候了。\BibleQuote{于是他们移开了死人躺卧的石头。耶稣举目向上说：父啊，我感谢祢，因为祢俯听了我。我本来知道祢常常俯听我，但是为了四周站立的群众，我才说了这话，为叫他们相信是祢派遣了我。}那么，我们来问问异端者：祂从祈祷中获得动力，从而复活了死人吗？那么，祂如何不祈祷也行其他奇迹呢？祂说：\BibleQuote{你这恶鬼，我命令你，从他身上出来}\BibleRef{谷 9:25}；和\BibleQuote{我愿意，你洁净了吧}\BibleRef{谷 1:41}；和\BibleQuote{起来，拿起你的床走吧}\BibleRef{若 5:8}；和\BibleQuote{你的罪赦了}\BibleRef{玛 9:2}；并对海说：\BibleQuote{安静，平息了吧}\BibleRef{谷 4:39}。总之，如果祂也凭借祈祷工作，祂比宗徒们多什么？或者我应该说，就连他们也不是都靠祈祷，而是常常不祈祷，只呼号耶稣的名字行事。现在，如果祂的名字有如此大的能力，祂怎能需要祈祷？如果祂需要祈祷，祂的名字就不会有效。当祂完全造人时，祂需要什么祈祷？那时岂不是有极大的荣耀平等吗？\BibleQuote{让我们造}，它说，\BibleQuote{人}\BibleRef{创 1:26}。如果祂需要祈祷，还有什么比这更显软弱的迹象？但让我们看看这祈祷是什么：\BibleQuote{我感谢祢，因为祢俯听了我。}有谁这样祈祷过？在说出任何祈祷之前，祂就说：\BibleQuote{我感谢祢，}表明祂不需要祈祷。\BibleQuote{我也知道祢常常俯听我。}祂这样说，不是好像祂自己无能为力，而是为显示祂的意愿与父的意愿是一个。但为什么祂采取祈祷的形式？听，不是我的，而是祂自己的话：\BibleQuote{为了四周站立的群众，叫他们相信是祢派遣了我。}祂没有说：\Quote{叫他们相信我是低劣的，我需要从上而来的动力，没有祈祷我什么都不能做；}而是说：\BibleQuote{是祢派遣了我。}而是摒弃了这一切，为使你们没有这样的怀疑，祂说出了祈祷的真正原因：\Quote{叫他们不认为我是天主的敌人；叫他们不说：祂不是从天主来的，好让我向他们显示这事是按祢的旨意成就的。}差不多是说：\Quote{如果我是天主的敌人，所成就的就不会成功，}但\BibleQuote{祢俯听了我，}这句话用在朋友和同等者之间。\BibleQuote{我也知道祢常常俯听我，}也就是说：\Quote{为使我的意愿成就，我不需要祈祷，除非为了说服人相信祢和我属于同一旨意。}\Quote{那么祢为什么祈祷？}是为了软弱和粗俗的人。\par

% structure: p0006 source=h2 target=subsection reason=relative-heading-rank; base=h2

\subsection{若 11:43}

为什么祂不说\Quote{因我父的名，出来}？或者为什么祂不说\Quote{父啊，使他复活}？为什么祂省略了所有这些说法，而在采取祈祷的姿态后，通过祂的行动显示祂独立的权柄？因为这也是祂智慧的一部分，用言语显示俯就，用行动显示能力。因为除了说祂不是从天主来的，他们再没有别的可以控告祂，并且他们用这种方式欺骗了许多人，为此祂通过祂所说的话最充份地证明这一点，并且以他们的软弱所需要的方式。因为祂本有能力用其他方式同时显示祂与父的一致和祂自己的尊威，但群众不能上升到那么高。祂说：\par

% structure: p0008 source=h2 target=subsection reason=relative-heading-rank; base=h2

\subsection{若 11:43}

3. 这正是祂所说的：\Quote{时候将到，死者要听见圣子的声音，听见的人就必活起来。} \BibleRef{若 5:28} 因为，为使你不以为祂从别人那里获得工作的能力，祂先前已教导你这件事，并以行动证明了，祂没有说\Quote{起来}，而是说\Quote{出来}，与死人谈话如同活人一般。这权柄谁能相比？如果祂不是凭自己的能力做的，那么祂比宗徒们有什么更多呢？宗徒们说：\Quote{你们为什么注视我们，好像我们凭自己的能力和虔诚使这人行走呢？} \BibleRef{宗 3:12} 因为如果祂不是凭自己的能力工作，祂没有加上宗徒们论及自己所说的话，那么他们在某种意义上比祂更有智慧，因为他们拒绝了光荣。又在别处说：\Quote{你们为什么做这些事？我们也是和你们一样有血肉的人。} \BibleRef{宗 14:15} 宗徒们既然自己不做什么，就这样说以说服人；但祂，当类似的想法在祂身上形成时，祂难道不会消除这怀疑吗？如果至少祂不是凭自己的权柄行事？谁能断言？但事实上基督做了相反的事，祂说：\Quote{因为站在周围的群众，我才说了这话，为叫他们相信}；所以如果他们相信了，就不需要祈祷了。如果祈祷并不有损祂的尊严，为什么祂把祂祈祷的原因归咎于他们？为什么祂不说：\Quote{我这样做，为叫他们相信我和祢并不平等}？因为祂本应因怀疑而说出这话。当祂被怀疑破坏法律时，祂使用了这说法，甚至在他们还没说什么时，祂说：\Quote{不要以为我来是废除法律} \BibleRef{玛 5:17}；但在这里，祂却证实了他们的怀疑。事实上，绕这么大的圈子，用这么隐晦的话，有什么必要呢？只要说\Quote{我不平等}就足够了，就摆脱了此事。有人说：\Quote{但祂不是说过：我不凭自己行事吗？}就连这一点，祂也是以隐秘的方式做的，适合他们的软弱，并且出于与祈祷同样的原因。但\Quote{祢已听从了我}是什么意思？意思是：\Quote{在我这方面，没有与祢相悖的事。}因此，正如\Quote{祢已听从了我}并不是宣告祂自己没有能力的话（因为如果是这样，那不仅是无能，也是无知，如果祂在祈祷前不知道天主会应允祈祷；如果祂不知道，为什么祂说\Quote{我去叫醒他}而不是\Quote{我去祈求我的圣父叫醒他}呢？）。因此，这个说法不是软弱的标记，而是旨意一致的标记；同样，\Quote{祢常常听从我}也是如此。那么我们必须说，要么是这样，要么是它针对他们的怀疑。如果祂既不是无知也不是软弱，那么显然祂说出这些谦卑的话，是为了使你们因它们过分的卑微而被说服，并不得不承认，这些话不符合祂的尊严，而是出于俯就。那么真理的敌人怎么说？有人说：\Quote{祂说这些话，不是为了听者的软弱，而是为了显示优越。}但这并不是显示优越，而是极大谦卑自己，并显示自己没有什么超过人。因为祈祷不属于天主，也不属于同享宝座的那一位。那么你看到祂这样做，除了他们的不信之外，没有别的原因吗？至少注意，行动见证了祂的权柄。\par

\Quote{祂叫了，死人便裹着出来。} 然后，为使此事不像是幻觉（因为他裹着出来，其奇迹不亚于复活），耶稣吩咐解开他，好让他们触摸并靠近他，看见那真是他。祂又说：\par

% structure: p0011 source=h2 target=subsection reason=relative-heading-rank; base=h2

\subsection{\BibleRef{若 11:43}}

你看祂多么不炫耀？祂没有领他走，也没有命他和自己同行，免得有人以为祂在展示他；祂深知如何保持适度。\par

这神迹完成后，有人惊讶，有人去告诉法利塞人。他们做了什么？他们本应惊奇并赞美祂，却商议要杀死那复活死人的。何等愚蠢！他们以为能把那在别人身上战胜死亡的交与死亡。\par

% structure: p0014 source=h2 target=subsection reason=relative-heading-rank; base=h2

\subsection{\BibleRef{若 11:47}}

这些人接受了祂神性的如此证明，却仍称祂为\Quote{人}。\Quote{我们做什么呢？}他们本应相信、事奉、朝拜祂，而不再视祂为一个凡人。\par

% structure: p0016 source=h2 target=subsection reason=relative-heading-rank; base=h2

\subsection{\BibleRef{若 11:48}}

他们商议要做什么？他们想煽动民众，仿佛他们自己会因为建立王国的嫌疑而陷入危险。\Quote{因为如果，}其中一人说，\Quote{罗马人得知这人带领群众，就会怀疑我们，并来摧毁我们的城。}请问，这是为什么？祂教导叛乱了吗？祂不是允许你们给凯撒纳税吗？你们不是想立祂为王而祂却躲避你们吗？祂不是过着卑微无华的生活，既无房屋也无其他类似之物吗？因此，他们这样说，并非出于任何这样的担心，而是出于恶意。然而，结果却与他们的预期相反，罗马人在他们杀死基督后夺取了他们的民族和城邑。因为祂所行的一切都远超任何怀疑。那医治病人、教导最卓越的生活方式、命令人服从统治者的主，不是在建立暴政，而是在拆毁暴政。\Quote{但是，}有人说，\Quote{我们从先前的（骗子）猜测。}但他们教导叛乱，而祂则教导相反。你看见这些话不过是借口吗？因为祂并未表现出任何这样的行为？祂身边有显赫的卫队吗？祂有车队随行吗？祂不是常往荒野去吗？但他们为了使自己不显得出于恶意，便说全城都处于危险之中，公共利益正受威胁，他们不得不担心最坏的情况。这些并不是你们被掳的原因，而是与之相反的原因；无论是本次被掳，还是巴比伦之囚，以及后来安条克时期的被掳：并非因为你们中有敬拜者，而是因为你们中行不义的人激怒了天主，这才导致你们被交给仇敌。嫉妒就是这样的事，一旦蒙蔽了灵魂，就使人看不见当看的事物。祂不是教导人温良吗？祂不是命令人当被打右颊时，连左颊也转过来由他打吗？祂不是命令人受伤害时忍受吗？要表现出比他人施害更大的忍耐吗？这些，请问，是建立暴政的迹象，还是拆毁暴政的迹象？\par

但是，如我所说，恶意是可怕的事，充满虚伪；这给世界带来了无数邪恶；通过这种病症，法庭充满；由此产生对名声和财富的贪欲，由此产生对统治的贪求和傲慢，因此道路上有了邪恶的强盗，海上有海盗，由此产生世上的谋杀，因此我们的种族分裂，无论你看到什么邪恶，你都会发现它起源于此。这甚至闯入了教会，从一开始就带来了无数可怕的事，这是贪婪之母，这种病症颠倒了万物，败坏了正义。因为\Quote{礼物，}它说，\Quote{能蒙蔽智慧人的眼目，又如嘴上的口罩，转离责斥。}\BibleRef{德 20:29}这使自由人变成奴隶，我们每天都在谈论这事，却没有益处，我们变得比野兽更坏；我们掠夺孤儿，剥削弱妇，欺压穷人，将祸患加于祸患。\BibleQuote{哀哉！义人已从地上灭绝了！}\BibleRef{米 7:1}从今以后，我们也应当哀哭，或者不如说，我们每天都需要说这话。我们从祈祷中一无所得，从劝告和劝诫中一无所得，因此我们只剩下哭泣。基督也是这样行的；在多次劝诫耶路撒冷的人之后，当他们毫无进益时，祂曾为他们的顽固而哭泣。先知们也是如此，现在我们也要这样行。如今是哀哭、流泪、哀悼的时期；对我们来说，现在也适合说：\BibleQuote{叫哀悼的妇女来，并召那些精于哭丧的，使她们大声哀号。}\BibleRef{耶 9:17}或许这样我们能驱除那些建造华丽房屋、以抢劫所得田产围困自己的人的病态。现在是哀哭的时候；但你们这些被剥夺受欺压的人，要与我一同哀哭，你们的哀哭引来我的泪水。但在哀哭时，我们要哀哭的不是我们自己，而是为他们；他们没有伤害你们，而是毁灭了自己；因为你们因所受的不义将得到天国，他们却因自己的利益将得到地狱。为此，受伤害比伤害人更好。让我们以并非人为的哀歌，而是以先知们也哀哭过的圣经哀歌为他们哀哭。让我们与依撒意亚一同哀哭，说：\BibleQuote{祸哉，那些以房接房，以地连地，以致不留余地，只想独居地上的人！许多华美的房屋，必成为荒凉，无人居住。}\BibleRef{依 5:8}\par

让我们与\Person{纳鸿}一同哀悼，并与他同声说：\BibleQuote{祸哉，那将自己的房屋建在高处的。} \BibleRef{耶 22:13} 或者不如让我们为他们哀悼，如同基督哀悼古时的人们一样。\BibleQuote{你们富有的是有祸的，因为你们已获得了你们的安慰。} \BibleRef{路 6:24} 我恳求你们，让我们不要停止这样哀哭，若这不失体统，甚至让我们为弟兄们的疏忽捶胸顿足。我们不要为已死的人哭泣，却要为那强暴、贪婪、吝啬、贪得无厌的人哭泣。我们为何要为死者哀悼？对他们而言，今后已无法成就任何事。让我们为那些仍能改变的人哀悼吧。但当我们哀叹时，他们或许会嘲笑。这本身就是值得哀哭的原因：他们该哀哭时却嘲笑。因为倘若他们被我们的忧伤所感动，我们本应因他们承诺悔改而停止忧伤；但因他们麻木不仁，让我们继续哭泣，不仅为富人，也为爱财者、贪婪者、强暴者。财富并非邪恶之物（因为当我们将其用于接济有需要的人时，可以正确使用它），但贪婪是邪恶的，它预备了永罚。那么让我们为他们哀哭吧；也许会有一些悔改；即便那些已陷入的人不能逃脱，至少其他人不会陷入危险，而会防范它。愿他们能摆脱恶疾，也愿我们中间无人陷入其中，使我们众人能共同获得所应许的福泽，藉着我们的主耶稣基督的恩宠和慈爱。愿光荣归于祂，直到永远。阿们。\par

\SourceRecord{Translated by Charles Marriott.；Nicene and Post-Nicene Fathers, First Series；Vol. 14.；Edited by Philip Schaff.；Buffalo, NY: Christian Literature Publishing Co.,；1889.；<http://www.newadvent.org/fathers/240164.htm>.}

% structure: p0001 source=h1 target=section reason=page-title

\section{\WorkTitle{若望福音讲道集第六十五篇}}

% structure: p0002 source=h2 target=subsection reason=relative-heading-rank; base=h2

\subsection{\BibleRef{若 11:49-50}}

1.  \BibleQuote{外邦人陷在自己所掘的坑中；他们的脚被自己暗设的网罗缠住。}\BibleRef{咏 9:15}犹太人的情形正是如此。他们说，他们要杀死耶稣，免得罗马人来夺去他们的地方和民族；而当他们杀了祂之后，这些事就临到他们身上。他们藉着杀死祂来逃脱，却终究没有逃脱。被杀的已在天上，杀人的却得到地狱为业。然而他们并不考虑这些事，而是怎样呢？经上说：\BibleQuote{从那一天起，他们就议决要杀害祂。}\BibleRef{若 11:53}因为他们说：\BibleQuote{罗马人必要来，夺去我们的民族；}其中一人，盖法，那一年做大司祭（比别人更无耻），\BibleQuote{你们什么都不懂。}别人犹豫不决、谨慎商议的事，这人却无耻、公开、大胆地喊叫。他说什么？\BibleQuote{你们什么都不懂，也不想想：一个人替百姓死，免得整个民族灭亡，为你们是有利的。}\par

% structure: p0004 source=h2 target=subsection reason=relative-heading-rank; base=h2

\subsection{\BibleRef{若 11:51}}

你看见大司祭权柄的力量何等大吗？或者，既然他无论如何被看为配得大司祭职分，尽管他不配，他仍预言了，却不知道自己在说什么；恩宠仅仅利用了他的口，却没有触动他被咒诅的心。事实上，许多其他不配的人也预言了未来的事，如拿步高、法郎、巴郎；理由都是显而易见的。但他的话说的是：\Quote{你们仍然静坐，漫不经心地看待这事，不知道为了团体的利益可以轻视一个人的安全。}请看圣神的能力何等伟大：祂能从邪恶的意念中引出充满奇妙预言的言语。福音作者称外邦人为\Quote{天主的子女}，是从将要发生的事说的；正如基督自己说：\BibleQuote{我另有羊}（\BibleRef{若 10:16}），也是从以后要发生的事称呼他们。\par

但\Quote{那一年做大司祭}是什么意思？这件事也和其余的一样腐败了；因为从官职成为可购买的东西时起，他们不再终身做司祭，而只做一年。尽管如此，在这情形中圣神仍然临在。但当他们举手反对基督时，圣神就离开了他们，转到宗徒们那里。这由圣所帐幔的撕裂和基督的话表明：\BibleQuote{看，你们的房屋将给你们留下荒凉。}\BibleRef{玛 23:38} 约瑟夫在稍后的时候记载，某些天使仍然与他们同在（要看他们是否改变行为），但后来离开了他们。当葡萄园还在时，一切都进行；但当他们杀了继承人之后，就不再如此，他们灭亡了。天主从犹太人那里夺取葡萄园，如同从一个无用的儿子身上脱去一件荣耀的衣服，把它交给了外邦人中正直的仆人，使前者荒凉赤身。此外，连一个敌人也能预言这事，不是小事。这可能吸引其他人。因为论到他的意愿，事情恰恰相反：当祂死了，信友们反而因此免于将来的惩罚。\BibleQuote{好聚集那些四散的}（\BibleRef{若 11:52}）是什么意思？祂使他们成为一个身体。住在罗马的人视印度人如同自己的肢体。有什么比这个\BibleQuote{聚集}更伟大的呢？万有的头就是基督。\par

% structure: p0007 source=h2 target=subsection reason=relative-heading-rank; base=h2

\subsection{\BibleRef{若 11:53}}

事实上，他们早已想要这样做；因为福音作者说：\BibleQuote{因此犹太人想要杀祂}\BibleRef{若 5:18}；又说：\BibleQuote{你们为什么想要杀我？}\BibleRef{若 7:19} 但那时他们只是想要，现在他们批准了他们的决定，并把那行动当作他们的事务。\par

% structure: p0009 source=h2 target=subsection reason=relative-heading-rank; base=h2

\subsection{\BibleRef{若 11:54}}

2. 祂又以人的方式拯救自己，祂常这样做。但我已提过祂常离开和退去的原因。这时祂住在厄弗辣大，靠近旷野，并与祂的门徒们同住。你想那些门徒看见祂像人一样拯救自己时，心里何等混乱？此后没有人跟随祂。因为节期临近，众人都跑到耶路撒冷；但他们，当别人都在欢乐、举行盛会时，却躲藏起来，处于危险中。然而他们仍然与祂同在。因为他们在逾越节和帐棚节时藏在加里肋亚；此后又在节期中，只有他们与他们的师傅一同逃亡和隐藏，显明了他们对祂的善意。因此路加记载祂说：\BibleQuote{我在你们的诱惑中与你们同在}；祂这样说，表明他们因祂的影响而坚强。\par

% structure: p0011 source=h2 target=subsection reason=relative-heading-rank; base=h2

\subsection{\BibleRef{若 11:55}}

% structure: p0012 source=h2 target=subsection reason=relative-heading-rank; base=h2

\subsection{\BibleRef{若 11:57}}

奇妙的净化，带着杀人的意愿、杀人的意图和染血的手！\par

% structure: p0014 source=h2 target=subsection reason=relative-heading-rank; base=h2

\subsection{\BibleRef{若 11:56}}

藉着逾越节他们阴谋反对祂，把欢庆的时间变成谋杀的时间，因为那时祂会落入他们手中，因为节期召唤祂。何等不敬！当他们需要更加谨慎，并应该赦免那些犯下最重罪行的人时，他们却企图陷害一位没有做过任何错事的人。然而这样行，他们不但一无所获，反而变得可笑。为此，祂不断来到他们中间，却逃脱，并当他们商议要杀祂时约束他们，使他们困惑，以显示祂的能力来刺痛他们；这样当他们抓住祂时，他们就知道所发生的事不是出于他们的能力，而是出于祂的许可。因为即使那时，他们也不能抓住祂，尽管伯达尼很近；而当他们真的抓住祂时，祂把他们打退。\par

% structure: p0016 source=h2 target=subsection reason=relative-heading-rank; base=h2

\subsection{\BibleRef{若 12:1-2}}

这是祂复活真实性的证明：过了许多天，祂仍然活着并进食。\Quote{玛尔大伺候}；由此可知饭是在她家里，因为他们接待耶稣如同爱人和被爱者。然而有人说，这事发生在另一人家。玛利亚没有伺候，因为她是个门徒。这里她又以更属灵的方式行事。因为她并非应邀伺候，也没有把服务提供给所有人。她只把尊荣归于祂一人，并接近祂，不将祂当作人，而是当作天主。为此她倾倒了香液，用头发擦祂的脚，这是那种对祂持有与别人不同看法之人的行为；然而犹达斯责备了她，假借关切之名。基督如何说？\Quote{她在我埋葬的事上做了一件好事。}为什么祂不在这妇人面前揭露那个门徒，也不对他说福音作者所宣告的话，即因他自己的偷窃而责备她？因着祂极大的忍耐，祂愿意将他引向更好的心。因为祂知道他是个叛徒，从一开始就多次责备他，说：\Quote{不是所有人都信}，和\Quote{你们中间有一个是魔鬼。}\BibleRef{若 6:64}祂向他们表明祂知道他是叛徒，却没有公开责备他，而是容忍他，想要挽回他。那么，另一位福音作者怎么说所有门徒都说了这些话？\BibleRef{玛 26:70}所有人都说了，他也说了，但别人不是出于同样的目的。如果有人问为什么祂把穷人的钱袋交给一个小偷，并让一个爱财的人做管家，我们会回答：天主知道秘密的理由；但如果我们可以猜测一下，那是为了切断他的一切借口。因为他不能说这事是出于爱财（因为他钱袋里有足够的钱满足他的欲望），而是出于极度的邪恶，基督想要遏制它，对他用了很多迁就。所以祂甚至没有责备他偷窃，虽然知道，阻止了他邪恶欲望的道路，并剥夺了他的一切借口。\Quote{由她吧，}祂说，\Quote{因为她做这事是为我埋葬之日。}祂又提到了叛徒，谈及祂的埋葬。但是责备却没有触及他，那话也没有软化他，尽管足以激发他的怜悯：好像祂说：\Quote{我是个负担和麻烦，但稍等一下，我就离开。}这也是祂说这话的意思，\par

% structure: p0018 source=h2 target=subsection reason=relative-heading-rank; base=h2

\subsection{若 12:8}

但这些事都不能使那个野蛮的疯子回头；事实上耶稣所说的和所做的远远超过这些，那晚祂洗了他的脚，让他分享同席和盐，这是通常连强盗的心也能约束的事，还说了其他足以融化石头的话，而且不是在很久以前，而是在那一天，以免时间使一切被遗忘。但他抗拒了一切。\par

3. 因为可怕，可怕的是贪财，它使双眼和双耳失灵，使人比野兽更难对付，不允许人考虑良心、友谊、共融或自己灵魂的得救，反而一下子从这一切中撤回，像严厉的女主人，使被它俘获的人成为奴隶。这种苦涩的奴役中可怕的是，它甚至说服他们为此感激；他们越成为奴隶，快乐就越增加；而正因如此，疾病变得不可医治，怪物变得难以征服。这使革阿齐从门徒和先知变成了癞病人；这毁灭了阿纳尼雅和她一起；这使犹达斯成了叛徒；这败坏了犹太人的首领，他们收受贿赂，成了贼的伙伴。这带来了千万场战争，使道路充满鲜血，城市充满哀号和哭泣。这使筵席变为不洁，筵席变为诅咒，使食物充满了过犯；因此保禄称它为\Quote{偶像崇拜}：\BibleRef{哥 3:5}，即便如此也没有阻止人们。为什么他称它为\Quote{偶像崇拜}？许多人拥有财富，却不敢使用，而是奉献它，原封不动地传下去，不敢碰触，好像是什么奉献之物。如果他们有时被迫这样做，就觉得好像做了非法的事。此外，如同希腊人精心照料他的雕像，你也把你的金子托付给门和闩；准备箱子代替神龛，把它存放在银器皿中。但你不像他拜雕像那样拜它？然而你对它显示了所有的关注。\par

再者，他宁愿放弃自己的眼睛或生命，也不愿放弃他的偶像。同样，那些爱慕黄金的人也是如此。有人说：\Quote{可是，我并不朝拜黄金。}他说，那人也不朝拜偶像，而是朝拜住在其中的魔鬼；同样，你虽不朝拜黄金，但你却朝拜那从黄金的景象和你的贪欲中扑向你灵魂的魔鬼。因为爱财的贪欲比恶神更可怕，许多人服从它胜过服从偶像。因为后者在许多事上不顺从，但在此事上，他们一切都顺从，无论它吩咐什么，他们都听从。它说什么？\Quote{与所有人争战，}它说，\Quote{与所有人为敌，不认识本性，蔑视天主，将你自己祭献给我，}而他们都服从。他们向偶像献上牛和羊，但贪欲说：把你的灵魂献给我，那人就听从。你看到它有怎样的祭台，接受怎样的祭献吗？贪婪的人不能继承天主的国，但即便如此，他们也不害怕。\BibleRef{格前 6:10}然而，这种欲望比所有其他欲望都更软弱，它不是天生的，也不是本性的（因为如果是，起初它就会被安置在我们里面）；但在起初没有黄金，也没有人渴望黄金。但若你愿意，我将告诉你这祸害是怎样进入的。每人嫉妒前人，人就加重了这疾病，而先得到的人激起那本无欲望的人。因为当人看到辉煌的房屋、广大的土地、成群的奴隶、银器以及大量的衣物时，他们用尽一切手段去超越别人；这样，第一批人成为第二批人的原因，而第二批人又是后来人的原因。如果他们能清醒持重，就不会成为别人的（邪恶）教师；然而这些人也没有借口。因为还有别人轻视财富。有人说：\Quote{谁轻视它们呢？}可怕的是，因为邪恶如此普遍，这似乎已变得不可能，甚至没有人相信人能行善。那么，我是否要提到许多在城市和山区的人呢？这又有什么益处？你不会因他们的榜样而变得更好。此外，我们的讲话现在没有这个目的，要你们空乏自己的财产：我希望你们能如此；然而，既然这担子对你们太重，我不强迫你们；我只劝你们不要贪图别人的东西，要分施一些自己的东西。我们会找到许多这样的人，他们满足于自己的东西，照顾自己的，并靠诚实的劳动生活。我们为什么不争相效仿他们呢？让我们想想那些已经过去的人。他们的产业不是仍然立着，只保存了他们的名字吗？某人的浴场，某人的郊区别墅和住所？当我们看到它们时，不是立刻叹息，想到他忍受了怎样的劳苦，进行了怎样的掠夺吗？现在他无处可见，而别人在他的产业中享乐，那些人他从未预料会如此，甚至可能是他的敌人，而他正受最严厉的惩罚。这些事也等待我们；因为我们必定会死，必定会面对同样的结局。告诉我，这些人招来了多少愤怒、多少花费、多少仇敌？而收益是什么？是永不止息的惩罚，没有安慰；不仅在活着时，而且去世后，被众人指责？什么？当我们看到许多人家里收藏的偶像时，我们不是更要哀哭吗？先知说得对：\Quote{的确，人生活，只是幻影，白白操劳。}\BibleRef{咏 38:11}因为对这类事的忧虑确实是操劳，是操劳和多余的烦恼。但在永久的居所却不是这样，在那些帐幕里不是这样。这里一个人劳苦，另一个人享受；但在那里，每个人将拥有自己的劳苦，并接受百倍的报酬。让我们奋力争取那产业吧，在那里为自己预备房屋，使我们能在基督耶稣我们的主内安息，愿光荣归于祂和父及圣神，世世无穷。阿们。\par

\SourceRecord{Translated by Charles Marriott.；Nicene and Post-Nicene Fathers, First Series；Vol. 14.；Edited by Philip Schaff.；Buffalo, NY: Christian Literature Publishing Co.,；1889.；<http://www.newadvent.org/fathers/240165.htm>.}

% structure: p0001 source=h1 target=section reason=page-title

\section{论\WorkTitle{若望福音}第六十六篇讲道}

% structure: p0002 source=h2 target=subsection reason=relative-heading-rank; base=h2

\subsection{\BibleRef{若 12:8}}

1. 财富常将那些不谨慎的人抛入毁灭，权势也是如此；前者引人贪婪，后者引人骄傲。试看，被统治的犹太群众是如何健全，而他们的统治者是如何败坏；因为群众中前者相信了基督，圣史们不断断言，说：\Quote{群众中有许多人信了祂}\EditorialNote{若 7:31-48}；但那些属于统治者的，却不信。他们自己说，不是群众，\Quote{难道有统治者信了祂吗？}但有一个人怎么说？\BibleQuote{不知道天主的群众是该受咒诅的}\BibleRef{若 7:49}；他们称信徒为咒诅的，却称自己这些杀人者为聪明。在此处也一样，看见奇迹后，许多人都信了；但统治者不仅不满足于自己的恶行，还试图杀害拉匝禄。假设他们试图杀害基督是因为祂违犯了安息日，因为祂使自己与圣父平等，并因为你们所提到的罗马人，那么他们对拉匝禄有什么指控，而要杀害他呢？难道领受恩惠是犯罪吗？你看见他们的意志是多么嗜杀？然而祂行了许多奇迹；但没有一个像这个奇迹那样激怒他们，不是瘫子，不是瞎子。因为这个奇迹在本质上更奇妙，并且在许多其他奇迹之后施行，看见一个死了四天的人行走说话，是一件奇异的事。诚然，在庆节上将庄严的集会与凶杀混在一起，真是光荣的行为。此外，在一种情况下，他们认为可以以安息日来指控祂，从而拉拢群众；但在这里，既然他们找不到祂的过错，便对那个被治好的人下手。因为在这里，他们甚至不能说祂反对圣父，因为祈祷堵住了他们的口。既然他们不断加给祂的指控被驳倒，而奇迹是明显的，他们就急忙去谋杀。所以，他们也会对瞎子这样做，若不是他们能在安息日上吹毛求疵。此外，那人默默无闻，他们把他赶出圣殿；但拉匝禄是个显赫之人，显然，因为许多人来安慰他的姊妹们；而且奇迹在众人眼前施行，最为奇妙。为此，所有的人都跑去看。这件事刺痛了他们，因为当庆节正在进行时，所有人都离开庆节去到伯达尼。因此，他们下手要杀他，并以为他们什么也不怕，他们是如此嗜杀。为此，法律在开始时就这样宣布：\BibleQuote{不可杀人}\BibleRef{出 20:13}；先知也以此指控他们：\BibleQuote{他们的手沾满了血}\BibleRef{依 1:15}。\par

然而，祂不再在犹太公开行走，而退到荒野之后，怎么又公开进入呢？祂通过退隐平息了他们的怒气，当他们平静时，祂来到他们当中。此外，前呼后拥的群众足以使他们陷入痛苦；因为没有什么神迹像拉匝禄的神迹那样吸引百姓。另一位圣史说，他们把衣服铺在祂脚下\BibleRef{玛 21:8}，并且\BibleQuote{全城都惊动了}\BibleRef{玛 21:10}；祂以如此大的荣耀进入。祂这样做，是预示了一项预言，又应验了另一项；同一个行为是一项预言的开始和另一项预言的结束。因为\BibleQuote{欢呼吧，因为你的君王来到你这里，是柔和的}\BibleRef{匝 9:9}，这属于祂，应验了预言；但骑坐驴驹却是预示未来事件的行为，即祂要使不洁的外邦民族服从于祂。\par

但其他圣史怎么说，祂派遣门徒，并说：\BibleQuote{解开驴和驴驹}\BibleRef{玛 21:2}，而若望却没有提到这类事，只说\Quote{找到一头小驴，就骑在上面}？因为很可能是两种情况都发生了：祂在驴驹被解开后，门徒正牵着它来的时候，找到（驴驹）并骑了上去。他们拿了棕榈树和橄榄树的枝条，把衣服铺在路上，表明他们现在对祂的看法比先知更高，并且说：\par

% structure: p0006 source=h2 target=subsection reason=relative-heading-rank; base=h2

\subsection{\BibleRef{若 12:13}}

你看见吗？这最令他们窒息，即所有人都确信祂不是天主的敌人？而祂说祂由圣父而来，这最使百姓分裂。但这是什么意思，\par

% structure: p0008 source=h2 target=subsection reason=relative-heading-rank; base=h2

\subsection{\BibleRef{若 12:15}}

因为他们的君王大多是不义贪婪之人，曾将他们交在敌人手中，败坏百姓，使他们臣服于仇敌；\Quote{鼓起勇气吧，}它说，\Quote{这位并非如此，而是柔和谦卑的}；正如驴驹所显示的，因为祂进入时没有军队随行，只有一头驴。\par

% structure: p0010 source=h2 target=subsection reason=relative-heading-rank; base=h2

\subsection{\BibleRef{若 12:16}}

2. 你看见吗？他们对许多事无知，因为祂没有启示给他们？因为当祂说：\BibleQuote{拆毁这座圣殿，三天之内我要把它重建起来}\BibleRef{若 2:19}，那时门徒也不明白。另一位圣史说：\BibleQuote{这话向他们隐藏了}\BibleRef{路 18:34}，他们不知道祂必须从死者中复活。如今，这话对他们隐藏是有理由的（因此另一位圣史说，当他们一次又一次听到时，他们忧愁沮丧，这是因为他们不明白关于复活的话）；隐藏起来是有理由的，因为这对他们太高深了：但为什么驴驹的事没有启示给他们？因为这也是件大事。但你应观察圣史的智慧，他如何不耻于显扬他们先前的无知。他们知道经上记载了，却不知道是记载关于祂的。因为如果祂是一位君王，将要遭受这样的事，并如此被出卖，这会冒犯他们。此外，他们不能立刻理解祂所说的王国；因为另一位圣史说，他们认为这话是关于现世的王国说的。\BibleRef{玛 20:21}\par

% structure: p0012 source=h2 target=subsection reason=relative-heading-rank; base=h2

\subsection{\BibleRef{若 12:17}}

因为如此多的人不会突然改变，除非他们相信了这奇迹。\par

% structure: p0014 source=h2 target=subsection reason=relative-heading-rank; base=h2

\subsection{\BibleRef{若 12:19}}

现在，在我看来，这话似乎是那些感觉正确却不敢直言的人说的，他们接着指出结果来约束别人，仿佛他们在尝试不可能的事。这里他们又称群众为\Quote{世界}。因为圣经惯于以\Quote{世界}之名既指受造物，也指活在邪恶中的人；一处说：\BibleQuote{他按数目引出他的万象}\BibleRef{依 40:26}；另一处说：\BibleQuote{世界不恨你们，却恨我}\BibleRef{若 7:7}。必须确切知道这些，免得因词语的含义而给异端者以把柄。\par

% structure: p0016 source=h2 target=subsection reason=relative-heading-rank; base=h2

\subsection{若望福音 12:20}

他们当时快要成为皈依者，正在过庆节。因此，当他们得知关于祂的传闻时，就说：\par

% structure: p0018 source=h2 target=subsection reason=relative-heading-rank; base=h2

\subsection{若望福音 12:21}

斐理伯让位给安德肋，因为安德肋在他之前，并将此事告知他。但他也没有立刻行使权柄；因为他曾听过那句话：\BibleQuote{不要走进外邦人的路}\BibleRef{玛 10:5}。因此，与那门徒商议后，他将此事呈报给他的老师。因为他们两人都对祂说了话。但祂说什么呢？\par

% structure: p0020 source=h2 target=subsection reason=relative-heading-rank; base=h2

\subsection{若望福音 12:23-24}

\BibleQuote{时刻已到}是什么意思？祂曾说过：\BibleQuote{不要走进外邦人的路}（从而消除了犹太人所有无知的借口），并约束了门徒。因此，当犹太人继续不顺服，而他人想要来到祂面前时，祂说：\Quote{现在}，祂说，\Quote{是时候进行我的苦难了，因为一切都已成就。因为如果我们继续等待那些不听话的人，而拒绝这些甚至愿意来的人，这将不符合我们的温情关怀。} 既然如此，祂即将允许门徒在十字架之后去外邦人那里，看见他们争先恐后，祂便说：\Quote{是时候走向十字架了。}因为祂不允许他们更早去，好为他们作证。直到犹太人因他们的行为拒绝了祂，直到他们钉死了祂，祂没有说：\BibleQuote{你们要去使万民成为门徒}\BibleRef{玛 28:19}，而是说：\BibleQuote{不要走进外邦人的路}\BibleRef{玛 10:5}，又：\BibleQuote{我被派遣，只是为了以色列家迷失的羊}\BibleRef{玛 15:24}，又：\BibleQuote{拿儿女的饼丢给小狗吃是不对的}\BibleRef{玛 15:26}。但当他们恨祂，恨到要杀祂时，在他们拒斥祂的情况下继续坚持就是多余的了。因为他们拒绝了祂，说：\BibleQuote{除了凯撒，我们没有君王}\BibleRef{若 19:15}。因此，当祂离开他们时，他们已离开了祂。所以祂说：\BibleQuote{我多少次愿意聚集你的子女，你却不愿意}\BibleRef{玛 23:37}。\par

\BibleQuote{一粒麦子如果不落在地里死了}是什么意思？祂指的是十字架，因为当他们看到，正当希腊人也来到祂面前时，祂却被杀，他们可能不会困惑。祂对他们说：\Quote{这件事特别促使他们前来，并将增加对我的宣讲。} 然后，既然祂不能很好地用言语说服他们，祂便通过实际经验来证明这一点，告诉他们谷粒的情况：当它死了时，它结出更多的果实。\Quote{现在}祂说，\Quote{如果种子是这样，那么对我来说更是如此。} 但门徒们不明白所说的话。因此，福音作者不断提到这一点，为他们后来的逃跑开脱。保禄在谈论复活时也提出了同样的论据。\par

3. 不信復活的人將有什麼藉口呢？因為這行動每天都在種子、植物和我們自身生育中實踐。首先種子必須死去，然後生育才會發生。但總之，當天主做任何事情時，推理是無用的；因為祂如何從無中造出我們呢？這話是對那些自稱相信\WorkTitle{聖經}的基督徒說的；但我也要從人的推理中說些別的話。有些人生活於惡習，有些人生活於美德；而在惡習中的人，許多在繁榮中達到了極高年壽，許多有美德的人卻忍受了相反的情況。那麼，每個人何時會得到他所應得的？在什麼時候呢？\Quote{是的，}有人說，\Quote{但沒有身體的復活。} 他們不聽\Person{保祿}說：\BibleQuote{這可朽壞的必須穿上不朽壞的。}\BibleRef{格前 15:53} 他所說的不是靈魂，因為靈魂是不朽壞的；況且，\Quote{復活}是指那跌倒的，跌倒的是身體。但你為什麼堅持沒有身體的復活呢？對天主來說不可能嗎？但這樣說完全是愚蠢的。是不體面嗎？為什麼不體面——那曾分擔勞苦和死亡的可朽壞者，也應分享冠冕？因為若是不體面，它起初就不會被創造，基督也不會重新取回肉身。但為證明祂重新取回了肉身並使祂復活，請聽祂所說的話：\BibleQuote{伸過你的手指來}\BibleRef{若 20:27}；以及\BibleQuote{你們看，鬼神是沒有骨頭和肉筋的。}\BibleRef{路 24:39} 但為什麼祂又復活了\Person{拉匝祿}，如果沒有身體復活更好？祂為什麼這樣做，將它列為奇蹟和恩惠？祂為什麼供給食物呢？所以，親愛的，不要被異端者欺騙：因為有復活，有審判，但那些不願為自己的行為交賬的人否認這些。因為這復活必須像基督的復活一樣，因為祂是初果，是死者中的首生者。但如果復活就是靈魂的淨化、脫離罪惡，而基督沒有犯罪，祂如何復活呢？如果祂也犯了罪，我們又如何脫離了詛咒呢？現在祂說：\BibleQuote{這世界的首領來了，在我身上一無所有}？\BibleRef{若 14:30} 這是宣稱祂無罪的話。因此，按照他們的說法，祂要麼沒有復活；要麼為了復活，祂在復活前犯了罪。但祂既復活了，也沒有犯罪。所以祂在身體中復活了，這些邪惡的道理只不過是虛榮的產物。讓我們逃避這疾病。因為經上說：\BibleQuote{不良的交往敗壞良好的習俗}\BibleRef{格前 15:33} 這些不是宗徒的道理；\Person{馬爾西翁}和\Person{瓦倫廷}是新近發明它們的。所以，親愛的，讓我們逃避它們，因為純潔的生活在道理敗壞時毫無益處；另一方面，若生活敗壞，純正的道理也無益。外邦人是這些觀念的父母，而那些異端者從外邦哲學家那裡接受後，將它們養育起來，斷言物質是未被造的，以及許多這類事情。因此，正如他們斷言除非有某種未被造的基底物質，否則不可能有工藝師一樣，他們也否認復活。但我們不要理會他們，因為知道天主的能力是全能的。我們不要理會他們。這話是對你們說的；因為\Emph{我們}不會拒絕與他們交戰。但一個赤手空拳的人，即使落在弱者中間，即使他更強壯，也容易被擊敗。如果你們曾留意\WorkTitle{聖經}，如果你們每天磨練自己，我就不會勸你們逃避與他們的戰鬥，而是會勸你們與他們搏鬥；因為真理是堅強的。但既然你們不知道如何使用\WorkTitle{聖經}，我害怕鬥爭，恐怕他們抓住你們毫無防備，將你們打倒。因為沒有什麼比失去聖神幫助的人更軟弱的了。如果這些異端者使用外邦人的智慧，我們不應欽佩，而應嘲笑他們，因為他們使用了愚蠢的教師。因為那些人關於天主或創造未能找到任何健全的道理，而我們中間的寡婦所知道的東西，\Person{畢達哥拉斯}還不知道，卻說靈魂變成灌木、魚或狗。對這些人，請問，你們應該聽從嗎？這樣做怎麼可能合理呢？他們在本地是了不起的人物，留著美麗的鬈髮，裹著斗篷；他們的哲學就到這裡為止；但如果你往裡面看，卻有塵土和灰燼，沒有健全的東西，反而\BibleQuote{他們的喉嚨是敞開的墳墓}\BibleRef{詠 5:9}，充滿了不潔和腐敗，他們所有的道理（充滿）了蟲子。例如，他們中的第一個人說水是天主，他的繼承者說是火，另一個人說是空氣，他們墮落到有形體的東西；請問，我們應該欽佩這些甚至從未想過無形體的天主的人嗎？即使他們後來獲得了這思想，也是在埃及與我們的人民交談之後。但為了不給你們帶來太多混亂，我們在此結束我們的講話。因為如果我們開始向你們陳述他們的道理，以及他們關於天主、關於物質、關於靈魂、關於身體所說的話，將會伴隨許多嘲笑。他們甚至不需要我們指控，因為他們互相攻擊；那個寫書反對我們關於物質的人，自己滅亡了。因此，為了不白白耽誤你們，也不編織語言的迷宮，我們離開這些事，勸你們緊緊抓住對\WorkTitle{聖經}的傾聽，不要無目的地爭論言詞；正如\Person{保祿}也勸勉\Person{弟茂德}\BibleRef{弟後 2:14}，儘管他充滿了許多智慧，擁有施行奇蹟的能力。現在讓我們服從他，離開無聊的事，抓住真正的工作，我指的是兄弟之愛和款待；讓我們重視施捨，好能得到所應許的美物，藉著我們的主耶穌基督的恩寵和慈愛，願光榮歸於祂，直到永遠。阿們。\par

\SourceRecord{Translated by Charles Marriott.；Nicene and Post-Nicene Fathers, First Series；Vol. 14.；Edited by Philip Schaff.；Buffalo, NY: Christian Literature Publishing Co.,；1889.；<http://www.newadvent.org/fathers/240166.htm>.}

% structure: p0001 source=h1 target=section reason=page-title

\section{若望福音讲道第六十七篇}

% structure: p0002 source=h2 target=subsection reason=relative-heading-rank; base=h2

\subsection{若望福音 12:25-26}

1. 现世生命是甘甜的，充满许多快乐，但并非对所有人，而是对那些执着于它的人。因为，若有人仰望天堂，看见那里的美好事物，他便会立即轻视这生命，不认为它有什么价值。正如当没有见到更美之物时，一物的美丽便受赞赏，但若有更美的出现，前者便遭轻视。因此，若我们选择注视那美丽，默观那里王国的光辉，我们便会很快从当前的锁链中解脱；因为这种对现世事物的同情，正是一种锁链。请听基督为吸引我们而说的：\BibleQuote{爱惜自己生命的，必要丧失生命；在现世憎恨自己生命的，必要保全生命直到永生；若有人服事我，就该跟随我}；以及\BibleQuote{我在哪里，我的仆人也在哪里。}这些话看似谜语，其实不然，而是充满极大的智慧。然而，\BibleQuote{爱惜自己生命的，必要丧失生命}是怎么回事呢？当他满足它不正当的欲望，当他在不该之处放纵它时。因此有人劝告我们说：\BibleQuote{不要随从你灵魂的欲望}\BibleRef{德 18:30}；因为这样你必毁灭它，因为它引人离开通向德行的道路；同样，相反地，\BibleQuote{在现世憎恨它的，必要拯救它。}但\BibleQuote{凡憎恨它的}是什么意思呢？就是当它命令有害之事时，不顺从它的人。主没有说\Quote{不顺从它的}，而是说\BibleQuote{凡憎恨它的}；因为正如我们甚至不能忍受我们所恨之人的声音，也不愿喜悦地看他们，同样，当灵魂命令违反天主喜悦之事时，我们也必须猛烈地远离它。因为主当时正要向他们谈论许多关于死亡，即他自己的死亡的事，看见他们沮丧和气馁，便很强烈地说：\Quote{我说什么？如果你们不能勇敢地忍受我的死亡？不，如果你们自己不死，你们将一无所获。}也要注意他如何缓和话语。被告知爱生命的人必要死，这是一件非常痛苦和悲伤的事。我又何必提及古代？因为即使在现在，我们也会发现许多人为了享受现世生命而甘愿忍受任何痛苦，即使他们对将来之事深信不疑；当他们看见建筑、艺术品和发明时，便哭泣，发出这样的感叹：\Quote{人发明了多少事物，却终究成为尘土！对现世生命的渴望是多么强大！}因此，为了解除这些锁链，基督说：\BibleQuote{在现世憎恨自己生命的，必要保守它直到永生。}为使你们知道他说话是为劝勉他们，驱散他们的恐惧，请听下文。\par

\BibleQuote{若有人服事我，就该跟随我。}\par

主谈到死亡，并要求以行为来跟随。因为服事者当然必须跟随被服事者。且注意他是在何时对他们说这些话；不是在他们受迫害时，而是在他们自信时；当他们因众人的尊崇与关注而自以为安全时，当他们可以振作起来而听到：\BibleQuote{背起自己的十字架，跟随我}\BibleRef{玛 16:24}；也就是说，他说：\Quote{要时刻准备好面对危险、面对死亡、面对离开此世。}然后，在他说了难以忍受的话之后，他也摆出了奖赏。这奖赏是什么？就是跟随他，并与他同在；表明复活将接续死亡。因为他说：\par

\BibleQuote{我在哪里，我的仆人也在哪里。}\par

但基督在哪里？在天上。因此，让我们甚至在复活之前，在灵魂和心意上移居那里。\par

\BibleQuote{若有人服事我，圣父必爱他。}\par

为何主不说是\Quote{我}？因为他们当时对主尚未持正确看法，而对圣父有更高的看法。因为他们连主将要复活都不知道，又怎能对他想象任何伟大的事呢？因此主对载伯德的儿子们说：\BibleQuote{不是我可以给的，而是为我父所预备的，就赐给谁}\BibleRef{谷 10:40}，然而审判的正是他。但在此处，他也确立了他真实的儿子身份。因为圣父接納他们，如同接納自己儿子的仆人一样。\par

% structure: p0010 source=h2 target=subsection reason=relative-heading-rank; base=h2

\subsection{若望福音 12:27}

\Quote{但这肯定不是鼓励他们甚至赴死之人的话。}不，这正是大大鼓励他们之人的话。因为恐怕他们说：\Quote{他既免于死亡的痛苦，便轻易地哲学化谈论死亡，且自己毫无危险地劝诫我们，}他强调，虽然他感受到死亡的痛苦，但因死亡有益，他并不拒绝。但这些属于经世，而非神性。因此他说：\BibleQuote{现在我心神烦乱}；因为若非如此，他所说的话与他说的\BibleQuote{父啊，救我脱离这时辰}有什么关联呢？而且他如此烦乱，甚至寻求从死亡中得解救，如果至少有可能逃脱的话。这些都是他人性的软弱。\par

2. \Quote{但}他说，\Quote{当我求释放时，我没有什么可说的。}\par

\BibleQuote{我正是为此来到这时辰。}\par

\Quote{好像他曾说：\InnerQuote{虽然我们困惑，虽然我们烦乱，但不要逃避死亡，因为即使我现在烦乱，也不说要逃避；因为应当承受将要来临的事。我不说：救我脱离这时辰，}但说什么呢？}\par

% structure: p0015 source=h2 target=subsection reason=relative-heading-rank; base=h2

\subsection{若望福音 12:28}

\Quote{雖然我的煩惱催逼我說這話，我卻說相反的話：\InnerQuote{光榮你的名，}意思是，從此以後領我走向十字架}；這極大顯示了祂的屬人性，以及一個不願死、卻戀棧現世生命的本性，證明祂並非不受人的情感影響。因為正如飢餓或睡覺並非過錯，同樣，戀慕現世生命也非過錯；基督確實有無罪的身體，卻並非沒有自然需要，否則就不是身體了。藉這些話祂也教導了另一件事。是甚麼呢？就是：倘若我們處於痛苦和恐懼中，即使那時也不應退避擺在我們面前的事；而藉著說\Quote{光榮你的名}，祂表明祂為真理而死，稱這行動為\Quote{天主的榮耀}。這事在十字架之後發生了。世界即將被歸化，承認天主的名，並事奉祂——不僅聖父的名，也包含聖子的名；但關於後者，祂仍保持沉默。\par

\Quote{因此有聲音從天上來說：\InnerQuote{我已經光榮了它，還要再光榮它。}}\par

祂在何時\Quote{光榮了它}？藉著先前所成就的事；而\Quote{我要再光榮它}是在十字架之後。那麼基督說了甚麼呢？\par

% structure: p0019 source=h2 target=subsection reason=relative-heading-rank; base=h2

\subsection{若望福音 12:30}

他們以為是打雷了，或有天使向祂說話。他們怎能這樣想呢？那聲音難道不清楚分明嗎？確實清楚，但對那些粗俗、屬肉體且遲鈍的人來說，它很快就飛逝了。有些人只捕捉到聲響，另一些人知道聲音是清晰的，卻不明白其意義。基督說：\Quote{這聲音不是為了我，而是為了你們的緣故。}祂為甚麼這樣說？祂是針對他們不斷斷言祂不出於天主的話而說的。因為那位被天主光榮的，怎能不是出於那藉著祂而被光榮的天主呢？實際上，這聲音正是為此而來。因此祂自己說：\Quote{這聲音不是為了我，而是為了你們的緣故，}\Quote{不是為了我可以從中學到我所不知的事（因為我知曉一切屬於聖父的事），而是為了你們的緣故。}當他們說：\Quote{有天使向祂說話}，或\Quote{打雷了}，而不留意祂時，祂說：\Quote{這是為你們的緣故}，好叫你們即使如此也被引導去探究這些話的意義。但他們受了激動，卻仍不探究，儘管他們聽聞此事與他們有關。因為對於一個不知為何發出這聲音的人來說，那聲音自然顯得模糊不清。\Quote{這聲音是為了你們的緣故。}你看見這些卑微的情況是為他們而發生，並非因為聖子需要幫助嗎？\par

% structure: p0021 source=h2 target=subsection reason=relative-heading-rank; base=h2

\subsection{若望福音 12:31}

這與\Quote{我已經光榮了，還要光榮}有何關聯？關係密切，且完全和諧。因為當天主說\Quote{我要光榮}時，祂表明了光榮的方式。是甚麼呢？就是那一位應當被推翻。但\Quote{這世界的審判}是甚麼？就好像祂說：\Quote{將有一個審判台和報應。}如何並以甚麼方式呢？\Quote{祂殺了第一個人，因為發現他有罪，（因為\InnerQuote{罪惡藉著死亡進入了}\BibleRef{羅 5:12}）但在祂身上，祂沒有找到罪。那麼，為甚麼祂撲向我，將我交於死亡呢？為甚麼祂將毀滅我的意念放入猶達斯心中呢？}（別告訴我這是天主的安排，因為這不屬於魔鬼，而是祂的智慧；目前先探究那惡者的意向。）\Quote{那麼，在我身上，世界如何被審判呢？}將要如此說，彷彿開庭一樣，對撒殫說：\Quote{你看，你殺了所有世人，因為你發現他們有罪。但你為甚麼殺了基督？難道不清楚你是枉然做的嗎？}因此，在祂身上，全世界將得到伸冤。但為使這點更清楚，我舉例說明。假設有一個殘忍的暴君，對所有落入他手中的人帶來無數災禍。如果這樣一個人與一位君王或君王的兒子交戰，並不義地殺了他，那麼他的死將有能力為其他人報仇。假設有一個人向他的債務人討債，毆打他們並把他們投入監獄；然後出於同樣的魯莽，他把一個不欠他的人也帶進同一監獄：這人將為他對別人所作的事而受罰。因為那個人將毀滅他。\par

3. 在聖子的事上也是如此；因為魔鬼對我們所作的那些事，將因他對基督所行的膽大妄為而受罰。為表明祂暗示這一點，請聽祂說的話：\Quote{現在這世界的元首將被趕出去，}\newline{}\Quote{藉著我的死亡。}\par

% structure: p0024 source=h2 target=subsection reason=relative-heading-rank; base=h2

\subsection{若望福音 12:32}

也就是說，\Quote{連外邦人也包括在內。}為使沒有人問：\Quote{如果祂比你還強，祂怎能被趕出去呢？}祂說：\Quote{祂並不更強；祂怎能強過那位吸引別人歸向祂的呢？}祂所說的不是復活，而是超越復活的事：\Quote{我要吸引眾人歸向我。}因為如果祂說\Quote{我要復活}，還不明確他們會相信；但藉著說\Quote{他們將相信}，兩者同時得到證明：既證明這件事，也證明祂必須復活。因為如果祂繼續死亡，並且只是一個凡人，沒有人會相信。\Quote{我要吸引眾人歸向我。}（\BibleRef{若 6:44}）那麼祂為何說聖父吸引人呢？因為當聖子吸引時，聖父也吸引。祂說：\Quote{我要吸引他們}，好像他們被暴君拘留，不能憑自己單獨接近祂，並逃脫那抓住他們之人的手。在另一處，祂稱之為\Quote{劫掠}：\Quote{沒有人能搶奪壯士的財物，除非先捆綁那壯士，然後才搶奪他的財物。}（\BibleRef{瑪 12:29}）祂說這話是為證明祂的能力；在那裡祂稱為\Quote{劫掠}的，在這裡稱為\Quote{吸引}。\par

知道这些事情以后，让我们振作起来，让我们光荣天主，不仅靠我们的信德，也靠我们的行为，因为否则就不是光荣，而是亵渎。因为一个不洁的外邦人并不像腐化的基督徒那样亵渎天主。所以我恳求你们做一切事，使天主得到光荣；因为经上说：\Quote{祸哉，} 经上说：\Quote{那个因他使天主的名受亵渎的仆人，}（而哪里有 \Quote{祸哉，} 每一刑罚和报复就立刻跟随而来，）\Quote{但那使这名得光荣的人是有福的。} 让我们不要像在黑暗中一样，而要避免一切罪过，特别是那些伤害他人的罪过，因为天主最因这些被亵渎。当我们被命令施舍给人，却抢夺别人的财物时，我们还有什么宽赦？我们还有什么救恩的希望？你如果没有喂养饥饿的人，就要受罚；但如果你还剥夺了一个有衣服穿的人，你将得到什么样的宽赦？这些事我将永不止息地说，因为今天没有听见的人也许明天会听见，明天不注意的人也许后天会被说服；即使有人如此固执而不被说服，但在审判时我们无需为他们交账。我们已尽了自己的本分；愿我们永不因我们的话而羞愧，你们也不要掩面，但愿众人都能勇敢地站在基督的审判台前，使我们也能因你们而欢喜，并且因你们在基督耶稣我们的主内被认可，而对我们自己的过失有些补偿。愿光荣与祂（基督）一同归于圣父及圣神，至于永远。阿们。\par

\SourceRecord{Translated by Charles Marriott.；Nicene and Post-Nicene Fathers, First Series；Vol. 14.；Edited by Philip Schaff.；Buffalo, NY: Christian Literature Publishing Co.,；1889.；<http://www.newadvent.org/fathers/240167.htm>.}

% structure: p0001 source=h1 target=section reason=page-title

\section{\WorkTitle{若望福音第六十八篇讲道}}

% structure: p0002 source=h2 target=subsection reason=relative-heading-rank; base=h2

\subsection{若 12:34}

1. 欺骗是极易识破且软弱无力的，即使外面涂抹万千色彩。因为正如那些粉刷朽墙的人不能通过粉刷使墙壁坚固一样，说谎者也容易被发现，正如当时犹太人的情形。当基督对他们说：\Quote{当我从地上被举起来时，便要吸引众人归向我}，他们中有人说：\Quote{我们从法律中听说，基督永远长存；你怎么说人子必须被举起来？这人子是谁？} 连他们也知晓基督是永生的那一位，生命无终。因此他们也明白他所指的含义；因为圣经中常在同一处提及苦难与复活。例如依撒意亚将二者并列，说：\BibleQuote{他被牵去待宰，有如羔羊}\BibleRef{依 53:7}，以及随后的一切。达味在第二篇圣咏及许多其他地方也将这两件事相连。族长雅各伯在说\BibleQuote{他卧下，如雄狮蹲伏}之后，又加上\BibleQuote{如幼狮，谁能唤醒他？}\BibleRef{创 49:9} 他同时显示了苦难与复活。但这些人想要使他沉默，并表明他不是基督，却因这事实而承认基督永远长存。请看他们的恶行；他们不是说\Quote{我们听说基督既不受苦也不被钉十字架}，而是说\Quote{他永远长存。} 然而即便所提及的这一点也并不构成真正的反对，因为苦难并不妨碍他的不朽。由此我们可见，他们明白许多疑点，却故意犯错。因为他先前曾论及死亡，如今他们在此听到\Quote{被举起来}，便猜出是指死亡。于是他们说：\Quote{这人子是谁？} 这也是他们诡诈的行径。有人说：\Quote{请别以为我们是针对你说这话，别断言我们是出于敌意反对你，看，我们不知你指的是谁，却仍陈述我们的意见。} 基督如何回应？为使他们沉默，并表明苦难并不妨碍他永远长存，他说：\par

% structure: p0004 source=h2 target=subsection reason=relative-heading-rank; base=h2

\subsection{若 12:35}

这表明他的死亡是一种转移；因为太阳的光不会消灭，而是隐退片刻后又重现。\par

\BibleQuote{当你们还有光的时候，应该行走。}\par

他这里指的是什么时期？是指整个现世生命，还是指受难之前？我个人认为两者都包括，因为因他不可言喻的慈爱，甚至受难后也有许多人相信。他说这些话是为催促他们信从，正如他先前所说：\BibleQuote{短短的时候，我和你们同在。}\BibleRef{若 7:33}\par

\BibleQuote{在黑暗中行走的，不知道往哪里去。}\par

例如，就连现在，犹太人做了多少事却不自知，仿佛在黑暗中行走？他们以为自己走对了路，实则背道而驰；守安息日，尊重法律和关于食物的规条，却不知往哪里走。因此他说：\par

% structure: p0010 source=h2 target=subsection reason=relative-heading-rank; base=h2

\subsection{若 12:36}

即\Quote{我的孩子们。} 然而起初福音作者说：\BibleQuote{他们不是由血统，也不是由肉情，也不是由男欲，而是由天主生的}\BibleRef{若 1:13}，即由圣父所生；而这里他自己被称为生养他们者，为叫你们明白圣父与圣子的行动是同一的。\BibleQuote{耶稣说完了这些话，}就离开他们，隐藏了。\par

为何他现在\Quote{隐藏}？他们没有拿起石头打他，也没有像先前那样毁谤他；那么他为何隐藏？他行走在人心深处，知道他们的怒气虽未言说却激烈，知道它沸涌并充满杀意，他没有等到它付诸行动，而是隐藏自己，以平息他们的恶意。请看福音作者如何暗示这种情绪；他紧接着说：\par

% structure: p0013 source=h2 target=subsection reason=relative-heading-rank; base=h2

\subsection{若 12:37}

2. 什么\Quote{这么多}？就是福音作者所省略的那些。这从下文也可清楚看出。因为他退去让步后，又回到他们身边，以谦卑的态度对他们说：\BibleQuote{信我的，不是信我，而是信那派遣我来的。}\BibleRef{若 12:44} 请看他的做法。他以谦卑温和的言辞开始，投靠圣父；然后他又提高语调，见他们恼怒便退去；然后再到他们那里，又以谦卑的话开始。他在何处这样做？岂不处处如此？例如，请看他在开始时说什么：\BibleQuote{我照我所听见的而审判。}\BibleRef{若 5:30} 然后以更高昂的语调说：\BibleQuote{就如圣父唤起死者并使他们复生，同样圣子也随己意使人们复生。}\BibleRef{若 5:21} 又说：\Quote{我不审判你们，另有审判者。} 然后他又退去。后来来到加里肋亚，他说：\BibleQuote{不要为那可朽坏的食粮劳碌}\BibleRef{若 6:27}，并在说了许多关于自己的大事——他从天降下、赐予永生——之后，又退隐。在帐棚节他也来临并做同样的事。人们可以见到他不断这样变化他的教导：时而临在，时而缺席，时而谦卑，时而高亢的言谈。他在这里也是如此。\BibleQuote{他虽然在他们面前行了这么多神迹，}却说\BibleQuote{他们仍然不信他。}\par

% structure: p0015 source=h2 target=subsection reason=relative-heading-rank; base=h2

\subsection{若 12:38}

% structure: p0016 source=h2 target=subsection reason=relative-heading-rank; base=h2

\subsection{若 12:39-41}

这里再次观察，\Quote{因为}和\Quote{说}并非指他们不信的原因，而是指事件本身。因为并非\Quote{因为}依撒意亚说了，他们才不信；而是因为他们将不信，他才说了。那么，为什么福音作者不这样表达，反倒使不信源于预言，而非预言源于不信呢？后面他更明确地说：\Quote{因此他们不能信，因为依撒意亚说过。}他希望借此用许多证据建立圣经无误的真实性，并使依撒意亚所预言的没有落空，而是照他所说的发生了。免得有人说：\Quote{基督为什么来？难道祂不知道他们不会听从祂吗？}他便引述也知道此事的先知。但祂来了，好使他们对自己的罪没有借口；因为先知所预言的事，是作为必定发生的事而预言的；如若不是必定发生，他就不能预言它们；而它们之所以必定发生，是因为这些人无可救药。\par

若\Quote{他们不能}是代替\Quote{他们不愿}，不必惊奇，因为祂在另一处也说：\BibleQuote{能领受的，就让他领受罢。}\BibleRef{玛 19:12}同样，在许多地方，祂惯常将选择称为能力。又说：\BibleQuote{世界不能恨你们，却是恨我。}\BibleRef{若 7:7}在普通谈话中也能见到这种用法；例如当一个人说：\Quote{我不能爱这个人或那个人}，他是将意志的力量称为能力。又说：\Quote{这个人或那个人不能成为好人。}先知怎么说？\BibleQuote{雇士人能改变他的皮肤吗？或豹子能改变它的斑点吗？}\BibleRef{耶 13:23}祂并不是说行善对他们是不可能的，而是因为他们不愿意，所以不能。福音作者的意思是说，先知不可能说谎；然而，他们不信并不是由此成为不可能。因为即使他们相信，他仍可以是真实的；因为如果他们将要相信，他就不会预言这些事。\Quote{那么，}有人说，\Quote{他为什么没有这样说？}因为圣经有某些这种惯用语，我们需要体谅其规律。\par

他说：\Quote{他看见了祂的荣耀时，就说了这些事。}谁的荣耀？圣父的。那么若望为什么论及圣子？保禄又为什么论及圣神？并非混淆位格，而是显示尊荣是同一的，他们才这样说。因为凡属圣父的，也属于圣子；凡属圣子的，也属于圣神。然而天主许多事由天使传达，却没有人说\Quote{正如天使所说}，而是说\Quote{正如天主所说}。因为通过天使职务所说的话，出自天主；但因此凡出自天主的，并不也出自天使。但在此处，若望说这些话是圣神的。\par

\Quote{并论及祂。}他说了什么？\BibleQuote{我看见主坐在崇高的宝座上，}\BibleRef{依 6:1}以及随后的话。因此，他在那里称那神视、烟雾、听闻不可言喻的奥迹、看见色辣芬、从宝座发出的闪电——那些神体不能直视的——为\Quote{荣耀}。\Quote{并论及祂。}他说了什么？他听见一个声音说：\BibleQuote{我要派遣谁呢？谁肯为我们去？我回答说：我在这里，请派遣我！祂说：你们听是听见，却不明白；看是看见，却不察觉。}\BibleRef{依 6:8-9}因为，\par

% structure: p0021 source=h2 target=subsection reason=relative-heading-rank; base=h2

\subsection{若 12:40}

这里又是一个问题，但若我们正确思考，便不是如此。因为正如太阳使弱者眼花，并非由于它本性的缘故，同样，那些不留意天主话语的人也是如此。这样，在法郎的事上，经上说天主硬了他的心，对一切抗拒天主话语的人也是如此。这是圣经的一种特殊说法，如同\Quote{天主任凭他们陷于邪恶的心思}\BibleRef{罗 1:28}和\Quote{将他们分散到列邦}，即允许、许可他们去。因为作者在这里并不是说天主亲自成就这些事，而是指出它们是通过他人的邪恶而发生的。因为当我们被天主舍弃时，我们就交于魔鬼；这样被交出去，我们就遭受无数可怕的事。因此，为恐吓听众，作者说：\Quote{祂硬了心}和\Quote{放弃了。}为显示天主不仅不放弃我们，甚至也不离开我们，除非我们愿意，请听祂说：\BibleQuote{是你们的罪恶使你们与我隔绝吗？}\BibleRef{依 59:2}。又说：\BibleQuote{远离祢的人必要灭亡。}\BibleRef{咏 72:27}欧瑟亚说：\BibleQuote{你忘记了你的天主的法律，我也要忘记你。}\BibleRef{欧 4:6}祂自己在福音中也说：\BibleQuote{我多少次愿意聚集你的子女——可是你们不愿意。}\BibleRef{路 13:34}依撒意亚又说：\BibleQuote{我来了，却没有人；我呼唤，却无人应允。}\BibleRef{依 50:2}祂说这些话，是显示我们开始离弃祂，成为自己丧亡的原因；因为天主不仅不愿离弃或惩罚我们，甚至当祂惩罚时，也是不情愿的；祂说：\BibleQuote{我不喜欢罪人丧亡，却喜欢他归正而生活。}\BibleRef{则 18:32}基督也为耶路撒冷的毁灭而哀恸，如同我们为朋友哀恸一样。\par

3. 知道这事，让我们尽一切努力，以免远离\Person{天主}，而让我们持守对我们灵魂的关心，并持守彼此之间的爱；不要撕裂我们自己的肢体（因为这是疯狂和失去理智者的行为），我们越看见有人心怀恶意，就越要善待他们。因为我们常常看到许多人在身体上患了难治或不治之症，却仍不停施治。有什么比脚或手的痛风更糟糕呢？难道因此就要截肢吗？绝不是，而是使用一切方法，使患者能得些安慰，因为我们无法消除疾病。对待我们的弟兄，也应如此；即使他们患了不治之症，我们仍要继续照顾他们，并彼此背负担子。这样，我们将满全\Person{基督}的法律，并获得所预许的美善，借着我们的主\Person{耶稣基督}的恩宠和慈爱，愿光荣归于他，及\Person{圣父}和\Person{圣神}，直到永远。阿们。\par

\SourceRecord{Translated by Charles Marriott.；Nicene and Post-Nicene Fathers, First Series；Vol. 14.；Edited by Philip Schaff.；Buffalo, NY: Christian Literature Publishing Co.,；1889.；<http://www.newadvent.org/fathers/240168.htm>.}

% structure: p0001 source=h1 target=section reason=page-title

\section{若望福音讲道集 69}

% structure: p0002 source=h2 target=subsection reason=relative-heading-rank; base=h2

\subsection{若 12:42-43}

1. 我们必须避免所有败坏灵魂的情欲，但尤其要避免那些本身滋生众多罪过的情欲。我指的是诸如贪爱金钱。它本身固然是一种可怕的疾病，但因为它是一切祸害的根源和母亲，就变得更加严重了。爱虚荣也是如此。看，例如，这些人如何因着爱荣誉而脱离了信仰。经上说：\Quote{纵然许多首领也信从了祂，但为了犹太人的缘故，并不明认，免得被逐出会堂。}正如祂先前对他们所说的：\Quote{你们既然彼此寻求光荣，而不寻求唯独来自天主的光荣，怎么能信从呢？}\BibleRef{若 5:44}所以，他们不是首领，而是极度奴隶般的奴隶。然而，这种恐惧后来被消除了，因为在宗徒时代，我们找不到他们被这种情感所支配，因为在他们那时，首领和司祭都信了。圣神的恩宠降临后，使他们比金刚石更坚固。既然这就是当时阻止他们信从的原因，请听祂所说的话。\par

% structure: p0004 source=h2 target=subsection reason=relative-heading-rank; base=h2

\subsection{若 12:44}

祂仿佛说：\Quote{你们为什么害怕信从我？信德通过我达于圣父，不信也是如此。}看，祂如何从各方面显示祂本性的不变性。祂没有说\Quote{信从\Emph{我的话}}，免得有人认为祂指的是祂的话语；这种情况在人身上或许可以这样说，因为信从宗徒的人并非信从他们，而是信从天主。但为使你们知道祂在这里所讲的是对祂本性的信从，祂没有说\Quote{谁信从我的话}，而是说\Quote{谁信从我}。\Quote{可是，}有人说，\Quote{为什么祂从未反过来说过\InnerQuote{凡信从圣父的，不是信从圣父，而是信从我}呢？}因为他们会回答说：\Quote{看，我们信从圣父，却不信从祢。}他们的心志仍然过于软弱。无论如何，在与门徒交谈时，祂的确这样说了：\Quote{你们信从圣父，也当信从我}\BibleRef{若 14:1}；但见这些人当时太软弱，听不进这样的话，祂便用另一种方式引导他们，表明不信从祂就不可能信从圣父。为使你们不认为这些话是出自人，祂又加上：\par

% structure: p0006 source=h2 target=subsection reason=relative-heading-rank; base=h2

\subsection{若 12:45}

那么，难道天主是身体吗？绝不然。祂这里所说的\Quote{看见}是指心智的看见，从而表明同质性\OriginalTerm{同质}{Consubstantiality}。而\Quote{谁信从我}是什么意思呢？就好像有人说：\Quote{谁从河里取水，不是从河里取，而是从源头取}；或者，这个比喻与我们现在谈论的事情相比，也显得太弱了。\par

% structure: p0008 source=h2 target=subsection reason=relative-heading-rank; base=h2

\subsection{若 12:46}

因为圣父在旧约和新约中都普遍被称为\Quote{光}，所以基督也使用了同样的名称；因此保禄也称呼祂为\Quote{光辉}\BibleRef{希 1:3}，这是从祂这里学来的。祂在此显示出祂与圣父的亲密关系，以及祂们之间没有分离，如果祂说信从祂不是信从祂自己，而是归于圣父的话。祂称自己为\Quote{光}，是因为祂拯救人脱离谬误，驱散心中的黑暗。\par

% structure: p0010 source=h2 target=subsection reason=relative-heading-rank; base=h2

\subsection{若 12:47}

2. 免得他们认为祂因无能而放过轻慢祂的人，因此祂说：\Quote{我来不是为审判世界。}然而，为了使他们不会因此变得更疏忽，当他们得知\Quote{信从的得救，不信的受罚}之后，请看祂如何又给他们摆出了一个可怕的审判庭，接着说道：\par

% structure: p0012 source=h2 target=subsection reason=relative-heading-rank; base=h2

\subsection{若 12:48}

\Quote{如果圣父不审判任何人，而祢又不是来审判世界的，那么谁审判他呢？} \Quote{我所讲的话，那话要审判他。}因为他们说过：\Quote{祂不是从天主来的。}祂说这话，是为了\Quote{他们那时就不能再说这些话了，而是我现在所说的话，要充当控告者，定他们的罪，断绝一切推诿之词。} \Quote{我所讲的话。}什么样的话？\par

% structure: p0014 source=h2 target=subsection reason=relative-heading-rank; base=h2

\subsection{若 12:49}

当然，这些话是为了他们的缘故而说的，好使他们没有任何借口。如果不是这样，那么祂比依撒意亚有什么更多呢？因为依撒意亚也说了同样的话：\Quote{主天主赐给我受教者的舌，使我知道怎样用言语扶助疲乏的人。}\BibleRef{依 50:4}比耶肋米亚更多呢？因为耶肋米亚受差遣时也受了默感。\BibleRef{耶 1:9}那么厄则克尔呢？因为他在吃了书卷之后，也这样说话。\BibleRef{则 3:1}否则，将要听祂话的人就会被认为是祂知识的来源。因为如果祂是在受差遣时才领受了该说什么的命令，那么你们就会争辩说，在受差遣之前祂什么也不知道。还有什么比这种论断更不敬的呢？如果人照此理解基督的话，而不明白祂卑微言语的原因的话。可是保禄说，他和那些成为门徒的人都知道\Quote{什么是天主良善、可悦纳和完美的旨意}\BibleRef{罗 12:2}，难道圣子直到领受命令才知道吗？这怎么可能合理呢？你们难道没有看见，祂将自己的表达降到一种极度的谦卑，为的是既吸引那些人，又堵住后来者的口？这就是为什么祂说出只适合凡人的话语，甚至这样也迫使我们避开言语的卑下，意识到这些话不属于祂的本性，而是适应听者的软弱。\par

% structure: p0016 source=h2 target=subsection reason=relative-heading-rank; base=h2

\subsection{若 12:50}

你看见这些话的谦卑吗？因为接受诫命的人并非自己的主人。然而祂说：\Quote{就如父使死人复活并赐予他们生命，同样，子也随意赐予生命。}\BibleRef{若 5:21}难道祂有权随意赐予生命，却没有权说祂愿意说的话吗？那么祂用这话所要表达的是：\Quote{这事没有自然的可能性，以至于祂说一套话，而我说另一套话。}\Quote{我知道祂的诫命是永生。}祂对那些称祂为欺骗者、声称祂来是要害人的人说了这话。但当祂说：\Quote{我不审判，}祂表明祂并非这些人丧亡的原因。借此祂几乎是明白地作证，当祂即将离开他们、不再与他们同在时，\Quote{我与你们交谈，说的不是出于我自己，而全是来自父。}为此，祂向他们讲话时，将自己限制在谦卑的言辞中，好能说：\Quote{直到末了，我向他们说了这最后的话。}那句话是什么？\Quote{父怎样对我说，我就怎样讲。}\Quote{如果我与天主对立，我就该说相反的话，即我不说任何中悦天主的事，好将光荣归于自己；但现在我已将所有事都归于祂，以至于不以任何事为我的。那么，当我说\InnerQuote{我接受了诫命}，并如此强烈地消除你们关于竞争的恶意猜疑时，你们为何还不信我呢？因为如同接受诫命的人，只要他们执行诫命、不伪造任何事，就除了他们的差遣者所愿意的以外，不能行或说任何事；同样，我也不能行或说任何事，除非按照我父的意愿。因为我所做的，祂也做，因为祂与我同在，并且\InnerQuote{父没有留下我独自一人}。}\BibleRef{若 8:29}你看见祂如何处处显示自己与生祂的那一位相连，且毫无分离吗？因为当祂说：\Quote{我不是凭自己来的，}祂这样说，并非剥夺自己的权能，而是消除一切疏离或对立。因为如果人是自己的主人，独生子更是如此。为表明这是真的，请听保禄所说：\Quote{祂空虚了自己，并为我们舍弃了自己。}\BibleRef{斐 2:7}但是，正如我所说，虚荣是一件可怕的事，非常可怕\BibleRef{弗 5:2}；这使这些人不信，并使另一些人信得不当，以至于那些出于仁爱为他们所说的话，他们竟转为不敬。\par

3. 那么，让我们永远逃避这个怪物：它多种多样，在财富、享乐、美貌上处处散布其独特的毒液。藉由它，我们处处超越必要的使用；藉由它，产生衣着的奢华和一大群仆役；藉由它，在房屋、衣着、餐桌上，必要的使用处处被轻视，而奢华占了上风。你愿享受荣耀吗？那就施舍，这样天使会称赞你，天主会接纳你。如今，赞美只止于金匠和织工，而你离去时没有冠冕，常常看到你受到诅咒。但如果你不把这些穿戴在身上，而是用来喂养穷人，那么从各方而来的喝彩将是巨大的，赞美将是巨大的。当你把财物给别人时，你才拥有它们；当你为自己保留时，你便没有它们。因为房屋是不可靠的宝库，但穷人双手是可靠的宝库。你为什么装饰你的身体，而你的灵魂却被忽视，为不洁所占据？你为什么不在你的灵魂上费像在身体上一样多的心思？你应当费更多的心思；但无论如何，亲爱的，我们应当对它给予同等的关怀。请告诉我，如果有人问你，你宁愿你的身体健康、体态良好、容貌超群，但穿着粗劣的衣服，还是愿身体畸形、百病缠身，却穿戴金银和华丽服饰？你岂不更喜欢那依赖于你本人本性的美，而非你所穿衣服的美吗？那么，你会为你的身体选择这样，而为你的灵魂却选择相反吗？当你拥有那丑陋、难看、黑暗的灵魂时，你以为从黄金饰品能获得什么吗？这是何等疯狂！将这种装饰移入内在，将这些项链挂在你的灵魂上。那些佩戴在身体上的东西对健康和美丽都无帮助，因为它们不能使黑变白，也不能使丑变美或好看。但如果你把它们挂在灵魂上，你会很快使它由黑变白，由丑变美，由难看变为可爱。这些话不是我的，而是主自己说的，祂说：\BibleQuote{尽管你们的罪如同朱红，我要使它们变得洁白如雪}\BibleRef{依 1:18}；以及：\BibleQuote{施舍——一切就都为你们洁净了}\BibleRef{路 11:41}；藉着这样的心意，你不仅美化你自己，也美化你的丈夫。因为如果他们看到你摘下这些外在的装饰，他们就不会有巨大的花费需求；没有这需求，他们就会避免一切贪恋，并更倾向于施舍，而你也能大胆地给他们适宜的劝告。目前，你被剥夺了所有这样的权威。因为你会用什么口舌谈论这些事？你会用什么眼睛正视你的丈夫，向他乞求施舍的钱财，而你却把大部分钱花在遮盖身体上？那么，当你放下你的金饰时，你就能大胆地与你的丈夫谈论施舍了。即使你毫无所成，你已经尽了你的全部本分；但我宁愿说，妻子不可能不赢得丈夫，当她用实际的行动说话时。\BibleQuote{因为女人，你怎能知道你是否会救你的丈夫？}\BibleRef{格前 7:16}。所以，正如现在你要为你自己和他交账，同样，如果你摆脱这一切虚荣，你将获得双重的冠冕，戴上你的冠冕，与你的丈夫一同在那些纯净的岁月中凯旋，享受永久的福乐。愿我们都能获得这福乐，借着我们的主耶稣基督的恩宠和慈爱，愿荣耀归于祂，直到永远。阿们。\par

\SourceRecord{Translated by Charles Marriott.；Nicene and Post-Nicene Fathers, First Series；Vol. 14.；Edited by Philip Schaff.；Buffalo, NY: Christian Literature Publishing Co.,；1889.；<http://www.newadvent.org/fathers/240169.htm>.}

% structure: p0001 source=h1 target=section reason=page-title

\section{《若望福音》讲道集第七十篇}

% structure: p0002 source=h2 target=subsection reason=relative-heading-rank; base=h2

\subsection{若 13:1}

1. 保禄说：\Quote{你们该效法我，如我效法基督一样。}\BibleRef{格前 11:1} 为此，祂也取了我们的血肉之体，好借此教导我们德行。因为经上说：\Quote{天主派遣了自己的儿子，成为罪恶肉身的形状，并为罪恶，在肉身上定了罪案。}\BibleRef{罗 8:3} 基督自己又说：\Quote{你们跟我学习，因为我是良善心谦的。}\BibleRef{玛 11:29} 祂教导这事，不仅用言语，也用行动。因为人们称祂为撒玛黎雅人，说祂附有魔鬼，是骗子，并向祂投石；有时法利塞人派遣差役捉拿祂，有时又派遣阴谋者陷害祂；他们自己也不断侮辱祂，尽管找不出祂的过错，反而不断蒙受恩惠。即便如此，祂仍不懈地以言以行善待他们。当一个仆役打了祂的脸，祂说：\Quote{我若说得不对，你指证那里不对；若说得对，你为什么打我？}\BibleRef{若 18:23} 这是对那些仇恨和阴谋反对祂的人。现在让我们看看祂对门徒做了什么，或者不如说祂对叛徒显示了怎样的行为。那个最应当被恨的人，因为是门徒，同席共盐，见了奇迹，蒙受了如此大的恩惠，却行了比任何人都更恶劣的事——不是投石，也不是侮辱，而是出卖和交出祂——请看祂以何等友善的态度接待这人，洗他的脚；因为祂甚至以此想阻止他的恶行。其实，祂若愿意，本可以像对待无花果树一样使他枯萎，像劈开岩石一样把他劈成两半，像撕裂帐幔一样把他裂开；但祂不愿用强迫，而愿用自愿的方式使他离开他的企图。因此祂洗了他的脚；但连这样也没有使那个不幸而可怜的人感到羞耻。\par

经上说：\Quote{逾越节前，耶稣知道自己的时辰到了。} 不是那时\Quote{知道}，而是说祂早已\Quote{知道}要做这事。\Quote{祂要离开世界。} 福音作者尊称祂的死为\Quote{离开}。\Quote{祂既爱了世上属于自己的人，就爱他们到底。} 你看见吗？当祂即将离开他们时，显示了更大的爱？因为\Quote{祂既爱了，就爱他们到底}，表明祂没有省略任何一个深切爱者所当做的事。那么，为什么祂不从一开始就做这事呢？祂在最后行了最伟大的事，为的是加深他们的依恋，并为他们预先储备安慰，以应对即将降临的可怕之事。圣若望称他们为\Quote{属于自己的人}，是就个人依恋而言，因为他也称其他人为\Quote{自己的}，是就创造工程而言；如他说：\Quote{自己的人却不接受祂。}\BibleRef{若 1:11} 但\Quote{在世上}是什么意思？因为已死的人也是\Quote{自己的}，如亚巴郎、依撒格、雅各伯及类似的人，但他们不在世上。你看见祂是旧约和新约的天主吗？\Quote{爱他们到底}是什么意思？意思就是\Quote{祂持续不断地爱他们}，福音作者提到这事，作为大爱的确证。的确，在别处祂提到另一证验：为朋友舍命；但那时尚未实现。为什么祂\Quote{现在}做这事？因为在众人眼中祂显得更荣耀的时候，这事更为奇妙。此外，当祂即将离开时，祂给了他们不小的安慰，因为他们将要大大忧伤，祂借此也为忧伤引入安慰。\par

% structure: p0005 source=h2 target=subsection reason=relative-heading-rank; base=h2

\subsection{若 13:2}

福音作者惊奇地说这话，表明耶稣洗了那已经决意出卖祂的人的脚。这也证明那人的极大邪恶，就连共盐（这是最能抑制邪恶的事）也未能阻止他；就连直到最后一天，他的主仍容忍他，也未能使他改变。\par

% structure: p0007 source=h2 target=subsection reason=relative-heading-rank; base=h2

\subsection{若 13:3}

这里福音作者甚至惊奇地说：一位如此伟大、极其伟大、从天主而来又往天主而去、掌管万有者，竟做了这事，并且不嫌弃承担如此行动。依我看来，圣若望所说的\Quote{交付}，是指信徒的得救。因为当祂说：\Quote{我父将一切交给了我}\BibleRef{玛 11:27}，祂是指这种交付；正如在另一处祂说：\Quote{他们本是祢的，祢将他们赐给了我}\BibleRef{若 17:6}；又说：\Quote{除非蒙父所吸引，谁也不能到我这里来}\BibleRef{若 6:44}；以及\Quote{除非从天上赐给他}\BibleRef{若 3:27}。那么，福音作者或指这事，或指基督不会因此行动而有所减损，因为祂从天主而来，往天主而去，拥有一切。但你听到\Quote{交付}时，不可按人的方式理解，因为它显示祂如何光荣圣父，以及与祂的一致。因为正如圣父将一切交给祂，祂也将一切交给圣父。保禄宣告这事说：\Quote{当祂将王国交给天主父时。}\BibleRef{格前 15:24} 但圣若望在这里以更人性化的方式说了这话，显示祂对他们的大关怀，并宣告祂不可言喻的爱，即祂如今关心他们如同自己的；教导他们一切美善之母——谦卑，祂称之为德行的开始和终结。加上\Quote{祂从天主而来，往天主而去}不是没有理由的，而是为使我们知道祂做了与从那里来、往那里去者相称的事，践踏一切骄傲。\par

% structure: p0009 source=h2 target=subsection reason=relative-heading-rank; base=h2

\subsection{若 13:4}

2. 注意谦卑如何不仅藉洗涤，也藉另一种方式表现出来。因为祂不是在躺下之前，而是在他们都坐好之后，才起身。其次，祂不仅洗他们的脚，而且还脱去外衣。祂甚至未止于此，还束上毛巾。祂还不以此为足，亲自倒水（在盆中），没有叫别人倒；祂亲自做了这一切，以此表明我们在行善时必须如此行，不仅是形式而已，而要全力以赴。现在祂似乎首先洗了叛徒的脚，根据经文所说：\par

% structure: p0011 source=h2 target=subsection reason=relative-heading-rank; base=h2

\subsection{\BibleRef{若 13:5}}

% structure: p0012 source=h2 target=subsection reason=relative-heading-rank; base=h2

\subsection{\BibleRef{若 13:6}}

\Quote{用这些手，就是祢用来开了瞎子的眼睛、洁净了癞病人、并使死人复活的手吗？}因为这个问题非常强调；因此祂不需要说多于\Quote{祢}这个字；因为单单这个字就足以表达一切。有人可以合理地问，为什么其他人都没有阻止祂，只有伯多禄阻止，这是不小的爱和尊敬的标志。那么原因是什么？我看祂似乎先洗了叛徒的脚，然后来到伯多禄那里，其他人后来从他的案例得到教导。祂在伯多禄之前先洗了别人，这从经文\Quote{但祂来到伯多禄那里}可以清楚看出。不过，福音作者并不是一个激烈的控告者，因为\Quote{开始}这个词暗示了这一点。即使伯多禄是第一个，但叛徒作为一个冒失的人，很可能在首领之前就已坐下。因为他的冒失也从另一件事上显出来：当他同主在一个盘子里蘸食，被揭露后仍无悔意；而伯多禄尽管只在先前一次被责斥，且是因出自爱而说的话，却如此羞愧，甚至苦恼颤抖，求别人替他问问题。但犹达斯虽然屡次被揭露，却感觉不到。\BibleRef{若 13:24}因此当祂来到伯多禄跟前，伯多禄对祂说：\Quote{主，祢洗我的脚吗？}\par

% structure: p0014 source=h2 target=subsection reason=relative-heading-rank; base=h2

\subsection{\BibleRef{若 13:7}}

也就是说\Quote{你将会知道从这事中所得的利益何等巨大，这教训的益处，以及它如何能引导我们进入完全的谦卑。}那么伯多禄做了什么？他仍然阻拦祂，说：\par

% structure: p0016 source=h2 target=subsection reason=relative-heading-rank; base=h2

\subsection{\BibleRef{若 13:8}}

\Quote{我若不洗你，你就与我无份了。}那么那个热心如火的人说什么？\par

% structure: p0018 source=h2 target=subsection reason=relative-heading-rank; base=h2

\subsection{\BibleRef{若 13:9}}

他在恳求时激烈，在默从时更加激烈；但两者都出于爱。因为祂为什么不说明祂如此行的原因，反而加上威胁呢？因为伯多禄不会被说服。因为如果祂说：\Quote{容忍吧，因为藉此我劝你谦卑}，伯多禄会千百次承诺，只要他的主不做这事。但现在祂说了什么？祂说伯多禄最害怕、最恐惧的事，就是与祂分离；因为他不停地问：\Quote{祢往哪里去？}\BibleRef{若 13:36}所以他也曾说：\Quote{我愿为祢舍命。}\BibleRef{若 13:37}并且，如果在听了\Quote{我现在所做的，你如今不知道，后来必明白}之后，他仍然坚持，那么他若知道了（行动的含意）就会更加坚持。所以祂说\Quote{后来必明白}，因祂知道，如果他立刻知道，他还会抗拒。伯多禄并没有说\Quote{告诉我，好让我容忍你}，而是（更为激烈）他甚至不愿知道，反而抗拒祂，说：\Quote{祢永远不可洗我的脚。}但祂一威胁，他立刻缓和了语气。但\Quote{后来必明白}是什么意思？\Quote{后来？}什么时候？\Quote{当你们以我的名驱魔，当你们看见我被接到天上，当你们从圣神那里得知我坐在祂右边时，那时你们就会明白现在所行的事。}那么基督说了什么？当伯多禄说：\Quote{不但我的脚，连手和头也要洗}时，祂回答：\par

% structure: p0020 source=h2 target=subsection reason=relative-heading-rank; base=h2

\subsection{\BibleRef{若 13:10-11}}

\Quote{他们既已洁净，为什么祂还洗他们的脚？}为了让我们学习谦逊。因此祂不是来到身体的其他部位，而是来到那被认为是比其余更不体面的部位。但\Quote{那洗过澡的}是什么意思？它代替\Quote{那洁净的}。难道他们还没有从罪中得到解救，没有被圣神所认可，而罪仍掌权，咒诅的字据还在，祭品尚未献上，他们就是洁净的吗？那么祂如何称他们为\Quote{洁净的}？为了使你们不会因他们从罪中得到解救而认为他们洁净，祂加添说：看，\Quote{你们因我向你们所讲的话，已是洁净的了。}那就是说：\Quote{就这样你们是洁净的；你们接受了光明，从犹太人的错误中得到释放。因为先知也说：\InnerQuote{你们要洗濯，自洁，从我眼前除掉你们的恶行}\BibleRef{依 1:16}；所以这样的人是洗过澡而洁净的。}既然这些人已从灵魂中除去一切邪恶，并以纯洁的心与祂交往，因此祂按照先知的话说：\Quote{那洗过澡的，已是洁净的。}因为在那处，它指的不是犹太人实行的水\Quote{洗}，而是良心的洁净。\par

3. 那么我们也当洁净；学习行善。但什么是\Quote{善}？\BibleQuote{为孤儿伸冤，为寡妇辩护；来，我们彼此辩论——主说。}\BibleRef{依 1:7} 圣经中屡次提及寡妇和孤儿，但我们却不以为意。然而，请思量赏报是多么巨大。\Quote{虽然，} 经上说，\Quote{你们的罪似朱红，我将使它们白如雪；虽红如丹颜，我将使它们白如羊毛。}因为寡妇是无保护者，所以祂十分照顾她们。因为她们即便有权再婚，却出于对天主的敬畏，忍受寡妇的艰辛。那么，我们所有的人，无论男女，都要向她们伸出援手，使我们永不遭受寡妇的悲伤；或者，如果我们不得不遭受，让我们为自己积存大量的仁慈。寡妇的眼泪并非渺小，它能开启天堂本身。因此，我们不要践踏她们，也不要加重她们的苦难，而是要用各种方式帮助她们。如果我们如此行，我们将在自己周围筑起极大的安全，无论今生还是来世。因为不仅在现世，而且在彼世，她们将成为我们的辩护者，因我们对她们的恩惠而削去我们大部分的罪，并使我们得以在基督的审判台前坦然无惧。但愿我们都能获得这一切，藉着我们的主耶稣基督的恩宠和慈爱，愿光荣归于祂，直到永远。阿们。\par

\SourceRecord{Translated by Charles Marriott.；Nicene and Post-Nicene Fathers, First Series；Vol. 14.；Edited by Philip Schaff.；Buffalo, NY: Christian Literature Publishing Co.,；1889.；<http://www.newadvent.org/fathers/240170.htm>.}

% structure: p0001 source=h1 target=section reason=page-title

\section{\WorkTitle{若望福音讲道集第七十一篇}}

% structure: p0002 source=h2 target=subsection reason=relative-heading-rank; base=h2

\subsection{若望福音第十三章第一节}

一、亲爱的诸位，陷入罪恶的深渊是一件可悲的事，非常可悲；因为那时灵魂就难以复原。因此我们应当尽一切努力，根本不被罪恶缠住；因为不跌进去比跌进去之后再爬起来更容易。比如，看看犹达斯，当他投身于罪恶之后，他得到了多么大的帮助，然而即使这样，他也没有被扶起来。基督向他说：\Quote{你们中有一个是魔鬼}\BibleRef{若 6:71}；祂又说：\Quote{不是所有的人都信}\BibleRef{若 6:65}；祂还说：\Quote{我不是指你们全体说的，} 并且 \Quote{我知道我所拣选的是谁}\BibleRef{若 13:18}；但这些话没有一句触动他。如今，当祂洗了他们的脚，穿上衣服，又坐下时，祂说：\Quote{你们明白我给你们所做的吗？} 祂不再只对伯多禄说话，而是对他们所有人说话。\par

% structure: p0004 source=h2 target=subsection reason=relative-heading-rank; base=h2

\subsection{若望福音第十三章第十三节}

\Quote{你们称我。} 祂引用他们的判断，然后，为了不致让人以为这些话是出于他们的客气，祂接着说道：\Quote{因为我实在是的。} 祂通过引入他们的话，使之不令人反感，又通过祂自己证实这从他们而来的话，使之不被怀疑。\Quote{因为我实在是的，} 祂说。你看到了吗？当祂与门徒交谈时，祂更清楚地揭示属于祂自己的事。正如祂说：\Quote{不要称地上的人为师傅，因为你们的导师只有一位}\BibleRef{玛 23:8}，同样地，\Quote{也不要称地上的人为父。} 但这里所说的\Quote{一位}和\Quote{一位}并非只指圣父，也指祂自己。因为如果祂说话时排除自己，那么祂怎么能说\Quote{为使你们成为光明之子}呢？再者，如果祂只称圣父为\Quote{师傅}，那么祂怎么说\Quote{因为我实在是的}呢？并且，\Quote{你们的导师只有一位，就是基督}\BibleRef{若 12:26}？\par

% structure: p0006 source=h2 target=subsection reason=relative-heading-rank; base=h2

\subsection{若望福音第十三章第十四至十五节}

然而这并不相同，因为祂是主和师傅，而你们彼此是同伴仆人。那么\Quote{如同}是什么意思呢？\Quote{以同样的热忱。} 为此，祂从更伟大的行为中选取例子，好让我们至少能完成较小的行为。正如学校的老师把字母写得非常漂亮，好让学生们即使模仿得差一些也能学会。如今，那些唾弃同伴仆人的人在哪里？那些要求受人尊敬的人在哪里？基督洗了叛徒、亵渎者、窃贼的脚，而且是在背叛临近之时，并且尽管他已是无可救药，还让他分享自己的餐桌；而你们却趾高气扬，扬起眉头吗？\Quote{那么，让我们彼此洗脚吧，} 有人说，\Quote{那我们还得洗仆人的脚了。} 我们即使洗自己仆人的脚，又算什么大事呢？在我们这里，\Quote{奴隶}和\Quote{自由人}只是字面上的区别；但在那里，却是实际上的区别。因为按照本性，祂是主，我们是仆人，但即使如此，祂当时也不拒绝做这事。如今，如果我们不把自由人当作买来的奴隶那样对待，就已是值得满足的事了。在那一天，我们将说什么呢？如果我们自己接受了如此忍耐的证明之后，一点也不效仿，反而背道而驰，骄傲自大，不尽这债？因为天主使我们彼此成为负债者，祂自己首先这样做了，并且使我们成为较小数额的负债者。因为祂是我们的主，但我们这样做（如果我们真的这样做的话）是对我们的同伴仆人，这件事祂自己通过说\Quote{我，你们的主和师傅——你们也照样做}暗示了。按理自然应该接着说\Quote{你们做仆人的，更应当如何呢}，但祂把这留给听众的良心去判断。\par

二、但祂为什么\Quote{现在}做这事呢？将来他们要享受一些更大的尊荣，一些更小的。为了不让他们彼此高抬自己，像以前那样说\Quote{谁最大}\BibleRef{玛 18:1}，也不彼此动怒，祂通过说\Quote{即使你们非常伟大，也不应对你们的弟兄存有高傲之心}来压下他们所有人的高傲思想。祂没有提到那更大的行动，即\Quote{如果我洗了叛徒的脚，你们彼此洗脚又算什么大事呢？}但祂通过行动树立了榜样，然后把这交予旁观者的判断。因此祂说：\Quote{无论谁遵守并教导，这人在天国里将称为大的}\BibleRef{玛 5:19}；因为这就是\Quote{教导}一件事，实际上是通过行动去做。什么骄傲不应当被除去呢？什么愚蠢和傲慢不应当被消灭呢！那坐在革鲁宾之上的，竟洗了叛徒的脚，而你，人啊，你是尘土、灰烬、尘埃，你竟自高自大，趾高气扬！你该受何等大的地狱刑罚呢！那么，如果你渴望一种高尚的心境，来，我指给你路；因为你甚至不知道什么是高尚的心境。一个把眼前事物看得很伟大的人，是心灵卑微的人；因此，没有心灵的伟大就没有谦卑，也没有骄傲不是出于心灵的渺小。正如小孩子热衷于琐事，贪恋球、铁环和骰子，却无法想象重要的事；同样，一个真正有智慧的人，会把眼前的事物当作虚无（因此他既不愿自己去获取它们，也不愿从别人那里接受它们）；但一个不这样的人，则会怀着相反的情感，专注于蛛网、阴影和比这些更不真实的梦境。\par

% structure: p0009 source=h2 target=subsection reason=relative-heading-rank; base=h2

\subsection{若望福音第十三章第十六至十八节}

祂先前所说的话，在此又重提，为要使他们羞愧。\Quote{因为仆人并不大于主人，受派遣的也不大于派遣他的，既然这事我作了，你们更应当作。}然后，免得有人说：\Quote{为什么你现在说这些话？难道我们不知道吗？}祂就加上这句话：\Quote{我不是对你们说，因为你们不知道，而是要用行动证明你们所听的。}因为\Quote{知道}属于众人，但\Quote{去做}却不属于众人。为此祂说：\Quote{你们若实行了，便是有福的。}因此我不断对你们说同样的话，虽然你们知道，为的是使你们付诸实行。因为连犹太人也\Quote{知道}，但他们却不是\Quote{有福的}，因为他们不实行他们所知道的。\par

祂说：\Quote{我不是说你们全体。}啊，何等容忍！祂尚未揭露背叛者，而是遮掩此事，给他悔改的机会。当祂这样说：\Quote{同我吃饼的人，举脚踢我。}祂是揭露又不揭露。在我看来，祂说\Quote{仆人并不大于他的主人}，也是为了这个目的：倘若有人有时从家人或任何下等人那里受伤害，他们不应跌倒；以犹达斯为例，他享受了万千好处，却以相反的事回报他的恩主。为此祂又加上\Quote{同我吃饼的人}，而略过其他一切，祂提出那最能约束和使他羞愧的话：\Quote{他是由我喂养的，}祂说，\Quote{并共享我的餐桌。}祂说这些话，是教导他们要善待那些作恶于他们的人，即使这样的人仍不悔改。\par

但祂说了\Quote{我不是说你们全体}之后，为不使恐惧加于多人以上，祂终于区分出出卖者，这样说：\Quote{同我吃饼的人。}因为\Quote{不是全体}并没有指向某一个人，因此祂加上\Quote{同我吃饼的人}；向那可怜的人表明祂不是无知的被捉拿，而是明明知道；这件事本身最能约束他。祂没有说\Quote{出卖我}，而是说\Quote{举脚踢我}，意欲表示欺骗、背信和阴谋的隐秘。\par

3. 这些事记载下来，是要我们不怀恨那些伤害我们的人，反而责备他们，为他们哭泣；因为哭泣的合适对象不是那受苦的人，而是那行恶的人。贪得无厌的人、诬告的人，以及任何作恶的人，都加害自己最大，而加给我们最大益处，如果我们不报复的话。有这样的情况：有人抢了你；你曾为那伤害感谢并光荣天主吗？因那感谢，你获得了无数的赏报，正如他为自己收集了不可言喻的火。但若有人说：\Quote{那么，如果我不能自卫抵抗那冤枉我的人，因为我较弱小呢？}我要说，你能付诸实行的是：不满、不耐烦（因为这些事在我们能力之内）、祈祷反对那令你忧伤的人、说出千万咒骂反对他、对众人说他坏话。因此，那没有做这些事的人，甚至要因不行自卫而得赏报，因为显明即使他有能力，他也不会做。受伤的人，当灵魂渺小之时，用任何到手的武器自卫，以咒骂、辱骂、阴谋来对抗那伤害他的人。你却不仅不要做这些事，甚至要为他祈祷；因为如果你不做这些，反而为他祈祷，你就变得像天主一样。因为，\Quote{祈祷，}经上说，\Quote{为那迫害你们的人，好使你们成为你们在天之父的子女。}\BibleRef{玛 5:44} 你看见我们从别人的侮辱中得了最大的益处吗？没有什么比不以恶报恶更令天主喜悦？但我说什么呢？不以恶报恶？我们确实被命令以相反的事回报，即恩惠、祈祷。因此基督也以一切相反的事回报那将要出卖祂的人。祂洗他的脚，私下责备他，有节制地斥责他，照顾他，允许他分享祂的餐桌和祂的吻，即使这样他也没有变好；然而（基督）继续做祂自己的部分。\par

但是来吧，让我们甚至从仆人的榜样教导你们，并且（为了使教训更强而有力）从旧约中的榜样，好叫你们知道，当我们怀怨记仇时，我们毫无辩护的理由。那么，你们愿意我告诉你们关于\Person{梅瑟}的事，还是我们更往前追溯？因为所举出的例子越古老，我们就越被超越。\Quote{为什么这样说？}因为那时德行更困难。那些人没有成文的诫命，没有生活的典范，但他们的本性赤手空拳独自作战，被迫无压舱物地向各方漂浮。因此，当称赞\Person{诺厄}时，天主并非简单地称他为义人，而是加上\BibleQuote{在他那一代}\BibleRef{创 7:1}；意思是\Quote{在那时}，当时有许多障碍，因为在他之后有许多人闪耀，但他仍不亚于他们；因为在他自己的时代，他是完美的。那么，在\Person{梅瑟}之前，谁是忍耐的？有福而高贵的\Person{若瑟}，他因贞洁而闪耀，也同样因忍耐而闪耀。他被出卖时没有做过任何错事，反而是在服侍他人，作仆人，履行所有家仆的职责。他们用恶毒的控告加害于他，他没有为自己辩护，尽管他有父亲站在他一边。不仅如此，他甚至去旷野给他们送食物，当找不到他们时，他没有绝望或折返（虽然如果他愿意，他本有理由这样做），而是停留在野兽和那些野蛮人附近，保持着真正的兄弟情谊。再者，当他住在监牢里，被问及原因时，他没有说他们的坏话，只是说：\Quote{我没有做什么事}和\Quote{我是从希伯来地被偷来的}；此后，当他被立为主人时，他供养他们，并把他们从万千危险中解救出来。如果我们保持清醒，邻人的邪恶并不足以将我们从自己的德行中抛出。但那些人不像他；他们剥了他的衣服，试图杀他，并用他的梦来责备他，尽管他们甚至从他那里得到食物，还计划剥夺他的生命和自由。他们吃着，却不关心躺在坑里赤身的兄弟。还有什么比这种残忍更糟糕？他们难道不比许多杀人犯更坏吗？在这之后，他们把他拉上来，却将他交给万千死亡，卖给了那些正前往野蛮人之地的野蛮残暴之人。然而当他成为统治者时，他不仅赦免了他们的惩罚，甚至至少就他自己而言，免去了他们的罪，称所发生的事是天主的上智安排，而非他们的恶行；他对他们所做的事，并非出于记恨，而是在这一切中他都掩饰了，为了他兄弟的缘故。此后，当他看到他们紧靠他时，他立刻丢掉面具，放声大哭，拥抱他们，仿佛他接受了最大的恩惠，他，从前被他们谋害的，却将他们全部带到埃及，并以万千恩惠回报他们。那么，如果我们在法律之后，在恩宠之后，在增添了如此多的天上智慧之后，甚至不努力去与那位生活在恩宠和法律之前的人相比，我们还有什么借口？谁能拯救我们脱离惩罚？因为没有，没有比记仇更令人痛苦的事。而那欠一万塔冷通的人就表明了这一点；从他的债务有时不被要求偿还，有时又被要求偿还；不被要求，是由于天主的慈爱；被要求，则是由于他自己的邪恶，以及他对同仆的恶意。知道这一切，让我们宽恕邻人的过犯，并用相反的行为回报他们，好使我们也能从天主获得怜悯，藉着我们的主\Person{耶稣基督}的恩宠和慈爱，愿荣耀和权能归于他，至于无穷世。阿们。\par

\SourceRecord{Translated by Charles Marriott.；Nicene and Post-Nicene Fathers, First Series；Vol. 14.；Edited by Philip Schaff.；Buffalo, NY: Christian Literature Publishing Co.,；1889.；<http://www.newadvent.org/fathers/240171.htm>.}

% structure: p0001 source=h1 target=section reason=page-title

\section{\WorkTitle{若望福音讲道集 第七十二篇}}

% structure: p0002 source=h2 target=subsection reason=relative-heading-rank; base=h2

\subsection{若望福音 13:20}

1. 照顾天主的仆人所得到的酬报是大的，而且它自身就为我们结出果实。因为，经上说：\BibleQuote{凡接待你们的，就是接待我；凡接待我的，就是接待那差遣我来的} \BibleRef{玛 10:40} 那么，还有什么能与接待基督和祂的父相比呢？但这与前面所说的有什么关系呢？它与祂所说的\BibleQuote{你们若实行这些事，就有福了}有什么共通之处呢？加上\BibleQuote{凡接待你们的}？联系紧密，且十分和谐。看是如何。当他们即将出去遭受许多可怕的事时，祂以两种方式安慰他们：一是出于祂自己，一是出于他人。因为祂说：\Quote{因为如果你们真有智慧，常常记念我，并把我所说所行的一切放在心上，你们就能轻易忍受可怕的事。不仅如此，你们还会从所有人那里得到极大的关注。}祂在说\BibleQuote{你们若实行这些事，就有福了}时，点明了第一点；在说\BibleQuote{凡接待你们的，就是接待我}时，点明了第二点。因为祂为众人打开了所有的房屋，使他们既从自己行为的明智，又从那些愿意照料他们的人的热忱中，得到双重的安慰。当祂给这些人这些指示，如同给那些将要跑遍全世界的人时，想到叛徒被剥夺了这两样，既不会享受劳苦中的忍耐，也不会享受好客者的服侍，祂再次感到烦乱。福音作者为表明这一点，并显示祂是因叛徒而烦乱，便加上说：\par

% structure: p0004 source=h2 target=subsection reason=relative-heading-rank; base=h2

\subsection{若望福音 13:21}

祂再次使众人恐惧，因祂没有点名（叛徒）。\par

% structure: p0006 source=h2 target=subsection reason=relative-heading-rank; base=h2

\subsection{若望福音 13:22}

耶稣若把一切放在一个人身上，本可消除他们的恐惧，但加上\BibleQuote{你们中间有一个人}，却使众人不安。那么怎样呢？其余的人彼此对看，但一向热心的\Person{伯多禄}向\Person{若望}\BibleQuote{点头}。因为他先前曾受责备，当基督要给他洗脚时他曾拦阻，又因为他在各处被发现虽出于爱心却受责备；因此他害怕，既不保持沉默，也不说话，而想藉\Person{若望}得知消息。但有一个问题值得追问：为什么当众人都愁苦战栗，他们的领袖害怕时，\Person{若望}却像安闲的人靠在耶稣的怀里，不但靠着，甚至（躺）在祂的胸膛上？这不仅值得探究，下面的事也值得。是什么呢？他论到自己所说的话：\BibleQuote{耶稣所爱的}。为什么没有别人这样说自己呢？其实别人也被爱，但他比任何人都更被爱。若没有别人说过关于他的这话，而是他自己说的，这并不奇怪。\Person{保禄}也是如此，当有机会时，就这样说：\BibleQuote{我认识一个人，十四年前}；事实上他还说过其他不小的自夸。你以为这是小事吗？当他听见\BibleQuote{跟随我}，就立刻撇下网和父亲，跟随了祂；并且基督只带他和\Person{伯多禄}一同上山，\BibleRef{玛 17:1} 另有一次祂进到一间屋子里？\BibleRef{路 8:51} 他自己又毫无隐瞒地称赞\Person{伯多禄}，告诉我们基督说：\BibleQuote{伯多禄，你爱我比这些更深吗？} \BibleRef{路 21:15} 并且处处显明\Person{伯多禄}对他热情而高尚；例如，当他问\BibleQuote{主，这人将来怎样？}时，是出于极大的爱。但为什么没有别人说关于他的这话呢？因为若不是到了这里，他自己也不会说。因为若是他在告诉我们\Person{伯多禄}向\Person{若望}点头询问之后，没有再加什么，就会引起很大的疑惑，迫使我们追问原因。所以为了亲自解决这个困难，他说：\BibleQuote{他靠在耶稣怀里。}你以为听到\BibleQuote{他靠着}，并且他们的主允许他们如此大胆，就知道了一件小事吗？你若想知道原因，这行为是出于爱；因此他说\BibleQuote{耶稣所爱的}。我想\Person{若望}这样做还有另一个理由，即想表明他无罪，所以他公开说话且充满信心。再者，为什么他不在其他任何时间，而只在宗徒之长点头时才用这话呢？为叫你不以为\Person{伯多禄}向他是因他更伟大，他说这事是由于（耶稣对他）的大爱。但他为什么甚至躺在祂的怀里呢？他们当时对祂还没有形成高超的想法；此外，祂这样平息了他们的沮丧；因为此时他们的脸上很可能布满了愁云。若他们心里不安，面容更是如此。所以祂用言语和问题安抚他们，预先铺路，并让他靠在祂的胸前。也要注意他的谦逊：他不提自己的名字，只说\BibleQuote{祂所爱的}。正如\Person{保禄}，当他这样说时：\BibleQuote{我认识一个人，大约十四年前。}这时耶稣第一次指明叛徒，但甚至还不是指名；而是怎样呢？\par

% structure: p0008 source=h2 target=subsection reason=relative-heading-rank; base=h2

\subsection{若望福音 13:26}

连这（责备）的方式也是要使他羞愧的。他不尊重筵席，虽然同吃饼；就算这样；但从主手里蘸饼，谁能不被感化呢？然而却不能感化他。\par

% structure: p0010 source=h2 target=subsection reason=relative-heading-rank; base=h2

\subsection{若望福音 13:27}

嘲笑他的无耻。只要他还属于门徒的团体，他就不敢攻击他，而是从外面袭击；但当基督揭露他并把他分离出来时，他就无所畏惧地攻击了。把这样一个长久不改的性格留在内部是不合适的。因此祂从此把他赶出去，然后那另一个（魔鬼）在他被割除时抓住他，他离开他们，在夜间出去了。\par

\BibleQuote{耶稣对他说：朋友，你所要做的，就快做吧！}\par

% structure: p0013 source=h2 target=subsection reason=relative-heading-rank; base=h2

\subsection{若望福音 13:28}

3. 何等奇妙的麻木不仁！他怎能既不软化也不羞愧，反而更加无耻地\BibleQuote{出去}了？\BibleQuote{你所要做的，快做吧}不是命令或劝告的表达，而是责备，并表明主渴望纠正他，但因他不可救药，就任他去了。圣史说，\BibleQuote{坐席的人没有一个知道这事。}有人可能在此发现相当大的困难：当门徒们问\BibleQuote{是谁？}而主回答\BibleQuote{我蘸一点饼给谁，就是谁}时，他们仍然不明白；除非祂是悄悄说的，不让任何人听见。因为若望正是为此，斜靠在主的怀里，几乎贴着祂的耳朵问，以免叛徒暴露；基督也以同样的方式回答，以至于即使那时也没有揭穿他。尽管祂加重语气说\BibleQuote{朋友，你所要做的，快做吧}，他们仍然不明白。祂这样说，是为表明祂曾对犹太人论及自己死亡的话是真实的。因祂曾对他们说：\BibleQuote{我有权舍弃我的生命，也有权再取回它}；以及\BibleQuote{没有人能夺去我的生命。}\BibleRef{若 10:18}因此，当祂愿意保留时，无人能取；但当祂交出来时，行动就变得容易了。这一切祂在说\BibleQuote{你所要做的，快做吧}时暗含了。然而即使那时祂也没有揭露他，因为其他人可能会把他撕碎，或者伯多禄可能会杀了他。为此，\BibleQuote{坐席的人没有一个知道。}连若望也不知道？连他也不知道：因为他不可能预料到一个门徒会堕落到如此邪恶的地步。因为他们自己远离这种不义，所以不会怀疑别人有这样的行为。正如之前祂曾告诉他们\BibleQuote{我不是指你们所有人}\BibleRef{若 13:18}，却没有揭露那个人；同样，这里他们以为祂所说的是关于另一件事。\par

圣史说，他出去时正是\BibleQuote{黑夜}。\Quote{你为什么要告诉我时间呢？}这是为叫你认识他的鲁莽，连时间也未能阻止他的企图。然而即便如此，也没有使他完全暴露，因为其他门徒此时惊慌失措，被恐惧和巨大的痛苦所占据，他们不知道所说之话的真正原因，却以为主这样说，是为了让犹达斯拿些东西周济穷人。因祂非常关心穷人，也教导我们要在这事上倍加努力。他们这样想并非没有原因，而是\BibleQuote{因为他掌管钱囊}。然而似乎没有人把钱带到祂那里；经上曾说女门徒用她们的财物供应祂，但这里却从未提及。\BibleRef{路 8:3} 但那位吩咐门徒不要带口袋、金钱、棍杖的主，自己却带着钱囊以赒济穷人，这是为什么呢？为叫你知道，甚至那极其贫穷且被钉十字架的人，也应当在这事上非常谨慎。因为祂做了许多事，都是为我们教导而作的安排。门徒们于是以为祂说这话，是要犹达斯拿些东西给穷人；但连这也未能使他羞愧，因为主不愿直到最后一天还公开羞辱他。我们也应当这样做，不要宣扬同伴的罪过，即使他们不可救药。因为即在此之后，祂还亲吻了那来出卖祂的人，忍受了那样的行为，然后转向更为大胆的事——十字架本身，直至羞辱的死亡，并在那里再次彰显了祂的慈爱。在此祂称它为\BibleQuote{光荣}，向我们显示，若按天主的旨意而行，就没有什么可羞耻和可指责的事不会使走向它的人更加光明。至少，在犹达斯出去出卖祂之后，祂说：\par

% structure: p0016 source=h2 target=subsection reason=relative-heading-rank; base=h2

\subsection{\BibleBook{若望}{福音} 13:31}

祂这样激励门徒沮丧的思想，劝他们不仅不要灰心，甚至要喜乐。为此祂首先责备了伯多禄，因为一个曾经历死亡的人胜过死亡，乃是极大的光荣。这正是祂论及自己所说的：\BibleQuote{当我被举起时，你们将知道我是那一位}\BibleRef{若 8:28}；又说：\BibleQuote{拆毁这座圣殿}\BibleRef{若 2:19}；又说：\BibleQuote{除了约纳的征兆外，没有征兆赐给你们}\BibleRef{玛 12:39}。因为死后还能做比生前更大的事，这岂非莫大的光荣吗？为使复活得以相信，门徒们行了更大的事。但除非祂活着且是天主，这些人怎能因祂的名行这样的事呢？\par

第32节。\BibleQuote{天主将光荣祂} \BibleRef{若 13:32}。\par

\BibleQuote{天主将光荣祂在自己内}是什么意思？就是\Quote{藉着祂自己，而非藉着另一个。}\par

\BibleQuote{并且立刻光荣祂。}\par

4. 这就是说，\Quote{与十字架同时。} \Quote{因为不会在很长时间以后，}祂说，\Quote{祂也不会等待复活的遥远时刻，那时才显示祂的光荣，而是就在十字架上，祂的光荣就要显现。} 于是太阳昏暗了，磐石崩裂，圣所的帐幔裂为两半，许多睡了的圣人的身体复活了，坟墓有封印，卫兵在旁，而石头压在尸体上时，尸体却复活了；四十天过去，圣神的恩赐来到，他们都立刻宣讲祂。这就是\BibleQuote{要光荣祂在自己内，并且立刻光荣祂}；不是藉着天使或总领天使，也不是藉着任何其他权能，而是藉着祂自己。但祂又怎样藉着自己光荣祂呢？藉着为圣子的光荣作一切。然而圣子却成就了一切。你看见了吗？祂将祂自己所做的事归于圣父。\par

第33节。\BibleQuote{孩子们，我同你们在一起的时候不多了——就如我曾对犹太人说过：我所去的地方，你们不能去；现在我也这样对你们说。} \BibleRef{若 13:33}。\par

祂现在在晚餐后开始说忧伤的话。因为犹达斯出去后，不再是傍晚，而是黑夜。但由于他们即将很快到来，有必要将所有事都摆在门徒面前，使他们能记在心中；或者更确切地说，圣神会提醒他们一切。因为他们很可能忘记许多事，因为是初次听见，而且即将经历如此大的考验。那些被睡意压垮的人（正如另一位福音传道者所说，\BibleRef{路 22:45}），被沮丧占据的人，正如基督自己所说：\Quote{因为我向你们说了这些话，忧愁充满了你们的心}\BibleRef{若 16:6}，他们怎能准确记住这一切呢？那么为什么还要说这些话呢？这对他们对基督的看法不无小补：在日后被提醒时，他们确实知道这些事他们早已从基督那里听到过。但为什么祂先使他们的心灵沮丧，说：\Quote{我还有不多的时候与你们同在}？\Quote{对犹太人这样说是有道理的，但为什么你把我们和那些顽固的人归入同一类呢？}祂绝不是这样做的。\Quote{那么为什么祂说：\InnerQuote{正如我对犹太人说的}？}祂提醒他们，祂现在并非因为患难临头才警告他们这些事，而是祂从一开始就预知了，并且他们就是听见祂向犹太人说过这些话的见证人。因此祂也加上了\Quote{孩子们}一词，为叫他们听见\Quote{正如我对犹太人说的}时，不会以为这话对他们也是同样的意思。因此祂这样说并非要使他们沮丧，而是要安慰他们，免得危险突然临到而过度困扰他们。\par

\Quote{我去的地方，你们不能来。}祂显示祂的死亡是一种迁移，是向一个更美好的地方的转变，那地方不容纳可朽坏的身体。祂说这话，既是为激发他们对祂的爱，使之更炽热。你们知道，当我们看到最亲爱的朋友离我们而去时，我们的感情最热烈，尤其当我们看到他们去一个我们甚至不可能去的地方时。因此，祂说这些话使犹太人恐惧，却在门徒心中燃起渴望。\Quote{那地方就是如此，不仅他们不能去，就连你们，我最心爱的，也不能去那里。}在这里祂也显示了自己的尊荣。\par

\Quote{所以现在我向你们说。}为什么是\Quote{现在}？\Quote{对他们是一种方式，对你们是另一种方式}；也就是说，\Quote{不与他们一样。}但犹太人何时寻找祂，门徒又何时寻找祂？门徒是在他们逃离犹太人时，在他们遭受无法忍受的、难以形容的苦难（在城陷之时），当天主的愤怒从各方冲击他们的时候。因此祂对犹太人说话是因为他们的不信，\Quote{但对你们现在说话，是免得患难突然临到你们。}\par

第34节。\BibleQuote{我给你们一条新诫命。}\BibleRef{若 13:34}\par

因为他们听到这些话时，很可能感到困扰，好像要被遗弃似的，所以祂安慰他们，赋予他们那作为一切福乐根源和保障的东西——爱。祂好像说：\Quote{你们为我的离去忧伤吗？不，如果你们彼此相爱，你们就会更坚强。}那么为什么祂没有这样说？因为祂说的比这更有益于他们。\par

第35节。\BibleQuote{凭此众人将认出你们是我的门徒。}\BibleRef{若 13:35}\par

5. 借此祂同时表明这团体绝不会灭绝，因为祂给了他们一个标志。祂说这话时，叛徒已被从他们中间剪除。但祂为何称那条也包含在旧（盟约）中的诫命为新的？祂亲自以方式使之成为新的；因此祂加了\Quote{如同我爱了你们。}\Quote{我并没有先偿还你们善行的债务，而是我自己开始了，}祂说，\Quote{同样你们也该祝福你们所爱的人，即使你们不欠他们什么}；祂不提及他们将行的奇迹，而是使爱成为他们的特征。为什么呢？因为爱最能显明人的圣洁；它是一切德行的根基；我们众人主要也是因爱得救。因为\Quote{这，}祂说，\Quote{就是做门徒；这样众人都会赞美你们，当他们看见你们效法我的爱。}那么，奇迹不更能显明这一点吗？绝非如此。因为\Quote{许多人会说：主，我们不是因你的名字赶过魔鬼吗？}\BibleRef{玛 7:22}又有，当他们因魔鬼服从他们而欢喜时，祂说：\Quote{不要因为魔鬼服从了你们而欢喜，却要因为你们的名字记录在天上而欢喜。}\BibleRef{路 10:20}确实，爱征服了世界，因为爱先于奇迹；没有爱，奇迹也不会持久。爱立刻使他们成全，使众人一心一意。但若他们彼此分离，一切就都会丧失。\par

现在祂说这话，不仅是对他们说的，也是对一切信从祂的人说的；因为直到如今，使外教人失足的原因没有别的，就是缺少爱德。有人说：\Quote{但，他们也以我们不行奇迹为由反对我们。}但并非如此。有人说：\Quote{但宗徒们在哪里显示了他们的爱德？}你没看见\Person{伯多禄}和\Person{若望}形影不离，一同上圣殿吗？\BibleRef{宗 3:1}你没看见\Person{保禄}对他们也怀着同样的心情吗？你还怀疑吗？如果他们已经得了其他的福分，更不用说他们得了万善之母了。因为这是出于有德之灵魂的事；但哪里有不义，那植物就枯萎了。因为经上说：\BibleQuote{当不法之事增多时，许多人的爱德就会冷淡。}\BibleRef{玛 24:12}引外教人归正，奇迹不如生活的榜样；而正当的生活，没有什么能比爱德更促成它。因为那些行奇迹的人，他们甚至常称之为骗子；但对于纯洁的生活，他们却无法指责。在福音的信息尚未广传的时候，奇迹理所当然地受人惊叹，但现在人们必须藉着生活来赢得钦佩。因为在外教人眼中，没有什么比德行更能赢得尊重，没有什么比恶行更能冒犯他们。这理所当然。当他们中有一个人看见贪婪的人、掠夺者，竟劝别人行相反的事，当他看见那被命令要爱仇敌的人，对待自己的骨肉如同禽兽，他便会说那些话是胡言。当他看见有人害怕死亡，他如何能接受不朽的言论呢？当他看见我们喜爱管辖，被其他情欲奴役，他就会更坚定地持守自己的教义，对我们没有高的评价。我们，我们正是他们仍留在错误中的原因。他们早已谴责了自己的教义，也同样钦佩我们的教义，但他们却被我们的生活方式所阻碍。在言谈中随从智慧是容易的，他们自己中间有许多人做到了；但他们要求以行为来证明。\Quote{那么，让他们看看我们教会的古人吧。}但关于古人，他们完全不信；他们询问的是现世活着的人。因为经上说：\BibleQuote{把你因行为而生的信德显示给我。}\BibleRef{雅 2:18}但事实并非如此；相反，他们看见我们撕咬邻舍比任何野兽更凶残，便称我们为世界的祸害。这些事阻止了外教人，不让他们投向我们。因此我们也将为这些事受罚；不仅为自己所做的错事，而且因为天主的名被亵渎。我们还要沉溺于财富、奢侈及其他情欲多久呢？今后我们要舍弃它们。请听先知论及某些愚昧人所说的：\BibleQuote{让我们吃喝吧，因为明天我们就要死了。}\BibleRef{依 22:31}但在我们这种情况下，我们甚至不能说这样的话，因为\Quote{许多人}把属于众人的东西都聚集到自己身边。因此先知也斥责他们说：\BibleQuote{你们要独自住在地上吗？}\BibleRef{依 5:8}因此我害怕恐怕有严重的灾祸发生，我们招致天主对我们沉重的惩罚。为了避免此事，让我们留心各样德行，好使我们获得未来的福分，藉着我们的主\Person{耶稣基督}的恩宠和慈爱，藉着祂，并与祂一起，归于圣父及圣神，愿光荣现在、永远，及无穷世。阿们。\par

\SourceRecord{Translated by Charles Marriott.；Nicene and Post-Nicene Fathers, First Series；Vol. 14.；Edited by Philip Schaff.；Buffalo, NY: Christian Literature Publishing Co.,；1889.；<http://www.newadvent.org/fathers/240172.htm>.}

% structure: p0001 source=h1 target=section reason=page-title

\section{\WorkTitle{若望福音讲道集73}}

% structure: p0002 source=h2 target=subsection reason=relative-heading-rank; base=h2

\subsection{若 13:36}

爱是伟大的，甚至比火本身更强，它直上高天；没有任何障碍能阻挡它的撕裂之力。因此，最热忱的伯多禄听到主说\Quote{我去的地方，你们不能来}时，他怎么说呢？他说\Quote{主，祢往哪里去？}他说这话，与其说是出于求知的愿望，不如说是出于跟随的渴望。他不敢公开说\Quote{我去}，而是问\Quote{祢往哪里去？}基督回答的不是他的话，而是他的心思。因为他的渴望从基督的话中可以显明：\Quote{我去的地方，你现在不能跟随我。}你看见他渴望跟随祂，因此才问这个问题吗？当他听到\Quote{你后来要跟随我}时，他仍无法抑制渴望，尽管已获得美好的希望，他仍急切地说：\par

% structure: p0004 source=h2 target=subsection reason=relative-heading-rank; base=h2

\subsection{若 13:37}

当他摆脱了成为叛徒的恐惧，并被显明是祂自己的人后，他随后在其他人沉默时，自己大胆地问：\Quote{你怎么说，伯多禄？祂说\InnerQuote{你不能}，而你说\InnerQuote{我能}？因此你要从这次试探中知道，没有从上而来的冲动，你的爱算不了什么。}由此清楚可见，出于对他的关心，祂允许了那次跌倒。祂确实想通过最初的话教导他，但当他继续坚持自己的激烈态度时，祂并没有将他推入或强迫他否认，而是任由他独自一人，好让他认识自己的软弱。基督曾说过祂必须被出卖；伯多禄回答说：\Quote{主，愿这事远离祢，这事不会临到祢身上。}\BibleRef{玛 16:22}他受了斥责，却没有得到教诲。相反，当基督想要洗他的脚时，他说：\Quote{祢永远不可洗我的脚。}\BibleRef{若 13:8}又，当他听到\Quote{你现在不能跟随我}时，他说：\Quote{即使所有人都否认祢，我也不会否认祢。}既然他常因争辩而变得愚妄，耶稣接下来教导他不要反对祂。路加也暗示了这一点，他告诉我们基督说：\Quote{我已为你祈求，使你的信德不致丧失}\BibleRef{路 22:32}；意思是\Quote{使你最终不致沉沦。}从各方面教导他谦卑，并证明人性本身微不足道。但是，由于极大的爱使他容易争辩，祂现在使他冷静，免得日后当他领受治理世界的职务时，仍受此困扰，而是记念他所受的苦，从而认识自己。并且看他跌倒的猛烈程度：不是一次两次，而是他几乎失控，在短时间内三次说出否认的话，好让他知道，他的爱不如他被爱的程度。然而，对这如此跌倒的人，祂又说：\Quote{你爱我比这些更深吗？}因此，否认并非因他的爱冷淡，而是因为他失去了从上而来的援助。祂接纳了伯多禄的爱，但杜绝了由此产生的争辩精神。\Quote{因为如果你爱，你就该服从你所爱的那一位。我对你和与你在一起的人说过\InnerQuote{你们不能}，你为什么还要争辩？你知道违背天主意味着什么吗？既然你不愿通过这种方式知道我的话必然应验，你将在否认中明白这一点。}而且，这事在你看来更为难以置信。因为这事你甚至不明白，而那事你心里明白。然而，那未曾预料的事还是发生了。\par

\Quote{我愿为祢舍命。}因为他曾听见\Quote{没有比这更大的爱}，他立刻奋起，不知满足地渴望达到德行的最高峰。但基督为了表明惟有祂自己有权柄应许这些事，就说：\par

% structure: p0007 source=h2 target=subsection reason=relative-heading-rank; base=h2

\subsection{若 13:38}

即\Quote{现在}；只有片刻时间。祂说话时已是深夜，第一更和第二更已经过去。\par

% structure: p0009 source=h2 target=subsection reason=relative-heading-rank; base=h2

\subsection{若 14:1}

祂这样说，是因为他们听见时很可能会感到不安。因为如果他们团体中那位极为热忱的领袖，被告知在鸡叫以前要三次否认他的主，他们很可能预期要遭受某种巨大的打击，足以使刚硬的心也屈服。既然他们想到这些事时很可能会惊惶，请看祂如何安慰他们，说：\Quote{你们心里不要愁烦。}通过这第一句话显示祂天主性的能力，因为他们心里所存的，祂知道并显露出来。\par

\Quote{你们信天主，也当信我。}意思是：\Quote{一切危险都将越过你们，因为对我和我父的信德比临到你们的事更强大，不容任何恶事得胜。}然后祂接着说：\par

% structure: p0012 source=h2 target=subsection reason=relative-heading-rank; base=h2

\subsection{若 14:2}

正如祂安慰困惑的伯多禄说\Quote{但你后来要跟随我}，祂也同样将这希望的一瞥赐给其他人。为免他们以为这应许只赐给他一人，祂说：\Quote{在我父的家里有许多住处。}\par

\Quote{若不是这样，我早已告诉你们，我去是为你们预备地方。}\par

即\Quote{接纳伯多禄的同一处地方也将接纳你们。}因为那里有极多的住所，不能说他们需要预备。当祂说\Quote{你们现在不能跟随我}时，为免他们以为自己最终被隔绝，祂接着说：\par

% structure: p0016 source=h2 target=subsection reason=relative-heading-rank; base=h2

\subsection{若 14:3}

\Quote{我对此事如此关切，若不是早已为你们预备，我早已被交付了。}显示他们应当非常勇敢和自信。然后，为免祂的话好似在诱惑他们，而是让他们相信事情的确如此，祂接着说：\par

% structure: p0018 source=h2 target=subsection reason=relative-heading-rank; base=h2

\subsection{若 14:4}

你们看见祂证明这些事不是无故说的吗？祂说这些话，是因为祂在自己内知道，他们的灵魂现在渴望学习这事。因为伯多禄说这话，不是为了学习，而是为了跟随。但当伯多禄受责斥，基督声明那在当时看来不可能的事成为可能，而表面的不可能使他渴望确切知道这事时，祂就对其他人说：\Quote{你们也知道那条路。}因为正如当祂说\Quote{你们要否认我}，在任何人说话之前，就洞察他们的心，说\Quote{不要惊慌}，同样，此处祂说\Quote{你们知道}，揭示了他们心中的渴望，并亲自给他们提问的借口。而伯多禄说\Quote{祢往哪里去？}是出于极深的爱，多默则是出于怯懦。\par

% structure: p0020 source=h2 target=subsection reason=relative-heading-rank; base=h2

\subsection{若望福音 14:5}

他说：\Quote{地方，我们不知道，又怎能知道通往那里的路呢？}请注意他说话多么谦卑；他不说\Quote{告诉我们那地方}，而是说\Quote{我们不知道祢往哪里去}；因为大家早已渴望听到这事。如果犹太人听到（祂离开）时就在彼此议论，尽管他们想摆脱祂，那么那些希望永不与祂分离的人，岂不更渴望学习呢？所以他们害怕问祂，却仍因极大的爱与焦虑问了祂。基督怎么说呢？\par

% structure: p0022 source=h2 target=subsection reason=relative-heading-rank; base=h2

\subsection{若望福音 14:6}

\Quote{那么，当伯多禄问祂\InnerQuote{祢往哪里去}时，为什么祂不直接说\InnerQuote{我到父那里去，但你们现在不能来}？为什么祂绕了这么多圈子，穿插问题和答案？祂不告诉犹太人，这是有理由的；但为什么不告诉这些人呢？}祂确实曾对这些人，也对犹太人说过，祂来自天主，并要到天主那里去；现在祂比从前更清楚地说了同样的事。此外，祂对犹太人说得不那么清楚；因为如果祂说\Quote{你们若不藉着我，就不能到父那里去}，他们立刻就会认为这是夸夸其谈；但现在祂隐藏这事，使他们陷入困惑。\Quote{但为什么，}有人说，\Quote{祂对门徒和伯多禄都这样说呢？}祂知道伯多禄极大的热诚，而且因此他会更加追问和困扰祂；因此为了引开他，祂隐藏了这事。然后，在藉着隐晦和掩盖祂的言语而实现了祂所愿的之后，祂又揭示了这事。在说了\Quote{在我父的家里有许多住处}之后，祂又说\Quote{除非藉着我，没有人能到父那里去}。起初祂不愿告诉他们这事，以免使他们陷入更大的沮丧；但既然祂安慰了他们，就告诉他们。因为藉着责斥伯多禄，祂消除了他们许多沮丧；他们害怕受到同样的对待，就更加克制了。\Quote{我就是道路。}这是\Quote{除非藉着我，没有人能到父那里去}的证明；而\Quote{真理，生命}则是\Quote{这些事必定成就}的证明。\Quote{那么，在我这里没有虚假，因为我是\InnerQuote{真理}；如果我也是\InnerQuote{生命}，就连死亡也不能阻止你们到我这里来。此外，如果我是\InnerQuote{道路}，你们就不需要任何人引导你们；如果我也是\InnerQuote{真理}，我的话就不是谎言；如果我也是\InnerQuote{生命}，即使你们死了，也必得到我所告诉你们的。}现在，他们说祂是\Quote{道路}，他们明白并承认，但其余的事他们不知道。他们的确不敢说出他们不知道的事。尽管如此，他们从祂是\Quote{道路}获得了极大的安慰。祂说：\Quote{如果我有独一的权威将你们带到父那里，你们必定能到那里；因为不可能藉着别的路去。}但祂先说\Quote{除非父吸引他，没有人能到我这里来}；又说\Quote{我从地上被举起来时，我要吸引众人归向我}\BibleRef{若 12:32}；又说\Quote{除非藉着我，没有人能到父那里去}\BibleRef{若 14:6}；祂显示自己与生祂者同等。但为何在说了\Quote{我往哪里去，你们知道，那条路你们也知道}之后，祂又加上：\par

% structure: p0024 source=h2 target=subsection reason=relative-heading-rank; base=h2

\subsection{若望福音 14:7}

祂并不自相矛盾；他们确实认识祂，但并非应当地认识。他们认识天主，但尚不认识父。因为后来，圣神降临到他们身上，在他们内完成了所有的知识。祂说的是这样：\Quote{如果你们认识我的\OriginalTerm{本体}{Essence}和\OriginalTerm{尊威}{Dignity}，你们也会认识父的；从今以后，你们将认识祂，并且已经看见祂}（后者属于将来，前者属于现在），也就是说，\Quote{藉着我}。所谓\Quote{看见}，祂是指藉着理智知觉的知识。因为对于被看见的人，我们可以看见而不认识；但对于被认识的人，我们不能认识而又不认识。因此祂说\Quote{你们已经看见祂}；正如经上说\Quote{也被天使看见}\BibleRef{弟前 3:16}。然而，本体本身是看不见的；但经上说祂\Quote{被看见}，意思是按他们可能看见的程度。这些话是为了让你知道，看见祂的人认识生祂者。但他们看见祂，不是在本然的本体中，而是披着肉身。祂在其他地方常以\Quote{看见}代替\Quote{认识}；如祂说\Quote{心里洁净的人是有福的，因为他们要看见天主}\BibleRef{玛 5:8}。所谓\Quote{洁净}，祂不仅指那些免于私通的人，而且指免于一切罪恶的人。因为每一种罪都在灵魂上带来污秽。\par

3. 那么，让我们用一切方法擦去污秽。但首先是圣洗池洁净，此后还有其他许多各式各样的途径。因为天主是仁慈的，在此之后仍赐给我们各种和好的方法，其中首要的是藉着哀矜。\Quote{藉着哀矜，}经上说，\Quote{与信德的行为，罪恶得以洁净。} \BibleRef{德 3:30}藉着哀矜，我并非指由不义所支撑的哀矜，因为那不是哀矜，而是野蛮与残忍。剥去一人的衣服去给另一人穿上，有何益处呢？因为我们的行动应始于慈悲，但这是残忍。如果我们把从别人那里得来的一切都施舍出去，这对我们并无益处。匝凯就证明了这一点，他当时说，他藉着偿还四倍所贪得的来悦乐天主。 \BibleRef{路 19:8}但我们，当我们无度地掠夺，却只给出一丁点，就以为能使天主和好，实则更激怒祂。请告诉我，如果你从路边或巷子里拖来一头腐烂的死驴，带到祭台前，众人岂不要用石头砸你，视你为被诅咒和不洁的吗？那么，若我证明藉掠夺所得的祭献比这更污秽，我们还能有什么辩护呢？让我们假设某物是掠夺来的，难道它的气味不比死驴更臭吗？你愿知道罪恶的腐烂有多大吗？请听先知说：\Quote{我的创伤溃烂流脓。} \BibleRef{咏 37:5}你一方面用言语恳求天主忘记你的过犯，另一方面却用你自己所做的——抢劫、攫取，并把自己的罪过放在祭台上——使祂不断记起它们？但现在，这还不是唯一的罪，还有比这更严重的，那就是你玷污了圣人们的灵魂。因为祭台不过是一块被祝圣的石头，而圣人们却常把基督本身带在身上；你竟敢将任何这样的不洁送到那里去吗？\Quote{不，}有人说，\Quote{不是同一笔钱，而是其他的。}这是嘲弄，是儿戏。难道你不知道，若一滴不义掉进大量的财富中，整个就都被污秽了吗？正如一个人将粪土扔进纯净的泉水，使整个泉水变脏；同样，在财富的事上，任何非法得来的东西进入其中，都会使它们因自身的气味而沾染恶臭。我们进入圣堂时洗手，但心却不然。为什么？难道我们的手会发出声音吗？是灵魂在发言；天主看的是灵魂；身体的洁净毫无用处，如果灵魂是污秽的。如果你擦干净外面的手，而里面的手却不洁，有何益处？因为可怕的事、败坏一切善的事是这样：我们为琐碎之事恐惧，却不顾重要之事。用未洗的手祈祷是无关紧要的事；但用未洁净的心祈祷，这是万恶之极。请听对那忙于外在污秽的犹太人所言：\Quote{洗净你心中的罪恶，恶念存于你心里要到几时？} \BibleRef{耶 4:14}让我们也来洗涤自己，不是用淤泥，而是用清水，用哀矜，而非贪财。先摆脱掠夺，再施行哀矜。让我们\Quote{避恶行善。} \BibleRef{咏 36:27}管住你们的手，远离贪婪，然后才施与哀矜。但若我们用同一双手剥去一些人的衣服，即便我们没有把从他们那里夺来的衣物给另一些人穿上，我们也逃脱不了惩罚。因为那本是和好基础的东西，却成了所有邪恶的基础。这样行哀矜，不如不行哀矜；因为加音当初若不奉献祭品反而更好。如果奉献太少的人尚且惹天主发怒，何况拿别人的东西来奉献呢？岂不更激怒祂？\Quote{我命令你，}祂会说，\Quote{不可偷盗，而你们竟用偷来的东西尊敬我吗？你以为我喜悦这些事吗？}然后祂要对你说：\Quote{你邪恶地以为我像你一样；我要斥责你，将你的罪恶摆在你面前。} \BibleRef{咏 49:21}但愿我们中间没有人听到这话，而是行了纯洁的哀矜，并举着点燃的灯，得以进入婚宴厅，藉着我们的主耶稣基督的恩宠与慈爱，愿光荣归于祂及圣父与圣神，世世无穷。阿们。\par

\SourceRecord{Translated by Charles Marriott.；Nicene and Post-Nicene Fathers, First Series；Vol. 14.；Edited by Philip Schaff.；Buffalo, NY: Christian Literature Publishing Co.,；1889.；<http://www.newadvent.org/fathers/240173.htm>.}

% structure: p0001 source=h1 target=section reason=page-title

\section{若望福音讲道集 第七十四篇}

% structure: p0002 source=h2 target=subsection reason=relative-heading-rank; base=h2

\subsection{\BibleRef{若 14:8-9}}

先知对犹太人说：\BibleQuote{你有了娼妓的面容，你向众人毫无羞耻。}\BibleRef{耶 3:3} 如今看来，使用这话不仅针对那座城，也针对所有无耻地对抗真理的人，似乎是合适的。因为当\Person{斐理伯}对基督说：\Quote{请将圣父显示给我们}，祂回答说：\BibleQuote{我与你们在一起这么久了，你还不认识我吗？斐理伯！}然而有些人甚至在这句话之后仍将圣父与圣子分开。你还需要比这更近的亲近吗？的确，从这句话本身，有人陷入了\Person{撒伯里乌}的病症。但是，让我们把这两方面的人都撇下，视之为陷入完全相反的错谬，来思考这话的精确含义。祂说：\BibleQuote{我与你们在一起这么久了，你还不认识我吗？斐理伯！}那么情况如何呢？\Person{斐理伯}回答说：\Quote{你是那个我所询问的圣父吗？}祂说：\Quote{不。}因此祂没有说：\Quote{你不认识祂}，而是说：\Quote{你不认识我}，无非是宣告：圣子不是别的，正是圣父所是的，然而仍然是圣子。但\Person{斐理伯}怎么会问这个问题呢？基督曾说过：\BibleQuote{你们若认识我，也就认识我的圣父}\BibleRef{若 14:7}，祂也多次对犹太人说过同样的话。既然\Person{伯多禄}和犹太人常常问祂：\Quote{圣父是谁？}，既然\Person{多默}问过祂，而没有人学到任何清楚的事，祂的话仍未被理解；\Person{斐理伯}为了避免显得纠缠不休，在犹太人之后接着提问而麻烦祂，他说：\Quote{请将圣父显示给我们}，并加上：\Quote{这为我们就够了}，\Quote{我们不再寻求别的了。}然而基督曾说过：\BibleQuote{你们若认识我，也就认识我的圣父}\BibleRef{若 14:7}，并亲自宣告了圣父。但\Person{斐理伯}颠倒了顺序，说：\Quote{请将圣父显示给我们}，好像已经确切认识基督。但基督不容忍他，而是将他引入正途，说服他通过自己获得对圣父的认识，而\Person{斐理伯}却渴望用这肉眼看见祂，大概曾听说过先知们\Quote{看见过天主}。但那些情况，斐理伯，是迁就之举。因此基督说：\BibleQuote{从来没有人见过天主}\BibleRef{若 1:18}；又说：\BibleQuote{凡听见父而学习的，都到我这里来}\BibleRef{若 6:45}。\BibleQuote{你们从来没有听过祂的声音，也没有见过祂的形状}\BibleRef{若 5:37}。在旧约中：\BibleQuote{没有人看见我的面而能存活}\BibleRef{出 33:20}。基督说了什么？祂非常责备地说：\BibleQuote{我与你们在一起这么久了，你还不认识我吗？斐理伯！}祂没有说：\Quote{你没有看见}，而是说：\Quote{你不认识我}。\Person{斐理伯}可能会说：\Quote{为什么，我想了解关于祢的事？现在我寻求看见祢的圣父，而祢对我说，你不认识我？}这与问题有何关联呢？当然是十分密切的；因为如果祂是圣父所是的，但仍然为圣子，祂就有理由在自己内显示那生祂的那一位。然后为了区分位格，祂说：\BibleQuote{看见我的，就是看见了圣父}，免得有人断言圣父与圣子是同一的。因为如果祂是圣父，祂就不会说：\BibleQuote{看见我的，就看见了祂}。那么为什么祂没有回答说：\Quote{你所求的是不可能的事，是人不可允许的；只有我能这样}？因为\Person{斐理伯}曾说：\Quote{这为我们就够了}，好像认识基督，祂表明他甚至连祂都没有看见过。因为如果他能够认识圣子，他必定会认识圣父。因此祂说：\BibleQuote{看见我的，就是看见了圣父}。\Quote{若有人看见了我，他也要看见祂。}祂所说的是这样的：\Quote{看见我和祂都是不可能的。}因为\Person{斐理伯}寻求的是凭视觉的认识，既然他认为自己已经如此看见了基督，便同样渴望看见圣父；但耶稣表明他连自己都没有看见过。如果有人在这里把认识称为看见，我不反对，因为基督说：\Quote{认识我的，就认识了圣父。}然而祂没有这样说，而是为了确立\OriginalTerm{同体}{Consubstantiality}，祂宣告：\Quote{认识我的\OriginalTerm{本质}{Essence}的，也认识圣父的本质。}\Quote{这是什么？}有人说：\Quote{因为认识受造物的人，也认识天主。}然而所有人都认识受造物，并看见了它，但并非所有人都认识天主。此外，让我们思考\Person{斐理伯}寻求看见的是什么。是圣父的智慧吗？是祂的良善吗？不，而是天主无论是什么，就是那\OriginalTerm{本质}{Essence}。因此基督回答：\BibleQuote{看见我的}。现在，看见受造物的人并没有看见天主的本质。祂说：\Quote{若有人看见了我，就看见了圣父。}现在，如果祂是不同本质的，祂就不会这样说。但使用一个更粗糙的论证：不知道金子是什么的人，无法在银子里辨别金子的实体。因为一个本性不能由另一个本性显示出来。因此祂恰当地责备了他，说：\BibleQuote{我跟你们在一起这么久？}你们享受了这样的教导，看见了以权威行的奇迹，以及所有属于\OriginalTerm{天主性}{Godhead}的事，这些事只有圣父能行：罪得赦免，秘密被公开，死亡退却，从泥土中创造万物，而你们还不认识我吗？因为祂穿着血肉，所以祂说：\Quote{你不认识我吗？}\par

2. 你已经看见了圣父；不要再寻求看见更多；因为在祂内你看见了我。如果你看见了我，不要过于好奇；因为你也在我内认识了祂。\par

% structure: p0005 source=h2 target=subsection reason=relative-heading-rank; base=h2

\subsection{\BibleRef{若 14:10}}

也就是说，\Quote{我在那\OriginalTerm{本质}{Essence}中被看见。}\par

\BibleQuote{我对你们所说的话，我不是凭自己说的，}\par

你看见那无比的亲近，以及那同一\OriginalTerm{本质}{Essence}的证明吗？\par

\BibleQuote{住在我内的圣父，祂做这些工作。}\par

祂如何从言语进到作为？因为按理接下来祂应该说\Quote{父说话}。但祂在这里提到两件事，既关乎道理，也关乎奇迹。或可解释为言语本身也是作为。那么祂如何做呢？在另一处祂说：\Quote{我若不作我父的作为，你们不要信我。}（\BibleRef{若 10:37}）那么祂在这里说父作这些事是什么意思？正是要表明这一点：父与子之间没有间隔。祂所说的是：\Quote{父不会用一种方式行事，而我用另一种。}的确，在另一处，祂和父都作工：\Quote{我父一直工作到现在，我也工作}（\BibleRef{若 5:17}）；前一处表明作为的一贯性，后一处表明同一性。若话语的明显含义表示谦卑，不必惊讶；因为祂先说了\Quote{你们不信吗？}然后才如此说，表明祂这样调整祂的话语是为了引他走向信德；因为祂行走在他们心中。\par

% structure: p0011 source=h2 target=subsection reason=relative-heading-rank; base=h2

\subsection{若望福音 14:11}

\Quote{你们听到\InnerQuote{父}和\InnerQuote{子}，不应再寻求别的来确立那关乎本质的关系；但如果这不足以向你们证明那\OriginalTerm{相称性}{Condignity}与\OriginalTerm{同体}{Consubstantiality}，你们甚至可以从这些作为中得知。}若\Quote{看见了我，就看见了父}这一话是指着作为说的，祂后来就不会说：\par

\Quote{不然，当为那些作为本身而信我。}然后，为表明祂不仅能行这些事，还能行比这些更大的事，祂以超越的方式表述。因为祂不是说\Quote{我可行比这些更大的事}，而是更为奇妙的是，\Quote{我还能赐给别人也行比这些更大的事。}\par

% structure: p0014 source=h2 target=subsection reason=relative-heading-rank; base=h2

\subsection{若望福音 14:12}

这就是说：\Quote{如今你们去行奇迹，因为我要去了。}然后，当祂完成了祂论证的目的时，祂说：\par

% structure: p0016 source=h2 target=subsection reason=relative-heading-rank; base=h2

\subsection{若望福音 14:13}

你看见又是祂在行事吗？祂说：\Quote{我必成就}，而不是\Quote{我要求父}，而是\Quote{为叫父因我受光荣}。在另一处祂说：\Quote{天主要在祂身上光荣祂}（\BibleRef{若 13:32}），但这里说\Quote{祂要光荣父}；因为当子以大能显现时，生祂的父将受光荣。但\Quote{因我的名}是什么意思？即宗徒们所说的：\Quote{因耶稣基督的名，起来行走罢。}（\BibleRef{宗 3:6}）因为他们所行的所有奇迹，都是祂在他们内行的，并且\Quote{主的手与他们同在。}（\BibleRef{宗 11:21}）\par

% structure: p0018 source=h2 target=subsection reason=relative-heading-rank; base=h2

\subsection{若望福音 14:14}

你看见祂的权柄了吗？借着别人所行的事，祂自己亲自行；难道祂自己行的事，祂反而没有能力，除非被父推动？谁能这样说呢？但祂为何将这放在第二位？为要证实祂自己的话，并表明先前的话是迁就的说法。但\Quote{我到父那里去}意思是：\Quote{我不会消亡，而是存留在我的固有尊严中，且在天上。}祂说这一切，是为安慰他们。因为很可能他们尚不理解祂有关复活的言论，会想象一些悲伤的事，所以祂在别的言论中许诺要赐予他们这些事，以各种方式安抚他们，表明祂持续存在；不仅存在，还要显示更大的德能。\par

3. 那么，让我们跟随祂，背起十字架。因为即使没有迫害，另一种死亡的季节仍与我们同在。经上说：\Quote{你们要致死属于地上的肢体。}\BibleRef{哥 3:5} 那么，让我们熄灭\OriginalTerm{私欲偏情}{concupiscence}，杀死愤怒，消除嫉妒。这就是\Quote{生活的祭献。}\BibleRef{罗 12:1} 这祭献不以灰烬告终，不消散在烟雾中，不需木柴、火或刀，因为它既有火又有刀，即圣神。用这把刀，割去你心中多余的、异己的部分；开启你封闭的耳朵，因为恶习和邪恶的欲望常会阻挡圣言的入口。贪财的欲望，当它摆在人面前时，就不允许人听关于施舍的圣言；而怨恨临到时，就筑起一道墙，阻挡关于爱德的教导；其他疾病相继袭来，使灵魂对一切更加迟钝。那么，让我们除去这些邪恶的欲望；只要愿意，它们就全部熄灭。因为，我恳求，不要因财富之爱的暴政而眼看它，而要明白暴政来自于我们自身的懈怠。许多人甚至说他们不知道金钱是什么。这欲望不是自然的；自然的欲望从一开始就植入我们内，但金银在很长一段时间里甚至不知道它们是什么。那么这种欲望从何而来？来自虚荣和极度的懈怠。因为欲望中，有些是必要的，有些是自然的，有些两者都不是。例如，那些若不满足就会毁灭受造物的欲望，既是自然的又是必要的，如饮食和睡眠的欲望；肉身的欲望是自然的，但不必要，因为许多人已克服它而未死。而财富的欲望既非自然也非必要，而是多余的；如果我们选择，甚至可以不让它开始。基督论及童贞时曾说：\Quote{谁能够领会，就让他领会吧。}\BibleRef{玛 19:12} 但论及财富却不是这样，而是如何？\Quote{若不撇下一切所有的，不能做我的门徒。}\BibleRef{路 14:33} 容易的祂予以推荐，而超越多数人的则留给自由选择。那么，为何我们剥夺自己的一切借口？被更暴虐的情欲所俘虏的人，不会受很重的惩罚；但被软弱的情欲所征服的人，则失去一切辩护。因为当祂说：\Quote{你看见我饿了，却不给我吃}\BibleRef{玛 25:42} 时，我们如何回答？我们还有什么托辞？我们当然会诉穷；但我们并不比那投了两个小钱而超越众人的寡妇更穷。因为天主所要求的不是奉献的多寡，而是心灵的分寸；祂这样做是出于祂的慈爱。那么，让我们景仰祂的慈爱，尽己所能奉献，好使我们今生和来世都丰盛地获得天主的慈爱，并能享有应许给我们的福祉，借着我们的主耶稣基督的恩宠和慈爱。愿光荣归于祂，直到永远。亚孟。\par

\SourceRecord{Translated by Charles Marriott.；Nicene and Post-Nicene Fathers, First Series；Vol. 14.；Edited by Philip Schaff.；Buffalo, NY: Christian Literature Publishing Co.,；1889.；<http://www.newadvent.org/fathers/240174.htm>.}

% structure: p0001 source=h1 target=section reason=page-title

\section{若望福音讲道集第七十五篇}

% structure: p0002 source=h2 target=subsection reason=relative-heading-rank; base=h2

\subsection{若望福音 14:15-17}

1. 我们到处都需要工作和行动，而不仅仅是言辞的虚饰。因为说起来和许诺对任何人而言都是容易的，但行动起来却不那么容易。我为何这样说？因为现在有许多人说他们敬畏并爱慕天主，但在行为上却表现出相反；然而，天主所要求的是藉行为显示的爱。因此祂对门徒们说：\Quote{如果你们爱我，就要遵守我的诫命。}因为祂曾告诉他们：\Quote{你们无论求什么，我必成就，}以免他们以为仅仅\Quote{祈求}就有效，祂又加上：\Quote{如果你们爱我，}\Quote{那么，}祂说，\Quote{我必成就。}既然他们听到\Quote{我往父那里去，}时很可能会困扰，祂告诉他们\Quote{现在困扰就不是爱，爱就是服从我的话。我曾给你们一条诫命，要你们彼此相爱，如同我爱了你们一样；这就是爱：服从我的这些话，并顺从你们所爱的那一位。}\par

\Quote{我要请求父，祂必赐给你们另一位护慰者。}祂的言论再次是俯就的。因为他们还不认识祂，很可能热切寻求与祂同在、听祂讲论、与祂在肉身中的同在，而当祂不在时，就会不接受任何安慰，祂说什么呢？\Quote{我要请求父，祂必赐给你们另一位护慰者，}即，\Quote{另一位像我一样的。}让那些患有\Person{撒伯里乌}病症、对圣神持有不当意见的人羞愧吧。因为这番言论的奇妙之处在于，它用同一击打倒了相互矛盾的异端。因为藉说\Quote{另一位，}祂显示了位格的差别，藉说\Quote{\OriginalTerm{护慰者}{Paraclete},}显示了本体的联系。但祂为什么说\Quote{我要请求父}呢？因为如果祂说\Quote{我要派遣祂，}他们就不会如此相信，而现在目的是要祂被相信。因为后来祂宣布祂亲自派遣祂，说：\BibleQuote{你们领受圣神吧}\BibleRef{若 20:22}；但在此处祂告诉他们说祂请求父，为要使祂的言论对他们显得可信。既然若望论祂说\BibleQuote{从祂的圆满中我们众人都领受了}\BibleRef{若 1:16}；但祂所拥有的，祂怎能从另一位领受呢？并且又说：\BibleQuote{祂要以圣神和火给你们施洗}\BibleRef{路 3:16}。\Quote{但祂比宗徒们多有什么呢？如果祂要请求祂的父将祂赐给他人，而宗徒们常常甚至不经祈祷就似乎这样做了？}并且，如果祂是根据来自父的请求而被派遣，祂如何自己降临呢？那充满万有的怎能被另一位派遣？祂\BibleQuote{随祂的意愿分给各人}\BibleRef{格前 12:11}，并且有权柄地说：\BibleQuote{你们为我分开保禄和巴尔纳伯}\BibleRef{宗 13:2}。那些执事是在服事天主，但祂仍有权柄地把他们召至祂自己的工程；并非召他们到不同的工程，而是为显示祂的能力。\Quote{那么，}有人说，\Quote{\InnerQuote{我要请求父？}是什么意思？}（祂这样说）为要显示祂来临的时刻。因为当祂藉祭献洁净了他们后，圣神才降在他们身上。\Quote{为什么祂与他们同在时，圣神没有来呢？}因为祭献尚未献上。但后来罪被解除，他们被派遣去面对危险，并为战斗而脱去（世俗），那时就需要那位傅油者来临。\Quote{但为什么圣神没有在复活后立即降临呢？}为使他们极其渴望祂，并以极大的喜乐领受祂。因为只要基督与他们同在，他们就不在患难中；但当祂离去后，他们变得没有防卫，陷入极大的恐惧，就会以极大的乐意领受祂。\par

\Quote{祂与你们同在。}这显示即使在死亡后，祂也不离去。但唯恐他们听到\Quote{\OriginalTerm{护慰者}{Paraclete},}时，会想象第二次降生成人，并期望用眼睛看见祂，祂纠正他们说：\Quote{世界不能接受祂，因为看不见祂。}\Quote{祂不会像我那样与你们同在，但将住在你们的灵魂内}；因为这正是\Quote{要在你们内。}祂称它为\Quote{真理之神}；从而解释了旧约中的预象。\Quote{为的是祂能与你们同在。}什么是\Quote{与你们同在}？就是祂自己所说的：\BibleQuote{我与你们同在}\BibleRef{玛 28:20}。此外，祂也暗示另一件事：关于圣神的情况不会像我的那样，祂永远不会离开你们。\Quote{世界不能接受祂，因为看不见祂。}\Quote{为什么，其他位格有什么是可见的呢？}没有；但祂在这里说的是认识；至少祂加上：\Quote{也不认识祂。}因为祂惯常在精确认识的情况下称之为\Quote{看见}；因为视觉比其他感官更清晰，祂常以此表示精确的认识。在此\Quote{世界}祂指的是\Quote{恶人}，这样也安慰了门徒，给了他们一份特别的恩赐。看，祂在多少方面提升关于祂的言论。祂说：\Quote{祂是另一位像我一样的}；祂说：\Quote{祂不会离开你们}；祂说：\Quote{祂独自来到你们这里，正如我也独自来到}；祂说：\Quote{祂留在你们内}；但即使这样，祂也没有驱散他们的沮丧。因为他们仍在寻求祂和与祂同在。为治疗这种情绪，祂说：\par

% structure: p0006 source=h2 target=subsection reason=relative-heading-rank; base=h2

\subsection{若望福音 14:18}

2. \Quote{不要害怕，}祂说，\Quote{我不是说我要给你们派遣另一位\OriginalTerm{护慰者}{Comforter}，好像我自己要永远离开你们；我是说祂与你们同在，好像我再也见不到你们一样。因为我自己也要到你们这里来，我不会留下你们做孤儿。}因为开始时祂说\Quote{孩子们，}所以这里祂也说\Quote{我不会留下你们做孤儿。}起初祂告诉他们：\Quote{你们要去我往的地方}；以及\Quote{在我父的家里有许多住处}；但在这里，由于时间很长，祂赐给他们圣神；当他们不明白祂所说的是什么而未能得着足够安慰时，祂说\Quote{我不会留下你们做孤儿，}；因为他们主要需要这个。但既然\Quote{我要到你们这里来}这句话是宣告一种\Quote{\OriginalTerm{临在}{presence}}，请看，为了不使他们再次寻求像以前那样的临在，祂没有明确告诉他们这件事，而是暗示了它；因为祂曾说过：\par

% structure: p0008 source=h2 target=subsection reason=relative-heading-rank; base=h2

\subsection{若望福音 14:19}

好像在说：\Quote{我确实到你们这里来，但不是像从前一样，每天都与你们同在。}并且免得他们说：\Quote{那么你怎么对犹太人说过\InnerQuote{从今以后你们不得再见我}呢？}祂解决了这个矛盾，说：\Quote{只对你们而言}；因为圣神的性质也是如此。\par

\Quote{因为我活着，你们也要活着。}\par

因为十字架并不能最终使我们分离，只是暂时隐藏片刻；而\Quote{生命}一词在我看来不单指今世，也指来世。\par

% structure: p0012 source=h2 target=subsection reason=relative-heading-rank; base=h2

\subsection{若望福音 14:20}

就父而言，这些话是指\OriginalTerm{本质}{Essence}；就门徒而言，是指心意合一和来自天主的帮助。\Quote{请告诉我，这合理吗？}有人说。请问，相反的就合理吗？因为基督与门徒之间的间隔是巨大且无穷的。如果同样的词语被使用，不必惊奇；因为圣经在用于天主和人时，常在同一词语上使用不同含义。所以我们被称为\Quote{神}和\Quote{天主的儿子}，然而这个词应用于我们和应用于天主时，力量不同。圣子被称为\Quote{\OriginalTerm{肖像}{Image}}和\Quote{\OriginalTerm{光荣}{Glory}}；我们也一样，但我们之间的间隔是巨大的。此外，\BibleQuote{你们是属于基督的，基督是属于天主的}\BibleRef{格前 3:23}，但不像基督属于天主那样，我们属于基督。但祂说的是什么？\Quote{当我复活后，}祂说，你们将知道我不与父分离，而是与祂有同等的权能，并且我不断与你们同在，当事实宣告你们从我而来的帮助，当你们的仇敌被压制，你们放胆说话，当危险从你们的路途被移除，当福音的宣讲日渐兴旺，当万物都屈服并让位于真宗教之言时。\BibleQuote{父怎样派遣了我，我也怎样派遣你们。}\BibleRef{若 20:21}你看，这里的词也不是同等的力量吗？因为如果我们把它当作同等的，宗徒们与基督就没有区别了。但为什么祂说\Quote{那时你们就知道}？因为那时他们看见祂复活了并与他们交谈，那时他们学会了正确的信仰；因为圣神的能力是伟大的，教导了他们一切。\par

% structure: p0014 source=h2 target=subsection reason=relative-heading-rank; base=h2

\subsection{若望福音 14:21}

仅仅拥有它们是不够的，还需要精确地遵守。但为什么祂常常对他们说同样的话？如：\BibleQuote{如果你们爱我，就要遵守我的命令}\BibleRef{若 14:15}；以及\Quote{那有了我的命令而遵守的}；以及\BibleQuote{如果有人听到我的言语而遵守，他就是爱我的——那不遵守我言语的，就不爱我。}\BibleRef{若 14:24}我想祂是暗示他们的沮丧；因为既然祂曾对他们讲过许多关于死亡的智慧言语，说：\BibleQuote{那在现世憎恨自己性命的，必要保存性命于永生}\BibleRef{若 12:25}；以及\BibleQuote{谁不背起自己的十字架跟随我，就不配是我的}\BibleRef{玛 10:38}；并且还要说其他事情，责备他们，祂说：\Quote{你们以为你们的忧愁是出于爱吗？不忧愁才是爱的记号。}因为祂一直想确立这一点，所以当祂继续时，祂将祂的言论总结于这一点：\BibleQuote{如果你们爱我，}祂说，\BibleQuote{你们就应当欢喜，因为——我往父那里去}\BibleRef{若 14:28}，但现在你们处于这种状态是由于懦弱。这样对待死亡的人不是那些记住我命令的人；因为你们应该被钉十字架，如果你们真的爱我，因为我的言语劝诫你们不要害怕那些杀害身体的人。那些如此的人，父和我都爱他们。并且我要向他显现我自己。于是犹达斯说：\par

% structure: p0016 source=h2 target=subsection reason=relative-heading-rank; base=h2

\subsection{若望福音 14:22}

你看见他们的灵魂被恐惧紧紧压抑吗？因为他心慌意乱，以为就像我们在梦中看见死人一样，祂也会被看见。因此，为了不使他们这样想，请听祂说什么。\par

% structure: p0018 source=h2 target=subsection reason=relative-heading-rank; base=h2

\subsection{若望福音 14:23}

祂几乎是在说：\Quote{正如圣父启示自己，我也如此启示自己。}不仅如此，祂还藉着说\Quote{我们要到他那里去，并与他同住}来消除疑虑，这事不属于梦境。但请看，我祈求你们，门徒们困惑了，不敢直说他们想说的话。因为他们没有说\Quote{我们有祸了，祢死了，并要像死人那样来到我们这里}；他们没有这样说；而是说\Quote{祢为什么只向我们显现，而不向世界显现呢？}于是耶稣说：\Quote{我接纳你们，因为你们遵守我的诫命。}为了使他们在以后看见祂时，不至于认为祂是幻象，祂事先说了这些话。为了不让他们认为祂会像我所说的那样显现，祂还告诉他们理由：\Quote{因为你们遵守我的诫命}；祂说圣神也会以同样的方式显现。如今，在陪伴祂这么长时间之后，他们尚且不能承受那\OriginalTerm{本质}{Essence}，或者更确切地说，甚至不能想象那本质，如果祂一开始就如此向他们显现，他们会怎样呢？为此，祂也同他们一起进食，以使那行为不像是幻象。因为如果他们在看见祂在水面上行走时就已这样想，尽管他们看见了祂惯常的形象，而且祂离得不远，那么如果他们突然看见祂复活了——他们曾看见祂被带走并被包裹起来——他们会想象什么呢？因此，祂不断地告诉他们祂要显现，以及为什么显现和如何显现，免得他们以为祂是幻象。\par

% structure: p0020 source=h2 target=subsection reason=relative-heading-rank; base=h2

\subsection{若 14:24}

\Quote{所以，那不听这些话的，不仅不爱我，也不爱圣父。}因为如果遵守诫命是爱的确凿证明，而这些诫命属于圣父，那么那听诫命的人就不仅爱圣子，也爱圣父。\Quote{那么\InnerQuote{你们的}和\InnerQuote{不是你们的}这话是什么意思呢？}意思是\Quote{我不是没有圣父就说话，也不是凭自己说任何不合祂心意的话。}\par

% structure: p0022 source=h2 target=subsection reason=relative-heading-rank; base=h2

\subsection{若 14:25}

由于这些话并不清楚，而且有些他们不明白，对大多数话感到疑惑，为了不使他们再次困惑并问\Quote{什么诫命？}，祂使他们从一切困惑中解脱出来，说：\par

% structure: p0024 source=h2 target=subsection reason=relative-heading-rank; base=h2

\subsection{若 14:26}

\Quote{也许现在这些事你们不明白，但祂是这些事的明师。}而\Quote{与你们同在}\BibleRef{若 14:17}这句话，暗示祂自己要离开。然后，为了不使他们悲伤，祂说，只要祂还与他们同在，而圣神还没有来，他们就无法理解任何伟大或崇高的事。祂这样说，是为预备他们高尚地承受祂的离别，因为这离别将为他们带来极大的福分。祂不断称祂为\OriginalTerm{护慰者}{Comforter}，是因为他们当时所承受的苦难。即使听到这些事，他们仍然感到困扰，想到忧伤、战争、祂的离别，请看祂如何再次安慰他们，说：\par

% structure: p0026 source=h2 target=subsection reason=relative-heading-rank; base=h2

\subsection{若 14:27}

祂几乎是在说：\Quote{世间的困扰对你们有何伤害，只要你们与我平安相处？因为这平安不同于那平安。那平安是外在的，常常有害无益，对拥有它的人毫无益处；但我赐给你们的平安是使你们彼此和睦，这使你们更坚强。}因为祂又说\Quote{我离开}，这是离开者的措辞，足以使他们困惑，因此祂再说：\par

\Quote{不要让你们的心烦乱，也不要害怕。}\par

你们看见吗，他们既出于爱，又出于恐惧？\par

% structure: p0030 source=h2 target=subsection reason=relative-heading-rank; base=h2

\subsection{若 14:28}

4. 这能给他们带来什么喜乐？什么安慰？那么这些话是什么意思？他们还不明白关于复活的事，对祂也没有正确的看法（因为他们尚且不知道祂会复活，又怎能明白呢？），他们认为圣父是全能的。于是祂说：\Quote{如果你们为我担心，认为我不能自卫，如果你们不相信我在受难后会再见你们，那么当你们听说我到圣父那里去时，你们应当欢喜，因为我到一位更大的、能解除一切危险的那里去。}\Quote{你们听见我怎样对你们说了。}祂为什么这样说？因为祂说：\Quote{我对于将要发生的事如此确信，以致我甚至预先告诉你们，我远非害怕。}这也是后面话的意思。\par

% structure: p0032 source=h2 target=subsection reason=relative-heading-rank; base=h2

\subsection{若 14:29}

好像祂说过：\Quote{如果我没有告诉你们，你们就不会知道。而我若没有信心，就不会告诉你们。}你们看见这都是迁就的话语吗？因为当祂说\Quote{你们以为我不能请求圣父，祂就立刻赐给我十二营以上的天使吗？}\BibleRef{玛 26:53}时，祂是针对听者隐秘的念头而说的；因为没有人，即使疯狂至极，会说祂不能自救而需要天使；但因他们把祂看作人，因此祂提到\Quote{十二营以上的天使}。但实际上祂只是问来捉拿祂的人一个问题，就把他们打倒在地了。\BibleRef{若 18:6}（若有人说圣父更大，因为祂是圣子的根源，我们不会反驳。但这绝不意味着圣子有不同的\OriginalTerm{本质}{Essence}。）但祂所说的是这样的：\Quote{只要我在这里，你们自然以为我处于危险中；但当我到了那里，你们可以放心，我安全了；因为没有人能胜过祂。}这些话都是针对门徒的软弱而说的，因为\Quote{我自己有信心，不在乎死亡。}因此，祂说：\Quote{在事情发生之前，我已经告诉了你们这些事}；\Quote{但因为，}祂说，\Quote{你们还不能接受关于这些事的话，我甚至从你们称为伟大的圣父那里带给你们安慰。}这样安慰了他们之后，祂又告诉他们忧伤的事，\par

% structure: p0034 source=h2 target=subsection reason=relative-heading-rank; base=h2

\subsection{若 14:30}

祂说\Quote{今世的领袖}，是指魔鬼，又将恶人也称为同一名字。因为他不统治天地，否则他早已颠覆并摧毁一切，但他统治那些顺从自己的人。因此祂称他为\Quote{这世界黑暗的领袖}，在此再次称恶行为\Quote{黑暗}。那么，魔鬼要杀害祢吗？断然不是；\Quote{他在我内一无所能}。那么他们如何杀害祢？因为我愿意，并且\par

% structure: p0036 source=h2 target=subsection reason=relative-heading-rank; base=h2

\subsection{\BibleRef{若 14:31}}

祂说：\Quote{因为并非隶属于死亡，也不是死亡的债户，我因对圣父的爱而忍受它。}祂这样说，为再次激励他们的心灵，使他们知道祂不是不情愿，而是甘愿走向此事，并且祂这样做是轻视魔鬼。祂说了\Quote{短短的时候，我同你们在一起}\BibleRef{若 7:33} 还不够，而是不断提起这个痛苦的主题（有充分理由），直到通过穿插愉快的事使它为他们所接受。因此，祂时而说：\Quote{我去，并且我再来}；以及\Quote{我在那里，你们也在那里}；以及\Quote{现在你们不能跟我去，但后来你们要跟我去}；以及\Quote{我到圣父那里去}；以及\Quote{圣父比我大}；以及\Quote{在事情发生以前，我已告诉你们}；以及\Quote{我不是因被迫而受这些，而是因对圣父的爱}。这样他们可以认为，这行动既非毁灭性，也非有害的，如果至少那位深爱祂且为祂所深爱者如此愿意。为此，在穿插这些愉快话语的同时，祂也经常说出痛苦的话，以操练他们的心智。因为\Quote{同你们在一起}\BibleRef{若 16:7} 以及\Quote{我的离去对你们有益}都是安慰的话。为此，祂事先说了许多关于圣神的话：\Quote{在你们内}，\Quote{世界不能领受}，\Quote{祂要使你们想起一切}，\Quote{真理之神}，\Quote{圣神}，\Quote{护慰者}，以及\Quote{这对你们有益}，为使他们在好像没有人站在他们面前帮助他们时不沮丧。\Quote{对你们有益}，祂说，表明圣神将使他们成为属神的人。\par

5. 至少，我们看到，这正是所发生的事。因为那些先前战栗恐惧的人，在领受了圣神之后，便跃入危险之中，赤身与刀剑、烈火、野兽、海洋以及各种刑罚斗争；这些没有学问、愚昧无知的人，竟如此大胆地宣讲，令听者惊奇。因为圣神将他们从泥人变为铁人，给他们插上翅膀，不容任何人间事物使他们跌倒。因为那恩宠正是如此：若遇到沮丧，它便驱散它；若遇到邪欲，它便焚毁它；若遇到怯懦，它便驱逐它，并且不允许那领受它的人此后仍是凡人，而是仿佛将他提升到天上一样，使他心中描绘出那里的一切。\BibleRef{宗 4:32} 因此没有一个人说他的财物中有什么是自己的，但他们恒心祈祷、赞美，并一心一意。因为圣神最要求的是这些，因为\Quote{圣神的效果是喜乐、平安——信德、温和。}\BibleRef{迦 5:22} 有人说：\Quote{然而属神的人也常常忧愁。}但那种忧愁比喜乐更甘甜。加音忧愁，却是世俗的忧愁；保禄忧愁，却是合乎天主的忧愁。凡属神的事都带来最大的益处，正如凡属世俗的事都带来极大的亏损。因此，让我们藉着遵守诫命，吸引圣神那无敌的助佑，这样我们就丝毫不亚于天使。因为天使之所以有这种品格，并非因为他们无形体；若是如此，就没有一个无形体者会变坏了；但在一切情况下，意志都是原因。因此，在无形体者中，有些被发现比有形体的人或无理之物更坏；而在有形体者中，有些比无形体者更好。例如，所有义人，无论行了什么正义的事，都是在居住在地上、拥有肉体时行的。因为他们居住在地上，如同客旅和寄居的；但在天上，如同公民。因此，你也不要说：\Quote{我披着肉体，不能得胜，也不能承担为了德性而付出的辛劳。}不要责怪造物主。因为如果披着肉体使德性不可能，那么过错就不在我们了。但肉体并不使德性不可能，这已由圣人的行列证明了。肉体的本性并没有阻止保禄成为他所是的，也没有阻止伯多禄领受天堂的钥匙；哈诺客也曾披着肉体，被接去而不复寻见。同样，厄里亚也是带着肉体被提升天。亚巴郎、依撒格和雅各伯也都披着肉体而光辉闪耀；若瑟在肉体中与那个放荡的女人搏斗。但我为何谈论肉体呢？因为即使你在肉体上加上锁链，也无妨碍。保禄说：\Quote{我虽被捆绑，}但\Quote{天主的话却未被捆绑。}\BibleRef{弟后 2:9} 我为何谈论捆锁和锁链呢？再加上监狱和铁栏，但这些对德性也毫无阻碍；至少保禄这样教导了我们。因为灵魂的束缚不是铁链，而是怯懦、贪财和种种情欲。这些捆绑我们，即使我们的身体是自由的。有人说：\Quote{但是，}有人说，\Quote{这些都是从身体而来的。}这是借口和虚伪的托词。因为如果它们是从身体产生的，那么人人都应承受它们。正如我们不能避免疲劳、睡眠、饥饿和口渴，因为它们属于我们的本性；同样，如果这些也属于同一类，它们就不会允许任何人免于它们的暴虐；但由于许多人避开了它们，显然这些事情是粗心灵魂的过失。因此，我们要停止这种说法，不要责怪身体，而要叫身体服从灵魂，好使我们因我们的主耶稣基督的恩宠和慈爱，在管辖之下享受永福。愿光荣归于祂，直到永远。阿们。\par

\SourceRecord{Translated by Charles Marriott.；Nicene and Post-Nicene Fathers, First Series；Vol. 14.；Edited by Philip Schaff.；Buffalo, NY: Christian Literature Publishing Co.,；1889.；<http://www.newadvent.org/fathers/240175.htm>.}

% structure: p0001 source=h1 target=section reason=page-title

\section{若望福音第七十六篇讲道}

% structure: p0002 source=h2 target=subsection reason=relative-heading-rank; base=h2

\subsection{\BibleRef{若 14:31-15:1}}

1. \Quote{无知}使灵魂胆怯而不刚毅，正如天上教理的训导使之伟大崇高。因为当灵魂未曾享受照顾时，它在某种意义上变得胆怯——不是出于本性，而是出于意志。因为当我看见一个从前勇敢的人如今变成懦夫时，我说这种感觉不再属于本性，因为本性的东西是不变的。再者，当我看见那些方才还是懦夫的人突然变得大胆时，我也作同样的判断，并把一切都归因于意志。既然连门徒们在学习他们应当知道的以及堪当圣神恩赐之前，也是极其恐惧的；但之后他们变得比狮子还要勇敢。所以伯多禄，他不能忍受一个使女的威胁，后来却被倒悬、受鞭打，虽经历万千危险，却缄默不语，而是忍受他所忍受的如同做梦，在这样的境况中放胆说话；但在受难之前却不是这样。因此基督说：\BibleQuote{起来，我们从这里去吧。}\Quote{但为什么，请告诉我？难道他不知道犹达斯要来抓他的时辰吗？或者他可能害怕他会来并抓住他们，害怕阴谋者在他完成最杰出的教训之前就扑向他。}断然不可！这些事与他的尊严相去甚远。\Quote{如果他不害怕，为什么他要移开他们，然后在完成他的讲论后，把他们带到一个犹达斯知道的园子里？即使犹达斯来了，他难道不能使他们眼瞎，正如他在叛徒不在时也曾做过的那样？他为什么要移开他们？}他让门徒们稍得喘息。因为他们很可能，由于在一个显眼的地方，会因时间和地点而战栗害怕（因为那时是深夜），并且不会留心他的话语，而是不断转身，想象他们听到了那些要攻击他们的人的声音；尤其是当他的师傅的言论使他们预料到坏事时。因为他说：\BibleQuote{还有一会儿，我就不与你们同在了}，以及\BibleQuote{这世界的首领来了}。既然现在当他们听到这些和类似的话时，他们感到困扰，好像他们一定会立刻被捉拿，他就带他们到另一个地方，好使他们自认为安全，并无忧惧怕地听他讲论。因为他们即将听到高深的教理。因此他说：\BibleQuote{起来，我们从这里去吧。}然后他加上并说：\BibleQuote{我是葡萄树，你们是枝条。}他借这个比喻要暗示什么？就是那不注意他话语的人不能有生命，并且即将发生的奇迹将由基督的能力完成。\BibleQuote{我父是园丁。}\Quote{那么怎样？圣子需要一个内在工作的能力吗？}断然不可！这个例子并不表示这个。注意他多么精确地贯彻这个比喻。他并没有说\Quote{根}享受园丁的照顾，而是说\Quote{枝条}。而且在这里引入脚目的用意无他，就是要他们知道，没有他的能力，他们什么也不能做，并且他们应当借着信德与他联合，如同枝条与葡萄树联合一样。\par

% structure: p0004 source=h2 target=subsection reason=relative-heading-rank; base=h2

\subsection{\BibleRef{若 15:2}}

这里他暗示生活方式，表明没有善工就不可能在他内。\par

\BibleQuote{凡结果实的枝条，他就修剪它。}\par

也就是说，\Quote{使它得到极大的照顾。}然而，根比枝条更需要照顾，比如松土、清理，但他在这里对此却只字未提，只谈枝条。表明他自给自足，而门徒们需要园丁的大量帮助，虽然他们是非常优秀的。因此他说：\BibleQuote{那结果实的，他就修剪它。}一根枝条，因为不结果实，甚至不能留在葡萄树上；但另一根，因为它结果实，他使它更加多产。有人可能断言，这话也与那时临到他们身上的迫害有关。因为\Quote{修剪}就是\Quote{修剪}，使枝条结得更好。由此可见，迫害反而使人变得更坚强。然后，免得他们问他说这些事是指谁说的，并免得他使他们再次陷入焦虑，他说：\par

% structure: p0008 source=h2 target=subsection reason=relative-heading-rank; base=h2

\subsection{\BibleRef{若 15:3}}

你们看见他如何介绍自己照顾枝条吗？他说：\BibleQuote{我已经洁净了你们}，然而上面他宣称父做这事。但父与子之间没有分离。\Quote{现在你们的部分也必须履行。}然后，为表明他这样做不是需要他们的服务，而是为了他们的长进，他加上：\par

% structure: p0010 source=h2 target=subsection reason=relative-heading-rank; base=h2

\subsection{\BibleRef{若 15:4}}

为使他们不被胆怯从他分离，他将他们因恐惧而松弛的灵魂紧固并粘合在自己身上，并向他们提出对未来的美好希望。因为根存留，但被移走或被留下，属于枝条。然后，在通过愉快和痛苦的事两方面激励他们之后，他首先要求我们方面应做的事。\par

% structure: p0012 source=h2 target=subsection reason=relative-heading-rank; base=h2

\subsection{\BibleRef{若 15:5}}

你们看见圣子对于门徒的照顾并不少于圣父吗？圣父修剪，但他把他们保留在自己内。留在根上使枝条结果实。因为那未修剪的，如果留在根上，会结果实，虽然或许不如所应有的那么多；但那不留在根上的，就完全不结果实。然而，\Quote{修剪}也被表明属于圣子，而\Quote{留在根上}则属于圣父，他也生了那根。你们看见一切如何是共有的，无论是\Quote{修剪}，还是享受从根而来的德能？\par

2. 现在，什么也不能做是一个重大的惩罚，但他在此并未止于惩罚，而是继续讲论。\par

% structure: p0015 source=h2 target=subsection reason=relative-heading-rank; base=h2

\subsection{\BibleRef{若 15:6}}

不再享受园丁之手的恩惠。\BibleQuote{就枯干了。}也就是说，如果他有什么从根来的，他就失去它；如果有任何恩宠，他就被剥夺了这恩宠，并失去了从根而来的帮助和生命。结局是什么呢？\BibleQuote{就扔在火里。}凡与他同住的，必不如此。然后他显示什么是\Quote{住在}，并说：\par

% structure: p0017 source=h2 target=subsection reason=relative-heading-rank; base=h2

\subsection{\BibleRef{若 15:7}}

你们看到我先前说得有理吗？祂要求凭行为来证明。因为当祂说了\Quote{你们无论求什么，我必成就}\EditorialNote{若14:14-15}，祂又加上了\Quote{如果你们爱我，就要遵守我的诫命。}这里也说：\Quote{如果你们住在我内，我的话也住在你们内。}\par

\Quote{你们无论求什么，必给你们成就。}\par

祂说这话是要表明，那些图谋反对祂的人将被焚毁，而\Quote{他们}要结果实。然后将恐惧从他们身上转移到他人身上，并显示他们将是不可战胜的，祂说：\par

% structure: p0021 source=h2 target=subsection reason=relative-heading-rank; base=h2

\subsection{\BibleRef{若 15:8}}

因此祂使祂的言论可信，因为结果实关乎圣父的光荣，祂不会忽略自己的光荣。\Quote{你们将成为我的门徒。}你们看到结果实的人如何就是门徒吗？但\Quote{圣父因此得光荣}是什么意思？\Quote{祂因你们住在我内、你们结果实而喜乐。}\par

% structure: p0023 source=h2 target=subsection reason=relative-heading-rank; base=h2

\subsection{\BibleRef{若 15:9}}

现在祂终于以更人性的方式说话，因为这话既是向人说的，便有其独特的力量。祂显出了何等大的爱！祂选择了死亡，使那些曾是祂的奴仆、仇恨祂、公开敌对祂的人配受如此光荣，并带领他们上升到天上！\Quote{因此，如果我爱你们，你们要勇敢；如果你们结果实是圣父的光荣，就不要想任何坏事。}然后，为了不使他们懈怠，请看祂如何再次激励他们：\par

\Quote{你们要常在我的爱内。}\par

\Quote{因为你们有能力这样做。}这如何成就呢？\par

% structure: p0027 source=h2 target=subsection reason=relative-heading-rank; base=h2

\subsection{\BibleRef{若 15:10}}

祂又以人的方式继续论说，因为立法者肯定不会受诫命约束。你们看到这里也正如我一直所说的，这是因为听者的软弱而宣告的吗？因为祂主要是对他们的疑虑说话，用各种方式向他们表明他们是安全的，他们的敌人正在丧亡，并且他们所有的一切都来自圣子，而且如果他们显示出纯洁的生命，就没有人能主宰他们。请注意祂以非常权威的方式与他们讲话，因为祂没有说\Quote{常在我父的爱内}，而是说\Quote{在我的爱内}。然后，免得他们说\Quote{祢使我们与众人为敌后，就离开我们、走了}，祂表明祂并没有离开他们，而是如果他们愿意，祂就与他们结合，好像葡萄枝在葡萄树上。然后，免得他们因自信而懈怠，祂没有说如果心思懈怠，祝福就不会失去。而且为了不将行动归于自己而使他们更容易跌倒，祂说：\Quote{圣父因此得光荣。}因为祂处处显示自己和圣父对他们的爱。因此，光荣的不是犹太人的事物，而是他们将要领受的事物。并且为了不让他们说\Quote{我们被赶出了祖先的产业，被抛弃，赤身露体、一无所有}，祂说：\Quote{看，}祂说，\Quote{看我。圣父爱我，但我仍然承受这些指定的事。因此我现在离开你们，不是因为我不爱你们。因为如果我被杀害，却不将此视为不被圣父所爱的证明，你们也不该烦恼。因为如果你们常在我的爱内，这些危险决不能因爱而对你们造成任何伤害。}\par

3. 既然爱是强而有力、不可抗拒的，不是空言，就让我们以行动彰显它。祂在我们仍是仇敌时就与我和好了，既已成为祂的朋友，就让我们继续如此。祂先走了这条路，至少让我们跟随；祂爱我们不是为了自己的益处，（因为祂一无所缺，）至少让我们为我们的益处而爱祂；祂爱了我们这些仇敌，至少让我们爱祂这位朋友。目前我们却反其道而行；因为每一天天主都因我们而受亵渎，因我们的掠夺，因我们的贪得无厌。也许你们中有人会说：\Quote{你每天讲论的都是贪财。} 但愿我也能每夜讲论它；但愿我能这样做，在市场上、在餐桌前跟着你们；但愿妻子、朋友、儿女、仆人、农夫、邻舍，甚至连铺路石和墙壁，都能不断高声说出这句话，好让我们或许稍得缓解。因为这病症已占据了全世界，占据了所有人的灵魂，\OriginalTerm{玛门}{Mammon}的暴政何等厉害。我们已被基督赎出，却成了黄金的奴隶。我们宣告一者的主权，却服从另一者。无论\Quote{他}命令什么，我们都欣然服从，并为他拒绝了家庭、友谊、本性、法律，以及一切。没有人仰望天堂，没有人思想未来的事。但有一个时候要来临，那时连这些话也没有益处。\BibleQuote{在坟墓里，} 经上说，\BibleQuote{谁能称颂祢？} 黄金是令人渴求之物，能为我们带来许多奢华，使我们受尊敬，但不如天堂那样。因为许多人甚至避开富人，恨恶他；但对于生活有德的人，他们却尊敬并荣耀他。\Quote{但是，}有人说，\Quote{穷人被嘲笑，即使他有德行。} 不在人中，而在畜类中。因此他不应理会他们。因为如果驴对我们嘶鸣，寒鸦对我们聒噪，同时所有智者称赞我们，我们也不该忽略后者而去理会畜类的喧嚷；因为那些羡慕眼前事物的人，就像寒鸦，甚至不如驴。再者，如果一位地上的君王赞许你，你就不在乎众人，即使他们都嘲笑你；但如果宇宙的主宰赞扬你，你还要寻求甲虫和蚊蚋的好话吗？因为这些人相比天主，就是如此，或者连这也不是，而是更卑贱的东西，如果有什么更卑贱的话。我们要在泥沼中打滚多久？我们要让懒惰者和肚腹之神作我们的审判官多久？他们能证明赌徒、醉汉、为肚腹而活的人，但关于德性和恶行，他们连做梦都想不到。如果有人嘲笑你不懂得怎样开凿水渠，你不会认为是什么可怕的事，反而会嘲笑那指摘你这种无知的人；而你，当你想实践德行时，却指定那些对此一无所知的人作审判官吗？因此我们永远达不到那技艺。我们把案件交给的不是有经验的人，而是无学识的人，他们判断的依据不是技艺的规则，而是自己的无知。因此，我劝你们，让我们鄙视众人；或者不如说，让我们既不渴求赞美，也不渴求财物，也不渴求财富，也不认为贫穷是什么坏事。因为贫穷是我们智慧、忍耐和一切真智的教师。这样，\Person{拉匝禄}生活在贫穷中，得了冠冕；\Person{雅各伯}只求得到饼；\Person{若瑟}处于极度贫穷中，他不仅是奴隶，还是囚犯；因此我们更钦佩他，我们赞美他分发粮食时不如赞美他住在地牢时；不是当他戴王冠时，而是当他在锁链中时；不是当他坐宝座时，而是当被人谋害和出卖时。因此，考虑到这一切，以及争战之后为我们编织的冠冕，让我们不要羡慕财富、尊荣、奢华和权力，而要羡慕贫穷、锁链、束缚和为德行而忍受。因为那些事的结局充满烦恼和混乱，它们的命运与此世相连；但这些事的果实却是天堂，以及天上的美物，是眼睛未曾看见，耳朵未曾听见的；愿我们都能获得，赖我们的主耶稣基督的恩宠和慈爱，愿光荣归于祂，直到永远。阿们。\par

\SourceRecord{Translated by Charles Marriott.；Nicene and Post-Nicene Fathers, First Series；Vol. 14.；Edited by Philip Schaff.；Buffalo, NY: Christian Literature Publishing Co.,；1889.；<http://www.newadvent.org/fathers/240176.htm>.}

% structure: p0001 source=h1 target=section reason=page-title

\section{\WorkTitle{若望福音讲道集第七十七篇}}

% structure: p0002 source=h2 target=subsection reason=relative-heading-rank; base=h2

\subsection{\BibleRef{若 15:11-12}}

1. 一切美善之事皆有赏报，当它们到达其适切终点时；但若中途被截断，便遭遇沉船之灾。如同一艘重载的巨轮，若未能及时入港，却在中途沉没，则漫长的航行毫无所获，反而因经历更多辛劳而加重灾难。同样，那些在辛劳将近终点时退缩、在战斗中昏厥的灵魂也是如此。因此，保禄曾说：\BibleQuote{荣耀、尊崇与平安，要归于那些恒心行善而奔跑到底的人。}这正是基督如今在门徒身上所成就的。\BibleRef{罗 2:7} 因为祂已接纳了他们，他们因祂而喜乐，然而突然来临的苦难及祂悲伤的言词很可能截断他们的喜乐；在祂以足够抚慰他们的话与他们交谈之后，祂补充说：\BibleQuote{这些事我已经对你们说了，是要叫我的喜乐留在你们内，并使你们的喜乐得以圆满。} 也就是说，\Quote{使你们不至于与我分离，不至于截断你们的行程。你们曾因我而喜乐，且极为喜乐，但沮丧已落在你们身上。因此我除去这沮丧，好使喜乐最终来临，表明你们目前的处境不是痛苦，而是愉悦的缘由。我看见你们跌倒，却没有轻视你们；我没有说：\InnerQuote{你们为何不坚持高贵？}但我对你们说了带来安慰的话。因此，我愿使你们常存于同一爱内。你们听到了关于天国的事，便喜乐了。为使你们的喜乐得以圆满，我向你们说了这些话。} 然而，\BibleQuote{诫命就是：你们要彼此相爱，如同我爱了你们一样。} 你看见吗？天主的爱与我们的爱交织在一起，如同一条锁链般相连。因此，有时说两条诫命，有时只说一条。因为那抓住第一条的人，不可能不也拥有第二条。因为祂时而说：\BibleQuote{法律和先知全系于此} \BibleRef{玛 22:40}；时而说：\BibleQuote{凡你们愿意人给你们做的，你们也要照样给人做：这就是法律和先知。} \BibleRef{玛 7:12} 以及：\BibleQuote{爱是法律的满全。} \BibleRef{罗 13:10} 这正是祂在此处所说的；因为若留在（天主内）出于爱，而爱出于遵守诫命，而诫命是我们要彼此相爱，那么留在天主内便出于彼此相爱。祂不仅谈论爱，还指明了爱的方式：\BibleQuote{如同我爱了你们一样。} 祂再次表明，祂的离去并非出于仇恨，而是出于爱。\Quote{为此，我更应受钦佩，因为我为你们舍命。} 然而，祂在此处却没有直接说这些话，而是在之前的段落中，通过描绘善牧的形像来暗示；而在本段，祂则通过劝勉门徒、显示祂爱的伟大以及祂自身是谁。但为何祂处处高举爱？因为这是门徒的标记，是德行的纽带。为此，保禄作为基督的真门徒，且亲身经历过这爱，才说了如此伟大的话。\par

% structure: p0004 source=h2 target=subsection reason=relative-heading-rank; base=h2

\subsection{\BibleRef{若 15:14-15}}

那么，祂怎么说：\BibleQuote{我还有许多事要告诉你们，然而你们现在不能担负。} \BibleRef{若 16:12} 祂藉着\Quote{一切}与\Quote{听见}所表明的，无非是祂所说的没有一样是外来的，唯独是从父来的。既然谈论隐秘之事似乎是友谊最有力的证明，祂说：\Quote{你们也被认为配得上这共融。} 然而当祂说\Quote{一切}时，祂的意思是：\Quote{凡他们宜于听见的一切事。} 然后祂又提出了友谊的另一个确证，绝非寻常。那是什么样的呢？\par

% structure: p0006 source=h2 target=subsection reason=relative-heading-rank; base=h2

\subsection{\BibleRef{若 15:16}}

就是说，我主动追求了你们的友谊。祂并未止于此，而是说：\par

\BibleQuote{我拣选了你们}（就是说，\Quote{我栽种了你们}），\BibleQuote{是要你们去}（祂仍沿用葡萄树的比喻），也就是说，\Quote{是要你们扩展}；\BibleQuote{并结出果实，且使你们的果实常存。}\par

\Quote{现在，如果你们的果实常存，那么你们将更甚。因为我不仅爱了你们，} 祂说，\Quote{而且赐予了你们最大的恩惠，使你们的枝子扩展到全世界。} 你看见祂以多少种方式显示祂的爱吗？藉着告诉他们隐秘之事，藉着首先主动寻求他们的友谊，藉着赐予他们最大的祝福，藉着为他们承受祂当时所受的苦难。在此之后，祂表明祂也常与那些结果实的人同在；因为需要依靠祂的助佑，才能结出果实。\par

\BibleQuote{使你们因我的名无论向父求什么，祂必赐给你们。}\par

然而，被求者履行为所请求之事本是常理；但若父是被求者，子如何成就这事？这是为使你们明白，子并不次于父。\par

% structure: p0012 source=h2 target=subsection reason=relative-heading-rank; base=h2

\subsection{\BibleRef{若 15:17}}

那便是：\Quote{我告诉你们，我为你们舍命，或是我跑向你们，这不是为了责备，而是为了引导你们进入友谊。}然后，由于被众人迫害和侮辱是一件痛苦且无法忍受的事，足以使高傲的灵魂谦卑，因此，基督在说了万千话语之后，便进入这个主题。他先安抚了他们的心灵，然后才论及这些事，表明这些事也对他们有极大的益处，正如他曾表明先前的事一样。因为他曾告诉他们不应忧伤，反而要喜乐，\Quote{因为我去圣父那里}（他这样做不是离弃他们，而是深爱他们），所以在这里他也表明他们应当喜乐，不应忧伤。请注意他是如何做到的。他没有说：\Quote{我知道这事是痛苦的，但为了我的缘故忍受吧，因为你们也是为了我的缘故而受苦}——这个理由还不足以安慰他们；所以他略过这一点，提出了另一个理由。那是什么呢？就是这事将成为他们先前德行的确凿证据。\Quote{相反，你们应当忧伤，不是因为现在被恨，而是因为你们可能会被爱}；他以此暗示这样说，\par

% structure: p0014 source=h2 target=subsection reason=relative-heading-rank; base=h2

\subsection{若望福音 15:19}

因此，如果你们被爱，就显然表明你们显出了邪恶的标记。然后，当他先说了这个却未达到目的时，他又继续讲述。\par

% structure: p0016 source=h2 target=subsection reason=relative-heading-rank; base=h2

\subsection{若望福音 15:20}

他表明在这件事上，他们将最像他的效仿者。因为当基督仍在肉身时，人们与他争战，但他被接升天后，争战就接着临到他们身上。然后，由于他们人数稀少，害怕将要面对如此众多群众的攻击，他提升他们的心灵，告诉他们被这些人所恨是一件特别喜乐的事；\Quote{因为这样你们将有分于我的受苦。因此你们不应忧愁，因为你们并不比我更好，}正如我先前告诉你们的：\Quote{仆人不能大过主人。}还有第三个安慰的源泉：圣父也一同受辱。\par

% structure: p0018 source=h2 target=subsection reason=relative-heading-rank; base=h2

\subsection{若望福音 15:21}

也就是说，\Quote{他们也侮辱了他。}除此以外，为了剥夺那些人的借口，并加入另一个安慰的源泉，他说：\par

% structure: p0020 source=h2 target=subsection reason=relative-heading-rank; base=h2

\subsection{若望福音 15:22}

表明他们做反对他和反对他们的事，都是不义的。\Quote{那么，你为什么带我们进入这样的灾祸？你难道没有预知战争和仇恨吗？}因此他又说：\par

% structure: p0022 source=h2 target=subsection reason=relative-heading-rank; base=h2

\subsection{若望福音 15:23}

由此也预先宣布了对他们的不小惩罚。因为他们一直假装是因圣父而迫害他，为了剥夺他们的这个借口，他说了这些话。\Quote{他们无可推诿。我给了他们言语的教导，又用行为加以补充，按照梅瑟法律，那法律吩咐所有人要听从一位讲话并行事如此的人，当他既要引领人走向虔敬，又显示最大的奇迹时。}他不仅简单地说\Quote{征兆}，而是说：\par

% structure: p0024 source=h2 target=subsection reason=relative-heading-rank; base=h2

\subsection{若望福音 15:24}

对此他们自己也是证人，这样说：\Quote{在以色列从未见过这样的事}（\BibleRef{玛 9:33}）；并说：\Quote{自古以来从未听说有人开了生来就是瞎子的眼睛}（\BibleRef{若 9:32}）；拉匝禄的事也属同类，所有其他事也一样，行奇迹的方式是新的，一切都是超乎想象的。\Quote{那么，}有人说，\Quote{为什么他们既迫害你也迫害我们？}\Quote{因为你们不属于世界。如果你们属于世界，世界会爱自己的。}（\BibleRef{若 15:19}）他首先提醒他们他曾对自己的弟兄也说过的话（\BibleRef{若 7:7}）；但在那里他更多是作为反思而说，以免冒犯他们，而在这里，他却揭示了所有。\Quote{怎么明显我们是因为这个而被恨？}\Quote{从我所遭的事就知道了。告诉我，我的言行中有什么可挑剔的，使他们不接受我？}然后，由于这事对我们来说会令人惊奇，他说明了原因：就是他们的邪恶。他也没有停在这里，而是引用了先知（\BibleRef{咏 35:19}），显示他早已预言并说：\par

% structure: p0026 source=h2 target=subsection reason=relative-heading-rank; base=h2

\subsection{若望福音 15:25}

3. 保禄也是如此做的。因为当许多人惊奇犹太人怎么不信时，他引用了先知早已预言并说明了原因：他们的邪恶和骄傲是不信的原因。\Quote{既然如此：如果他们不遵守你的话，也不会遵守我们的话；如果他们迫害了你，因此他们也会迫害我们；如果他们看见了奇迹，是别人从未行过的；如果他们听见了言语，是别人从未说过的，却没有得益；如果他们恨你的圣父和你也一起恨，那么，}有人说，\Quote{你为什么差遣我们到他们中间？此后我们如何能被人相信？我们的亲族中谁会听从我们？}为了不让他们被这些想法困扰，请看他又加上了什么样的安慰。\par

% structure: p0028 source=h2 target=subsection reason=relative-heading-rank; base=h2

\subsection{若望福音 15:26-27}

\Quote{他将是值得相信的，因为他是真理之神。}因此他称他不是\Quote{圣神}，而是\Quote{真理之神}。而\Quote{从圣父所发}表明他知道一切事，正如基督也论自己说：\Quote{我知道我从哪里来，往哪里去}（\BibleRef{若 8:14}），在那里也论及真理。\Quote{他将差遣。}看，不再是只有圣父，而是圣子也差遣。\Quote{你们也是，}他说，\Quote{有权利被相信，因为你们与我同在，不是听别人说。}确实，宗徒们自信地依据这一情形说：\Quote{我们曾与他一同吃喝。}（\BibleRef{宗 10:41}）为了表明这不仅是为取悦而说，圣神为所说的话作证（\BibleRef{宗 10:44}）。\par

% structure: p0030 source=h2 target=subsection reason=relative-heading-rank; base=h2

\subsection{若望福音 16:1}

也就是说，\Quote{当你们看见许多人不信，而你们自己受虐待时。}\par

% structure: p0032 source=h2 target=subsection reason=relative-heading-rank; base=h2

\subsection{若望福音 16:2}

（因为\Quote{犹太人已经商定，若有人明认基督，就把他逐出会堂}（\BibleRef{若 9:22}））\par

\BibleQuote{是的，时候将到，凡杀害你们的，还以为是事奉天主。}\par

\BibleQuote{他们必追求杀害你们，如同追求一件虔诚中悦天主的举动。}祂再次加上安慰的话，\par

% structure: p0036 source=h2 target=subsection reason=relative-heading-rank; base=h2

\subsection{\BibleRef{若 16:3}}

\BibleQuote{你们为我的缘故，也为父的缘故忍受这些事，足以安慰你们。}在这里祂提醒他们起初所说的真福：\BibleQuote{几时人为了我的缘故辱骂你们，迫害你们，捏造一切坏话毁谤你们，你们是有福的。你们欢喜踊跃吧！因为你们在天上的赏报是丰厚的。}\BibleRef{玛 5:11}\par

% structure: p0038 source=h2 target=subsection reason=relative-heading-rank; base=h2

\subsection{\BibleRef{若 16:4}}

\BibleQuote{因此，从这些话判断，你们也应认为其余的话是可信的。因为你们不能说，我只阿谀地告诉你们那些使你们高兴的事，也不能说这些话是欺诈之言；因为一个有意欺骗的人，不会预先告诉你们那些可能使你们离开的事。所以，我预先告诉了你们，好使这些事不会猝不及防地落在你们身上，扰乱你们；还有一个其他的原因，就是你们不能说，我不知道这些事会发生。那么，请记住我告诉过你们的话。}的确，外邦人总是以他们的邪恶为借口，掩盖对你们的迫害，将他们当作败坏者驱逐出去；但这并没有困扰那些预先听到并且知道自己为何受苦的门徒。所发生之事的原因足以激发他们的勇气。因此，祂处处这样处理，说：\BibleQuote{他们没有认识我}；以及\BibleQuote{他们为了我的缘故要这样做}；以及\BibleQuote{为了我的名和父的缘故}；以及\BibleQuote{我先受了苦}；以及\BibleQuote{他们无故敢做这些事。}\par

4. 让我们也在我们的诱惑中思考这些事情，当我们因恶人遭受任何事时，\BibleQuote{仰望我们信德的创始者和完成者}（\BibleRef{希 12:2}），并思考这是出于恶人，且是为了德行，并为了祂的缘故。因为如果我们思考这些，一切都会变得非常容易和可忍受。既然一个人为他所爱的人受苦甚至会为此感到骄傲，那么为天主的缘故受苦的人又会对可怕之事有何感受呢？因为祂为了我们，甚至将那可耻的十字架称为\BibleQuote{光荣}（\BibleRef{若 13:31}），我们更应当如此心态。如果我们能如此轻视苦难，就更能够轻视财富和贪吝了。那么，当我们即将忍受任何不愉快之事时，我们不应想到劳苦，而应想到花冠；因为如同商人不仅考虑海洋，也考虑利润，我们也应想到天堂和对天主的信赖。如果获取更多似乎令人愉快，请想到基督不愿如此，那么它立刻就会显得不愉快。如果施舍穷人令您痛苦，不要止步于支出，而要立即将思绪转至播种所产生的收获；当轻视陌生女子的爱恋难以做到时，请想到奋斗之后的花冠，您就能轻易承受奋斗。因为如果恐惧能使一个人远离不当之事，基督的爱岂不更要如此？德行是艰难的；但让我们用未来应许的伟大来装点她的形式。事实上，那些有德行的人，即使没有这些应许，也看到她本身之美，并为此追求她，并工作，因为这在天主眼中为善，而不是为了报酬；他们认为清醒自持是一件大事，不是为了不受惩罚，而是因为天主命令了它。但如果有人过于软弱，就让他想到奖赏。让我们在施舍方面也如此行，让我们怜悯我们的同胞，我恳求，不要在他们因饥饿而垂死时忽视他们。否则，我们坐在餐桌前又笑又享受，而听到别人在街上哭泣经过，竟然不因他们的哭声回头，反而对他们生气，称他们为\Quote{骗子}，这岂不不妥？\Quote{你什么意思，人啊？有人会为一块面包计划骗局吗？}\Quote{是的，}有人说。那么在这种情况下，他更应被怜悯；在这种情况下，他更应被解除他的需要。或者如果您不打算施舍，至少不要侮辱；如果您不愿拯救遇难者，就不要将他推入深渊。因为请思考，当您推开来到您身边的穷人时，您将来呼求天主时会是怎样。\BibleQuote{你们用什么尺度量，也要用什么尺度量给你们。}（\BibleRef{玛 7:2}）请思考他离开时，垂头丧气，哀叹不已；除了贫穷，还受到了您傲慢的打击。因为如果您认为乞讨是一件可咒之事，那么要想乞讨一无所获，反而受辱离去，这会造成多么大的风暴？我们要像野兽一样多久，因为贪婪而不认识本性本身？许多人为这些话叹息；但我希望他们不是现在，而是永远都有这种怜悯之心。请思考，我祈祷您，那日我们将站在基督的审判台前，我们将乞求怜悯，而基督会把他们带上前来，说：\Quote{为了一块面包，一个铜钱，你们就在这些灵魂中掀起了如此大的风浪！}我们将如何回答？我们将如何辩护？要表明祂会把他们带上前来，请听祂所说的：\BibleQuote{你们没有给这些最小中的一个做，就是没有给我做。}（\BibleRef{玛 25:45}）他们将不再对我们说什么，但天主会代替他们责备我们。因为那个富人看见拉匝禄，拉匝禄也没有对他说什么，但亚巴郎为他辩护；同样，在我们现在鄙视的穷人身上也会如此。我们不会看到他们悲惨地伸出双手，而是安息；我们将继承他们的处境（但愿只是那种处境，而不是更糟的）作为惩罚。因为那个富人也不希望在那里用碎屑填饱肚子，而是被灼烧和严厉折磨，并被告知：\BibleQuote{你活着时享尽了福，拉匝禄同样也受尽了苦。}（\BibleRef{路 16:25}）因此，我们不要认为财富是什么大事；如果我们不留意，它将帮助我们走向惩罚，正如如果我们留意，贫穷也会成为我们享受和安息的增添。因为如果我们以感恩之心承受它，我们就会脱去罪过，并在天主面前获得极大的胆量。\par

5. 那么，我们不要在此世永远寻求安逸，好使我们在那里得以享受安逸；但要接受为德行而付出的辛劳，割舍多余之物，只求所需，并将我们所有的财物都施舍给需要的人。因为当天主应许我们天堂时，我们连一块饼也不肯给祂，那还有什么借口呢？当祂确实为你使太阳升起，并供应受造物的一切服务时，你却不给祂一件衣服，也不容祂分享你的屋顶？但我为什么说太阳和月亮呢？祂将自己的身体摆在你面前，将祂的宝血赐给你，难道你连自己的杯也不分给祂吗？但你这样做过一次吗？这不是怜悯；只要你有能力却不帮助，你就还没有尽到全部责任。如此，那些拿灯的童女有油，却不够充足。唉，即使你从自己的东西里施舍，也不该如此吝啬，但现在你施舍的是属于你主的东西，为何还要计较每一个小份呢？你们愿意我告诉你们这种不人道的缘由吗？当人们通过贪婪聚敛财富时，这些人施舍起来就很迟缓；因为一个学会如此获取的人，不知道如何花费。因为一个惯于掠夺的人，怎能适应其反面呢？那拿走别人东西的人，又怎能把自己的东西交给别人呢？一条惯于食肉的狗不能守护羊群；因此牧羊人会杀死这样的狗。为了不让我们遭遇同样的命运，让我们避免这样的盛宴。因为这些人也在\Quote{食肉}，当他们使人因饥饿而死时。难道你没有看见天主允许我们一切都共有吗？如果在财富中祂容许有人贫穷，那是为了富人的安慰，好使他们能通过向他们显示慈悲来除去自己的罪。但你即使在这件事上也残忍而不人道；由此可见，如果你在更大的事情上得到同样的权力，你肯定会犯下千万次谋杀，并剥夺人们的光明，乃至生命本身。为了防止这种情况发生，必需之物限制了在这方面贪得无厌。\par

如果你们听到这些事感到痛苦，那么我看到它们发生时就更加痛苦。你要富有到何时，而那人贫穷到何时？到傍晚，却不再更长；因为生命如此短暂，万物如此接近终点，此后一切都在门口，所以一切都只能被视为一小会儿。你需要胀满的仓库、众多的仆役和管家做什么？为什么你没有一万个宣扬你善举的人呢？仓库不发出声音，却会为你招来许多盗贼；但穷人的仓库却会上升到天主自己那里，并使你今世的生活甘甜，除掉你所有的罪，你将从天主获得荣耀，从人获得尊荣。那么，你为什么还要吝惜这些好东西呢？因为当你施惠于穷人时，你为穷人所作的好事，并不及你为自身所作的好事。你改善他们目前的状况；但为你自己，你却为将来的荣耀和信心预先储备。愿我们众人都能藉着我们的主耶稣基督的恩宠和慈爱获得这一切，愿光荣与权能归于祂及圣父和圣神，直到永远。阿们。\par

\SourceRecord{Translated by Charles Marriott.；Nicene and Post-Nicene Fathers, First Series；Vol. 14.；Edited by Philip Schaff.；Buffalo, NY: Christian Literature Publishing Co.,；1889.；<http://www.newadvent.org/fathers/240177.htm>.}

% structure: p0001 source=h1 target=section reason=page-title

\section{若望福音讲道集第七十八篇}

% structure: p0002 source=h2 target=subsection reason=relative-heading-rank; base=h2

\subsection{若 16:4-6}

\BibleQuote{这些事我起初没有告诉你们，因为我和你们在一起。但现在我往那派我来的那里去，你们中没有人问我：祢往哪里去？但因为我告诉你们了这些事，忧愁充满了你们的心。}\par

1. 沮丧的暴政是强大的，我们需要极大的勇气才能勇敢地抵挡这种感觉，从其中获取有益的部分，而让多余的部分离去。它也有某些有益之处；因为当我们自己或他人犯罪时，那时悲伤才是好的；但当我们遭遇人性的变迁时，沮丧便是无用的。现在，当沮丧压倒那些尚未成全的门徒时，请看基督如何藉着责备再次扶起他们。在此之前，他们曾向祂提出无数问题（因为伯多禄说：\BibleQuote{祢往哪里去？}（\BibleRef{若 13:36}）；多默说：\BibleQuote{我们不知道祢往哪里去，\newline{}怎么能知道那条路呢？}（\BibleRef{若 14:5}、\BibleRef{若 14:8}）；斐理伯说：\BibleQuote{请将父显示给我们看}）；这些人，我说，如今听到\BibleQuote{他们要把你们赶出会堂}、\BibleQuote{要恨你们}、\BibleQuote{凡杀害你们的，还以为是服侍天主}，便如此沮丧以至于哑口无言，什么也不对祂说。于是祂以此责备他们，说：\BibleQuote{这些事我起初没有告诉你们，因为我与你们同在；但现在我往那派我来的那里去，你们中没有人问我：祢往哪里去？但因为我告诉你们了这些事，忧愁充满了你们的心。}因为过度的悲伤是可怕的事，可怕且致死的。因此保禄说：\BibleQuote{免得这样的人被过多的悲伤所吞噬。}（\BibleRef{格后 2:7}）\par

\Quote{这些话}，祂说，\Quote{我起初没有告诉你们。}为什么祂起初没有告诉呢？为免有人说祂是从事件的通常进程猜度的。祂为什么又着手这样不愉快的事呢？\Quote{这些事我知道}，祂说，\Quote{从起初就知道，但没有说；不是因为我不知道，而是\InnerQuote{因为我与你们同在}。}这又是按人的方式说的，好像祂说：\Quote{因为你们处于安全之中，随时可以问我，所有风暴都落在我身上，所以起初告诉你们这些事是多余的。}但祂没有告诉他们这事吗？祂不是召叫了十二人，对他们说：\InnerQuote{你们要为了我的缘故被带到总督和君王面前}，以及\InnerQuote{他们要在会堂里鞭打你们}？\BibleRef{玛 10:18}。祂又怎么说：\InnerQuote{我起初没有告诉你们}？因为祂先前已预告过鞭打和被带到王子面前，但还没有说到他们的死会显得如此可敬，以至于行为甚至被视为对天主的服务。因为没有什么比这更足以使他们恐惧：他们将被视为不敬和败坏之人。也可以这样说，在那里祂说的是他们将从外邦人所受的苦，但在这里祂以更强的方式加上了犹太人的行为，并告诉他们这些事已近在咫尺。\par

\BibleQuote{但现在我往那派我来的那里去，你们中没有人说：祢往哪里去？但因为我告诉你们了这些事，忧愁充满了你们的心。}得知祂知道他们沮丧的程度，这对他们是不少的安慰。因为他们因祂离去的痛苦和等待将要来临的可怕之事而失魂落魄，因为他们不知道是否能勇敢地承受它们。\Quote{那么，此后祂为什么没有告诉他们已蒙赐予圣神呢？}为使你得知他们是非常有德行的。因为如果他们尚未蒙赐圣神时，虽被忧伤压倒却没有退缩，那么想想他们在享受恩宠之后会成为怎样的人。如果那时他们听了这话并因而忍受了，我们就会将一切归于圣神，但现在这完全是他们自己心境的果实，是他们爱基督的清晰表现，祂对他们的心灵（尚无助）施用了试金石。\par

% structure: p0007 source=h2 target=subsection reason=relative-heading-rank; base=h2

\subsection{若 16:7}

请注意祂如何再次安慰他们。\Quote{我说这些话，不是为了讨你们喜欢，尽管你们愁苦万分，仍必须听对你们有益的话；你们确实喜欢我与你们同在，但对你们有益的却是不同。而关心他人的人，不应在关涉他人利益的事上对朋友过分温和，或引导他们偏离有益的事。}\par

\BibleQuote{因为我若不去，护慰者就不会来。}\par

对此，那些对圣神持不当观点的人怎么说？难道是\Quote{有益}于师父离去而仆人来到吗？你看到圣神的尊荣何等大吗？\par

\BibleQuote{但我若离去，我将派遣祂到你们这里来。}那么益处是什么呢？\par

% structure: p0012 source=h2 target=subsection reason=relative-heading-rank; base=h2

\subsection{若 16:8}

那就是，\Quote{祂来以后，他们不会不受惩罚。因为已经发生的事足以堵住他们的口；但当这些事也由祂成就时，当教义更完善、奇迹更伟大时，他们看到因我名而行的事——这使得复活的证明更确实——便更当被定罪。因为现在他们能说：\InnerQuote{这是木匠的儿子，我们认识他的父母}；但当他们看到死亡的锁链被解开，邪恶被驱逐，天生的瘸腿痊愈，魔鬼被驱散，圣神丰沛地赐予，这一切都因呼求我名而实现时，他们将说什么？圣父曾为我作证，圣神也将作证。}然而圣神在起初已经作证。是的，祂现在也将这样做。但 \Quote{将指证}\par

% structure: p0014 source=h2 target=subsection reason=relative-heading-rank; base=h2

\subsection{若 16:9}

\Quote{是指将切断他们所有的借口，表明他们犯了不可宽恕的过犯。}\par

% structure: p0016 source=h2 target=subsection reason=relative-heading-rank; base=h2

\subsection{若 16:10}

这就是说：\Quote{我展示了无可指摘的生活，这证据就是\InnerQuote{我往父那里去}。}因为犹太人不断用这一点攻击祂，说祂不是出于天主，并称祂为罪人和违犯法律者，祂说，圣神也要从他们那里夺去这个借口。\Quote{因为若我之被认为不是出于天主，表明我是违犯法律者，那么当圣神显示我往那里去，不仅是暂时的，而是常在那里（因为\InnerQuote{你们不再看见我}正是宣告这一点的话），那时他们还有什么可说呢？}请注意，这两件事如何消除了他们恶意的猜疑：因为行奇迹不属于罪人（罪人不能行奇迹），而常与天主同在也不属于罪人。\Quote{所以你们再不能说\InnerQuote{这人是罪人}，\InnerQuote{这人不是出于天主}。}\par

% structure: p0018 source=h2 target=subsection reason=relative-heading-rank; base=h2

\subsection{\BibleRef{若 16:11}}

祂在这里再次提出关于正义的论证，即祂胜过了对手。假如祂是罪人，祂就不能胜过对手；这是一件连任何义人也无力做到的事。\Quote{但那些将来践踏他（魔鬼）的人，以及清楚知道我的复活——这正是审判他的标记——必晓得他因我而受审判。因为他不能扣留我。至于他们说我有魔鬼，我是骗子，这些事将来也要显为虚假；因为如果我是有罪的，我就不能胜过他；但现在他受了审判并被赶了出去。}\par

% structure: p0020 source=h2 target=subsection reason=relative-heading-rank; base=h2

\subsection{\BibleRef{若 16:12}}

\Quote{所以我去是为你们有益，若我不去，你们那时就不能担当；但若我去了，你们将来就能担当。} \Quote{究竟发生了什么事？难道圣神比祢更大，以致我们现在不能担当，而祂却能使我们担当？难道祂的工作更有能力、更完全？} \Quote{并非如此；因为祂也要讲我的话。}因此祂说：\par

% structure: p0022 source=h2 target=subsection reason=relative-heading-rank; base=h2

\subsection{\BibleRef{若 16:13-15}}

因为祂曾告诉他们：\Quote{\InnerQuote{祂要教导你们，并使你们记起}，并在你们的苦难中安慰你们}\BibleRef{若 14:26}（这一点祂自己并未做），又\Quote{\InnerQuote{我去是为你们有益}，且祂要来，以及\InnerQuote{现在你们不能担当}，但那时你们将能担当}\BibleRef{若 16:7}\BibleRef{若 16:12}，并且\Quote{\InnerQuote{祂要引导你们进入一切真理}}\BibleRef{若 16:13}；恐怕他们听了这些就以为圣神更大，从而陷入极不敬虔的见解，因此祂说：\Quote{祂要从我领受，}也就是说，\Quote{凡我对你们所说的一切，祂也要对你们说。}当祂说：\Quote{祂不凭自己说话，}意思是：\Quote{没有与我的话相反的话，没有自己的、与我的话对立的话。}正如祂论到自己说：\Quote{我不凭自己说话}\BibleRef{若 14:10}，意思是祂不说任何与父不同的话，不说任何反对或不同于父的话；同样，论到圣神也是如此。但\Quote{从我}，意思是\Quote{从我所知道的}，\Quote{从我自己的知识}；\Quote{因为我和圣神的知识是同一的。}\par

\Quote{祂要告诉你们未来之事。}祂激发了他们的心思，因为人类再也没有比知道未来更贪婪的了。例如，他们不断地问祂：\Quote{祢往哪里去？} \Quote{哪条路？}因此，为了解除他们的焦虑，祂说：\Quote{祂要预先告诉你们一切事，使你们不至于毫无预备地遭遇它们。}\par

\Quote{祂要光荣我。}如何？\Quote{祂要藉我的名施行祂的内在工程。}因为圣神降临后，弟子们将要行更大的奇迹，因此，祂再次引入尊荣的平等，说：\Quote{祂要光荣我。}\par

祂说\Quote{一切真理}是什么意思？因为祂也为圣神作证说：\Quote{祂要引导我们进入一切真理。}\BibleRef{若 16:13} 因为祂披戴着肉体，且不愿显得是在论说自己，加上他们对复活尚未明白，也过于不成熟，又因犹太人恐怕他们会以为他们惩罚祂如同惩罚违犯法律者；因此祂没有常常说伟大的话，也没有清楚地将他们从法律中拉出来。但当门徒与他们分离，以后又在外面；且当许多人将要相信并脱离罪过时；当有别的人讲论祂时，祂有理由不谈论自己的伟大。\Quote{所以这并非出于我的无知，}祂说，\Quote{而是出于听者的软弱，我才没有告诉你们我本该告诉你们的事。}为此，说过\Quote{祂要引导你们进入一切真理}之后，祂又加上\Quote{祂不凭自己说话}。因为为表明圣神不需要教导，请听保禄说：\Quote{同样，除了天主的圣神，也没有人知道天主的事。}\BibleRef{格前 2:11} \Quote{因此，正如人的灵，不须从别人学习，就知道；同样，圣神\InnerQuote{要从我领受}，也就是说，\InnerQuote{要与我的话一致地说话}。}\par

\Quote{凡父所有的，都是我的。} \Quote{既然那些是属我的，祂要从父的方面说话，祂就要从我说话。}\par

3. \Quote{但为什么圣神在祂离开之前没有来？} 因为诅咒尚未除去，罪恶尚未解除，一切都仍在惩罚之下，祂不能来。 \Quote{那么，必须除去敌意}，祂说，\Quote{与天主和好，然后接受那恩赐。} 但为何祂说，\BibleQuote{我要派遣祂}？\BibleRef{若 16:7} 意思是：\Quote{我要事先预备你们，好接受祂。} 因为，那无所不在的，怎能被\Quote{派遣}？此外，祂也显示了位格的分别。祂如此说，有两个原因：也是因为他们难以从祂被吸引开，祂劝勉他们紧紧持守圣神，好使他们珍爱祂。因为祂自己本可行这些事，却将施行奇迹的工作让给圣神，为使他们明白祂的尊威。正如圣父本可使万物存在，但圣子却如此行，为使我们明白祂的大能，此处亦然。为此，祂自己降生成人，将内在的工作保留给圣神，堵住那些将祂不可言喻的爱情作为不虔借口之人的口。因为当他们说圣子降生成人是因为祂低于圣父时，我们要回答他们：\Quote{那么，关于圣神，你们要说什么？} 祂没有取了肉身，但你们当然不会因此称祂大于圣子，也不称圣子低于祂。因此，在圣洗圣事中，也包含天主圣三。圣父能成就一切，圣子亦然，圣神亦然；然而，既然对圣父无人质疑，但质疑在于圣子和圣神，他们被包含在这礼仪中，为使他们借着共同分施那些不可言喻的恩惠，也得以完全认识他们在尊威上的共融。因为圣子能独自完成在圣洗中与圣父同能的事，圣神亦然，请听这些话被明说。因为祂对犹太人说：\BibleQuote{为叫你们知道人子在地上有赦罪的权柄} \BibleRef{谷 2:10}；又：\BibleQuote{为叫你们成为光明之子} \BibleRef{若 12:36}；以及：\BibleQuote{我赐给他们永生} \BibleRef{若 10:28} 随后：\BibleQuote{为使他们获得生命，且获得更丰富的生命} \BibleRef{若 10:10} 现在让我们看圣神也行同样的事。我们在哪里能看到呢？经上说：\Quote{但圣神显示，}经上说，\Quote{赐予每人，全是为人的好处} \BibleRef{格前 12:7}；那么赐予这些恩赐的，岂不更能赦免罪过？又：\BibleQuote{是圣神使人活}；以及：\BibleQuote{祂必藉那住在你们内的圣神，使你们复活} \BibleRef{罗 8:11}；以及：\BibleQuote{圣神因正义而成为生命} \BibleRef{罗 8:10}；以及：\BibleQuote{如果你们随从圣神引导，就不在法律权下} \BibleRef{迦 5:18} \BibleQuote{因为你们没有领受奴隶的圣神，以致再恐惧；但你们领受了义子的圣神} \BibleRef{罗 8:15} 他们当时所行的所有奇迹，都是在圣神降临时所行的。保禄写信给格林多人说：\BibleQuote{但你们已经洗净了，已经圣化了，并因我们的主耶稣基督之名，并因我们天主的圣神} \BibleRef{格前 6:11} 既然他们听过许多关于圣父的事，也见过圣子行了许多事，但对圣神还一无所知，因此那位圣神行奇迹，带来完美的知识。但（如我先前所说）为使祂不致因此被认为更大，基督如此说：\Quote{凡祂所听到的，祂都要讲出来，并要将未来的事告诉你们} 既然，若非如此，如果祂那时才要听见，并因那些正在成为门徒的人而听见，那岂不荒谬？因为按照你们的说法，祂甚至那时也不会知道，除非为了那些要听的人。还有什么比这话更不合法呢？此外，祂需要听见什么？祂不是曾借着众先知说了这一切吗？因为如果祂要教导关于法律的废除，这早已说过；如果关于基督、祂的神性和救恩计划，这些也已说过。此后祂还能说得更清楚吗？\par

\Quote{并将要显示给你们未来的事。} 这里基督最显示了祂的尊威，因为预言未来之事尤其属于天主的特性。现在，如果祂也从别人学习，祂将不比先知更多，但这里基督宣告了一种与天主精确一致的知识，即祂不可能说别的。但 \Quote{将从我领受} 意思是 \Quote{将从我领受，或是进入我血肉的恩宠，或是我也有的知识，并非需要它，也不是从别人学习，而是因为它是同一个。} \Quote{那么为什么祂这样说，而不是别样？} 因为他们还不明白关于圣神的圣言，所以祂只预备一件事：圣神应被他们相信和接受，而他们不应跌倒。因为祂曾说过：\BibleQuote{你们的导师只有一位，就是基督} \BibleRef{玛 23:10}，为使他们不认为在服从圣神时就是不服从祂，祂说：\Quote{祂的教导和我的教导是同一个；我所应教导的，那些事祂也要说。不要以为祂的话与我的话不同，因为那些话是我的，并证实我的意见。因为圣父、圣子和圣神的旨意是同一个。} 祂也愿意我们如此，当祂说：\BibleQuote{愿他们合而为一，如同你和我合而为一} \BibleRef{若 17:11}\par

4. 没有什么能与同心合意相比；因为这样，一人就成了多人。如果两个人或十个人同心合意，那一个人就不再是一个人了，而是每个人都被倍增到十倍，你会在十个人中看到一个人，也在一个人中看到十个人；如果他们有一个敌人，那攻击这一个人的，就如同攻击了十个人，必被击败；因为他不是只向一个人，而是向十个对手暴露目标。一个人会缺乏吗？不，他不会缺乏，因为他大部分是富足的，即在九个人中；而缺乏的那一小部分，被富足的较大部分遮盖了。每个人都会有二十只手、二十只眼和同样多的脚。因为他不仅用自己的眼睛看，也用别人的眼睛看；他不仅用自己的脚走，也用别人的脚走；他不仅用自己的手工作，也用别人的手工作。他有十个灵魂，因为他不仅为自己着想，那些灵魂也为他着想。如果它们变成一百个，情况也一样，他们的力量会扩展。你看到爱德的极致了吗？它如何使一个人既不可抗拒又变得众多，如何一个人甚至可以在许多地方，同在波斯和罗马，而自然所不能做到的，爱德却能做到？因为他的一部分会在这里，一部分在那里，或者更确切地说，他将完全在这里，也完全在那里。那么，如果他有一千或两千个朋友，请再想想他的力量会扩展到何处。你看到爱德是多么增长的事物吗？因为神奇的是，它使一人变成一千人。那么，我们为什么不获得这种力量，使自己处于安全中呢？这胜过一切权力或财富，胜过健康，胜过光明本身，它是勇气的根基。我们还要把爱局限于一个人或两个人多久呢？也考虑相反的情况吧。假设一个人没有朋友，这是极端愚昧的标志（因为愚昧人会这样说：\Quote{我没有朋友}），这样的人会过什么样的生活呢？尽管他极其富有，在充裕和奢华之中，拥有一万样好东西，但他却是荒凉和赤贫的。但在朋友的情况下则不然；尽管他们是穷人，却比富人更有供应；一个人不敢为自己说的话，朋友会为他说；他无法为自己获得的任何满足，将能通过另一个人获得，而且更多。这对我们将是一切享受和安全的基础，因为一个被如此多护卫保护的人不会受到伤害。因为国王的卫兵在严格上不能与这些相比。前者出于强迫和恐惧而放哨，后者出于善意和爱德；而爱德远比恐惧强大。国王害怕自己的卫兵；朋友却更信任他们胜过信任自己，并且因着他们，不惧怕任何阴谋反对他的人。那么，让我们从事这种交易吧：穷人，好在他贫困中有安慰；富人，好安全地拥有财富；统治者，好安全地统治；被统治者，好有仁慈的统治者。这是仁慈的源头，这是温柔的根基；因为即使在野兽中，那些不合群的也是最凶猛、最难以驯服的。因此我们住在城市中，有公共场所，以便彼此交往。\Person{保禄}也这样命令，说：\BibleQuote{不可放弃我们的聚会} \BibleRef{希 10:25}；因为没有比孤独、没有团结和交往的状态更大的邪恶。有人会说：\Quote{那么，}有人会说，\Quote{那些隐修士和占据山顶的人呢？} 他们也并非没有朋友；他们确实逃离了世俗的纷扰，但他们有许多同心的人，彼此紧密相连；他们退隐是为了更好地实现此事。因为事务的竞争引起许多争端，所以他们远离人群，极其精确地培养爱德。有人会说：\Quote{但如果一个人独处，}有人会说，\Quote{他怎么能有一万个朋友呢？} 就我而言，我希望，若可能的话，人知道如何彼此共处；但目前让友谊的特性保持不变。因为不是地方造就朋友。例如，他们有许多敬佩他们的人；现在这些敬佩者若不是爱他们，就不会敬佩。此外，他们为整个世界祈祷，这是友谊最大的证明。因此，我们在圣事中彼此问候，好使众多的人成为一体；对于未入门者，我们的祈祷是共同的，为病患、为世界的出产、为陆地海洋祈求。你看到爱德的一切力量了吗？在祈祷中，在圣事中，在劝勉中？这就是造成一切好事的原因。如果我们谨慎持守此爱德，我们将既妥善管理眼前的事物，也获得天国；愿我们众人因我们的主\Person{耶稣基督}的恩宠和慈爱，透过祂并与祂一起，归于\Person{圣父}及\Person{圣神}，光荣永无穷尽。阿们。\par

\SourceRecord{Translated by Charles Marriott.；Nicene and Post-Nicene Fathers, First Series；Vol. 14.；Edited by Philip Schaff.；Buffalo, NY: Christian Literature Publishing Co.,；1889.；<http://www.newadvent.org/fathers/240178.htm>.}

% structure: p0001 source=h1 target=section reason=page-title

\section{若望福音讲道集 79}

% structure: p0002 source=h2 target=subsection reason=relative-heading-rank; base=h2

\subsection{若望福音 16:16-17}

没有什么比不断思量那些带来痛苦的话语，更能使那焦虑且深陷沮丧的灵魂沮丧的了。那么，为什么基督在说了\Quote{我去}，以及\Quote{今后我不再与你们说话}之后，仍然继续谈论同一主题，说\Quote{等不多时，你们就不得见我，因为我往那派遣我的圣父那里去}呢？当祂用关于圣神的话使他们恢复之后，又再次打击他们的勇气。祂为何这样做？祂考验他们的情感，使他们更经得起考验，并使他们习惯于通过听到悲伤的事而勇敢地忍受与祂的分离。因为那些在言语上操练过这一点的人，很可能在行动上也容易在日后承受它。若有人仔细探究，这话本身就是一种安慰，就是那\Quote{我到圣父那里去}这句话。因为这是那位宣告祂不会灭亡，而是祂的结局是一种迁移的人的表达。祂又加上另一种安慰；因为祂不仅说\Quote{等不多时，你们就不得见我}，还说\Quote{等不多时，你们还要见我}；表明祂将再次回到他们那里，他们的分离只是暂时的，而祂与他们同在则是持续的。然而，他们不明白这一点。因此，人有理由惊奇，他们多次听见这些话后，仍然疑惑，好像什么都没听见一样。那么，他们为何不明白呢？依我看，或是由于悲伤，因为悲伤将所说的话从他们的理解中驱散了；或是由于言语的隐晦。因为祂在他们看来提出了两个矛盾，而其实并不矛盾。他们中一人说：\Quote{我们若看见祢，祢往哪里去呢？祢若去，我们怎能看见祢呢？}因此他们说：\Quote{我们不知道祂在说什么。}祂即将离去，他们知道；但他们不知道祂会很快回到他们那里。为此，祂责备他们，因为他们不明白祂的话。因为想要向他们灌输关于祂死亡的道理，祂说了什么？\par

% structure: p0004 source=h2 target=subsection reason=relative-heading-rank; base=h2

\subsection{若望福音 16:20}

因为他们不愿祂死亡，所以很快就相信祂不会死，然后当他们听到祂会死时，就不知所措，不明白那\Quote{片刻}是什么意思，祂就说：\Quote{你们将要痛哭哀号。}\par

\Quote{但你们的忧愁要变为喜乐。}然后，祂在表明忧愁之后有喜乐，忧愁产生喜乐，忧愁短暂而喜乐无穷之后，便转而举一个常见的例子；祂说什么呢？\par

% structure: p0007 source=h2 target=subsection reason=relative-heading-rank; base=h2

\subsection{若望福音 16:21}

祂用了一个比喻，先知们也常用这比喻，将沮丧比作分娩的剧痛。但祂所说的意思是：\Quote{产痛要抓住你们，但分娩的剧痛是喜乐的原因}；既证实祂关于复活的话，又表明离开此地如同从子宫进入光天化日。好像祂说过：\Quote{不要惊奇我带你们经过这样的痛苦而获得益处，因为即使一个母亲要成为母亲，也同样经过痛苦。}这里祂也暗示某种奥秘，即祂解除了死亡的产痛，并使之生出新人。祂不单说痛苦会过去，而且说\Quote{她甚至不记得}，因为随后而来的喜乐是如此之大；对圣人们也将如此。然而，那妇女并不是因为\Quote{一个男人来到了世上}而喜乐，而是因为给她生了一个儿子；因为若是这样，没有甚么能阻止不生育的妇女因别人生育而喜乐。那么祂为何这样说呢？因为祂引入这个比喻只是为了这个目的：表明忧愁是暂时的，而喜乐是持久的；并表明（死亡）是转入生命；以及表明他们的痛苦有极大的益处。祂没有说\Quote{一个孩子出生了}，而是说\Quote{一个人}。因为依我看来，祂在此暗示祂自己的复活，并暗示祂将不是为了那产生产痛的死亡而生，而是为了国度而生。因此祂没有说\Quote{一个孩子为她出生了}，而是说\Quote{一个人降生到世界上来了。}\par

% structure: p0009 source=h2 target=subsection reason=relative-heading-rank; base=h2

\subsection{若望福音 16:22-23}

祂又用这些话证明别的事，只证明祂来自天主。\Quote{因为那时你们将来要明白一切。}但\Quote{你们将不求问我}是什么意思？\Quote{你们不需要中保，只要呼求我的名就足够了，这样就能获得一切。}\par

\Quote{我实实在在告诉你们：你们因我的名无论向圣父求什么。}\par

祂显示祂之名的权能，既然至少既不被看见也不被呼求，而仅仅被提及，就使我们蒙父悦纳。但这在何时发生呢？在他们说\Quote{主，请看他们的恐吓，赐祢的仆人们以大胆讲论祢的道}\BibleRef{宗 4:29}、\Quote{并因祢的名行奇事}的时候。\Quote{他们所在的地方就震动了。}\par

% structure: p0013 source=h2 target=subsection reason=relative-heading-rank; base=h2

\subsection{若望福音 16:24}

2. 所以祂表明祂离开是有益的，如果至今他们什么也没有求，而那时凡他们所求的，都必得着。\Quote{因为你们不要以为我不再与你们同在，你们就被遗弃了；我的名要给你们更大的胆量。}既然祂所用的话语一直是隐晦的，祂就说：\par

% structure: p0015 source=h2 target=subsection reason=relative-heading-rank; base=h2

\subsection{若望福音 16:25}

\Quote{时候要到，那时你们要明明白白知道一切。}祂是指复活的时候。\Quote{那时，}\par

\Quote{我要明明地告诉你们关于圣父的事。}\par

（因为祂和他们在一起，四十天之久向他们显现，讲论天主国的事\BibleRef{宗 1:3}）——\Quote{因为现在你们因害怕而不留意我的话；但那时你们见我复活了，与我交谈，你们就能明明白白知道一切，因为当你们对我的信德坚定时，圣父自己也要爱你们。}\par

% structure: p0019 source=h2 target=subsection reason=relative-heading-rank; base=h2

\subsection{若望福音 16:26}

\Quote{你们对我的爱足以成为你们的辩护士。}\par

% structure: p0021 source=h2 target=subsection reason=relative-heading-rank; base=h2

\subsection{若望福音 16:27-28}

因为祂关于复活的言论，加上听闻\Quote{我出自天主，也往那里去}，给了他们非同寻常的安慰，祂便不断处理这些事。祂首先给了保证，他们相信祂是正确的；其次，他们应得安全。因此当祂说\Quote{片刻，你们将不见我；再过片刻，你们将见我}\BibleRef{若 16:17}，他们理所当然不明白祂。但现在不再如此。那么，\Quote{你们将不再问我}？\Quote{你们将不说：\InnerQuote{请将圣父显示给我们}，以及\InnerQuote{你往哪里去？}因为你们将知道一切知识，圣父将如我一样对待你们。}正是这一点特别使他们重获生气，得知他们将作圣父的朋友，因此他们说：\par

% structure: p0023 source=h2 target=subsection reason=relative-heading-rank; base=h2

\subsection{若 16:30}

你看见祂回答了暗暗藏在他们心里的问题吗？\par

\Quote{你不需要任何人问你。}\par

就是说，\Quote{在听到之前，你已知道那些使我们绊倒的事，并给了我们安息，因为你曾说：\InnerQuote{圣父爱你们，因为你们爱了我。}}经过这么多且这么大的事，他们说，\Quote{现在我们知道了。}你看见他们处在多么不完善的状态？然后，当他们好像对祂施恩似的说\Quote{现在我们知道了}，祂回答说\Quote{你们还需要许多其他事物才能达至成全；你们尚未成就任何事。你们即将把我出卖给我的仇敌，恐惧将抓住你们，甚至你们不能彼此退避，但由此我将不受任何可怕之事。}你看见祂的言语多么俯就？而且祂确实以此责备他们，他们不断需要俯就。因为当他们说\Quote{看，如今你明明白白说话，不说比喻}\BibleRef{若 16:29}，\Quote{因此我们相信你}，祂向他们表明，现在当他们相信时，他们还不相信，祂也不接纳他们的话。祂这样说，是指引他们到另一个时节。但，\par

% structure: p0027 source=h2 target=subsection reason=relative-heading-rank; base=h2

\subsection{若 16:32}

祂再次将其计入他们的账上；因为这是他们处处想知道的。然后，为显示祂说这些并非给予他们完全的知识，而是为了使他们理智不至于反叛（因为他们很可能形成一些人的想法，以为他们不应享有祂的任何助佑），祂说：\par

% structure: p0029 source=h2 target=subsection reason=relative-heading-rank; base=h2

\subsection{若 16:33}

就是说，\Quote{你们不要把我从你们的思想中驱除，而要接待我。}那么，不要让人把这些话扯入教义；它们是为了我们的安慰和爱而说的。\Quote{因为即使我们遭受了我提到的那样的事，你们的患难也不会止步于此，但你们在世上的时候将有忧愁，不仅现今我被出卖时，以后也是如此。但要振奋你们的心灵，因为你们不受任何可怕之事。当主人战胜了他的仇敌，门徒不可沮丧。}\Quote{那么如何，}告诉我，\Quote{你\InnerQuote{征服了世界}？}我已经告诉过你们，我推翻了它的统治者，但以后你们会知道，当万物都顺从并让位给你们时。\par

3. 但是，我们也可以得胜，仰望我们信德的创始者，行走在祂为我们开辟的路上。如此，死亡也不能胜过我们。\Quote{那么，我们就不会死吗？}有人说。这岂不正是清楚地表明它不能胜过我们吗？真正的斗士那时正光荣，不是当他还没有与对手交锋，而是当他交锋后没有被对方制服。所以我们不是可死的，因为我们与死亡争战，而是不朽的，因为我们得胜；如果我们永远留在死亡中，那才是可死的。正如我不称最长寿的动物为不朽的，即使它们长期免于死亡，同样，我也不称那死后复活的人为可死的，因为他曾被死亡化解。因为，请告诉我，如果一个人脸微微发红，我们是否就说他一直是红润的？不是的，因为动作不是一种习惯。如果一个人变得苍白，我们是否称他为黄疸病？不是的，因为这种状态只是暂时的。同样，你也不会称那只在死亡手中短暂停留的人为可死的。因为我们可以这样说那些睡着的人，他们可以说是死的，没有活动。但是死亡腐蚀我们的身体吗？那又怎样？这并非要它们停留在朽坏中，而是要使它们变得更好。那么，让我们征服世界，让我们奔向不朽，让我们跟随我们的君王，让我们也竖立凯旋碑，让我们轻看世界的快乐。我们不需要辛苦就能做到；让我们将我们的灵魂转移到天堂，整个世界就被征服了。如果你不渴望它，它就被征服了；如果你嘲笑它，它就失败了。我们是旅客和寄居者，那么让我们不为世上任何痛苦的事而悲伤。因为如果你出身于一个著名的国家，来自显赫的祖先，却到了一个遥远的土地，无人认识你，没有仆人，没有财富，然后有人侮辱你，你不会像在家乡受苦那样悲伤。因为清楚地知道你在一个陌生的异乡，就会说服你轻松地承受一切，轻视饥饿、口渴和任何痛苦。现在也要这样想：你是旅客和寄居者，在这异乡不要让任何事烦扰你；因为你有一座城，其建设者和创造者是天主，而寄居本身只是短暂微小的时间。让谁想打、侮辱、毁谤吧；我们在异乡，生活卑微；可怕的是，在自己的家乡，在同胞面前这样受苦，那才是最大的不体面和损失。因为如果一个人处于没有人认识他的地方，他容易忍受一切，因为侮辱因施暴者的意图而变得更痛苦。例如，如果一个人侮辱总督，明知他是总督，那么这侮辱是苦涩的；但如果他在不知道的情况下侮辱，以为他是一个普通人，那么他甚至无法触及那个受侮辱的人。我们也要这样推理。因为毁谤我们的人不知道我们是什么，即我们是天堂的公民，登记在那个天上的国度，是革鲁宾的同唱者。那么我们不要悲伤，也不要认为他们的侮辱是侮辱；如果他们知道，就不会侮辱我们了。他们认为我们贫穷卑微吗？我们也不要把这当作侮辱。因为请告诉我，如果一个旅客先于他的仆人到达，坐在客栈里等他们，然后客栈老板或一些旅客对他粗暴无礼，辱骂他，他岂不会嘲笑那人的无知？他们的错误岂不反而给他快乐？他岂不感到一种满足，仿佛被侮辱的不是他，而是另一个人？我们也要这样行。我们也是坐在客栈里，等待与我们同路的友人；当大家都聚集时，他们就会知道他们侮辱的是谁。这些人那时将垂头丧气；那时他们要说：\Quote{这就是我们}愚人\Quote{所讥笑的对象。}\BibleRef{智 5:3}\par

4. 因此，让我们用这两点来自我安慰：我们之所以不受侮辱，是因为他们不知道我们是谁；而且，如果我们想要得到满足，他们将来会给我们一个极痛苦的满足。但愿天主禁止任何人拥有如此残忍和无人性的灵魂。\Quote{那么，如果我们被自己的亲属侮辱呢？因为这才是沉重之事。}不，这反而是轻松之事。\Quote{请问为什么？}因为我们忍受我们所爱的人对我们的侮辱，与忍受我们不认识的人的方式不同。例如，在安慰那些受伤的人时，我们常说：\Quote{是一个兄弟伤害了你，要高尚地忍受；是父亲；是叔伯。}但如果\Quote{父亲}和\Quote{兄弟}的名分使你羞愧，那么我若向你提及一种比这些更亲密的关系，那就更甚了；因为我们不仅彼此是兄弟，而且还是肢体，同一个身体。现在，如果兄弟的名分使你羞愧，那么肢体的名分更甚。你们没有听过那个外邦人的谚语吗？它说：\Quote{应该容忍朋友的缺点。}你们没有听过\Person{保禄}说：\Quote{你们要彼此分担重担}吗？你们没有看到相爱的人吗？因为我被迫，既然不能从你们中取例，就将我的言论引到那个论证基础上。\Person{保禄}也这样做，这样说：\BibleQuote{此外，我们有生身的父管教我们，我们尚且敬重他们。}\BibleRef{希 12:9}或者更恰当的是他写给罗马人的话：\Quote{你们从前怎样将肢体献给不洁和不法为奴，以至于不法，现今也要照样将肢体献给正义为奴。}为此，让我们有信心地把握住这个比喻。现在，你们没有看到那些恋爱中的人吗？当他们被对娼妓的欲望点燃时，遭受何等痛苦：被掌掴、被打、被嘲笑，忍受一个转身离开并用千万种方式侮辱他们的娼妓；然而只要他们看到任何甜蜜或温柔之事，一切就都安好，所有先前的事都消失了，一切顺风顺水，无论是贫穷、疾病，还是其他任何事。因为他们把自己的生活视为悲惨或幸福，完全取决于他们所爱之人对他们的态度。他们不知道尘世的荣誉或耻辱，即使有人侮辱，他们也因从她那里获得的巨大快感和愉悦而轻松承受；即使她辱骂，即使她朝他们脸上吐唾沫，他们在忍受这些时，还以为是被抛洒玫瑰。这有什么奇怪呢？如果他们对她的感情如此，那么对于她的房子，他们也会认为比任何宫殿都更辉煌，哪怕只是泥屋，哪怕即将倒塌。但我何必说墙壁呢？当他们看到那些她晚上常去的地方时，他们会激动。现在请允许我接下来引用宗徒的话。正如他所说：\Quote{你们从前怎样将肢体献给不洁为奴，现今也要照样将肢体献给正义为奴}；同样地，我现在说：\Quote{正如我们爱过这些女子，让我们彼此相爱，我们就不会觉得遭受任何可怕之事。}为什么我说\Quote{彼此}呢？让我们如此爱天主吧。当你们听到我要求对天主的爱如同我们对娼妓所表现的那样时，你们战栗吗？但我战栗的是，我们甚至没有表现出这么多。如果你们愿意，让我们继续论证，尽管所说的话很痛苦。那个被爱的女人向她的情人不承诺任何好事，只承诺耻辱、羞愧和无礼。因为侍奉娼妓使人变得可笑、可耻、受辱。但天主向我们承诺天堂和天堂中的一切美善；祂使我们成为儿子，成为那独生子的兄弟，并在我们活着时赐给我们成千上万的事物，在我们死时赐予复活，并承诺赐予我们甚至无法想象的美善，使我们受到尊荣和敬畏。再者，那个女子迫使她的情人将全部财产耗费在坑中和毁灭中；但天主吩咐我们播种在天上，并赐予百倍收成和永生。再者，她像奴隶一样使用她的情人，发出比任何暴君更严苛的命令；但天主说：\BibleQuote{我不再称你们为仆人，而是朋友。}\BibleRef{若 15:15}\par

5. 你们岂不见此处的邪恶和彼处的福乐都过度了吗？接下来是什么呢？为了这个妇人的缘故，许多人彻夜不眠，她命令什么，他们就欣然服从；放弃房屋、父亲、母亲、朋友、金钱、庇护，把属于他们的一切都留在匮乏和荒凉中；但为了天主，或者更确切地说，为了我们自己的缘故，我们常常连财产的第三部分都不愿花费；我们看见饥饿的人，却忽视他，从赤身露体的人身边跑过，甚至连一句话也不肯对他说。但那些情人们，如果看见他们所爱之人的一个小婢女，而且她还是个外邦女子，他们就会站在市场中央，与她交谈，仿佛以此为荣、以此为乐，没完没了地说个不停；为了她的缘故，他们视自己的全部生活为无物，视统治者和统治为无物（所有经历过这种病态的人都知道这一点），她命令他们时，他们比其他人服侍时更加感激她。难道地狱不是理所应当的吗？难道万种刑罚不是理所应当的吗？那么，让我们清醒过来，让我们将别人供给娼妓的那么多、或一半、甚至三分之一，用于事奉天主。也许你们又颤栗了；我自己也是如此。但我希望你们不仅在言语上颤栗，而且在行为上颤栗；事实上，此刻我们的心确实被整理好了，但我们一出去就把一切都抛弃了。那么有什么益处呢？因为在那种情况下，如果需要花钱，没有人哀叹自己的贫困，甚至还会借钱来给，也许是在被打动的时候。但在这里，我们只要一提施舍，他们就以孩子、妻子、房屋、庇护和成千上万的借口来搪塞我们。有人说：\Quote{但是，} \Quote{那里的快乐很大。}这正是我哀叹和悲伤的。如果我证明这里的快乐更大，那又怎样呢？因为在那里，羞耻、侮辱和花费大大削减了快乐，此外还有争吵和敌意；但在这里，根本没有这些。告诉我，有什么能与这快乐相比：期待天堂和那里的天国、圣人们的光荣、和永恒的生命？有人说：\Quote{但这些事，} \Quote{是在期望中，那些是在经验中。}什么样的经验？你愿意我告诉你这里同样由经验而来的快乐吗？想想你享有何等自由，当你与德行相伴时，你如何不惧怕任何人，没有敌人、阴谋者、告密者、名誉或爱情上的竞争者、嫉妒的人、贫穷、疾病，或任何其他属人的事物。但在那里，即使万事顺遂你的心意，财富如泉水般涌来，但与竞争者的战争、阴谋和埋伏，将使与那些女人厮混之人的生活比任何人都更悲惨。因为当那可憎者傲慢无礼时，你必然要挑起争吵来奉承她。因此这比千万次死亡更痛苦，比任何刑罚更难以忍受。但这里根本没有这些。因为经上说：\BibleQuote{圣神的果实，就是爱、喜乐、平安。} \BibleRef{迦 5:22}这里没有争吵，没有不合时宜的金钱花费，没有耻辱和花费；即使你只给一个小钱、一个饼、或一杯凉水，祂也会非常感激你，祂不会做任何使你痛苦或忧伤的事，反而一切都使你光荣，使你脱离一切羞耻。因此，如果我们离开这些事，投身于相反的事，并自愿跳入燃烧的火炉中，我们还有什么可辩护的？还能得到什么宽恕？因此，我劝告那些患有此病的人，要恢复自己，回到健康，不要让自己陷入绝望。因为那个儿子也处于比这更悲惨的境地，但当他回到父亲家时，他恢复了昔日的尊荣，比那个一直讨父亲喜悦的儿子更显荣耀。让我们也效法他，回到我们的圣父面前，即使已经晚了，让我们离开那种束缚，转移到自由中，这样我们就能享受天国，借着我们的主耶稣基督的恩宠和慈爱。愿光荣归于祂，偕同圣父及圣神，世世无尽。阿们。\par

\SourceRecord{Translated by Charles Marriott.；Nicene and Post-Nicene Fathers, First Series；Vol. 14.；Edited by Philip Schaff.；Buffalo, NY: Christian Literature Publishing Co.,；1889.；<http://www.newadvent.org/fathers/240179.htm>.}

% structure: p0001 source=h1 target=section reason=page-title

\section{\WorkTitle{若望福音讲道集 80}}

% structure: p0002 source=h2 target=subsection reason=relative-heading-rank; base=h2

\subsection{若 17:1}

1. \Quote{那既实行又教导的，}经上说，\Quote{在天国里要称为大的。} 这话很有道理；因为用言语显示真智慧是容易的，但用行为来证明，却是有德性且伟大之人的事。因此基督在谈论忍受邪恶时，也以自己为例，教导我们效法祂。为此，在这训诫之后，祂转向祈祷，教导我们在诱惑中要放下一切，投奔天主。因为祂曾说：\Quote{在世上你们将有苦难}，这动摇了他们的心灵，祂便以祈祷来重新振作他们。当时他们仍把祂当作凡人来看待；祂为他们这样做，正如祂在拉匝禄的事上所做的那样，并在那里说明了原因：\BibleQuote{我说这话，是为了四周站着的群众，使他们相信是祢派遣了我。}\BibleRef{若 11:42} \Quote{是的}，有人说，\Quote{这事在犹太人身上发生得有理；但在门徒身上又为何如此？} 在门徒身上也同样有理。他们听了这么多话、见了这么多事之后，却说：\BibleQuote{现在我们晓得祢知道}\BibleRef{若 16:30}，他们最需要得到坚固。此外，福音作者甚至不把这事称为祈祷；他说的是：\Quote{祂举目向天}，而说这是与圣父交谈。若在别处提到祈祷，一面显示祂屈膝跪下，另一面显示祂举目向天，你切莫困惑；因为通过这些，我们被教导在恳求时应有的热忱，应站立举目，不仅用肉身的眼睛，更要用心灵的眼睛，并要屈膝，破碎自己的心。因为基督来不仅是为彰显自己，也是为教导那不可言说的德性。但教师应当不仅用言语，更要用行动来教导。让我们听祂在此处所说的话。\par

\BibleQuote{父啊，时辰来到了，请光荣祢的子，好使子也光荣祢。}\par

祂再次向我们显示，祂并非不情愿地走向十字架。因为一个祈祷这事成就、且称这行动为\Quote{光荣}——不仅为被钉的祂自己，也为圣父——的人，怎会不情愿呢？事实正是如此，因为不仅圣子，圣父也受了光荣。因为钉十字架之前，连犹太人也未曾认识祂；经上说：\BibleQuote{以色列却不认识我}\BibleRef{依 1:3}；但钉十字架之后，普世都奔向祂。然后祂也论及这光荣的方式，以及祂将如何光荣圣父。\par

% structure: p0006 source=h2 target=subsection reason=relative-heading-rank; base=h2

\subsection{若 17:2}

因为恒常行善，就是天主的荣耀。但\BibleQuote{祢赐给祂权柄掌管一切血肉}是什么意思？祂现在显示，这宣讲的对象不仅限于犹太人，而是扩展到普世，并预先为外邦人设下初步的邀请。既然祂曾说过：\BibleQuote{外邦人的路你们不要走}\BibleRef{玛 10:5}，此后又要说：\BibleQuote{你们要去使万民成为门徒}\BibleRef{玛 28:19}，祂便指明圣父也愿意如此。因为这事大大得罪了犹太人，也得罪了门徒；事实上此后他们也不容易忍受去抓住外邦人，直到他们领受了圣神的教导；因为这对犹太人成了不小的绊脚石。因此，当伯多禄在圣神如此显现后回到耶路撒冷时，他几乎无法借着叙述那块布的异象来免除对他的指控。但\BibleQuote{祢赐给祂权柄掌管一切血肉}是什么意思？我要问异端者：\Quote{祂何时领受了这权柄？是在祂形成他们之前，还是之后？} 祂自己说是在祂被钉十字架并复活之后；至少那时祂说：\BibleQuote{天上地下的一切权柄都给了我}\BibleRef{玛 28:18}，以及\BibleQuote{你们去使万民成为门徒}\BibleRef{玛 28:19}。那么，难道祂对自己的工程没有权柄吗？祂创造了它们，难道在创造之后对它们没有权柄吗？然而，我们看到祂在古时已经行使一切，惩罚某些罪人（因为经上说：\BibleQuote{我岂能隐瞒我的仆人亚巴郎我要做的事}\BibleRef{创 18:17}），并尊荣其他义人。那么，祂当时就有这权柄，到这时却失去了，后又再次领受？哪个魔鬼会如此断言？但若祂的权柄当时和现在都一样（因为祂说：\BibleQuote{就如父唤起死者并使他们复生，同样，子也使他所愿意的人复生}\BibleRef{若 5:21}），那么这些话有什么意义呢？祂将要派遣他们到外邦人那里；为使他们不认为这是一项革新，因为祂曾说过：\BibleQuote{我奉派遣，只到以色列家迷失的羊那里去}\BibleRef{玛 15:24}，祂便显示这事也中悦圣父。若祂用非常卑屈的言词说这话，这并不奇怪。因为祂如此建立当时的人以及后来的人；正如我先前所说，祂常借过度的卑屈使他们坚定相信这些话是迁就人的话。\par

2. 但\Quote{一切血肉}是什么意思？当然不是所有人都相信了。然而，就祂而言，所有人都相信了；若人们不听祂的话，错不在教师，而在那些不接受的人身上。\par

\BibleQuote{好使祂将永生赐给祢所交给祂的人。}\par

若在此处祂也以更人性化的方式说话，请不要惊奇。因为祂这样做，既是为了我已给出的理由，也是为了避免说任何关于自己的大话；因为这对听众来说是个绊脚石，因为他们那时对祂还没有什么伟大的想象。例如，\Person{若望}以自己的身份说话时，就不是这样，而是将他的言辞引向更高的崇高境界，说：\Quote{万物是藉着祂造成的；凡受造的，没有一样不是藉着祂造成的}\EditorialNote{若 1:3, 4, 9, 11}；并且祂是\Quote{生命}，是\Quote{光}，祂\Quote{来到自己的地方}；他没有说，若不是领受了，祂就没有权能，而是说祂也赐给别人\Quote{成为天主的子女的权能}。\Person{保禄}也同样称祂与天主同等。但祂自己以更人性化的方式发问，这样说：\Quote{叫祂将永生赐给祢所赐给祂的人。}\BibleRef{斐 2:6}\par

% structure: p0011 source=h2 target=subsection reason=relative-heading-rank; base=h2

\subsection{\BibleBook{若望}{福音} 17:3}

\Quote{唯一真天主}，祂说，是为了与那些不是神的有所区别；因为祂即将派遣他们到外邦人那里去。但如果他们不允许这一点，却因这\Quote{唯一}一词而拒绝承认子是真实的天主，那么按照他们的方法，他们也拒绝承认祂是天主。因为祂也说：\Quote{你们不寻求那来自唯一的天主的荣耀。}\BibleRef{若 5:44}那么，难道子就不是天主吗？但若子是天主，子是被称为唯一天主的父的子，那么显然祂也是真实的，且是被称为唯一真天主的父的子。为什么当\Person{保禄}说\Quote{只有我和\Person{巴尔纳伯}}\BibleRef{格前 9:6}时，他排除了\Person{巴尔纳伯}？绝不是；因为\Quote{只有}是为了与其他人区别而说的。而如果祂不是真天主，祂又怎么是\Quote{真理}呢？因为真理远超乎真实。如果祂不是\Quote{真实}的人，我们该称之为什么，告诉我？难道我们不称之为根本不是人吗？所以，如果子不是真天主，祂又怎么是天主？祂又怎么使我们成为神和子女，如果祂不是真实的？但关于这些事，我们在别处已更详细地讨论过；因此，让我们专注以下的内容。\par

% structure: p0013 source=h2 target=subsection reason=relative-heading-rank; base=h2

\subsection{\BibleBook{若望}{福音} 17:4}

\Quote{我完成了祢所交给我的工作，要我做的。}\par

然而，这行动才刚刚开始，或者甚至还没有开始。那么祂怎么说\Quote{我完成了}呢？要么祂的意思是\Quote{我已尽了我的一份}；要么祂说未来之事，如同已经发生；或者，最可以说的是，一切都已生效，因为福泽的根基已经奠定，其果实必定会随之而来，并且因为祂亲身参与并协助那些此后将要发生的事。为此，祂又以迁就的方式说：\Quote{祢所交给我的。}因为如果祂真的等待聆听和学习，那就有损于祂的光荣。因为祂是出于自己的意愿来做这事，从许多经文中可以清楚看到。如\Person{保禄}所说：\Quote{祂如此爱了我们，以至为我们舍弃了自己}\BibleRef{弗 5:2}；以及\Quote{祂空虚了自己，取了奴仆的形体}\BibleRef{斐 2:7}；以及\Quote{父怎样爱了我，我也怎样爱了你们。}\BibleRef{若 15:9}\par

% structure: p0016 source=h2 target=subsection reason=relative-heading-rank; base=h2

\subsection{\BibleBook{若望}{福音} 17:5}

那荣耀在哪里？因为即使祂在人类中不被尊荣是合理的，因为包裹祂的遮盖物；祂为何寻求与父同受光荣？那么祂在这里说什么？这话指的是降生奥迹；因为祂的肉性尚未得荣耀，尚未享受不朽，也尚未分享王座。因此祂没有说\Quote{在地上}，而是说\Quote{在祢面前}。\par

3. 按我们的尺度，我们也将享有这光荣，只要我们醒寤。因此保禄说：\BibleQuote{只要我们与祂一同受苦，也必与祂一同受光荣。}\BibleRef{罗 8:17} 那么，面对如此荣耀的赏报，却因懈怠和昏睡而自绝于它的人，即使没有地狱，也当受万千眼泪；因为在他们有能力与天主子一同为王、一同受光荣时，竟剥夺了自己如此伟大的福乐。倘若需要被肢解、死千万次、每日舍弃千万条性命和千万具身体，难道我们不应当为这样的光荣承受这一切吗？但现在我们甚至轻视金钱，这金钱日后我们虽不情愿也要留下；这金钱给我们带来千万祸患，留在此世，不属于我们。因为我们不过是那不属我们之物的管家，尽管我们从祖先那里接受它。然而，除了地狱、不死的虫、不灭的火和咬牙切齿的哀号之外，我们又将如何承受这一切？我们何时才肯清醒地看清，将全部精力耗费于日常争斗、口角和无益的闲谈，喂养、耕种土地，养肥身体而忽略灵魂，不为必要之事操心，却为多余无益之事焦虑？我们建造华丽的坟墓，购置昂贵的房屋，随从成群结队的仆役，设立各式各样的管家，任命土地、房屋、钱财的管理者，甚至为这些管理者再设管理者；而对于我们那孤寂的灵魂，却毫不关心。这要到何时为止？我们填满的不正是一个肚腹吗？我们穿戴的不正是一具身体吗？这般忙碌究竟为何？为何要为服侍这些事物而将交托给我们的独一灵魂切割撕裂，为自己制造沉重的奴役？因为需要许多事物的人便是许多事物的奴仆，尽管他看似是它们的主人。主人甚至成为他仆役的奴仆，引入另一种更沉重的服役方式；从另一意义上他也是他们的奴仆，因为没有他们他不敢进入广场、浴场或田野，而他们却常常在无他陪伴时四处走动。那看似主人的人，若没有仆人在场，便不敢出门；若无人随从而只是将头探出屋外，便以为被人嘲笑。或许有人说这话时会嘲笑我们，但正因如此，他们才当受万千眼泪。为证明这是奴役，我愿问你们：你们希望有人将食物送到你们嘴里，将杯盏递到你们唇边吗？你们不认为这种服侍值得流泪吗？如果你们需要常有人搀扶才能行走，你们不认为自己可怜，在这一方面比任何人都更悲惨吗？那么现在你们也应当如此感受。因为无论是被无理性之物如此对待，还是被人如此对待，都没有区别。\par

请告诉我，天使与我们不同之处不正在于此吗？他们不像我们一样需要如此多的事物。因此，我们需要得越少，就越接近他们；我们需要得越多，就越沉沦于这朽坏的生命。为使你们明白确是如此，请问那些年长者：他们认为哪种生活最幸福？是受欲望支配的时候，还是如今能主宰这些事物的时候？我们提到这些人，因为那些被青春冲昏头脑的人甚至不知道他们奴役的极端程度。那么，那些发高烧的人呢？他们在极度口渴时大量饮水、越发需要时，自认为是幸福的吗？还是在康复之后、不再有这种欲望时？你们看到，在任何情形下，需求多都是可悲的，远非真智慧，反而加剧了奴役和欲望。那么，我们为何要自愿增加自己的苦恼？请告诉我，如果有可能不受伤害地生活在没有屋顶和墙壁的地方，你们难道不会选择这样吗？为何还要增加自己软弱的标志？我们岂不正是因为亚当什么都不需要——没有房屋、没有衣服——而称他为幸福的吗？有人说：\Quote{是的，但现在我们需要这些东西。}那么，为什么我们还要使自己的需求更大？如果许多人缩减了实际所需的许多东西（我是指仆役、房屋和金钱），我们若超出必需，还有什么借口？你们越把东西围在身边，就越成为奴隶；因为你们需要得越多，就越侵犯了自己的自由。绝对自由是什么也不需要；其次是需要得少；天使和他们的效仿者尤其拥有这一点。而人尚在血肉之躯中能达到这一点，想想这是何等大的赞美！这也正是保禄写给格林多人时所说的：\BibleQuote{但我要宽待你们}和\BibleQuote{免得这样的人在肉身方面受苦。}\BibleRef{格前 7:28} 财富被称为\Quote{可用之物}，是为叫我们善\Quote{用}它们，而非保存或埋藏；因为这并非拥有它们，而是被它们所拥有。如果我们以此为目标——如何增加财富，而非如何善用——那么秩序就颠倒了，是它们拥有我们，而非我们拥有它们。让我们摆脱这沉重的束缚，终于获得自由吧。为什么我们要为自己造出千万条锁链？难道自然本性的束缚、生活的必需和万千事务的纷扰还不够吗？你们还要为自己编织别的网罗，缠住双脚？你们何时才能抓住天堂，站立在那高处？因为有一件大事，一件大事，即使是斩断所有这些绳索，你们才能抓住那天上的城邑。还有如此多的阻碍；为战胜这一切，让我们持守中等的生活［并除去多余之物，持守必需之物］。这样，我们必能得着永生，赖我们的主耶稣基督的恩宠和慈爱，愿光荣归于祂，至于无穷之世。阿们。\par

\SourceRecord{Translated by Charles Marriott.；Nicene and Post-Nicene Fathers, First Series；Vol. 14.；Edited by Philip Schaff.；Buffalo, NY: Christian Literature Publishing Co.,；1889.；<http://www.newadvent.org/fathers/240180.htm>.}

% structure: p0001 source=h1 target=section reason=page-title

\section{关于若望福音的讲道第八十一篇}

% structure: p0002 source=h2 target=subsection reason=relative-heading-rank; base=h2

\subsection{\BibleRef{若 17:6}}

1. \Quote{伟大的谋士的使者}\BibleRef{依 9:6}，天主子被称为如此，是因为祂所教导的其他事，尤其因为祂将圣父启示给世人，正如现在祂所说：\Quote{我已将你的名显示给这些人。}因为在说了\Quote{我完成了你的工作}之后，祂接着详细解释，说明是什么样的工作。这名原是众所周知的，因为依撒意亚说过：\Quote{你要指着真实的天主起誓。}\BibleRef{依 65:16}。但我屡次告诉你们的事，现在仍告诉你们：虽然这名已知，但仅限于犹太人，而且并非所有犹太人都知道；但现在祂是论及外邦人。祂不仅声明这一点，还表明他们认识祂就是圣父。因为知道祂是创造主，与知道祂有圣子，并非同一回事。而祂\Quote{显示了他的名}，既通过言语，也通过行动。\par

\Quote{你从世上赐给我的人。}正如祂以上所说：\Quote{除非蒙赐予，没有人能到我这里来}\BibleRef{若 6:65}；以及\Quote{除非我的父吸引他}\BibleRef{若 6:64}；同样在这里：\Quote{你赐给我的人。}\BibleRef{若 14:6}现在祂自称是\Quote{道路}；由此可见，祂在此所说的话确立了这两件事：祂与圣父并非对立，并且圣父的旨意是将他们托付给圣子。\par

\Quote{他们原属于你，你将他们赐给了我。}在此祂愿意教导我们，祂深受圣父所爱。因为祂并不需要接收他们，这一点很清楚：祂创造了他们，祂不断照顾他们。那么祂如何接收他们呢？正如我先前所说，这表明祂与圣父同心合意。若有人愿意以人的方式，并按字面来探讨此事，那么他们就不再属于圣父。因为如果当圣父拥有他们时，圣子没有他们，那么显然当圣父将他们赐给圣子时，祂便放弃了对他们的主权。此外，还有一个更不体面的结论：他们会发现在他们尚与圣父同在时是不完美的，而到了圣子那里才变得完美。但这样说简直是戏谑。那么祂借此声明什么呢？\Quote{圣父也认为他们应该相信圣子。}\par

\Quote{他们也遵守了你的话。}\par

% structure: p0007 source=h2 target=subsection reason=relative-heading-rank; base=h2

\subsection{\BibleRef{若 17:7}}

他们如何\Quote{遵守了你的话}？\Quote{因着相信我，并不理会犹太人。因为相信祂的人，经上说：\InnerQuote{已经印上印，证明天主是真实的。}}\BibleRef{若 3:33}有些经文读作：\Quote{现在我知道，凡你所赐给我的，都是从你而来的。}但这并无理由；因为圣子岂会不知道圣父的事？不，这话是对门徒说的。\Quote{从我告诉他们这些事的时候，}祂说，\Quote{他们就知道，凡你所赐给我的，都是来自你；没有什么外在于你，没有什么是我独有的。}（因为凡是独有的，大多会使事物成为外来的。\Quote{因此他们知道，凡我所教导的，都是你的道理和教训。}）\Quote{他们从哪里学到这个？}从我的话；因为我就是这样教导他们的。而且我不仅教导他们这个，还教导他们\Quote{我是从你出来的。}因为通过整个福音，祂热切地要证明这一点。\par

% structure: p0009 source=h2 target=subsection reason=relative-heading-rank; base=h2

\subsection{\BibleRef{若 17:9}}

\Quote{祢说什么？} \Quote{祢教导圣父，好像祂不知道？祢对祂说话，好像对人不知道？} \Quote{那么这区别是什么意思？}你难道看不出，这祈祷无非是为了使他们理解祂对他们的爱？因为祂不仅将自己所有的赐予他们，还呼求另一位也这样做，这显出了更大的爱。那么，\Quote{我为他们祈求}是什么意思？\Quote{不是为整个世界，}祂说，而是\Quote{为你所赐给我的人。}祂不断提到\Quote{所赐给}，好让他们知道，这也中悦圣父。然后，因为祂不断说：\Quote{他们属于你}，和\Quote{你把他们赐给了我}，为了消除任何恶意的猜疑，免得有人以为祂的权柄是后来的，是现在才领受的，祂说了什么呢？\par

% structure: p0011 source=h2 target=subsection reason=relative-heading-rank; base=h2

\subsection{\BibleRef{若 17:10}}

你看到尊荣的平等吗？因为免得你听到\Quote{你把他们赐给了我}，就认为他们脱离了圣父的权柄，或在此之前脱离了圣子的权柄，祂通过这样说消除了这两个难题。祂好像在说：\Quote{当你听到\InnerQuote{你把他们赐给了我}时，不要以为他们脱离了圣父，因为凡是我的，也是祂的；当你听到\InnerQuote{他们原属于你}时，不要以为他们与我是疏远的，因为凡祂的，也是我的。}因此，\Quote{你赐给}的说法，只是为了屈就；因为圣父所有的，也是圣子的，圣子所有的，也是圣父的。但即使按人的方式，也不能对儿子这样说，因为他们处于更大的尊荣平等中。因为凡属于较小的，也属于较大的，这是人人都明白的，但反过来则不成立；但在这里，祂转换了这些术语，而这种转换表明了平等。在另一处，祂说：\Quote{凡父所有的，都是我的}，以此声明这一点，这是指知识。而\Quote{你赐给我}以及类似的表达，是为了表明祂并非作为外人而来，将他们吸引到自己身边，而是作为自己的财产接收他们。然后祂陈述了原因和证据，说：\Quote{我在他们身上受了光荣}，意思是：\Quote{我对他们有权威}，或者\Quote{他们将荣耀我，相信你和信我，并将同样荣耀我们。}但如果祂没有在他们身上同样受光荣，那么圣父所有的就不再是祂的了。因为没有人会在他没有权威的人身上受光荣。然而，祂如何同样受光荣呢？所有人都同样为祂而死，如同为圣父而死；他们传扬祂，如同传扬圣父；并且正如他们说一切事都是因圣父的名而成就，同样也是因圣子的名。\par

% structure: p0013 source=h2 target=subsection reason=relative-heading-rank; base=h2

\subsection{\BibleRef{若 17:11}}

那就是：\Quote{虽然我不再以肉身显现，却藉这些事受光荣。}但祂为何连续地说：\Quote{我不在这世界上了}；以及\Quote{因为我离开他们，我把他们托付给祢}；以及\Quote{当我在世界上的时候，我保守了他们}？因为若人按字面理解这些话，就会产生许多荒谬的结论。因为说祂不再在世界上了，并且当祂离去时将门徒托付给另一位，这怎能合理？这不过是一个凡人永远离开他们的话语。你看出祂多半像人一样说话，并且以一种适应他们心境的方式吗？因为他们以为，由于祂的同在，他们享有更大的安全。因此祂说：\Quote{当我与他们在一起时，我保守了他们。}\BibleRef{若 14:28}然而祂又告诉他们：\Quote{我往你们那里去}；以及\Quote{我同你们天天在一起，直到今世的终结。}\BibleRef{玛 28:20}那么祂怎么说这些话，仿佛要离开他们似的？正如我先前所说，祂是迁就他们的思想，好让他们听见祂这样说话并把他们托付给父的照顾时，稍得喘息。因为，既然他们听了祂许多劝诫后仍不信服，祂便转而与父交谈，显明祂对他们的爱。仿佛祂说：\Quote{既然祢召我归祢自己，请保护这些人安全；因为我往祢那里去。}\Quote{祢说什么？祢不能保守他们吗？}\Quote{是的，我能。}\Quote{那么祢为何这样说？}\Quote{为使他们拥有我的喜乐得以圆满}\BibleRef{若 17:13}；也就是说，\Quote{不致因不完全而受困扰。}祂以此表明，祂说这一切是为了使他们得到安宁和喜乐。因为这话看起来矛盾：\Quote{现在我不再在世界上了，而他们却在世界上。}这正是他们所猜疑的。因此，祂一时俯就他们，因为倘若祂说：\Quote{我保守他们}，他们就不会如此相信；因此祂说：\Quote{圣父啊，请因祢的名保守他们}；也就是说，\Quote{藉祢的帮助。}\par

% structure: p0015 source=h2 target=subsection reason=relative-heading-rank; base=h2

\subsection{若望福音 17:12}

祂再次像人和先知一样说话，因为祂从未藉天主的名做过任何事。\par

\BibleQuote{祢所赐给我的人，我保守了他们，其中除了丧亡之子没有丧失一个，为应验经上的话。}\par

在另一处祂说：\BibleQuote{凡祢所赐给我的，我必不失落任何一个。}\BibleRef{若 6:39}然而不仅他失落了，以后还有许多人也失落了；那么祂怎么说\Quote{我决不失落}呢？\Quote{就我而言，我不会失落。}所以在另一处，祂更清楚地说明此事，说道：\BibleQuote{我决不把他抛弃在外面。}\BibleRef{若 6:37}\Quote{并非由于我的过失，不是因为我的唆使或抛弃；但如果他们自行离开，我并不强留他们。}\par

% structure: p0019 source=h2 target=subsection reason=relative-heading-rank; base=h2

\subsection{若望福音 17:13}

你看出这言论是以一种更属人的方式构成的吗？因此，若有人想从这些话中贬抑圣子，他也必贬抑圣父。注意，为证明这一点，祂如何从起初一面好像是在告知并解释给祂，一面是命令。告知，如祂说：\Quote{我不为世界祈求}；命令，如\Quote{我至今保守了他们}，\Quote{他们没有丧失一个}；并且\Quote{所以现在求祢保守他们}，祂说。又说：\Quote{他们原是祢的，祢把他们赐给了我}；以及\Quote{当我在世界上的时候，我保守了他们。}但这一切的解释是：这些话是针对他们的软弱而说的。\par

但在说了\BibleQuote{其中除了丧亡之子没有丧失一个}之后，祂又加上\BibleQuote{为应验经上的话。}祂说的是什么经上的话？是那个预言了许多关于祂的事的经上的话。并不是祂因那话而丧亡，为使经上的话得以应验。但关于这一点，我们以前已详细谈论过，即这是圣经的特有笔法，它把照着它发生的事，仿佛是由于它而发生。并且，如果我们不想得出错误结论，就必须准确地考察一切，包括说话者的方式、祂的论证以及圣经的法则。因为：\BibleQuote{弟兄们，在心思上不要作小孩子。}\BibleRef{格前 14:20}\par

3. 这事须仔细思量，不仅为明了圣经，也为在生活方式上奋发努力。因为幼童不向往大事，却常羡慕那些毫无价值之物；他们喜欢看陶土制成的战车、马匹、赶骡者与轮子；但若看见一位君王端坐于战车，配以一对白骡与盛大威仪，他们甚至不屑一顾。他们用同样的材料打扮玩偶如新娘，却对真实美丽的新娘视而不见；在其他许多事上也是如此。如今许多成年人也遭遇此景；因为当听见天上的事，他们毫不留意，却如孩童般热切于泥土之物，愚蠢地崇拜地上的财富，推崇现世的光荣与奢华。然而这些与孩童的玩具无异；前者却是生命、光荣与安息之源。但如同孩童被夺去玩具便哭泣，不知甚至渴望真实之物，许多貌似成人者亦然。为此经上记载：\BibleQuote{不要作小孩子在心智上。}\BibleRef{格前 14:20} 告诉我，你贪求财富，难道不是渴望那长存的财富，而是孩童的玩物？若看见一人欣赏一枚铅币并弯腰拾起，你必断言其极度贫穷；而你，收集比这更无价值之物的人，竟自视为富足？这岂合理？我们要称那轻视现世一切的人为富足。因为除非渴求更大的事，没有人会选择嘲笑这些小事——金银与其他炫耀之物；正如一个人若无金币，便不会轻视铅币。因此，当你看见一人弃绝一切世俗事物时，你要认为他如此行别无他因，只因他向往更大的世界。农夫轻视几粒麦子，是因期望更大的丰收。若希望不确定时我们尚且轻视眼前之物，何况期望确凿之时？因此我恳求并劝诫你们：不要自招损失，不要紧握泥泞，却夺去自己天上的宝藏，使你们的船满载稻草与糠秕进入港口。任凭各人如何论说我们，对我们不断的劝诫发怒，称我们为愚拙、烦琐、令人厌烦，我们仍不会停止在这些事上持续劝勉你们，并不断向你们重复先知的话：\Quote{\BibleQuote{你要借施舍除去你的罪，借怜悯穷人除去你的不义} \BibleRef{达 4:27}，并将它们系在你的颈上。} 不要今天如此行，明天便作罢。因为身体需要每日饮食，灵魂亦然，甚至更需要；若不供给，它就变得更软弱、更卑下。我们切莫在其衰亡窒息之时忽略它。灵魂每日受到许多创伤：因情欲、愤怒、懒惰、辱骂、报复、嫉妒。因此也必须为它预备良药，而施舍绝非小药，可以敷于每一创伤。因为经上说：\BibleQuote{把你们所有的施舍给人，看，一切对你们都是洁净的。}\BibleRef{路 11:41} \Quote{施舍}，而非贪婪，因为出于贪婪的施舍不能持久，即使你给予需要者。因为施舍是远离一切不义的事，\Quote{这}使一切洁净。这事甚至胜过禁食与卧地；后者可能更苦更劳，但前者更为有益。它光照灵魂，使之润泽、美丽而健壮。橄榄之果不能如此扶持运动员，如同这油恢复虔诚的斗士。让我们涂抹双手，好能有力地举起对抗仇敌。那实行怜悯需要者的人，不久将停止贪婪；那持续周济穷人者，不久将停止愤怒，甚至不会自高。因为如同医生不断照顾伤员，因他人灾祸而看见人性，便容易清醒；同样，我们若从事救助穷人的工作，必容易成为真正智慧，不羡慕财富，不看重现世之事，反轻视一切，高飞向天，必轻易获得永恒的福乐，借着我们的主耶稣基督的恩宠与慈爱；愿光荣归于祂，及圣父与圣神，至于无穷世。阿们。\par

\SourceRecord{Translated by Charles Marriott.；Nicene and Post-Nicene Fathers, First Series；Vol. 14.；Edited by Philip Schaff.；Buffalo, NY: Christian Literature Publishing Co.,；1889.；<http://www.newadvent.org/fathers/240181.htm>.}

% structure: p0001 source=h1 target=section reason=page-title

\section{\WorkTitle{若望福音讲道集第八十二篇}}

% structure: p0002 source=h2 target=subsection reason=relative-heading-rank; base=h2

\subsection{\BibleRef{若 17:14}}

1. 当我们修德成善却遭受恶人迫害，或当我们渴慕德行而被恶人嘲笑时，我们不可心烦意乱或动怒。因为这是事物的自然规律，美德处处都会招致恶人的憎恨。那些不愿正当地生活的人，嫉妒那些愿意善度生活的人，并为自己寻找借口，以为若能推翻他人的美誉，就能为自己开脱；他们憎恨这些善人，因为他们的追求与自身相反，并想尽办法羞辱他们的生活方式。但我们不要忧伤，因为这是美德的标志。因此，基督也说：\BibleQuote{你们若属于世界，世界必爱属于自己的人。}（\BibleRef{若 15:19}）在另一处又说：\BibleQuote{几时众人都夸赞你们，你们是有祸的。}（\BibleRef{路 6:26}）因此祂在此处也说：\BibleQuote{我将祢的话赐给了他们，世界却憎恨了他们。}祂又说明他们为何堪当蒙受圣父极大的眷顾：\BibleQuote{为了祢的缘故，}祂说，\BibleQuote{他们被憎恨，也为了祢言语的缘故。}因此他们应得一切眷顾。\par

% structure: p0004 source=h2 target=subsection reason=relative-heading-rank; base=h2

\subsection{\BibleRef{若 17:15}}

祂再次使祂的话变得简单；祂再次使之更加清晰。这正表明，祂如此精确地为他们恳求，无非是要显明，祂对他们怀有极其温柔的关怀。然而祂自己曾告诉过他们，凡他们所求的，圣父必会全部成就。那么祂在此为何还为他们祈祷呢？正如我说过的，无非是为了显示祂的爱。\par

% structure: p0006 source=h2 target=subsection reason=relative-heading-rank; base=h2

\subsection{\BibleRef{若 17:16}}

可是祂在另一处说：\BibleQuote{祢从世上赐给我的人，他们原是祢的}（\BibleRef{若 17:6}）？那里祂论及他们的本性；此处则论及邪恶的行为。祂又特意列举了他们的众多优点：首先，\BibleQuote{他们不属于世界}；其次，\BibleQuote{圣父将他们赐给了祂}；第三，\BibleQuote{他们遵守了祂的话}；第四，他们因此\BibleQuote{被憎恨}。如果说：\BibleQuote{如同我不属于世界}，请不要困扰，因为这里的\Quote{如同}并非表示严格等同。因为在论及祂与圣父时，用\Quote{如同}表明极大的平等，这是由于本性上的关系；而在论及我们与祂时，差距是极大的，因为各自的本性之间存在着巨大而无限的差距。因为祂\BibleQuote{没有犯过罪，口中也找不到诡诈}（\BibleRef{伯前 2:22}），宗徒们又怎能与祂相等呢？那么祂所说的\BibleQuote{他们不属于世界}是什么意思？\BibleQuote{他们向往另一个世界，与大地毫无共同之处，而成了天上的公民。}祂以这些话显示祂的爱，因为祂将他们托付给圣父，并把他们交托给那生祂的一位。当祂说\BibleQuote{保守他们}时，不仅是说使他们免于危险，也是指他们持守信德。为此祂接着说：\par

% structure: p0008 source=h2 target=subsection reason=relative-heading-rank; base=h2

\subsection{\BibleRef{若 17:17}}

\BibleQuote{祢的话就是真理。}\par

这就是说，\BibleQuote{其中毫无虚假，凡在其中所说的，必定成就}；又说，它不表示任何预象或肉体的事物。如圣保禄论及教会时所说，祂藉着圣言祝圣了教会（\BibleRef{弗 5:26}）。而且，\BibleQuote{求祢以真理祝圣他们}，在我看来还意味着别的事，如同：\BibleQuote{将他们分别为圣，用于圣言和宣讲。}这从下文可以显明。因为祂说：\par

% structure: p0011 source=h2 target=subsection reason=relative-heading-rank; base=h2

\subsection{\BibleRef{若 17:17}}

如圣保禄也说：\BibleQuote{将和好的话放在了我们心中}（\BibleRef{格后 5:19}）。因为基督来临的目的，也正是宗徒们占据世界的目的。在此，\Quote{如同}同样不是指祂与宗徒们之间的相似，因为人怎能以别的方式被派遣呢？但祂的习惯是将未来之事说成已经发生。\par

% structure: p0013 source=h2 target=subsection reason=relative-heading-rank; base=h2

\subsection{\BibleRef{若 17:19}}

\BibleQuote{我为自己祝圣}是什么意思？就是\Quote{我向祢奉献祭献}。所有祭献都被称为\Quote{圣的}，那些为天主保留的祭品特别称为\Quote{圣物}。因为古时在预象中，祝圣是通过羊羔完成的，但现在不是在预象中，而是藉着真理本身，因此祂说：\BibleQuote{为叫他们因真理而被祝圣。}因为\Quote{我将他们奉献给祢，并将他们作为祭品}；祂这样说，或是因为他们的元首正在成为圣的，或是因为他们自己也被祭献了；因为经上说：\BibleQuote{将你们的身体献上，当作生活、圣洁的祭品}（\BibleRef{罗 12:1}）；又说：\BibleQuote{我们被看作待宰的羊}（\BibleRef{咏 43:22}）。祂使他们不经死亡而成为祭品和奉献；因为当祂说\BibleQuote{我祝圣}时，指的是祂自己的祭献，这从下文可以清楚看出。\par

% structure: p0015 source=h2 target=subsection reason=relative-heading-rank; base=h2

\subsection{\BibleRef{若 17:20}}

2. 既然祂为他们而死，并说\BibleQuote{我为他们祝圣自己}，免得有人认为祂仅是为宗徒们这样做，祂便加上：\BibleQuote{我不但为他们祈求，也为那些因他们的话而信我的人祈求。}祂以此再次振作他们的心灵，表明门徒将会众多。因为祂使那原本专属门徒的恩赐变为共有，祂安慰他们说，他们正成为他人救恩的原因。\par

在如此论及他们的救恩，以及他们藉信德和祭献而被祝圣之后，祂接着论及和睦，并以之结束祂的讲论。祂一开始就提到和睦，也以此结尾。因为在开始祂说：\BibleQuote{我给你们一条新命令}（\BibleRef{若 13:34}）；而此处：\par

% structure: p0018 source=h2 target=subsection reason=relative-heading-rank; base=h2

\subsection{\BibleRef{若 17:21}}

这里的\Quote{如同}同样不是表示他们与祂的完全相似（因为他们不可能达到如此程度），而是指人所能达到的限度。正如祂说：\BibleQuote{你们应当慈悲，如同你们的父那样慈悲}（\BibleRef{路 6:36}）。\par

但\Quote{在我們內}是什麼意思？就是在我們所信的那信德內。因為沒有什麼比分裂更得罪所有的人，所以祂使他們合而為一。有人說：\Quote{那麼，祂成就了這事嗎？}當然，祂成就了。因為所有通過宗徒相信的人都是一體的，雖然有些人從他們中間被撕裂了。這事沒有逃脫祂的知識，祂甚至預言了它，並表明它出自人的懈怠。\par

\Quote{為使世界相信是禰派遣了我。}\par

正如祂起初所說：\Quote{眾人因此就可認出你們是我的門徒，如果你們彼此相愛。}他們如何因此而相信呢？祂說：\Quote{因為禰是平安的天主。}因此，如果他們效法他們所學習的那些人，他們的聽眾就會由門徒認出師傅；但如果他們爭吵，人們就會否認他們是平安天主的門徒，也不會承認我——既然我不是和平的——是禰所派遣的。你看，祂如何直到末了，證明祂與聖父的一致？\par

% structure: p0023 source=h2 target=subsection reason=relative-heading-rank; base=h2

\subsection{若 17:22}

藉著奇蹟，藉著教導，並使他們心意合一；因為這就是光榮：他們合而為一，甚至比奇蹟更大。正如人們欽佩天主，因為在那本性中沒有爭鬥或不和，這是祂最大的光榮，同樣，祂說：\Quote{願他們也因此成為光榮的。}有人說：\Quote{祂怎麼向聖父祈求賜予他們這個，而祂又說祂自己賜予它呢？}不論祂的言論是關於奇蹟、合一或平安，祂都被看見親自將這些賜給了他們；由此顯明，這祈禱是為了他們的安慰。\par

% structure: p0025 source=h2 target=subsection reason=relative-heading-rank; base=h2

\subsection{若 17:23}

\Quote{祂如何賜予光榮？}藉著在他們內，並有聖父與祂同在，從而使他們結合在一起。但在另一處，祂不是這樣說的；祂沒有說聖父藉著祂而來，而是說：\Quote{祂和聖父要來，並要與他一同住下}，此處消除了\Person{撒柏略}的嫌疑，此處消除了\Person{亞流}的嫌疑。\par

\Quote{使他們完全合一，並使世界知道是禰派遣了我。} \BibleRef{若 14:23}\par

祂緊接在另一句話之後說了這些話，以表明平安比奇蹟更有吸引人的力量；因為正如分裂的本性是分離，同樣，合一的性質是結合。\par

\Quote{我愛了他們，如同禰愛了我一樣。}\par

這裡的\Quote{如同}是指，就如人所能被愛的程度；祂的愛的確實證明是祂為他們捨棄了自己。告訴了他們他們將處於安全中，不會被翻倒，他們將成為聖的，有許多人將通過他們相信，他們將享受極大的光榮，不僅祂自己愛他們，聖父也愛了他們之後，祂接著告訴他們，在他們今世旅居之後的事，關於為他們所預備的獎賞和冠冕。\par

% structure: p0031 source=h2 target=subsection reason=relative-heading-rank; base=h2

\subsection{若 17:24}

\Quote{那麼，禰通過祈禱獲得了，而禰還沒有擁有他們不斷詢問的那事麼？他們說：\InnerQuote{禰往哪裡去？}禰說什麼？那麼禰對他們說：\InnerQuote{你們要坐在十二個寶座上}？ \BibleRef{瑪 19:28} 禰怎樣應許了其他更大更多的事？}你看，祂一切都是在自貶中說的？因為祂怎能說：\Quote{你們以後要跟隨}？ \BibleRef{若 13:36} 但祂這樣說，是為了更充分的說服和證明祂的愛。\par

\Quote{使他們看見禰所賜給我的光榮。}\par

這又是祂與聖父同心的一個標记，比前面的更大，因為祂說：\Quote{在世界奠基以前}，卻也有某種自貶；因為祂說：\Quote{禰給了我。}如果事實不是這樣，我很樂意問反駁者一個問題：給予的人，是給予一個已經存在的人；那麼聖父是先生了聖子，然後才賜給祂光榮，之前讓祂沒有光榮？這怎麼合理？你看見\Quote{祂給}就是\Quote{祂生了}嗎？\par

3. 但祂為什麼不說：\Quote{使他們分享我的光榮}，而說：\Quote{使他們看見我的光榮}？這裡祂暗示，那安息的全部就是注視天主聖子。當然，這正是使他們受榮耀的原因；正如\Person{保祿}所說：\Quote{我們以揭開的臉面反映主的光榮。} \BibleRef{格後 3:18} 因為正如那些注視陽光並享受非常清晰空氣的人，從他們的視覺中獲得享受，同樣，那時也將如此，且更大大地，這將給我們帶來喜樂。同時，祂也表明，他們所注視的不是那時所見的身體，而是某種可畏的本體。\par

% structure: p0036 source=h2 target=subsection reason=relative-heading-rank; base=h2

\subsection{若 17:25}

這是什麼意思？有什麼關聯？祂在此表明，除了那些認識聖子的人以外，沒有人認識天主。祂所說的是這一類：\Quote{我希望所有人都是這樣，但他們卻不認識禰，雖然他們對禰沒有怨言。}因為這就是\Quote{公義的聖父}的意義。在我看來，祂說這些話，似乎因他們不願認識如此公義良善的一位而煩惱。因為猶太人曾說他們認識天主，而祂不認識祂，祂針對這一點說：\Quote{因為在世界奠基以前，禰就愛了我。}從而為自己辯護，反駁猶太人的指控。因為祂既得了光榮，在世界奠基以前就蒙愛，又願意他們為那光榮作見證，祂怎能反對聖父呢？\Quote{那麼猶太人所說的不對，他們認識禰，而我不認識禰；相反，我認識禰，而他們不認識禰。}\par

\Quote{這些人都知道是禰派遣了我。}\par

你看，祂提及那些說祂不是出自天主的人，一切總結起來，都是為了反駁這個論點？\par

% structure: p0040 source=h2 target=subsection reason=relative-heading-rank; base=h2

\subsection{若 17:26}

\Quote{然而你說，完美的知識來自聖神。}\Quote{但聖神的事是我的。}\par

\Quote{使禰愛我的愛，存留在他們內，我也在他們內。}\par

\Quote{因为如果他们认识到祢是谁，就会知道我与祢并不分离，而是极为蒙爱的一位，是真正的圣子，与祢紧密相连。那些正确信服此事的人，将持守对我的信德和完美的爱德；当他们按所应该的去爱时，我就住在他们内。}你们看见祂如何抵达了美好的终点，以爱德——诸福之母——结束了这篇讲论吗？\par

4. 让我们相信天主并爱祂，免得有人这样评论我们：\Quote{他们自称认识天主，但在行为上却否认祂。} \BibleRef{铎 1:16} 又说：\Quote{他背弃了信德，比不信的人更坏。} \BibleRef{弟前 5:8} 因为当那人帮助他的家人、亲属和外人时，而你却连与你同族的亲戚都不扶持，那么今后你有什么借口呢？那时天主因你而受到亵渎和侮辱。想想天主赐给了我们多少行善的机会。祂说：\Quote{你要怜悯这个人，}祂说，\Quote{如同对待亲属；另一个人，如同对待朋友；另一个人，如同对待邻人；另一个人，如同对待公民；另一个人，如同对待人。} 如果这些都不能约束你，反而你打破了所有纽带，那么就听听保禄的话，你\Quote{比不信的人更坏}；因为他从未听说过施舍或天上的事，却在爱人方面超过了你；而你，被命令要爱你的仇敌，却把朋友看作仇敌，更加关心你的钱财胜过他们的身体。然而钱财花费了不会受损，但你的兄弟若被忽略却会丧亡。那么，关心钱财而不关心亲属，是何等的疯狂？这贪求财富的欲望是怎样冲进我们中间的？这不人道和残忍又是从何而来？因为如果有人能像坐在剧院最高处一样，俯瞰整个世界——或者更确切地说，如果你愿意，我们暂且拿一座城市举例——如果一个人坐在高处，能一眼看到那里所有人的所作所为，想想他会谴责何等愚蠢，会流下何等眼泪，会发出何等笑声，会怀有何等憎恨；因为我们所做的这类行为，既该受嘲笑和愚蠢的指控，也当受眼泪和憎恨。一个人养狗捕捉野兽，自己却沉沦于残暴；另一个人养牛和驴运石头，却忽略正在饥饿消瘦的人；并且毫无节制地花费黄金去制造石人，却忽略真实的人，这些人因邪恶状态变得像石头一样。另一个人费尽心力收集金矿，将其围绕在墙上，但当他看见穷人赤裸的肚腹时，无动于衷。有些人又在衣服上加添衣服，而他们的兄弟甚至没有东西遮盖赤裸的身体。又有人在法庭上吞并了另一个人；另一个人将钱花在女人和寄生虫上；另一个人花在戏子和戏班子上；另一个人花在辉煌的建筑物、购买田地和房屋上。又有人计算利息，另一个人计算复利；另一个人组成充满死亡的团伙，甚至在夜间也不得休息，躺着不眠谋害他人。然后天亮时，他们奔跑：一个去谋取不义之财，另一个去恣意挥霍，其他人去公开抢劫。对于多余和禁止之事物，他们极其认真，但对于必要之事却毫不理会；那些裁决法律问题的人，虽然名义上是陪审员，实际上是盗贼和凶手。如果有人查问诉讼和遗嘱，他会在那里再次发现无数的祸害、欺诈、抢劫、阴谋；所有时间都花费在这类事上；但对属灵的事却毫不关心，他们都为了仅仅给人看而妨碍了教会。但这并非所要求的；我们需要行为，和纯洁的心灵。但如果你终日忙于攫取财富，然后进来念几句经文，你不仅没有平息天主的怒气，反而更加激怒了他。你想与你的主和解吗？要显出行为，使自己熟知众多的苦难，看顾赤裸的人、饥饿的人、受冤屈的人；祂已为你开辟了千万条显示对人的爱的途径。那么我们不要漫无目的地生活而自欺欺人，也不要因我们现在健康而自以为是；但要记住，我们常常在患病时，达到极度虚弱，因恐惧和展望未来之事而濒死，让我们期待再次陷入同样状态，再次拥有同样的恐惧，从而成为更好的人；因为现在的所作所为该受无限的谴责。因为在法庭上的人如同狮子和狗；在公共场所的人如同狐狸；而那些过着悠闲生活的人，甚至没有按应有的方式使用闲暇，将所有时间花费在戏院和由此产生的祸患上。没有人谴责正在发生的事；但有许多人嫉妒，并且因自己不在同样境况而烦恼，以致这些人虽未实际作恶，却反过来受罚。因为他们\Quote{不仅做这些事，而且还喜欢那些做这些事的人。}因为属于他们意愿的同样败坏；由此显然，意图也将受罚。这些事我每天都说，并且不会停止说。因为如果有人听，就是获益；但如果无人留心，你们将来就会听到这些事，那时对你们毫无益处，你们会责怪自己，而我们则可免于过失。但但愿我们永远不会只有这个借口，而是你们在基督的审判台前成为我们的夸耀，使我们能一同享受福分，借我们的主耶稣基督的恩宠和慈爱，愿光荣偕同祂归于圣父及圣神，至于世世。阿们。\par

\SourceRecord{Translated by Charles Marriott.；Nicene and Post-Nicene Fathers, First Series；Vol. 14.；Edited by Philip Schaff.；Buffalo, NY: Christian Literature Publishing Co.,；1889.；<http://www.newadvent.org/fathers/240182.htm>.}

% structure: p0001 source=h1 target=section reason=page-title

\section{\WorkTitle{若望福音讲道集第八十三篇}}

% structure: p0002 source=h2 target=subsection reason=relative-heading-rank; base=h2

\subsection{若望福音 18:1}

死亡是一件可怕的事，充满恐怖，但对于那些学习了天上真正智慧的人并非如此。因为对来世一无所知，只以为死亡是生命的某种消散和终结的人，自然会战栗害怕，如同归于虚无。但我们靠着天主的恩宠，已经学到了祂智慧中隐藏和秘密的事，认为死亡是迁往另一个地方，就不应当战栗，反而应当喜乐欢欣，因为离开这必朽的生命，我们将前往一个更美好、更光明且没有终结的生命。基督以祂的行动教导了这一点，祂前往受难，不是出于强迫和必须，而是出于自愿。经上说：\BibleQuote{耶稣说了这些话，就同门徒出去，到了\Place{克德龙溪}的对岸。那里有一个园子，祂和门徒进去了。}\par

% structure: p0004 source=h2 target=subsection reason=relative-heading-rank; base=h2

\subsection{若望福音 18:2}

祂在半夜出发，渡过一条河，急忙来到叛徒所知的地方，减轻了那些图谋对付祂的人的辛劳，免去他们一切麻烦；并向门徒表明，祂是自愿去受难（这件事最足以安慰他们），并把自己置身于园中，如同在监牢里。\par

\BibleQuote{耶稣对他们说了这些话。} \Quote{你说什么呢？祂一定是在与圣父交谈，一定是在祈祷。那么你为什么不说：\InnerQuote{祂结束了祈祷之后，就去了那里}呢？} 因为那不是祈祷，而是为了门徒的缘故所说的一番话。\BibleQuote{门徒就进了园子。}祂已如此消除了他们的恐惧，使他们不再抵抗，而与祂一同进了园子。但是犹达斯怎么来到那里？他来的时候从哪里得到消息？由此可知，耶稣通常是在户外过夜。因为如果祂惯常在家过夜，犹达斯就不会到旷野来，而会到家里去，以为会在那里找到祂睡觉。免得你们听到\Quote{园子}就以为耶稣隐藏了自己，经上补充说，\BibleQuote{犹达斯知道那地方}；不仅如此，还记载祂\BibleQuote{屡次同门徒在那里聚集}。因为祂常与门徒单独在一起，谈论必要的事，以及那些不许别人听的事。祂这样做，尤其是在山上和园子里，找一个不受干扰的地方，好让他们的注意力不被听道所分散。\par

% structure: p0007 source=h2 target=subsection reason=relative-heading-rank; base=h2

\subsection{若望福音 18:3}

这些人以前也多次派人去捉拿祂，但未能成功；由此可见，此时祂是自愿交出自已的。他们又是怎样说服那队兵的呢？他们是兵士，素来习惯为钱做事。\par

% structure: p0009 source=h2 target=subsection reason=relative-heading-rank; base=h2

\subsection{若望福音 18:4}

这就是说，祂不是等到他们来了才得知这事，而是毫不慌乱地说话行事，因为祂知道这一切。\Quote{但他们为什么带着武器来捉拿祂呢？}他们害怕祂的跟随者，因此深夜前来。\BibleQuote{祂就上前对他们说：你们找谁？}\par

% structure: p0011 source=h2 target=subsection reason=relative-heading-rank; base=h2

\subsection{若望福音 18:5}

你看到祂不可战胜的能力了吗？祂站在他们中间，却使他们的眼睛昏花。因为黑暗不是他们不认识祂的原因，圣史已表明，他们还拿着火把。即使没有火把，他们至少也该从声音认出祂；或者，如果他们不认识祂的声音，那一直与祂在一起的犹达斯又怎能不认识呢？因为他也和他们站在一起，却和他们一样不认识祂，反而和他们一起跌倒在地。耶稣这样做，是为了表明，除非祂允许，他们不仅不能捉拿祂，甚至连站在中间的祂也看不见。\par

% structure: p0013 source=h2 target=subsection reason=relative-heading-rank; base=h2

\subsection{若望福音 18:7}

% structure: p0014 source=h2 target=subsection reason=relative-heading-rank; base=h2

\subsection{若望福音 18:8}

请看圣史的宽容，他没有辱骂叛徒，而是叙述所发生的事，只想证明一件事：这一切都是出于祂自己的同意。然后，免得有人说，是祂自己把他们引到这一步，因为祂将自己交到他们手中，并向他们显明了自己；在向他们显示了足以使他们退却的一切之后，当他们仍执迷不悟、毫无借口时，祂就将自己交在他们手中，说：\par

\BibleQuote{所以，你们若找我，就让这些人走吧。}\par

\Quote{如果，}祂说，\Quote{你们要我，就不要管这些人，因为，看，我把自己交出来。}\par

% structure: p0018 source=h2 target=subsection reason=relative-heading-rank; base=h2

\subsection{若望福音 18:9}

\Quote{祂在此所说的\Quote{丧失}并非指死亡的丧失，而是指永远的丧失；不过圣史在目前的情况下也包括了前者。也许有人会奇怪，他们为什么不连同祂一起捉拿门徒，并把他们剁成肉酱，尤其是伯多禄对仆人所做的事激怒了他们。那么是谁阻止了他们呢？不是别的，正是那使他们跌倒的能力。因此，圣史为了表明这不是出于他们的意图，而是出于那被捉拿者的能力和命令，便补充说：\Quote{这应验了祂所说的话，}即\Quote{没有丧失一个等。}（\BibleRef{若 17:12}）}\par

2. 因此，伯多禄从祂的声音和已发生的事中鼓起勇气，武装自己以对付攻击者。有人说：\Quote{祂既吩咐不可带钱囊、不可有两件衣裳，怎么又有一把剑？}我想他早已预备了这剑，正是惧怕所发生的事。但若你说：\Quote{那连用手击打都不被允许的人，如何成了杀人者？}他确实曾被命令不可自卫，但此处他并非自卫，而是为他的主。况且，他们那时尚未成全或完备。但若你渴望看到伯多禄被赋予天上的智慧，此后你将见他受伤却温顺地忍受，遭受无数可怕之事而不动怒。然而耶稣在此又行了一个奇迹，既表明我们应当善待那些对我们作恶的人，也彰显祂自己的权能。因此祂恢复了仆人的耳朵，并对伯多禄说：\Quote{凡持剑的，必死于剑下}\BibleRef{玛 26:52}；正如在洗脚盆的事上，祂以威胁缓和了伯多禄的激烈，此处亦然。福音作者记下那仆人的名字，是因为所行之事非常伟大，不仅因为祂医治了他，更因祂医治了一个前来攻击祂、不久后还要掌掴祂的人，并且祂平息了因这事而可能对门徒燃起的战争。为此，福音作者记下名字，好让当时的人能仔细查问这些事是否真的发生。他并非无故提到\Quote{右耳}，而是我认为他渴望显示那位宗徒的冲动，几乎直攻头部。然而耶稣不仅以威胁制止他，还用其他话语安抚他，说：\par

% structure: p0021 source=h2 target=subsection reason=relative-heading-rank; base=h2

\subsection{\BibleRef{若 18:11}}

表明所发生的事并非出于他们的能力，而是出于祂的同意，并宣称祂并非反对天主，而是顺服圣父以至于死。\par

% structure: p0023 source=h2 target=subsection reason=relative-heading-rank; base=h2

\subsection{\BibleRef{若 18:12-13}}

为何解到亚纳斯那里？他们在喜悦中炫耀所行之事，仿佛立了一座凯旋碑。\par

\Quote{他是盖法的岳父。}\par

% structure: p0026 source=h2 target=subsection reason=relative-heading-rank; base=h2

\subsection{\BibleRef{若 18:14}}

为何福音作者再次提醒我们他的预言？为表明这些事是为了我们的救恩而发生的。真理的力量如此强大，连仇敌也事先宣告了这些事。免得听者听到捆绑的事就惊慌，他提醒那预言，即耶稣的死是世界的救恩。\par

% structure: p0028 source=h2 target=subsection reason=relative-heading-rank; base=h2

\subsection{\BibleRef{若 18:15}}

那另一个门徒是谁？是作者本人。\Quote{他为何不提名自己？当他躺在耶稣怀里时，他有理由隐瞒名字；但现在为何如此？}出于同样理由，因为此处他也提到一件大善行，即当所有人都逃散时，他却跟随了。因此他隐藏自己，并将伯多禄置于他之前。他不得不提到自己，好让你明白他比其余人更精确地叙述了厅内发生的事，因为他自己曾在里面。但请注意他如何贬低自己的称赞；免得有人问：\Quote{众人退却时，此人如何比西满更深入？}他说他\Quote{是大司祭认识的}。这样就没有人会惊讶他跟随了，也不会称赞他的勇气。但令人惊奇的是伯多禄的事，他既如此害怕，却来到厅前，而其余人已退却。他的来到出于爱，他的未进入则出于痛苦和恐惧。因为福音作者记下这些事，是为给否认的事开辟辩解之路；论到自己，他不将认识大司祭当作什么大事，但因他说唯独他与耶稣一同进去，免得你以为这行动出于崇高情感，他也说明了原因。并且，若伯多禄被允许，他也会进去，这由下文显明：当他出去并吩咐看门的使女领伯多禄进来时，伯多禄立即进来了。但他为何不亲自领他进来？他依恋基督并跟随祂；为此他吩咐那女人领他进来。那么那女人说什么？\par

% structure: p0030 source=h2 target=subsection reason=relative-heading-rank; base=h2

\subsection{\BibleRef{若 18:17}}

你说什么，伯多禄？你不是刚刚宣告：\Quote{我情愿为祢舍命，我必舍命}？如今怎么了，你竟连一个看门女子的盘问都受不了？是士兵盘问你吗？是那些捉拿祂的人之一吗？不，是一个卑微低贱的看门女子，而且盘问并非粗暴。她没说：\Quote{你是那骗子和败坏者的门徒吗？}而是\Quote{那人的}，这更像是怜惜和心软的措辞。但伯多禄连这些话都受不了。\Quote{你不也是}这样说，是因为若望在里面。那女子说得如此温和。但他丝毫没有察觉，也不放在心上，无论是第一次、第二次还是第三次，直到鸡叫；甚至这也没有使他醒悟，直到耶稣以苦涩的目光看他。他站着与大司祭的仆人们一起取暖，而基督却被捆绑在里面。我们说这些并非指责伯多禄，而是表明基督所说之事的真实性。\par

% structure: p0032 source=h2 target=subsection reason=relative-heading-rank; base=h2

\subsection{\BibleRef{若 18:19}}

3. 何等的邪恶！虽然他一直听见祂在圣殿里公开讲论和教训，现在却想得到指示。因为他们没有可控告的，便询问有关祂门徒的事，或许他们在哪里，为何召集他们，以及出于什么意图和条件。他说这话，是想证明祂是煽动者和革新者，因为除了他们以外无人听从祂，仿佛祂的团体是某种邪恶工场。那么基督说什么？为推翻这指控，祂说：\par

% structure: p0034 source=h2 target=subsection reason=relative-heading-rank; base=h2

\subsection{\BibleRef{若 18:20}}

\Quote{那么，祂没有在暗中说过什么吗？}祂说了，但并非如他们所想，出于恐惧或制造阴谋，而是因为有时祂的话对众人来说过于高深。\par

% structure: p0036 source=h2 target=subsection reason=relative-heading-rank; base=h2

\subsection{\BibleRef{若 18:21}}

这些话不是出于傲慢之言，而是出于对祂所说真理的自信。因此，祂在起初所说的：\BibleQuote{我若为自己作证，我的证据不足凭信}\BibleRef{若 5:31}，现在祂暗示了这一点，渴望使祂的见证更加可信。因为当亚纳斯提及门徒时，祂说什么？\Quote{你问我有关我的事吗？问我的敌人，问那些图谋陷害我、捆绑我的人，让他们说吧。}这是一个无可置疑的真理证明，即一个人召来他的敌人为他所说的话作证。那么，大司祭做了什么？本应如此进行查问的事，那人却没有这样做。\par

% structure: p0038 source=h2 target=subsection reason=relative-heading-rank; base=h2

\subsection{若望福音 18:22}

还有什么比这更狂妄呢？诸天啊，颤抖吧！大地啊，惊愕吧！面对主的容忍和仆人的愚昧！然而，祂说了什么？祂并没有说\Quote{你为什么问我？}好像拒绝说话，而是希望消除一切无理行为的借口；并且在此情况下被打耳光，虽然祂能够震动、毁灭或移走万物，但祂没有做任何这些事，而是说出足以缓和任何残暴的话。\par

% structure: p0040 source=h2 target=subsection reason=relative-heading-rank; base=h2

\subsection{若望福音 18:23}

这就是说：\Quote{你若能抓住我的话柄，就指出来；若不能，为什么打我？}你看见那审判厅充满了喧闹、麻烦、激情和混乱吗？大司祭欺诈而阴险地提问，基督直截了当地回答，恰如其分。那么接下来该做什么？要么反驳，要么接受祂的话。但这并没有发生，而是一个仆人打了祂。这根本不是法庭，而是阴谋和暴力的行为。然而即便如此，他们也没有进一步发现什么，就把祂捆绑着送到盖法那里。\par

% structure: p0042 source=h2 target=subsection reason=relative-heading-rank; base=h2

\subsection{若望福音 18:25}

奇妙啊，当耶稣被带走时，那个激愤如火的人竟被何等昏睡所控制！经过这些事情之后，他并没有动，仍然在取暖，好叫你明白，如果天主舍弃，我们的本性是多么软弱。并且，被问到时，他又否认了。\par

% structure: p0044 source=h2 target=subsection reason=relative-heading-rank; base=h2

\subsection{若望福音 18:26}

但是，那花园没有使他记起所发生的事，耶稣在那里凭那些话所彰显的极大关爱也没有，但所有这些都因焦虑的压力而被他抛之脑后。然而，为什么圣史们一致地写到关于他的事？不是要控告那门徒，而是渴望教导我们，不完全交托于天主、却信赖自己是何等大的恶。但你要钦佩祂主人对祂的关爱，祂虽被囚捆绑，却为祂的门徒伯多禄极大操心，用祂的一瞥扶起跌倒的他，并把他投入泪海之中。\par

\BibleQuote{他们于是从盖法那里把耶稣带到总督衙门。}\par

这样做，是为了使审判祂的人数，甚至违反他们自己的意愿，显示祂的真理经受了充分的考验。\BibleQuote{那时，天还早。}鸡叫之前，祂被带到盖法那里；清早，被带到比拉多那里。因此，圣史表明，被盖法审问了大半夜，祂没有在任何事上被证明有罪；所以盖法把他解送到比拉多那里。但是，把这些事留给其他圣史去叙述，若望接着讲述下面的事。请看犹太人荒谬的行为。他们抓住了无辜者，并拿起了武器，却不进入总督府，\BibleQuote{以免受玷污}。告诉我，踏进一个审判厅，那里恶人受到惩罚，这算什么玷污？他们连薄荷、茴香都献上十分之一，却不认为自己因图谋不义杀人而受玷污，却认为因踏进法庭而玷污了自己。\Quote{他们为什么不杀死祂，却把祂带到比拉多那里？}首先，他们统治和权威的大部分已被剥夺，当他们的政务置于罗马人权力之下时；此外，他们害怕日后被祂控告和惩罚。\Quote{但是，\InnerQuote{好能吃逾越节晚餐}是什么意思？}他要么称整个节日为逾越节，要么意思是，他们那时正在过逾越节，而祂提前一天把逾越节传给祂的门徒，将自己的祭献留到预备日，就是古时逾越节庆祝的日子。但他们虽然拿起了武器——这是违法的——并且流了血，却对地方顾忌，把比拉多带到他们面前。\par

% structure: p0048 source=h2 target=subsection reason=relative-heading-rank; base=h2

\subsection{若望福音 18:29}

4. 你看见他没有傲慢和恶意吗？因为看见耶稣被绑着、被这么多人带来，他并不认为他们有确凿的证据来指控，却质问他们，认为奇怪的是他们为自己拿走了审判，然后把惩罚没有任何审判地交给他。那么，他们说什么？\par

% structure: p0050 source=h2 target=subsection reason=relative-heading-rank; base=h2

\subsection{若望福音 18:30}

啊，疯狂！为什么你们不提及祂的恶行，反而隐瞒它们？为什么你们不证明邪恶？你看见他们处处回避直接控告，无话可说吗？亚纳斯问祂关于祂的道理，听完后，把祂送到盖法那里；盖法又问了祂，发现什么都没有，就把祂送到比拉多那里。比拉多说：\BibleQuote{你们对这个人提出什么控告？}这里他们也无话可说，又使用某些猜测。比拉多为此困惑，说：\par

% structure: p0052 source=h2 target=subsection reason=relative-heading-rank; base=h2

\subsection{若望福音 18:31-32}

\Quote{那么，\InnerQuote{我们不可把人处死}这句话如何表明这事？}要么圣史的意思是祂不仅要被犹太人杀死，也要被外邦人杀死，要么是说他们无权钉十字架。但是，如果他们说\BibleQuote{我们不可把人处死}，那是针对那个季节说的。因为他们的确杀了人，并且以不同的方式杀了人，斯德望被石头砸死就是明证。但他们渴望钉祂十字架，以便展示祂死的方式。比拉多希望摆脱麻烦，没有将祂释放以进行长期审判，而是，\par

% structure: p0054 source=h2 target=subsection reason=relative-heading-rank; base=h2

\subsection{若望福音 18:33-34}

基督为何问这事？因为祂要揭露犹太人的恶意。比拉多从多人听到这说法，且控告者无话可说，为免审讯过长，他想提出那不断被传的事。但当他对他们说：\Quote{你们按你们的法律审判祂}，欲表明祂的过犯并非犹太式的，他们回答说：\Quote{我们是不合法的。}\Quote{祂没有违反我们的法律，而是普遍的控告。}比拉多察觉此事后，说（身为可能受害的人）：\Quote{你是犹太人的王吗？}然后耶稣，并非出于无知，而是希望犹太人甚至被他控告，问他说：\Quote{是别人告诉你的吗？}关于此点，比拉多表明自己，回答说：\par

% structure: p0056 source=h2 target=subsection reason=relative-heading-rank; base=h2

\subsection{\BibleRef{若 18:35}}

在此他欲为自己开脱。然后因为他曾说：\Quote{你是王吗？}耶稣责备他，回答：\Quote{这话你是从犹太人听来的。你为何不仔细查问？他们说我是一个恶人；问他们我做了什么恶事。但你不这样做，只是简单地构陷我。}\Quote{耶稣回答他说：这话是你自己说的，}还是从别人那里听来的？比拉多不能立即说他听说了，而是随从民众，说：\Quote{他们把你交给了我。}\Quote{所以我必须问你你做了什么。}基督说了什么呢？\par

% structure: p0058 source=h2 target=subsection reason=relative-heading-rank; base=h2

\subsection{\BibleRef{若 18:36}}

祂引导比拉多向上，这比拉多并非一个十分邪恶的人，也不像他们那样，祂欲表明祂不只是一个凡人，而是天主和天主子。祂说了什么呢？\par

\Quote{我的国若属于这世界，我的臣仆必要作战，使我不致被交给犹太人。}\par

祂消除了比拉多一度害怕的那件事，即夺取王权的嫌疑。\Quote{那么祂的国不也属于这世界吗？}当然属于。\Quote{那祂为何说\InnerQuote{不属于}？}不是因为祂不在这里治理，而是因为祂的权柄来自天上，且不是人的，远比这更伟大、更辉煌。\Quote{既然更伟大，祂怎会被那另一个所俘？}因祂同意并交付自己。但祂此时并未明示这点，而是说：\Quote{我若属于这世界，我的臣仆必要作战，使我不致被交出。}在此祂显示了人间王权的软弱，即其力量在于臣仆；而属天的权柄自足，无需他物。异端者由此抓住机会说，祂不同于造物主。那么，当经上说\Quote{祂来到自己的地方}时（\BibleRef{若 1:11}），当祂自己说\Quote{他们不属于这世界，如同我不属于这世界}时（\BibleRef{若 17:14}），又意味着什么？同样，祂说祂的国不来自此处，并非剥夺世界对祂的眷顾和护理，而是如我所说，表明祂的权柄不是人的或可朽的。比拉多于是说了什么？\par

% structure: p0062 source=h2 target=subsection reason=relative-heading-rank; base=h2

\subsection{\BibleRef{若 18:37}}

那么，如果祂生来就是王，祂的一切其他属性都出于\OriginalTerm{受生}{Generation}，祂没有领受任何附加的东西。因此，当你听到：\Quote{圣父怎样在自己有生命，也照样赐给圣子在自己有生命}（\BibleRef{若 5:26}），不要想别的，只想到祂的受生，其余亦然。\par

\Quote{我为此而来，为给真理作证。}\par

\Quote{就是说，\InnerQuote{我要说这件事本身，并教导它，劝服众人。}}\par

5. 但是，人啊，当你听到这些事，看到你的主被捆绑、被牵引时，你要视当前之事为虚无。因为如果基督为你承受了如此的事，而你却常常连言语都无法忍受，这岂不是奇怪吗？祂受唾污，而你却用衣饰和戒指装扮自己；若你得不到众人的好评，就认为人生不堪忍受吗？祂受侮辱，忍受戏弄和打在脸颊上的讥讽耳光；而你却希望处处受人尊敬，并且不忍受基督的辱骂吗？你难道没有听见保禄说：\BibleQuote{你们该效法我，如同我效法基督一样}？\BibleRef{格前 11:1} 因此，当有人戏弄你时，要记住你的主，他们曾在戏弄中向祂屈膝，用言语和行为折磨祂，用许多讥讽对待祂；但祂不但不为自己辩护，反而以相反的方式回报他们，即温良和柔和。如今让我们效法祂；这样，我们就能甚至摆脱一切侮辱。因为不是侮辱者使侮辱行为生效并使之刺人，而是那心胸狭隘、为之痛苦的人。如果你不感到痛苦，你就没有被侮辱；因为受伤害的痛苦不在于施加者，而在于承受者。你为何要忧伤呢？如果有人无端侮辱你，在这种情况下你当然不该忧伤，反而要可怜他；若是有原因的，更该缄默。因为如果有人称呼你，一个穷人，好像你是富人，他话中的赞美对你毫无意义，他的赞美反倒是嘲弄；同样，如果有人侮辱你，说了不真实的话，那指责对你也是毫无意义。但如果良心抓住了所说的话，不要因言语忧伤，而要用行动改正。我这是对真正的侮辱而言。因为如果有人以贫穷或出身低贱责备你，就嘲笑他。这些事不是对听者的责备，而是对说者的责备，因为他不知道真正的智慧。有人说：\Quote{但是，当这些话在许多不明真相的人面前说出时，伤口就变得无法忍受。}不，这是最可忍受的，当你在场有听众作见证，他们赞美和鼓掌欢迎你，而嘲笑和戏弄他。因为不是为自己辩护的人，而是什么也不说的人，得到明智者的称赞。如果在场的人中没有明智者，那么更要嘲笑他，并以上天的听众为乐。因为在那边，众人都会赞美、鼓掌、欢迎你。因为一位天使抵得上整个世界。但我为何提及天使呢？当主亲自称赞你时？让我们用这些推理来操练自己。因为被侮辱时沉默并非损失，相反，被侮辱时为自己辩护才是损失。因为如果默默忍受所说的话是过错，基督就绝不会告诉我们：\BibleQuote{有人打你的右颊，把另一面也转给他。}\BibleRef{玛 5:39} 因此，如果我们的仇敌说了不真实的话，我们更应为此可怜他，因为他招致告发者的惩罚和报复，甚至不配阅读圣经。因为天主对罪人说：\BibleQuote{你为什么传述我的规诫，口里提及我的盟约？你坐著毁谤你的兄弟。}\BibleRef{咏 49:16} 如果他说的是真话，也同样值得可怜；因为连法利塞人也说了真话；但他对听者没有造成伤害，反而有益，而他自己却剥夺了无数的祝福，因这控告而遭受沉船之苦。所以无论哪种情况，都是他受害，而不是你；但你若清醒，却会获得双重的益处：既因你的沉默而悦乐天主，又变得更加谨慎，从所说的话中获得了改正所行之事的机遇，并轻视人间的光荣。因为这是我们痛苦的根源：许多人渴求人的意见。如果我们有意这样真正明智，我们就会清楚知道人的事物算不了什么。那么，让我们学习，在数算自己的过失后，及时加以改正，并决心这个月改正一个，下个月改正另一个，再下个月改正第三个。如此，如同踏着阶梯，让我们藉着雅各伯的梯子登上天堂。因为那梯子在我看来以那异象隐晦地表示借由德行逐步上升，使我们能从地上登上天堂，不是用物质阶梯，而是品德的改善和修正。因此，让我们把握这离别和上升的方法，好使我们获得天堂后，也能享受那里的一切福祉，借着我们的主耶稣基督的恩宠和慈爱；愿光荣归于祂，至于无穷之世。阿们。\par

\SourceRecord{Translated by Charles Marriott.；Nicene and Post-Nicene Fathers, First Series；Vol. 14.；Edited by Philip Schaff.；Buffalo, NY: Christian Literature Publishing Co.,；1889.；<http://www.newadvent.org/fathers/240183.htm>.}

% structure: p0001 source=h1 target=section reason=page-title

\section{若望福音讲道集第八十四篇}

% structure: p0002 source=h2 target=subsection reason=relative-heading-rank; base=h2

\subsection{若 18:37}

1. 忍耐是一件奇妙的事，它使灵魂如同安处于宁静的港湾，免于颠簸和恶灵的侵扰。这一点，基督处处教导了我们，尤其在祂受审判、被拖曳、被带领的时候。因为当祂被带到\Person{亚纳斯}面前时，祂以极大的温良回答，并对那打祂的仆人说了一句足以挫败其全部傲慢的话；从此处去到\Person{盖法}那里，然后到\Person{比拉多}那里，整夜都在这些场景中度过，祂始终表现出自己的温和。当他们说祂是作恶者却又无法证明时，祂沉默不语；但当祂被问及关于国的事时，祂便对\Person{比拉多}说话，教导他，引导他进入更高的事理。但为什么\Person{比拉多}不在他们面前，而是私下进入审判厅询问呢？他心中对祂有些重大的猜测，不愿受犹太人干扰，想准确地了解一切。然后，当他说\Quote{你做了什么？}时，耶稣对此没有回答；但关于\Person{比拉多}最想听的事，即祂的国，祂回答说：\Quote{我的国不属于这世界。}也就是说：\Quote{我确实是一个王，但不是你所怀疑的那种，而是更为荣耀的。}藉着这些话以及随后的话，表明祂没有做过任何恶事。因为说\Quote{我为此而生，也为此来到世界，为给真理作证}的人，表明祂没有做过任何恶事。然后当祂说\Quote{凡属于真理的，必听我的声音}时，祂藉此吸引他，并说服他成为话语的聆听者。\Quote{因为，}祂说，\Quote{如果有人是真实的，并且渴望这些事，他必会听我。}事实上，祂藉着这些简短的话如此抓住了他，以至于他说：\par

% structure: p0004 source=h2 target=subsection reason=relative-heading-rank; base=h2

\subsection{若 18:38}

但此刻他处理的是紧急的事，因为他知道这个问题需要时间，并且他想要从犹太人的暴力中解救祂。因此他出去，说了什么？\par

\Quote{我在这人身上查不出什么罪过。}\par

请看他是多么审慎行事。他没有说：\Quote{既然他犯了罪，该死，就因庆节宽恕他吧}，而是首先宣告祂完全无罪，然后进一步请求他们，即使不愿意释放无罪之人，至少因时节宽恕有罪之人。因此他补充说：\par

% structure: p0008 source=h2 target=subsection reason=relative-heading-rank; base=h2

\subsection{若 18:39-40}

哦，可诅咒的决定！他们要求释放与自己同类的人，却让有罪者离开；而命令他惩罚那无辜者。因为这是他们自古以来的习惯。但你要在此一切中观察主在那些处境中的慈爱。\Person{比拉多}鞭打祂，或许是为了消磨并平息犹太人的狂怒。因为当他之前无法通过先前的措施解救祂时，他急于在此刻止住恶事，便鞭打了祂，并允许了所做的事，即给祂穿上袍子和戴上冠冕，以缓和他们的怒气。因此他也把祂带出来给他们看，祂戴着荆棘冠冕\BibleRef{若 19:5}，希望看到祂所受的侮辱后，他们能稍微从激动中平静，吐出毒液。\Quote{兵士们怎么可能这样做，如果不是他们长官的命令？}为了取悦犹太人。因为他们起初夜间进入并非出于他的命令，而是为了取悦犹太人；他们为了钱什么都敢做。但祂，当如此多且如此的事发生时，仍然沉默，正如在审讯中所做的那样，什么也不回答。并且你不要仅仅听这些事，而要常常记在心里，当你看见世界和众天使的君王被兵士嘲笑，被言语和行为戏弄，并默默忍受一切时，你要自己用行为效法祂。因为当\Person{比拉多}称祂为犹太人的王，而他们给祂披上嘲弄的服饰时，\Person{比拉多}便领祂出来，说：\par

% structure: p0010 source=h2 target=subsection reason=relative-heading-rank; base=h2

\subsection{若 19:4-5}

但即使如此，他们的怒火也未熄灭，反而喊叫说：\par

% structure: p0012 source=h2 target=subsection reason=relative-heading-rank; base=h2

\subsection{若 19:6}

那时\Person{比拉多}见一切都是徒劳，就说：\par

\Quote{你们自己把他带去，钉在十字架上罢。}\par

由此可见，他先前允许所做的事，是因为他们的疯狂。\par

\Quote{因为我查不出他有什么罪过。}他说。\par

2. 请看法官以多少方式为祂辩护，不断地宣告祂无罪；但这些事没有一件使那些狗羞愧而放弃他们的意图。因为\Quote{你们自己把他带去，钉在十字架上罢}是推卸责任之言，并将他们推向前去行一件不被允许的事。因此他们带来祂，是为了使事情通过总督的判决而完成；但相反的事发生了，祂反而被总督的判决宣告无罪，而不是定罪。然后，因为他们感到羞愧，\par

% structure: p0018 source=h2 target=subsection reason=relative-heading-rank; base=h2

\subsection{若 19:7}

既然如此，当审判官说\InnerQuote{你们自己把他带去，按照你们的法律审判他}时，你们回答说\InnerQuote{我们不可杀人}，而在这里你们却求助于法律？请考虑这指控：\InnerQuote{祂使自己成为天主子。}请告诉我，这算是一种控告吗？那行了天主子所行的事，却称自己为天主子，这算罪过吗？那时基督做什么呢？当他们彼此这样对话时，祂保持了沉默，应验了先知的话：\BibleQuote{祂不开口；在祂受屈辱时，祂的审判被夺去了。}\BibleRef{依 53:7}\par

然后比拉多从他们听说祂使自己成为天主子，就惊慌，害怕这断言可能是真的，而他似乎要逾越；但这些从祂的言行得知这事的人，并不战栗，反而为了他们本该朝拜祂的种种理由而将祂处死。为此他不再问祂\Quote{你做了什么？}而是被恐惧摇撼，重新开始查问说\Quote{你是基督吗？}但祂没有回答。因为那个曾听见\Quote{我为此而生，也为此而来}以及\Quote{我的国不属于这世界}的人，在他应当反对祂的仇敌并释放祂的时候，却没有这样做，反而助长了犹太人的狂怒。于是他们既在各方面被驳倒，便将他们的呼喊变成一项政治指控说\Quote{凡自称为王的，就是反对凯撒。}\BibleRef{若 19:12}比拉多本应仔细查证祂是否曾图谋王权并着手将凯撒逐出王国。但他没有精确查问，因此基督没有回答他什么，因为祂知道他所问的一切都是徒劳。此外，既然祂的作为为祂作证，祂不愿凭言语取胜，也不作任何辩护，表明祂是自愿来到这状态。当祂沉默时，比拉多说：\par

% structure: p0021 source=h2 target=subsection reason=relative-heading-rank; base=h2

\subsection{若 19:10}

你们看他如何预先定了自己的罪；因为\Quote{如果一切都在你手中，你既然查不出祂有什么罪，为什么不释放祂？}那时比拉多已说出对自己不利的判词，然后祂说：\par

% structure: p0023 source=h2 target=subsection reason=relative-heading-rank; base=h2

\subsection{若 19:11}

显示他也有罪。然后，为了挫败他的骄傲和狂妄，祂说：\par

\Quote{除非由上赐给你，你绝无权柄。}\par

显示这事的发生并非仅仅按平常的事件次序，而是奥秘性地成就。然后，免得当你听见\Quote{除非由上赐给你}时，你会认为比拉多免于一切责难，因此祂说：\Quote{所以那把我交付给你的，罪过更大。}\Quote{然而，如果是赐给你的，他和他们都不应受任何指控。}\Quote{你徒然反对；因为这里的\InnerQuote{given}意思是什么是\InnerQuote{allowed}；好像祂曾说：\InnerQuote{祂容许了这些事发生，然而你并不因此而清白无过。}}\par

% structure: p0027 source=h2 target=subsection reason=relative-heading-rank; base=h2

\subsection{若 19:12}

因为他们从自己法律提出的指控毫无效果，就邪恶地转向外部法律说：\par

\Quote{凡自称为王的，就是反对凯撒。}\par

这人何处显出暴君的模样？你从何证明？凭那紫红袍？凭那王冠？凭那衣饰？凭那些兵士？祂岂不总是独行，除了祂的十二门徒，在各方面都遵循谦卑的生活方式，无论是饮食、衣着还是居所？啊，何等的无耻和不合时宜的懦弱！因为比拉多认为他若忽略这些话，现在会招致危险，便出来仿佛要查问此事（因为\Quote{坐下}显示了这点），但并未作任何查问，就把祂交给他们，以为能使他们羞愧。为证明他为此目的而行，请听他说：\par

% structure: p0031 source=h2 target=subsection reason=relative-heading-rank; base=h2

\subsection{若 19:14-15}

他们自愿服刑；因此天主也放弃了他们，因为他们首先将自己从祂的眷顾和照管中排除出去；既然他们异口同声地拒绝祂的主权，祂就任凭他们按自己的票决堕落。尽管如此，所说的话本应足以平息他们的激情，但他们害怕唯恐祂被释放后会再次吸引群众，就竭尽全力阻止这事。因为统治欲是可怕的事，可怕且能毁灭灵魂；正是由于这个原因，他们从未听从祂。然而比拉多因几句话就想释放祂，但他们紧逼说\Quote{钉死祂}。为什么他们要用这种方式杀死祂？这是一种耻辱的死。因此他们怕今后会有什么对祂的纪念，就想要把祂带到被诅咒的刑罚上，不知道真理因阻碍而被高举。为证明他们有这种猜疑，请听他们所说：\Quote{我们听见那骗子说过：\InnerQuote{三天后我要复活}}\BibleRef{玛 27:63}；为此他们大加鼓噪，把事情弄得天翻地覆，以便毁坏日后的事。而那被统治者败坏的混乱民众，不断呼喊\Quote{钉死祂！}。\par

3. 但我们不仅要阅读这些事，还要铭记于心：茨冠、紫袍、芦苇、掌击、颊打、唾面、讥讽。这些事若不断默想，足以消除一切愤怒；若我们被嘲笑、若我们遭受不公，仍要说：\BibleQuote{仆人不能大于主人}\BibleRef{若 13:16}，并要提起犹太人在狂乱中所说的话：\BibleQuote{你是个撒玛黎雅人，并附有魔鬼}\BibleRef{若 8:48}，以及\BibleQuote{祂靠贝耳则布驱魔}\BibleRef{路 11:15}。因为祂为此忍受了这一切，为使我们能跟随祂的足迹，并忍受那些比任何其他侮辱更令人不安的嘲笑。然而，祂不仅忍受了这些，甚至想尽一切办法拯救那些行恶的人，并使他们免于预定的惩罚。因为祂连宗徒也差遣去为他们得救，至少你们听他们说：\BibleQuote{我们知道你们是出于无知才做了这事}\BibleRef{宗 3:17}，藉此引导他们悔改。我们也要效法这一点；因为没有什么比爱仇敌、善待那些恶待我们的人，更能使天主满意。当有人侮辱你时，不要看他，而要看他背后的魔鬼，将你所有的怒气都倾泻在魔鬼身上，却要怜悯那个被魔鬼驱动的人。因为如果说谎来自魔鬼，那么无缘无故的发怒就更是如此。当你看见有人嘲笑他人时，要知道是魔鬼在驱动他，因为嘲笑不属于基督徒。因为那被吩咐要哀恸，且听过\BibleQuote{你们欢笑的人有祸了}\BibleRef{路 6:25}，而之后却侮辱、戏弄、激动的人，应该引起我们的惋惜，而非责难，因为基督想起\Person{犹达斯}时也曾忧伤。因此，这一切我们都要付诸行动，因为若我们不在这些事上行得正直，我们来到世界便是徒劳无益，甚至有害。因为信德不足以为人带来天国，不，它甚至更能定那些生活邪恶之人的罪；因为那\BibleQuote{知道主人的旨意，却不去行的，必受许多鞭打}\BibleRef{路 12:47}；又说：\BibleQuote{我若没有来，也没有对他们说话，他们就没有罪了}\BibleRef{若 15:22}。那么，我们还有什么借口呢？我们被安置在殿内，被认为配得屈身进入圣所，且成为释放的奥迹的分享者，却比那些从未分享过这一切的希腊人更糟糕？因为他们为了虚浮的荣耀尚且显出如此多的真智慧，我们更该追求一切美德，因为这中悦天主。但现在我们甚至不轻视财富；而他们常常不惜生命，在战争中为魔鬼的疯狂而舍弃子女，并为他们的魔鬼而轻视本性，但我们却不为基督轻视金钱，也不为天主的旨意而克制愤怒，反而怒火中烧，与发热的病人无异。正如他们被疾病缠身时浑身灼热，我们也被某种火窒息，在欲望上毫无节制，既增愤怒又增贪婪。为此，我羞愧惊异，当我看到希腊人中有人轻视财富，而我们中却个个疯狂。因为即使我们找到一些轻视财富的人，也会发现他们被其他恶习掳获，如情欲或嫉妒；要发现毫无瑕疵的真智慧实在困难。但原因是，我们不认真从圣经中获取药方，也不以痛悔、哀伤和叹息来应用圣经，而是漫不经心，若有闲暇才偶尔一读。因此，当世俗事务的洪流涌来时，便淹没一切；若曾有过任何益处，也毁掉了。因为如果有人受伤，敷上膏药后，却不扎紧，任其脱落，使伤口暴露于潮湿、灰尘、炎热以及成千上万能刺激它的东西，他便毫无益处；这不是因为药方无效，而是因为他的疏忽。这也常发生在我们身上，当我们对天主圣言不甚留意，却将自己完全且不断地投身于今生俗事时；这样，所有种子都被窒息，一无结果。为避免如此，让我们稍加留意，举目望天，俯身看逝者的坟墓和棺椁。因为同样的结局等着我们，同样的离世必要常在我们傍晚之前降临。因此，让我们为这远行准备；旅途需要许多补给，因为那里酷热、干旱、荒凉。自此以后，无法在客栈歇息，也无法购买任何东西，除非从今世带足一切。至少听听童女们的话：\BibleQuote{你们到卖油的人那里去吧}\BibleRef{玛 25:9}；但去了的人却没有找到。听听\Person{亚巴郎}的话：\BibleQuote{在我们与你们之间隔着深渊}\BibleRef{路 16:26}。听听\Person{厄则克耳}论及那日子所说的：\Person{诺厄}、\Person{约伯}和\Person{达尼尔}决不能救他们的子女\BibleRef{则 14:14}。但愿我们永远不必听到这些话，而是从今世带足通往永生的路粮，使我们能坦然无惧地瞻仰我们的主耶稣基督，愿光荣、权能、尊崇归于祂和圣父及圣神，从现在直到永远，万世无疆。阿们。\par

\SourceRecord{Translated by Charles Marriott.；Nicene and Post-Nicene Fathers, First Series；Vol. 14.；Edited by Philip Schaff.；Buffalo, NY: Christian Literature Publishing Co.,；1889.；<http://www.newadvent.org/fathers/240184.htm>.}

% structure: p0001 source=h1 target=section reason=page-title

\section{关于若望福音的讲道集 第八十五讲}

% structure: p0002 source=h2 target=subsection reason=relative-heading-rank; base=h2

\subsection{若望福音 19:16-18}

1. 成功具有可怕的力量，能使不小心的人跌倒或偏离正道。因此，犹太人起初享受天主的眷顾，却向外邦人要求王的法律，在旷野中，获得玛纳后，又想念韭菜。同样，这里他们拒绝基督的王国，却为自己召来了凯撒的王国。因此，天主按照他们自己的决定，给他们立了一位君王。当比拉多听到这话时，就将祂交出去钉十字架。这完全是没有理由的。因为当他应当查问基督是否图谋王权时，他却只因恐惧而宣判。但为了不让这事临到祂，基督事先说过：\Quote{我的国不属于这世界}；但他完全沉溺于眼前的事物，不会实践多大的智慧。然而，他妻子的梦本应足以吓倒他；但这一切都没有使他变好，他没有仰望上天，而是将祂交出。现在他们将十字架放在祂身上，如同对待一个罪犯。因为他们甚至厌恶那木头，连碰都不愿碰。在预像中也是如此：依撒格背了木头。但那时事情止于他父亲的意愿，因为那是预像；而这里却付诸行动，因为这是实体。\par

\Quote{祂来到一个名叫髑髅的地方。}有人说亚当死在那里，并葬在那里；而耶稣就在这死亡掌权的地方，也竖起了凯旋纪念碑。因为祂背着十字架出去，如同对抗死亡暴政的凯旋纪念碑；正如得胜者所做的那样，祂将胜利的象征扛在肩上。犹太人做这些事怀着不同的意图，这有什么关系呢？他们也把祂与两个强盗同钉，这也在无意中应验了预言；因为他们出于侮辱所做的一切，反而促成了真理，使你们明白它的力量有多大，因为先知早已预言说：\Quote{祂被列于罪犯之中。}\BibleRef{依 53:12} 因此，魔鬼想要掩盖所发生的事，却未能做到；因为三人被钉，但唯独耶稣是光荣的，使你们明白，祂的大能成就了一切。然而，奇迹发生在三人被钉在十字架上时；但没有人将所发生的一切归于其他两人，只归于耶稣一人；如此，魔鬼的阴谋完全落空，一切都回到了他自己头上。因为这两个人中，有一个得救了。所以，他并没有侮辱十字架的光荣，反而为其增色不少。因为在十字架上改变一个强盗，并把他带入乐园，并不比震动岩石更小的事。\par

% structure: p0005 source=h2 target=subsection reason=relative-heading-rank; base=h2

\subsection{若望福音 19:19}

同时，这也报复了犹太人，并为基督作了辩护。因为既然他们将祂当作无用之人交出，并试图让祂与强盗同受惩罚来证实这判决，为的是将来没有人能对祂提出恶毒的指控，或指责祂是无用邪恶之人，并堵住他们以及所有可能想要控告祂之人的口，并表明他们起而反对自己的君王，比拉多便将那些文字放在上面，如同放在一个凯旋纪念碑上，这些文字发出清晰的声音，彰显祂的胜利，宣告祂的王国，尽管不是完全的。而且他用三种语言揭示此事，并非单用一种语言；因为节期期间，犹太人中有混杂的人群，为了不让任何人不知道这辩护，他用所有语言公开记录了犹太人的疯狂。因为他们甚至在祂被钉时仍对祂怀有恶意。\Quote{但这对你们有何害处？毫无。因为如果祂是必死的、软弱的，即将消灭，你们为何害怕那些宣称祂是犹太人之王的文字呢？}他们要求什么？\Quote{你说\InnerQuote{他说}。因为现在这是一个断言，一个一般的判决，但如果加上\InnerQuote{他说}，指控就显示为出于他自己的鲁莽和傲慢。}但比拉多没有动摇，坚持他最初的判决。从这一情况中，所成就的并非小事，而是整件事。因为既然十字架的木头被埋了，因为没有人小心地把它取出来，因为恐惧紧迫，而信徒们正忙于其他紧急事务；而且后来它要被寻找，三个十字架很可能放在一起，为使主的十字架不被混淆，它首先通过位于中间而显明给众人，然后通过那块牌子。因为强盗的十字架没有牌子。\par

2. 兵丁分了衣服，却没有分那件外衣。看，预言由他们的恶行一一应验；因为这事也早已被预言；然而三人被钉，但预言在祂身上应验。为什么他们不对其他两人这样做，而唯独对祂？请你们也思考预言的精确性。因为先知不仅说他们\Quote{分了}，而且说他们\Quote{没有分}。因此，其余的衣服他们分了，外衣却没有分，而是将此事交给抽签决定。而\Quote{由上而下织成的}\BibleRef{若 19:23} 不是没有目的写下的；但有人说，这宣告了一个比喻性的断言，即被钉者不仅仅是人，而且也有从上而来的天主性。另有人说，福音作者描述了外衣的样式。因为在巴勒斯坦，他们将两幅布条缝合在一起，然后这样织成衣服，若望为了表明外衣是这种样式，就说\Quote{由上而下织成的}；在我看来，他说这话是暗指衣服的简陋，并且正如在其它一切事上，在服饰上祂也遵循简单的样式。\par

% structure: p0008 source=h2 target=subsection reason=relative-heading-rank; base=h2

\subsection{若望福音 19:24}

然而祂在十字架上，将祂的母亲托付给门徒，教导我们即使到最后一口气，也要对父母尽一切关怀。诚然，当她不恰当地烦扰祂时，祂说：\BibleQuote{女人，这与我和你有什么关系？}\BibleRef{若 2:4} 以及 \BibleQuote{谁是我的母亲？}\BibleRef{玛 12:48} 但在此祂流露出深切的慈爱，将她托付给祂所爱的门徒。若望再次因谦逊而隐晦自己；因为他若想夸耀，必也会列出被爱的缘由，那或许是一件伟大而奇妙的事。但祂为何不与若望谈及其他，也不在他沮丧时安慰他？因为那时不是用言语安慰的时候；此外，能获此殊荣并得到坚忍的赏报，对他而言绝非小事。但我请你思考，祂如何在十字架上毫不慌乱地成就一切：与门徒谈论祂的母亲，应验预言，向盗贼宣告美好的希望。然而在被钉之前，祂却显得出汗、痛苦、恐惧。这究竟意味着什么？没有什么难解，没有什么可疑。那里显示了本性的软弱，这里却展示了能力的超越。此外，透过这两件事，祂教导我们：即使在可怕的事面前感到困扰，也不因此而退缩；但一旦投入争战，就要认为万事皆有可能且容易。因此，让我们不要畏惧死亡。我们的灵魂本性爱惜生命，但我们可以松开本性的束缚，使这欲望变弱；或者紧握它，使欲望更加暴虐。正如我们有性交的欲望，但当我们实践真智慧时，便使这欲望变弱；生命之事也是如此。天主将肉欲结合于生儿育女，以维系我们的延续，却并未禁止我们走上更崇高的节欲之路；同样，祂在我们心中植入了对生命的爱，禁止我们自毁，但不妨碍我们轻视现世的生命。我们知晓这些，就当持守适度，既不可在任何时候自愿赴死，即使有千万可怕之事缠身；也不可在为悦乐天主而被引向死亡时，退缩畏惧，而要勇敢地投入，宁愿选择来世而非现世。\par

但妇女们站在十字架旁，那较软弱的性别此时却显得更为刚强\BibleRef{若 19:25}；从此一切都完全转变了。\par

% structure: p0011 source=h2 target=subsection reason=relative-heading-rank; base=h2

\subsection{若望福音 19:26-27}

3. 祂将自己的母亲托付给若望，说：\BibleQuote{看，你的儿子。}\BibleRef{若 19:26} 啊，何等的光荣！祂以何等的尊荣荣耀了这位门徒！当祂自己即将离去时，祂将母亲托付给门徒照顾。因为作为祂的母亲，她很可能会悲伤并需要保护，所以祂有理由将她委托给所爱的门徒。对那门徒说：\BibleQuote{看，你的母亲。}\BibleRef{若 19:27} 祂这样说，将他们在爱中结合起来；门徒明白了，就接她到自己家里。\Quote{但为何祂没有提及任何其他妇女，尽管另一位也站在那里？} 为教导我们给予母亲超乎寻常的尊敬。因为正如当父母在属灵的事上反对我们时，我们甚至不可认他们；同样，当他们不拦阻我们时，我们理当给予他们一切应有的尊敬，并超过他人，因为他们生了我们，养育了我们，为我们承受了千万可怕的事。祂用这些话堵住了\Person{马西昂}的无耻之言；因为祂若不是按肉身所生，也没有母亲，为何唯独为她如此操心呢？\par

% structure: p0013 source=h2 target=subsection reason=relative-heading-rank; base=h2

\subsection{若望福音 19:28}

这就是说：\Quote{救恩计划一无所缺。} 因为祂处处渴望表明，这死亡是崭新的一类，既然一切都在死者的权能之下，死亡并未在祂愿意之前临到身体；祂是在成就万事之后才愿意的。因此祂也说：\BibleQuote{我有权捨弃我的生命，也有权再取回它。}\BibleRef{若 10:18} 所以知道一切都成就了，祂说：\par

% structure: p0015 source=h2 target=subsection reason=relative-heading-rank; base=h2

\subsection{若望福音 19:28}

此处再次应验一项预言。但我请你思考，旁观者的可咒骂本性。虽然我们有千万仇敌，并曾忍受他们不堪的对待，但当我们看见他们灭亡时，仍会心软；然而他们甚至不因此与祂和解，也未因所见而被驯服，反而变得更加野蛮，更增添讥讽；他们将醋蘸在海绵上，如同给囚犯一般，递到祂嘴边；正因如此，牛膝草也被用上了。\par

% structure: p0017 source=h2 target=subsection reason=relative-heading-rank; base=h2

\subsection{若望福音 19:30}

你看见祂如何安然且有力地成就一切吗？下文就表明了这一点。因为当一切都完成了，\par

\BibleQuote{祂低下头（这部位未被钉住），便交付了灵魂。}\par

这就是说：\Quote{死了。} 然而气绝通常是在低头之后，但这里却相反。因为祂不是像我们一样，在断气之后才低头，而是先低下头，然后才断气。藉着这一切，福音作者表明了祂是万有的主。\par

但犹太人，他们吞下了骆驼却滤出了蠓虫，行了如此残暴的事，却对日子斤斤计较。\par

% structure: p0022 source=h2 target=subsection reason=relative-heading-rank; base=h2

\subsection{若望福音 19:31}

你是否看见真理是多么坚强？借着他们热心追求的那些事本身，预言便得以应验，因为借着这些事，这个与他们无关的明确预言得到了实现。兵士们到来时，打断了其他人的腿，却没有打断基督的腿。但这些人为了取悦犹太人，用长枪刺透了他的肋旁，如今还侮辱尸体。啊，可憎可咒的意图！然而，亲爱的，你不要困惑，不要沮丧；因为这些人的邪恶意图所行的事，却为真理而战。因为有一个预言说：（借着此事）\BibleQuote{他们要仰望他们所刺透的。}\BibleRef{匝 12:10}不仅如此，这胆大妄为的行为，对于那些后来会怀疑的人——如\Person{多默}及像他一样的人——也是一个信德的证明。与此一起，还有一个不可言喻的奥迹得以完成。因为\BibleQuote{流出了水和血。}这些泉源流出，并非无目的或偶然，而是因为教会由这两者共同构成。受过洗的人知道这一点：借着水而重生，并由圣血和圣体滋养。因此，圣事由此开始；当你接近那可畏的杯爵时，你便如此接近，如同从那肋旁饮用一样。\par

% structure: p0024 source=h2 target=subsection reason=relative-heading-rank; base=h2

\subsection{若望 19:35}

也就是说：\BibleQuote{我不是从别人听见，而是我亲自在场并看见了，这见证是真实的。}可以这样认为。因为他讲述了一个侮辱的行为；他并没有讲述什么伟大而可赞叹的事，以致你会怀疑他的叙述；而是要堵住异端者的口，并大声预告将来要发生的圣事，并注视其中所藏的宝藏，他对所发生的事非常精确。那预言也应验了，\par

% structure: p0026 source=h2 target=subsection reason=relative-heading-rank; base=h2

\subsection{若望 19:36}

因为即使这句话是针对犹太人的羔羊说的，但仍是预像先行，为的是那实体，而在他内预言更加圆满地应验了。因此，福音作者引用了先知的话。因为他若不断以自己为证人，就会显得不可信，所以他请来\Person{梅瑟}帮助他，并说，这件事也并非无目的发生，而是古时早已写下的。这句话的意思是：\BibleQuote{他的骨头一根也不可折断。}他又以自己的见证证实了先知的话。他说：\Quote{这些事我告诉你们，是为了叫你们知道预像与实体之间的联系是多么紧密。}你看见他费了多少苦心，使那看似可耻羞辱的事被相信吗？因为士兵连尸体都侮辱，远比被钉十字架更糟。他说：\Quote{尽管如此，我还是将这些事告诉了你们，并且非常恳切地告诉你们，\InnerQuote{为使你们相信。}}\BibleRef{若 19:35}但愿没有人不信，也不要因羞愧而损害我们的事业。因为那些看似最可耻的事，正是我们福泽的极其可敬的记录。\par

% structure: p0028 source=h2 target=subsection reason=relative-heading-rank; base=h2

\subsection{若望 19:38}

不是宗徒中的一位，或许是七十位中的一位。因为那时他们以为犹太人的愤怒因十字架而平息了，便毫无恐惧地前来，负责他的葬礼。于是\Person{若瑟}前来，向\Person{比拉多}请求恩准，\Person{比拉多}准许了；他为什么不呢？\Person{尼苛德摩}也协助他，提供了昂贵的殓葬。因为他们仍然倾向于把他当作一个普通人看待。他们带来了那些具有特殊性质的香料，能长久保存尸体，不让它迅速腐朽，这是那些对他并不存有任何伟大想法的人的行为；但无论如何，他们表现出了极其热爱的情感。但为什么十二位中没有人来，无论是\Person{若望}，\Person{伯多禄}，还是其他更杰出的门徒？作者也没有隐瞒这一点。若有人说是因为害怕犹太人，这些人也处于同样的恐惧中；因为\Person{若瑟}，经上也说，是\BibleQuote{因怕犹太人，暗暗地做门徒。}没有人可以说\Person{若瑟}这样做是因为他极其藐视他们，而是虽然他自己害怕，仍然来了。但\Person{若望}在场，亲眼看见他断气，却没有做类似的事。在我看来，\Person{若瑟}是一个地位很高的人（从葬礼便可看出），并且为\Person{比拉多}所认识，因此他获得了恩准；然后他埋葬了他，不是像埋葬一个罪犯，而是慷慨地，按照犹太人的方式，仿佛埋葬一位伟大而可敬的人。\par

4. 由于时间紧迫（因为死亡发生在第九时辰，很可能在去见比拉多和取下尸体之间，傍晚就要降临，那时不许工作），他们把他安放在附近的坟墓里。这是天主眷顾的安排：他被放在一个新坟墓里，里面从未葬过人，这样他的复活就不会被认为是与同葬的另一个人一同复活；也为了使门徒们能够轻易前来目睹所发生的事，因为那地方很近；并且不仅让门徒们见证他的埋葬，也让他的敌人见证，因为封上坟墓的印痕和士兵坐在那里看守，都是见证埋葬的人的行为。因为基督热切地希望这件事被宣认，不亚于复活。因此，门徒们也非常热心地证明他死了。因为复活会被后来的一切时代所证实，但死亡，如果在那时被部分掩盖或没有非常显明地展示，可能会损害复活的叙述。他被安放在附近，不仅为了这些原因，也为了使关于尸体被盗的说法不攻自破。\par

% structure: p0031 source=h2 target=subsection reason=relative-heading-rank; base=h2

\subsection{若望 20:1}

因为祂复活时，石头和封印仍覆在祂身上；但因为必须使其他人完全确信，坟墓在复活后被打开，所发生的事便得到证实。这正是感动\Person{玛利亚}的原因。她对主充满热爱之情，安息日一过，她便不能安息，而是清早早早地来到，渴望从这地方得些安慰。但她看见那地方和石头被挪开，既没有进去也没有弯腰，却因极大的渴望跑到门徒那里；因为她急切地渴望，希望尽快知道身体怎么样了。这就是她跑来的含义，她的话也说明了这一点。\par

% structure: p0033 source=h2 target=subsection reason=relative-heading-rank; base=h2

\subsection{\BibleRef{若 20:2}}

你看到她尚未对复活有清楚的了解，以为身体被挪移了，便简单地把一切告诉了门徒。圣史没有剥夺这妇人的如此赞美，也不以为耻：他们是从她——那彻夜守望的人——首先得知这些事的。如此，他性情中热爱真理的本性处处闪耀。当她来并说了这些事，他们听见，就极其急切地走近坟墓，看见殓布放着，这是复活的记号。因为，如果人挪走了身体，他们不会事先剥去衣服；如果有人偷了它，也不会费事解开头巾、卷好、另放一处。而是会直接取走身体。为此，\Person{若望}预先告诉我们，埋葬时用了很多没药，这使殓布紧贴身体，牢固如铅；为使你听到头巾另放时，不会容忍那些说祂被偷的人。因为一个小偷不会如此愚蠢，在多余的事上费这么多功夫。他为何要解开衣服？如果做了，怎能逃脱侦测？他可能会花很多时间，并因拖延逗留被发现。但为什么衣裳分开放着，头巾却单独卷好？使你明白这不是人混乱仓促的行为——把一些放一处，一些放另一处，并卷在一起。由此他们相信了复活。为此，基督后来在他们因所见而被说服时向他们显现。再请注意圣史的无自夸之心：他如何为\Person{伯多禄}查看的仔细作证。因为他自己先于\Person{伯多禄}，看见殓布后没有进一步查问，就离开了；但那位热忱的门徒更深入进去，仔细察看了一切，看到更多，然后另一位也被叫来看。因为他在\Person{伯多禄}之后进入，看见殓布放着，且是分开的。分开，并将一件放一处，另一件卷起后另放一处，是某人细心行事之举动，并非偶然或仓促慌乱的样子。\par

5. 但是，当你听到你的主赤裸地复活了，就应从你对丧葬的疯狂中醒悟过来；因为那种多余而无益的花费有什么意义呢？它给哀悼者带来许多损失，对逝者却毫无益处，或者（如果我们必须说它带来什么的话）反而有害？因为昂贵的葬礼往往导致坟墓被盗，使得被精心安葬的人被赤裸地抛出来，不得安葬。唉，虚荣啊！它甚至在悲哀中显示出多大的暴虐！这是何等的愚蠢！许多人为了防止这种事发生，将那些华服剪碎，填满许多香料，以便对侮辱死者的人毫无用处，然后才埋入地下。这不是疯子的行为吗？不是神志不清者的行为吗？为了炫耀自己的野心，然后又毁掉它？有人说：\Quote{是的，我们使用所有这些计策，是为了让它们与死者一同安全地躺卧。}那么，如果盗贼得不到它们，蛀虫和蛆虫就不会得到它们吗？如果蛀虫和蛆虫得不到它们，时间和腐烂的湿气就不会毁掉它们吗？但让我们假设，既没有盗墓者，也没有蛀虫、蛆虫、时间或其他任何东西毁掉坟墓中的东西，而身体本身一直保持完好直到复活，这些衣料也保持崭新、鲜亮、精致；这对逝者有什么益处呢？当身体赤裸地复活，而这些东西留在这里，对我们必须交出的账目毫无帮助时？那么，有人说，\InnerQuote{为什么在基督的事上也这样做呢？}首先，不要将这些与人的事相比，因为那个妓女甚至将香液倒在他的圣足上。但如果我们必须谈论这些事，我们说，这些事情发生时，做这些事的人还不知道复活的道理；因此经上说：\Quote{如同犹太人的习俗。}因为尊敬基督的人不是那十二宗徒，而是那些并不十分尊敬祂的人。那十二宗徒不是这样尊敬祂，而是通过为祂而死、被屠杀和危险来尊敬祂。那另一种尊敬固然是尊敬，但远不如我所说的这种。此外，正如我开始时所说的，我们现在谈论的是人，但那时这些事是为主而做的。为使你知道基督并不看重这些事，祂说：\BibleQuote{我饿了，你们给了我吃的；我渴了，你们给了我喝的；我赤身露体，你们给了我穿的}\BibleRef{玛 25:35}；但祂从未说过：\Quote{我死了，你们埋葬了我。}我说这话，并不是要废除埋葬的习俗（绝无此意），而是要制止其铺张和不合时宜的虚荣。有人说：\InnerQuote{然而，对逝者的感情、悲哀和同情使人这样做。}这种作法并非出于对逝者的同情，而是出于虚荣。因为如果你真的想同情死者，我会向你展示另一种哀悼的方式，并教你给他穿上那将与他一同复活、使他荣耀的衣裳。因为这些衣裳不会被蛆虫蛀蚀，不会被时间磨损，也不会被盗墓者偷走。那么这些是什么衣裳呢？就是施舍的衣裳；因为这是一种将与他一同复活的袍子，因为施舍的印记与他同在。穿着这些衣裳的人，那时将听到：\BibleQuote{你们饥饿时，给了我吃的。}这些衣裳使人出众，这些衣裳使他们荣耀，这些衣裳使他们安全；而现今所用的那些，只是给蛀虫提供食物，给蛆虫提供餐桌。我说这些，并不是禁止丧葬礼仪，而是劝你们适度操办，只求遮盖身体，不要使之赤裸地还土。因为如果活着时，祂要我们仅以够遮盖为满足，何况死后呢？因为死人的身体比活人呼吸时更不需要衣裳。因为活着时，因寒冷和端庄的缘故，我们需要衣裳遮盖；但死后，并不因这些理由需要殓衣，而是为了不让身体赤裸；而比殓衣更好的，我们有大地，它是最美妙的遮盖物，更适合于我们这类身体的本性。因此，既然在这么多需求下，我们不应追求任何多余的东西，那么在没有这种必要的情况下，虚荣更是不得体的。\par

6. 有人说：\Quote{旁观者会嘲笑。}确实，若有嘲笑，我们无需太在意如此极其愚昧的人；但如今，许多人反倒钦佩并接受我们的真智慧。因为这些事不值得嘲笑，而是我们现在所做的——哭泣、哀号、与逝者同葬——才该被嘲笑和惩罚。然而，在这些事上表现真智慧，连同服饰的朴素，会为我们预备冠冕和赞美，众人都将赞赏我们，并钦佩基督的大能，说：\Quote{何等奇妙！被钉者的大能何其伟大！他说服了那些将要灭亡和朽坏的人，死亡并非死亡；因此他们不行将亡人之事，而是如同将死者先送到遥远更美好的居所。他说服了他们，这可朽坏、属地的身体，将披上比丝绸或金缕衣更荣耀的衣服，即不朽的衣服；因此他们不十分忧虑葬礼，而是将德善的生活视为可钦佩的裹尸布。}他们若看见我们显出真智慧，便会说这些话；但若见我们被忧伤压倒，如妇人般啼哭，围绕一群女哭丧者，他们就会嘲笑、讥讽，并以种种方式指责，拆穿我们愚昧的花费和徒劳的劳苦。我们常听见众人指责这些事，而且非常合理。因为当我们装饰一个被腐败和虫子蚕食的尸体，却忽略饥渴、赤身露体、作客旅的基督时，我们还有何借口？让我们停止这徒劳的烦恼吧。为逝者举行葬礼，要像对我们和他们都有益、为光荣天主那样去行：我们要为他们的缘故多行施舍，将最好的路费与他们一同送出。因为，如果可钦佩之人的纪念，即使他们已死，也能保护活着的人（因为经上说：\Quote{我为自己，并为我的仆人达味，必要保卫这座城}\BibleRef{列下 19:34}），那么施舍更能成就此事；因为施舍甚至曾使死人复活，就如寡妇们站在伯多禄周围，将多尔卡同她们在一起时所做的衣服指给他看一般。\BibleRef{宗 9:39}因此，当一个人快要死时，让那将死之人的朋友预备葬礼，并劝说临终者留些什么给穷人。让他用这些衣服送他进坟墓，立基督为他的继承人。因为，若那些将君王立为继承人的人，给他们的亲属留下稳妥的份额，那么当一个人立基督与他的子女同为继承人时，请想，他将为自己和所有人带来多大的益处！这些是正确的葬礼，这些有益于留下的人和离去的人。如果我们这样被埋葬，我们在复活之时将是荣耀的。但若我们只顾身体而忽略灵魂，我们将遭受许多可怕的事，并招致许多嘲笑。因为，死时不披戴德行，并非寻常的不体面；而且，身体即使被抛弃而无坟墓，也不如灵魂在那一天显露无德行那般耻辱。让我们穿上德行，让我们用它包裹自己：最好在我们一生中都这样做；但若我们在此生疏忽了，至少要在临终时醒悟，嘱咐我们的亲属在我们离世时以施舍帮助我们；这样，彼此相助，我们就能获得极大的信心，赖我们的主耶稣基督的恩宠和慈爱，愿光荣、权能、尊崇归于他及圣父与圣神，自今至永远，及世之世。阿们。\par

\SourceRecord{Translated by Charles Marriott.；Nicene and Post-Nicene Fathers, First Series；Vol. 14.；Edited by Philip Schaff.；Buffalo, NY: Christian Literature Publishing Co.,；1889.；<http://www.newadvent.org/fathers/240185.htm>.}

% structure: p0001 source=h1 target=section reason=page-title

\section{\WorkTitle{若望福音讲道集第八十六篇}}

% structure: p0002 source=h2 target=subsection reason=relative-heading-rank; base=h2

\subsection{\BibleRef{若 20:10-11}}

1. 女性常充满温情，更倾向于怜悯。我说这话，免得你们希奇玛利亚为何在坟墓前痛哭，而伯多禄却未如此感动。因为经上说：\BibleQuote{\OriginalTerm{门徒}{disciples}，} \BibleQuote{各自回家去了}；但她却站在那里流泪。这是因为女性的性体柔弱，而她又尚未确切明了复活的道理；门徒们看到殓布就信了，于是惊异地各自回家去了。他们为什么不立刻往加里肋亚去，正如在受难前领受的命令那样呢？他们或许是等候其他人，并且还处在极大的惊异之中。于是门徒们走了，但她却仍在原处站立；因为正如我所说，仅看到坟墓就给了她极大的安慰。你看她，为了减轻悲伤，弯下腰来，想要看清安放尸体的地方。因此她因这极大的热诚而获得了不小的赏报。因为门徒们没有看到的，这女人先看到了：天使坐着，一个在脚边，一个在头边，穿着白衣；连那穿戴也充满了荣光和喜乐。由于这女人的心思尚未提升到足以从巾帕的证据接受复活，于是有更多的事发生了，她看见了更多的事：天使穿着发光的衣服坐着，以便如此将她从强烈的悲伤中提升起来，并安慰她。但他们关于复活却什么也没有对她说，然而她温良地被引导进入这一教义中。她看到了光明而异常的容貌，她看到了发光的衣服，她听到了同情的说话声。因为天使说了什么？\par

% structure: p0004 source=h2 target=subsection reason=relative-heading-rank; base=h2

\subsection{\BibleRef{若 20:13}}

借着这一切情况，仿佛为她开了一道门，她逐渐地被引导认识复活。他们坐着的姿态也邀请她发问，因为他们显示出知道所发生的事；为此他们也不是坐在一起，而是彼此分开。因为似乎她不敢立刻请问他们，于是他们借着问她和她坐着的姿态，引导她交谈。那么她说了什么呢？她说话十分热情而深情：\par

\BibleQuote{他们把我的主搬走了，我不知道他们把祂放在了哪里。}\par

\Quote{你在说什么？你还不知道有关复活的事，仍然想着祂被安放吗？}你看她是如何尚未接受崇高的教义？\par

% structure: p0008 source=h2 target=subsection reason=relative-heading-rank; base=h2

\subsection{\BibleRef{若 20:14}}

那么这是怎样的顺序呢？她对他们说了话，还没有从他们那里听到什么，就转过身去了？我想，当她说话的时候，基督突然出现在她背后，使天使们感到敬畏；他们看见了自己的主，立刻以姿态、神情、动作表明他们看见了主；这引起了女人的注意，使她向后转身。祂这样向他们显现，却不像这样向女人显现，为的是不立即惊吓她，而是以较卑微平凡的形象，从她以为祂是园丁这一点就可以清楚看出。应当以温和而非突然的方式，将如此卑微心思的人引向崇高的事。因此祂反过来问她：\par

% structure: p0010 source=h2 target=subsection reason=relative-heading-rank; base=h2

\subsection{\BibleRef{若 20:15}}

这显示出祂知道她想要问什么，并引导她回答。女人明白了这一点，就不再提耶稣的名字，而是好像她的问者知道她询问的主题，回答说：\par

\BibleQuote{先生，若是你把他搬走了，请告诉我你把他放在哪里，我便去取回他来。}\par

她又谈论安放、搬走、拿走，好像谈论一具尸体。但她的意思是：\Quote{若是你因怕犹太人而把他搬走了，请告诉我，我去取回他。}这女人的恩情和挚爱是伟大的，但还没有什么崇高之处。因此主现在不是借着形貌，而是借着声音将此事摆在她面前。因为正如祂有时被犹太人认识，有时虽在场却不被察觉；同样在说话时，祂在祂选定的时刻，就使人认出祂来；正如祂曾对犹太人说：\BibleQuote{你们找谁？}，他们既不认识脸庞也不认识声音，直到祂愿意。这里也是如此。祂只叫了她的名字，责备并责怪她竟对一位活着的人存有这种想法。但是，这是怎么回事呢？\par

% structure: p0014 source=h2 target=subsection reason=relative-heading-rank; base=h2

\subsection{\BibleRef{若 20:16}}

% structure: p0015 source=h2 target=subsection reason=relative-heading-rank; base=h2

\subsection{\BibleRef{若 20:17}}

2. 有些人断言，她请求了属神的恩宠，因为她曾听见祂与门徒在一起时说：\BibleQuote{我若往父那里去，我要求父，父就给你们另派一位护慰者来。}\EditorialNote{见若14:3,16} 但她当时并不在场与门徒同在，怎能听见这话呢？此外，这种想象与这里的意思相去甚远。并且，祂尚未往父那里去，她怎能祈求呢？那么意思是什么呢？我想，她仍愿像从前一样与祂交谈，而且在她的喜乐中，她对祂并没有感受到什么伟大的事，虽然祂在肉身中已成为远更卓越的。因此，为了将她从这种想法中领出，并使她以极大的敬畏与祂说话（因为此后祂与门徒的相处也不像从前那样亲密），祂提升了她的思想，使她更恭敬地注意祂。若是说：\Quote{不要像从前那样接近我，因为情况不同了，我将不再以同样的方式与你们同在}，那会是粗暴而高亢的；但说：\par

\BibleQuote{我还没有升到父那里去}，虽然听来不刺耳，但仍是宣告同一件事的话。因为借着说\BibleQuote{我还没有升上去}，祂显示出祂急忙赶往那里；而且，一个即将离开那里，不再与世人交往的人，不应被以从前同样的情感来看待。后面的经文也表明这是实情。\par

\BibleQuote{你去告诉我的弟兄们，说：我升到我的父和你们的父那里去，升到我的天主和你们的天主那里去。}\par

但祂不会立即如此做，而是在四十天之后。那么祂为何这样说呢？为了提升他们的心神，并说服他们祂是升上天庭。而\Quote{到我父和你们父那里去，到我的天主和你们的天主那里去}这句话，属于\OriginalTerm{救恩计划}{Dispensation}，因为\Quote{上升}也属于祂的\OriginalTerm{血肉}{Flesh}。因为祂是对一个没有高尚思想的人说这些话。\Quote{那么，父是祂的父，与是我们的父，方式不同吗？}确实不同。因为如果祂是义人的天主，与祂是其他人的天主方式不同，那么在圣子与我们的情形下更是如此。因为祂曾说过：\Quote{告诉我的弟兄们}，免得他们由此想象任何平等，祂便显示了这个区别。祂将坐在父的宝座上，而他们则要站在一旁。所以，虽然祂按血肉的\OriginalTerm{本体}{Subsistence}成了我们的兄弟，但在尊荣上祂远超我们，其程度甚至无法言说。\par

% structure: p0020 source=h2 target=subsection reason=relative-heading-rank; base=h2

\subsection{\BibleRef{若 20:18}}

恒心与忍耐是何等大的善。但为何他们不再因祂即将离去而忧伤，也不像先前那样说话呢？那时他们如此忧苦，是以为祂将要死去；但现在既然祂已复活，他们还有什么理由忧伤呢？况且，玛利亚报告了祂的显现和祂的话，这些足以安慰他们。既然门徒们听了这些事，很可能要么不相信那妇人，要么相信了却因祂没有认为他们配得显现而忧伤，尽管祂应许在加里肋亚会见他们；为使他们不因思虑这事而不安，祂不让一天过去，而是借着他们知道祂已复活，以及从妇人那里听到的事，使他们处于渴慕状态，当他们极欲见祂，且极其恐惧时（这本身尤其使他们的渴望更强烈），祂便在晚上显现在他们面前，并且极其奇妙。为什么祂在\Quote{晚上}显现呢？因为那时他们很可能特别恐惧。但奇妙的是，为何他们不以为祂是幽灵；因为祂\Quote{当门关着时}进入，并且是突然的。主要原因在于，那妇人预先在他们心中产生了极大的信德；此外，祂向他们显示容光焕发、温和慈祥的面容。祂不是在白天来的，为叫众人都聚集在一起。因为惊异是巨大的；祂没有敲门，而是突然站在中间，并显示祂的肋旁和双手。同时，祂也以声音平息他们翻腾的思绪，说道：\par

% structure: p0022 source=h2 target=subsection reason=relative-heading-rank; base=h2

\subsection{\BibleRef{若 20:19}}

就是说：\Quote{不要不安}，同时提醒他们在受难前祂对他们说的话：\Quote{我把平安留给你们}\BibleRef{若 14:27}；以及：\Quote{在我内你们有平安，但在世界上你们将有苦难。}\BibleRef{若 16:33}\par

% structure: p0024 source=h2 target=subsection reason=relative-heading-rank; base=h2

\subsection{\BibleRef{若 20:20}}

你看见话语在行为中实现了吗？因为祂在受难前曾说：\Quote{我要再见你们，你们的心要喜乐，你们的喜乐没有人能夺去}\BibleRef{若 16:22}，现在祂在行为上实现了这一切；但所有这些事都引导他们达到最确切的信德。因为既然他们与犹太人进行着无休止的战争，祂就不断重复\Quote{愿平安与你们同在}，赐给他们平安以平衡战争。因此这是祂复活后对门徒说的第一句话（为此保禄也常说：\Quote{愿恩宠与平安归于你们}），而对妇女，祂则报喜讯，因为那个性别正处忧愁中，并以此作为第一个诅咒。因此祂分别给予适合的喜讯：对男人赐予平安，因为他们在战争中；对女人赐予喜乐，因为她们在忧愁中。然后，在除去一切痛苦之后，祂讲述十字架的胜利，这些就是\Quote{平安}。\Quote{既然一切阻碍都已除去，}祂说，\Quote{我已使我的胜利荣耀，一切都已成就，}（然后祂接着说，）\par

% structure: p0026 source=h2 target=subsection reason=relative-heading-rank; base=h2

\subsection{\BibleRef{若 20:21}}

\Quote{你们没有困难，因为既成之事和派遣你们的那一位的尊荣。}在这里祂提升他们的灵魂，并向他们显示他们极大的信心依据，既然他们要承担祂的工作。并且祂不再诉诸父，而是以权威赋予他们能力。因为：\par

% structure: p0028 source=h2 target=subsection reason=relative-heading-rank; base=h2

\subsection{\BibleRef{若 20:22-23}}

正如一位君王派遣总督，赐予他们下监和释放的权力，同样，基督派遣他们时，也赐予他们这同样的权力。但是祂怎么说：\BibleQuote{我若不去，祂就不会来}\BibleRef{若 16:7}，却仍然赐给他们圣神呢？有人说，祂并没有赐予圣神，而是使他们准备好去接受祂，通过向他们吹气。因为如果达尼尔看见天使时就害怕，那么当他们领受那不可言喻的恩赐时，除非祂先使他们成为学习者，否则他们将何等受苦呢？因此祂没有说：\Quote{你们已领受了圣神}，而是说：\Quote{领受圣神吧}。然而，人若断言他们当时也领受了一些属灵的能力和恩宠，也不算错；但不是为了复活死人或行奇迹，而是为了赦罪。因为圣神的恩赐有多种；因此祂补充说：\Quote{你们赦免谁的罪，谁的罪就赦免了}，表明祂所赐予的是何种权能。但在另一种情况下，四十天之后，他们领受了行奇迹的能力。因此祂说：\BibleQuote{当圣神降临于你们身上时，你们将领受能力，并要在耶路撒冷及全犹太地作我的见证人。}\BibleRef{宗 1:8}他们因着奇迹而成为见证人，因为圣神的恩宠是不可言喻的，恩赐是多种形式的。但这发生，是为叫你们知道圣父、圣子及圣神的恩赐和权能是唯一的。因为那些似乎专属圣父的事，也可见于圣子和圣神。有人说：\Quote{那么，没有人能到圣子那里，\InnerQuote{除非父吸引他}？}\BibleRef{若 6:44}然而，这事同样显示也属于圣子。\Quote{我}，祂说，\Quote{就是道路；除非藉着我，没有人能到父那里去。}\BibleRef{若 14:6}并且注意，这也属于圣神；因为\BibleQuote{除非受圣神感动，没有人能说耶稣基督是主。}\BibleRef{格前 12:3}此外，我们看到，宗徒们有时由圣父赐给教会，有时由圣子，有时由圣神；并且\BibleQuote{恩赐的多样}\BibleRef{格前 12:4}属于圣父、圣子及圣神。\par

4. 那么，让我们尽一切努力使圣神与我们同在，并极其尊敬那些受托执行其运作的人。因为司铎的尊位是伟大的。经上说：\BibleQuote{你们赦免谁的罪，谁的罪就赦免。}因此保禄也说：\BibleQuote{你们要听从那些管理你们的人，并要服从。}\BibleRef{希 13:17}并且要非常尊敬他们。因为你确实关心自己的事务，如果你安排得好，你不必为他人交账；但司铎即使自己的生活安排得正确，如果他不焦虑地关心你以及周围所有人的事务，他将与恶人一同进入地狱；而且常常即使他不被自己的行为所出卖，他也会因你的行为而丧亡，如果他没有完全尽到自己的本分。因此，知道这危险之大，你们要慷慨地对待他们；保禄也暗示了这一点，他说：\BibleQuote{他们为你们的灵魂警醒，}并且不只是简单地说，而是\BibleQuote{如同那些将要交账的人。}\BibleRef{希 13:17}所以他们应当从你们那里得到极大的关注。但如果你们与其他人一起践踏他们，那么你们的事务也不会顺利。因为当舵手保持勇气时，船员也会安全；但如果他因他们的辱骂和敌意而疲惫不堪，他就不能同样警醒，也不能保持他的技艺，从而无意中将他们抛入无数的灾难之中。同样，司铎若从你们那里得到尊敬，就能好好安排你们的事务；但如果你们使他们沮丧，你们就削弱了他们的手，使自己和他们都容易成为风浪的猎物，即使他们非常勇敢。想想基督关于犹太人的话：\BibleQuote{经师和法利塞人坐在梅瑟的座位上，凡他们所吩咐你们的，你们都要遵守。}\BibleRef{玛 23:2}现在我们不必说\Quote{司铎坐在梅瑟的座位上}，而是\Quote{坐在基督的座位上}；因为他们相继接受了祂的教导。因此保禄也说：\BibleQuote{我们为基督作大使，好像天主借着我们劝你们一样。}\BibleRef{格后 5:20}难道你们没有看见，在外邦统治者面前，所有人都向他们鞠躬，甚至常常是那些在家族、生活、智慧上都优于审判者的人？然而，由于任命他们的人，他们都不考虑这些，而是尊重他们的治理者的决定，无论谁被赋予统治权。那么，当人任命时就有如此恐惧，而当天主任命时，我们却藐视被任命者，辱骂他，用无数的责备玷污他，并且虽然被禁止判断我们的弟兄，我们却磨尖舌头攻击司铎？这怎么能得原谅，当我们看不见自己眼中的大梁，却苛刻地挑剔别人眼中的木屑？难道你们不知道，这种判断会使你们自己的判断更严厉吗？我说这些，并不是赞同那些不配执行司铎职分的人，而是深深怜悯和为他们哭泣；然而我因此并不认为他们应当被他们管辖的人判断是合理的。即使他们的生活受到很多指责，如果你谨慎自守，你在天主所委托给他们的事情上绝不会受到任何损害。因为如果天主借着驴发出声音，并借着巴郎这个占卜者的愚昧口和不洁舌，为冒犯的犹太人赐予属灵祝福，那么为了你们这些正直人的缘故，即使司铎极其卑劣，祂也会完成一切属于祂的工作，并派遣圣神。因为不是纯洁的人因自己的纯洁引来圣神，而是恩宠成就一切。因为经上说：\BibleQuote{一切都是为你们的，无论是保禄，或是阿颇罗，或是刻法。}\BibleRef{格前 3:22}因为置于司铎手中的事物，只有天主能赐予；无论人的智慧达到多高，它都比那恩宠低下。我说这些，不是要我们放任自己的生活，而是当你们中的一些上司生活疏忽时，你们被治理者不必常常为自己堆积邪恶。但我何必提司铎？即使天使或总领天使也不能对天主所赐予的任何事物做什么；而是圣父、圣子和圣神安排一切，而司铎借出他的舌头，提供他的手。因为不应当因他人的邪恶，使那些怀着信心来到救恩标记前的人受到伤害。知道这一切之后，让我们敬畏天主，尊敬祂的司铎，对他们表示一切尊敬；这样，因我们自己的善行和对他们的关注，我们将在我们的主耶稣基督的恩宠和慈爱下，从天主那里获得丰厚的回报；愿光荣、权能和尊崇归于祂和圣父及圣神，自今至永远，直到世世代代。阿们。\par

\SourceRecord{Translated by Charles Marriott.；Nicene and Post-Nicene Fathers, First Series；Vol. 14.；Edited by Philip Schaff.；Buffalo, NY: Christian Literature Publishing Co.,；1889.；<http://www.newadvent.org/fathers/240186.htm>.}

% structure: p0001 source=h1 target=section reason=page-title

\section{若望福音讲道词第八十七篇}

% structure: p0002 source=h2 target=subsection reason=relative-heading-rank; base=h2

\subsection{若 20:24-25}

1. 轻易而随意地相信，出自过于随和的性情；而过分好奇与好管闲事，则表明一种极其粗笨的理解。为此，多默受到责备。因为他没有相信宗徒们所说\Quote{我们见了主}的话；他与其说是不信任他们，不如说是认为那件事不可能，即从死者中复活。因为他不是说\Quote{我不信你们}，而是说\Quote{除非我用手摸——我不信}。但那时，当众人都聚集时，为何唯独他缺席？大概是因为不久前的分散，他那时尚未回来。但你，当你看到门徒的不信时，要思量主的慈爱，祂为了一个灵魂，带着祂的伤痕显现，并且前来拯救那一个人，尽管他比其余人更为迟钝；正因如此，他要求从最粗重的感官中寻求证据，甚至不愿相信他的眼睛。因为他不说\Quote{除非我看见}，而是说\Quote{除非我摸到}，以免他看见的或许是一种幻象。然而告诉他的那些门徒当时是可信的，许下诺言的那位也是可信的；可是，既然他渴望更多，基督连这一点也没有拒绝他。\par

祂为何没有立即显现给他，而是等到\Quote{八天后}？\BibleRef{若 20:26} 为的是在这期间他不断受门徒的教导，听到同一件事，从而更加渴望，以后更易于相信。但他从哪里知道祂的肋旁被刺开了？是从门徒那里听说的。那么他如何一部分相信，一部分不相信？因为这件事非常奇异奇妙。但请注意，我恳求你，宗徒们是何等真诚，他们不隐藏任何过失，无论是自己的还是别人的，而是以极大的真实性记录下来。\par

耶稣再次向他们显现，并没有等待多默的请求，也没有听到他任何这样的话，而是在他说话之前，祂自己就先满足了他的愿望；表明当他对门徒说那些话时，祂已经临在。因为祂用了同样的话，以一种严厉责备且为将来教训的方式。因为祂说了：\par

% structure: p0006 source=h2 target=subsection reason=relative-heading-rank; base=h2

\subsection{若 20:26}

\Quote{不要作不信的，而要作相信的。}\par

你看见他的怀疑出自不信吗？但这是在领受圣神之前；此后就不再如此，而今后他们得以成全。\par

耶稣不仅这样责备他，也通过随后的事：当他完全满足后，舒了口气，大声呼喊：\par

% structure: p0010 source=h2 target=subsection reason=relative-heading-rank; base=h2

\subsection{若 20:28}

% structure: p0011 source=h2 target=subsection reason=relative-heading-rank; base=h2

\subsection{若 20:29}

因为信德在于接受不可见的事；因为\BibleQuote{信德是所希望之事的实质，是未看见之事的确证。}\BibleRef{希 11:1} 在此祂不仅宣示门徒是有福的，也包括他们之后会相信的人。\Quote{然而，}有人说，\Quote{门徒看见了就相信了。}是的，但他们没有寻求任何此类的事，而是从殓布的证明立刻接受了关于复活的话，并在看见身体之前就表现了全部信德。因此，现今若有人说：\Quote{我希望活在那些时代，看见基督行奇迹}，让他们想想：\Quote{那些没有看见而相信的人是有福的。}\par

值得探究的是，一个不朽坏的身体如何显示出钉痕，并且能被凡人的手触摸？但你不要苦恼；所发生的事是屈尊俯就。因为那如此微妙轻盈，甚至当门关着时也能进入的身体，是毫无密度的；但这个奇迹被显现出来，是为了使人相信复活，并让人们知道是祂自己，被钉十字架的那位，而不是另一个人代替祂复活了。为此，祂带着十字架的标记复活，也为此祂吃东西。至少宗徒们到处以此为复活的标记，说：\BibleQuote{我们这些与祂一同吃喝的人。}\BibleRef{宗 10:41} 因此，正如当我们在受难前看见祂在水面上行走时，我们不说那身体有另一种性质，而是我们自己的；同样，复活后，当我们看见祂有钉痕时，我们也不会因此说祂是可朽坏的。因为祂为了门徒的缘故而显出这些迹象。\par

% structure: p0014 source=h2 target=subsection reason=relative-heading-rank; base=h2

\subsection{若 20:30}

2. 既然这位福音作者比其他福音作者提到得更少，他告诉我们，其他人也没有都提到所有的事，而是只提到足以使听者信服的事。因为，经上记载：\Quote{如果}它说，\BibleQuote{他们每一样都写下来，我想连世界本身也容不下那些书。}\BibleRef{若 21:25} 由此清楚可见，他们所写下的并非为炫耀，而只是为了有益的事。因为他们既然略过了大部分，又为何写下这些以炫耀呢？但为什么他们没有全部写完？主要是因为数量众多；此外，他们也认为，那不相信他们已经写下的，也不会因更多数目而相信；而接受这些的，就不必为相信而需要其他了。这里他似乎是暂时谈论复活后的奇迹。因此祂说：\par

\Quote{在祂门徒面前。}\par

因为正如在复活前需要行许多奇迹，好使他们相信祂是天主子，同样在复活后也需要，好使他们承认祂复活了。还有一个原因他加了\Quote{在祂门徒面前}，因为祂在复活后只与他们单独交谈；为此祂也说：\BibleQuote{世界不再看见我了。}\BibleRef{若 14:19} 然后，为叫你明白所行的一切只是为了门徒的缘故，他加上：\par

% structure: p0018 source=h2 target=subsection reason=relative-heading-rank; base=h2

\subsection{若 20:31}

这是对众人说话，并表明祂不是向被相信的那一位，而是向我们自己，赐予一份极大的恩惠。\Quote{因祂的名}，即\Quote{藉着祂}；因为祂就是生命。\par

% structure: p0020 source=h2 target=subsection reason=relative-heading-rank; base=h2

\subsection{若 21:1}

你看，祂不继续同他们在一起，也不像先前那样常与他们同在？例如，祂晚上显现，然后离去；八天之后再次显现，然后再次离去；此后在海边，又带着极大的惊惶。但\Quote{显示}是什么意思？由此可知，除非祂俯就，否则人看不见祂，因为祂的身体从此是不朽的，且是纯粹无杂质的。但为何作者要提到地点？为表明祂现已除去他们大部分的恐惧，使他们敢于走出居所，四处走动。因为他们不再关在家里，而是去了\Place{加里肋亚}，躲避来自犹太人的危险。\Person{西满}于是去打鱼。因为既然祂不常与他们同在，圣神尚未赐下，那时他们也没有受托任何事，无事可做，便重操旧业。\par

% structure: p0022 source=h2 target=subsection reason=relative-heading-rank; base=h2

\subsection{\BibleRef{若 21:2}}

他们既然无事可做，便去打鱼，且在夜间做这事，因为十分害怕。\Person{路加}也提到这事，但这不是同一场合，而是另一次。其他门徒也跟随，因为他们从此彼此相连，同时想看捕鱼，并善用空闲。当他们正劳苦疲惫时，\Person{耶稣}出现在他们面前，并不立即显示自己，好使他们与祂交谈。于是祂对他们说：\par

% structure: p0024 source=h2 target=subsection reason=relative-heading-rank; base=h2

\subsection{\BibleRef{若 21:5}}

有一阵子祂说话颇带人性，仿佛要向他们买些东西。但当他们示意没有时，祂命他们将网撒在右边；撒下后便得了许多鱼。当他们认出祂时，门徒\Person{伯多禄}和\Person{若望}再次表露了他们各自的性情特点。一个更热忱，一个更超拔；一个更敏锐，一个更明澈。因此\Person{若望}先认出\Person{耶稣}，伯多禄先到祂那里。因为所发生的事并非寻常征兆。是什么征兆？首先，捕到这么多鱼；其次，网没有破；然后，在他们上岸之前，已发现有炭火，上面有鱼和饼。因为祂不再从已有的物质中造物，正如受难前出于某种安排所做的那样。因此，当伯多禄认出祂时，便放下一切——鱼和网——束上外衣。你看到他的尊敬和爱吗？然而他们只相隔二百肘；但即使如此，伯多禄也不能等船到祂那里，而是游泳上岸。耶稣怎么办呢？\par

% structure: p0026 source=h2 target=subsection reason=relative-heading-rank; base=h2

\subsection{\BibleRef{若 21:12}}

因为他们不再有同样的胆量，也不那么自信，现在不再用言语接近祂，而是怀着沉默、极大的敬畏和尊敬，坐下注视祂。\par

\BibleQuote{因为他们知道是主。}\par

因此他们不问祂：\Quote{你是谁？}但见祂的容貌改变，充满极大的威严，便大为惊愕，想要问些什么；但恐惧，以及知道祂并非他人，而是同一位，便阻止了询问，他们只吃了祂用比先前更大的权能所创造的。因为在这里祂不再举目望天，也不做那些人性行为，显示祂先前所做的也是出于俯就。为表明祂不继续常与他们同在，也不像先前那样，经上说：\par

% structure: p0030 source=h2 target=subsection reason=relative-heading-rank; base=h2

\subsection{\BibleRef{若 21:14}}

祂吩咐他们\BibleQuote{把鱼拿些来}，为表明他们所见的并非幻象。但这里并没有说祂与他们同吃，而\Person{路加}在另一处说祂吃了；因为\BibleQuote{祂同他们一起进食。}\BibleRef{宗 1:4}但\Quote{如何}不是我们可说的；因为这些事发生得太奇特，并非因为祂的本性如今需要食物，而是出于俯就的行动，为证明复活。\par

3. 或许当你听到这些事时，你心中炽热，称那些当时与他同在的人为有福，也称那些将在普遍复活之日与他同在的人为有福。那么，让我们竭尽全力，好能得见那令人赞叹的圣容。因为如果我们现在仅凭听闻就如此火热，渴望曾身处他在世上的那些日子，听到他的声音，看见他的面容，并得以亲近、触摸、事奉他；那么，想象一下，不再以可朽的身体、不行人性之事，而是由天使护卫、我们自己也在纯洁无瑕的形态中，瞻仰他，并享受那无可言喻的福乐，这是何等的大事！因此，我恳求，让我们用尽一切方法，以免错失如此的光荣。因为只要我们愿意，就没有难事；只要我们留意，就没有重担。\BibleQuote{如果我们忍受，也必与他一同为王。}\BibleRef{弟后 2:12}那么，\BibleQuote{如果我们忍受}是什么意思？即如果我们承受磨难，承受迫害，如果我们行走在窄路上。因为窄路就其本性而言是艰辛的，但因着我们的意愿，以及对未来希望的期盼，它变得轻省。\BibleQuote{因为我们这暂时轻微的苦难，为我们成就极重无比、永远的荣耀；我们不是顾念所见的，乃是顾念所不见的。}\BibleRef{格后 4:17}因此，让我们把目光转向天上，常常思念\Quote{那些}事，并注视它们。因为如果我们常与他们同在，就不会被世上享乐的欲望所动摇，也不会觉得忍受世上的忧愁是难的；反而会嗤笑这些以及类似的事，只要我们一心向往那里，并注视那爱。何必说我们不会为眼前的苦难而忧愁？我们甚至将不再看见它们。强烈的欲望就是如此。例如，那些此刻不与我们同在，但因被爱而远离的人，我们每天都在心中描绘他们。因为爱的主权是强大的，它使灵魂远离一切其他事物，并锁链于所渴望的对象。如果我们这样爱基督，这里的一切都将如影、如像、如梦。我们也将说：\BibleQuote{谁能使我们与基督的爱隔绝？是患难吗？是困苦吗？}\BibleRef{若 8:35}他没有说：\Quote{金钱、财富或美貌}（这些是极卑微可鄙的），而是列举了那些看似令人痛苦的事：饥荒、迫害、死亡。他甚至唾弃这些，视若无物；但我们却为了金钱远离生命，自绝于光明。而且保禄确实认为\Quote{无论是死，是生，是现在的事，是将来的事，是别的受造之物}，都不能与对基督的爱相比；但我们若见到一小块金子，就被点燃，并践踏他的法律。如果这些话令人无法容忍，那么行为更是如此。因为可怕的是：我们听到时战栗，但行出来却不战栗；我们轻易发誓并背誓，掠夺，放高利贷，毫不关心节制，不坚持祈祷的规范，违反大部分诫命，并且为了金钱不顾自己的肢体。因为爱财富的人会向邻人施行万般恶事，也加害于自己。他容易对人发怒，辱骂人，称人为愚人，发誓并背誓，甚至不遵守旧法律的尺度。因为爱金子的人不会爱邻人；然而，我们为了天国，奉命甚至要爱仇敌。如果我们遵守旧诫命尚不能进入天国，除非我们的义德超过并超越它们；那么当我们连这些也违反时，我们还能得到什么宽恕呢？爱金钱的人，不仅不会爱仇敌，甚至会视朋友如仇敌。\par

4. 但为何我谈到朋友？贪爱钱财的人常常连本性本身也置之不理。这样的人不认亲属，不记得友谊，不尊敬年龄，没有朋友，反而对所有人都心怀恶意，尤其对自己更甚，不仅毁灭自己的灵魂，而且用万般的忧虑、劳苦和悲伤折磨自己。因为他会忍受远行、仇恨、危险、阴谋，任何事，只为了在家中拥有那万恶之根，积攒许多金子。那么，还有什么比这种病更可悲的呢？它缺乏任何奢侈或快乐（人常为了这些而犯罪），也缺乏荣誉或荣耀。因为贪爱钱财的人怀疑自己拥有万千，实际上有许多人指责他、嫉妒他、诽谤他、图谋他。被他冤枉的人因受虐待而恨他；那些尚未受害的人，因害怕可能受害，并同情那些受害的人，也表现出同样的敌意；而那些更强大、更有权势的人，因卑微者的事而刺痛和愤慨，同时也嫉妒他，成为他的敌人和恨他的人。何必谈到人呢？因为当一个人也使天主成为他的敌人时，那时他还有什么希望？什么安慰？什么慰藉？爱财富的人永远无法使用财富；他将是财富的奴隶和看守，而不是主人。因为总是渴望增加财富，他永远不会愿意花费；反而会克扣自己，比任何穷人都贫穷，因为他的欲望永无止境。然而，财富被创造不是要我们保存，而是要我们使用；如果我们要为别人埋藏它们，那么还有谁比我们更悲惨呢？我们四处奔走，渴望聚集所有人的财产，好把它们关在屋里，与公共使用隔绝？但还有一种病不亚于这个。有些人把钱埋在地里，有些人埋在自己的肚腹里，在享乐和醉酒中；他们不仅不义，还为自己增加了放纵的惩罚。有些人用他们的财产服侍寄生虫和奉承者，有些人用来掷骰子和嫖妓，有些人用于其他类似的消费，为自己开辟了千万条通往地狱的道路，却离开了那正确而蒙准许的、通往天堂的道路。然而，这条正路不仅有着更大的益处，而且有比我们提到的那些事物更大的快乐。因为给娼妓的人是可笑可耻的，会招来许多争吵和短暂的快乐；甚至不是短暂的，因为无论他给那些他所迷恋的女人什么，她们都不会感谢他；因为\Quote{外人的房屋是漏水的箱。}\BibleRef{箴 23:27}此外，那类人脸皮厚，\Person{撒罗满}将他们的爱情比作坟墓；只有当她们看到爱慕者被剥光一切时才停止。或者更确切地说，这样的女人即使在那时也不停止，反而更加打扮自己，在他跌倒时践踏他，引起许多嘲笑，对他造成如此大的伤害，甚至无法用言语描述。得救之人的快乐并非如此；因为那里没有竞争者，所有人都欢欣喜悦，无论是接受祝福的人还是旁观者。没有愤怒、沮丧、羞耻、耻辱围攻这样的人的灵魂，而是良心上极大的喜乐，对将来事物极大的希望；他的荣耀光明，他的尊荣伟大；最重要的是来自天主的恩宠与安全，没有一个悬崖，没有猜疑，而是无波的海港与平静。因此，考虑这一切，将快乐与快乐相比，让我们选择更好的，好使我们获得未来的美好事物，藉着我们的主\Person{耶稣基督}的恩宠与慈爱，愿光荣与权能归于他，直到永远。阿们。\par

\SourceRecord{Translated by Charles Marriott.；Nicene and Post-Nicene Fathers, First Series；Vol. 14.；Edited by Philip Schaff.；Buffalo, NY: Christian Literature Publishing Co.,；1889.；<http://www.newadvent.org/fathers/240187.htm>.}

% structure: p0001 source=h1 target=section reason=page-title

\section{若望福音释义第八十八篇}

% structure: p0002 source=h2 target=subsection reason=relative-heading-rank; base=h2

\subsection{若 21:15}

有许多事的确能使我们在天主面前有胆量，并显为光明和可赞许，但最能使我们从上获得恩宠的，是对邻人的关怀。因此基督要求伯多禄这样做。因为当他们用完餐后，耶稣对西满伯多禄说：\Quote{西满，若纳的儿子，你爱我比这些更深吗？他对祂说：主，是的，祢知道我爱祢。}\par

他对他说：\Quote{你牧放我的羊。}\par

为什么祂绕过了其他的人，而和伯多禄谈论这些事？因为他是宗徒中的被选者，门徒们口中的发言人，团体的首领；为此，保禄也曾一度上耶路撒冷求问他，而非求问别人。同时，为了向他表明，既然否认的事已经过去，他现在应该放心了，耶稣就把弟兄中的首席权交在他手里；祂没有提出否认的事，也没有责怪所发生的事，却说：\Quote{你如果爱我，就管理你的弟兄们；你从前始终表现的那种热烈的爱，以及你曾以之为乐的爱，现在展示出来罢；你曾说愿意为我舍命，现在为我的羊群舍命罢。}\par

当祂问了一次、两次之后，伯多禄便请祂作察知人心隐秘的证人，然后第三次被问时，他感到不安，害怕重蹈覆辙（因为当时他自信断言，后来却被证实有罪），因此他又转向祂。因为这句话，\par

% structure: p0007 source=h2 target=subsection reason=relative-heading-rank; base=h2

\subsection{若 21:17}

\Quote{你知道一切}，意思是\Quote{现在的事和将来的事}。你看到他变得多么好、多么清醒了，不再固执己见或顶撞吗？因此他感到不安，\Quote{恐怕我认为我爱，而其实并不爱，就像从前我自以为爱并断言时，最终被证实有罪一样。}但耶稣第三次问他，第三次给他同样的命令，以表明祂多么重视对自己羊群的照顾，而这正是爱祂的标记。在对他说了关于爱祂的话之后，祂预言他将要受的殉道，表明祂对他说这些话并非出于不信任，而是大大信任他；此外，祂还希望指出爱祂的证明，并教导我们应当如何特别地爱祂。因此祂说：\par

\SourceRecord{Translated by Charles Marriott.；Nicene and Post-Nicene Fathers, First Series；Vol. 14.；Edited by Philip Schaff.；Buffalo, NY: Christian Literature Publishing Co.,；1889.；<http://www.newadvent.org/fathers/240188.htm>.}
